% SPDX-FileCopyrightText: 2026 Will Cook
% SPDX-License-Identifier: CC-BY-4.0
%
% Full editorial replacement, 18 September 2026. Numbered results and
% pinned formal-source destinations are retained; prose and hypotheses are
% aligned with the cited statements where this revision corrects a mismatch. Apply the prose changes to
% the authored parts before regenerating this file.
% Originally GENERATED by scripts/assemble_reasoning_surfaces.py from the authored files
% under paper/reasoning-parts/erdos249. Parts are inlined, never \input, so
% that the Pandoc projection cannot silently drop a section.
\documentclass[11pt]{article}
% SPDX-FileCopyrightText: 2026 Will Cook
% SPDX-License-Identifier: CC-BY-4.0
%
% Shared preamble for the two per-problem reasoning-surface papers.
% These documents are written to be read in full as working context: every
% result appears with its exact hypotheses, its evidence band, the scale at
% which it is proved, and the Lean declaration that carries it.
\usepackage{paper-house-style}
\usepackage{array}
\usepackage{adjustbox}
\usepackage{booktabs}
\usepackage{longtable}
\usepackage{enumitem}
\setlist{topsep=0.35em,itemsep=0.2em}
\usepackage{xurl}
\usepackage{xspace}
\usepackage[colorlinks=true,linkcolor=linkinternal,urlcolor=linkexternal,
  citecolor=linkinternal]{hyperref}
\hypersetup{bookmarksnumbered=true,bookmarksopen=true,bookmarksopenlevel=1,
  pdftitle={The Binary Totient Series: results and remaining arithmetic conditions for Erdos Problem 249},
  pdfauthor={Will Cook},pdflang={en-GB}}
% Lean declaration names are full of underscores and are typeset in text mode.
% Category codes are fixed when TeX READS a line, not when a macro expands, so
% no amount of \catcode juggling inside \lean can rescue an underscore that has
% already been tokenised as a subscript. This package makes _ print as itself in
% text while leaving it as a subscript in maths, which is exactly the split we
% need. It also protects any stray underscore in authored prose.
\usepackage{underscore}

% ===========================================================================
% Theorem environments.
% ===========================================================================
\theoremstyle{plain}
\newtheorem{thm}{Theorem}[section]
\newtheorem{lem}[thm]{Lemma}
\newtheorem{prop}[thm]{Proposition}
\newtheorem{cor}[thm]{Corollary}
\theoremstyle{definition}
\newtheorem{defn}[thm]{Definition}
\theoremstyle{remark}
\newtheorem{rem}[thm]{Remark}
\newtheorem{obs}[thm]{Observation}

% ===========================================================================
% Citation of a Lean declaration. The declaration name is the authority; the
% file:line is where it sat at the pinned formal-source checkpoint. Underscores
% are handled locally so authors may write names verbatim.
% ===========================================================================
\newcommand{\PK}{Erdos249257}
\newcommand{\commit}{99f4bf47422abbd8757cbb22b50ba079d764d3a7}
\newcommand{\repobase}{https://github.com/wcook04/plectis-erdos/blob/\commit}
% Sources that postdate \commit are cited with their lean/ path and resolve at \latecommit,
% whose tree keeps the Lean roots under lean/.
\newcommand{\latecommit}{e4ba6841dfdbf3a025c2a3a8c5cb98bcc0457ce8}
\newcommand{\laterepobase}{https://github.com/wcook04/plectis-erdos/blob/\latecommit}
% \lean declaration file:line forms a hyperlink to the pinned public
% formal-source checkpoint. The second argument is a machine coordinate,
% never typeset: \allowbreak and spaces are stripped, an ErdosProblems/
% prefix links the ErdosProblems tree directly (otherwise the \PK tree is
% used), a trailing :N becomes a GitHub #LN anchor, :N-M and :N--M become
% range anchors #LN-LM, and comma-separated lists link the first line.
\ExplSyntaxOn
\tl_new:N \l__leanlink_arg_tl
\tl_new:N \l__leanlink_url_tl
\seq_new:N \l__leanlink_split_seq
\cs_new_protected:Npn \leanlink #1
  {
    \tl_set:Nx \l__leanlink_arg_tl { \tl_to_str:n {#1} }
    \regex_replace_all:nnN { \\allowbreak } { } \l__leanlink_arg_tl
    \regex_replace_all:nnN { \s+ } { } \l__leanlink_arg_tl
    \exp_args:NnV \regex_match:nnTF { \A lean/ } \l__leanlink_arg_tl
      { \tl_set:Nx \l__leanlink_url_tl { \laterepobase } }
      {
        \tl_set:Nx \l__leanlink_url_tl { \repobase }
        \exp_args:NnV \regex_match:nnTF { \A (ErdosProblems|Erdos249257)/ } \l__leanlink_arg_tl
          { }
          { \tl_put_right:Nx \l__leanlink_url_tl { / \PK } }
      }
    \exp_args:NnV \regex_extract_once:nnNTF { (.*\.lean) :? (.*) } \l__leanlink_arg_tl \l__leanlink_split_seq
      {
        \tl_put_right:Nx \l__leanlink_url_tl { / \seq_item:Nn \l__leanlink_split_seq { 2 } }
        \tl_set:Nx \l_tmpa_tl { \seq_item:Nn \l__leanlink_split_seq { 3 } }
        \regex_replace_all:nnN { -- } { - } \l_tmpa_tl
        \exp_args:NnV \regex_extract_once:nnNTF { ([0-9]+) (- ([0-9]+))? } \l_tmpa_tl \l__leanlink_split_seq
          {
            \tl_put_right:Nx \l__leanlink_url_tl { \c_hash_str L \seq_item:Nn \l__leanlink_split_seq { 2 } }
            \str_if_eq:eeF { \seq_item:Nn \l__leanlink_split_seq { 4 } } { }
              { \tl_put_right:Nx \l__leanlink_url_tl { -L \seq_item:Nn \l__leanlink_split_seq { 4 } } }
          }
          { }
      }
      { \tl_put_right:Nx \l__leanlink_url_tl { / \tl_use:N \l__leanlink_arg_tl } }
    \href { \tl_use:N \l__leanlink_url_tl } { \textnormal{\textsc{Lean~source}} }
  }
\newcommand{\lean}[2]{\leanlink{#2}}
\ExplSyntaxOff

% ===========================================================================
% Statement-level proof links.  The declarations cited under the statements
% resolve at their own immutable pin, \ledgercommit.  The note printed under
% each statement is generated.  Do not edit a \leannote by hand.
% ===========================================================================
\newcommand{\ledgercommit}{181078b6b009d905809cf2e007ad309386a3d1e0}
\newcommand{\ledgerbase}{https://github.com/wcook04/plectis-erdos/blob/\ledgercommit}
\newcommand{\lproof}[3]{\href{\ledgerbase/lean/#1\#L#2}{\texttt{#3}}}
\newcommand{\leannote}[1]{%
  \par\nobreak\vspace{-0.3\baselineskip}%
  \noindent{\small #1\par}\vspace{0.15\baselineskip}}

% Declaration names are written \mathtt{...} in maths and, unavoidably, in
% running prose too. Rather than police every occurrence, make the command
% mode-agnostic: monospace is what is meant in both places.
\AtBeginDocument{%
  \let\origmathtt\mathtt
  \renewcommand{\mathtt}[1]{\ifmmode\origmathtt{#1}\else\texttt{#1}\fi}%
  \let\origmathrm\mathrm
  \renewcommand{\mathrm}[1]{\ifmmode\origmathrm{#1}\else\textrm{#1}\fi}%
}

% One authored part uses \m{...} as an inline-maths shorthand. Defining it is
% cheaper and safer than rewriting the prose around every occurrence.
\providecommand{\m}[1]{\ensuremath{#1}}

% ===========================================================================
% Evidence band. These are never blurred: a reader must always be able to see
% what a kernel checked, what an exact computation checked, what was proved in
% ordinary mathematics, what is quoted from the literature, and what is open.
% ===========================================================================
\newcommand{\ev}[1]{{\small\textsf{[#1]}}\xspace}

% Scale at which a statement is proved. This is the load-bearing axis of both
% documents: the corpus proves at fixed or bounded scale, and both open
% producers need cofinal scale.
\newcommand{\scale}[1]{{\small\textsf{[#1]}}\xspace}

% The representation a statement is expressed in. Obstructions are relative to
% a coordinate, never to the object, so every obstruction carries one.
% \nolinkurl manipulates catcodes, so it cannot be written into the
% hyperref bookmark file. \coord appears in subsubsection titles, so it
% must supply a plain-text PDF string alongside the typeset form.
\newcommand{\coord}[1]{} % Source classification retained, not printed.
\newcommand{\codeid}[1]{\nolinkurl{#1}}

% ===========================================================================
% Mathematical shorthand.
% ===========================================================================
\newcommand{\N}{\mathbb{N}}
\newcommand{\Npos}{\mathbb{N}_{>0}}
\newcommand{\Z}{\mathbb{Z}}
\newcommand{\Q}{\mathbb{Q}}
\newcommand{\R}{\mathbb{R}}
\newcommand{\C}{\mathbb{C}}
\newcommand{\ph}{\varphi}
\newcommand{\EB}{E}

\allowdisplaybreaks
\emergencystretch=2em

% These working records contain a small number of intentionally exhaustive
% formulas. Keep ordinary displays at their natural size, but bound any display
% wider than the classical text measure rather than letting it disappear past
% the page edge.
\long\def\[#1\]{%
  \leavevmode\par\vspace{\abovedisplayskip}%
  \noindent\makebox[\linewidth][c]{%
    \begin{adjustbox}{max width=\linewidth}%
      \(\displaystyle #1\)%
    \end{adjustbox}}%
  \par\vspace{\belowdisplayskip}%
}

\runninghead{Erd\H{o}s \#249: the binary totient series}
\title{The Binary Totient Series}
\subtitle{Results, proof dependencies and remaining arithmetic conditions\newline for Erd\H{o}s Problem~\#249}
\author{Will Cook}
\date{July 2026; revised 18 September 2026}

\begin{document}
% ---- part a249_front ----
\maketitle

\paperprovenance

\begin{abstract}
The dyadic sections of Euler's totient through level $e\ge1$ have an
explicit rational basis of size $2^e+1$. This record gives the detailed
statements behind that result and examines what coefficient structure can
say about $S=\sum_{n\ge1}\varphi(n)2^{-n}$. Truncating a scaled tail
difference produces an integer whose residue certifies nonintegrality
when its distances from both endpoints of the residue interval exceed
an explicit error bound. We develop the resulting finite
denominator exclusions, equivalent residue criteria and sufficient
exponential-sum estimates. Rational comparison sequences identify the
limitations of particular size, parity, rank and finite-prefix arguments.
The short paper supplies the all-base basis theorem and its integral
relations. Here each conditional result is accompanied by the arithmetic
hypothesis still needed to apply it; these hypotheses are not established
in the form required to prove irrationality of $S$.

\end{abstract}

\section*{Reading guide}
The short paper gives a direct proof of the basis theorem and its integral
normal form. This longer record is organised for reference. The opening
section introduces the finite residue test and its quantified form,
Definition~\ref{defn:sep}. Section~\ref{sec:wall} gives rational comparison
sequences, and Section~\ref{sec:survivors} states four possible arithmetic
conditions for irrationality. The detailed formulas then begin with the
series identities and denominator bound in Section~\ref{ssec:farey}.
The basis and relation modules are revisited in Section~\ref{sec:mahler-defect}.

Two distinctions recur. ``Cofinal'' means that a condition holds at some
index beyond each threshold; it does not mean that it holds at every
sufficiently large index. A counterexample excludes only the hypotheses
it satisfies. Agreement with a finite totient prefix, for example, need
not preserve multiplicativity.

The annotations identify the kind of support: \ev{Lean} refers to a
formal statement and its pinned source or supplied checking record;
\ev{Cert} to a finite computation; \ev{Math} to an ordinary argument;
\ev{Cited} to a literature result; and \ev{Open} to an unproved assertion.
The Lean kernel was not rerun for this revision. In particular, a checked
implication does not establish its hypothesis.

\setcounter{tocdepth}{1}
\clearpage
\tableofcontents
\clearpage
% ; - part a249_front ; -
\section{Totient sections and tail differences}

Let $\varphi$ be Euler's totient function, with $\varphi(0)=0$, and use
$\N=\{0,1,2,\ldots\}$. The dyadic sections
$n\mapsto\varphi(2^jn+r)$, for $0\le j\le e$ and $0\le r<2^j$,
have rational rank $2^e+1$ when $e\ge1$. A basis consists of
$n\mapsto\varphi(n)$, $n\mapsto\varphi(2n)$, and the sections with odd
$r$ at levels $1\le j\le e$. For example,
$\varphi(4n)=2\varphi(2n)$ and $\varphi(4n+2)=\varphi(2n+1)$ recover the
omitted sections at level two. The short paper proves the corresponding
rank $k^e+1$ in every integer base $k\ge2$, together with the integral
coordinates and relations. For composite $k$, the retained condition is
$k\nmid r$, not $\gcd(k,r)=1$.

The question about the series concerns more than these coefficient
relations. Put
\[
  S \;:=\; \sum_{n \ge 1} \frac{\varphi(n)}{2^{n}}
  \;=\; \tfrac12 + \tfrac14 + \tfrac{2}{8} + \tfrac{2}{16} + \tfrac{4}{32} + \cdots
  \;=\; 1.3676308019850223\ldots
\]
The series converges absolutely because $\varphi(n) \le n$. Erd\H{o}s Problem
\#249 asks whether $S$ is irrational, which remains open. The exact rank
theorem controls linear relations among coefficient sections. Irrationality
requires information about their weighted infinite sum. The countermodels
below explain several ways that a passage from coefficient structure to the
sum can fail. Each excludes its stated assumptions; they do not establish a
single obstruction theorem covering all possible methods.

\begin{rem}[Earlier coefficient results]
Erd\H{o}s and Graham record the irrationality question for
$S$~\cite[p.~61]{erdosgraham1980}. An integer sequence is $k$-regular
when the $\Z$-module generated by its sections $n\mapsto c(k^jn+r)$,
$j\ge0$ and $0\le r<k^j$, is finitely generated. Coons proved that
$\varphi$ is not $k$-regular for any $k\ge2$~\cite[Theorem~3.2]{coons}.
Martin's
Theorem~1 implies that the functions $n\mapsto\varphi(a_in+b_i)$ are linearly
independent for pairwise nonproportional affine forms with positive integer
slopes~\cite[Theorem~1]{martin-phi-inequalities}. The exact ranks in
\S\ref{sec:mahler-defect}, together with the all-base normal form in the
short paper, identify the retained sections and every relation among them.
The independence itself is not claimed as new. The CRT--Dirichlet separation
method is already present in Yazdani's Theorem~2 and proof, which explicitly
credit Shallit~\cite[pp.~652--653]{yazdani2001}; our evaluation matrix uses
one such separation for each column. Yazdani also proves non-automaticity of
$\varphi(n)\bmod m$ for every $m\ge3$ in every base~\cite[Cor.~4]{yazdani2001}.
Allouche, Shallit and Yassawi give a survey account~\cite[Thm.~3,
Ex.~4 and Rem.~7]{allouche-shallit-yassawi}.

\end{rem}

The following identity relates the series to finite sums. It also explains
which finite inequalities would suffice for irrationality. Write $2^{N} S = \Phi_N + R_N$
with $\Phi_N := \sum_{n=1}^{N} \varphi(n) 2^{N-n} \in \N$ and
$R_N := \sum_{j \ge 0} \varphi(N+1+j)/2^{j+1} \in [0,\infty)$
(\lean{Erdos249257.TotientTailPeriodKiller.totientTail}{TotientTailPeriodKiller.lean:57},
\lean{two_pow_mul_totient_series_eq}{TotientTailPeriodKiller.lean:150}).
If $S=a/(2^cv)$ with $c\ge0$ and $v\ge1$ odd, Euler's theorem gives
$h=\varphi(v)>0$ and $N_0=c$ such that $R_{N+h}-R_N\in\Z$ for every
$N\ge N_0$
(\lean{Erdos249257.eventual\_period\_of\_not\_irrational}{TotientTailPeriodKiller.lean:358}).
Indeed,
\[
 R_{N+h}-R_N\equiv 2^N(2^h-1)S\pmod\Z,
\]
and the right side is integral for those choices.

Thus it suffices, for every positive shift $h$, to find arbitrarily large
$N$ for which $R_{N+h}-R_N$ is not an integer. A finite test is obtained by
truncating this difference. Throughout, $a\bmod m$ denotes the least
nonnegative residue, and an unqualified index in a sum is a nonnegative
integer. For $h,N,L\in\N$, write
\[
  {D}(h,N,L)
  \;:=\; \sum_{j<L}\bigl(\varphi(N{+}h{+}1{+}j) - \varphi(N{+}1{+}j)\bigr)\,2^{\,L-1-j} \in \Z,
\]
the depth-$L$ truncation of $2^{L}(R_{N+h} - R_N)$, and
\[
  \mathcal{C}(h,N,L)
  \;:\Longleftrightarrow\;
  (N{+}h{+}L{+}2) \;<\; {D}(h,N,L) \bmod 2^{L}
  \;<\; 2^{L} - (N{+}h{+}L{+}2),
\]
a decidable condition
(\lean{Erdos249257.TotientTailPeriodKiller.windowDiscrepancy}{TotientTailPeriodKiller.lean:66},
\lean{Erdos249257.TotientTailPeriodKiller.certifiedKill}{TotientTailPeriodKiller.lean:72}).

The error bound $N{+}h{+}L{+}2$ comes from
$\varphi(m) \le m$, so a certificate is sound: it forces
$R_{N+h} - R_N \notin \Z$, and in fact certificates are \emph{complete} for
non-integrality
(\lean{Erdos249257.exists\_certifiedKill\_iff\_tail\_diff\_notMem\_int}{LcmConeFlatness.lean:316}).
The resulting quantified condition is the following.

\begin{defn}[The quantified residue condition]
\label{defn:sep}
\[
  \mathrm{Sep}
  \;:\Longleftrightarrow\;
  \forall h \ge 1\ \forall N_0\ \exists N \ge N_0\ \exists L,\;
  \mathcal{C}(h,N,L).
\]
$\mathrm{Sep} \Longleftrightarrow S\notin\mathbb Q$
(\lean{irrational\_totient\_series\_iff\_certificate\_supply}{LcmConeFlatness.lean:412}).
The forward implication is also
\lean{Erdos249257.irrational\_totient\_series\_of\_certificate\_supply}{TotientTailPeriodKiller.lean:394}.
The equivalence identifies the missing quantified input; it supplies no
certificate by itself.
\coord{binary-digit} \scale{cofinal} \ev{Lean}
\end{defn}

The order of the quantifiers matters. A certificate at one $(h,N,L)$
excludes those hypothetical rational values whose denominator divides
$2^N(2^h-1)$; it does not exclude all denominators. The condition above asks
for certificates beyond every threshold for each positive $h$, and is
therefore exactly as strong as the original irrationality assertion.
Its depth must grow: the two strict residue inequalities imply
$2(N+h+L+2)<2^L$. Keeping $L$ fixed cannot provide certificates as $N$ grows.
Neither this necessary size condition nor a finite list of successful tests
supplies the required residues.
\lean{Erdos249257.certifiedKill\_depth\_floor}{TotientTailPeriodKiller.lean:79}


\subsection{Unconditional results}

The following results require no irrationality hypothesis. The finite
denominator bound concerns rational values of $S$; the rank and positivity
results concern its coefficients and tails. The latter do not by themselves
exclude rationality.

\begin{thm}[Denominator exclusion from a fixed Farey window]
\label{thm:denom}
If $S \in \Q$ then its reduced denominator exceeds
$Q_0 := 79\,639\,646\,646\,701\,375\,323\,355\,774\,875\,831\,053 \approx 7.96 \times 10^{34}$.
Equivalently, $S$ differs from every rational number whose reduced
denominator is at most $Q_0$.
The bound is sharp for this window: $q = Q_0 + 1$ is the exact first failing
denominator. It is the denominator of the mediant of two explicit
unimodular Farey neighbours.
\coord{farey} \scale{bounded} \ev{Lean}
\lean{Erdos249257.tsum\_totient\_div\_pow\_two\_ne\_ratCast\_of\_den\_le\_79639646646701375323355774875831053}{CertificateKernel.lean:18384}
\lean{Erdos249257.gap\_check\_window\_1\_240\_first\_failure}{GapFareyBound.lean:225}
\end{thm}

This is obtained from a classical Stern--Brocot gap lemma
(\lean{farey\_gap}{GapFareyBound.lean:51}) applied at window $K = 240$, the
largest of the recorded bounds
$4838$, $2^{22}$, approximately $2.49\times10^{17}$, and $Q_0$. It is
proved without assuming $\mathrm{Sep}$ or using the tail-residue tests.
To deduce irrationality by this argument, the analogous denominator bounds
would have to be unbounded as the window grows; this has not been proved.
\ev{Open}

\begin{prop}[Finite certificate computations]
\label{prop:deposits}
$\mathcal{C}$ has been verified at: the $28$ diagonal instances of
the least-common-multiple diagonal through $t = 64$
(\lean{Erdos249257.certifiedKill\_diagonal\_all\_imported\_through\_t64}{DiagonalPincerCertificatesT64.lean:1967},
endpoint \lean{Erdos249257.certifiedKill\_diagonal\_t64}{DiagonalPincerCertificatesT64.lean:1928});
all shifts $h \in [1,16]$ simultaneously at $(N,L) = (14,9)$, by
\texttt{decide}
(\lean{Erdos249257.certifiedKill\_all\_upto\_sixteen}{CarrySurvivorExtinction.lean:574});
and eight further period examples at $N = 300$. These are historical
subsets of the examples. The supplied source also proves a complete diagonal band for
$1\le t\le82$, with no omitted scales
(\lean{exists\_diagonalKill\_le\_82}{ErdosProblems/Skip/LadderT67.lean:71264}).
This is bounded coverage, not a cofinal diagonal family.
Note the quantifier order: the simultaneous small-shift example
uses one pair $N,L$ for every $1\le h\le16$. In symbols, it proves
$\exists N\,\exists L\,\forall h\,(1\le h\le16\Rightarrow\mathcal C(h,N,L))$.
In contrast, $\mathrm{Sep}$ requires arbitrarily large $N$ for each
positive $h$; neither the shift range nor the basepoint threshold is bounded.
\coord{binary-digit} \scale{fixed} \ev{Lean}
\end{prop}
\leannote{Lean: \lproof{Erdos249257/DiagonalPincerCertificatesT64.lean}{1967}{certifiedKill_diagonal_all_imported_through_t64}, \lproof{Erdos249257/DiagonalPincerCertificatesT64.lean}{1928}{certifiedKill_diagonal_t64}, \lproof{Erdos249257/CarrySurvivorExtinction.lean}{574}{certifiedKill_all_upto_sixteen}, \lproof{ErdosProblems/Erdos249/PeriodMultipleEscape.lean}{471}{certifiedKill_67_300}, and 8 further declarations in the \hyperref[sec:coverage]{coverage section}.}

\begin{prop}[Positivity does not exclude an integer]
\label{prop:sign}
Put $H_a=\operatorname{lcm}(1,\ldots,2^a)$. For every $a\ge8$ and
$J\ge0$ with $J+(a+6)<2\cdot2^a$,
\[
 R_{2H_a+J}-R_{H_a+J}>0.
\]
This needs no irrationality hypothesis
(\lean{Erdos257PeriodNoncollapse.actualLcmTailDiff\_shift\_pos}{TotientActualLcmOrbitSign.lean:39},
\lean{Erdos257PeriodNoncollapse.actualLcmTailOrbit\_pos}{TotientActualLcmOrbitSign.lean:144}).
If this difference is an integer and $K\ge0$ satisfies
\[
 J+K+(a+6)<2\cdot2^a,\qquad 2H_a+J+K+2<2^K,
\]
then
\[
 D(H_a,H_a+J,K)\bmod2^K
   =2^K-\bigl(R_{2H_a+J+K}-R_{H_a+J+K}\bigr)
\]
lies strictly between $2^K-(2H_a+J+K+2)$ and $2^K$
(\lean{Erdos257PeriodNoncollapse.actualLcm\_integral\_forces\_topEdgeResidue}{TotientActualLcmOrbitSign.lean:211}). The later tail difference is a positive integer.
Its negative, not the carry itself, is the representative near zero.
Thus positivity locates the residue near the upper endpoint; it does not
establish the central-residue inequalities. Both restrictions on $K$ are
part of this conclusion.
\coord{LCM-diagonal} \scale{cofinal} \ev{Lean}
\end{prop}
\leannote{Lean: \lproof{Erdos249257/TotientActualLcmOrbitSign.lean}{39}{actualLcmTailDiff_shift_pos}, \lproof{Erdos249257/TotientActualLcmOrbitSign.lean}{211}{actualLcm_integral_forces_topEdgeResidue}, \lproof{Erdos249257/CarrySurvivorExtinction.lean}{393}{carryOrbit_eq_tail_diff}.}

For the next two propositions, a rational value $S=a/v$, with integers
$a$ and $v\ge1$, would give the integer carry
\[
 u_N=vR_N=2^Na-v\Phi_N,\qquad
 u_{N+1}=2u_N-v\varphi(N+1),\qquad 0\le u_N\le v(N+2).
\]
The formal sources call this a tempered carry. Here $v$ is a positive
multiplier, not necessarily odd. Its dyadic sections are the sequences
$n\mapsto u_{2^jn+r}$.

\begin{prop}[Rationality forces unbounded carry rank]
\label{prop:rank}
If $S$ is rational then, for every $e$, the carry sections
$n\mapsto u_{2^jn+r}$ with $1\le j\le e$ and $0\le r<2^j$ span a rational
vector space of dimension at least $2^{e}-1$
(\lean{Erdos249257.not\_irrational\_totientSeries\_implies\_unbounded\_carryRank\_unconditional}{TotientCarryKernelRigidity.lean:300}).
The lower bound holds at every depth. It comes from the linear independence
of the $2^e+1$ retained dyadic totient sections for $e\ge1$, proved using
the Chinese remainder theorem and Dirichlet's theorem
(\lean{Erdos249257.linearIndependent\_canonicalTotientKernelFamily}{TotientMahlerDefect.lean:935}),
so the full family spans an infinite-dimensional space
(\lean{Erdos249257.not\_finiteDimensional\_span\_fullTotientKernel}{TotientMahlerDefect.lean:1145}),
the case $k=2$ of Coons's non-regularity theorem (\S\ref{sec:mahler-defect}).
\coord{carry-rank} \scale{uniform} \ev{Lean}
\end{prop}

\begin{prop}[Periodicity modulo an integer does not bound rational rank]
\label{prop:period-not-rank}
If $S$ is rational, the same carry $u$ has unbounded rational section rank
and one eventual period modulo $v$ valid for all its dyadic sections.
More precisely, there are $h\ge1$ and $N_0$ such that, for every $j,r\ge0$
and $n\ge N_0$,
\[
 u_{2^j(n+h)+r}\equiv u_{2^jn+r}\pmod v.
\]
The rank and periodicity assertions hold together
(\lean{Erdos257PeriodNoncollapse.not\_irrational\_totientSeries\_implies\_mod\_period\_and\_unbounded\_rank}{TotientTailCarryPeriod.lean:224}).
This conditional theorem alone is not a counterexample to a general
periodicity-to-rank implication: its antecedent is not established.
The separate $5/4$ control supplies a concrete counterexample to the generic
rationality-driven rank ceiling; see
 the detailed comparison in Section~\ref{sec:mahler-defect}.
\coord{carry-rank} \scale{uniform} \ev{Lean}
\end{prop}

\begin{prop}[Equivalent certificate conditions]
\label{prop:iffs}
Put $H_t=\operatorname{lcm}(1,\ldots,t)$. Several variations of
Definition~\ref{defn:sep} are equivalent to $S\notin\mathbb Q$. One may allow a positive multiple of each prescribed
shift, still requiring certificates beyond every basepoint threshold
(Theorem~\ref{catalogue:cert:b2};
\lean{periodMultipleKillSupply\_iff\_irrational}{ErdosProblems/Erdos249/PeriodMultipleEscape.lean:432}).
Alternatively, one may use the original quantified condition
(\lean{irrational\_totient\_series\_iff\_certificate\_supply}{LcmConeFlatness.lean:412})
or restrict to $h=N=H_t$ at arbitrarily large $t$
(Theorem~\ref{catalogue:cert:b3};
\lean{irrational\_totient\_series\_iff\_lcm\_diagonal\_certificate\_supply}{LcmConeFlatness.lean:426}).
Replacing the symmetric residue test by
\[
 N+L+2\le D(h,N,L)\bmod2^L\le2^L-(N+h+L+2)
\]
also gives equivalent quantified conditions, including its restriction to
$h=N=H_t$ at arbitrarily large $t$
(\lean{Erdos257PeriodNoncollapse.CofinalDirectedLcmCertificateSupply}{TotientTailCarryPeriod.lean:866}).
Finally, the counted phase-separation condition in
Proposition~\ref{prop:b2}(b) is equivalent to irrationality
(\lean{Erdos249257.dtwWindowSeparatedPairs\_iff\_irrational\_totient\_series}{PivotAntiReconstruction.lean:1765}).
These equivalences change the form of the arithmetic question, not its
logical strength.
\coord{binary-digit} \scale{cofinal} \ev{Lean}
\end{prop}
\leannote{Lean: \lproof{ErdosProblems/Erdos249/PeriodMultipleEscape.lean}{432}{periodMultipleKillSupply_iff_irrational}, \lproof{Erdos249257/LcmConeFlatness.lean}{412}{irrational_totient_series_iff_certificate_supply}, \lproof{Erdos249257/LcmConeFlatness.lean}{426}{irrational_totient_series_iff_lcm_diagonal_certificate_supply}, \lproof{Erdos249257/TotientTailCarryPeriod.lean}{875}{irrational_totientSeries_iff_cofinalDirectedLcmCertificateSupply}, and 1 further declaration in the \hyperref[sec:coverage]{coverage section}.}

\begin{prop}[A rational sequence preserving size bounds, parity and aperiodicity]
\label{prop:parity}
There is $c : \N \to \N$ with $c(n) \le 6$ and $c(n) \le n$ for all $n$,
$c(n) \equiv \varphi(n) \pmod 2$ for \emph{every} $n$, and $c$ not eventually
periodic; indeed for every $N, G, K$ there are $K$ explicit carry pulses
beyond $N$, pairwise separated by more than $G$; and yet
$\sum_n c(n)/2^{n} = 3/2 \in \Q$
(\lean{Erdos257PeriodNoncollapse.exists\_totientParity\_arbitrarilyManySeparatedCarry\_rational\_countermodel}{TotientParityCoboundaryCountermodel.lean:637};
sum at \texttt{:359}, aperiodicity at \texttt{:487}, parity match at \texttt{:194};
\texttt{\#print axioms} clean).
\coord{coefficient-word} \scale{uniform} \ev{Lean}
\end{prop}
\leannote{Lean: \lproof{Erdos249257/TotientParityCoboundaryCountermodel.lean}{637}{exists_totientParity_arbitrarilyManySeparatedCarry_rational_countermodel}, \lproof{Erdos249257/TotientParityCoboundaryCountermodel.lean}{359}{tsum_parityCoboundaryWeight_eq_three_halves}.}

These propositions separate several logical obligations. Positivity leaves
an upper-endpoint condition; a large rational carry rank is not a contradiction;
and exact equivalences still require a quantified certificate condition. The stated size
bounds, parity and aperiodicity also hold for a rational control. Each
conclusion concerns its stated hypotheses, not all possible arithmetic
information about $\varphi$.

\section{Rational comparison sequences}
\label{sec:wall}

The countermodels distinguish information that determines a finite calculation
from information that forces nonintegrality at unbounded indices. A fixed
prefix, parity at every index, and a specified finite collection of linear
relations are different hypotheses. The useful conclusion in each case is the
precise information a rational control sequence can preserve.

\subsection{What finite information does not determine}

The comparison sequences below preserve different parts of the totient
data. Theorem~\ref{thm:gamma} preserves an arbitrary fixed prefix.
Proposition~\ref{prop:parity} preserves parity at every index, together with
size bounds and aperiodicity. Propositions~\ref{prop:b5} and~\ref{prop:b6}
address particular carry descriptions and linear constructions, not every
finite-state or finite-dimensional argument.

Two other limitations do not require a rational comparison sequence.
A congruence construction can increase the index so much that its tail
error exceeds the residue it produces (Proposition~\ref{prop:b4}). An exact
reformulation or a finite list of examples still needs the quantified
arithmetic assertion (Proposition~\ref{prop:b2} and Remark~\ref{prop:b3}).

The finite certificates in Proposition~\ref{prop:deposits} and the $K=240$
Farey calculation in Theorem~\ref{thm:denom} have bounded ranges.
Theorem~\ref{thm:gamma} excludes
an inference from one such fixed range uniformly over the whole coefficient
class. It does not exclude a proof using the arithmetic of $\varphi$ at
arbitrarily large horizons. Likewise, the cofinal positivity theorem does
not by itself separate the residue from the top edge: that is the additional
information identified by Proposition~\ref{prop:sign}.

\subsection{A rational sequence with any prescribed totient prefix}

A finite prefix leaves the remaining coefficients free. We construct a
rational continuation satisfying the same linear size bound, then show why
it cannot satisfy the quantified residue condition. First we record the
two tail facts needed for a general coefficient sequence. For
$c : \N \to \N$ with $c(n) \le n$, write
$T_c := \sum_{n \ge 1} c(n)/2^{n}$,
$R^{c}_{N} := \sum_{j \ge 0} c(N{+}1{+}j)/2^{\,j+1}$,
$D_c(h,N,L) := \sum_{j<L}(c(N{+}h{+}1{+}j) - c(N{+}1{+}j))2^{\,L-1-j}$,
and let $\mathcal C_c(h,N,L)$ and $\mathrm{Sep}_c$ be
$\mathcal{C}$ and $\mathrm{Sep}$ with $D_c$ in place of
${D}$. For $c = \varphi$ these are the objects of
Definition~\ref{defn:sep}.

\begin{lem}[A residue certificate excludes an integral tail difference]
\label{lem:gsound}
Let $c:\N\to\N$ satisfy $c(n)\le n$ for all $n$. For $h,N,L\in\N$,
$\mathcal C_c(h,N,L) \Rightarrow R^{c}_{N+h} - R^{c}_{N} \notin \Z$.
\coord{binary-digit} \scale{uniform} \ev{Lean}
\end{lem}
\leannote{Lean: \lproof{ErdosProblems/Erdos249/PaperCompleteR20/GenericTailCertificates.lean}{66}{certificate_sound}, \lproof{ErdosProblems/Erdos249/PaperCompleteR20/GenericTailCertificates.lean}{48}{scaled_difference}, \lproof{ErdosProblems/Erdos249/PaperCompleteR20/GenericTailCertificates.lean}{55}{truncation_error_bound}.}

\begin{proof}
Splitting both tails after $L$ terms gives
\[
 2^L(R^c_{N+h}-R^c_N)-D_c(h,N,L)
   =R^c_{N+h+L}-R^c_{N+L}.
\]
Both tails on the right are nonnegative, and
$R^c_M\le\sum_{j\ge1}(M+j)2^{-j}=M+2$. Their difference therefore has
absolute value at most $N+h+L+2$. The corresponding totient bound is
formalised at
\lean{Erdos249257.tail\_after\_le}{TotientTailPeriodKiller.lean:198} and used
there to prove
\lean{Erdos249257.tail\_diff\_notMem\_int\_of\_certifiedKill}{TotientTailPeriodKiller.lean:262}.
If $R^{c}_{N+h} - R^{c}_{N} = k \in \Z$ then $2^{L}k \equiv 0 \pmod{2^{L}}$
and $D_c$ lies within $N{+}h{+}L{+}2$ of $2^{L}k$, so
$D_c \bmod 2^{L}$ lies within $N{+}h{+}L{+}2$ of $0$ or of $2^{L}$,
contradicting $\mathcal C_c$.
\end{proof}

\begin{lem}[Generic tail-period law]
\label{lem:gperiod}
Let $c:\N\to\N$ satisfy $c(n)\le n$ for all $n$. If
$T_c = p/(2^{e}m)$ with $p\in\Z$, $e\ge0$ and $m$ a positive odd integer,
and if $h\ge1$ satisfies
$m\mid 2^h-1$, then
$R^{c}_{N+h} - R^{c}_{N} \in \Z$ for every $N \ge e$.
\coord{binary-digit} \scale{cofinal} \ev{Lean}
\end{lem}
\leannote{Lean: \lproof{ErdosProblems/Erdos249/PaperCompleteR20/GenericTailCertificates.lean}{98}{generic_tail_period}, \lproof{ErdosProblems/Erdos249/PaperCompleteR20/GenericTailCertificates.lean}{33}{scaled_tail_split}.}

\begin{proof}
Since $2^NT_c-R^c_N=\sum_{n=1}^N c(n)2^{N-n}$ is an integer,
\[
 R^c_{N+h}-R^c_N\equiv2^N(2^h-1)T_c\pmod\Z.
\]
The right side is integral because $2^e\mid2^N$ and $m\mid2^h-1$.
\end{proof}

\begin{rem}
Lemmas~\ref{lem:gsound} and~\ref{lem:gperiod} are also checked for every
coefficient sequence satisfying $c(n)\le n$. The formal proofs use the
true-tail recurrence and the bound $0\le R^c_N\le N+2$; the residue
certificate keeps the full finite discrepancy. The tail-period proof uses
$m\mid2^h-1$ directly, so oddness of $m$ is only needed when choosing a
period from the denominator. The following rational continuation and its
finite-prefix consequence are also checked.
\lean{ErdosProblems.Erdos249.PaperCompleteR20.GenericTailCertificates.certificate\_sound}{lean/ErdosProblems/Erdos249/PaperCompleteR20/GenericTailCertificates.lean:66}
\lean{ErdosProblems.Erdos249.PaperCompleteR20.GenericTailCertificates.generic\_tail\_period}{lean/ErdosProblems/Erdos249/PaperCompleteR20/GenericTailCertificates.lean:98}
\end{rem}

\begin{thm}[A rational sequence agreeing with any finite totient prefix]
\label{thm:gamma}
Let $B\ge1$ and $P>B$ be integers. Define $\gamma : \N \to \N$ by
\[
  \gamma(n) := \varphi(n) \ \ (n \le B), \qquad
  \gamma(n) := \begin{cases} n-1, & P \mid n \\ n, & P \nmid n \end{cases}
  \ \ (n > B).
\]
Then:
\begin{enumerate}[label=(\roman*),leftmargin=*]
\item $\gamma(n) \le n$ for all $n$, so $\gamma$ lies in the same coefficient
  class as $\varphi$;
\item $\gamma(n) = \varphi(n)$ for every $n \le B$, and consequently
  $D_\gamma(h,N,L) = D_\varphi(h,N,L)$ for every $(h,N,L)$ with
  $N + h + L \le B$;
\item $T_\gamma = 2-\sum_{n\le B}(n-\varphi(n))/2^n-1/(2^P-1)\in\Q$,
  and the odd part of its reduced denominator is \emph{exactly} $2^P-1$;
\item $\mathrm{Sep}_\gamma$ is false; indeed it fails already at $h = P$.
\end{enumerate}
Hence for every $B$ there is a coefficient sequence in the same class,
agreeing with $\varphi$ on all of $[1,B]$, whose quantified condition is false.
\coord{binary-digit} \scale{uniform} \ev{Lean}
\end{thm}
\leannote{Lean: \lproof{ErdosProblems/Erdos249/PaperCompleteR20/FinitePrefixCountermodel.lean}{15}{gamma_le}, \lproof{ErdosProblems/Erdos249/PaperCompleteR20/FinitePrefixCountermodel.lean}{21}{gamma_prefix}, \lproof{ErdosProblems/Erdos249/PaperCompleteR20/FinitePrefixCountermodel.lean}{24}{discrepancy_prefix}, \lproof{ErdosProblems/Erdos249/PaperCompleteR20/FinitePrefixCountermodel.lean}{108}{exact_series}, and 4 further declarations in the \hyperref[sec:coverage]{coverage section}.}

\begin{proof}
(i) is immediate; (ii) holds because $D_\varphi(h,N,L)$ reads $\varphi$ only at
indices $\le N{+}h{+}L$. For (iii), since $P > B$ every multiple of $P$ exceeds
$B$, so
\[
  T_\gamma
  = \sum_{n\ge1}\frac{n}{2^{n}} - \sum_{n \le B}\frac{n - \varphi(n)}{2^{n}}
    - \sum_{k \ge 1} 2^{-kP}
  = 2 - \sum_{n \le B}\frac{n-\varphi(n)}{2^{n}} - \frac{1}{2^{P}-1}.
\]
The first two terms have denominator dividing $2^B$. On putting the
expression over $2^B(2^P-1)$, its numerator is congruent to $-2^B$ modulo
$2^P-1$. Since $\gcd(2^B,2^P-1)=1$, no odd factor cancels. Thus the
reduced denominator has odd part exactly $2^P-1$ and its power-of-two part
divides $2^B$. For (iv), apply Lemma~\ref{lem:gperiod} with $h=P$:
$R^{\gamma}_{N+P} - R^{\gamma}_{N} \in \Z$ for every $N \ge B$, so by
Lemma~\ref{lem:gsound} no $\mathcal C_\gamma(P,N,L)$ holds for any
$N \ge B$ and any $L$. Taking $N_0 = B$ refutes the inner existential of
$\mathrm{Sep}_\gamma$ at $h = P$.
The formalisation checks the series value, the exact reduced denominator,
and the failure of the separation condition for this same $\gamma$.
\lean{ErdosProblems.Erdos249.PaperCompleteR20.FinitePrefixCountermodel.exact\_series}{lean/ErdosProblems/Erdos249/PaperCompleteR20/FinitePrefixCountermodel.lean:108}
\lean{ErdosProblems.Erdos249.PaperCompleteR20.FinitePrefixCountermodel.exact\_denominator}{lean/ErdosProblems/Erdos249/PaperCompleteR20/FinitePrefixCountermodelEndpoint.lean:39}
\lean{ErdosProblems.Erdos249.PaperCompleteR20.FinitePrefixCountermodel.gamma\_not\_separation}{lean/ErdosProblems/Erdos249/PaperCompleteR20/FinitePrefixCountermodelEndpoint.lean:69}
\end{proof}

\begin{cor}[The limit of a finite-prefix argument]
\label{cor:b1}
No proof rule uniform over all $c:\N\to\N$ with $c(n)\le n$ can
establish $\mathrm{Sep}$ from a single fixed prefix
$\{c(n):n\le B\}$: Theorem~\ref{thm:gamma} supplies a rational countermodel
with that same prefix.  This does \emph{not} invalidate an argument that uses
the fixed arithmetic sequence $\varphi$ together with compatible information
at arbitrarily large horizons; the theorem gives a different $\gamma_B$ for
each $B$, not one sequence agreeing with $\varphi$ at every $B$.
\end{cor}
\leannote{Lean: \lproof{ErdosProblems/Erdos249/PaperCompleteR20/FinitePrefixCountermodelEndpoint.lean}{74}{no_uniform_prefix_rule}, \lproof{ErdosProblems/Erdos249/PaperCompleteR20/FinitePrefixCountermodelEndpoint.lean}{69}{gamma_not_separation}.}
\lean{ErdosProblems.Erdos249.PaperCompleteR20.FinitePrefixCountermodel.no\_uniform\_prefix\_rule}{lean/ErdosProblems/Erdos249/PaperCompleteR20/FinitePrefixCountermodelEndpoint.lean:74}

\begin{rem}[Four distinct finite restrictions]
\label{rem:b1-scope}
Theorem~\ref{thm:gamma} concerns the information in a finite prefix.
It does not identify $\varphi$ with one of the comparison sequences. It also
does not exclude a proof that combines finite computation with an arithmetic
theorem applying at all larger indices. Four different restrictions must
be kept separate: a bound on the indices inspected, a bound on the depth
$L$, a bound on the carry states retained, and a bound on matrix rank.
The depth inequality $2(N+h+L+2)<2^L$, for example, concerns the second
restriction, not all four.
\end{rem}

Proposition~\ref{prop:parity} gives another rational comparison sequence,
this time preserving parity and specified aperiodicity properties. The two
examples test different assumptions and must not be combined into a single
stronger example.


\subsection{Limits of specific reductions and estimates}

The following results concern specified assumptions, not all possible
proof methods. Statements labelled as an audit describe the finite set of
arguments examined in that audit. They are not mathematical impossibility
theorems.


\begin{prop}[Three particular equivalences]
\label{prop:b2}
\begin{enumerate}[label=(\alph*),leftmargin=*]
\item \emph{Certificate completeness.} The complete residue tests considered here are
  equivalent to the corresponding nonintegrality assertions. Rewriting the
  quantified condition using one of these equivalences does not weaken it
  (\lean{periodMultipleKillSupply\_iff\_irrational}{ErdosProblems/Erdos249/PeriodMultipleEscape.lean:432},
  forward at \texttt{:132}, converse at \texttt{:143}).
\item \emph{A selectable two-point sample.} The condition is as follows.
  For every $h\ge1$ and $X_0$, choose $X,L\in\N$ with
  $X\ge\max(X_0,1)$ and $16(2X+h+L+2)\le2^L$, a nonempty set
  $T\subseteq[X,2X)\cap\N$, a set $P\subseteq T\times T$ of ordered
  pairs, and a real number $\delta\ge0$. With
  $E_N=\exp(2\pi iD(h,N,L)/2^L)$, require
  \[
   |E_i-E_j|\ge\delta\quad((i,j)\in P),\qquad
   \frac{2|T|^2}{5}\le |P|\delta^2.
  \]
  This condition is equivalent to irrationality. For the converse,
  irrationality and the doubling map allow
  $T=\{N,N+1\}$ and $P=\{(N,N+1),(N+1,N)\}$ at arbitrarily large $N$,
  with $\delta=9/10$. The numerical requirement is then
  $8/5\le2(9/10)^2$. The freedom to select $T$ matters: this is not an
  estimate on a sample prescribed in advance, such as all prime positions.

\noindent Formal sources: \lean{Erdos249257.dtwWindowSeparatedPairs\_iff\_irrational\_totient\_series}{PivotAntiReconstruction.lean:1765}

\item \emph{Four integral tail differences.} For $H\ge0$ and positive
  integers $p,q$, the conjunction
  \[
   R_{2kH}-R_{kH}\in\Z\qquad(k\in\{1,p,q,pq\})
  \]
  is equivalent to $R_{2H}-R_H\in\Z$ alone. The affine transport identity
  for $H\mapsto kH$ preserves integrality, so the three additional
  conditions add no restriction. Neither $p$ nor $q$ need be prime
  (\lean{Erdos249257.fullTarget\_primeAdjunction\_diamond\_iff\_root}{FullTargetPrimeAdjunctionNoGo.lean:196}).
\end{enumerate}
These three equivalences explain why these particular reformulations
retain the original arithmetic question. They do not rule out useful weaker
intermediate lemmas. The full-block and fixed-margin conditions in
\S\ref{sec:survivors} ask for additional quantitative information that the
three arguments above do not supply.
\coord{binary-digit} \scale{cofinal} \ev{Lean}
\end{prop}
\leannote{Lean: \lproof{ErdosProblems/Erdos249/PaperCompleteR21/ThreeParticularEquivalences.lean}{33}{three_particular_equivalences}, \lproof{ErdosProblems/Erdos249/PaperCompleteR21/ThreeParticularEquivalences.lean}{55}{two_point_sample_numerical_requirement}.}

\begin{rem}
Second-difference certificates need not occur at a smaller depth than
first-difference certificates. At $(h,N)=(1,8)$, the first-difference test
holds at depth $8$, whereas no second-difference certificate exists at a
depth at most $8$ (certificate kernel source:18762, \ev{Lean}). This one
example refutes a uniform claim that second differences always save depth;
it does not give a uniform ordering in the opposite direction.

Separately, unbounded growth of the denominator exclusions produced by the
Farey windows would imply irrationality. The record does not establish a
converse for that specific sequence of windows, so the proposed growth
condition is used only as a sufficient condition. \ev{Math}
\end{rem}

\begin{rem}[A finite audit of proposed extensions]
\label{prop:b3}
This is a report of a finite audit. The audit listed
$56$ candidate deductions for four stated goals and examined $18$ proposed
extensions. Its recorded verdict was $16$ not accepted as routine
extensions and $2$ accepted only as restatements for a wider target family.
One of the $16$, labelled NM-12 in the audit, was not verified because the
endpoint proofs were not read. Thus it is not correct to say that all $18$
proof bodies were checked or that all $16$ negative verdicts were proved.
The recorded categories were hypothesis strength ($6$), parameter range
($5$), representation ($4$), and multiple issues ($3$). Neither accepted
extension supplied the arithmetic assertion needed for irrationality.
No conclusion about arguments outside this finite audit follows.
\coord{audit} \scale{n/a} \ev{Cert}
\end{rem}

\begin{obs}[What the finite audit observed]
\label{obs:pointwise}
Several unsuccessful arguments in the audit sought one certificate at
a specially chosen large index, whereas the exponential-sum criterion
seeks cancellation across an entire block. This is a difference in the
statements being attempted, not evidence that one approach is the only
possible one. The older list ending at $t=64$ is not the current finite
record: the supplied sources certify every diagonal scale $1\le t\le82$.
Neither list proves the required assertion for arbitrarily large scales.
\end{obs}

\begin{prop}[Failure of a specified two-adic congruence construction]
\label{prop:b4}
A certificate requires the residue to lie farther than $N+h+L+2$
from either endpoint modulo $2^L$. In the specific construction below,
a pulse places the residue at $2^{K-1}$ modulo $2^K$, but the defining
congruence is $p\equiv1+2^{K-1}\pmod{2^K}$. Hence
$p\ge1+2^{K-1}$. At $N=p-K$, $h=H$, and $L=K$, the error bound is
$p+H+2>2^{K-1}$, so the residue does not satisfy the certificate inequalities.
This calculation defeats this construction at every depth. It does not
prove a corresponding statement for every use of the Chinese Remainder
Theorem or every prescribed totient pattern.
\coord{two-adic-pulse} \scale{cofinal} \ev{Lean}
\lean{Erdos257PeriodNoncollapse.exists\_prime\_deltaTotient\_twoAdic\_pulseBlock}{TotientTwoAdicPulseBlock.lean:211}
\lean{Erdos257PeriodNoncollapse.windowDiscrepancy\_modEq\_half\_of\_twoAdic\_pulse}{TotientTwoAdicPulseBlock.lean:262}
\end{prop}
\leannote{Lean: \lproof{ErdosProblems/Erdos249/PaperCompleteR21/TwoAdicPulseCertificateFailure.lean}{29}{twoAdic_pulse_construction_never_certifies}, \lproof{ErdosProblems/Erdos249/PaperCompleteR21/TwoAdicPulseCertificateFailure.lean}{60}{twoAdic_pulse_defining_congruence}, \lproof{ErdosProblems/Erdos249/PaperCompleteR21/TwoAdicPulseCertificateFailure.lean}{85}{twoAdic_pulse_error_bound}.}

\begin{rem}[Two further prescribed patterns]
\label{rem:b4}
Two related results concern particular prescribed patterns. Under the
stated room condition, the complete terminal divisibility pattern is
impossible: its last totient difference would be positive, below a modulus,
and divisible by that modulus. For the pattern with the last condition
removed, the penultimate difference is forced to be $2^{m-1}$, with
$2^m<2(2H+J+K+2)$.

A different sufficient condition would prescribe the accumulated sum
rather than an individual difference:
${D}(h,N,L)\equiv2^{L-1}\pmod{2^L}$ and
$2^{L-1}>N+h+L+2$. These two inequalities would give a certificate.
They have not been supplied at the required quantified scope. The failure
of the earlier patterns does not prove that this replacement is necessary.
\ev{Math}

\noindent Formal sources: \lean{Erdos257PeriodNoncollapse.not\_actualLcmTerminalDyadicStaircase\_of\_room}{TotientActualLcmTopEdgeStaircase.lean:251}

\end{rem}

\begin{prop}[Information lost by specific carry descriptions]
\label{prop:b5}
\begin{enumerate}[label=(\roman*),leftmargin=*]
\item Fix $m\ge2$ and write $R=\lfloor(m+1)/2\rfloor$. For
  $r=0,\ldots,R$, let $c_r$ vanish except at
  $c_r(m)=R-r$ and $c_r(m+1)=2r$. Each sequence satisfies
  $0\le c_r(n)\le n$, has the same binary sum $R2^{-m}$, and has
  the same coefficients and scaled tails before position $m$.
  Its scaled tail at position $m$ is $r$. The common history therefore
  does not determine that tail. A finite set of labels that determines
  the tail for each member of this family must contain at least
  $R+1=\lfloor(m+1)/2\rfloor+1$ elements, one for each value of $r$
  (\lean{Erdos249257.balancedPulse\_no\_autonomous\_decoder}{GenericTailOrbitRigidity.lean:247}).
\item Reduction modulo $2^L$ forgets the initial value after the same
$L$ recurrence steps: two
  affine binary orbits with different seeds satisfy
  $\mathrm{orbit}_u(L) - \mathrm{orbit}_v(L) = 2^{L}(u_0 - v_0)$, so the
  endpoint residue mod $2^{L}$ is independent of the initial carry
  (\lean{Erdos249257.affineBinaryOrbit\_mod\_twoPow\_eq}{GenericTailOrbitRigidity.lean:307}).
\item Fix a precision $u\ge1$, a finite list of nonnegative valuations
  $\nu_j$ and odd integers $a_j$, and an initial integer $e_0$.
  There are integers $z_j$ such that
  \[
   c_j=2^{\nu_j}(a_j+2^u z_j),\qquad
   e_{j+1}=2e_j+c_j,\qquad |e_{j+1}|\le2^{\nu_j+u-1}.
  \]
  At each step, choose the centred representative of
  $2e_j+2^{\nu_j}a_j$ modulo $2^{\nu_j+u}$
  (\lean{Erdos249257.fixedPrecisionTropicalNoGo}{TropicalCurvatureCarry.lean:137}).
  The unrestricted integers $z_j$ are essential: the assertion does not
  prescribe the full coefficients $c_j$.
\end{enumerate}
These statements exclude decoders using only the specified common
history, distinctions using only the endpoint residue, and contradictions
using only the stated local symbols. They do not prove that the actual
totient expansion cannot be studied by automata. A proposed application
must verify that it uses no additional arithmetic information that the
comparison families fail to preserve.
\coord{carry-orbit} \scale{uniform} \ev{Lean}
\end{prop}
\leannote{Lean: \lproof{ErdosProblems/Erdos249/PaperCompleteR21/CarryDescriptionInformationLoss.lean}{93}{balancedPulse_common_history}, \lproof{ErdosProblems/Erdos249/PaperCompleteR21/CarryDescriptionInformationLoss.lean}{30}{balancedPulse_tail_at}, \lproof{ErdosProblems/Erdos249/PaperCompleteR21/CarryDescriptionInformationLoss.lean}{76}{balancedPulse_series}, \lproof{ErdosProblems/Erdos249/PaperCompleteR21/CarryDescriptionInformationLoss.lean}{114}{balancedPulse_label_lower_bound}, and 3 further declarations in the \hyperref[sec:coverage]{coverage section}.}

\begin{prop}[Four limits of particular linear constructions]
\label{prop:b6}
The following four constructions have different limitations. They must
not be read as a claim that every finite family of totient sections is
independent: the full family has the explicit relations described in the
short paper.
\begin{enumerate}[label=(\roman*),leftmargin=*]
\item \emph{Dyadic sections and an integer identity.} The retained
  family of $2^e+1$ dyadic sections is linearly independent over $\Q$
  for $e\ge1$ (Proposition~\ref{prop:rank}). Separately, positive integers
  $Q,v$ and integers $A,b$ cannot satisfy
  \[
     A\ne0,\qquad QvA=b,\qquad |b|<Qv,
  \]
  because $|A|\ge1$ gives $|b|\ge Qv$
  (\lean{Erdos249257.false\_of\_compressedAdjointCertificate}{TotientMahlerDefect.lean:1183}).
  This elementary incompatibility concerns the displayed requirements;
  it does not exclude every argument using an adjugate matrix.
\item \emph{M\"obius incidence.} $U_N(i,j) = \mu((i{+}1)/(j{+}1))$ when
  $(j{+}1) \mid (i{+}1)$ and $0$ otherwise is lower triangular with unit
  diagonal, so $\det U_N = 1$ for every $N$. Hence multiplication by
  $U_N$ is injective; the corresponding coefficient transformation cannot
  create a nonzero vector in its kernel
  (\lean{Erdos249257.incidenceMobius\_det\_eq\_one}{IncidenceQuotientHermitePade.lean:56},
  injectivity of the companion map at \texttt{:77}).
\item \emph{Adjugate reconstruction.} Let $w_i\in\Q$ and
  $x_i\in\N$ be finite families with $\sum_i w_i\varphi(x_i)=1$.
  Reconstructing each value from two tails and bounding them separately
  gives the error bound $\sum_i|w_i|(3x_i+4)$. It is at least $3$, since
  \[
   1\le\sum_i|w_i|\varphi(x_i)\le\sum_i|w_i|x_i.
  \]
  Consequently this particular bound cannot be less than $1$, whatever
  the size of the evaluation matrix
  (\lean{Erdos257PeriodNoncollapse.three\_le\_totientAdjugateTailCost}{TotientTailCarryPeriod.lean:266},
  closure at \texttt{:307}).
\item \emph{Finite combinations of shifts.} The synthetic sequence in
  Observation~\ref{prop:B4b-kill} has the prescribed differences
  $a_{(q-1)H-1}=\varphi(H)$ for $2\le q<t$, where
  $H=\operatorname{lcm}(1,\ldots,t)$. It is of the form
  $a_i=2c_i-c_{i+1}$, and every finite integer combination of its shifts
  has the same form with a correspondingly shifted state. Thus these
  linear operations alone do not remove the compatible carry recurrence.
  Their uniform bounds depend on the absolute coefficient sum, as made
  explicit in that observation. The construction does not assert that
  $a_i$ equals the actual totient difference at other indices.
\end{enumerate}
A further limitation concerns quotients of the finite sums
$t(Y,r)=\sum_{d=1}^{Y}\mu(d)/(2^d-1)^r$. For $e\ge1$ and $Y\ge4$,
\[
 \frac{t(Y,e+2)^2}{t(Y,2e+2)}-(S-\tfrac12)>\frac1{480}.
\]
The proof bounds each infinite sum
$\sum_{d\ge1}\mu(d)/(2^d-1)^r$, $r\ge3$, between $1429/1512$ and $1$,
and its truncation error after $Y\ge4$ terms by $1/3584$.
Thus the lower bound holds for every stated pair $e,Y$, not just a finite
list of computed examples
(\lean{rankOneSubrankQuotient\_sub\_theta\_two\_gt}{ErdosProblems/Erdos249/RankOneSubrankObstruction.lean:238}).
These results concern the listed matrices, estimates, and comparison
sequence. The fourth item has the explicit construction and source theorem given
in Observation~\ref{prop:B4b-kill}; it has not been replayed here. A different finite-dimensional argument,
or one using further arithmetic assumptions, is not excluded.
Proposition~\ref{prop:period-not-rank} is conditional on rationality of
$S$ and is not, by itself, a counterexample. The separate rational $5/4$
comparison supplies the counterexample to the generic rank bound.
\coord{carry-rank} \scale{uniform} \ev{Lean}
\end{prop}
\leannote{Lean: \lproof{ErdosProblems/Erdos249/PaperCompleteR21/FourLinearConstructionLimits.lean}{72}{b6_retained_dyadic_sections_independent}, \lproof{ErdosProblems/Erdos249/PaperCompleteR21/FourLinearConstructionLimits.lean}{80}{b6_compressed_adjoint_identity_impossible}, \lproof{ErdosProblems/Erdos249/PaperCompleteR21/FourLinearConstructionLimits.lean}{101}{b6_mobius_incidence_unimodular_and_injective}, \lproof{ErdosProblems/Erdos249/PaperCompleteR21/FourLinearConstructionLimits.lean}{129}{b6_adjugate_tail_cost_floor}, and 6 further declarations in the \hyperref[sec:coverage]{coverage section}.}

\begin{prop}[The coefficient properties of the rational example]
\label{prop:b7}
The rational sequence in Proposition~\ref{prop:parity} satisfies uniform
boundedness, $c(n)\le n$, agreement with $\varphi(n)$ modulo $2$ at every
index, and the stated separated-carry form of aperiodicity. Consequently,
those properties alone cannot imply irrationality of a dyadic series.
The example does not exclude arguments using further totient identities,
including its multiplicative relations.
\coord{coefficient-word} \scale{uniform} \ev{Lean}
\end{prop}
\leannote{Lean: \lproof{ErdosProblems/Erdos249/PaperCompleteR21/RationalParityCountermodelProperties.lean}{28}{exists_rational_parityComparison}, \lproof{ErdosProblems/Erdos249/PaperCompleteR21/RationalParityCountermodelProperties.lean}{37}{parityComparisonProperties_do_not_imply_irrational}.}

\subsection{What the counterexamples establish}

The counterexamples isolate several distinctions. Finite prefix data are
not a statement about all larger indices. A congruence construction can
produce a large residue while producing an even larger error bound.
Periodicity modulo an integer does not control rational rank. An exact
reformulation leaves the original arithmetic assertion to be proved.

These are limitations of specified arguments, not an exhaustive
classification of possible methods. In particular, an averaging estimate is
one possible source of the missing information, not the only possible
source. A new arithmetic identity or a propagation theorem could also be
useful. The next section describes the precise assumptions of several
approaches considered here.

\section{Conditions that would imply irrationality}
\label{sec:survivors}

Four conditions considered below would imply irrationality. The first
three ask for quantitative estimates; the fourth is an exact reformulation
of irrationality. We also record a generic rank assertion that is false,
to prevent its being mistaken for an open sufficient hypothesis.
None of this is an exhaustive list of approaches. For each condition the
question is whether its particular arithmetic hypothesis can be proved.

\subsection{Cancellation on a full block}

\begin{defn}[The full-block norm condition]
\label{defn:fh}
Write $e(x) := \exp(2\pi i x)$ and
${E}(h,N,L) := e\bigl(({D}(h,N,L) \bmod 2^{L})/2^{L}\bigr)$.
The block norm condition asserts that
\[
  \forall h \ge 1\ \forall X_0\ \exists X, L : \quad
  \max(X_0,1) \le X, \quad
  16(2X + h + L + 2) \le 2^{L}, \quad
  \Bigl\lVert \sum_{N \in [X,\,2X)} {E}(h,N,L) \Bigr\rVert
  \le \tfrac{21}{25}X .
\]
It implies $S\notin\mathbb Q$.
\coord{first-harmonic} \scale{cofinal} \ev{Lean}
\lean{Erdos249257.DTWFirstHarmonicNormGap}{FirstHarmonicPivot.lean:75}
\lean{Erdos249257.irrational\_totient\_series\_of\_first\_harmonic\_norm\_gap}{FirstHarmonicPivot.lean:83}
\end{defn}

A bound on the real part also suffices: it is enough that
$\sum_{N \in [X,2X)} \operatorname{Re}E(h,N,L) \le \tfrac{9}{10}X$
(\lean{Erdos249257.exists\_certifiedKill\_of\_first\_harmonic\_gap}{FirstHarmonicGap.lean:128},
the elementary real-part estimate; the implication from $21/25$ to $9/10$ is
\lean{Erdos249257.exists\_certifiedKill\_of\_first\_harmonic\_norm\_gap}{FirstHarmonicPivot.lean:53}).
Unpacked, the required object is a constant-saving cancellation estimate for
the dyadically weighted totient-difference exponential sum
\[
  \sum_{X \le N < 2X}
  e\!\left(\frac{\sum_{j<L}\bigl(\varphi(N{+}h{+}1{+}j)-\varphi(N{+}1{+}j)\bigr)2^{\,L-1-j}}{2^{L}}\right),
  \qquad L \text{ subject to the displayed room inequality}.
\]
No theorem in the supplied record establishes this estimate with the
required quantifiers. The displayed room inequality imposes a lower
bound on $L$; it does not require $L=\log_2X+O(1)$. A depth of that
order would be a possible regime for the estimate, not an additional
proved conclusion. A finite instance would not suffice.
The norm hypothesis requires cancellation: if all phases are $1$, the
norm is $X$ and the bound fails; a balanced set of opposite phases has
sum zero. If instead every phase is $-1$, the real part is $-X$ and
satisfies the real-part bound, while the norm is $X$ and fails the norm
bound. Thus a real-part estimate is strictly weaker for general phase
families; these examples do not distinguish the two quantified conditions
for the actual totient sequence. The same depth $L$ must work across the
whole block. \ev{Open}

\paragraph{Comparison with the preceding counterexamples.} The earlier counterexamples do not establish or refute this block
estimate. The following distinctions explain why those particular arguments
do not settle it.


\begin{itemize}[leftmargin=*]
\item The parameters grow. The room condition forces the truncation depth
  to increase with $X$, so this is not a test of a fixed prefix.
\item The sample is prescribed: all $N\in[X,2X)$ are summed at one depth.
  Choosing two favourable indices does not establish this block estimate.
\item The sum uses the accumulated totient differences. It does not
  prescribe individual values through a congruence construction.
\item The estimate concerns the actual phases. Neither a bound on carry
  rank nor agreement only in parity gives it.
\item The proved implication uses the real-part threshold $9/10$ and the
  room constant $16$. The stronger norm bound $21/25$ implies that
  real-part bound. These constants do not assert any distribution theorem
  for the totient phases.
\end{itemize}

Formal sources for these implications: \lean{Erdos249257.windowFirstCos\_gt\_of\_not\_certifiedKill}{FirstHarmonicGap.lean:44}
\lean{Erdos249257.exists\_certifiedKill\_of\_first\_harmonic\_gap\_subset}{FirstHarmonicGap.lean:104}


\begin{rem}[The finite-prefix example does not choose a method]
\label{rem:b1-not-selector}
The finite-prefix example does not select an analytic method.
It shows only that a proof valid for every sequence $0\le c(n)\le n$
cannot deduce the required assertion from one fixed prefix. An argument
using additional identities of $\varphi$, a theorem propagating a finite
base case, a Farey estimate, or a suitable exponential sum could provide
information not contained in that prefix. The earlier finite audit cannot
rank all such possibilities.
\end{rem}

\begin{rem}[Finite certificate observations]
In the historical table of $28$ diagonal certificates through $t=64$,
the reported depths exceed the necessary depth bound by a small additive
constant. A separate enumeration at the indices where the least common
multiple increases records a closest margin from the central interval's
boundary of approximately $0.000221$ of the modulus, at $t=100$.
Both are finite observations. Neither establishes nor refutes
equidistribution or the estimate at arbitrarily large scales. \ev{Cert}
\end{rem}

\subsection{A four-term bound for the real part}\label{sub:four-sums}

The next condition separates one totient value from each phase, then
subtracts a mean on each group with the same cofactor. The definitions
below make all four sums explicit.

Fix $h,s\ge1$, $L\ge s+h$, $X\ge1$ and $0<\eta<1$, and put
$t=L-s+1$. For $X\le N<2X$, let $p_N$ be the largest prime factor of
$N+t$ and put $m_N=(N+t)/p_N$. Assign the index $N$ to this factorisation
when
\[
 m_N\le\left\lfloor\frac{\sqrt X}{2}\right\rfloor,
 \qquad p_N>2\lfloor\sqrt X\rfloor.
\]
Write $\mathcal A$ for the assigned indices and
$\mathcal G=\{N\in\mathcal A:\varphi(m_N)\ge\eta m_N\}$.
The two inequalities ensure $p_N>m_N\ge1$, so
$\varphi(N+t)=\varphi(m_N)(p_N-1)$.

The offset $t$ lies in the overlap of the two truncated tails. Their net
coefficient at $\varphi(N+t)$ is $(2^h-1)2^{-t}$. On the assigned indices,
set
\[
 z_N=e\!\left(\frac{(2^h-1)\varphi(m_N)(p_N-1)}{2^t}\right),
 \qquad w_N=\frac{E(h,N,L)}{z_N}.
\]
Both factors have absolute value $1$. Let $\bar z_m$ be the average of
$z_N$ over the assigned indices with cofactor $m$, taking the average to
be zero when the group is empty. The four contributions are then defined
by the identity
\[
\begin{aligned}
 \sum_{N=X}^{2X-1}E(h,N,L)
 &=\underbrace{\sum_{N\in\mathcal G}w_N(z_N-\bar z_{m_N})}
                 _{\text{mean-subtracted contribution}}
   +\underbrace{\sum_{N\in\mathcal G}w_N\bar z_{m_N}}
                 _{\text{mean contribution}}\\
 &\quad+\underbrace{\sum_{N\in\mathcal A\smallsetminus\mathcal G}E(h,N,L)}
                 _{\text{excluded-cofactor contribution}}
   +\underbrace{\sum_{\substack{X\le N<2X\\N\notin\mathcal A}}E(h,N,L)}
                 _{\text{unassigned contribution}}.
\end{aligned}
\]
This follows by partitioning the indices and using $E(h,N,L)=w_Nz_N$.
Subtracting the group mean is exact algebra; it does not make the two
factors independent. The chosen prime need not be confined to the selected
argument: with $X=16$, $h=s=1$, $L=20$ and $N=18$, the argument is
$N+t=38=2\cdot19$, which satisfies both assignment inequalities, but the
same prime also divides $N+1=19$. Thus the factorisation at $N+t$ alone
gives no separation from the other totient values in the phase.


\begin{defn}[Four simultaneous inequalities]
\label{defn:pivot}
For every $h > 0$ there are $s > 0$ and $\eta \in (0,1)$ such that for every
$X_0$ there are $X, L$ with $\max(X_0,1) \le X$, $h \le L - s$,
$16(2X{+}h{+}L{+}2) \le 2^{L}$, and all four of
\[
\begin{aligned}
 \operatorname{Re}\sum_{N\in\mathcal G}w_N(z_N-\bar z_{m_N})
   &\le \tfrac{14}{25}X,\\
 \left\lVert\sum_{N\in\mathcal G}w_N\bar z_{m_N}\right\rVert
   &\le \tfrac{1}{100}X,\\
 \left\lVert\sum_{N\in\mathcal A\smallsetminus\mathcal G}E(h,N,L)\right\rVert
   &\le \tfrac{1}{100}X,\\
 \left\lVert\sum_{\substack{X\le N<2X\\N\notin\mathcal A}}E(h,N,L)\right\rVert
   &\le \tfrac{8}{25}X.
\end{aligned}
\]
It implies $S\notin\mathbb Q$.
\coord{first-harmonic} \scale{cofinal} \ev{Lean}
\lean{Erdos249257.DTWPivotResidualDecorrelation}{FirstHarmonicPivot.lean:569}
\lean{Erdos249257.irrational\_totient\_series\_of\_pivotResidualDecorrelation}{FirstHarmonicPivot.lean:577}
\end{defn}

The four bounds must hold for the same $X,L,s,\eta$, not along four
unrelated sequences of blocks. Their constants add to
$14/25+1/100+1/100+8/25=9/10$, so they imply the real-part criterion.
They do not establish the stronger $21/25$ norm criterion.
The exact partition proves an identity, not any one of these estimates;
in particular, grouping by a prime factor does not justify statistical
independence between the mean-subtracted phase and its weight.

Like the block criterion, this condition concerns actual totient values
at arbitrarily large indices and at a common truncation depth. The four
bounds are a sufficient way to prove the needed real-part estimate, not
an equivalent restatement of the norm estimate.


The decomposition is an exact identity. The supplied formal record
contains its unconditional proof and an axiom report
(\lean{Erdos249257.windowFirstExp\_sum\_eq\_pivot\_decomposition}{FirstHarmonicPivot.lean:514},
budget implication
\lean{Erdos249257.first\_harmonic\_gap\_of\_pivotBudgetAt}{FirstHarmonicPivot.lean:549}).
The mean being subtracted is the mean of $z_N$, not of the residual
weight $w_N$. On each cofactor group, $\sum_N(z_N-\bar z_m)=0$, but
$\sum_Nw_N(z_N-\bar z_m)$ need not vanish because $w_N$ varies with $N$.
Thus centring the phase does not estimate the weighted correlation. Only its real part is required:
$14X/25$ suffices in place of the earlier proposed norm bound $X/2$.
For $X\ge4$ and $s=L-h$, the assigned indices \emph{with cofactor
$m_N=1$} are exactly those for which $N+h+1$ is prime. This describes
one cofactor group, not all assigned indices. The source proves this
membership equality
(\lean{Erdos249257.mem\_pivotFiber\_one\_overlap\_iff}{FirstHarmonicPivot.lean:423}),
using $\varphi(mp)=\varphi(m)(p-1)$ and $p>m$.
This special choice has $s=L-h$ varying with $L$; it illustrates the
arithmetic of the groups but does not establish the condition with $s$
fixed before $X_0$.

The four inequalities are not proved with their required common parameters.
Propositions~\ref{prop:dickman} and~\ref{prop:badcof} below do establish the
unassigned and excluded-cofactor bounds at minimal admissible depth, with
one sufficiently small fixed $\eta$. The mean and mean-subtracted estimates
remain to be established at that same depth and for that same $\eta$.
Only the real part of the mean-subtracted contribution is required.


A nonzero determinant alone does not prove the needed correlation
estimate. The comparison in the linked source preserves such a determinant
while fixing an entire row of weighted phases to $1$. A successful
application must distinguish that comparison, for example by additional
arithmetic information, rather than only by column rescaling. The exact
four-term identity supplies no prime-distribution estimate. \ev{Open}
\lean{Erdos249257.locked\_reconstruction\_preserves\_nonzero\_minor}{ResidualGaugeObstruction.lean:94}


\subsection{A fixed separation at prime positions}

\begin{defn}[A fixed gap at prime positions]
\label{defn:primegap}
For every $h \ge 1$ and every $N_0$ there is a prime $p$ with
$\max(N_0{+}h{+}1,\, h{+}5) \le p$ and
$\mathrm{Re}\,e(R_{p-1}-R_{p-h-1}) < \tfrac{9}{10}$,
where $e(x)=\exp(2\pi i x)$. Equivalently:
cofinally many primes $p$ at which the totient tail difference across the shift
$h$ stays a \emph{fixed} distance
$> \arccos(9/10)/(2\pi) \approx 0.0718$ from every integer.
It implies $S\notin\mathbb Q$.
\coord{prime-pivot} \scale{cofinal} \ev{Lean}
\lean{DTWNaturalPrimeTailOrbitStrictGap}{ErdosProblems/Erdos249/TotientStrictPrimeEscape.lean:157}
\lean{irrational\_totient\_series\_of\_naturalPrimeTailOrbitStrictGap}{ErdosProblems/Erdos249/TotientStrictPrimeEscape.lean:524}
\end{defn}

This asks for more than nonintegrality. A nonintegral real number can
be arbitrarily close to an integer, whereas the displayed condition requires
a fixed positive distance, along prime positions and for each shift $h$.
For example, phases at a half-integer satisfy the numerical gap and phases
converging to an integer eventually fail it. These are illustrations of the
hypothesis, not assertions about the totient tails. No estimate for these tail differences at prime positions
with the stated quantifiers is proved here. Certificate completeness alone
does not supply it. \ev{Open}

\subsection{Certificates with depth equal to the shift}

\begin{defn}[Depth equal to a shift multiple]
\label{defn:apfde}
For every $d \ge 1$ and every $N$ there is $t \ge 1$ with
$\mathcal{C}(td,\,N,\,td)$; unpacked,
\[
  N + 2td + 2 \;<\; {D}(td,N,td) \bmod 2^{td}
  \;<\; 2^{td} - (N + 2td + 2).
\]
It is equivalent to $S\notin\mathbb Q$.
\coord{period-ray} \scale{cofinal} \ev{Lean}
\lean{ApFullDepthEscape}{ErdosProblems/Erdos249/PeriodMultipleEscape.lean:441}
\lean{apFullDepthEscape\_iff\_irrational}{ErdosProblems/Erdos249/FullDepthRayAmplifier.lean:406}
\end{defn}

For the actual totient series, the condition in Definition~\ref{defn:apfde} is equivalent to
irrationality, by the full-depth amplification argument of the short note
(Theorem~3.1 there;
\lean{apFullDepthEscape\_iff\_irrational}{ErdosProblems/Erdos249/FullDepthRayAmplifier.lean:406}).
For a fixed pair $(d,N)$, a nonintegral tail difference supplies a certificate
in every sufficiently late adjacent pair of multipliers. The missing
arithmetic input is a nonintegral tail difference for every relevant pair,
not a converse implication for the condition with depth equal to the shift. This remains an exact
reformulation of the open problem, not a proof of it. An older one-direction
lemma
(\lean{irrational\_totient\_series\_of\_apFullDepthEscape}{ErdosProblems/Erdos249/PeriodMultipleEscape.lean:450})
is still in the cited source; it is not the current classification.

The finite examples establish the hypothesis for individual pairs
$(d,N)$, not for all pairs. Thus the universally quantified condition is
not a mild regularity assumption: it has the full strength of the
irrationality question. \ev{Open}

\noindent Formal sources: \lean{windowDiscrepancy\_self\_eq\_totientBlock\_sub}{ErdosProblems/Erdos249/PeriodMultipleEscape.lean:67}

\subsection{A false generic rank bound}

\begin{defn}[A counterfactual rank criterion]
\label{defn:rankupper}
The proposed assertion is a constant $C$ bounding the rank for every
$c:\N\to\N$ with $c(n)\le n$ and $T_c\in\Q$, every integer $v>0$,
and every integer sequence $u$ satisfying
\[
 u(n+1)=2u(n)-v c(n+1),\qquad u(n)/2^n\longrightarrow0.
\]
Specifically, the proposed bound is
\[
 \dim_{\Q}\operatorname{span}_{\Q}
 \{n\mapsto u(2^jn+r):1\le j\le e,\ 0\le r<2^j\}\le C
 \qquad(e\ge0).
\]
The alternative replaces $C$ by a function $g(e)$ growing strictly slower
than $2^e-1$, with the same quantifiers. Either assertion would contradict
the lower bound of Proposition~\ref{prop:rank} under rationality of $S$.
Both assertions are false, as the following comparison shows.
\coord{carry-rank} \scale{uniform}
\end{defn}

The generic assertion in Definition~\ref{defn:rankupper} is false, not a
live open proposition. The $5/4$ control in
Section~\ref{sec:mahler-defect} has $0\le c(n)\le n$ and a
tempered integral carry of rank at least $2^e-1$ at every depth
(the unbounded-rank theorem for the rational comparison sequence, line~478).
It refutes both proposed generic ceilings. A bound restricted by additional
totient-specific hypotheses would be a different assertion.

Periodicity modulo $v$ alone does not provide the additional rank bound.
\lean{Erdos249257.false\_of\_compressedAdjointCertificate}{TotientMahlerDefect.lean:1183}


\section{Guide to the detailed statements}

The remaining sections are reference material. They restate the
identities and implications in enough detail to compare their hypotheses.
A repeated statement has the same meaning as before; its repetition does
not supply a missing assumption.

\paragraph{Quantifiers and hypotheses.} The direct certificate computations cover individual parameters or a
finite range. Other results hold uniformly or at arbitrarily large scales,
but their conclusions are different: positivity, for example, is not
nonintegrality. The finite-prefix example, rank bounds, and residue estimates
must each be used only under their stated hypotheses.

\paragraph{Dependencies.} Sections~\ref{sec:wall} and~\ref{sec:survivors} explain the principal
counterexamples and sufficient conditions. The following catalogue supplies
the associated formulas and source references. The later tables compare
logical implications and the limits of particular estimates; their source
labels are not extra mathematical hypotheses.
\section{Series identities and finite exclusions}\label{sec:series}

The identities in this section express $S$ through finite prefixes,
scaled tails and M\"obius sums. The finite calculation in
Section~\ref{ssec:farey} then excludes a specified range of rational
denominators. Equivalent series representations do not by themselves
extend that finite exclusion to all denominators.

\subsection{The series}

\begin{defn}[The Erd\H{o}s--249 constant]
\label{defn:S}
Let $\varphi$ denote Euler's totient function. Define
\[
S \;:=\; \sum_{n \ge 1} \frac{\varphi(n)}{2^{n}} \;=\; \sum_{n : \mathbb{N}} \frac{\varphi(n)}{2^{n}},
\]
the two forms coinciding because $\varphi(0)=0$; the second, $\mathbb{N}$-indexed
form is the one carried in the Lean source. The series converges absolutely
since $\varphi(n)\le n$. The irrationality of this value is the question
under consideration.
\end{defn}

\begin{rem}[Status]
No proof or disproof of $S\notin\mathbb Q$ exists anywhere in the record,
formal or informal. What exists is: one large unconditional finite denominator
exclusion (\S\ref{ssec:farey}), several exact reformulations of $S$ that
relocate the same open question onto different coordinates without touching
its truth value (\S\ref{ssec:coprime}--\S\ref{ssec:lambert}), and a certificate
apparatus whose cofinal quantified condition is the precise open target
(Definition~\ref{defn:sep}; only its finite, checked instances
are catalogued here, \S\ref{ssec:certtable}).
\end{rem}

\subsection{The tail-shift identity}

\begin{prop}[Digit-shift identity]
\label{prop:shift}
For every $N : \mathbb{N}$,
\[
2^{N} \cdot S \;=\; \Phi_N + R_N,
\]
where $\Phi_N := \sum_{n \le N} \varphi(n) \cdot 2^{N-n} \in \mathbb{N}$ (the
integer prefix) and $R_N := \sum_{j \ge 0} \varphi(N+1+j)/2^{j+1}$ (the
scaled tail, \ensuremath{{R}} in Lean).
\coord{binary-digit} \scale{uniform} \ev{Lean}
\lean{two_pow_mul_totient_series_eq}{TotientTailPeriodKiller.lean:150}
\end{prop}

\begin{proof}[Consequence for rationality]
Subtracting the prefix identities at $N+h$ and $N$ gives
\[
 R_{N+h}-R_N\equiv2^N(2^h-1)S\pmod\Z.
\]
If $S=a/(2^cv)$ with $v\ge1$ odd, Euler's theorem supplies
$h=\varphi(v)>0$ and $N_0=c$, for which the difference is integral for
all $N\ge N_0$
(\lean{eventual_period_of_not_irrational}{TotientTailPeriodKiller.lean:358},
\ev{Lean}). Conversely, an integral difference for even one $h>0,N\ge0$
forces $S$ to be rational, because $2^N(2^h-1)$ is a nonzero integer.
Thus the unresolved question is not this converse, but whether every such
tail difference is nonintegral. Definition~\ref{defn:sep} expresses that
question through finite residues.
\end{proof}

\subsection{The finite denominator bound}
\label{ssec:farey}

The finite calculation gives a lower
bound on its denominator should it happen to be rational. It is obtained by a
Farey mediant argument that does not assume the quantified residue
condition or use the tail-period test.

\begin{lem}[Farey gap, fully general]
\label{lem:farey}
For integers $a,b,c,d,r,s$ with $b>0$, $d>0$, $bc-ad=1$ (i.e.\ $a/b$ and
$c/d$ are unimodular Farey neighbours), and $as < rb$, $rd < cs$ (i.e.\ $r/s$
lies strictly between them): $b+d \le s$.
\coord{farey} \scale{n/a} \ev{Lean}
\lean{farey_gap}{GapFareyBound.lean:51}
\end{lem}

Indeed, the two strict inequalities make $cs-dr$ and $br-as$ positive
integers, and
\[
 s=b(cs-dr)+d(br-as)\ge b+d,
\]
using $bc-ad=1$. This also proves $s>0$, which the formulation did not
need to assume separately. The lemma is the classical Stern--Brocot
denominator bound. The mediant $(a+c)/(b+d)$ lies between the two
neighbours and is reduced: any common divisor of $a+c$ and $b+d$
divides $b(a+c)-a(b+d)=bc-ad=1$. It attains denominator $b+d$,
so the general interval bound is sharp. A particular tail interval
still requires its own endpoint and first-failure checks.

\begin{prop}[The Farey certificate at $K=240$]
\label{prop:gapwindow}
For the window $(N,K)=(1,240)$, put
\[
 V=\left(\sum_{r=1}^{240}\varphi(1+r)2^{240-r}\right)\bmod2^{240}.
\]
For every $q : \mathbb{N}$ with
\[
0 < q \;\le\; Q_0 := 79\,639\,646\,646\,701\,375\,323\,355\,774\,875\,831\,053 \;\;(\approx 7.96\times 10^{34}),
\]
\[
(q \cdot V) \bmod 2^{240} \;+\; 243\,q \;<\; 2^{240}.
\]
This bound is \emph{sharp}: $q = Q_0+1 = 79\,639\,646\,646\,701\,375\,323\,355\,774\,875\,831\,054$
is the exact first failing denominator, obtained as the denominator of
the mediant of two explicit unimodular Farey neighbours.
\coord{farey} \scale{bounded} \ev{Lean}
\lean{gap_check_window_1_240_le_79639646646701375323355774875831053}{GapFareyBound.lean:176}
\lean{gap_check_window_1_240_first_failure}{GapFareyBound.lean:225}
\end{prop}
\leannote{Lean: \lproof{Erdos249257/CertificateKernel.lean}{18310}{totient_carry_residue_window_1_240_eq}, \lproof{Erdos249257/GapFareyBound.lean}{176}{gap_check_window_1_240_le_79639646646701375323355774875831053}, \lproof{Erdos249257/GapFareyBound.lean}{225}{gap_check_window_1_240_first_failure}.}

\begin{thm}[Denominator exclusion]
\label{thm:denom-record}
For every $p\in\Q$ whose reduced denominator is at most $Q_0$,
\[
S \;\neq\; p.
\]
Equivalently: \emph{if $S$ is rational, its reduced denominator exceeds}
$Q_0 \approx 7.96 \times 10^{34}$.
\coord{farey} \scale{bounded} \ev{Lean}
\lean{tsum_totient_div_pow_two_ne_ratCast_of_den_le_79639646646701375323355774875831053}{CertificateKernel.lean:18384}
\end{thm}

\begin{rem}[What this does and does not say]
Theorem~\ref{thm:denom-record} is a complete, unconditional finite fact: no rational
of denominator at most $Q_0$ equals $S$. A single such bound does not
exclude arbitrarily large denominators. Proving unbounded exclusions by
the same mechanism would, however, establish irrationality. The docstring behind
Proposition~\ref{prop:gapwindow} states the question this leaves:
whether $\sup_K (b+d)(K)$ (the growing analogue of $Q_0$ as the window $K$
grows) is unbounded as $K \to \infty$; which would close \#249 through this
theorem's implication, entirely independently of the certificate-supply
obligation of Proposition~\ref{prop:shift}. \ev{Open}; not claimed or proved
anywhere in the record.
\end{rem}

\begin{rem}[Earlier finite denominator bounds]
The recorded Farey computations gave bounds $4838$, $2^{22}$,
approximately $2.49\times10^{17}$ at $K=120$, and $Q_0$ at $K=240$.
Each applies the finite mechanism of Lemma~\ref{lem:farey} and
Proposition~\ref{prop:gapwindow}. These values do not prove that the
bounds are unbounded as $K$ increases.
\end{rem}

\subsection{Equivalent expressions and their denominator bounds}
\label{ssec:coprime}

The exact finite Farey record behind Theorem~\ref{thm:denom-record} is not tied to
the $\varphi(n)/2^n$ presentation of $S$; it transfers verbatim to two other
exact reformulations of the same constant, at exactly half the bound.
Amiram Eldar posted both reformulations to OEIS A256936 on 15 March 2026: an
equivalent coprimality interpretation of Proposition~\ref{prop:coprime} in
revision~28, and the M\"obius-square formula of Proposition~\ref{prop:mobsq}
in revision~31~\cite{eldar2026oeis}.  Steve Fan posted the M\"obius-square
formula on the problem's forum thread on 16 May 2026~\cite{fan2026totient}.
Both identities follow here from M\"obius inversion and geometric sums, with
the proofs recorded in the linked declarations.

\begin{prop}[Fair-coin coprimality form]
\label{prop:coprime}
Let $X,Y$ be independent random variables with $\Pr(X=n)=\Pr(Y=n)=2^{-n}$
for $n \ge 1$ (independent fair-coin waiting times). Then
\[
S \;=\; \tfrac12 \;+\; \Pr\bigl(\gcd(X,Y)=1\bigr)
\;=\; \tfrac12 \;+\; \sum_{\substack{a,b\ge 1\\ \gcd(a,b)=1}} 2^{-(a+b)}.
\]
Equivalently, on the visible lattice: summing $2^{-(a+b)}$ over the half-open
coprime pairs ($a\ge1$, $b\ge0$, $\gcd(a,b)=1$) recovers $\sum_n \varphi(n)/2^n$
exactly, with no boundary correction, because the visible-point count on the
half-open antidiagonal at height $n$ equals $\varphi(n)$ for every $n$,
including $n=0,1$.
\coord{probability} \scale{n/a} \ev{Lean}
\lean{tsum_visible_coprime_pairs_eq_totient_series}{CertificateKernel.lean:18544}
\lean{totient_series_eq_half_add_visible_coprime_pairs}{CertificateKernel.lean:18557}
\end{prop}

\begin{prop}[The gcd-layer normalisation]
\label{prop:gcdlayer}
For independent fair-coin waiting times as above, $\sum_{g\ge1}\Pr(\gcd(X,Y)=g)=1$
exactly; and for every $d>0$, $\Pr(d\mid X \wedge d\mid Y) = 1/(2^d-1)^2$.
\coord{probability} \scale{uniform} \ev{Lean}
\lean{tsum_pos_coprime_inv_mersenne_eq_one}{GcdMomentCalculus.lean:349}
\lean{tsum_pos_pair_both_dvd_half_eq_inv_mersenne_sq}{GcdMomentCalculus.lean:266}
\end{prop}
\leannote{Lean: \lproof{Erdos249257/GeometricCoprimality.lean}{480}{tsum_gcd_layer_pos_coprime_half_eq_one}, \lproof{Erdos249257/GcdMomentCalculus.lean}{266}{tsum_pos_pair_both_dvd_half_eq_inv_mersenne_sq}.}

\begin{thm}[Denominator exclusion for the coprimality-probability form]
\label{thm:denomcoprime}
Let
\[
Q_1 := \left\lfloor \frac{Q_0}{2} \right\rfloor = 39\,819\,823\,323\,350\,687\,661\,677\,887\,437\,915\,526.
\]
For every $a\in\mathbb{Z}$, $d\in\mathbb{N}$ with $0<d\le Q_1$: the visible
coprime-pair probability $\Pr(\gcd(X,Y)=1)$ is not equal to $a/d$.
\coord{farey} \scale{bounded} \ev{Lean}
\lean{tsum_visible_coprime_pairs_ne_int_div_of_den_le_39819823323350687661677887437915526}{CertificateKernel.lean:18572}
\end{thm}

\begin{thm}[Denominator exclusion for the M\"obius-square form]
\label{thm:denommobsq}
With $Q_1$ as in Theorem~\ref{thm:denomcoprime}: for every $a\in\mathbb{Z}$,
$d\in\mathbb{N}$ with $0<d\le Q_1$, the signed series
$T := \sum_{m\ge1} \mu(m)/(2^m-1)^2 = S - \tfrac12$ (see \S\ref{ssec:mobius})
is not equal to $a/d$.
\coord{farey} \scale{bounded} \ev{Lean}
\lean{tsum_moebius_div_two_pow_sub_one_sq_ne_int_div_of_den_le_39819823323350687661677887437915526}{CertificateKernel.lean:18487}
\end{thm}

\begin{rem}[Why the bound halves, and why this is not new information]
$Q_1 = \lfloor Q_0/2 \rfloor$ arithmetically: $Q_0$ is odd
($Q_0 = 79\,639\,646\,646\,701\,375\,323\,355\,774\,875\,831\,053$), so
$Q_0/2 = 39\,819\,823\,323\,350\,687\,661\,677\,887\,437\,915\,526.5$ and
$Q_1$ is its floor. The halving is the cost of transporting the known bound
through the affine shift
$S = \tfrac12 + T$ (resp.\ $S = \tfrac12 + \Pr(\gcd(X,Y)=1)$) on a
denominator-exclusion statement: excluding all denominators $\le Q_1$ for $T$
follows from excluding all denominators $\le Q_0$ for $S$ (a denominator-$d$
value of $T$ with $d\le Q_1$ yields a
denominator dividing $2d\le Q_0$ for $S$). Theorems~\ref{thm:denomcoprime}
and~\ref{thm:denommobsq} are therefore the \emph{same} finite Farey record as
Theorem~\ref{thm:denom-record}, transported through Proposition~\ref{prop:coprime}
and Proposition~\ref{prop:mobsq} respectively; not independent evidence.
The converse finite implication is not asserted: subtracting $1/2$ can double
a denominator, so the $Q_1$ exclusion for $T$ alone need not recover the full
$Q_0$ exclusion for $S$.
\end{rem}

\subsection{The M\"obius-weighted identity}
\label{ssec:mobius}

\begin{prop}[M\"obius-square reduction]
\label{prop:mobsq}
\[
S \;=\; \sum_{n\ge1}\frac{\varphi(n)}{2^n} \;=\; \frac12 \;+\; \sum_{d\ge1}\frac{\mu(d)}{(2^d-1)^2},
\]
where $\mu$ is the M\"obius function, so $\mu(d)\in\{-1,0,1\}$ for every $d$.
Consequently \emph{Erd\H{o}s \#249 $\iff$ $T:=\sum_{d\ge1}\mu(d)/(2^d-1)^2 \notin \mathbb{Q}$},
since adding a rational number does not change irrationality.
\coord{mobius-mersenne} \scale{n/a} \ev{Lean}
\lean{totientSeries_eq_pnat_half_pow}{SquaredMersenneDiagonalEnclosure.lean:72}
\lean{tsum_totient_half_pow_eq_half_add_moebius_sq}{MersenneLambertLadder.lean:665}
\end{prop}
\leannote{Lean: \lproof{ErdosProblems/Erdos249/PaperCompleteR20/MobiusSquareReduction.lean}{10}{mobius_square_reduction}, \lproof{ErdosProblems/Erdos249/PaperCompleteR20/MobiusSquareReduction.lean}{17}{irrational_totient_iff_mobius_square}, \lproof{ErdosProblems/Erdos249/PaperCompleteR20/MobiusSquareReduction.lean}{25}{moebius_three_values}.}

\begin{rem}[Why the bounded coefficients matter; the Erd\H{o}s-1948 regime]
The identity of Proposition~\ref{prop:mobsq} rewrites $S$ (equivalently $T$)
as a Lambert-type series with squared denominators and bounded
M\"obius weights. To see the constant $1/2$ directly, expand
\[
 \sum_{d\ge1}\frac{\mu(d)}{2^d-1}
   =\sum_{n\ge1}\frac{\sum_{d\mid n}\mu(d)}{2^n}=\frac12.
\]
Combining this identity with $\varphi=\mu*\mathrm{Id}$ and
$2^d=(2^d-1)+1$ proves the formula for $T$; absolute convergence
justifies the rearrangements. Thus the corresponding first-power sum
$L(\mu):=\sum_d\mu(d)/(2^d-1)$ is rational, whereas
$L(1)=\sum_{d\ge1}1/(2^d-1)=E$ is the Erd\H{o}s--Borwein constant.
Erd\H{o}s proved its irrationality in 1948~\cite{erdos1948lambert} and the supplied formal source also proves it irrational
(\lean{irrational_erdosBorwein_series}{CertificateKernel.lean:8335}, \ev{Lean}).

\par\medskip

In contrast, the coefficients of the original power series satisfy
$0\le\varphi(n)\le n$ and are unbounded. Bounded M\"obius weights do
not make the new denominators powers of a fixed base, so they do not
turn this representation into a bounded-digit expansion.  The function $\varphi$ has average
order $6n/\pi^2$, equivalently
$\sum_{k\le x}\varphi(k)\sim 3x^2/\pi^2$, but there is no pointwise estimate
$\varphi(n)=\Theta(n)$. The near-integer irrationality criterion
(formalised generically as
\lean{irrational_of_int_mul_near_int}{CertificateKernel.lean:6120} and its
base-power specialisation
\lean{irrational_of_pow_mul_near_int}{CertificateKernel.lean:6149}, both
\ev{Lean}) applies to either representation if its nonzero approximation
errors can be established. Boundedness of the coefficients does not itself
establish those errors.
Erd\H{o}s's 1948 argument is itself a proof for the specific Lambert series:
it chooses arguments by a system of congruences so that the divisor counts
along a block are divisible by increasing powers of the base, and a tail bound
then forces a long run of zero digits~\cite[pp.~63--66]{erdos1948lambert}.
Boundedness of a signed weight supplies none of that divisibility input.

\par\medskip

The rational identity $L(\mu)=1/2$ is already a warning: bounded signed
coefficients alone do not imply irrationality. The later denominator
calculations retain the M\"obius terms whose indices do not divide the
chosen LCM; those terms cannot be discarded when transferring an estimate.
\end{rem}

\subsection{A general squared-Lambert identity}
\label{ssec:lambert}

\begin{prop}[The squared-Lambert identity]
\label{prop:lambertengine}
For $w:\mathbb{N}\to\mathbb{R}$ with $|w(d)|\le d$ for all $d>0$, and $0\le r<1$:
\[
\sum_{d\ge1} w(d)\left(\frac{r^d}{1-r^d}\right)^2 \;=\; \sum_{n\ge1}\left(\sum_{e\mid n} w(e)\Bigl(\tfrac{n}{e}-1\Bigr)\right) r^n.
\]
\coord{mobius-mersenne} \scale{uniform} \ev{Lean}
\lean{tsum_lambert_linear_weight_sq_pure}{GcdMomentCalculus.lean:105}
\end{prop}

\noindent The identity follows by expanding $r^d/(1-r^d)$ as a geometric
series and interchanging absolutely convergent sums.  It is the squared
analogue of the classical divisor-sum Lambert identity
$\sum_na_nq^n/(1-q^n)=\sum_m\bigl(\sum_{d\mid m}a_d\bigr)q^m$~\cite[(1),
p.~2]{merca-schmidt}.

At $r=1/2$, the choices $w=1$ and $w=\varphi$ give the following two
identities. They help distinguish the role of the weight from that of the
squared denominator.

\begin{prop}[The divisor-sum identity]
\label{prop:zetaq}
\[
\sum_{d\ge1} \frac{1}{(2^d-1)^2} \;=\; \sum_{n\ge1} \frac{\sigma(n)-\tau(n)}{2^n} \;=\; \zeta_q(2)-\zeta_q(1) \text{ at } q=\tfrac12,
\]
where $\sigma$ is the sum-of-divisors function and $\tau$ the number-of-divisors
function. The \emph{identity} is machine-checked (\ev{Lean}).  With
$\zeta_q(s)=\sum_{n\ge1}n^{s-1}q^n/(1-q^n)$, Postelmans and Van Assche prove
that $1,\zeta_q(1),\zeta_q(2)$ are linearly independent over $\mathbb{Q}$ for
$q=1/p$ with an integer $p\ge2$~\cite[Theorem~1.3, p.~3]{postelmans-vanassche};
at $q=\tfrac12$ this gives the irrationality of the \emph{value}
$\zeta_q(2)-\zeta_q(1)$ (\ev{Cited}, \emph{not} formalised in this corpus).
Two irrational numbers can have a rational difference, so the separate
irrationality of the two values would not suffice.
\coord{mobius-mersenne} \scale{n/a}
\lean{tsum_one_div_mersenne_sq_eq_sigma_sub_tau_series}{GcdMomentCalculus.lean:216}
\end{prop}
\leannote{Lean: \lproof{ErdosProblems/Erdos249/PaperCompleteR21/SquaredMersenneDivisorIdentities.lean}{36}{divisor_sum_identity}, \lproof{ErdosProblems/Erdos249/PaperCompleteR21/QZetaAnchor.lean}{17}{qZeta}, \lproof{ErdosProblems/Erdos249/PaperCompleteR21/QZetaAnchor.lean}{81}{qZeta_half_two_sub_one}, \lproof{ErdosProblems/Erdos249/PaperCompleteR21/QZetaAnchor.lean}{143}{irrational_qZeta_half_difference_of_linearIndependent}, and 1 further declaration. Conditional on the linear independence theorem of Postelmans and Van Assche; see the \hyperref[sec:coverage]{coverage section}.}

\begin{prop}[The gcd-moment identity]
\label{prop:pillai}
\[
\sum_{d\ge1} \frac{\varphi(d)}{(2^d-1)^2} \;=\; \sum_{n\ge1} \bigl(P(n)-n\bigr)\cdot 2^{-n} \;=\; \mathbb{E}[\gcd(X,Y)],
\]
where $P(n) := \sum_{e\mid n}\varphi(e)\cdot(n/e) = (\varphi * \mathrm{Id})(n)$
is Pillai's gcd-sum function $P(n)=\sum_{k\le n}\gcd(k,n)$~\cite[\S1,
(1)--(2)]{toth-gcd} and $X,Y$ are the independent fair-coin waiting
times of Proposition~\ref{prop:coprime}. The identity itself is
machine-checked (\ev{Lean}); irrationality of $\mathbb{E}[\gcd(X,Y)]$ is not proved here. This is a
different series from $S$.
\coord{mobius-mersenne} \scale{n/a}
\lean{tsum_totient_div_mersenne_sq_eq_gcd_moment_series}{GcdMomentCalculus.lean:235}
\end{prop}
\leannote{Lean: \lproof{ErdosProblems/Erdos249/PaperCompleteR21/PillaiGcdExpectation.lean}{334}{gcd_moment_identity_three_members}, \lproof{ErdosProblems/Erdos249/PaperCompleteR21/PillaiGcdExpectation.lean}{307}{tsum_pos_pair_gcd_half_eq_totient_div_mersenne_sq}, \lproof{ErdosProblems/Erdos249/PaperCompleteR21/PillaiGcdExpectation.lean}{77}{sum_gcd_Icc_eq_pillaiP}, \lproof{ErdosProblems/Erdos249/PaperCompleteR21/PillaiGcdExpectation.lean}{67}{pillaiP_eq_totient_mul_id}, and 1 further declaration in the \hyperref[sec:coverage]{coverage section}.}

\begin{rem}[Related series]
Writing $L(f):=\sum_d f(d)/(2^d-1)$ and
$L_2(f):=\sum_d f(d)/(2^d-1)^2$, the preceding identities give
$L(\mu)=\tfrac12$ (rational), $L(1)=E$ (Erd\H{o}s--Borwein,
irrational, \ev{Lean}), $L(\varphi)=2$ (rational); $L_2(\mu)=S-\tfrac12$
(\ev{Open}, this is \#249), $L_2(1)=\zeta_q(2)-\zeta_q(1)$ (\ev{Cited}
irrational, Proposition~\ref{prop:zetaq}), $L_2(\varphi)=\mathbb{E}[\gcd(X,Y)]$
(\ev{Open}, Proposition~\ref{prop:pillai}). These examples show why the denominator shape alone does not determine
irrationality: changing the weight can produce either rational or
irrational values.
\end{rem}

\subsection{Diagonal certificates through scale 82}
\label{ssec:certtable}

This table records finite instances of the residue test introduced in
Definition~\ref{defn:sep}. The aggregate result covers every integer scale
$t\le82$. It does not assert an instance at scale $83$ or an unbounded
family. The individual source links identify the computations used.

\paragraph{Definition of the finite residue test} For $h,N,L : \mathbb{N}$, recall the truncated
difference $D(h,N,L) := \sum_{j<L} (\varphi(N+h+1+j)-\varphi(N+1+j))\cdot 2^{L-1-j} \in \mathbb{Z}$
(\lean{windowDiscrepancy}{TotientTailPeriodKiller.lean:66}, \ev{Lean}), and
\[
\mathcal{C}(h,N,L) \;:\equiv\; (N+h+L+2) < D(h,N,L) \bmod 2^{L} < 2^{L}-(N+h+L+2),
\]
(\lean{certifiedKill}{TotientTailPeriodKiller.lean:72}, \ev{Lean}, decidable). By
completeness (\lean{exists_certifiedKill_iff_tail_diff_notMem_int}{LcmConeFlatness.lean:316},
\ev{Lean}), 
\[
(\exists L,\ \mathcal{C}(h,N,L)) \iff R_{N+h}-R_N \notin \mathbb{Z}
\]

for every $h,N$. Requiring such certificates beyond every basepoint
threshold for each fixed positive $h$ is precisely $\mathrm{Sep}$ and is
equivalent to $S\notin\mathbb Q$. The table records finite instances,
indexed either by $(h,N,L)$ or by $t$; it does not establish that quantified
condition.

\small
\begin{longtable}{@{}
  >{\raggedright\arraybackslash}p{0.17\linewidth}
  >{\raggedright\arraybackslash}p{0.40\linewidth}
  >{\raggedright\arraybackslash}p{0.13\linewidth}
  >{\raggedright\arraybackslash}p{0.20\linewidth}
  @{}}
\textbf{Anchor} & \textbf{Exact statement of what was checked} & \textbf{Scale} & \textbf{Site} \\
\hline
\endhead
Fixed-window example, depth 16
&
$\mathcal{C}(h,12,16)$ holds for every $h \in \{1,\dots,8\}$ (by \texttt{decide},
$256$ evaluations with totient arguments below $37$); consequently $S \ne r$ for every
$r\in\mathbb{Q}$ with $1\le h\le 8$ and $r.\mathrm{den} \mid 2^{12}(2^h-1)$.
&
\scale{fixed}
&
\lean{certifiedKill_all_small}{TotientTailPeriodKiller.lean:404}
\lean{totient_series_ne_rat_of_den_dvd}{TotientTailPeriodKiller.lean:416}
\\
Fixed-window example, depth 9, wider net
&
$\mathcal{C}(h,14,9)$ holds for every
$h\in\{1,\dots,16\}$: $S \ne r$ for every $r\in\mathbb{Q}$ with $1\le h\le16$
and $r.\mathrm{den}\mid 2^{14}(2^h-1)$.
&
\scale{fixed}
&
\lean{totient_series_ne_rat_of_den_dvd_pow_two_mul_mersenne_upto_sixteen}{CertificateKernel.lean:18890}
\\
Diagonal-pincer certificates, base 12 scales
&
$\mathcal C(H_t,H_t,L_t)$, $H_t := \mathrm{lcm}(1,\dots,t)$,
holds (by explicit \texttt{decide}/\texttt{norm_num} on checked
\texttt{Nat.totient} values, via factored prime-power blocks with
Lucas-primality certificates) for $t \in \{1,2,3,4,5,7,8,9,11,13,16,17\}$
with certificate depths $L_t \in \{6,5,7,7,9,14,15,14,21,22,23,26\}$
respectively (in the same order).
&
\scale{fixed}
&
\lean{certifiedKill_diagonal_all_imported}{DiagonalPincerCertificates.lean:2870}
\lean{diagonalPincerCertificateScales}{DiagonalPincerCertificates.lean:2852}
\lean{diagonalPincerKillDepth}{DiagonalPincerCertificates.lean:2854}
\\
Diagonal-pincer certificates, extended scales through $t=64$
&
The same predicate $\mathcal{C}(H_t,H_t,L_t)$ is additionally checked, one
sibling module per scale, at
$t \in \{19,\allowbreak23,\allowbreak25,\allowbreak27,\allowbreak29,\allowbreak31,
\allowbreak32,\allowbreak37,\allowbreak41,\allowbreak43,\allowbreak47,\allowbreak49,
\allowbreak53,\allowbreak59,\allowbreak61,\allowbreak64\}$
(16 further explicit values). Together with the base 12 scales above this is
\textbf{28 explicit values of $t$ in total}, the complete list being
$\{1,\allowbreak2,\allowbreak3,\allowbreak4,\allowbreak5,\allowbreak7,
\allowbreak8,\allowbreak9,\allowbreak11,\allowbreak13,\allowbreak16,\allowbreak17,
\allowbreak19,\allowbreak23,\allowbreak25,\allowbreak27,\allowbreak29,\allowbreak31,
\allowbreak32,\allowbreak37,\allowbreak41,\allowbreak43,\allowbreak47,\allowbreak49,
\allowbreak53,\allowbreak59,\allowbreak61,\allowbreak64\}$.
This does \emph{not} establish $\mathcal{C}(H_t,H_t,\cdot)$ for infinitely
many $t$; these 28 examples are a historical strict subset of the contiguous
band in the next row.
&
\scale{fixed}
&
Per-scale Lean modules; endpoint
\lean{certifiedKill_diagonal_t64}{DiagonalPincerCertificateT64Endpoint.lean:1928}
\\
Contiguous lcm-diagonal band through $t\le82$
&
For every natural $t\le82$ there exists a depth $L$ with
$\mathcal{C}(H_t,H_t,L)$.  This closes every scale in the finite interval,
including plateau transfers, with no holes.  The next lcm jump is at the prime
$83$; no certificate at $t=83$ is claimed, and any such certificate must have
depth at least $125$.
&
\scale{bounded}
&
\lean{ErdosProblems.Skip.LadderT67.exists_diagonalKill_le_82}{ErdosProblems/Skip/LadderT67.lean:71264}
\lean{ErdosProblems.Skip.LadderT67.t83_depth_floor}{ErdosProblems/Skip/LadderT67.lean:71294}
\\
Farey denominator floor
&
Every $q\in\mathbb{N}$ with \adjustbox{max width=\linewidth}{$0<q\le Q_0=79\,639\,646\,646\,701\,375\,323\,355\,774\,875\,831\,053$}
satisfies the window-$(N,K)=(1,240)$ gap certificate; $q=Q_0+1$ is the exact
first failure.
&
\scale{bounded}
&
\lean{gap_check_window_1_240_le_79639646646701375323355774875831053}{GapFareyBound.lean:176}
\\
Coprimality-probability / M\"obius-square Farey floor
&
Every $d\in\mathbb{N}$ with \adjustbox{max width=\linewidth}{$0<d\le Q_1=39\,819\,823\,323\,350\,687\,661\,677\,887\,437\,915\,526$}
(exactly $\lfloor Q_0/2\rfloor$) is excluded as a denominator of
$\Pr(\gcd(X,Y)=1)$ and, separately, of $T=\sum_d\mu(d)/(2^d-1)^2$.
&
\scale{bounded}
&
\lean{tsum_visible_coprime_pairs_ne_int_div_of_den_le_39819823323350687661677887437915526}{CertificateKernel.lean:18572}
\lean{tsum_moebius_div_two_pow_sub_one_sq_ne_int_div_of_den_le_39819823323350687661677887437915526}{CertificateKernel.lean:18487}
\\
\end{longtable}
\normalsize

\begin{rem}[Scope of the table]
No row above is cofinal in its indexing parameter ($h$, $t$, or the
denominator bound), and no row is claimed to extrapolate. The historical
diagonal-pincer depths are irregular and nonmonotone
($6,5,7,7,9,14,15,14,21,22,23,26,\dots$); no closed-form growth rate for
$L_t$ is proved or conjectured in the record. This table records the
historical 28-scale bank through $t=64$; the later aggregate theorem for every
$t\le82$ is stated above and is not itemised row by row here.  Neither bounded
record supplies the cofinal obligation stated in Definition~\ref{defn:sep}.
\end{rem}
% ; - part a249_p1a ; -
\section{Detailed statements}

The following statements collect the dependencies used in the earlier
arguments. We begin with unconditional identities, rank bounds and finite
examples. The next subsection supplies the definitions and elementary
implications, followed by the quantified conditions that would be needed
to prove irrationality. A source link supports the indicated statement,
not an unproved hypothesis appearing in it.

\subsection{Unconditional structure and finite examples}

The distinction between a finite example and a family at arbitrarily
large scales is important here. The rank and denominator statements
below hold with their displayed quantifiers; the residue computations
establish only the explicitly stated finite ranges.

\begin{thm}[Rationality gives an eventual tail period]\label{catalogue:cert:a9}
If $S=a/(2^cv)$ with $a\in\mathbb Z$, $c\in\mathbb N$ and $v$ a
positive odd integer, then
\[
 R_{N+h}-R_N\in\mathbb Z\qquad(N\ge c),\qquad h=\varphi(v)>0.
\]
Euler's theorem gives $v\mid2^h-1$, and the prefix-tail identity then
makes $2^N(2^h-1)S$ integral. Thus rationality supplies a positive
shift for which all sufficiently late tail differences are integral.
\par\noindent\ev{Lean} \scale{uniform}
\lean{eventual\_period\_of\_not\_irrational}{TotientTailPeriodKiller.lean:358}
\end{thm}
\leannote{Lean: \lproof{ErdosProblems/Erdos249/PaperCompleteR20/SpecifiedTailPeriod.lean}{44}{specified_euler_tail_period}.}

\begin{thm}[Rationality forces unbounded carry rank]\label{catalogue:cert:d5}
If $S\in\mathbb Q$, there are an integer $v>0$ and an integer sequence
$u$ such that
\[
 u(N+1)=2u(N)-v\varphi(N+1),\qquad u(N)/2^N\longrightarrow0,
\]
and, for every $e\ge0$, the sequences
$ n\mapsto u(2^jn+r)$ with $1\le j\le e$ and $0\le r<2^j$
span a space of $\mathbb Q$-dimension at least $2^e-1$.

The recurrence is the one obtained from the scaled tails after clearing
a hypothetical denominator of $S$. The formal term \emph{tempered}
means exactly $u(N)/2^N\to0$ here, not a separate assumption of
subexponential growth. This rank lower bound is a
necessary consequence of rationality, not a contradiction. An upper
bound using additional properties of the actual totient coefficients
would be needed; the later rational comparison rules out such a bound
for arbitrary rational coefficient series.
\par\noindent\ev{Lean} \scale{uniform}
\lean{not\_irrational\_totientSeries\_implies\_unbounded\_carryRank\_unconditional}{TotientCarryKernelRigidity.lean:300}
\end{thm}

\begin{thm}[Independence of the retained dyadic sections]\label{catalogue:cert:d4}
For every integer $e\ge0$, the family
\[
 \varphi(n),\quad\varphi(2n),\quad
 \varphi(2^jn+r)\quad(1\le j\le e,\ 0<r<2^j,\ r\text{ odd})
\]
is linearly independent over $\mathbb Q$ and has $2^e+1$ members.
At $e=0$ this family still contains $\varphi(2n)$; it is not the full
level-zero truncation, which consists only of $\varphi(n)$. For $e\ge1$
it is the basis of the full truncated dyadic family.

The proof separates one section at a time: the Chinese remainder theorem
forces suitable prime divisors in the other affine arguments, while
Dirichlet's theorem makes the selected argument prime. This proves
independence, not just the displayed cardinality. In particular, the
span of all dyadic sections is infinite-dimensional; that conclusion
alone says nothing about rationality of their weighted sum.
\par\noindent\ev{Lean} \scale{uniform}
\lean{card\_totientCanonicalIndex}{TotientMahlerDefect.lean:61}
\lean{linearIndependent\_canonicalTotientKernelFamily}{TotientMahlerDefect.lean:935}
\end{thm}
\leannote{Lean: \lproof{ErdosProblems/Erdos249/PaperCompleteR8/FullKernelAssemblies.lean}{177}{displayed_canonical_and_full_dyadic}.}

\begin{prop}[Denominator divisibility forces tail integrality]\label{catalogue:cert:a8}
If $S=a/q$ in lowest terms with $q>0$ and
$q\mid2^N(2^h-1)$, then $R_{N+h}-R_N\in\mathbb Z$.
Indeed, the prefix identity expresses that difference as
$2^N(2^h-1)S$ minus an integer. This is the contradiction used by a
nonintegrality certificate to exclude the rational value $a/q$.
\par\noindent\ev{Lean} \scale{uniform} \coord{binary-digit} \lean{tail\_diff\_int\_of\_den\_dvd}{TotientTailPeriodKiller.lean:327}
\end{prop}

The denominator results use the following polynomial, defined for $r\ge1$:
\[
 P_r(X)=\sum_{d\mid r}\mu(d)\frac rd
                 \sum_{j=0}^{r/d-1}X^{dj}.
\]
For example, $P_p(X)=p(1+X+\cdots+X^{p-1})-1$ when $p$ is prime.
Thus $P_p(2)=p(2^p-1)-1$ is coprime to $2^p-1$.
For general $r\ge1$ and $H\ge1$, write
\[
 b_r=\frac{P_r(2)}{2^r-1},\qquad
 \beta_H=\frac{b_{\operatorname{rad}(H)}}{\operatorname{rad}(H)},
 \qquad r_t=\operatorname{rad}(H_t).
\]
Here $\operatorname{rad}(H)$ is the product of the distinct prime
factors of $H$, and $\operatorname{den}(x)$ denotes the positive
reduced denominator of $x\in\mathbb Q$. These rational numbers are
finite rational contributions to the M\"obius decomposition; no convergence
to $S$ is asserted for them.
Also put
$\mathcal P_t=\{p\text{ prime}:t/2<p\le t\}$.
The hypotheses below specify which factors of $2^{r_t}-1$ remain in
the reduced denominator after cancellation.

\begin{prop}[A cyclotomic factor remains after cancellation]\label{catalogue:mob:b5}
If $r$ is squarefree, then
\[
 \gcd\bigl(|P_r(2)|,|\Phi_r(2)|\bigr)=1,
 \qquad |\Phi_r(2)|\mid\operatorname{den}(b_r).
\]
Here $\Phi_r$ is the $r$-th cyclotomic polynomial. Proposition~\ref{catalogue:mob:b4}
gives $P_r\equiv\mu(r)\pmod{\Phi_r}$, and $\mu(r)=\pm1$ because
$r$ is squarefree. This proves that the indicated cyclotomic factor
cannot cancel from $b_r$.
\par\noindent\ev{Lean} \scale{uniform} \coord{cyclotomic} \lean{mobiusNumerator\_gcd\_cyclotomicValue}{CyclotomicProjectionOfShadow.lean:350} \lean{cyclotomicValue\_dvd\_baseMobiusShadow\_den}{CyclotomicProjectionOfShadow.lean:389}
\end{prop}
\leannote{Lean: \lproof{Erdos249257/RepunitMobiusNumerator.lean}{445}{mobiusNumeratorPolynomial_eval_two}, \lproof{Erdos249257/CyclotomicProjectionOfShadow.lean}{350}{mobiusNumerator_gcd_cyclotomicValue}, \lproof{Erdos249257/CyclotomicProjectionOfShadow.lean}{389}{cyclotomicValue_dvd_baseMobiusShadow_den}.}

\begin{prop}[Primes in the upper half remain after cancellation]\label{catalogue:mob:b6}
For $t\ge5$ the whole product
\[
 \prod_{p\in\mathcal P_t}(2^p-1)
       \ \mid\ \operatorname{den}(H_t\beta_{H_t})
\]
remains after cancellation. More generally, the source gives the same
product-divisibility conclusion for a multiplier coprime to that product;
it also supplies a sufficient condition on the multiplier's prime factors.
For the multiplier $H_t/r_t$, those prime factors are at most $t$, and
the LCM result above follows.

\par\noindent\ev{Lean} \scale{uniform} \coord{cyclotomic} \lean{upperHalfPrimes}{MersenneShadowCyclotomicNoncollapse.lean:914} \lean{upperHalfChannel\_product\_dvd\_den\_of\_scale\_primeFactors\_le}{MersenneShadowCyclotomicNoncollapse.lean:809} \lean{lcmHeight\_upperHalf\_product\_dvd\_den}{MersenneShadowCyclotomicNoncollapse.lean:983}
\end{prop}
\leannote{Lean: \lproof{Erdos249257/MersenneShadowCyclotomicNoncollapse.lean}{983}{lcmHeight_upperHalf_product_dvd_den}, \lproof{Erdos249257/MersenneShadowCyclotomicNoncollapse.lean}{795}{upperHalfChannel_product_dvd_den_of_coprime_scale}, \lproof{Erdos249257/MersenneShadowCyclotomicNoncollapse.lean}{809}{upperHalfChannel_product_dvd_den_of_scale_primeFactors_le}.}

The restriction $t\ge5$ belongs to the sufficient prime-factor condition,
not just to the LCM example. At $t=3$, one has $r_t=6$,
$b_6=298/63$ and $\mathcal P_3=\{2,3\}$. The multiplier $9$ has only
prime factors at most $3$, but $\operatorname{den}(9b_6)=7$ is not
divisible by $(2^2-1)(2^3-1)=21$. Thus the prime-factor condition alone
does not imply survival of the product at every scale.
This is a denominator statement about a finite rational contribution,
not a nonintegrality assertion about a real tail.


\begin{prop}[A lower bound for the reduced denominator]\label{catalogue:mob:b7a}
For every integer $t\ge5$,
\[
 2^{\lfloor t/2\rfloor}
 \le \prod_{p\in\mathcal P_t}(2^p-1)
 \le \operatorname{den}(H_t\beta_{H_t}).
\]
Bertrand's postulate supplies a prime in $\mathcal P_t$; its factor is
at least $2^{\lfloor t/2\rfloor}$. Proposition~\ref{catalogue:mob:b6}
then gives the second inequality. The complementary term in the real
tail decomposition can still cancel this rational contribution, so
denominator growth alone does not prove irrationality of $S$.
\par\noindent\ev{Lean} \scale{uniform} \coord{mobius-mersenne} \lean{upperHalfMersenneProduct\_lower\_bound}{MersenneShadowDenominatorGrowth.lean:60}
\end{prop}
\leannote{Lean: \lproof{ErdosProblems/Erdos249/PaperCompleteR20/DenominatorBounds.lean}{10}{upper_half_product_denominator_bounds}.}

\begin{prop}[The exact reduced denominator]\label{catalogue:mob:b7b}
At every integer scale $t\ge0$,
\[
 \operatorname{den}(H_t\beta_{H_t})
 =\frac{2^{r_t}-1}
 {\gcd\!\left(2^{r_t}-1,
       \dfrac{H_t}{r_t}
       \prod_{\substack{p\mid r_t\\p\ \mathrm{odd\ prime}}}(p^2-1)\right)}.
\]
The product is $1$ when there is no odd prime divisor. This gives the
exact cancellation, including the multiplier $H_t/r_t$, rather than
merely a lower bound for the denominator. For example, $H_5=60$ and
$r_5=30$, and the denominator is $(2^{30}-1)/3$.
\par\noindent\ev{Lean} \scale{uniform} \coord{mobius-mersenne} \lean{lcmHeight\_scaledMobiusShadow\_den\_exact}{MersenneShadowDenominatorGrowth.lean:147}
\end{prop}
\leannote{Lean: \lproof{Erdos249257/MersenneShadowDenominatorGrowth.lean}{147}{lcmHeight_scaledMobiusShadow_den_exact}, \lproof{Erdos249257/MersenneShadowDenominatorGrowth.lean}{187}{lcmHeight_five_scaledMobiusShadow_den_exact}.}

\begin{prop}[Nonvanishing of a signed dyadic sum]\label{catalogue:mob:d3}
Let $I$ be finite, let $u_i\in\mathbb Z$ and $e_i\in\mathbb N$, and
suppose that $m\in I$ is the unique index with maximal exponent $e_m$.
If $u_m$ is odd, then
\[
 2^{e_m}\sum_{i\in I}\frac{u_i}{2^{e_i}}
       =\sum_{i\in I}u_i2^{e_m-e_i}\equiv1\pmod2.
\]
Indeed, every summand other than $u_m$ is even. The sum is therefore
nonzero. In the determinant application, the rectangular
Cauchy--Binet formula produces such finite signed sums from truncated
moments. The result tests a specified configuration; it does not show
that a suitable maximal-exponent term exists at arbitrarily large scales.
\par\noindent\ev{Lean} \scale{bounded} \coord{p-adic} \lean{scaled\_dyadic\_sum\_ne\_zero}{SignedQMomentObstruction.lean:96} \lean{det\_mul\_rectangular}{SignedQMomentObstruction.lean:29}
\end{prop}
\leannote{Lean: \lproof{ErdosProblems/Erdos249/PaperCompleteR20/SignedDyadicClearing.lean}{8}{signed_dyadic_clearing}, \lproof{ErdosProblems/Erdos249/PaperCompleteR20/SignedDyadicClearing.lean}{28}{signed_dyadic_sum_ne_zero}, \lproof{Erdos249257/SignedQMomentObstruction.lean}{78}{scaled_dyadic_sum_odd}.}

\begin{prop}[Small certificates]\label{catalogue:cert:a11}
For each integer $1\le h\le8$, the finite test $\mathcal C(h,12,16)$
holds. Each discrepancy uses two 16-term windows. Across all eight
shifts, only the 24 distinct totient values at $13\le n\le36$ are
needed, since the windows overlap. The source verifies the eight
integer residue inequalities by exact computation.

Consequently, if $a/b$ is a reduced fraction with $b>0$ and
$b\mid2^{12}(2^h-1)$ for at least one $1\le h\le8$, then $S\ne a/b$.
This follows from Propositions~\ref{catalogue:cert:a6}
and~\ref{catalogue:cert:a8}.
\par\noindent\ev{Cert} \scale{fixed} \coord{other:binary-window} \lean{certifiedKill\_all\_small}{TotientTailPeriodKiller.lean:404} \lean{totient\_series\_ne\_rat\_of\_den\_dvd}{TotientTailPeriodKiller.lean:416}
\end{prop}
\leannote{Lean: \lproof{ErdosProblems/Erdos249/PaperCompleteR20/FiniteCertificateBatch.lean}{33}{small_certificate_windows}, \lproof{ErdosProblems/Erdos249/PaperCompleteR20/FiniteCertificateBatch.lean}{45}{small_certificates_and_exclusions}.}

\begin{prop}[A common certificate for sixteen shifts]\label{catalogue:cert:a12}
For each integer $1\le h\le16$, one has $\mathcal C(h,14,9)$.
The common basepoint is $14$ and the depth is $9$; these parameters
have different roles. The two windows for each shift use only the
25 distinct totient values at $15\le n\le39$ across the whole family.
Thus $S\ne a/b$ for every reduced fraction with $b>0$ such that
$b\mid2^{14}(2^h-1)$ for at least one $1\le h\le16$.
\par\noindent\ev{Cert} \scale{fixed} \coord{other:binary-window} \lean{totient\_series\_ne\_rat\_of\_den\_dvd\_pow\_two\_mul\_mersenne\_upto\_sixteen}{CertificateKernel.lean:18890}
\end{prop}
\leannote{Lean: \lproof{ErdosProblems/Erdos249/PaperCompleteR20/FiniteCertificateBatch.lean}{39}{sixteen_certificate_windows}, \lproof{ErdosProblems/Erdos249/PaperCompleteR20/FiniteCertificateBatch.lean}{51}{sixteen_certificates_and_exclusions}.}

\begin{prop}[Historical diagonal examples and the complete band through 82]\label{catalogue:cert:b11}
Let $H_t=\operatorname{lcm}(1,\ldots,t)$ and let
\[
P(t)\quad:\Longleftrightarrow\quad \exists L,\ \mathcal{C}(H_t,H_t,L).
\]
The checked certificate table proves $P(t)$ at 28 explicit indices, ending at $t=64$.
For the initial indices
\[
t=1,2,3,4,5,7,8,9,11,13,16,17,\ldots,
\]
the corresponding depths begin
$6,5,7,7,9,14,15,14,21,22,23,26,\ldots$.
Each entry is a finite kernel computation on explicit totient values. The factorisations use
checked prime-power blocks and Lucas primality certificates. Thus the table proves a finite list
of instances of $P(t)$; it does not prove that $P(t)$ holds for infinitely many $t$.
\par\noindent\ev{Cert} \scale{fixed} \coord{other:lcm-diagonal} \lean{certifiedKill\_diagonal\_all\_imported}{DiagonalPincerCertificates.lean:2870} \lean{diagonalPincerKillDepth}{DiagonalPincerCertificates.lean:2854} \lean{certifiedKill\_diagonal\_t64}{DiagonalPincerCertificateT64Endpoint.lean:1928}

The historical 28 examples are now a strict subset of the aggregate theorem
$\forall t\le82,\ P(t)$.  That theorem closes the finite interval without
holes but supplies neither $P(83)$ nor a cofinal family.
\par\noindent\ev{Cert} \scale{bounded} \coord{other:lcm-diagonal}
\lean{ErdosProblems.Skip.LadderT67.exists_diagonalKill_le_82}{ErdosProblems/Skip/LadderT67.lean:71264}
\end{prop}
\leannote{Lean: \lproof{ErdosProblems/Erdos249/PaperCompleteR20/FiniteCertificateBatch.lean}{57}{historical_table_size_and_initial_depths}, \lproof{ErdosProblems/Erdos249/PaperCompleteR20/FiniteCertificateBatch.lean}{64}{historical_table_and_complete_band}.}

\subsection{Definitions and elementary identities}

The finite test rests on two identities: the prefix-tail decomposition
and the truncation of a tail difference. We state those first, then give
the soundness and completeness arguments, the LCM specialisation and the
corresponding rational-approximation criteria. The subsequent M\"obius
identities describe the same series with different coefficient weights.

\subsubsection{Initial implications}

\begin{defn}[The scaled tail]\label{catalogue:cert:a1}
For $N\in\mathbb N$,
\[
 R_N=\sum_{j\ge0}\frac{\varphi(N+1+j)}{2^{j+1}}.
\]
The series converges because $\varphi(n)\le n$. It is the tail of $S$
multiplied by $2^N$, not the fractional part of $2^NS$; indeed, $R_N$
need not lie in $[0,1)$.
\par\noindent\ev{Lean}
\lean{totientTail}{TotientTailPeriodKiller.lean:57}
\end{defn}

\begin{defn}[The truncated tail difference]\label{catalogue:cert:a3}
For $h,N,L\in\N$, define the integer
\[
 D(h,N,L)=\sum_{j=0}^{L-1}
 \bigl(\varphi(N+h+1+j)-\varphi(N+1+j)\bigr)2^{L-1-j}.
\]
The sum is empty when $L=0$. Dividing it by $2^L$ gives the first
$L$ terms of $R_{N+h}-R_N$. Computing the finitely many totients
therefore determines $D(h,N,L)$ exactly; the residue inequality below
is the decidable condition.
\par\noindent\ev{Lean} \scale{n/a} \coord{other:binary-window} \lean{windowDiscrepancy}{TotientTailPeriodKiller.lean:66}
\end{defn}

\begin{defn}[The central-residue inequality]\label{catalogue:cert:a4}
For $h,N,L\in\mathbb N$, the finite condition $\mathcal C(h,N,L)$ is
\[
 N+h+L+2<D(h,N,L)\bmod2^L<2^L-(N+h+L+2).
\]
It puts the residue outside both endpoint intervals of width $N+h+L+2$.
After division by $2^L$, the excluded width tends to zero as $L$ grows
with $h,N$ fixed. No analytic hypothesis is part of this finite test;
the separate condition $\mathrm{Sep}$ quantifies over such tests.
\par\noindent\ev{Lean}
\lean{certifiedKill}{TotientTailPeriodKiller.lean:72}
\end{defn}

\begin{prop}[The prefix-tail identity]\label{catalogue:cert:a2}
For every $N\in\mathbb N$,
\[
 2^NS=\Phi_N+R_N,
 \qquad \Phi_N=\sum_{n\le N}\varphi(n)2^{N-n}\in\mathbb Z.
\]
Thus $R_N$ and $2^NS$ have the same fractional part. Subtracting the
identities at $N+h$ and $N$ gives the tail-difference congruence used in
the irrationality criterion.
\par\noindent\ev{Lean} \scale{uniform}
\lean{two\_pow\_mul\_totient\_series\_eq}{TotientTailPeriodKiller.lean:150}
\end{prop}
\leannote{Lean: \lproof{ErdosProblems/Erdos249/PaperCompleteR20/TailDepthCorrespondence.lean}{8}{prefix_fractional_part}, \lproof{Erdos249257/TotientTailPeriodKiller.lean}{150}{two_pow_mul_totient_series_eq}, \lproof{Erdos249257/LcmConeFlatness.lean}{327}{tail_diff_mem_int_iff_scaled_series_mem_int}.}

\begin{lem}[The necessary depth inequality]\label{catalogue:cert:a5}
For $h,N,L\in\N$, the condition $\mathcal C(h,N,L)$ implies
\[
 2(N+h+L+2)<2^L.
\]
Indeed, its lower residue bound must be smaller than its upper bound.
Thus $L>1+\log_2(N+h+L+2)$ whenever a certificate holds. In
particular, a fixed depth cannot accommodate unbounded $N+h$.
This is a necessary depth bound, not an upper bound for finding a
certificate.
\par\noindent\ev{Lean} \scale{uniform} \coord{other:binary-window} \lean{certifiedKill\_depth\_floor}{TotientTailPeriodKiller.lean:79}
\end{lem}
\leannote{Lean: \lproof{ErdosProblems/Erdos249/PaperCompleteR20/TailDepthCorrespondence.lean}{20}{certificate_logarithmic_depth}, \lproof{ErdosProblems/Erdos249/PaperCompleteR20/TailDepthCorrespondence.lean}{30}{fixed_depth_bounds_indices}, \lproof{Erdos249257/TotientTailPeriodKiller.lean}{79}{certifiedKill_depth_floor}.}

\begin{prop}[A certificate implies nonintegrality]\label{catalogue:cert:a6}
For all $h,N,L\in\mathbb N$,
\[
 \mathcal C(h,N,L)\ \Longrightarrow\ R_{N+h}-R_N\notin\mathbb Z.
\]
The scaled truncation error satisfies
$|2^L(R_{N+h}-R_N)-D(h,N,L)|\le N+h+L+2$.
If the tail difference
were integral, $D(h,N,L)$ would therefore lie within that distance of a
multiple of $2^L$, contrary to the two strict residue inequalities.
\par\noindent\ev{Lean} \scale{uniform}
\lean{tail\_diff\_notMem\_int\_of\_certifiedKill}{TotientTailPeriodKiller.lean:262}
\end{prop}
\leannote{Lean: \lproof{ErdosProblems/Erdos249/PaperCompleteR20/TailDepthCorrespondence.lean}{13}{totient_scaled_truncation_error}, \lproof{Erdos249257/TotientTailPeriodKiller.lean}{262}{tail_diff_notMem_int_of_certifiedKill}.}

\begin{thm}[Nonintegrality gives a certificate at some depth]\label{catalogue:cert:a7}
For every $h,N\in\mathbb N$,
\[
 (\exists L\in\mathbb N,\ \mathcal C(h,N,L))
 \quad\Longleftrightarrow\quad R_{N+h}-R_N\notin\mathbb Z.
\]
For the converse to Proposition~\ref{catalogue:cert:a6}, fix a positive
distance from the nonintegral tail difference to the nearest integer.
The normalised error $(N+h+L+2)/2^L$ tends to zero, so a sufficiently
large truncation preserves that distance. This is a pointwise
equivalence; irrationality requires the stated quantifiers over the
shift and basepoint.
\par\noindent\ev{Lean} \scale{uniform}
\lean{exists\_certifiedKill\_iff\_tail\_diff\_notMem\_int}{LcmConeFlatness.lean:316}
\end{thm}

\begin{defn}[The least common multiple of the first integers]\label{catalogue:cert:b1}
Put $H(0)=1$ and $H(t+1)=\operatorname{lcm}(H(t),t+1)$. Then
$H(t)=\operatorname{lcm}(1,\ldots,t)$, $t\le H(t)$, and every positive
integer $h\le t$ divides $H(t)$. This is why a single LCM parameter can
absorb any fixed hypothetical period once $t$ is large enough.
\par\noindent\ev{Lean} \scale{uniform}
\lean{periodLcm}{CarrySurvivorExtinction.lean:515}
\lean{le\_periodLcm}{LcmDiagonalReduction.lean:81}
\end{defn}

\begin{prop}[The Farey gap lemma]\label{catalogue:cert:c1}
Let $a,b,c,d,r,s\in\mathbb Z$, with $b,d>0$ and $bc-ad=1$. If
\[
 as<rb\quad\hbox{and}\quad rd<cs,
\]
then $b+d\le s$. The two cross-products $rb-as$ and $cs-rd$ are positive
integers, and
\[
 s=d(rb-as)+b(cs-rd)\ge d+b.
\]
In particular, a rational $r/s$ strictly between the Farey neighbours
$a/b$ and $c/d$, with positive denominator, has denominator at least
$b+d$.
\par\noindent\ev{Lean} \scale{uniform}
\lean{farey\_gap}{GapFareyBound.lean:51}
\end{prop}

\begin{prop}[An irrationality criterion from rational approximations]\label{catalogue:cert:d1}
Let $x\in\mathbb R$ and let $u_k=p_k/q_k\in\mathbb Q$ be in lowest
terms, with $q_k>0$. If $u_k\ne x$ for every sufficiently large $k$ and
$q_k|x-u_k|\to0$, then $x$ is irrational. Indeed, if $x=a/b$ with
$a\in\mathbb Z$ and $b\ge1$, a nonzero integer numerator gives
$q_k|x-u_k|=|a q_k-b p_k|/b\ge1/b$, a contradiction.
This elementary criterion requires both nonvanishing and the scaled
error estimate. Convergence $u_k\to x$ alone is insufficient, and the
approximants need not be continued-fraction convergents.
\par\noindent\ev{Lean} \scale{n/a} \coord{n/a} \lean{irrational\_of\_den\_mul\_abs\_sub\_tendsto\_zero}{CertificateKernel.lean:5371}
\end{prop}

\begin{prop}[An irrationality criterion from near integers]\label{catalogue:cert:d2}
Let $\xi\in\mathbb R$. Suppose that for every integer $q\ge1$ there
are integers $m,z$ with
\[
 0<|m\xi-z|<1/q.
\]
Then $\xi$ is irrational: if $\xi=a/b$ with $b\ge1$, every nonzero
such difference is at least $1/b$. Restricting the multiplier to powers
$m=b_0^n$ of a fixed integer base $b_0\ge2$ gives a sufficient special
case, not a hypothesis satisfied by every irrational number.
The strict lower bound excludes exact integer hits; the upper bound must
be available for arbitrarily large $q$.

For a binary example, form a number by concatenating the blocks $10$ at
square indices $k\ge1$ and $01$ at the other indices. Its binary expansion
is not eventually periodic, since the block sequence has increasingly
long gaps between the square indices, so the number is irrational. There
are no three consecutive equal digits. Every fractional part after a
binary shift therefore lies in $[1/8,7/8]$, and the multiples $2^n\xi$
do not approach the integers. This verifies that the restriction to
base powers is genuinely stronger. This example is an ordinary mathematical
argument, not an additional claim about the linked formalisation.
\par\noindent\ev{Lean} \scale{n/a} \coord{n/a} \lean{irrational\_of\_int\_mul\_near\_int}{CertificateKernel.lean:6120} \lean{irrational\_of\_pow\_mul\_near\_int}{CertificateKernel.lean:6149}
\end{prop}
\leannote{Lean: \lproof{Erdos249257/CertificateKernel.lean}{6120}{irrational_of_int_mul_near_int}, \lproof{Erdos249257/CertificateKernel.lean}{6149}{irrational_of_pow_mul_near_int}, \lproof{ErdosProblems/Erdos249/PaperCompleteR21/NearIntegerIrrationalityCriterion.lean}{16}{irrational_of_near_integer_multiples}, \lproof{ErdosProblems/Erdos249/PaperCompleteR21/NearIntegerIrrationalityCriterion.lean}{25}{irrational_of_near_integer_base_powers}, and 16 further declarations in the \hyperref[sec:coverage]{coverage section}.}

\begin{prop}[A nonzero evaluation minor gives independence]\label{catalogue:cert:d3}
Let $I$ be a finite index set and let $f_j:\mathbb N\to\mathbb Q$
for $j\in I$. If there are evaluation points $n_i\in\mathbb N$ such that
\[
 \det\bigl(f_j(n_i)\bigr)_{i,j\in I}\ne0,
\]
then the family $(f_j)_{j\in I}$ is linearly independent over $\mathbb Q$.
A relation among the functions, evaluated at the $n_i$, gives a vector
in the kernel of this nonsingular matrix, so all coefficients vanish.
Finiteness is required to form the displayed determinant. In the totient
application the arithmetic work is to construct these evaluation points;
the implication itself is ordinary linear algebra.
\par\noindent\ev{Lean} \scale{n/a} \coord{n/a} \lean{SeparatedMinorCertificate}{TotientMahlerDefect.lean:83} \lean{linearIndependent\_of\_separatedMinorCertificate}{TotientMahlerDefect.lean:91}
\end{prop}

\begin{prop}[The gap between distinct rational numbers]\label{catalogue:cert:d9}
If $a/b<c/d$ are reduced rational numbers with $b,d>0$, then
\[
 \frac cd-\frac ab=\frac{bc-ad}{bd}\ge\frac1{bd}.
\]
The numerator $bc-ad$ is a positive integer. Thus an upper bound on a
positive rational error gives a lower bound on the product of the reduced
denominators. A large denominator in an unreduced displayed expression
need not be the reduced denominator appearing in this conclusion.
\par\noindent\ev{Lean} \scale{n/a} \coord{n/a} \lean{positive\_rational\_difference\_lower\_bound}{PrimitiveRationalGapSupply.lean:31}
\end{prop}
\leannote{Lean: \lproof{ErdosProblems/Erdos249/PaperCompleteR20/RationalSpacingCorrespondence.lean}{6}{rational_difference_exact}, \lproof{ErdosProblems/Erdos249/PaperCompleteR20/RationalSpacingCorrespondence.lean}{17}{rational_cross_numerator_positive}, \lproof{ErdosProblems/Erdos249/PaperCompleteR20/RationalSpacingCorrespondence.lean}{27}{rational_error_denominator_bound}, \lproof{Erdos249257/PrimitiveRationalGapSupply.lean}{31}{positive_rational_difference_lower_bound}.}

\begin{prop}[The M\"obius identity for $S$]\label{catalogue:mob:a1a}
The absolutely convergent series satisfy
\[
 S=\sum_{n\ge1}\frac{\varphi(n)}{2^n}
  =\sum_{d\ge1}\frac{\mu(d)2^d}{(2^d-1)^2}
  =\frac12+\sum_{d\ge1}\frac{\mu(d)}{(2^d-1)^2}.
\]
For the first rearrangement use
$\varphi(n)=\sum_{d\mid n}\mu(d)n/d$. For the second use
$2^d=(2^d-1)+1$ and the classical identity
$\sum_{d\ge1}\mu(d)/(2^d-1)=1/2$.
These are identities of values, not transfers of irrationality from a
Lambert series with a different weight.
\par\noindent\ev{Lean} \scale{n/a} \coord{mobius-mersenne} \lean{tsum\_totient\_half\_pow\_eq\_half\_add\_moebius\_sq}{MersenneLambertLadder.lean:665} \lean{totientSeries\_eq\_pnat\_half\_pow}{SquaredMersenneDiagonalEnclosure.lean:72} \lean{mobius\_square\_series\_eq\_partial\_add\_tail}{SquaredMersenneDiagonalEnclosure.lean:94}
\end{prop}
\leannote{Lean: \lproof{Erdos249257/MersenneLambertLadder.lean}{491}{tsum_moebius_lambert_sq}, \lproof{Erdos249257/CertificateKernel.lean}{18454}{totient_series_eq_half_add_moebius_mersenne_square}.}

\begin{cor}[An equivalent irrationality question]\label{catalogue:mob:a1b}
Subtracting the rational number $1/2$ gives
\[
 S\notin\mathbb Q
 \quad\Longleftrightarrow\quad
 \sum_{d\ge1}\frac{\mu(d)}{(2^d-1)^2}\notin\mathbb Q.
\]
This is the same irrationality question in a different series
representation. The notation $L_2(\mu)$ for the right-hand series is
introduced in Definition~\ref{catalogue:mob:a3}.
\par\noindent\ev{Lean} \scale{n/a} \coord{mobius-mersenne} \lean{tsum\_totient\_half\_pow\_eq\_half\_add\_moebius\_sq}{MersenneLambertLadder.lean:665}
\end{cor}
\leannote{Lean: \lproof{ErdosProblems/Erdos249/PaperCompleteR7/ArithmeticAssemblies.lean}{40}{irrational_totient_iff_moebius_square}.}

\begin{prop}[The squared-Lambert identity]\label{catalogue:mob:a2}
Let $w:\mathbb N\to\mathbb R$ satisfy $|w(d)|\le d$ for $d\ge1$,
and let $0\le r<1$. Then
\[
 \sum_{d\ge1}w(d)\left(\frac{r^d}{1-r^d}\right)^2
 =\sum_{n\ge1}\left(\sum_{d\mid n}w(d)\left(\frac nd-1\right)\right)r^n.
\]
Expand $(x/(1-x))^2=\sum_{j\ge1}(j-1)x^j$ and group terms by
$n=dj$. Absolute convergence follows from the bound on $w$ and
geometric decay. The cases $w=\mu$, $w=1$ and $w=\varphi$ give the next
three series; the same algebra applies to each, but their arithmetic
values differ.
\par\noindent\ev{Lean} \scale{uniform} \coord{mobius-mersenne} \lean{tsum\_lambert\_linear\_weight\_sq\_pure}{GcdMomentCalculus.lean:105}
\end{prop}

\begin{defn}[The squared-Lambert series with weight $\mu$]\label{catalogue:mob:a3}
Put
\[
 L_2(\mu)=\sum_{d\ge1}\frac{\mu(d)}{(2^d-1)^2}=S-\frac12.
\]
The equality follows from Proposition~\ref{catalogue:mob:a1a} and
Proposition~\ref{catalogue:mob:a2} with $w=\mu$.
Subtracting $1/2$ preserves rationality, so the irrationality of this
series is exactly the question about $S$, not a separate easier claim.
\par\noindent\ev{Cited} \scale{n/a} \coord{mobius-mersenne} \lean{tsum\_totient\_half\_pow\_eq\_half\_add\_moebius\_sq}{MersenneLambertLadder.lean:665}
\end{defn}

\begin{prop}[Weight one and divisor sums]\label{catalogue:mob:a4}
$\sum_{d:\mathbb{N}^+}' 1/(2^d-1)^2 = \sum_{n:\mathbb{N}^+}' \big(\sigma(n)-\tau(n)\big)\cdot(1/2)^n = \zeta_q(2) - \zeta_q(1)$ at $q=1/2$. The displayed identity is formalised. Irrationality of its value follows
from the cited linear independence result of Postelmans and Van Assche,
as explained in Proposition~\ref{prop:zetaq}; that literature result is
not formalised here. Replacing the weight $1$ by $\mu$ changes the value
to that of Definition~\ref{catalogue:mob:a3}. Observation~\ref{catalogue:mob:a6}
compares the two choices of weight.
\par\noindent\ev{Lean} \scale{n/a} \coord{mobius-mersenne} \lean{tsum\_one\_div\_mersenne\_sq\_eq\_sigma\_sub\_tau\_series}{GcdMomentCalculus.lean:216}
\end{prop}
\leannote{Lean: \lproof{ErdosProblems/Erdos249/PaperCompleteR21/SquaredMersenneDivisorIdentities.lean}{36}{divisor_sum_identity}, \lproof{ErdosProblems/Erdos249/PaperCompleteR21/QZetaAnchor.lean}{17}{qZeta}, \lproof{ErdosProblems/Erdos249/PaperCompleteR21/QZetaAnchor.lean}{81}{qZeta_half_two_sub_one}, \lproof{ErdosProblems/Erdos249/PaperCompleteR21/QZetaAnchor.lean}{143}{irrational_qZeta_half_difference_of_linearIndependent}, and 1 further declaration. Conditional on the linear independence theorem of Postelmans and Van Assche; see the \hyperref[sec:coverage]{coverage section}.}

\begin{prop}[Totient weight and gcd moments]\label{catalogue:mob:a5}
$\sum_{d:\mathbb{N}^+}' \varphi(d)/(2^d-1)^2 = \sum_{n:\mathbb{N}^+}' (P(n)-n)\cdot(1/2)^n$, where $P = \varphi * \mathrm{Id}$ (Pillai's gcd-sum function). It also equals $\mathbb E[\gcd(X,Y)]$ when $X,Y$ are independent and
$\mathbb P(X=n)=\mathbb P(Y=n)=2^{-n}$ for $n\ge1$; see
Proposition~\ref{prop:pillai}.
This is a different weighted series. Its rationality is not settled by
the identities proved here.
\par\noindent\ev{Lean} \scale{n/a} \coord{mobius-mersenne} \lean{tsum\_totient\_div\_mersenne\_sq\_eq\_gcd\_moment\_series}{GcdMomentCalculus.lean:235}
\end{prop}
\leannote{Lean: \lproof{ErdosProblems/Erdos249/PaperCompleteR21/PillaiGcdExpectation.lean}{334}{gcd_moment_identity_three_members}, \lproof{ErdosProblems/Erdos249/PaperCompleteR21/PillaiGcdExpectation.lean}{307}{tsum_pos_pair_gcd_half_eq_totient_div_mersenne_sq}, \lproof{ErdosProblems/Erdos249/PaperCompleteR21/PillaiGcdExpectation.lean}{67}{pillaiP_eq_totient_mul_id}, \lproof{ErdosProblems/Erdos249/PaperCompleteR21/PillaiGcdExpectation.lean}{77}{sum_gcd_Icc_eq_pillaiP}, and 1 further declaration in the \hyperref[sec:coverage]{coverage section}.}

\begin{obs}[Comparison of weighted series]\label{catalogue:mob:a6}
For $L(f)=\sum_{d\ge1}f(d)/(2^d-1)$, the three weights give
\[
 L(\mu)=\frac12,\qquad L(1)=E,\qquad L(\varphi)=2,
\]
where $E$ is the irrational Erd\H{o}s--Borwein constant.
For the squared denominators, they give
\[
 L_2(\mu)=S-\frac12,\qquad
 L_2(1)=\zeta_{1/2}(2)-\zeta_{1/2}(1),\qquad
 L_2(\varphi)=\mathbb E[\gcd(X,Y)].
\]
The first equality is Definition~\ref{catalogue:mob:a3}; the second and
third are Propositions~\ref{catalogue:mob:a4} and~\ref{catalogue:mob:a5}.
The middle value is irrational by the cited linear independence theorem.
The other two rationality questions are not settled in this record.
These examples show that the choice of weight matters: an identity with
the same denominators does not transfer an irrationality theorem from one
weight to another. They do not classify all weighted Lambert series.
\par\noindent\ev{Lean} \scale{n/a} \coord{mobius-mersenne}
\end{obs}

\begin{prop}[The weight of pairs divisible by a fixed integer]\label{catalogue:mob:a7}
Let $X,Y$ be independent random variables with
$\mathbb P(X=n)=\mathbb P(Y=n)=2^{-n}$ for $n\ge1$. For every $d\ge1$,
\[
 \mathbb P(d\mid X,\ d\mid Y)=\frac1{(2^d-1)^2}.
\]
This follows by multiplying the two geometric sums
$\sum_{k\ge1}2^{-dk}=1/(2^d-1)$. It explains the squared denominator
in the preceding M\"obius identity.
\par\noindent\ev{Lean} \scale{uniform} \coord{probability} \lean{tsum\_pos\_pair\_both\_dvd\_half\_eq\_inv\_mersenne\_sq}{GcdMomentCalculus.lean:266}
\end{prop}
\leannote{Lean: \lproof{ErdosProblems/Erdos249/PaperCompleteR21/SquaredMersenneDivisorIdentities.lean}{62}{pair_divisibility_mass}, \lproof{ErdosProblems/Erdos249/PaperCompleteR21/SquaredMersenneDivisorIdentities.lean}{79}{tsum_geometric_multiples}, \lproof{Erdos249257/GcdMomentCalculus.lean}{266}{tsum_pos_pair_both_dvd_half_eq_inv_mersenne_sq}.}

\begin{prop}[The sum over coprime directions]\label{catalogue:mob:a8}
The sum over positive coprime pairs satisfies
\[
 \sum_{\substack{a,b\ge1\\\gcd(a,b)=1}}
                  \frac1{2^{a+b}-1}=1.
\]
To see the normalisation, write every pair of positive integers uniquely
as $(ka,kb)$ with $\gcd(a,b)=1$, and sum $2^{-k(a+b)}$ over $k\ge1$.
The total is $(\sum_{n\ge1}2^{-n})^2=1$. This base-two normalisation
is used in Proposition~\ref{catalogue:mob:a9a}.
\par\noindent\ev{Lean} \scale{n/a} \coord{probability} \lean{tsum\_pos\_coprime\_inv\_mersenne\_eq\_one}{GcdMomentCalculus.lean:349}
\end{prop}

\begin{prop}[A Stern--Brocot recursion with stopping]\label{catalogue:mob:a9a}
For positive integers $a,b$, put
\[
 M(a,b)=\frac1{(2^a-1)(2^b-1)}.
\]
Then
\[
 M(a,b)=\frac1{2^{a+b}-1}+M(a+b,b)+M(a,a+b).
\]
Multiplication by the common denominator proves the identity. At the
root $(1,1)$, the three terms on the right are each $1/3$, and
$M(1,1)=1$.

The first term is a mass retained at the current node, not passed to
either child. Thus the probabilistic interpretation is a branching
process with stopping: after division by $M(a,b)$, the stopping and two
transition probabilities are, respectively,
\[
 \frac{(2^a-1)(2^b-1)}{2^{a+b}-1},\qquad
 \frac{2^a-1}{2^{a+b}-1},\qquad
 \frac{2^b-1}{2^{a+b}-1}.
\]
They sum to one. The stopping probability is at least $1/3$, which gives
the next proposition's geometric error bound. Also
$M(a,b)=\mathbb P(a\mid X)\mathbb P(b\mid Y)$ for the independent
geometric variables used above. The identity does not describe a
conservative mass split between the two children alone.
\par\noindent\ev{Lean} \scale{uniform} \coord{other:stern-brocot} \lean{cylinderMass\_split}{GcdMomentCalculus.lean:474}
\end{prop}
\leannote{Lean: \lproof{ErdosProblems/Erdos249/PaperCompleteR21/SternBrocotStoppingRecursion.lean}{42}{divisibility_mass}, \lproof{ErdosProblems/Erdos249/PaperCompleteR21/SternBrocotStoppingRecursion.lean}{86}{cylinderMass_eq_divisibility_mass_mul}, \lproof{ErdosProblems/Erdos249/PaperCompleteR21/SternBrocotStoppingRecursion.lean}{95}{cylinder_mediant_split}, \lproof{ErdosProblems/Erdos249/PaperCompleteR21/SternBrocotStoppingRecursion.lean}{103}{cylinder_root_values}, and 3 further declarations in the \hyperref[sec:coverage]{coverage section}.}

\begin{prop}[Convergence with an explicit error]\label{catalogue:mob:a9b}
In the splitting identity of Proposition~\ref{catalogue:mob:a9a}, the
sum of the two child terms is at most $2/3$ of the parent term.
If $M_d(a,b)$ is the sum of the contributions removed during the first
$d$ levels, then
\[
 |M(a,b)-M_d(a,b)|\le(2/3)^d M(a,b).
\]
In particular, $M_d(a,b)\to M(a,b)$. This convergence statement is
separate from identifying the limit with a sum over any independently
defined infinite set of descendants; no additional identification is
asserted here.
\par\noindent\ev{Lean} \scale{uniform} \coord{other:stern-brocot} \lean{sternBrocotDepthMass\_error}{GcdMomentCalculus.lean:525} \lean{tendsto\_sternBrocotDepthMass}{GcdMomentCalculus.lean:560}
\end{prop}
\leannote{Lean: \lproof{Erdos249257/GcdMomentCalculus.lean}{514}{cylinderMass_children_le}, \lproof{Erdos249257/GcdMomentCalculus.lean}{525}{sternBrocotDepthMass_error}, \lproof{Erdos249257/GcdMomentCalculus.lean}{560}{tendsto_sternBrocotDepthMass}.}

\begin{prop}[A numerator polynomial and its coefficients]\label{catalogue:mob:b1}
For squarefree $r\ge1$, the polynomial defined above has coefficients
\[
 [X^k]P_r(X)=
 \begin{cases}
 \dfrac r{\gcd(r,k)}\varphi(\gcd(r,k)),&0\le k<r,\\
 0,&k\ge r.
 \end{cases}
\]
For $k<r$, extracting the coefficient gives
$\sum_{d\mid\gcd(r,k)}\mu(d)r/d$; the usual divisor formula for
$\varphi$ evaluates this sum. All coefficients in the range $0\le k<r$
are positive. Thus a signed sum of geometric polynomials has an explicit
positive-coefficient expression. The squarefree hypothesis is retained
as in the cited formal statements.
\par\noindent\ev{Lean} \scale{n/a} \coord{cyclotomic} \lean{mobiusNumeratorPolynomial\_eq\_gcdWord}{RepunitMobiusNumerator.lean:205} \lean{mobiusNumeratorPolynomial\_coeff}{RepunitMobiusNumerator.lean:217} \lean{mobiusNumeratorPolynomial\_coeff\_pos}{RepunitMobiusNumerator.lean:234}
\end{prop}

\begin{prop}[Evaluation at two]\label{catalogue:mob:b2}
For squarefree $r\ge1$, evaluation of $P_r$ at two gives the integer
\[
 P_r(2)=\sum_{d\mid r}\mu(d)\frac rd
                         \frac{2^r-1}{2^d-1}.
\]
Each quotient is an integer because $d\mid r$. The formal definition
sums over subsets of the prime divisors of $r$; for squarefree $r$
these are exactly its divisors. This identifies the polynomial
calculation with the numerator used in $b_r$.
\par\noindent\ev{Lean} \scale{uniform} \coord{cyclotomic} \lean{mobiusNumeratorPolynomial\_eval\_two}{RepunitMobiusNumerator.lean:445} \lean{mobiusNumerator}{RadicalMobiusShadow.lean:101}
\end{prop}
\leannote{Lean: \lproof{Erdos249257/RepunitMobiusNumerator.lean}{445}{mobiusNumeratorPolynomial_eval_two}, \lproof{ErdosProblems/Erdos249/PaperCompleteR20/NumeratorEvaluation.lean}{32}{numerator_eval_two_divisors}.}

\begin{prop}[Decomposition by the radical]\label{catalogue:mob:b3}
For $H\ge1$ and $r\ge1$,
\[
 H\beta_H=\frac H{\operatorname{rad}(H)}b_{\operatorname{rad}(H)},
 \qquad
 \operatorname{den}(b_r)
   =\frac{2^r-1}{\gcd(|P_r(2)|,2^r-1)}.
\]
The first equality follows from the definition of $\beta_H$; the
second is ordinary reduction of an integer fraction. Multiplication
by $H/\operatorname{rad}(H)$ can cause further cancellation. No
coprimality assumption is included in either identity.
\par\noindent\ev{Lean} \scale{uniform} \coord{mobius-mersenne} \lean{scaledMobiusShadow\_eq\_radicalBase}{RadicalMobiusShadow.lean:125} \lean{baseMobiusShadow\_den}{RadicalMobiusShadow.lean:150}
\end{prop}
\leannote{Lean: \lproof{ErdosProblems/Erdos249/PaperCompleteR20/RadicalDecomposition.lean}{50}{radical_decomposition}.}

\noindent The constant in Proposition~\ref{catalogue:mob:b4} is the value of $P_r$ at a primitive
$m$-th root of unity. Since the $X^k$ coefficient of $P_r$ depends only on
$\gcd(k,r)$, the value is a discrete Fourier coefficient of a function of the
greatest common divisor, which Schramm expresses through Ramanujan
sums~\cite[Theorem and (2), p.~2]{schramm-gcd}; Proposition~\ref{catalogue:mob:b4} specialises this to the
present numerator polynomial and proves the integral cyclotomic consequence.

\begin{prop}[A cyclotomic congruence]\label{catalogue:mob:b4}
Let $r$ be squarefree and $m\mid r$. In $\mathbb Z[X]$,
\[
 P_r(X)\equiv\mu(m)J_2(r/m)\pmod{\Phi_m(X)},
\]
where $J_2(n)=n^2\prod_{p\mid n}(1-p^{-2})$ is the second Jordan
totient. Evaluation at two gives
\[
 \Phi_m(2)\mid P_r(2)-\mu(m)J_2(r/m).
\]
\par\noindent\ev{Lean} \scale{uniform} \coord{cyclotomic} \lean{cyclotomic\_dvd\_mobiusNumeratorPolynomial\_sub}{CyclotomicProjectionOfShadow.lean:299} \lean{mobiusNumerator\_mod\_cyclotomicEval}{CyclotomicProjectionOfShadow.lean:310}
\end{prop}
\leannote{Lean: \lproof{Erdos249257/CyclotomicProjectionOfShadow.lean}{299}{cyclotomic_dvd_mobiusNumeratorPolynomial_sub}, \lproof{Erdos249257/CyclotomicProjectionOfShadow.lean}{310}{mobiusNumerator_mod_cyclotomicEval}, \lproof{Erdos249257/CyclotomicProjectionOfShadow.lean}{67}{jordanTotientTwo_eq_prod_primeFactors}.}

\begin{proof}
To prove the congruence, let $\zeta$ be a primitive $m$-th root of
unity. Since $m\mid r$, the geometric sum in the definition of $P_r$
is $r/d$ when $m\mid d$ and zero otherwise. Hence
\[
 P_r(\zeta)=\sum_{\substack{d\mid r\\m\mid d}}\mu(d)(r/d)^2.
\]
Write $r=mu$. Squarefreeness gives $\gcd(m,u)=1$, so substituting
$d=me$ reduces this to
$\mu(m)\sum_{e\mid u}\mu(e)(u/e)^2=\mu(m)J_2(u)$.
The minimal polynomial $\Phi_m$ therefore divides the difference in
$\mathbb Q[X]$; because it is monic, polynomial division gives the
quotient in $\mathbb Z[X]$. Taking $m=r$ gives the coprimality in
Proposition~\ref{catalogue:mob:b5}.
\end{proof}

Squarefreeness cannot simply be dropped from this constant formula.
For $r=4$ and $m=2$,
$P_4(X)=2+4X+2X^2+4X^3$, so $P_4(-1)=-4$, whereas
$\mu(2)J_2(2)=-3$. Here $P_4(2)=50$ and $\Phi_4(2)=5$;
the latter factor cancels from $b_4=50/15=10/3$. Positive polynomial
coefficients alone therefore do not prevent denominator cancellation.


\begin{prop}[Adjoining a prime: new divisors]\label{catalogue:mob:b8a}
Let $r\ge1$, let $p\nmid r$ be prime, and let $m\mid r$. Then
\[
 P_{rp}(X)\equiv-P_r(X^p)\pmod{\Phi_{mp}(X)}.
\]
\par\noindent\ev{Lean} \scale{uniform} \coord{cyclotomic} \lean{cyclotomic\_dvd\_primeJump\_new\_fibre}{PrimePowerJumpDynamics.lean:281} \lean{cyclotomic\_dvd\_primeJump\_new\_fibre\_constant}{PrimePowerJumpDynamics.lean:347}
\end{prop}

\begin{proof}
Splitting the divisors of $rp$ into $d$ and $pd$, with $d\mid r$,
gives the exact polynomial identity
\[
 P_{rp}(X)=p\left(\sum_{j=0}^{p-1}X^{jr}\right)P_r(X)-P_r(X^p).
\]
At a primitive $mp$-th root $\zeta$, the number $\zeta^r$ is a
nontrivial $p$-th root, because $m\mid r$ and $p\nmid r$.
The geometric factor vanishes, proving the stated congruence.
No squarefree hypothesis on $r$ is needed for this identity.
If $r$ is squarefree, $\zeta^p$ is a primitive $m$-th root, and
Proposition~\ref{catalogue:mob:b4} gives the constant remainder
$-\mu(m)J_2(r/m)$, as in the second linked theorem.
\end{proof}


\begin{prop}[Adjoining a prime: existing divisors]\label{catalogue:mob:b8b}
Let $r$ be squarefree, let $p\nmid r$ be prime, and let $m\mid r$.
For the existing divisors $m$, the corresponding identity is
\[
 P_{rp}(X)\equiv(p^2-1)P_r(X)\pmod{\Phi_m(X)}.
\]
\par\noindent\ev{Lean} \scale{uniform} \coord{cyclotomic} \lean{cyclotomic\_dvd\_primeJump\_old\_fibre}{PrimePowerJumpDynamics.lean:310}
\end{prop}

\begin{proof}
Apply Proposition~\ref{catalogue:mob:b4} to both $r$ and $rp$.
Since $p\nmid r/m$, multiplicativity gives
$J_2(rp/m)=(p^2-1)J_2(r/m)$, which proves the identity.
The squarefree hypothesis is retained as in the linked theorem.
This compares the remainders at the existing divisors $m$; the
preceding proposition treats the new divisors $mp$.
\end{proof}


\begin{prop}[Lambert-series identities]\label{catalogue:cert:d7}
For an arithmetic function $f$ for which the sums converge absolutely,
write $L(f)=\sum_{n\ge1}f(n)/(2^n-1)$. Expanding each geometric series gives
\[
 L(f)=\sum_{m\ge1}\frac{(f*1)(m)}{2^m}.
\]
In particular, if $\alpha=\varphi*\mu$, then $\alpha*1=\varphi$ and
$L(\alpha)=S$. The table distinguishes this value from four other
Lambert-series values. The identity changes the coefficients and the
form of the denominator; it does not transfer an irrationality theorem
for one weight to another.
\par\noindent\ev{Lean} \scale{n/a} \coord{mobius-mersenne} \lean{tsum\_totient\_div\_pow\_two\_eq\_pnat\_half\_pow}{CertificateKernel.lean:18415}
\end{prop}
\leannote{Lean: \lproof{ErdosProblems/Erdos249/PaperCompleteR21/LambertDivisorTransform.lean}{42}{lambertValue_eq_divisor_sum_series}, \lproof{ErdosProblems/Erdos249/PaperCompleteR21/LambertDivisorTransform.lean}{101}{alpha_divisor_sum_eq_totient}, \lproof{ErdosProblems/Erdos249/PaperCompleteR21/LambertDivisorTransform.lean}{109}{lambertValue_alpha_eq_totientSeries}.}

\begin{longtable}{@{}
  >{\raggedright\arraybackslash}p{0.14\linewidth}
  >{\raggedright\arraybackslash}p{0.33\linewidth}
  >{\raggedright\arraybackslash}p{0.43\linewidth}
  @{}}
\caption{Five Lambert-series values and the support for their stated arithmetic status.}\\
\toprule
\papertablehead
Weight $f$ & Value $L(f)$ & Status \\
\midrule
\endfirsthead
$\mu$ & $L(\mu) = 1/2$ & rational, trivial \\
$\varphi$ & $L(\varphi) = 2$ & rational \\
$\mathrm{Id}$ & $L(\mathrm{Id}) = \sum_m \sigma(m)/2^m$ & transcendental, since $1-24L(\mathrm{Id})$ is Ramanujan's $P(1/2)$ (Nesterenko~\cite[Cor.~2, p.~1320]{nesterenko1996}, \ev{Cited}, not formalised) \\
$1$ & $L(1) = E$, Erd\H{o}s--Borwein constant & irrational, \ev{Lean} (\nolinkurl{irrational_erdosBorwein_series}, \codeid{CertificateKernel.lean:8007}); matches \#257's full-support case \\
$\alpha = \varphi*\mu$ & $L(\alpha) = S$ & \textbf{OPEN}; this is \#249 \\
\bottomrule
\end{longtable}

\begin{rem}[How the identities differ]
The coefficients matter: a result for the constant weight $1$ does not
settle the series with weight $\alpha=\varphi*\mu$. A theorem expressed
in Lambert-series form applies here only after its assumptions have
been checked for that weight.

The supplied source review checked the equality between the two
indexing conventions for $S$ directly:
\lean{tsum\_totient\_div\_pow\_two\_eq\_pnat\_half\_pow}{CertificateKernel.lean:18415}. For the other rows it consulted the summary in the source
rather than every proof body. The table retains the distinction between
a cited theorem, an ordinary identity, and a named formal declaration.
\end{rem}



\begin{prop}[An affine combination annihilates the specified divisor terms]\label{catalogue:mob:e3}
Let $H\ge1$ and $d\mid H$, with $d>0$. The $d$th term in the
M\"obius expansion of $R_{mH}$ is $\mu(d)K_d(mH)$, where
\[
 K_d(mH)=\frac{mH}{d(2^d-1)}+\frac{2^d}{(2^d-1)^2}.
\]
This is affine in $m$. Consequently, for any finite family of real
coefficients $c_i$ and nonnegative integers $m_i$ with
$\sum_i c_i=\sum_i c_i m_i=0$, one has
$\sum_i c_i K_d(m_iH)=0$.

For the four multipliers $(1,3,5,15)$ and coefficients $(4,-3,-2,1)$,
\[
 K_d(15H)-3K_d(3H)-2K_d(5H)+4K_d(H)=0.
\]
The zeroth and first moments are zero, while the second is
$15^2-3\cdot3^2-2\cdot5^2+4=152$.
Equivalently, $XY-3X-2Y+4$ takes the values $0,0,152$ at
$(1,1),(3,5),(9,25)$. This explains the coefficients: they cancel the
affine contribution of each divisor of $H$, without cancelling a general
quadratic function of the multiplier.

For a finite truncation of depth $L\ge0$, write
\[
 U=\sum_{j=1}^{L}
 \bigl(\varphi(15H+j)-3\varphi(3H+j)
       -2\varphi(5H+j)+4\varphi(H+j)\bigr)2^{L-j}.
\]
The difference between $2^L(R_{15H}-3R_{3H}-2R_{5H}+4R_H)$ and $U$
has absolute value at most $19H+5L+5$. Indeed, the error equals
\[
 (R_{15H+L}+4R_{H+L})-(3R_{3H+L}+2R_{5H+L}),
\]
and both parenthesised terms lie between $0$ and $19H+5L+5$ by
$0\le R_n\le n+1$ for $n\ge1$.
Thus
\[
 19H+5L+5<U\bmod2^L<2^L-(19H+5L+5)
\]
is sufficient for this four-tail combination to be nonintegral.
The bound is derived from these tail enclosures; no optimality for the
actual totient tails is asserted. Obtaining such certificates at the
required unbounded family of LCM heights remains unproved, as in
Theorems~\ref{catalogue:mob:e2} and~\ref{catalogue:cert:b6}.
\par\noindent\ev{Lean} \scale{uniform} \coord{other:lcm-diagonal} \lean{oldChannel\_affine\_moment\_annihilation}{JointExponentTransport.lean:36} \lean{joint35\_oldChannel\_zero}{JointExponentTransport.lean:71} \lean{sharpJoint35ConeRadius}{JointExponentTransport.lean:124}
\end{prop}
\leannote{Lean: \lproof{ErdosProblems/Erdos249/PaperCompleteR21/AffineDivisorAnnihilation.lean}{28}{transportResidueKernel_eq_mobiusTermKernel}, \lproof{ErdosProblems/Erdos249/PaperCompleteR21/AffineDivisorAnnihilation.lean}{47}{mobiusTermKernel_affine_in_multiplier}, \lproof{ErdosProblems/Erdos249/PaperCompleteR21/AffineDivisorAnnihilation.lean}{99}{joint35ConeWindow_eq}, \lproof{ErdosProblems/Erdos249/PaperCompleteR21/AffineDivisorAnnihilation.lean}{119}{totientTail_enclosure}, and 5 further declarations in the \hyperref[sec:coverage]{coverage section}.}

\begin{prop}[The error under composite dilation]\label{catalogue:mob:f1}
For $A\subseteq\mathbb N$ and $n\ge1$, let
$d_A(n)=\#\{d\ge1:d\mid n,\ d\in A\}$. If $a\in A$ and $a,x\ge1$, then
\[
 d_A(ax)=d_A(x)+\mathbf 1_{a\nmid x}
       +\#\{d\mid ax:d\in A,\ d\nmid x,\ d\ne a\}.
\]
Indeed, partition the divisors of $ax$ into those already dividing $x$,
the possible new divisor $a$, and all other new divisors. If every element
of $A$ is prime, the last set is empty, since a prime dividing $ax$ but
not $x$ must equal the prime $a$. For composite $a$ it can be nonempty:
at $A=\mathbb N$, $a=6$, $x=1$, the other new divisors are $2$ and $3$.

This is an unweighted counting identity. The associated Lambert series
satisfies
\[
 \sum_{\substack{a\in A\\a\ge1}}\frac1{2^a-1}
 =\sum_{n\ge1}\frac{d_A(n)}{2^n}.
\]
For the totient series the required divisor-convolution weight is instead
$\alpha=\varphi*\mu$, since $\alpha*1=\varphi$, as in
Proposition~\ref{catalogue:cert:d7}. Weighting all divisors by $\varphi$
would give $\sum_{d\mid n}\varphi(d)=n$, not $\varphi(n)$; already at
$n=2$ these values are $2$ and $1$. A weighted version must retain the
weights of the new divisors, so it cannot be obtained by substituting
$A=\mathbb N$ in this unweighted formula. Also $\alpha(p)=p-2$ for primes
$p$, so $\alpha$ is unbounded and is not a periodic weight.
\par\noindent\ev{Lean} \scale{uniform} \coord{other:support-divisor-counting} \lean{supportCoeff\_mul\_eq\_add\_defect}{CompositeDilationDefect.lean:30} \lean{compositeDilationDefect\_eq\_zero\_of\_prime\_support}{CompositeDilationDefect.lean:103} \lean{supportCoeff\_mul\_prime\_support}{CompositeDilationDefect.lean:119}
\end{prop}
\leannote{Lean: \lproof{ErdosProblems/Erdos249/PaperCompleteR21/CompositeDilationIdentity.lean}{27}{composite_dilation_divisor_count}, \lproof{ErdosProblems/Erdos249/PaperCompleteR21/CompositeDilationIdentity.lean}{36}{composite_dilation_defect_eq_zero_of_prime_support}, \lproof{ErdosProblems/Erdos249/PaperCompleteR21/CompositeDilationIdentity.lean}{42}{composite_dilation_divisor_count_prime_support}, \lproof{ErdosProblems/Erdos249/PaperCompleteR21/CompositeDilationIdentity.lean}{49}{composite_dilation_defect_univ_six_one}, and 8 further declarations in the \hyperref[sec:coverage]{coverage section}.}

\paragraph{Iterated prime dilations of divisor counts.}
Multiplying an argument by a new prime adds exactly one new family of
divisors. For the divisor count $d_A$ above and $n\ge1$, put
$\Delta_p f(n)=f(pn)-f(n)$. If $p_1,\ldots,p_r$ are distinct primes,
$P=p_1\cdots p_r$ and $\gcd(n,P)=1$, then
\[
 \Delta_{p_1}\cdots\Delta_{p_r}d_A(n)
       =\#\{d\mid n:Pd\in A\}.
\]
For one prime, the divisors of $pn$ not dividing $n$ are exactly $pd$
with $d\mid n$, because $p\nmid n$. Iteration gives the displayed
identity: distinctness ensures that each new prime is coprime to the
arguments produced at the earlier steps. Equivalently, inclusion--exclusion
cancels every divisor not containing all the new prime factors.

Let $\Omega(a)$ count the prime factors of $a\ge1$ with multiplicity,
and let $K\ge0$ be an integer. If $\Omega(a)\le K$ for every positive $a\in A$ and $r>K$, the set on
the right is empty, since $\Omega(Pd)=r+\Omega(d)>K$.
This proves the stated vanishing with all its hypotheses.
\lean{Erdos249257.SupportDilationDifferences.iteratedDilationDifference_eq_zero_of_cardFactors_le}{SupportDilationDifferences.lean:250}
The source also proves vanishing of the right-hand divisor count for a
list with repeated primes. That is a statement about the condition
$Pd\in A$, not an identity with repeated dilation differences.

For example, with $A=\{2\}$ and $n=1$,
\[
 \Delta_2^2d_A(1)=d_A(4)-2d_A(2)+d_A(1)=-1,
 \qquad \#\{d\mid1:4d\in A\}=0.
\]
Thus the bound $\Omega(a)\le1$ and two prime factors do not by themselves
make the repeated difference vanish. The strict inequality $r>K$ is
also material: $\Delta_2d_{\{2\}}(1)=1$ when $r=K=1$.
Finally, at $n=2$ the one-prime identity without the coprimality hypothesis
would give $\Delta_2d_{\{2\}}(2)=0$ on the left and
$\#\{d\mid2:2d\in\{2\}\}=1$ on the right.
This last example concerns the identity, not the necessity of coprimality
for every possible vanishing theorem. These divisor counts belong to
the unweighted support problem; they are not the totient coefficients.

\subsubsection{Equivalent quantified certificate conditions}

Each statement in this block relates a quantified certificate condition
to irrationality. The implications are proved; the assertions that
certificates occur beyond every threshold remain unproved. Completeness
gives equivalences for the base condition and the LCM-diagonal condition.
The one-directional implications remain useful proof components, but do
not by themselves describe the full logical relation.

\begin{thm}[The quantified certificate condition is equivalent to irrationality]\label{catalogue:cert:a10}
$\big(\forall h:\mathbb{N},\ 0<h \to \forall N_0:\mathbb{N},\ \exists N\ge N_0,\ \exists L,\ \mathcal{C}(h,N,L)\big) \leftrightarrow S\notin\mathbb Q$. The quantified-condition side is exactly $\mathcal{C}(h,N,L)$ quantified as $\forall h\ge1\ \forall N_0\ge0\ \exists N\ge N_0\ \exists L$. The universally quantified assertion remains unproved. The equivalence identifies the finite witnesses that suffice; it does not establish their existence beyond every threshold.
\par\noindent\ev{Lean} (equivalence proved; supply \ev{Open}) \scale{cofinal} \coord{other:binary-window} \lean{irrational\_totient\_series\_iff\_certificate\_supply}{LcmConeFlatness.lean:412}
\end{thm}

\begin{thm}[It suffices to use multiples of a period]\label{catalogue:cert:b2}
$\big(\forall h_0>0,\ \forall N_0,\ \exists m>0,\ \exists N\ge N_0,\ \exists L,\ \mathcal{C}(m\cdot h_0, N, L)\big) \to S\notin\mathbb Q$. The shift may be any positive multiple of a prescribed period. This
enlarges the choice of finite witness: Theorem~\ref{catalogue:cert:a10}
gives the condition with $m=1$. Conversely, rationality would make all
such tail differences integral after a fixed starting index. The
displayed implication therefore makes this quantified condition
equivalent to irrationality; its truth remains unproved.
\par\noindent\ev{Lean} (implication proved; hypothesis \ev{Open}) \scale{cofinal} \coord{other:lcm-period-multiple} \lean{irrational\_totient\_series\_of\_multiple\_certificate\_supply}{CarrySurvivorExtinction.lean:502}
\end{thm}
\leannote{Lean: \lproof{ErdosProblems/Erdos249/PaperCompleteR21/PeriodMultipleAndSecondDifference.lean}{25}{irrational_of_period_multiple_certificate_supply}, \lproof{ErdosProblems/Erdos249/PaperCompleteR21/PeriodMultipleAndSecondDifference.lean}{33}{period_multiple_certificate_at_one}, \lproof{ErdosProblems/Erdos249/PaperCompleteR21/PeriodMultipleAndSecondDifference.lean}{55}{rational_forces_period_multiple_integrality}, \lproof{ErdosProblems/Erdos249/PaperCompleteR21/PeriodMultipleAndSecondDifference.lean}{41}{period_multiple_certificate_supply_iff}.}

\begin{thm}[The diagonal condition is equivalent to irrationality]\label{catalogue:cert:b3}
$\big(\forall t_0:\mathbb{N},\ \exists t\ge t_0,\ \exists L,\ \mathcal{C}({H}(t), {H}(t), L)\big) \leftrightarrow S\notin\mathbb Q$. For a fixed hypothetical rational value, sufficiently large $t$ makes both its dyadic denominator and its odd-part period admissible at $N=h={H}(t)$. Conversely, irrationality and pointwise completeness supply a witness for every prescribed $t$. This is an equivalent condition with one scale parameter; the verified cases $t\le82$ do not establish it at arbitrarily large scales.
\par\noindent\ev{Lean} (equivalence proved; supply \ev{Open}) \scale{cofinal} \coord{other:lcm-diagonal} \lean{irrational\_totient\_series\_iff\_lcm\_diagonal\_certificate\_supply}{LcmConeFlatness.lean:426}
\end{thm}

\begin{lem}[Nondivisors in a short LCM window]\label{catalogue:cert:b4}
Let $t\ge1$ and $1\le j<2t$ be integers. If $j\nmid H(t)$, then
$j=p^a>t$ for a prime $p$ and an integer $a\ge1$. Indeed, a prime-power
divisor of $j$ must exceed $t$, and $j<2t$ leaves no room for a
cofactor larger than $1$. This classifies the exceptional offsets in
this short window; it does not bound their contribution to a weighted
sum of totient differences.

\par\noindent\ev{Lean} \scale{uniform} \coord{other:lcm-window} \lean{eq\_prime\_pow\_of\_not\_dvd\_periodLcm}{LcmDiagonalReduction.lean:137}
\end{lem}

\begin{prop}[Totient factorisation on an LCM progression]\label{catalogue:cert:b5}
Let $t\ge1$, $j\ge1$ and $q\ge0$ be integers. Suppose $j\mid H(t)$
and every prime divisor of $j$ also divides $H(t)/j$. Then
\[
 \varphi(qH(t)+j)=\varphi(j)\,
                 \varphi\bigl(q(H(t)/j)+1\bigr).
\]
The two factors in $qH(t)+j=j\bigl(q(H(t)/j)+1\bigr)$ are coprime:
the second is $1$ modulo each prime dividing $j$. Totient
multiplicativity gives the identity. The extra prime-divisor hypothesis
is not automatic from $j\mid H(t)$; for example, it fails at $t=j=2$,
where $\varphi(H(2)+2)=2$ but the proposed product would be $1$.

\par\noindent\ev{Lean} \scale{uniform} \coord{other:lcm-window-multiplicative} \lean{totient\_periodLcm\_ray\_split}{LcmDiagonalReduction.lean:208}
\end{prop}
\leannote{Lean: \lproof{ErdosProblems/Erdos249/PaperCompleteR20/LcmGridCorrespondence.lean}{67}{clean_lcm_ray_factorisation}, \lproof{ErdosProblems/Erdos249/PaperCompleteR20/LcmGridCorrespondence.lean}{76}{unclean_lcm_ray_counterexample}.}

\begin{thm}[Rationality forces tail integrality on an LCM grid]\label{catalogue:cert:b6}
If $S\in\mathbb Q$, there is $t_1\in\mathbb N$ such that for all
$t\ge t_1$, $q\ge1$ and $m\ge0$,
\[
 R_{(q+m)H(t)}-R_{qH(t)}\in\mathbb Z.
\]
Thus all tails at positive multiples of $H(t)$ have the same fractional
part once $t\ge t_1$. A certificate at one fixed scale does not refute
rationality: the threshold $t_1$ depends on the hypothetical rational
value. At a shift $mH(t)>0$ and basepoint $qH(t)$, such a certificate
excludes rational values whose reduced denominator divides
$2^{qH(t)}(2^{mH(t)}-1)$. Certificates at arbitrarily large scales,
as required in Theorem~\ref{catalogue:cert:b7}, exclude every hypothetical
rational value.

\par\noindent\ev{Lean} \scale{uniform} \coord{other:lcm-cone} \lean{rational\_totient\_series\_forces\_lcm\_cone\_flatness}{CertificateKernel.lean:19002}
\end{thm}
\leannote{Lean: \lproof{ErdosProblems/Erdos249/PaperCompleteR20/LcmGridCorrespondence.lean}{8}{lcm_grid_flatness}, \lproof{ErdosProblems/Erdos249/PaperCompleteR20/LcmGridCorrespondence.lean}{15}{lcm_grid_fractional_parts}, \lproof{ErdosProblems/Erdos249/PaperCompleteR20/LcmGridCorrespondence.lean}{27}{certificate_denominator_exclusion}, \lproof{ErdosProblems/Erdos249/PaperCompleteR20/LcmGridCorrespondence.lean}{40}{lcm_grid_supply_iff}.}

\begin{thm}[A sufficient nonintegrality condition on the grid]\label{catalogue:cert:b7}
Suppose that for every $t_0\in\N$ there are integers
$t\ge t_0$, $q\ge1$, $m\ge0$ and $L\ge0$ such that
\[
 \mathcal C(mH_t,qH_t,L).
\]
Then $S\notin\mathbb Q$. Rationality would make the corresponding
tail difference integral at every sufficiently large scale, contradicting
Proposition~\ref{catalogue:cert:a6}. A successful certificate necessarily
has $m>0$, since the discrepancy vanishes at shift zero.

The diagonal choice is $q=m=1$; consecutive grid points have $m=1$,
and for $p\ge2$ the pair $(H_t,pH_t)$ has $(q,m)=(1,p-1)$. The diagonal equivalence
in Theorem~\ref{catalogue:cert:b3} also gives the converse from
irrationality to the displayed quantified condition.
\par\noindent\ev{Lean} (implication proved; hypothesis \ev{Open}) \scale{cofinal} \coord{other:lcm-cone} \lean{irrational\_totient\_series\_of\_lcm\_cone\_window\_kill\_supply}{CertificateKernel.lean:19014}
\end{thm}
\leannote{Lean: \lproof{ErdosProblems/Erdos249/PaperCompleteR20/LcmGridCorrespondence.lean}{40}{lcm_grid_supply_iff}, \lproof{ErdosProblems/Erdos249/PaperCompleteR20/LcmGridCorrespondence.lean}{50}{lcm_grid_multiplier_positive}.}

\begin{cor}[Equivalent conditions stated without certificates]\label{catalogue:cert:b8}
The series $S$ is irrational if and only if
\[
 R_{2H_t}-R_{H_t}\notin\mathbb Z
 \qquad\text{for arbitrarily large integers }t.
\]
Equivalently, for arbitrarily large $t$ there are integers $q\ge1$
and $m\ge1$ with $R_{(q+m)H_t}-R_{qH_t}\notin\mathbb Z$.
Apply the pointwise equivalence of Theorem~\ref{catalogue:cert:a7}
to Theorems~\ref{catalogue:cert:b3} and~\ref{catalogue:cert:b7},
respectively. These are equivalent statements of the irrationality
question; the required unbounded sets of scales are not established.
\par\noindent\ev{Lean} (implication proved; hypothesis \ev{Open}) \scale{cofinal} \coord{other:real-analytic-nonintegrality} \lean{irrational\_totient\_series\_of\_lcm\_diagonal\_nonintegrality\_supply}{CertificateKernel.lean:19034} \lean{irrational\_totient\_series\_of\_lcm\_cone\_nonintegrality\_supply}{CertificateKernel.lean:19046}
\end{cor}
\leannote{Lean: \lproof{Erdos249257/LcmConeFlatness.lean}{386}{irrational_totient_series_iff_all_tail_diffs_nonintegral}, \lproof{Erdos249257/LcmConeFlatness.lean}{439}{periodLcm_diagonal_kill_iff_tail_diff_notMem_int}, \lproof{Erdos249257/LcmConeFlatness.lean}{451}{irrational_totient_series_of_lcm_diagonal_nonintegrality_supply}, \lproof{Erdos249257/LcmConeFlatness.lean}{462}{irrational_totient_series_of_lcm_cone_nonintegrality_supply}.}

\begin{prop}[Soundness of a second-difference certificate]\label{catalogue:cert:b9a}
Let $h,N,L\in\mathbb N$. If
\[
\begin{aligned}
 2(N+2h+L+2)
 &<\bigl(D(h,N+h,L)-D(h,N,L)\bigr)\bmod2^L\\
 &<2^L-2(N+2h+L+2),
\end{aligned}
\]
then $R_{N+2h}-2R_{N+h}+R_N\notin\mathbb Z$.
The numerator is the second difference of the finite windows.
Subtracting the two tail identities leaves an error of absolute value
at most $2(N+2h+L+2)$, which explains the radius in the hypothesis.
At $(h,N)=(1,8)$, the first-difference test holds at depth $8$,
whereas the second-difference test fails at every depth $L\le8$
and holds at $L=9$. A historical probe over $t\le20$ reported that the
first-difference test was at least as shallow in $30$ of $40$ sampled
cases. This finite comparison is not a uniform ordering of the depths.
\par\noindent\ev{Lean} \scale{uniform} \coord{other:binary-window}
\lean{second\_diff\_notMem\_int\_of\_certifiedRank2Kill}{LcmConeFlatness.lean:493} (the certified $(1,8)$ comparison: \lean{totient\_tail\_rank\_two\_kill\_sound\_but\_not\_shallower\_cell}{CertificateKernel.lean:19090}).
\end{prop}
\leannote{Lean: \lproof{ErdosProblems/Erdos249/PaperCompleteR21/PeriodMultipleAndSecondDifference.lean}{65}{abs_tail_diff_scaled_sub_window_le}, \lproof{ErdosProblems/Erdos249/PaperCompleteR21/PeriodMultipleAndSecondDifference.lean}{104}{second_difference_error_bound}, \lproof{ErdosProblems/Erdos249/PaperCompleteR21/PeriodMultipleAndSecondDifference.lean}{121}{second_difference_certificate_sound}, \lproof{ErdosProblems/Erdos249/PaperCompleteR21/PeriodMultipleAndSecondDifference.lean}{136}{second_difference_cell_one_eight}.}


For $H,L\in\mathbb N$ and a finite set $Q$ of multipliers, put
\[
 A_q=\sum_{j=1}^{L}\varphi(qH+j)2^{L-j},\qquad
 B_q=qH+L+2\quad(q\in Q).
\]
The finite-grid test compares these integer numerators modulo $2^L$.

\begin{thm}[A finite-grid condition gives a nonintegral pair]\label{catalogue:cert:b10a}
Let $Q\subseteq\mathbb N_{>0}$ be finite and nonempty. Suppose that
$B_q<2^L$ for every $q\in Q$ and that
\[
 \forall q_i\in Q\ \exists q_j\in Q,\qquad
 B_{q_j}<(A_{q_i}-A_{q_j})\bmod2^L.
\]
Then $R_{q_jH}-R_{q_iH}\notin\mathbb Z$ for some $q_i,q_j\in Q$.
\par\noindent\ev{Lean} \scale{uniform} \coord{other:lcm-cone-menu} \lean{exists\_nonintegral\_pair\_of\_coneNonflatCert}{LcmConeNonflat.lean:126}
\end{thm}
\leannote{Lean: \lproof{ErdosProblems/Erdos249/PaperCompleteR20/FiniteGridCorrespondence.lean}{12}{paperGridNumerator_eq}, \lproof{ErdosProblems/Erdos249/PaperCompleteR20/FiniteGridCorrespondence.lean}{36}{finite_grid_nonintegral_pair}.}
\begin{proof}
The tail identity and its elementary bound give
\[
 2^LR_{qH}=A_q+R_{qH+L},\qquad 0\le R_{qH+L}\le B_q.
\]
Suppose all pairwise tail differences are integers. Choose $q_*\in Q$
minimising $R_{qH+L}$. For every $q\in Q$, the difference
$R_{qH+L}-R_{q_*H+L}$ is then an integer in $[0,B_q]$, congruent to
$A_{q_*}-A_q$ modulo $2^L$. Since $B_q<2^L$, it is the least
nonnegative residue. This contradicts the displayed condition at $q_*$.
\end{proof}

For example, take $H=840$, $L=12$ and $Q=\{1,2,3,4\}$.
The following exact values verify one failed arc containment in each row:
\begin{center}
\begin{tabular}{rrrrr}
$q_i$ & $A_{q_i}\bmod4096$ & $B_{q_i}$ & $q_j$ & $(A_{q_i}-A_{q_j})\bmod4096$\\
\hline
1 & 2920 & 854  & 3 & 3484\\
2 & 2024 & 1694 & 1 & 3200\\
3 & 3532 & 2534 & 4 & 3832\\
4 & 3796 & 3374 & 1 & 876
\end{tabular}
\end{center}
The last column exceeds $B_{q_j}$, not necessarily $B_{q_i}$.
On the circle $\mathbb R/4096\mathbb Z$, every two arcs
$[A_q,A_q+B_q]$ intersect, but the four arcs have no common point.
Thus a test using all four points can exclude a common fractional part
without any pair of these arcs being disjoint. The theorem does not
require pairwise compatibility, or require $|Q|\ge3$.

The size hypothesis is $2^L>\max_{q\in Q}B_q$. In contrast, the
symmetric test $\mathcal C$ for a pair $q_i<q_j$ needs
$2^L>2B_{q_j}$ before its two-sided interval can contain a residue.
These are inequalities for the modulus $2^L$. They do not say that the
new test uses half as many binary digits, or give a uniform comparison
of the first depths at which the tests succeed.

\begin{thm}[A sufficient quantified finite-grid condition]\label{catalogue:cert:b10b}
Suppose that for every $t_0$ there are $t\ge t_0$, a depth $L$, and a
finite nonempty $Q\subseteq\mathbb N_{>0}$ such that, with $H=H_t$,
all the hypotheses of Theorem~\ref{catalogue:cert:b10a} hold.
Then $S\notin\mathbb Q$. Indeed, rationality would make all tail
differences on every sufficiently large LCM grid integral, whereas the
finite-grid theorem supplies a nonintegral pair on such a grid.
The example above verifies one grid; it does not establish the
quantified hypothesis.
\par\noindent\ev{Lean} (implication proved; hypothesis \ev{Open}) \scale{cofinal} \coord{other:lcm-cone-menu} \lean{irrational\_totient\_series\_of\_lcm\_cone\_nonflat\_supply}{CertificateKernel.lean:19140}
\end{thm}
\leannote{Lean: \lproof{ErdosProblems/Erdos249/PaperCompleteR20/FiniteGridCorrespondence.lean}{12}{paperGridNumerator_eq}, \lproof{ErdosProblems/Erdos249/PaperCompleteR20/FiniteGridCorrespondence.lean}{49}{finite_grid_supply_irrational}.}

\begin{prop}[A finite carry test implies nonintegrality]\label{catalogue:cert:b12}
For each integer $z$ with $|z|\le N+h+1$, define
\[
 c_z(0)=z,\qquad
 c_z(i+1)=2c_z(i)-\bigl(\varphi(N+h+i+1)-\varphi(N+i+1)\bigr).
\]
Suppose that for every such $z$ there is an $i\le K$ with
\[
 |c_z(i)|\ge N+i+h+2.
\]
Then $R_{N+h}-R_N\notin\mathbb Z$. Indeed, an integral value of the
tail difference would be one of the $2(N+h+1)+1$ initial candidates.
Its recurrence would satisfy $c_z(i)=R_{N+h+i}-R_{N+i}$ and hence
$|c_z(i)|<N+i+h+2$ at every step, contrary to the finite test.
The excluded strip is open: equality at its boundary already suffices.
This proposition proves soundness for a specified $h,N,K$; an
irrationality argument still needs the quantified family of tests.
\par\noindent\ev{Lean} \scale{uniform} \coord{other:carry-orbit} \lean{survivorKill}{CarrySurvivorExtinction.lean:412} \lean{tail\_diff\_notMem\_int\_of\_survivorKill}{CarrySurvivorExtinction.lean:428}
\end{prop}
\leannote{Lean: \lproof{ErdosProblems/Erdos249/PaperCompleteR20/FiniteCarryCorrespondence.lean}{8}{finite_carry_test_sound}, \lproof{ErdosProblems/Erdos249/PaperCompleteR20/FiniteCarryCorrespondence.lean}{22}{finite_carry_candidate_count}, \lproof{ErdosProblems/Erdos249/PaperCompleteR20/FiniteCarryCorrespondence.lean}{27}{finite_carry_true_orbit}.}

\begin{thm}[The block norm condition implies irrationality]\label{catalogue:mob:e1}
Suppose that for every integer $h\ge1$ and every threshold $X_0\in\N$
there are $X,L\in\N$ with
\[
 X\ge\max(X_0,1),\qquad 16(2X+h+L+2)\le2^L,
\]
and
\[
 \left\lVert\sum_{X\le N<2X}E(h,N,L)\right\rVert\le\frac{21}{25}X.
\]
Then $S\notin\mathbb Q$. Indeed, the norm bound implies the real-part
bound $\operatorname{Re}\sum E(h,N,L)\le(9/10)X$. Under the displayed
size condition, this latter bound already gives a finite residue
certificate, by $\cos(\pi/8)>9/10$ and averaging. The argument uses this
specific numerical gap; an unspecified saving below $X$ is not the
stated hypothesis.
\par\noindent\ev{Lean} (conditional theorem proved; hypothesis \ev{Open}) \scale{cofinal} \coord{binary-digit} \lean{irrational\_totient\_series\_of\_first\_harmonic\_norm\_gap}{FirstHarmonicPivot.lean:83} \lean{exists\_certifiedKill\_of\_first\_harmonic\_gap}{FirstHarmonicGap.lean:128}
\end{thm}
\leannote{Lean: \lproof{ErdosProblems/Erdos249/PaperCompleteR21/HarmonicGapAndFourTail.lean}{71}{irrational_of_first_harmonic_norm_gap}, \lproof{ErdosProblems/Erdos249/PaperCompleteR21/HarmonicGapAndFourTail.lean}{40}{first_harmonic_re_bound_of_norm_bound}, \lproof{ErdosProblems/Erdos249/PaperCompleteR21/HarmonicGapAndFourTail.lean}{54}{exists_certificate_of_first_harmonic_real_bound}, \lproof{ErdosProblems/Erdos249/PaperCompleteR21/HarmonicGapAndFourTail.lean}{61}{exists_certificate_of_first_harmonic_norm_bound}, and 1 further declaration in the \hyperref[sec:coverage]{coverage section}.}

\begin{thm}[A four-tail residue criterion]\label{catalogue:mob:e2}
For nonnegative integers $H,p,L$, define the four-tail combination,
its finite numerator and its error bound by
\[
\begin{aligned}
 J(H,p)&=R_{2pH}-R_{pH}-pR_{2H}+pR_H,\\
 W(H,p,L)&=D(pH,pH,L)-pD(H,H,L),\\
 B(H,p,L)&=3pH+(p+1)(L+2).
\end{aligned}
\]
Then
\[
 |2^LJ(H,p)-W(H,p,L)|\le B(H,p,L).
\]
Consequently the finite condition
\[
 B(H,p,L)<W(H,p,L)\bmod2^L<2^L-B(H,p,L)
\]
implies $J(H,p)\notin\mathbb Z$. If, for every $t_0\in\mathbb N$,
there are $t\ge t_0$, $p\ge1$ and $L\ge0$ satisfying this condition
with $H=H(t)$, then $S\notin\mathbb Q$. Rationality would instead make
both diagonal differences in $J(H(t),p)$ integral for every sufficiently
large $t$.

The error bound groups the remainder as
\[
 (R_{2pH+L}+pR_{H+L})-(R_{pH+L}+pR_{2H+L}).
\]
Each parenthesis is nonnegative and at most $B(H,p,L)$, by
$0\le R_n\le n+2$. Their difference therefore has absolute value at
most $B(H,p,L)$. This proves the bound without asserting optimality for
the actual totient tails. Primality of $p$ is not required by the
criterion; a prime $p$ is useful when comparing consecutive LCM heights.
The checked instance is $(H,p,L)=(12,5,15)$, for which
$W(H,p,L)\bmod2^L=18834$ and $B(H,p,L)=282$.

\par\noindent\ev{Lean} (conditional theorem + \ev{Cert} witness; supply \ev{Open}) \scale{cofinal} \coord{other:lcm-diagonal} \lean{irrational\_totient\_series\_of\_primeJumpSharpKill\_supply}{PrimeJumpWindow.lean:193} \lean{primeJumpSharpKill\_twelve\_five}{PrimeJumpWindow.lean:186}
\end{thm}
\leannote{Lean: \lproof{ErdosProblems/Erdos249/PaperCompleteR21/HarmonicGapAndFourTail.lean}{103}{four_tail_combination_eq}, \lproof{ErdosProblems/Erdos249/PaperCompleteR21/HarmonicGapAndFourTail.lean}{93}{four_tail_window_eq}, \lproof{ErdosProblems/Erdos249/PaperCompleteR21/HarmonicGapAndFourTail.lean}{81}{windowDiscrepancy_diagonal_eq}, \lproof{ErdosProblems/Erdos249/PaperCompleteR21/HarmonicGapAndFourTail.lean}{113}{four_tail_error_bound}, and 5 further declarations in the \hyperref[sec:coverage]{coverage section}.}
% ; - part a249_p1b ; -
\subsection{Consequences}

The next conditions concern a single LCM-diagonal tail difference. For
$t,j,a\in\mathbb N$, write
\[
 \delta_t(j)=\varphi(2H(t)+j)-\varphi(H(t)+j),
 \qquad \Omega_a=R_{2H(2^a)}-R_{H(2^a)}.
\]
For $h,N\in\mathbb N$ and $z\in\mathbb Z$, define the integer recurrence
\[
 c_{h,N,z}(0)=z,\qquad
 c_{h,N,z}(k+1)=2c_{h,N,z}(k)
   -\bigl(\varphi(N+h+k+1)-\varphi(N+k+1)\bigr).
\]
If $z=R_{N+h}-R_N$, then $c_{h,N,z}(k)=R_{N+h+k}-R_{N+k}$.
This separates an arbitrary integer initial value from the actual tail.

Two finite quantities will also recur. With $H=H_{2^a}$, put
\[
 \rho_{a,q}=\frac{D(H,H,2q+1)+\delta_{2^a}(2q+2)}{2^{2q+1}},
 \qquad \varepsilon_{a,q}=\frac{2H+2q+3}{2^{2q+1}}.
\]
The prescribed odd depth used below is $2q_a+1$, where
\[
 q_a=\left\lfloor\frac{\lfloor\log_2 H\rfloor+10}{2}\right\rfloor.
\]
Thus $2q_a+1$ is the least odd integer at least
$\lfloor\log_2 H\rfloor+10$. The error bound is
$|\Omega_a-\rho_{a,q}|<\varepsilon_{a,q}$.
For $a\ge2$, let $u_{a,q}$ be the centred integer representative of
\[
 \frac{D(H,H,2q+1)+\delta_{2^a}(2q+2)}2\pmod{4^q},
 \qquad -\frac{4^q}2<u_{a,q}\le\frac{4^q}2.
\]
The numerator is even, since the totient arguments exceed $2$; at the
midpoint this convention chooses the positive representative. The
normalised block $\rho_{a,q}$ approximates $\Omega_a$, whereas
$u_{a,q}$ records a residue. They must not be confused with the real
tail difference or with the integer recurrence.

\begin{prop}[An equivalent diagonal nonintegrality condition]\label{prop:NI-01}
The following statements are equivalent:
\[
 S\notin\mathbb Q
 \quad\Longleftrightarrow\quad
 \forall a_0\in\mathbb N\ \exists a\ge a_0,\ \Omega_a\notin\mathbb Z.
\]
This is an exact reformulation of irrationality, not a weaker theorem
about a finite set of scales. The later sufficient conditions seek a
finite inequality that establishes the right-hand side.
\par\noindent\ev{Lean} (equivalence)
\lean{irrational_totientSeries_iff_actualLcmOrbitNonintegralitySupply}{TotientActualLcmOrbitNonintegrality.lean:37}
\end{prop}
\leannote{Lean: \lproof{ErdosProblems/Erdos249/PaperCompleteR21/ActualLcmDiagonalConditions.lean}{29}{irrational_iff_diagonal_orbit_nonintegrality}.}

\begin{prop}[A sufficient short-window condition]\label{prop:AR-07}
Suppose that for every $a_0\in\mathbb N$ there are $a\ge a_0$
and $L<2\cdot2^a$ with $\mathcal C(H_{2^a},H_{2^a},L)$. Then
$S\notin\mathbb Q$. Writing out the certificate, the required inequality is
\[
 2H_{2^a}+L+2
 <D(H_{2^a},H_{2^a},L)\bmod2^L
 <2^L-(2H_{2^a}+L+2).
\]
It places the residue away from both endpoints of the modulus interval.
The examples at $a=4,6$ establish the bounded statement in
Proposition~\ref{prop:SK-02}, not the condition for every threshold.
Pointwise completeness allows a sufficiently large depth but does not
supply a depth below $2\cdot2^a$.

\scale{cofinal}\ \ev{Open}\ \coord{mobius-mersenne}\\
\lean{PowerTwoActualLcmShortArithmeticKillSupply}{TotientActualLcmOrbitArithmetic.lean:2107}
\end{prop}
\leannote{Lean: \lproof{ErdosProblems/Erdos249/PaperCompleteR21/ActualLcmDiagonalConditions.lean}{56}{irrational_of_short_window_diagonal_supply}, \lproof{ErdosProblems/Erdos249/PaperCompleteR21/ActualLcmDiagonalConditions.lean}{40}{diagonal_certificate_unfolded}, \lproof{ErdosProblems/Erdos249/PaperCompleteR21/ActualLcmDiagonalConditions.lean}{68}{pointwise_completeness_supplies_some_depth}.}

\begin{prop}[A sufficient approximation condition]\label{prop:SEP-03}
Use the prescribed index $q_a$ and the error
$\varepsilon_{a,q_a}$ defined above. Suppose
\[
 \forall a_0\in\mathbb N\ \exists a\ge\max(2,a_0)\ \forall z\in\mathbb Z,
 \qquad |\Omega_a-z|\ge\frac1{32}+\varepsilon_{a,q_a}.
\]
Then $S\notin\mathbb Q$. The estimate
$|\Omega_a-\rho_{a,q_a}|<\varepsilon_{a,q_a}$ implies
$|\rho_{a,q_a}-z|>1/32$ for every integer $z$, and hence the finite
residue separation used by the linked implication.
This is a fixed separation condition along an unbounded family of LCM
scales. It is not an irrationality-measure theorem for arbitrary
rational approximants, and nonintegrality alone does not give the
stated margin. The depth is prescribed, not an additional free witness.

\scale{cofinal}\ \ev{Open}\ \coord{mobius-mersenne}\\
\lean{irrational_totientSeries_of_actualLcmOrbitSeparationSupply}{TotientActualLcmOrbitSeparation.lean:305}
\end{prop}
\leannote{Lean: \lproof{ErdosProblems/Erdos249/PaperCompleteR21/ActualLcmDiagonalConditions.lean}{145}{irrational_of_diagonal_orbit_separation_supply}, \lproof{ErdosProblems/Erdos249/PaperCompleteR21/ActualLcmDiagonalConditions.lean}{77}{prescribedOddIndex}, \lproof{ErdosProblems/Erdos249/PaperCompleteR21/ActualLcmDiagonalConditions.lean}{80}{oddGuarded_depth_eq_prescribed}, \lproof{ErdosProblems/Erdos249/PaperCompleteR21/ActualLcmDiagonalConditions.lean}{98}{abs_orbit_sub_rawApprox_lt}, and 1 further declaration in the \hyperref[sec:coverage]{coverage section}.}

\begin{prop}[A sufficient upper-endpoint separation]\label{prop:TE-04}
Let $a,J,K,m\in\mathbb N$, $a\ge8$, and $H=H(2^a)$. Assume
\[
 J+K+(a+6)<2\cdot2^a.
\]
Define the finite condition
\[
 \begin{aligned}
 \mathcal G_a(J,K,m)\quad:\Longleftrightarrow\quad
 &m\le K,\qquad 2H+J+K+2<2^m,\\
 &D(H,H+J,K)\bmod2^m\le2^m-(2H+J+K+2).
 \end{aligned}
\]
Then $\mathcal G_a(J,K,m)$ implies $R_{2H+J}-R_{H+J}\notin\mathbb Z$.
In particular, irrationality follows if for every $a_0$ there are
$a\ge\max(a_0,8)$ and $K,m$ satisfying $K+(a+6)<2\cdot2^a$ and
$\mathcal G_a(0,K,m)$.

The sign theorem forces an integral tail to give a residue in the upper
endpoint interval. This one-sided test excludes that interval; it does
not require a lower residue bound. Both the sign-range hypothesis and
the strict modulus bound are essential premises of this implication.
The required unbounded family is not established.
\par\noindent\ev{Lean} (implication; quantified hypothesis \ev{Open})
\lean{PowerTwoActualLcmTopEdgeResidueGapSupply}{TotientActualLcmTopEdgeStaircase.lean:1325}
\end{prop}
\leannote{Lean: \lproof{ErdosProblems/Erdos249/PaperCompleteR21/TopEdgeStaircaseConditions.lean}{29}{upper_endpoint_condition_iff}, \lproof{ErdosProblems/Erdos249/PaperCompleteR21/TopEdgeStaircaseConditions.lean}{38}{upper_endpoint_gap_nonintegral}, \lproof{ErdosProblems/Erdos249/PaperCompleteR21/TopEdgeStaircaseConditions.lean}{50}{integral_tail_forces_upper_endpoint_residue}, \lproof{ErdosProblems/Erdos249/PaperCompleteR21/TopEdgeStaircaseConditions.lean}{66}{irrational_of_upper_endpoint_gap_supply}.}

\begin{prop}[Relations among the sufficient conditions]\label{prop:TE-05}
The sufficient conditions do not form a single linear chain. The proved
implications are
\[
 \begin{aligned}
 \text{guarded odd-prefix band}
 &\ \Longleftrightarrow\ \text{guarded centred-magnitude bound},\\
 \text{guarded centred-magnitude bound}
 &\ \Longrightarrow\ \text{flexible centred-magnitude bound},\\
 \text{flexible centred-magnitude bound}
 &\ \Longrightarrow\ \text{adjacent-suffix band}
   \ \Longrightarrow\ \text{upper-endpoint test}.
 \end{aligned}
\]
There is also a separate sufficient argument: terminal dominance and
flexible centred magnitude each imply the two-sided exclusion in
Proposition~\ref{prop:TE-06}, which implies diagonal nonintegrality.
These statements refer to the quantified versions, with the bounds on
$a,q$ and the depth retained. The detailed conditions and sources appear
in Proposition~\ref{prop:te-chain}. No reverse implication between
terminal dominance and flexible magnitude is asserted, and none of the
required unbounded families is proved here.
\par\noindent\ev{Lean} (implications; hypotheses \ev{Open})
\lean{powerTwoActualLcmTopEdgeResidueGapSupply\_of\_adjacentSuffixMidband}{TotientActualLcmTopEdgeStaircase.lean:2108}
\end{prop}
\leannote{Lean: \lproof{ErdosProblems/Erdos249/PaperCompleteR21/TopEdgeStaircaseConditions.lean}{80}{te_chain_relations}.}

\begin{prop}[An exact endpoint formula]\label{prop:TE-06}
Let $a,q\in\mathbb N$, $a\ge8$, and $H=H(2^a)$. Assume
\[
 2q+2+(a+6)<2\cdot2^a,\qquad 2(H+q+2)\le4^q.
\]
If $z\in\mathbb Z$ and $z=\Omega_a$, then
\[
 2u_{a,q}=\delta_{2^a}(2q+2)-c_{H,H,z}(2q+1).
\]
Consequently either of the inequalities
\[
 \begin{aligned}
 2u_{a,q}&\le\delta_{2^a}(2q+2)-(2H+2q+3),\\
 \delta_{2^a}(2q+2)&\le2u_{a,q}
 \end{aligned}
\]
is sufficient for $\Omega_a\notin\mathbb Z$. Under integrality the
recurrence equals a positive tail difference strictly smaller than
$2H+2q+3$, and the identity rules out both displayed inequalities.
The identity is exact; the resulting nonintegrality criterion is only
sufficient. Its converse is not asserted.
\par\noindent\ev{Lean}
\lean{two_mul_actualOddHalfCenteredLift_eq_terminal_sub_trueCarry}{TotientActualLcmTopEdgeStaircase.lean:1678}
\end{prop}
\leannote{Lean: \lproof{ErdosProblems/Erdos249/PaperCompleteR21/TopEdgeStaircaseConditions.lean}{142}{endpoint_identity}, \lproof{ErdosProblems/Erdos249/PaperCompleteR21/TopEdgeStaircaseConditions.lean}{158}{integral_carry_strictly_between}, \lproof{ErdosProblems/Erdos249/PaperCompleteR21/TopEdgeStaircaseConditions.lean}{187}{endpoint_criterion_nonintegral}, \lproof{ErdosProblems/Erdos249/PaperCompleteR21/TopEdgeStaircaseConditions.lean}{116}{oddHalfCenteredLift_spec}, and 1 further declaration in the \hyperref[sec:coverage]{coverage section}.}

\begin{prop}[A sufficient extension of the examples through exponent 6]\label{prop:SK-02}
The supplied finite result is
\[
 \forall a_0\le6\ \exists a,L\in\mathbb N,\qquad
 a\ge a_0,\quad L<2\cdot2^a,\quad
 \mathcal C(H_{2^a},H_{2^a},L).
\]
The witnesses at $(a,L)=(4,23)$ and $(6,93)$ give this bounded family.
A witness with $a\ge7$ would extend the verified range, but any finite
extension would still leave the assertion for arbitrarily large
thresholds in Proposition~\ref{prop:AR-07} unproved.

\scale{bounded}\ \ev{Lean (finite prefix only)}\ \coord{mobius-mersenne}\\
\lean{powerTwoActualLcmShortArithmeticKillSupply_through_six}{TotientActualLcmShortKill.lean:54}
\end{prop}
\leannote{Lean: \lproof{ErdosProblems/Erdos249/PaperCompleteR21/ActualLcmDiagonalConditions.lean}{161}{short_window_diagonal_through_six}, \lproof{ErdosProblems/Erdos249/PaperCompleteR21/ActualLcmDiagonalConditions.lean}{169}{short_window_diagonal_witnesses}.}

\begin{prop}[A sufficient extremal-order condition]\label{prop:FR-01}
Let $H\ge1$ and $j\ge0$. If
\[
 \varphi(2H+j)<\min\{\varphi(H+j),\varphi(3H+j)\},
\]
then
\[
 \varphi(3H+j)-2\varphi(2H+j)+\varphi(H+j)>0.
\]
If instead the middle value is strictly larger than both outer values,
the displayed second difference is negative. In each case the claim
follows by adding the two strict differences. No lower bound on their
sizes is needed.

For $H=H_{2^a}$, this reduces nonvanishing of one second difference to an
ordering of three totient values. For a prescribed $j$, the remaining
question is whether either strict ordering occurs for arbitrarily large
$a$. The elementary implication does not establish those occurrences.
\scale{uniform}\ \ev{Lean}\ \coord{other:fixed-rank-curvature}\\
\lean{MiddleRankTotientExtremal}{TotientFixedRankLcmAsymptotic.lean:267}
\end{prop}
\leannote{Lean: \lproof{ErdosProblems/Erdos249/PaperCompleteR21/ExtremalOrderDirectedAndPulse.lean}{27}{extremal_order_curvature_pos}, \lproof{ErdosProblems/Erdos249/PaperCompleteR21/ExtremalOrderDirectedAndPulse.lean}{36}{extremal_order_curvature_neg}, \lproof{ErdosProblems/Erdos249/PaperCompleteR21/ExtremalOrderDirectedAndPulse.lean}{45}{extremal_order_curvature_ne_zero}.}

\begin{prop}[A directed certificate condition]\label{prop:CP-06}
For fixed $h,N\in\N$, a depth $L\in\N$ satisfying
\[
 N+L+2\ \le\ D(h,N,L)\bmod2^L
 \ \le\ 2^L-(N+h+L+2)
\]
exists if and only if $R_{N+h}-R_N\notin\mathbb Z$.
The endpoint inequalities are non-strict.
\end{prop}
\leannote{Lean: \lproof{ErdosProblems/Erdos249/PaperCompleteR21/ExtremalOrderDirectedAndPulse.lean}{59}{directed_certificate_iff}, \lproof{ErdosProblems/Erdos249/PaperCompleteR21/ExtremalOrderDirectedAndPulse.lean}{71}{directed_certificate_example}.}

\begin{proof}
The exact truncation identity leaves an error $R_{N+h+L}-R_{N+L}$ strictly between
$-(N+L+2)$ and $N+h+L+2$. An integral tail difference would therefore
put the residue in $[0,N+L+2)$ or
$(2^L-(N+h+L+2),2^L)$, contrary to the displayed test.
Conversely, a nonintegral difference has a symmetric certificate by
Theorem~\ref{catalogue:cert:a7}, and every symmetric certificate
satisfies this wider interval condition.
\end{proof}

At $H_3=6$ and $L=6$, the discrepancy is $D(6,6,6)=270$,
with residue $14$ modulo $64$. The directed interval is $[14,44]$,
so its lower endpoint succeeds. The symmetric interval is $(20,44)$
and fails; its first successful depth is $7$. Thus the wider test
can improve the first successful depth, but no uniform depth saving
is asserted. Replacing $\mathcal C$ by the displayed inequality in
Theorem~\ref{catalogue:cert:b3} gives the equivalent directed diagonal
condition, whose occurrence at arbitrarily large scales remains unproved.
\noindent\scale{uniform}\ \ev{Lean}\ \coord{seam-integer}\\
\lean{CofinalDirectedLcmCertificateSupply}{TotientTailCarryPeriod.lean:866}
\lean{tail_diff_notMem_int_of_directedCertifiedKill}{TotientTailCarryPeriod.lean:777}
\lean{exists_directedCertifiedKill_iff_tail_diff_notMem_int}{TotientTailCarryPeriod.lean:829}

\begin{prop}[A sufficient condition at the prescribed mod-four pulses]\label{prop:CP-07}
Suppose that for every positive integer $h$ and every $B\in\mathbb N$
there are a prime $p>B$ and $K\in\mathbb N$ such that
\[
 \varphi(p+4h)-\varphi(p)\equiv2\pmod4
\]
and every integer $z$ with $|z|\le p+4h+1$ and $z\equiv2\pmod4$
has an index $0\le i\le K$ satisfying
\[
 |c_{4h,p,z}(i)|\ge p+i+4h+2.
\]
Then $S\notin\mathbb Q$. Under a hypothetical eventual period $h$,
the proved congruence argument at a sufficiently large such prime
forces $R_{p+4h}-R_p$ to be an integer in the tested residue class.
The finite test excludes every candidate in that class, a contradiction.

The mod-four congruence is supplied by
Proposition~\ref{prop:CP-05-inv}; exclusion at the same prime is not.
Restricting to $2\pmod4$ retains one residue class out of four, not
necessarily exactly one quarter of the candidates in a finite interval.
The quantifier over every bound $B$ cannot be replaced by one fixed
prime for each $h$.

\scale{cofinal}\ \ev{Lean}\ \coord{mobius-mersenne}\\
\lean{irrational_totientSeries_of_cofinal_modFourPulseSurvivorKill}{TotientTailCarryPeriod.lean:1066}
\end{prop}
\leannote{Lean: \lproof{ErdosProblems/Erdos249/PaperCompleteR21/ExtremalOrderDirectedAndPulse.lean}{105}{irrational_of_modFour_pulse_supply}, \lproof{ErdosProblems/Erdos249/PaperCompleteR21/ExtremalOrderDirectedAndPulse.lean}{88}{rational_forces_pulse_class_integrality}.}

\begin{prop}[The full quantified condition]\label{prop:A10}
\[
\big(\forall h\ge 1,\ \forall N_0,\ \exists N\ge N_0,\ \exists L,\ \mathcal{C}\ h\ N\ L\big)
\iff S\notin\mathbb Q.
\]
This is exactly the quantifier structure $\forall h\ge1\ \forall N_0\ \exists N\ge N_0\ \exists L\ \mathcal{C}(h,N,L)$.
The quantified condition is not supplied here. The equivalence specifies
exactly what a residue-based proof would need beyond the finite examples.
\scale{cofinal}\ \ev{Lean} (equivalence; supply \ev{Open})\ \coord{mobius-mersenne}\\
\lean{irrational_totient_series_iff_certificate_supply}{LcmConeFlatness.lean:412}
\end{prop}

\begin{prop}[One diagonal parameter suffices]\label{prop:B3}
The following condition is equivalent to irrationality:
\[
 \bigl(\forall t_0\in\mathbb N\ \exists t\ge t_0\ \exists L\ge0,
       \ \mathcal C(H(t),H(t),L)\bigr)
 \quad\Longleftrightarrow\quad S\notin\mathbb Q.
\]
For a hypothetical rational value, let $h_0$ and $N_0$ be its eventual
tail period and starting index. Taking $t\ge\max(h_0,N_0)$ ensures
$h_0\mid H(t)$ and $H(t)\ge N_0$. The corresponding diagonal difference
is then integral, contradicting a certificate at that scale.
Conversely, irrationality and pointwise completeness provide a witness
at every prescribed $t$. The reduction uses one scale parameter; it
neither establishes the certificate condition nor makes one fixed scale
sufficient for all rational values.

\scale{cofinal}\ \ev{Lean} (equivalence; supply \ev{Open})\ \coord{mobius-mersenne}\\
\lean{irrational_totient_series_iff_lcm_diagonal_certificate_supply}{LcmConeFlatness.lean:426}
\end{prop}

\begin{prop}[A sufficient condition on the LCM grid]\label{prop:B7}
\[
\big(\forall t_0,\ \exists t\ge t_0,\ \exists q\,m\,L,\ 0<q \wedge \mathcal{C}(m\cdot{H}\,t)(q\cdot{H}\,t)\ L\big)
\implies S\notin\mathbb Q.
\]
One residue certificate \emph{anywhere} on the two-multiplier LCM cone, at arbitrarily large $t$,
suffices; the diagonal choice $q=m=1$ recovers Proposition~\ref{prop:B3}.
\scale{cofinal}\ \ev{Lean (implication; premise unproved)}\ \coord{mobius-mersenne}\\
\lean{irrational_totient_series_of_lcm_cone_window_kill_supply}{CertificateKernel.lean:19014}
\end{prop}


\begin{prop}[A finite-grid certificate condition]\label{prop:B10}
Let $H,L\in\mathbb N$ and let $Q\subseteq\mathbb N_{>0}$ be finite and
nonempty. Define $A_q=\sum_{j=1}^{L}\varphi(qH+j)2^{L-j}$ and
$B_q=qH+L+2$. If $B_q<2^L$ for every $q\in Q$ and
\[
 \forall q_i\in Q\ \exists q_j\in Q,\qquad
 B_{q_j}<(A_{q_i}-A_{q_j})\bmod2^L,
\]
then $R_{q_jH}-R_{q_iH}\notin\mathbb Z$ for some $q_i,q_j\in Q$.
Such a test at $H=H_t$ for arbitrarily large $t$ implies
$S\notin\mathbb Q$. Neither $|Q|\ge3$ nor compatibility of every
pair is a hypothesis. The proof, by choosing a minimum tail in the
finite set, is given in Theorem~\ref{catalogue:cert:b10a}.
\scale{cofinal}\ \ev{Lean (implication; premise unproved)}\ \coord{mobius-mersenne}\\
\lean{irrational_totient_series_of_lcm_cone_nonflat_supply}{CertificateKernel.lean:19140}
\end{prop}
\leannote{Lean: \lproof{ErdosProblems/Erdos249/PaperCompleteR20/FiniteGridCorrespondence.lean}{12}{paperGridNumerator_eq}, \lproof{ErdosProblems/Erdos249/PaperCompleteR20/FiniteGridCorrespondence.lean}{36}{finite_grid_nonintegral_pair}, \lproof{ErdosProblems/Erdos249/PaperCompleteR20/FiniteGridCorrespondence.lean}{49}{finite_grid_supply_irrational}.}

\begin{prop}[Unbounded Farey bounds imply irrationality]\label{prop:C2sup}
The Farey-gap denominator bound at window $K$ is currently $\sim 7.96\times10^{34}$ at $K=240$
(Prop.~\ref{prop:C2-inv}). If
\[
\sup_K\, (b+d)(K) = \infty
\]
(that is, the proved exclusion bounds are arbitrarily large), then \#249 follows via the denominator-exclusion implication, without assuming the certificate supply in Proposition~\ref{prop:A10}. No unboundedness claim is proved or attempted in the record.
\scale{cofinal}\ \ev{Open}\ \coord{other:farey-gap}\\
\lean{gap\_check\_window\_1\_240\_le\_79639646646701375323355774875831053}{GapFareyBound.lean:176}
\end{prop}
\leannote{Lean: \lproof{ErdosProblems/Erdos249/PaperCompleteR21/FareyGapDenominatorExclusion.lean}{41}{gapCertificate_window_1_240}, \lproof{ErdosProblems/Erdos249/PaperCompleteR21/FareyGapDenominatorExclusion.lean}{29}{irrational_totientSeries_of_fareyGapExclusionUnbounded}, \lproof{ErdosProblems/Erdos249/PaperCompleteR21/FareyGapDenominatorExclusion.lean}{51}{gapFareyBound_window_1_240}.}

\begin{prop}[The lower bound and a false proposed upper bound]\label{prop:D5cons}
By Prop.~\ref{prop:D4-inv}, the canonical dyadic totient-kernel family is unconditionally
$(2^e+1)$-dimensional at every level $e\ge1$. Rationality of $S$ forces an associated tempered
carry orbit with $\mathbb Q$-rank $\ge 2^e-1$ at every level (Prop.~\ref{prop:CP-02}). An upper bound independent of $e$ for these same carry ranks would
contradict the lower bound. More generally, a bound $g(e)$ with
$g(e)<2^e-1$ at some level $e\ge1$ would suffice, provided it applies to the
actual carry under the hypothetical rationality of $S$. A bound that
merely grows with $e$ need not contradict anything. The generic assertion
for arbitrary rational coefficient series is false: the $5/4$ comparison
sequence in Section~\ref{sec:mahler-defect} has an integer carry satisfying
the growth condition and the same rank lower bound. The incompatible
integer identity and the finite-shift counterexample (Observation~\ref{prop:B4b-kill}) remain
separate counterexample results. None is a generic rationality-driven rank ceiling.
The evidence label below applies to the proved rank floor and counterexample results,
not to the counterfactual upper bound.
\scale{uniform}\ \ev{Lean}\ \coord{other:carry-kernel-rank}\\
\lean{not_irrational_totientSeries_implies_unbounded_carryRank_unconditional}{TotientCarryKernelRigidity.lean:300}
\end{prop}
\leannote{Lean: \lproof{ErdosProblems/Erdos249/PaperCompleteR21/GenericCarryRankCeilingCounterexample.lean}{143}{rank_floor_and_false_proposed_carryRank_ceiling}, \lproof{ErdosProblems/Erdos249/PaperCompleteR21/GenericCarryRankCeilingCounterexample.lean}{64}{fiveQuarter_comparison_rational_with_carryRank_floor}, \lproof{ErdosProblems/Erdos249/PaperCompleteR21/GenericCarryRankCeilingCounterexample.lean}{100}{no_generic_rationality_carryRank_ceiling}.}

\begin{prop}[Soundness of the finite carry test]\label{prop:B12cons}
For fixed $h,N,K\in\mathbb N$, use the integer recurrences
$c_{h,N,z}$ defined above. Test each of the $2(N+h+1)+1$ candidates
$z\in\mathbb Z$ with $|z|\le N+h+1$. If every candidate has some
$i\le K$ for which
\[
 |c_{h,N,z}(i)|\ge N+i+h+2,
\]
then $R_{N+h}-R_N\notin\mathbb Z$.
An integral tail difference would give one of the initial candidates
and would remain in the \emph{open} strip
$|c_{h,N,z}(i)|<N+i+h+2$ at every step. Equality at a boundary already
excludes a candidate. This is the soundness implication for a specified
finite test; no bound for a successful search depth is asserted.

\scale{n/a}\ \ev{Lean}\ \coord{other:carry-orbit-dynamics}\\
\lean{tail_diff_notMem_int_of_survivorKill}{CarrySurvivorExtinction.lean:428}
\end{prop}
\leannote{Lean: \lproof{ErdosProblems/Erdos249/PaperCompleteR20/FiniteCarryCorrespondence.lean}{8}{finite_carry_test_sound}, \lproof{ErdosProblems/Erdos249/PaperCompleteR20/FiniteCarryCorrespondence.lean}{22}{finite_carry_candidate_count}, \lproof{ErdosProblems/Erdos249/PaperCompleteR20/FiniteCarryCorrespondence.lean}{27}{finite_carry_true_orbit}.}

\subsection{Further exact identities}

The identities in this subsection hold independently of the proposed
residue tests. Each gives either an exact expression or a bound that
can be used when studying those tests.

\begin{prop}[The exact totient difference on an LCM progression]\label{prop:AR-04-inv}
Put $H=H(t)$. For every $j\ge0$, the arithmetic expression in the
linked source equals $\delta_t(j)=\varphi(2H+j)-\varphi(H+j)$.
For a divisor $j\mid H$, this assertion contains a useful product
formula, not merely a change of notation. Set
\[
 a=H/j,\qquad g_1=\gcd(j,a+1),\qquad g_2=\gcd(j,2a+1).
\]
Then
\[
 \delta_t(j)=\varphi(j)\left(
 \frac{g_2\varphi(2a+1)}{\varphi(g_2)}
 -\frac{g_1\varphi(a+1)}{\varphi(g_1)}\right).
\]
Here $j\ge1$ because $H>0$, so both denominators are nonzero.
The identity follows from
$\varphi(jx)=\varphi(j)\varphi(x)\gcd(j,x)/\varphi(\gcd(j,x))$,
which retains all primes shared by $j$ and $x$.
If every prime dividing $j$ also divides $a$, then $g_1=g_2=1$,
and the formula reduces to
$\varphi(j)\bigl(\varphi(2a+1)-\varphi(a+1)\bigr)$.
For example, at $H=j=2$ the two gcds are $g_1=2$ and $g_2=1$;
the exact difference is $\varphi(6)-\varphi(4)=0$, whereas discarding
the gcd factors would give $1$.
For $j\nmid H$, the source retains the literal totient difference.
Thus no coprimality assumption is imposed on the full expression.

\scale{uniform}\ \ev{Lean}\ \coord{mobius-mersenne}\\
\lean{lcmRayArithmeticLetter_eq_deltaTotient}{TotientActualLcmOrbitArithmetic.lean:298}
\end{prop}
\leannote{Lean: \lproof{ErdosProblems/Erdos249/PaperCompleteR21/ActualLcmShortWindowArithmetic.lean}{56}{lcmRayArithmeticLetter_eq_totient_difference}, \lproof{ErdosProblems/Erdos249/PaperCompleteR21/ActualLcmShortWindowArithmetic.lean}{70}{totient_mul_eq_totient_mul_gcd_div_totient_gcd}, \lproof{ErdosProblems/Erdos249/PaperCompleteR21/ActualLcmShortWindowArithmetic.lean}{86}{lcmRay_divisor_denominators_pos}, \lproof{ErdosProblems/Erdos249/PaperCompleteR21/ActualLcmShortWindowArithmetic.lean}{104}{lcmRay_divisor_product_formula}, and 3 further declarations in the \hyperref[sec:coverage]{coverage section}.}

\begin{prop}[A lower bound for the totient of a rough integer]\label{prop:AR-03-inv}
Let $a\ge8$, set $t=2^a$, and let $n>0$ be an integer all of
whose prime factors exceed $t$. If $n<2^{2t}$, then $n$ has fewer
than $t/4$ distinct prime factors and
\[
 \varphi(n)>\frac34 n.
\]
If $n$ has $r$ distinct prime factors, then
$n\ge t^r=2^{ar}$, so $ar<2t$ and $r<t/4$ because $a\ge8$.
For the second, use
$\varphi(n)/n=\prod_{p\mid n}(1-1/p)\ge1-\sum_{p\mid n}1/p>3/4$.
These are bounds for rough integers in the stated size range; they are
not pointwise bounds for arbitrary totient values.

\scale{uniform}\ \ev{Lean}\ \coord{mobius-mersenne}\\
\lean{three_quarters_mul_lt_totient_of_rough_lt_two_pow}{TotientActualLcmOrbitArithmetic.lean:768}
\end{prop}
\leannote{Lean: \lproof{ErdosProblems/Erdos249/PaperCompleteR21/ActualLcmShortWindowArithmetic.lean}{215}{rough_integer_prime_count_and_totient_bound}.}

\begin{prop}[Positivity of the short-window differences]\label{prop:AR-05-inv}
For $a\ge8$ and $1\le j<2\cdot2^a$,
\[
 \delta_{2^a}(j)=\varphi(2H_{2^a}+j)-\varphi(H_{2^a}+j)>0.
\]
The proof treats divisor offsets and the exceptional prime-power
offsets separately, using multiplicativity and the rough-integer
estimate of Proposition~\ref{prop:AR-03-inv}. No rationality hypothesis
is used. Positivity of these finitely many coefficients contributes to
the later sign theorem; a bound for the remaining infinite tail is
still needed to obtain that theorem.

\scale{uniform}\ \ev{Lean}\ \coord{mobius-mersenne}\\
\lean{lcmRayArithmeticLetter_pos_of_lt_two_mul}{TotientActualLcmOrbitArithmetic.lean:1703}
\end{prop}
\leannote{Lean: \lproof{ErdosProblems/Erdos249/PaperCompleteR21/ActualLcmShortWindowArithmetic.lean}{229}{shortWindow_totient_difference_pos}.}

\begin{prop}[The accumulated sum is the discrepancy]\label{prop:AR-06-inv}
For all $t,L\in\mathbb N$, the weighted diagonal sum is exactly
\[
 \sum_{j=1}^{L}\delta_t(j)2^{L-j}=D(H(t),H(t),L).
\]
Thus its two-sided residue inequality is precisely
$\mathcal C(H(t),H(t),L)$. The equality follows term by term from the
definition of $\delta_t$. A separate restriction on $L$, when present,
remains part of the hypothesis; this identity does not supply such a
restricted-depth witness.

\scale{n/a}\ \ev{Lean}\ \coord{mobius-mersenne}\\
\lean{lcmDiagonalArithmeticKill_iff_certifiedKill}{TotientActualLcmOrbitArithmetic.lean:2045}
\end{prop}
\leannote{Lean: \lproof{ErdosProblems/Erdos249/PaperCompleteR21/ActualLcmShortWindowArithmetic.lean}{241}{weighted_shortWindow_sum_eq_windowDiscrepancy}, \lproof{ErdosProblems/Erdos249/PaperCompleteR21/ActualLcmShortWindowArithmetic.lean}{262}{weighted_shortWindow_band_iff_certifiedKill}.}

\begin{prop}[A global-to-local identity]\label{prop:SEP-01-inv}
\[
\Omega_a \;=\; 2^H(2^H-1)\Big(\textstyle\sum_n' \varphi(n)/2^n\Big) - \big(\Phi_{2H}-\Phi_H\big),\quad H={H}(2^a).
\]
Thus a separation estimate for $\Omega_a$ is an estimate for this
integer translate of a multiple of $S$. The identity does not require
other rational-approximation arguments to use this particular translate.
\scale{n/a}\ \ev{Lean}\ \coord{mobius-mersenne}\\
\lean{actualLcmTailOrbit_eq_scaled_totientSeries_sub_prefix}{TotientActualLcmOrbitSeparation.lean:68}
\end{prop}
\leannote{Lean: \lproof{ErdosProblems/Erdos249/PaperCompleteR21/ActualLcmSeparationAndSign.lean}{24}{actualLcmTailOrbit_global_to_local}.}

\begin{prop}[A rational approximation with an error bound]\label{prop:SEP-02-inv}
For all $a,q\in\mathbb N$, with $H=H_{2^a}$ and the finite
$\rho_{a,q}$ defined above,
\[
 |\Omega_a-\rho_{a,q}|<\varepsilon_{a,q}
     =\frac{2H+2q+3}{2^{2q+1}}.
\]
The triangle inequality transfers a separation estimate between the
real tail and its rational approximation at the cost of this error.
The bound tends to zero with $q$ for fixed $a$; no assertion of
optimality is made.
\scale{uniform}\ \ev{Lean}\ \coord{mobius-mersenne}\\
\lean{abs_actualLcmTailOrbit_sub_rawApprox_lt}{TotientActualLcmOrbitSeparation.lean:141}
\end{prop}
\leannote{Lean: \lproof{ErdosProblems/Erdos249/PaperCompleteR21/ActualLcmSeparationAndSign.lean}{39}{abs_actualLcmTailOrbit_sub_rawApprox_lt_explicit}, \lproof{ErdosProblems/Erdos249/PaperCompleteR21/ActualLcmSeparationAndSign.lean}{49}{tendsto_actualLcmRawErrorRadius_atTop_nhds_zero}.}

\begin{prop}[Unconditional positivity]\label{prop:SGN-01}
For $a\ge8$, $J+(a+6)<2\cdot2^a$:
\[
0 < R_{2H+J} - R_{H+J},\qquad H={H}(2^a),
\]
with no irrationality hypothesis; the true, infinite, real translated tail difference is
strictly positive throughout almost the entire short window, proved unconditionally from
Prop.~\ref{prop:AR-05-inv} plus a directed one-sided tail bound. Any argument about the sign of
the actual orbit (not merely its residue mod $2^L$) must agree with this; specialised at $J=0$
this gives $0<\Omega_a$.
\scale{uniform}\ \ev{Lean}\ \coord{mobius-mersenne}\\
\lean{actualLcmTailDiff_shift_pos}{TotientActualLcmOrbitSign.lean:39}
\end{prop}
\leannote{Lean: \lproof{ErdosProblems/Erdos249/PaperCompleteR21/ActualLcmSeparationAndSign.lean}{82}{actualLcm_tailDiff_shift_pos_paper}, \lproof{ErdosProblems/Erdos249/PaperCompleteR21/ActualLcmSeparationAndSign.lean}{89}{actualLcmTailOrbit_pos}.}

\begin{prop}[The residue forced by integrality]\label{prop:SGN-03}
Let $a,J,K\in\mathbb N$, $a\ge8$, $H=H(2^a)$, and assume
\[
 J+K+(a+6)<2\cdot2^a,\qquad 2H+J+K+2<2^K.
\]
If $R_{2H+J}-R_{H+J}\in\mathbb Z$, then
\[
 \begin{aligned}
 D(H,H+J,K)\bmod2^K
   &=2^K-\bigl(R_{2H+J+K}-R_{H+J+K}\bigr),\\
 0<R_{2H+J+K}-R_{H+J+K}&<2H+J+K+2.
 \end{aligned}
\]
The later tail difference is a positive integer. Its negative is the
representative near zero, while the least nonnegative residue lies near
$2^K$. Positivity alone therefore does not prove nonintegrality; a
separate estimate must exclude this upper interval.
\par\noindent\ev{Lean} \scale{uniform}
\lean{actualLcm_integral_forces_topEdgeResidue}{TotientActualLcmOrbitSign.lean:211}
\end{prop}
\leannote{Lean: \lproof{ErdosProblems/Erdos249/PaperCompleteR21/ActualLcmSeparationAndSign.lean}{105}{actualLcm_integral_forces_topEdgeResidue_paper}.}


\begin{prop}[The penultimate term of a partial divisibility pattern]\label{prop:TE-02-inv}
Let $a,J,K,m\in\mathbb N$, $a\ge8$, $H=H_{2^a}$ and
$B=2H+J+K+2$. Suppose that
\[
\begin{gathered}
 0<m\le K,\qquad J+K+a+6<2\cdot2^a,\\
 B<2^m,\qquad \delta_{2^a}(J+K)\le2^m-B,\\
 2^{r+1}\mid\delta_{2^a}(J+K-m+r+1)
       \quad(0\le r,\ r+1<m).
\end{gathered}
\]
Then
\[
 \delta_{2^a}(J+K-1)=2^{m-1},\qquad 2^m<2B.
\]
Indeed, $B\ge4$ and $B<2^m$ force $m\ge3$. The penultimate
divisibility, positivity, and the bound
$\delta_{2^a}(J+K-1)<B<2^m$ leave $2^{m-1}$ as its only possible
value. Thus the modulus lies in the strict interval $B<2^m<2B$;
there is at most one possible power of two. The value fixed at
$2^{m-1}$ is the \emph{penultimate} difference. The terminal
difference $\delta_{2^a}(J+K)$ instead satisfies the separate upper
bound in the hypotheses. These restrictions are necessary for the
specified pattern; the proposition does not construct such patterns
at arbitrarily large $a$.
\scale{bounded}\ \ev{Lean}\ \coord{mobius-mersenne}\\
\lean{puncturedDyadicStaircase_penultimate_eq_half}{TotientActualLcmTopEdgeStaircase.lean:1187}
\end{prop}
\leannote{Lean: \lproof{ErdosProblems/Erdos249/PaperCompleteR21/PenultimateStaircaseAndRankCurvature.lean}{37}{penultimate_shortWindow_difference_eq_half}, \lproof{ErdosProblems/Erdos249/PaperCompleteR21/PenultimateStaircaseAndRankCurvature.lean}{80}{dyadicScale_unique_in_open_interval}.}

\begin{prop}[An equivalent test using two residue bits]\label{prop:TE-03-inv}
For every $h,N\in\mathbb N$, existence of a certificate at some
depth is equivalent to the following explicit condition:
\[
 \exists s,b\in\mathbb N,\quad
 \mathcal C(h,N+s,b+1)\quad\text{or}\quad
 \begin{cases}
  N+s+h+b+4<2^b,\\
  2^b\le D(h,N+s,b+2)\bmod2^{b+2}<3\cdot2^b.
 \end{cases}
\]
The last interval says that the two leading bits are $01$ or $10$.
The proof uses the first successful certificate of depth $L$ and may
take $b=\lfloor\log_2(N+h+L+2)\rfloor+1$. Since $b$ depends on that
unknown $L$, the equivalence is not an a priori bound on the search
depth. The proof uses the finite tail recurrence and its size bound;
it supplies no totient-specific assertion that a witness exists.
\scale{uniform}\ \ev{Lean}\ \coord{seam-integer}\\
\lean{exists_certifiedKill_iff_guardCylinderWitness}{TotientActualLcmTopEdgeStaircase.lean:844}
\end{prop}
\leannote{Lean: \lproof{ErdosProblems/Erdos249/PaperCompleteR21/PenultimateStaircaseAndRankCurvature.lean}{129}{exists_certifiedKill_iff_twoBitResidueTest}, \lproof{ErdosProblems/Erdos249/PaperCompleteR21/PenultimateStaircaseAndRankCurvature.lean}{104}{dyadicMixedGuard_iff_twoBitBand}.}

\begin{prop}[The factor in a second difference]\label{prop:FR-02-inv}
Let $a\ge4$ and $j\ge1$ be integers with $j^2\le2^a$, and put
$H_a=H(2^a)=\operatorname{lcm}(1,\ldots,2^a)$. Then
\[
 2\varphi(j)\mid
 \varphi(3H_a+j)-2\varphi(2H_a+j)+\varphi(H_a+j).
\]
To see the factor, set $A=H_a/j$. The square bound ensures that $j\mid H_a$
and that every prime divisor $p$ of $j$ also divides $A$: indeed,
$pj\le j^2\le2^a$ implies $pj\mid H_a$. Consequently
$\gcd(j,qA+1)=1$ for $q=1,2,3$, and the displayed difference equals
\[
 \varphi(j)\bigl(\varphi(3A+1)-2\varphi(2A+1)+\varphi(A+1)\bigr).
\]
Also $2j\le2^a$, so $2j\mid H_a$ and $A\ge2$ is even. The three
arguments $qA+1$ are therefore odd and greater than two, and all three
totients are even. This proves the extra factor of two.
The divisibility is a lower bound, not an exact valuation at each LCM
height. The different family in Observation~\ref{prop:FR-03-kill} does
not satisfy these LCM and square-window hypotheses, so it does not
establish sharpness within this restricted family.
\scale{bounded}\ \ev{Lean}\ \coord{other:fixed-rank-curvature}\\
\lean{two_mul_totient_dvd_fixedRankSecondDifference}{TotientFixedRankLcmAsymptotic.lean:519}
\end{prop}
\leannote{Lean: \lproof{ErdosProblems/Erdos249/PaperCompleteR21/PenultimateStaircaseAndRankCurvature.lean}{219}{two_mul_totient_dvd_totient_second_difference}, \lproof{ErdosProblems/Erdos249/PaperCompleteR21/PenultimateStaircaseAndRankCurvature.lean}{169}{fixedRank_cleanWindow_structure}.}

\begin{prop}[Carry displacement and tail integrality]\label{prop:CP-01-inv}
For a positive integer $v$ and a tempered integral totient carry $u$,
\[
 v\mid u(N+k)-u(N)\quad\Longleftrightarrow\quad R_{N+k}-R_N\in\Z.
\]
Indeed, temperedness identifies $u(N)=vR_N$, so the displacement is
$v(R_{N+k}-R_N)$. This is a divisibility test for that displacement,
not a dimension bound for the carry sections. The same proof works for
another coefficient sequence once its tempered carry has been identified
with its scaled tail; those hypotheses must be checked in each application.
\scale{uniform}\ \ev{Lean}\ \coord{seam-integer}\\
\lean{carryShift_dvd_iff_tailDiff_mem_int}{TotientTailCarryPeriod.lean:77}
\end{prop}
\leannote{Lean: \lproof{ErdosProblems/Erdos249/PaperCompleteR21/TailCarryPeriodAndRankFloor.lean}{31}{carryShift_dvd_iff_tailDiff_integral}, \lproof{ErdosProblems/Erdos249/PaperCompleteR21/TailCarryPeriodAndRankFloor.lean}{20}{temperedCarry_eq_scaledTail_and_shift}.}

\begin{prop}[A consequence of rationality and its limitation]\label{prop:CP-02}
If $S\in\mathbb Q$, there are a positive integer $v$ and a tempered
integer carry orbit $u$ such that its retained sections through every
level $e$ have rational rank at least $2^e-1$. At the same time, those
sections are uniformly eventually periodic modulo $v$.
Periodicity after reduction modulo $v$ is a statement in a finite
quotient; it is not a rank upper bound over $\mathbb Q$.
Proposition~\ref{prop:D5cons} gives the lower bound, while the rational
comparison sequence explains why a general rank upper bound cannot
be deduced from these recurrence assumptions alone.
\scale{uniform}\ \ev{Lean}\ \coord{other:carry-kernel-rank}\\
\lean{not_irrational_totientSeries_implies_mod_period_and_unbounded_rank}{TotientTailCarryPeriod.lean:224}
\end{prop}
\leannote{Lean: \lproof{ErdosProblems/Erdos249/PaperCompleteR21/TailCarryPeriodAndRankFloor.lean}{46}{rationality_forces_mod_period_and_unbounded_rank}.}

\begin{prop}[A totient difference congruent to two modulo four]\label{prop:CP-05-inv}
For every positive integer $h$ and every $B\in\mathbb N$, there is
 a prime $p>B$ such that
\[
 \varphi(p+4h)-\varphi(p)\equiv2\pmod4.
\]
To construct $p$, set $H=4h$, choose a prime $r>H$ with
$r\equiv1\pmod4$, and use Dirichlet's theorem in the progression
$p\equiv3r-H\pmod{4r}$. This progression is reduced: its residue is
odd and is $-H\not\equiv0\pmod r$. Then $p\equiv3\pmod4$, so
$\varphi(p)=p-1\equiv2\pmod4$, whereas $r\mid p+H$ gives
$4\mid\varphi(p+H)$. Subtracting gives the stated congruence.
This supplies the arithmetic congruence at arbitrarily large primes;
the additional finite exclusion in Proposition~\ref{prop:CP-07}
remains a separate hypothesis.

\scale{cofinal}\ \ev{Lean}\ \coord{other:dirichlet-crt}\\
\lean{exists_prime_deltaTotient_four_mul_mod_four_two}{TotientTailCarryPeriod.lean:384}
\end{prop}
\leannote{Lean: \lproof{ErdosProblems/Erdos249/PaperCompleteR21/TwoAdicPulseBlockAndMobiusInversion.lean}{28}{exists_prime_totient_shift_four_mul_congr_two_mod_four}.}

\begin{prop}[An arbitrarily long zero prefix followed by a two-adic pulse]\label{prop:TA-inv}
Let $K\ge2$ and $H>K$ be integers. For every $B\in\mathbb N$ there
is a prime $p>\max(B,H+K)$ such that
\[
\begin{aligned}
 p&\equiv1+2^{K-1}\pmod{2^K},&
 2^K&\mid\varphi(p+H),\\
 2^K&\mid\varphi(p-j),&
 2^K&\mid\varphi(p-j+H)\qquad(1\le j<K).
\end{aligned}
\]
The construction uses $2K-1$ distinct auxiliary primes congruent to
$1\pmod{2^K}$: one for $p+H$, and two for each of the $K-1$
preceding positions. Choose them larger than $\max(H+K,2^K)$.
For the first prime impose $p\equiv-H$; for the pair at position $j$
impose $p\equiv j$ and $p\equiv j-H$, respectively. These nonzero
residues, together with the displayed odd residue modulo $2^K$,
combine by the Chinese remainder theorem into a reduced progression.
Dirichlet's theorem then supplies arbitrarily large prime values of $p$.
Each auxiliary prime contributes the required factor $2^K$ to the
corresponding totient.

Thus the first $K-1$ totient differences in the window starting at
$p-K$ vanish modulo $2^K$, while the last is $2^{K-1}$ modulo $2^K$.
Consequently
\[
 D(H,p-K,K)\equiv2^{K-1}\pmod{2^K}.
\]
If $R_{N+H}-R_N\in\mathbb Z$ for every $N\ge N_0$, choose such a
prime with $p\ge N_0+K$. Iterating the carry recurrence from $p-K$
then gives
\[
 R_{p+H}-R_p\in\mathbb Z,\qquad
 R_{p+H}-R_p\equiv2^{K-1}\pmod{2^K}.
\]
The initial integer is multiplied by $2^K$ and disappears only after
reduction modulo $2^K$. This does not bound its higher quotient.
The condition $H>K$ is part of the construction; it does not provide
arbitrarily large $K$ for one fixed shift $H$.

\scale{cofinal}\ \ev{Lean}\ \coord{p-adic}\\
\lean{exists_prime_totient_twoAdic_pulse_divisors}{TotientTwoAdicPulseBlock.lean:54}
\end{prop}
\leannote{Lean: \lproof{ErdosProblems/Erdos249/PaperCompleteR21/TwoAdicPulseBlockAndMobiusInversion.lean}{81}{exists_prime_twoAdic_pulse_block}, \lproof{ErdosProblems/Erdos249/PaperCompleteR21/TwoAdicPulseBlockAndMobiusInversion.lean}{41}{pulse_delta_of_divisor_data}, \lproof{ErdosProblems/Erdos249/PaperCompleteR21/TwoAdicPulseBlockAndMobiusInversion.lean}{100}{exists_prime_integral_tailDiff_half_pulse}.}

\begin{prop}[A M\"obius-inversion formula for the tail]\label{prop:MP-01-inv}

For $N\in\mathbb N$ and $d\ge1$, let
$r_d(N)=d-(N\bmod d)$ and $q_d(N)=\lfloor N/d\rfloor+1$.
Thus $r_d(N)$ is the distance to the next strictly larger multiple of
$d$, so $1\le r_d(N)\le d$. M\"obius inversion gives
\[
 R_N=\sum_{d\ge1}\mu(d)2^{d-r_d(N)}
 \left(\frac{q_d(N)}{2^d-1}+\frac{1}{(2^d-1)^2}\right).
\]
Indeed, substitute $\varphi(n)=\sum_{d\mid n}\mu(d)n/d$ and group
by $d$. Its multiples after $N$ have quotients $q_d(N)+\ell$ and
forward distances $r_d(N)+d\ell$, for $\ell\ge0$; summing the two
geometric series gives the formula. Absolute convergence justifies the
regrouping, since the sum of the absolute divisor contributions at $n$
is at most $\sum_{d\mid n}n/d\le n^2$, and
$\sum_{n>N}n^2 2^{N-n}<\infty$.
This identity is available for later transformations; it does not
assert that every M\"obius or Lambert-series argument must use it.

\scale{uniform}\ \ev{Lean}\ \coord{cyclotomic}\\
\lean{totientTail_eq_tsum_shiftedMobiusPulse}{TotientShiftedMobiusPulse.lean:547}
\end{prop}
\leannote{Lean: \lproof{ErdosProblems/Erdos249/PaperCompleteR21/TwoAdicPulseBlockAndMobiusInversion.lean}{114}{forwardMultiple_spec}, \lproof{ErdosProblems/Erdos249/PaperCompleteR21/TwoAdicPulseBlockAndMobiusInversion.lean}{125}{forwardMultipleShift_least}, \lproof{ErdosProblems/Erdos249/PaperCompleteR21/TwoAdicPulseBlockAndMobiusInversion.lean}{139}{forwardMultiple_enumeration}, \lproof{ErdosProblems/Erdos249/PaperCompleteR21/TwoAdicPulseBlockAndMobiusInversion.lean}{151}{totientTail_eq_tsum_mobius_inversion}.}

\begin{prop}[Explicit finite examples]\label{prop:SK-01-inv}
The finite certificates
$\mathcal C(H_{16},H_{16},23)$ and
$\mathcal C(H_{64},H_{64},93)$ imply
$\Omega_4,\Omega_6\notin\mathbb Z$. Both satisfy the short-window
restriction, since $23<32$ and $93<128$. They use the pre-existing
certificates \mathtt{certifiedKill\_diagonal\_t16}/\mathtt{t64}.
These are two explicit instances of the desired condition, not the
start of a proved induction or a claim that no other finite instances
can be checked.

\scale{fixed}\ \ev{Cert}\ \coord{mobius-mersenne}\\
\lean{actualLcmTailOrbit_four_and_six_notMem_int}{TotientActualLcmShortKill.lean:35}
\end{prop}
\leannote{Lean: \lproof{ErdosProblems/Erdos249/PaperCompleteR21/RationalTailPeriodWitnesses.lean}{36}{shortWindow_certificates_kill_omega_four_and_six}, \lproof{ErdosProblems/Erdos249/PaperCompleteR21/RationalTailPeriodWitnesses.lean}{49}{shortWindow_inequalities}, \lproof{ErdosProblems/Erdos249/PaperCompleteR21/RationalTailPeriodWitnesses.lean}{27}{actualLcmTailOrbit_eq_tail_difference}.}

\begin{prop}[The necessary depth]\label{prop:A5-inv}
Every certificate $\mathcal C(h,N,L)$ satisfies
\[
 2(N+h+L+2)<2^L.
\]
The two endpoint margins must fit inside an interval of length $2^L$.
In particular, successful depths grow at least logarithmically with
$N+h$. This is a necessary lower bound, not an upper bound for the
first successful depth. The later two-bit reformulation applies at its
stated modulus and does not reverse this implication.

\scale{n/a}\ \ev{Lean}\ \coord{seam-integer}\\
\lean{certifiedKill_depth_floor}{TotientTailPeriodKiller.lean:79}
\end{prop}
\leannote{Lean: \lproof{ErdosProblems/Erdos249/PaperCompleteR20/TailDepthCorrespondence.lean}{20}{certificate_logarithmic_depth}, \lproof{ErdosProblems/Erdos249/PaperCompleteR20/TailDepthCorrespondence.lean}{30}{fixed_depth_bounds_indices}, \lproof{Erdos249257/TotientTailPeriodKiller.lean}{79}{certifiedKill_depth_floor}.}

\begin{prop}[Rationality gives an eventual tail period]\label{prop:A9-inv}
Suppose $S=a/(2^c v)$ in lowest terms, where $a\in\mathbb Z$,
$c\ge0$ and $v\ge1$ is odd. Take $h=\varphi(v)$, including
$h=1$ when $v=1$. Euler's theorem gives $v\mid2^h-1$, so
\[
 2^N(2^h-1)S\in\mathbb Z\qquad(N\ge c).
\]
The identity relating $S$ to its scaled tails therefore gives
$R_{N+h}-R_N\in\mathbb Z$ for every $N\ge c$.
These are explicit witnesses for the eventual period asserted in
Proposition~\ref{prop:A10}. To use the finite residue test against a
hypothetical rational value, one must obtain a nonintegral difference
for its period at an index beyond its preperiod. Soundness proves that
the quantified certificate condition rules out rationality. Conversely,
the prefix-tail identity makes every positive-shift tail difference
nonintegral when $S$ is irrational, and pointwise completeness
(Theorem~\ref{catalogue:cert:a7}) then supplies a depth.
Proposition~\ref{prop:AR-06-inv} only identifies the weighted diagonal
sum with $D$; it is not the completeness argument and supplies no
depth bound.
\scale{uniform}\ \ev{Lean}\ \coord{other:euler-theorem}\\
\lean{eventual_period_of_not_irrational}{TotientTailPeriodKiller.lean:358}
\end{prop}
\leannote{Lean: \lproof{ErdosProblems/Erdos249/PaperCompleteR21/RationalTailPeriodWitnesses.lean}{65}{rational_tail_period_explicit_witnesses}, \lproof{ErdosProblems/Erdos249/PaperCompleteR21/RationalTailPeriodWitnesses.lean}{56}{totient_one_eq_one}, \lproof{ErdosProblems/Erdos249/PaperCompleteR21/RationalTailPeriodWitnesses.lean}{111}{eventual_tail_period_of_not_irrational}, \lproof{ErdosProblems/Erdos249/PaperCompleteR21/RationalTailPeriodWitnesses.lean}{119}{irrational_of_certificate_supply}, and 2 further declarations in the \hyperref[sec:coverage]{coverage section}.}

\begin{prop}[Tail integrality on an LCM grid]\label{prop:B6-inv}
If $S\in\mathbb Q$, there is $t_1\in\mathbb N$ such that for all
$t\ge t_1$, $q\ge1$ and $m\ge0$,
\[
 R_{(q+m)H(t)}-R_{qH(t)}\in\mathbb Z.
\]
The threshold may depend on the hypothetical rational value. A
certificate at one fixed scale therefore gives a finite denominator
exclusion, not irrationality. Proposition~\ref{prop:B7} requires
certificates at arbitrarily large scales, so that one lies beyond the
threshold supplied by any such rational value.

\scale{n/a}\ \ev{Lean}\ \coord{mobius-mersenne}\\
\lean{rational_totient_series_forces_lcm_cone_flatness}{CertificateKernel.lean:19002}
\end{prop}
\leannote{Lean: \lproof{ErdosProblems/Erdos249/PaperCompleteR20/LcmGridCorrespondence.lean}{8}{lcm_grid_flatness}, \lproof{ErdosProblems/Erdos249/PaperCompleteR20/LcmGridCorrespondence.lean}{27}{certificate_denominator_exclusion}, \lproof{ErdosProblems/Erdos249/PaperCompleteR20/LcmGridCorrespondence.lean}{40}{lcm_grid_supply_iff}.}

\begin{prop}[Independence of the retained dyadic family]\label{prop:D4-inv}
For every $e\ge0$, the auxiliary canonical dyadic family has
$2^e+1$ indexed sections, and these sections are linearly independent
over $\mathbb Q$. The CRT--Dirichlet argument is explained in
Section~\ref{sec:mahler-defect}; the declaration below counts its index
set. For $e\ge1$ the family is a basis for all sections through level
$e$. At $e=0$ the auxiliary family still contains both $\varphi(n)$ and
$\varphi(2n)$, so it is not the actual level-zero truncation, whose
dimension is one. No rationality hypothesis on $S$ is used.

\scale{n/a}\ \ev{Lean}\ \coord{other:carry-kernel-rank}\\
\lean{card_totientCanonicalIndex}{TotientMahlerDefect.lean:61}
\end{prop}
\leannote{Lean: \lproof{ErdosProblems/Erdos249/PaperCompleteR21/CanonicalDyadicSectionRank.lean}{21}{card_canonical_dyadic_index}, \lproof{ErdosProblems/Erdos249/PaperCompleteR21/CanonicalDyadicSectionRank.lean}{26}{linearIndependent_canonical_dyadic_family}, \lproof{ErdosProblems/Erdos249/PaperCompleteR21/CanonicalDyadicSectionRank.lean}{33}{canonical_family_basis_through_level}, \lproof{ErdosProblems/Erdos249/PaperCompleteR21/CanonicalDyadicSectionRank.lean}{46}{canonical_level_zero_channels}, and 1 further declaration in the \hyperref[sec:coverage]{coverage section}.}

\begin{prop}[The exact range of the stated Farey-gap inequality]\label{prop:C2-inv}
Let $V$ be the totient residue at $(N,K)=(1,240)$ defined in
Section~\ref{ssec:farey}. For every integer $q$ satisfying
\[
 1\le q\le79639646646701375323355774875831053,
\]
one has
\[
 (qV)\bmod2^{240}+243q<2^{240}.
\]
The first failure occurs at
\[
 q=79639646646701375323355774875831054,
\]
obtained from the stated Farey mediant. Thus the range is exact for
this particular inequality.
It is not an optimality claim among all denominator-exclusion tests.
The preceding range gives the exclusion in
Proposition~\ref{prop:C3-inv}, approximately $7.96\times10^{34}$.
\scale{fixed}\ \ev{Cert}\ \coord{other:farey-gap}\\
\lean{gap_check_window_1_240_first_failure}{GapFareyBound.lean:225}
\end{prop}
\leannote{Lean: \lproof{ErdosProblems/Erdos249/PaperCompleteR21/FareyExactRangeAndDenominatorExclusion.lean}{26}{farey_window_1_240_exact_range}, \lproof{ErdosProblems/Erdos249/PaperCompleteR21/FareyExactRangeAndDenominatorExclusion.lean}{39}{farey_window_1_240_range_magnitude}.}

\begin{prop}[The resulting denominator exclusion]\label{prop:C3-inv}
For every reduced fraction $a/q$, with $a\in\mathbb Z$ and $q\ge1$,
\[
 q\le79639646646701375323355774875831053
 \quad\Longrightarrow\quad S\ne a/q.
\]
This is an unconditional finite denominator exclusion. It does not
assume certificates at arbitrarily large scales, as required by
Proposition~\ref{prop:A10}. Conversely, excluding this finite range
does not rule out rational values with larger reduced denominators.
\scale{fixed}\ \ev{Lean}\ \coord{other:farey-gap}\\
\lean{tsum_totient_div_pow_two_ne_ratCast_of_den_le_79639646646701375323355774875831053}{CertificateKernel.lean:18384}
\end{prop}
\leannote{Lean: \lproof{ErdosProblems/Erdos249/PaperCompleteR21/FareyExactRangeAndDenominatorExclusion.lean}{49}{totient_series_ne_reduced_fraction_of_small_denominator}, \lproof{ErdosProblems/Erdos249/PaperCompleteR21/FareyExactRangeAndDenominatorExclusion.lean}{57}{totient_series_ne_int_div_of_small_denominator}.}

\begin{prop}[Two general irrationality criteria]\label{prop:D1D2-inv}
Let $x\in\mathbb R$. One sufficient condition is a sequence of
reduced fractions $u_j=a_j/q_j$, $q_j\ge1$, with $u_j\ne x$ for
all sufficiently large $j$ and
\[
 q_j|x-u_j|\longrightarrow0.
\]
Another is that, for every integer $Q\ge1$, there exist $m,z\in\mathbb Z$
such that
\[
 0<|mx-z|<1/Q.
\]
Indeed, if $x=a/b$ is reduced with $b>0$, every nonzero error in the
first condition is at least $1/(bq_j)$, and every nonzero error in the
second is at least $1/b$. The strict lower bounds exclude exact
rational hits. Requiring $m=b_0^n$ for a fixed integer base $b_0\ge2$
is a stronger sufficient condition, not a reformulation valid for
every irrational number. The binary example in
Proposition~\ref{catalogue:cert:d2} explains this distinction.
The linked declaration proves the first criterion.
\scale{n/a}\ \ev{Lean}\ \coord{other:dirichlet-approximation}\\
\lean{irrational_of_den_mul_abs_sub_tendsto_zero}{CertificateKernel.lean:5371}
\end{prop}
\leannote{Lean: \lproof{Erdos249257/CertificateKernel.lean}{5371}{irrational_of_den_mul_abs_sub_tendsto_zero}, \lproof{Erdos249257/CertificateKernel.lean}{5342}{one_div_den_mul_den_le_abs_sub}, \lproof{Erdos249257/CertificateKernel.lean}{6120}{irrational_of_int_mul_near_int}, \lproof{Erdos249257/CertificateKernel.lean}{6149}{irrational_of_pow_mul_near_int}, and 8 further declarations in the \hyperref[sec:coverage]{coverage section}.}

\begin{prop}[Lambert identities involving $S$]\label{prop:D7-inv}
For $L(f)=\sum_{n\ge1}f(n)/(2^n-1)$, the identities are
$L(\mu)=1/2$, $L(\varphi)=2$, $L(1)=E$ and $L(\varphi*\mu)=S$.
The last identity rewrites the same unknown value; it does not deduce its
irrationality from that of $E$. The value
$L(\mathrm{Id})=\sum_{m\ge1}\sigma(m)/2^m$ is transcendental by
Nesterenko~\cite[Cor.~2, p.~1320]{nesterenko1996} (cited, not formalised).
These examples show that Lambert-series form alone does not determine
arithmetic status. The accompanying source reference checks the equality
between the two indexing conventions for $S$, not every status claim in
this paragraph; the individual comparisons are in
Proposition~\ref{catalogue:cert:d7} and its table.
\scale{n/a}\ \ev{Lean, Cited}\ \coord{mobius-mersenne}\\
\lean{tsum_totient_div_pow_two_eq_pnat_half_pow}{CertificateKernel.lean:18415}
\end{prop}
\leannote{Lean: \lproof{ErdosProblems/Erdos249/PaperCompleteR21/LambertSigmaRungAndNesterenko.lean}{38}{lambert_ladder_four_values}, \lproof{ErdosProblems/Erdos249/PaperCompleteR21/LambertSigmaRungAndNesterenko.lean}{53}{primWeight_is_totient_conv_moebius}, \lproof{ErdosProblems/Erdos249/PaperCompleteR21/LambertSigmaRungAndNesterenko.lean}{58}{totient_series_index_bridge'}, \lproof{ErdosProblems/Erdos249/PaperCompleteR21/LambertSigmaRungAndNesterenko.lean}{80}{lambert_id_rung_eq_sigma_series}, and 2 further declarations. Conditional on Nesterenko's algebraic independence theorem; see the \hyperref[sec:coverage]{coverage section}.}

\begin{prop}[A general rational gap bound]\label{prop:D9-inv}
If $a/b<c/d$ are reduced fractions with $b,d>0$, then
\[
 \frac{1}{bd}\le\frac cd-\frac ab,
\]
since $bc-ad$ is a positive integer. In particular, if $S=a/b$ and
$p_N/q_N<S$ is a reduced rational approximation with positive error
at most $\varepsilon_N$, then
\[
 b\ge\frac{1}{q_N\varepsilon_N}.
\]
These lower bounds tend to infinity precisely when
$q_N\varepsilon_N\to0$; a subsequence with this property is enough
to contradict a fixed $b$. A tail bound by itself does not imply this. For example, the ordinary dyadic
prefix has denominator dividing $2^N$ and tail at most
$(N+2)2^{-N}$. Using only these two bounds yields
$b\ge1/(N+2)$, which is vacuous. The rational-spacing lemma can
convert sharper error and denominator information into an exclusion;
it does not automatically improve the Farey bound in
Proposition~\ref{prop:C2-inv}.
\scale{n/a}\ \ev{Lean}\ \coord{other:farey-gap}\\
\lean{positive_rational_difference_lower_bound}{PrimitiveRationalGapSupply.lean:31}
\end{prop}
\leannote{Lean: \lproof{ErdosProblems/Erdos249/PaperCompleteR21/GeneralIrrationalityCriteriaAndGapBounds.lean}{138}{rational_gap_lower_bound}, \lproof{ErdosProblems/Erdos249/PaperCompleteR21/GeneralIrrationalityCriteriaAndGapBounds.lean}{146}{den_lower_bound_of_positive_error}, \lproof{ErdosProblems/Erdos249/PaperCompleteR21/GeneralIrrationalityCriteriaAndGapBounds.lean}{164}{denominator_bound_tendsto_atTop_iff}, \lproof{ErdosProblems/Erdos249/PaperCompleteR21/GeneralIrrationalityCriteriaAndGapBounds.lean}{189}{irrational_of_den_mul_error_product_tendsto_zero}, and 4 further declarations in the \hyperref[sec:coverage]{coverage section}.}

\begin{prop}[Diagonal certificates through $82$]\label{prop:B11-inv}
The current aggregate establishes
\[
 \forall t\in\mathbb N,\quad t\le82\ \Longrightarrow\ P\,t,
\]
where $P\,t$ is the diagonal certificate predicate in
Proposition~\ref{prop:B3}. The earlier aggregate covered 28 explicit
cases through $64$, including
$1,2,3,4,5,7,8,9,11,13,16,17$; the later theorem fills the gaps
and extends the range through $82$. Each instance reduces to exact
finite arithmetic with the displayed totient values, using the
prime-power factorisations and Lucas primality certificates in the
source. This proves neither $P\,83$ nor infinitely many instances,
and therefore does not establish the quantified conditions in
Propositions~\ref{prop:A10}, \ref{prop:B3} and \ref{prop:B7}.
\scale{bounded}\ \ev{Cert}\ \coord{mobius-mersenne}\\
\lean{certifiedKill_diagonal_all_imported}{DiagonalPincerCertificates.lean:2870}
\lean{ErdosProblems.Skip.LadderT67.exists_diagonalKill_le_82}{ErdosProblems/Skip/LadderT67.lean:71264}
\end{prop}
\leannote{Lean: \lproof{ErdosProblems/Erdos249/PaperCompleteR20/FiniteCertificateBatch.lean}{57}{historical_table_size_and_initial_depths}, \lproof{ErdosProblems/Erdos249/PaperCompleteR20/FiniteCertificateBatch.lean}{64}{historical_table_and_complete_band}.}

\begin{prop}[A coprime-pair expression for $S$]\label{prop:C1-inv}
For $n\ge0$, the number of integer pairs $a\ge1$, $b\ge0$ with
$a+b=n$ and $\gcd(a,b)=1$ is $\varphi(n)$, with
$\varphi(0)=0$. Indeed, $b=n-a$ and
$\gcd(a,n-a)=\gcd(a,n)$. The boundary case $n=1$ contributes
$(a,b)=(1,0)$; it disappears when both coordinates are required to
be positive. Consequently,
\[
 \sum_{\substack{a,b\ge1\\\gcd(a,b)=1}}2^{-(a+b)}=S-\frac12.
\]
For independent random variables $X,Y$ with
$\Pr(X=n)=\Pr(Y=n)=2^{-n}$, $n\ge1$, the left side is
$\Pr(\gcd(X,Y)=1)$. Thus $S-1/2$ has a coprimality-probability
interpretation. This identity does not supply an estimate for the
binary digits or tail residues used elsewhere in the paper.
\scale{n/a}\ \ev{Lean}\ \coord{other:visible-lattice-points}\\
\lean{tsum_pos_coprime_pair_pow}{GeometricCoprimality.lean:182}
\end{prop}
\leannote{Lean: \lproof{ErdosProblems/Erdos249/PaperCompleteR21/GeneralIrrationalityCriteriaAndGapBounds.lean}{218}{card_visible_antidiagonal}, \lproof{ErdosProblems/Erdos249/PaperCompleteR21/GeneralIrrationalityCriteriaAndGapBounds.lean}{225}{totient_zero_eq_zero}, \lproof{ErdosProblems/Erdos249/PaperCompleteR21/GeneralIrrationalityCriteriaAndGapBounds.lean}{228}{visible_antidiagonal_one}, \lproof{ErdosProblems/Erdos249/PaperCompleteR21/GeneralIrrationalityCriteriaAndGapBounds.lean}{234}{positive_antidiagonal_one}, and 2 further declarations in the \hyperref[sec:coverage]{coverage section}.}

\subsection{Counterexamples and limitations}

Each example below is a counterexample to a stated deduction, not to
irrationality of $S$. The assumptions that fail to suffice are given with
the example, so that its conclusion is not extended to other arguments.


\begin{obs}[The sign of the actual endpoint representative]\label{prop:SGN-02}
Assume $a\ge8$, $H=H_{2^a}$ and
$J+K+a+6<2\cdot2^a$, as in
Proposition~\ref{prop:SGN-01}. If
$d=R_{2H+J}-R_{H+J}$ is an integer, define
\[
 c_0=d,\qquad
 c_{j+1}=2c_j-\bigl(\varphi(2H+J+j+1)-\varphi(H+J+j+1)\bigr).
\]
The tail recurrence gives
$c_K=R_{2H+J+K}-R_{H+J+K}>0$. Its negative is an integer
representative satisfying
\[
 |{-c_K}|\le2H+J+K+2,\qquad
 -c_K\equiv D(H,H+J,K)\pmod{2^K}.
\]
Thus the representative corresponding to the actual tail is negative.
The stated size bound allows it to remain in the endpoint interval,
so its sign alone does not contradict integrality. Other integers in
that interval with the same residue need not represent the actual tail.
Proposition~\ref{prop:SGN-03} gives the further residue separation
that would exclude this case.
\scale{uniform}\ \ev{Lean}\ \coord{mobius-mersenne}\\
\lean{actualLcm_trueEndpointSurvivor_neg}{TotientActualLcmOrbitSign.lean:172}
\end{obs}

\begin{obs}[The complete divisibility pattern is impossible]\label{prop:TE-01-killer}
Let $a\ge8$, $H=H_{2^a}$ and $J,K,m\in\mathbb N$. Assume
\[
 m\le K,\qquad J+K+a+6<2\cdot2^a,\qquad
 2H+J+K+2<2^m.
\]
The simultaneous divisibilities
\[
 2^{r+1}\mid\delta_{2^a}(J+K-m+r+1)\qquad(0\le r<m)
\]
cannot hold. The terminal difference would be positive by
Proposition~\ref{prop:AR-05-inv}, smaller than $2^m$, and divisible by
$2^m$, a contradiction. The obstruction applies to this complete
pattern under the stated size bounds. Omitting its last divisibility
condition leads to the necessary restrictions in
Proposition~\ref{prop:TE-02-inv}; the residue inequality in
Proposition~\ref{prop:TE-04} is another sufficient condition. Neither
statement here supplies such witnesses at arbitrarily large scales.
\scale{uniform}\ \ev{Lean}\ \coord{mobius-mersenne}\\
\lean{not_actualLcmTerminalDyadicStaircase_of_room}{TotientActualLcmTopEdgeStaircase.lean:251}
\end{obs}

\begin{obs}[An exact valuation in a dyadically scaled family]\label{prop:FR-03-kill}
For every integer $n\ge0$, take $H=2^{n+2}$ and $j=3\cdot2^{n+1}$.
Then
\[
 \varphi(3H+j)-2\varphi(2H+j)+\varphi(H+j)=-2^{n+1}.
\]
Indeed, the three arguments are $9\cdot2^{n+1}$,
$7\cdot2^{n+1}$ and $5\cdot2^{n+1}$. Their totients are respectively
$6\cdot2^n$, $6\cdot2^n$ and $4\cdot2^n$, which give the identity.
Thus the difference has $2$-adic valuation exactly $n+1$.
This example limits a valuation bound for arbitrary even heights and
scaled odd offsets. Here $j>H$, whereas the offset in
Proposition~\ref{prop:FR-02-inv} is at most $2^{a/2}$ and the height is
$H(2^a)$. The example therefore does not exclude additional divisibility
on that LCM family.
\scale{uniform}\ \ev{Lean}\ \coord{other:fixed-rank-curvature}\\
\lean{fixedRankSecondDifference_dyadic_fixture_exact}{TotientFixedRankLcmAsymptotic.lean:462}
\end{obs}

\begin{obs}[A rational sequence with matching parity]\label{prop:CM-01}

Start with the coefficient sequence $0,1,1,2,4,4,4,\ldots$, indexed
from $0$. For each $m=2^{k+3}$, $k\ge0$, replace the pair at positions
$m,m+1$ from $(4,4)$ to $(6,0)$, leaving all other coefficients unchanged.
These pairs are disjoint. Each replacement preserves the binary sum because
\[
 2\cdot2^{-m}-4\cdot2^{-(m+1)}=0.
\]
The resulting sequence has the same sum $3/2$ and the same parity at
every index as $\varphi$. It is not eventually periodic, since the gaps
between its occurrences of $6$ tend to infinity. This construction uses
the particular initial parity pattern and the binary weights; it does
not establish an analogous result for arbitrary prescribed parities or
other summand weights.
\scale{n/a}\ \ev{Lean}\ \coord{other:binary-digit}\\
\lean{parityCoboundaryWeight}{TotientParityCoboundaryCountermodel.lean:59}
\end{obs}

\begin{obs}[What the parity example disproves]\label{prop:CM-02}
The sequence $c$ in Observation~\ref{prop:CM-01} satisfies
\[
 0\le c(n)\le\min(n,6),\qquad
 c(n)\equiv\varphi(n)\pmod2,\qquad
 \sum_{n\ge0}c(n)2^{-n}=\frac32,
\]
and is not eventually periodic. More precisely, for every
$N,G,K\in\mathbb N$, it has $K$ of the displayed $(6,0)$ pairs
starting beyond $N$, with successive starting positions separated
by more than $G$. Choose sufficiently late powers of two in the
construction to obtain these pairs.
Hence the growth bound, the exact parity pattern, nonperiodicity and
these arbitrarily separated occurrences are compatible with a rational
sum. They cannot by themselves prove irrationality. The construction
does not preserve the complete totient sequence or its
multiplicativity, and does not specify which additional arithmetic
hypothesis would suffice.
\lean{exists_totientParity_arbitrarilyManySeparatedCarry_rational_countermodel}{TotientParityCoboundaryCountermodel.lean:637}
\end{obs}

\begin{obs}[A Mersenne factor does not force tail-difference divisibility]\label{prop:CP-03-kill}
For the completely multiplicative sequence $c(n)=n$, the scaled tail is
\[
 T_N=\sum_{j\ge1}(N+j)2^{-j}=N+2,
 \qquad T_{N+K}-T_N=K.
\]
If $q>K>0$ and $q\mid2^K-1$, then every such tail difference is
integral but is not divisible by $q$. For example, take $K=3$ and
$q=7$; the multiplicative order of $2$ modulo $7$ is $3$.
Thus even this prime divisor of the homogeneous recurrence multiplier
does not force divisibility of the inhomogeneous tail difference.
An application to the totient sequence needs a further relation between
the prime and the coefficient residues or the tail's size. This
example supplies no general obstruction for Problem~\#257 merely
because that problem also involves Mersenne denominators. Compare the
specific totient construction in Proposition~\ref{prop:CP-05-inv}.
\scale{uniform}\ \ev{Lean}\ \coord{mobius-mersenne}\\
\lean{idCoeff_mersenneModulus_does_not_annihilate_tailShift}{TotientTailCarryPeriod.lean:342}
\end{obs}

\begin{obs}[A lower bound for the stated reconstruction cost]\label{prop:CP-04-kill}
Let $I$ be finite, $w_i\in\mathbb Q$ and $x_i\in\mathbb N$, and
assume $\sum_{i\in I}w_i\varphi(x_i)=1$. Then
\[
 \sum_{i\in I}|w_i|\bigl(2(x_i+1)+(x_i+2)\bigr)\ge3.
\]
Indeed,
$1\le\sum_i|w_i|\varphi(x_i)\le\sum_i|w_i|x_i$,
and the displayed cost is $3\sum_i|w_i|x_i+4\sum_i|w_i|$.
For $x\ge1$, the reconstruction
$\varphi(x)=2R_{x-1}-R_x$ and the separate estimates
$R_M\le M+2$ give exactly this cost. Terms at $x=0$ contribute
nothing to the normalisation, since $\varphi(0)=0$.
This particular termwise triangle bound cannot be smaller than one.
The conclusion is about that bound, not the true combined tail error;
cancellation or a different estimate is not excluded. The same elementary
calculation applies to any nonnegative sequence $c(n)\le n$ under
the corresponding normalisation $\sum_iw_i c(x_i)=1$.
\scale{uniform}\ \ev{Lean}\ \coord{other:adjugate-linear-algebra}\\
\lean{three_le_totientAdjugateTailCost}{TotientTailCarryPeriod.lean:266}
\end{obs}

\begin{obs}[A second-difference test need not use less depth]\label{prop:B9cert-kill}
For $h,N,L\in\mathbb N$, put
\[
 D_2=D(h,N+h,L)-D(h,N,L),\qquad B_2=2(N+2h+L+2).
\]
The second-difference certificate is the finite inequality
\[
 B_2<D_2\bmod2^L<2^L-B_2.
\]
It implies $R_{N+2h}-2R_{N+h}+R_N\notin\mathbb Z$.
The two first-difference truncation errors have bounds
$N+2h+L+2$ and $N+h+L+2$ after multiplication by $2^L$;
their sum is at most the stated $B_2$. This explains the radius
used by the linked test, without claiming it is optimal.
At $(h,N)=(1,8)$, the first-difference test holds at depth $8$,
whereas the second-difference test fails at every depth at most $8$
and first holds at $9$. The earlier 40-cell probe for $t\le20$
reported the first test at least as shallow in 30 cells; that aggregate
is a reported computation, not part of the single-cell theorem linked
below. These comparisons do not give a uniform ordering of the depths.
\scale{n/a}\ \ev{Cert}\ \coord{mobius-mersenne}\\
\lean{totient_tail_rank_two_kill_sound_but_not_shallower_cell}{CertificateKernel.lean:19090}
\end{obs}

\begin{obs}[The prime-power restriction in the unit-gap refinement]\label{prop:D8cert-kill}
The ordinary interval test excludes a rational denominator when its
possible numerator interval contains no integer. Reduced fractions
permit a refinement: every numerator is coprime to its denominator,
so an interval containing only nonunits can also be excluded.
For a candidate $S=a/(2^c u)$ in lowest terms, with $u>0$ odd
and $N\ge c$, the integer $uR_N=2^{N-c}a-u\Phi_N$ is coprime
to $u$. For such an $N$ and any $K\ge0$, define
\[
 \begin{aligned}
 V&=\left(\sum_{r=1}^{K}\varphi(N+r)2^{K-r}\right)\bmod2^K,\\
 uV&=J2^K+W,\qquad 0\le W<2^K,\\
 L&=\left\lfloor\frac{W+u(N+K+2)}{2^K}\right\rfloor.
 \end{aligned}
\]
After subtracting an integer multiple of $u$, the finite prefix
of $uR_N$ is $J+W/2^K$ and its remaining tail lies in
$(0,u(N+K+2)/2^K]$. Its possible integer values are therefore
$J+1,\ldots,J+L$. The refined condition is
$\gcd(J+t,u)>1$ for all $1\le t\le L$, contradicting the required
coprimality of $uR_N$.
For a prime $p$ and $e\ge1$, nonunits modulo $p^e$ are exactly the
multiples of $p$. Since two consecutive integers cannot both be
multiples of $p$, the condition for $u=p^e$ is equivalent to
\[
 L=0\quad\text{or}\quad L=1\text{ and }p\mid J+1.
\]
Thus a \emph{successful} refined certificate at a prime-power denominator
has at most one candidate integer. This is not a bound for an arbitrary
window: $(u,N,K)=(3,3,3)$ gives $V=2$, $J=0$, $W=6$ and $L=3$,
and the test fails because $1,2,3$ are not all nonunits modulo $3$.
The prime-power restriction also matters. With the actual totient
window $(u,N,K)=(15,13,8)$, one gets
\[
 V=148,\quad J=8,\quad W=172,\quad L=2.
\]
Both candidates $9,10$ are nonunits modulo $15$. This interval passes
the refined test but not the empty-interval test. The prime-power bound
therefore neither applies to all composite denominators nor bounds the
possible improvement in the final denominator cutoff.
\scale{bounded}\ \ev{Math; prime-power criterion in Lean}\ \coord{other:farey-gap}\\
\lean{reducedDenominatorCandidateCount_prime_pow_le_one}{PrimitiveWeightCertificate.lean:111}
\end{obs}

\begin{obs}[A synthetic sequence preserving the stated LCM differences]\label{prop:B4b-kill}
Here is the recurrence behind the comparison. Fix $t\ge3$, put
$H=\operatorname{lcm}(1,\ldots,t)$ and $A=\varphi(H)$, and define
\[
 c_i=\begin{cases}-A,&i=(q-1)H\text{ for some }2\le q<t,\\
                   0,&\text{otherwise},\end{cases}
 \qquad a_i=2c_i-c_{i+1}.
\]
The index $i$ here corresponds to the argument $H+i+1$ of a totient
difference. For $2\le q<t$ this gives
\[
 a_{(q-1)H-1}=\varphi((q+1)H)-\varphi(qH)=A.
\]
The equality follows because all prime divisors of $q$ and $q+1$ divide
$H$: hence $\varphi(qH)=qA$ and $\varphi((q+1)H)=(q+1)A$.
Every $a_i$ is divisible by $A$, and $|c_i|\le A$, $|a_i|\le2A$.
For finitely supported integer weights $(\lambda_j)_{j\ge0}$, put
$d_i=\sum_j\lambda_j c_{i+j}$ and $b_i=\sum_j\lambda_j a_{i+j}$.
Linearity and telescoping give
\[
 b_i=2d_i-d_{i+1},\qquad
 \sum_{r=0}^{L-1}2^{L-1-r}b_{i+r}=2^Ld_i-d_{i+L}.
\]
The triangle inequality gives the precise uniform bounds
\[
 |d_i|\le A\sum_j|\lambda_j|,\qquad
 |b_i|\le2A\sum_j|\lambda_j|.
\]
These are uniform in $i$, not in the weights. The factor
$\sum_j|\lambda_j|$ cannot be omitted: taking only $\lambda_0=2$
doubles the nonzero state and coefficient maxima. Divisibility by $A$
and by $\varphi(d)$ for $d\mid H$ is preserved, since
$\varphi(d)\mid\varphi(H)$.

Finite shift polynomials therefore preserve the recurrence, but not the
original numerical bounds independently of their coefficients. This
comparison does not identify $a_i$ with actual totient differences away
from the prescribed indices, or bound the terms discarded by a proposed
projection. The construction and the weight-dependent bounds are in
\lean{lcm_factorIdeal_finiteRank_shiftAlgebra_not_sufficient}{LcmFactorIdealPulseObstruction.lean:798}.
The original source introduction and the decomposition of the actual
totient into divisor contributions remain at
\lean{LcmFactorIdealPulseObstruction.lean}{LcmFactorIdealPulseObstruction.lean:1-24}
and \lean{totient\_eq\_sum\_mobiusTotientChannel}{LcmFactorIdealPulseObstruction.lean:328}.
\scale{uniform}\ \ev{Math; formal source cited}\ \coord{mobius-mersenne}
\end{obs}

\begin{obs}[Compatible carries for fixed-precision local symbols]\label{prop:B5-kill}
For example, prescribing an odd coefficient modulo $8$ permits any
coefficient of the form $a+8z$ with $a$ odd and $z\in\mathbb Z$.
Whatever the incoming integer carry $e$, one can choose $z$ so that
$2e+a+8z$ lies between $-4$ and $3$.
More generally, fix $u\ge1$ and a finite list of integers
$v_j\ge0$ and odd integers $a_j$. Starting from any integer $e_0$,
there are integers $z_j$ and coefficients
\[
 c_j=2^{v_j}(a_j+2^u z_j),\qquad
 e_{j+1}=2e_j+c_j,
 \qquad |e_{j+1}|\le2^{v_j+u-1}.
\]
At each step, reduce $2e_j+2^{v_j}a_j$ to its centred representative
modulo $2^{v_j+u}$, choosing the interval
\[
 [-2^{v_j+u-1},2^{v_j+u-1}).
\]
Solving for $z_j$ gives the stated next carry. Iterate this choice. The bounds apply to the resulting
states, not to the arbitrary initial state $e_0$.
The unrestricted higher quotients $z_j$ are essential: the construction
prescribes a valuation and a finite unit residue, not the full
coefficients or their signs. It therefore does not realise a prescribed
totient-difference sequence. A contradiction using only these local
symbols cannot exclude all compatible integer carries; further
restrictions on the coefficients need a separate argument.
\scale{n/a}\ \ev{Lean}\ \coord{p-adic}\\
\lean{fixedPrecisionTropicalNoGo}{TropicalCurvatureCarry.lean:137}
\end{obs}

\begin{obs}[Correction-free factorisation does not ensure nonvanishing]\label{prop:B8-kill}
Consider the two finite differences
\[
 \varphi(65)-\varphi(65)=0,\qquad
 \varphi(78)-\varphi(70)=24-24=0.
\]
They arise from the square-congruence choice
$52\equiv52\pmod{13^2}$ and $52\equiv2\pmod{5^2}$.
The four arguments factor as $13\cdot5$, $5\cdot13$,
$13\cdot6$ and $5\cdot14$; in each product the prime and cofactor
are coprime. Thus all four totients have the correction-free form
$\varphi(pm)=(p-1)\varphi(m)$, yet both differences vanish.
Nonvanishing is also possible: the choices
$27\equiv0\pmod{3^2}$ and $27\equiv2\pmod{5^2}$ give
\[
 \varphi(30)-\varphi(30)=0,\qquad
 \varphi(33)-\varphi(35)=20-24=-4,
\]
again with coprime factors at every term.
The square-CRT construction removes the non-coprime multiplicativity
correction on the prescribed finite range. That property alone does
not imply nonvanishing. These are explicit witnesses of the two
possibilities, not a claim that the first is a smallest counterexample.
\scale{n/a}\ \ev{Lean (docstring-sourced)}\ \coord{other:square-crt-cocycle}\\
\lean{squareCRT_clean_block_can_vanish}{SquareCRTCube.lean:459}
\end{obs}

\begin{obs}[No fixed integer clears every primitive denominator]\label{prop:D6-kill}
Let $A=\varphi*\mu$ be the Dirichlet convolution used in
Proposition~\ref{prop:D7-inv}. At prime powers,
\[
 A(p)=p-2,\qquad
 A(p^r)=p^{r-2}(p-1)^2\quad(r\ge2).
\]
These follow by subtracting consecutive prime-power values of
$\varphi$. For every odd prime $p$, the reduced denominator of
$A(p)/p$ is $p$. For every prime $p$, the reduced denominator of
$A(p^2)/p^2$ is $p^2$, including the denominator $4$ at $p=2$.
Consequently any positive integer $D$ for which
$DA(n)/n\in\mathbb Z$ for all $1\le n\le N$ must be divisible
by every odd prime $p\le N$, every prime square $p^2\le N$, and
by $4$ once $N\ge4$. No fixed $D$ can do this for all $n$, since
there are infinitely many odd primes.
In periodic-point counting, the quantities $A(n)/n$ would count
orbits of length $n$; their integrality and nonnegativity characterise
realisable periodic-point sequences~\cite[Lemma~3.1, p.~5]{everest-periodic}.
Here only the denominator obstruction is used. It excludes a fixed
integer clearing all these particular coordinates, not every possible
Euler-product representation or an argument with a multiplier growing
with $N$. It is not an irrationality proof for $S$.
\scale{n/a}\ \ev{Lean}\ \coord{mobius-mersenne}\\
\lean{no_fixed_integral_primitive_euler_index}{PrimitiveDeterminantLift.lean:169}
\end{obs}

\begin{obs}[A finite exclusion is not a membership test]\label{prop:B11cert-kill}
A successful residue certificate proves the exclusion stated in its
soundness theorem. Failure of that sufficient condition at finitely
many depths does not prove integrality of the tail difference. In
particular, a finite unsuccessful search cannot be substituted for
the eventual-period assertion needed to prove rationality.
The same logical distinction occurs for the linked greedy test in a
separate Mersenne achievement-set problem: its finite inequality
implies nonmembership. The linked object is a definition, with a
soundness theorem in the same file, not a theorem asserting a universal
impossibility of finite rationality proofs for Problem~\#249.
Membership or rationality established by some additional exact identity
is not excluded by this observation. Likewise, the successful finite
ranges in Propositions~\ref{prop:SK-01-inv} and \ref{prop:B11-inv}
remain finite ranges until an unbounded supply is proved.
\scale{n/a}\ \ev{Math; related Lean definition}\ \coord{n/a}\\
\lean{CertifiedGreedyMersenneDeath}{GreedyAchievementSet.lean:1756}
\end{obs}
% ; - part a249_newdecls ; -
\subsection{Related structural results}

The following results concern the M\"obius-weighted series, integer
carry recurrences, rational approximations, and factorizations used above.
Each statement retains its pinned formal-source reference. The statements
are grouped by their mathematical objects, rather than by the order in
which files were added to the repository.

\subsubsection{Log-concavity of the M\"obius-weighted series}

The series $\Theta_r=\sum_{n\ge0}\mu(n+1)/(2^{n+1}-1)^r$ includes
$S-1/2$ at $r=2$. The next result proves strict log-concavity for every
integer $r\ge1$. This is an order property of these values, not an
irrationality criterion.


\begin{defn}
For each integer $r\ge1$, put
\[
 \Theta_r=\sum_{n\ge0}\frac{\mu(n+1)}{(2^{n+1}-1)^r}
 =1-3^{-r}+\sum_{n\ge0}\frac{\mu(n+3)}{(2^{n+3}-1)^r}.
\]
The first two summands isolate the dominant contribution. The remaining
absolute sum is at most $1/(2^r(2^r-1))$, by comparison with
$\sum_{n\ge2}2^{-rn}$. In particular, $\Theta_r>0$ for $r\ge1$:
$1-3^{-r}\ge2/3$ and this tail bound is at most $1/2$.
\lean{Erdos249257.SignedQMomentObstruction.mobiusMersenneTheta}{SignedQMomentObstruction.lean:136}
\quad \lean{Erdos249257.SignedQMomentObstruction.mobiusMersenneTwoAtom}{SignedQMomentObstruction.lean:140}
\quad \lean{Erdos249257.SignedQMomentObstruction.mobiusMersenneTailAfterTwo}{SignedQMomentObstruction.lean:144}
\quad \ev{Lean} \quad \scale{uniform} \quad \coord{mobius-mersenne}.
\end{defn}

\begin{thm}[The values at exponents one and two]
The first two values are $\Theta_1=1/2$ and $\Theta_2=S-1/2$.
\lean{Erdos249257.SignedQMomentObstruction.mobiusMersenneTheta_one}{SignedQMomentObstruction.lean:197}
\quad \lean{Erdos249257.SignedQMomentObstruction.mobiusMersenneTheta_two_eq_totient_offset}{SignedQMomentObstruction.lean:214}
\quad \ev{Lean} \quad \scale{fixed} \quad \coord{mobius-mersenne}.
\end{thm}
\leannote{Lean: \lproof{ErdosProblems/Erdos249/PaperCompleteR21/MobiusMersenneLadderLogConcavity.lean}{23}{mobiusMersenneTheta_one_and_two}.}

\noindent The second equality is the M\"obius-square identity credited in
\S\ref{ssec:coprime}; whether $\Theta_2$ is rational is \#249 itself.

\begin{thm}[The first two summands give a positive Hankel gap]
$(1-3^{-(r+1)})^2 - (1-3^{-r})(1-3^{-(r+2)}) = 4/3^{r+2}$,
so the rational sequence $1-3^{-r}$ is strictly log-concave for every integer $r\ge1$.
\lean{Erdos249257.SignedQMomentObstruction.mobiusMersenneTwoAtom_hankelGap}{SignedQMomentObstruction.lean:311}
\quad \ev{Lean} \quad \scale{uniform} \quad \coord{mobius-mersenne}.
\end{thm}
\leannote{Lean: \lproof{ErdosProblems/Erdos249/PaperCompleteR21/MobiusMersenneLadderLogConcavity.lean}{35}{twoAtom_hankel_gap}, \lproof{ErdosProblems/Erdos249/PaperCompleteR21/MobiusMersenneLadderLogConcavity.lean}{43}{twoAtom_strict_logConcave}.}

\begin{thm}[Strict log-concavity for all integer $r\ge1$]
For every integer $r\ge1$,
\[
 \Theta_r\Theta_{r+2}<\Theta_{r+1}^2.
\]
Equivalently, the determinant of
$\left(\begin{smallmatrix}\Theta_r&\Theta_{r+1}\\
\Theta_{r+1}&\Theta_{r+2}\end{smallmatrix}\right)$ is negative.
The two-summand comparison gives an elementary proof, detailed below;
the linked declarations prove the same strict inequality.
\lean{Erdos249257.SignedQMomentObstruction.mobiusMersenneTheta_strict_logConcave}{SignedQMomentObstruction.lean:730}
\quad \lean{Erdos249257.SignedQMomentObstruction.mobiusMersenneTheta_strict_logConcave_of_five_le}{SignedQMomentObstruction.lean:534}
\quad \lean{Erdos249257.SignedQMomentObstruction.mobiusMersenneTheta_hankel_two_neg}{SignedQMomentObstruction.lean:740}
\quad \ev{Lean} \quad \scale{uniform} \quad \coord{mobius-mersenne}.
\end{thm}
\leannote{Lean: \lproof{ErdosProblems/Erdos249/PaperCompleteR21/MobiusMersenneLadderLogConcavity.lean}{52}{theta_strict_logConcave}, \lproof{ErdosProblems/Erdos249/PaperCompleteR21/MobiusMersenneLadderLogConcavity.lean}{58}{theta_hankel_det_neg}, \lproof{ErdosProblems/Erdos249/PaperCompleteR21/MobiusMersenneLadderLogConcavity.lean}{66}{theta_hankel_two_neg}.}
\begin{proof}[A quantitative comparison]
For this proof put $a_r=1-3^{-r}$ and $\varepsilon_r=1/(2^r(2^r-1))$.
The preceding bound gives $|\Theta_r-a_r|\le\varepsilon_r$.
Since $0<a_r<1$, expansion of the two products gives
\[
 \Theta_{r+1}^2-\Theta_r\Theta_{r+2}
 \ge \frac4{3^{r+2}}-B_r,
 \qquad
 B_r=2\varepsilon_{r+1}+\varepsilon_{r+1}^2
       +\varepsilon_r+\varepsilon_{r+2}
       +\varepsilon_r\varepsilon_{r+2}.
\]
Here $\varepsilon_{r+1}<\varepsilon_r/4$, so $B_{r+1}<B_r/4$.
The rational comparison
\[
 B_5=\frac{50106593}{32001951744}<\frac4{3^7}
\]
therefore proves the inequality for all $r\ge5$: the error decreases
faster than the positive main term.
For $r=1,2,3,4$, truncate each $\Theta_j$ at $n=8$. Its omitted
absolute sum is at most
\[
 \frac{1}{(2^9-1)^j(1-2^{-j})}.
\]
Using $\Theta_1=1/2$ and these rational intervals, the four gaps
$\Theta_{r+1}^2-\Theta_r\Theta_{r+2}$ are respectively greater than
$1/4$, $1/16$, $1/64$ and $1/200$. These finite rational comparisons
complete the ordinary proof.
\end{proof}

\begin{obs}
The inequality excludes a representation
$\Theta_r=\int_0^\infty x^r\,d\nu(x)$ by a positive measure with these
finite moments: Cauchy--Schwarz would give
$\Theta_{r+1}^2\le\Theta_r\Theta_{r+2}$. It does not by itself imply
irrationality of $\Theta_2$. Indeed, the rational sequence $1-3^{-r}$
satisfies the same strict log-concavity inequalities. Thus the determinant
sign supplies an obstruction to a positive-moment representation, not a
rational-approximation estimate.
\end{obs}

The accompanying formal results include a rectangular determinant
expansion derived from the Leibniz formula
(\lean{Erdos249257.SignedQMomentObstruction.det_mul_rectangular}{SignedQMomentObstruction.lean:29},
\ev{Lean}, \scale{uniform}, \coord{other:hankel-determinant}), and the
following elementary parity fact: in a finite sum of dyadic fractions,
if the largest denominator exponent occurs once and its numerator is odd,
clearing that denominator gives an odd numerator
(\lean{Erdos249257.SignedQMomentObstruction.scaled_dyadic_sum_odd}{SignedQMomentObstruction.lean:78},
\ev{Lean}, \scale{uniform}, \coord{other:dyadic-parity}).

\subsubsection{Integer recurrences and rational binary series}

For a binary series with integer coefficients, rationality can be
expressed by an integer recurrence whose growth is smaller than $2^N$.
The next equivalence is elementary algebra and convergence; no novelty is
claimed for it.


\begin{thm}[Integer recurrence criterion for a binary series]
\label{thm:binary-carry-criterion}
Let $c:\N\to\N$ satisfy $c(n)\le n$. Then
$\sum_{n\ge1}c(n)2^{-n}$ is rational if and only if there are an integer
$v>0$ and an integer sequence $u$ such that
\[
 u(N+1)=2u(N)-v\,c(N+1),\qquad \frac{u(N)}{2^N}\longrightarrow0.
\]
For each fixed $v$, such a sequence is unique and equals
\[
 u(N)=vR^c_N,\qquad R^c_N=\sum_{j\ge1}\frac{c(N+j)}{2^j}.
\]
\lean{Erdos249257.binaryCoeffSeries_rational_iff_exists_temperedBinaryOrbit}{GenericTailOrbitRigidity.lean:426}
\quad \lean{Erdos249257.temperedBinaryOrbit_eq_scaledTail}{GenericTailOrbitRigidity.lean:339}
\quad \ev{Lean} \quad \scale{uniform} \quad \coord{other:tail-orbit-rigidity}.
\end{thm}

\begin{proof}
Iterating the recurrence gives the finite identity
\[
 \frac{u(N)}{2^N}=u(0)-v\sum_{n=1}^N\frac{c(n)}{2^n}.
\]
The boundary condition therefore forces the series to equal $u(0)/v$;
subtracting the prefix gives $u(N)=vR^c_N$. Conversely, when the series
is rational, choose an integer $v>0$ making its product with $v$ an
integer, use this product as $u(0)$, and iterate the recurrence.
The displayed identity identifies the resulting integer sequence with
$vR^c_N$. Since $0\le R^c_N\le N+2$, its ratio to $2^N$ tends to zero.
\end{proof}

Without the boundary condition, any integer $u(0)$ gives an integer
solution of the recurrence, whether the series is rational or irrational.
Two solutions with the same $v$ differ by $2^N$ times their initial
difference. Thus the limit, not the recurrence alone, determines the
initial value. Changing $v$ rescales the canonical solution; uniqueness
is asserted only for a fixed multiplier.

\begin{prop}[Uniqueness under the growth condition]
A real sequence $d$ with $d(N{+}1)=2d(N)$ and $d(N)=o(2^N)$ is identically
zero.
\lean{Erdos249257.doublingOrbit_eq_zero_of_tempered}{GenericTailOrbitRigidity.lean:315}
\quad \ev{Lean} \quad \scale{uniform} \quad \coord{other:tail-orbit-rigidity}.
\end{prop}
\leannote{Lean: \lproof{ErdosProblems/Erdos249/PaperCompleteR21/TemperedOrbitAndSquaredMersenneTail.lean}{28}{doubling_tempered_sequence_eq_zero}.}

\begin{obs}[The specified finite-state decoder is impossible]
Fix $m\ge2$, put $R=\lfloor(m+1)/2\rfloor$, and for
$0\le r\le R$ define
\[
 c_{m,r}(n)=
 \begin{cases}
  R-r,&n=m,\\
  2r,&n=m+1,\\
  0,&\text{otherwise}.
 \end{cases}
\]
The two nonzero terms have total contribution
$(R-r)2^{-m}+2r2^{-(m+1)}=R2^{-m}$, independent of $r$.
All coefficients and scaled tails before $m$ are therefore the same,
whereas the scaled tail at $m$ is $r$.
No map from a state common to this whole family can recover that tail.
More generally, exact recovery of all its possible values requires at
least $R+1$ distinct labels.
\lean{Erdos249257.balancedPulse_no_autonomous_decoder}{GenericTailOrbitRigidity.lean:247}
\quad \lean{Erdos249257.balancedPulse_label_card_lower_bound}{GenericTailOrbitRigidity.lean:228}
\quad \ev{Lean} \quad \scale{uniform} \quad \coord{other:tail-orbit-rigidity}.
The sequences in this family have identical histories before the isolated
change but different subsequent tails. Therefore that finite history alone
cannot determine the exact tail in this class of sequences. This is
not a conclusion about the totient sequence.
\end{obs}

There is also an infinite family with a common rational value.
Let $a:\N\to\{0,1\}$ satisfy $a(0)=0$, and put
$c(2k)=k-a(k)$ and $c(2k+1)=2a(k)$. The adjacent terms contribute
$k4^{-k}$, so
\[
 \sum_{n\ge1}\frac{c(n)}{2^n}
   =\sum_{k\ge1}\frac{k}{4^k}=\frac49,
\]
regardless of the binary sequence $a$. The cited result also gives an
integer recurrence satisfying the growth condition
(\lean{Erdos249257.globalBalancedCoeff_value}{GenericTailOrbitRigidity.lean:545},
\ev{Lean}, \scale{uniform}, \coord{other:tail-orbit-rigidity}).
The assumptions $a(0)=0$ and $a(k)\in\{0,1\}$ are part of that
formal statement, not an arbitrary bounded-input hypothesis.
For comparison, the identity coefficients $c(n)=n$ have the explicit
multiplier-one recurrence $u(N)=N+2$
(\lean{Erdos249257.idCoeff_temperedOrbit}{GenericTailOrbitRigidity.lean:564},
\ev{Lean}, \scale{fixed}, \coord{other:tail-orbit-rigidity}).

\subsubsection{Rational approximation with a squared-Mersenne remainder}

The first-order Lambert term in the diagonal tail difference can be
summed exactly. The approximation below leaves only a squared-Mersenne
tail to estimate.


\begin{prop}[An exact rational approximation formula]
For integers $H,D\ge0$, with the integer prefixes $\Phi_N$ defined above,
\[
\begin{aligned}
 R_{2H}-R_H
 &-\left[\Phi_H-\Phi_{2H}
   +2^H(2^H-1)\left(\frac12+\sum_{d=1}^{D}\frac{\mu(d)}{(2^d-1)^2}\right)\right]\\
 &=2^H(2^H-1)\sum_{d>D}\frac{\mu(d)}{(2^d-1)^2}.
\end{aligned}
\]
This follows by substituting the squared-denominator identity for $S$
into $R_{2H}-R_H=2^H(2^H-1)S+\Phi_H-\Phi_{2H}$. There is no required
ordering between $D$ and $H$.
\lean{Erdos249257.SquaredMersenneDiagonalEnclosure.scaleDiagonalTailDifference_sub_lambertProjectedDiagonal}{SquaredMersenneDiagonalEnclosure.lean:138}
\quad \ev{Lean} \quad \scale{uniform} \quad \coord{other:squared-mersenne-tail}.
\end{prop}
\leannote{Lean: \lproof{ErdosProblems/Erdos249/PaperCompleteR21/TemperedOrbitAndSquaredMersenneTail.lean}{118}{tailDifference_sub_rationalApproximation}, \lproof{ErdosProblems/Erdos249/PaperCompleteR21/TemperedOrbitAndSquaredMersenneTail.lean}{88}{tailDifference_eq_coefficient_mul_series}, \lproof{ErdosProblems/Erdos249/PaperCompleteR21/TemperedOrbitAndSquaredMersenneTail.lean}{99}{totientSeries_pnat_form}, \lproof{ErdosProblems/Erdos249/PaperCompleteR21/TemperedOrbitAndSquaredMersenneTail.lean}{106}{totientSeries_eq_half_add_moebius_sq}.}

\begin{prop}[A geometric tail bound]
For every integer $D\ge0$,
\[
 \left|\sum_{d>D}\frac{\mu(d)}{(2^d-1)^2}\right|
 \le \frac{4}{3(2^{D+1}-1)^2}.
\]
Indeed, $|\mu(d)|\le1$ and
$2^{D+1+j}-1\ge2^j(2^{D+1}-1)$ for $j\ge0$. The sum of the resulting
geometric majorant is $\sum_{j\ge0}4^{-j}=4/3$.
\lean{Erdos249257.SquaredMersenneDiagonalEnclosure.abs_mobiusSquareTail_le}{SquaredMersenneDiagonalEnclosure.lean:167}
\quad \ev{Lean} \quad \scale{uniform} \quad \coord{other:squared-mersenne-tail}.
\end{prop}
\leannote{Lean: \lproof{ErdosProblems/Erdos249/PaperCompleteR21/TemperedOrbitAndSquaredMersenneTail.lean}{155}{abs_mobiusSquareTail_le_paper}, \lproof{ErdosProblems/Erdos249/PaperCompleteR21/TemperedOrbitAndSquaredMersenneTail.lean}{138}{abs_moebius_cast_le_one}, \lproof{ErdosProblems/Erdos249/PaperCompleteR21/TemperedOrbitAndSquaredMersenneTail.lean}{144}{mersenne_geometric_shift}, \lproof{ErdosProblems/Erdos249/PaperCompleteR21/TemperedOrbitAndSquaredMersenneTail.lean}{149}{tsum_quarter_geometric}.}

\begin{obs}[One-sided and symmetric error bounds]
Suppose $n>D$, $\mu(d)=0$ for $D<d<n$, and $|\mu(n)|=1$. Removing
that first nonzero summand gives the sharper enclosure
\[
 \left|\mu(n)\sum_{d>D}\frac{\mu(d)}{(2^d-1)^2}
                 -\frac1{(2^n-1)^2}\right|
       \le\frac4{3(2^{n+1}-1)^2}.
\]
The radius is smaller than the positive centre, so the tail has sign
$\mu(n)$. The interval's width, divided by the width of the preceding
symmetric enclosure, is
\[
 \left(\frac{2^{D+1}-1}{2^{n+1}-1}\right)^2<\frac14.
\]
The cited theorem gives this signed enclosure
(\lean{Erdos249257.SquaredMersenneDiagonalEnclosure.mobiusSquareTail_directed_of_first_support}{SquaredMersenneDiagonalEnclosure.lean:313},
\ev{Lean}, \scale{uniform}, \coord{other:squared-mersenne-tail}).
Multiplication by $2^H(2^H-1)$ and translation by the rational
approximation above enclose $R_{2H}-R_H$. If the resulting interval
contains no integer, the tail difference is nonintegral. For the symmetric
bound, the linked statement requires the distance from its rational centre
to every integer to exceed its error radius
(\lean{Erdos249257.SquaredMersenneDiagonalEnclosure.scaleFullTarget_miss_of_lambert_projected_separation}{SquaredMersenneDiagonalEnclosure.lean:426}).
The enclosure does not establish that strict separation. In particular,
shrinking an interval around a rational approximation does not, without a
denominator comparison, show that the interval misses every integer.
\end{obs}

\subsubsection{Divisor terms and their complement}

At $H=6$ and $s=2$, the difference
$\varphi(2H+s)-\varphi(H+s)=\varphi(14)-\varphi(8)$ is $2$, whereas
the terms with divisor index $d\mid H$ sum to $3$. The other indices
contribute $-1$. The following decomposition keeps both contributions.
It holds for every positive integer $H$, not only a least common multiple;
an estimate is needed before its second sum can be discarded.


\begin{prop}[The sum over divisor indices]
For integers $H>0$ and $s\ge0$, the part of the M\"obius expansion
indexed by divisors of $H$ is
\[
 \sum_{\substack{d\mid H\\d\mid s}}\mu(d)\frac Hd
       =H\frac{\varphi(\gcd(H,s))}{\gcd(H,s)}.
\]
This is the standard identity
$\sum_{d\mid g}\mu(d)/d=\varphi(g)/g$ at $g=\gcd(H,s)>0$.
The formal statement writes the same value after extracting
$H/\operatorname{rad}(H)$, where
$\operatorname{rad}(H)=\prod_{p\mid H}p$.
\lean{Erdos249257.DiagonalFreshLossBridge.oldMobiusIncrement_eq_scale_gcdWordCoeff}{DiagonalFreshLossBridge.lean:653}
\quad \ev{Lean} \quad \scale{uniform} \quad \coord{other:gcd-word}.
\end{prop}
\leannote{Lean: \lproof{ErdosProblems/Erdos249/PaperCompleteR21/DivisorChannelSplitAndSeamDoubling.lean}{68}{sum_divisorIndices_mobius}, \lproof{ErdosProblems/Erdos249/PaperCompleteR21/DivisorChannelSplitAndSeamDoubling.lean}{98}{sum_divisors_moebius_div_eq_totient_div}, \lproof{ErdosProblems/Erdos249/PaperCompleteR21/DivisorChannelSplitAndSeamDoubling.lean}{121}{sum_divisorIndices_radical_form}, \lproof{ErdosProblems/Erdos249/PaperCompleteR21/DivisorChannelSplitAndSeamDoubling.lean}{117}{squarefreeKernel_eq_prod_primeFactors}.}

\begin{prop}[The divisor sum and its complement]
For integers $H>0$ and $s\ge0$,
\[
\begin{aligned}
 \varphi(2H+s)-\varphi(H+s)
 &=H\frac{\varphi(\gcd(H,s))}{\gcd(H,s)}\\
 &\quad+\sum_{\substack{1\le d\le2H+s\\d\nmid H}}
 \mu(d)\left(\frac{2H+s}{d}\,\mathbf1_{d\mid2H+s}
           -\frac{H+s}{d}\,\mathbf1_{d\mid H+s}\right).
\end{aligned}
\]
Here $\mathbf1$ is the indicator of the stated divisibility condition.
The identity follows by expanding both totients as
$\varphi(n)=\sum_{d\mid n}\mu(d)n/d$ and separating the indices
$d\mid H$. Such an index divides either endpoint precisely when it
divides $s$, and its difference is then $\mu(d)H/d$.
\lean{Erdos249257.DiagonalFreshLossBridge.diagonalHeightIncrement_eq_oldMobius_add_finiteForeign}{DiagonalFreshLossBridge.lean:829}
\quad \ev{Lean} \quad \scale{uniform} \quad \coord{other:gcd-word}.
\end{prop}
\leannote{Lean: \lproof{ErdosProblems/Erdos249/PaperCompleteR21/DivisorChannelSplitAndSeamDoubling.lean}{217}{totientDifference_eq_divisorPart_add_complement}, \lproof{ErdosProblems/Erdos249/PaperCompleteR21/DivisorChannelSplitAndSeamDoubling.lean}{134}{totient_eq_mobius_divisor_sum}, \lproof{ErdosProblems/Erdos249/PaperCompleteR21/DivisorChannelSplitAndSeamDoubling.lean}{141}{divisorIndex_endpoint_behaviour}, \lproof{ErdosProblems/Erdos249/PaperCompleteR21/DivisorChannelSplitAndSeamDoubling.lean}{159}{complementSummand_eq_phaseTerm}.}

\begin{thm}[Doubling the full totient difference]
For nonnegative integers $H,r$ with $H$ even,
\[
 \varphi(4H+2r)-\varphi(2H+2r)=
 \begin{cases}
 2\bigl(\varphi(2H+r)-\varphi(H+r)\bigr),&r\text{ even},\\
 \varphi(2H+r)-\varphi(H+r),&r\text{ odd}.
 \end{cases}
\]
This applies to the full difference, not just its divisor contribution.
It follows from $\varphi(2n)=2\varphi(n)$ for even $n$ and
$\varphi(2n)=\varphi(n)$ for odd $n$.
\lean{Erdos249257.DiagonalFreshLossBridge.diagonalHeightIncrement_two_mul_even}{DiagonalFreshLossBridge.lean:308}
\quad \lean{Erdos249257.DiagonalFreshLossBridge.diagonalHeightIncrement_two_mul_odd}{DiagonalFreshLossBridge.lean:323}
\quad \ev{Lean} \quad \scale{uniform} \quad \coord{other:gcd-word}.
\end{thm}
\leannote{Lean: \lproof{ErdosProblems/Erdos249/PaperCompleteR21/DivisorChannelSplitAndSeamDoubling.lean}{262}{totientDifference_doubling_seam}, \lproof{ErdosProblems/Erdos249/PaperCompleteR21/DivisorChannelSplitAndSeamDoubling.lean}{250}{totient_two_mul_even}, \lproof{ErdosProblems/Erdos249/PaperCompleteR21/DivisorChannelSplitAndSeamDoubling.lean}{255}{totient_two_mul_odd}.}

\begin{obs}[Squarefree support of the complementary terms]
A nonzero summand in the complementary sum must have a squarefree
index $d$ dividing exactly one of $2H+s$ and $H+s$. Indeed, otherwise
$\mu(d)=0$ or both indicators vanish. Both endpoints cannot be divisible
by $d$: their difference is $H$, and $d\nmid H$.
\lean{Erdos249257.DiagonalFreshLossBridge.squarefree_of_foreignChannelPhaseTerm_ne_zero}{DiagonalFreshLossBridge.lean:909}
\quad \lean{Erdos249257.DiagonalFreshLossBridge.foreign_endpoint_support_disjoint}{DiagonalFreshLossBridge.lean:852}
\quad \ev{Lean} \quad \scale{uniform} \quad \coord{other:gcd-word}.
The restriction identifies which indices can contribute. It does not
control the sign or magnitude of their sum.
\end{obs}

\begin{prop}[A doubling identity for two portions of the sum]
Let $H>0$, $s\ge0$ and $d\ge1$, with $d\nmid H$ and $d\mid H+s$.
The $d$-summands in the complementary sums at offsets $s$ and $2s$ are,
respectively,
\[
 -\mu(d)\frac{H+s}{d}
 \qquad\text{and}\qquad
 2\mu(d)\frac{H+s}{d}.
\]
For the second expression, $d\mid2H+2s$ and $d\nmid H+2s$; otherwise
$d$ would divide their difference $H$. Thus the second contribution is
$-2$ times the first, including the zero case $\mu(d)=0$. The two
terms have different binary weights in a window sum, so this identity
alone is not a cancellation of their weighted contributions.
\lean{Erdos249257.DiagonalFreshLossBridge.foreignChannelPhaseTerm_low_double_echo}{DiagonalFreshLossBridge.lean:941}
\quad \ev{Lean} \quad \scale{uniform} \quad \coord{other:gcd-word}.
\end{prop}
\leannote{Lean: \lproof{ErdosProblems/Erdos249/PaperCompleteR21/DivisorChannelSplitAndSeamDoubling.lean}{289}{complementSummand_low_double_echo}.}

\subsubsection{A numerator polynomial with explicit positive coefficients}
\label{repunit-gcd-word}

\begin{thm}[Positive coefficients of the numerator polynomial]
For a squarefree integer $r\ge1$, the numerator polynomial has the
explicit expression
\[
 \sum_{d\mid r}\mu(d)\frac rd\sum_{j=0}^{r/d-1}X^{dj}
   =\sum_{k=0}^{r-1}
        \frac r{\gcd(r,k)}\varphi(\gcd(r,k))X^k.
\]
Every coefficient for $0\le k<r$ is positive, and the higher coefficients
are zero. Extracting $X^k$ on the left gives
$\sum_{d\mid\gcd(r,k)}\mu(d)r/d$, which proves the formula.
The squarefree assumption identifies this divisor sum with the
subset-of-primes definition in the formal source.
\lean{Erdos249257.RepunitMobiusNumerator.mobiusNumeratorPolynomial_eq_gcdWord}{RepunitMobiusNumerator.lean:205}
\quad \lean{Erdos249257.RepunitMobiusNumerator.gcdWordCoeff}{RepunitMobiusNumerator.lean:41}
\quad \ev{Lean} \quad \scale{uniform} \quad \coord{other:gcd-word}.
\end{thm}
\leannote{Lean: \lproof{ErdosProblems/Erdos249/PaperCompleteR21/NumeratorPolynomialAndMersenneRemainder.lean}{67}{paperNumerator_eq_gcdWordForm}, \lproof{ErdosProblems/Erdos249/PaperCompleteR21/NumeratorPolynomialAndMersenneRemainder.lean}{82}{paperNumerator_coeff_pos}, \lproof{ErdosProblems/Erdos249/PaperCompleteR21/NumeratorPolynomialAndMersenneRemainder.lean}{88}{paperNumerator_coeff_eq_zero}, \lproof{ErdosProblems/Erdos249/PaperCompleteR21/NumeratorPolynomialAndMersenneRemainder.lean}{94}{paperNumerator_coeff_eq_gcd_divisor_sum}, and 1 further declaration in the \hyperref[sec:coverage]{coverage section}.}

\begin{cor}
For the same squarefree $r$, evaluation at $X=2$ gives the integer
\[
 \sum_{d\mid r}\mu(d)\frac rd\frac{2^r-1}{2^d-1}.
\]
Each quotient is an integer because $d\mid r$.
\lean{Erdos249257.RepunitMobiusNumerator.mobiusNumeratorPolynomial_eval_two}{RepunitMobiusNumerator.lean:445}
\quad \ev{Lean} \quad \scale{uniform} \quad \coord{other:gcd-word}.
\end{cor}
\leannote{Lean: \lproof{ErdosProblems/Erdos249/PaperCompleteR21/NumeratorPolynomialAndMersenneRemainder.lean}{111}{paperNumerator_eval_two}, \lproof{ErdosProblems/Erdos249/PaperCompleteR21/NumeratorPolynomialAndMersenneRemainder.lean}{106}{mersenne_dvd_of_dvd}, \lproof{ErdosProblems/Erdos249/PaperCompleteR21/NumeratorPolynomialAndMersenneRemainder.lean}{145}{paperNumerator_eval_two_primeSubsetForm}.}

\subsubsection{Exact dyadic comparisons and tail bounds}

The next two formulas concern exact comparisons with a dyadic rational
and an explicit tail bound. They are useful when a residue is computed
using a finite prefix: the error must remain smaller than its distance
from the boundary.


\begin{defn}
For $p\in\mathbb Z$, $n\ge0$ and an integer $L>0$, the difference
between $p/(2L)$ and the dyadic point $2^{-(n+1)}$ is
\[
 \frac{p}{2L}-2^{-(n+1)}=\frac{2^np-L}{2^{n+1}L}.
\]
Its sign is therefore the sign of the integer numerator $2^np-L$.
For fixed $p,L$, the next numerator satisfies
$2^{n+1}p-L=2(2^np-L)+L$.
\lean{Erdos249257.nextDyadicExcessIntNumerator}{DyadicPrefixCompression.lean:156}
\quad \lean{Erdos249257.nextDyadicExcessIntNumerator_succ}{DyadicPrefixCompression.lean:161}
\quad \ev{Lean} \quad \scale{uniform} \quad \coord{other:dyadic-carry}.
\end{defn}

\begin{prop}[The geometric remainder bound]
The remainder after the first two geometric terms of $1/(2^n-1)$ is
\[
 \frac1{2^n-1}-2^{-n}-4^{-n}
   =\frac{8^{-n}}{1-2^{-n}}
   \le \frac43\,8^{-n}\qquad(n\ge2).
\]
Indeed, $1-2^{-n}\ge3/4$. Summing over $n>m$, for an integer $m\ge1$, gives
\[
 \sum_{n>m}\left(\frac1{2^n-1}-2^{-n}-4^{-n}\right)
 \le\frac43\sum_{n>m}8^{-n}=\frac4{21}\,8^{-m}.
\]
\lean{Erdos249257.mersenneWeightRemainder_le_four_thirds}{DyadicPrefixCompression.lean:1162}
\quad \lean{Erdos249257.mersenneWeightRemainderTail_le_four_twentyone}{DyadicPrefixCompression.lean:1179}
\quad \ev{Lean} \quad \scale{uniform} \quad \coord{other:dyadic-carry}.
\end{prop}
\leannote{Lean: \lproof{ErdosProblems/Erdos249/PaperCompleteR21/NumeratorPolynomialAndMersenneRemainder.lean}{166}{mersenneRemainder_identity_and_bound}, \lproof{ErdosProblems/Erdos249/PaperCompleteR21/NumeratorPolynomialAndMersenneRemainder.lean}{156}{one_sub_half_pow_ge}, \lproof{ErdosProblems/Erdos249/PaperCompleteR21/NumeratorPolynomialAndMersenneRemainder.lean}{208}{mersenneRemainderTail_le}, \lproof{ErdosProblems/Erdos249/PaperCompleteR21/NumeratorPolynomialAndMersenneRemainder.lean}{182}{tsum_eighth_pow_tail}, and 1 further declaration in the \hyperref[sec:coverage]{coverage section}.}

\subsubsection{Squared distances between complex phases}

For a finite set of unit complex phases, distance from $1$ is related
exactly to the real part of their sum. A second elementary identity bounds
the total squared distance between pairs. These facts turn suitable
separation estimates into the real-part inequality; neither supplies that
separation for the totient phases.


\begin{prop}[Squared distance from the phase one]
\[
\begin{aligned}
\sum_{N\in T}\|{E}(h,N,L)-1\|^2
  &= 2|T|-2\sum_{N\in T}\operatorname{Re}E(h,N,L).
\end{aligned}
\]
\lean{Erdos249257.TotientTailPeriodKiller.firstHarmonicAnchorDefect_eq_sum_norm_sq}{PivotAntiReconstruction.lean:562}
\quad \lean{Erdos249257.TotientTailPeriodKiller.firstHarmonicAnchorDefect_eq}{PivotAntiReconstruction.lean:550}
\quad \ev{Lean} \quad \scale{uniform} \quad \coord{other:pivot-fiber}.
\end{prop}
\leannote{Lean: \lproof{ErdosProblems/Erdos249/PaperCompleteR21/PhaseEnergyAndForeignResidueProjection.lean}{31}{sum_sq_dist_from_phase_one}.}

\begin{lem}[Squared-distance bound for separated pairs]
For a finite family $z:T\to\mathbb{C}$, a real number $\delta\ge0$, and any set of pairs $P\subseteq
T\times T$ each separated by $\ge\delta$, $|P|\cdot\delta^2 \le
\sum_{i,j\in T}\|z_i-z_j\|^2$.
\lean{Erdos249257.TotientTailPeriodKiller.card_mul_sq_le_sum_pairwise_norm_sq_of_separatedPairs}{PivotAntiReconstruction.lean:720}
\quad \ev{Lean} \quad \scale{uniform} \quad \coord{other:pivot-fiber}.
\end{lem}
\leannote{Lean: \lproof{ErdosProblems/Erdos249/PaperCompleteR21/PhaseEnergyAndForeignResidueProjection.lean}{40}{card_mul_sq_le_pairwise_energy}.}

\subsubsection{Approximating the nondivisor terms}\label{sec:nondivisor-enclosure}

The terms with $d\mid H$ form a finite rational contribution to
$R_{2H}-R_H$. To use it, we must also approximate the terms with
$d\nmid H$. The approximation and its error are explicit; for a cutoff
$D\ge2H$, only separation of the resulting rational centre from the
integers remains to be checked.

For $d\ge1$ and $N\ge0$, let $a_d(N)=d-(N\bmod d)$ be the least
positive shift from $N$ to a multiple of $d$. Summing the multiples of
$d$ in the totient tail gives the rational term
\[
 \kappa_d(N)=\mu(d)2^{d-a_d(N)}
 \left(\frac{N+a_d(N)}{d(2^d-1)}+\frac1{(2^d-1)^2}\right),
 \qquad R_N=\sum_{d\ge1}\kappa_d(N).
\]
Indeed, substitute $\varphi(n)=\sum_{d\mid n}\mu(d)n/d$ into $R_N$.
For a fixed $d$, the admissible shifts are $a_d(N)+kd$, $k\ge0$;
their contribution is
$\mu(d)\sum_{k\ge0}((N+a_d(N))/d+k)2^{-a_d(N)-kd}$,
which is the displayed expression. The interchange of sums is justified by
\[
 \sum_{j\ge1}2^{-j}\sum_{d\mid N+j}|\mu(d)|\frac{N+j}{d}
 \le \sum_{j\ge1}(N+j)^2 2^{-j}<\infty.
\]
Thus subtraction and finite splitting of these series are valid.
The terms are signed: for example, $\kappa_2(0)=-4/9$.
This is the kernel identity already present in the cited source
(\lean{Erdos249257.TotientShiftedMobiusForeignBridge.totientTail_eq_tsum_positiveForeignResidueKernel}{TotientShiftedMobiusForeignBridge.lean:61}).

\begin{prop}[The finite divisor sum]
For an integer $H>0$, the contribution of the divisor indices to
$R_{2H}-R_H$ is
\[
 A_H=H\sum_{d\mid H}\frac{\mu(d)}{d(2^d-1)}.
\]
For $d\mid H$, both $a_d(H)$ and $a_d(2H)$ equal $d$. Hence
$\kappa_d(2H)-\kappa_d(H)=H\mu(d)/(d(2^d-1))$.
Summing over the divisors gives $A_H$, not the whole tail difference:
the nondivisor terms must still be included.
\lean{Erdos249257.ActualForeignResidueProjection.scaleExplicitShadow_eq_sum_divisors_residueIncrement}{ActualForeignResidueProjection.lean:338}
\quad \ev{Lean} \quad \scale{uniform} \quad \coord{other:foreign-residue-projection}.
\end{prop}
\leannote{Lean: \lproof{ErdosProblems/Erdos249/PaperCompleteR21/PhaseEnergyAndForeignResidueProjection.lean}{69}{divisorChannels_sum_eq}, \lproof{ErdosProblems/Erdos249/PaperCompleteR21/PhaseEnergyAndForeignResidueProjection.lean}{50}{residueOffset_of_dvd}, \lproof{ErdosProblems/Erdos249/PaperCompleteR21/PhaseEnergyAndForeignResidueProjection.lean}{58}{residueKernel_increment_of_dvd}, \lproof{ErdosProblems/Erdos249/PaperCompleteR21/PhaseEnergyAndForeignResidueProjection.lean}{84}{scaleExplicitShadow_eq_divisorChannels}.}

For $H>0$ and an integer cutoff $D\ge0$, set
\[
\begin{aligned}
 P_{H,D}&=\sum_{\substack{1\le d\le D\\d\nmid H}}
                    \bigl(\kappa_d(2H)-\kappa_d(H)\bigr),\\
 \varepsilon_{H,D}
   &=2^H(2^H-1)\left(\frac2{2^D}+\frac4{3\cdot4^D}\right).
\end{aligned}
\]
Here $A_H$ includes \emph{all} divisor indices, even when $D<H$.
The following implication therefore keeps the error estimate as an
explicit hypothesis; it will be proved below when $D\ge2H$.

\begin{prop}[Separation larger than the error implies exclusion]
Suppose
\[
 \bigl|(R_{2H}-R_H)-(A_H+P_{H,D})\bigr|\le\varepsilon_{H,D},
 \qquad
 |A_H+P_{H,D}-z|>\varepsilon_{H,D}\quad\hbox{for every }z\in\Z.
\]
Then $R_{2H}-R_H\notin\Z$.
\lean{Erdos249257.ActualForeignResidueProjection.scaleFullTarget_miss_of_projected_separation}{ActualForeignResidueProjection.lean:414}
\quad \ev{Lean} \quad \scale{uniform} \quad \coord{other:foreign-residue-projection}.
\end{prop}
\leannote{Lean: \lproof{ErdosProblems/Erdos249/PaperCompleteR21/PhaseEnergyAndForeignResidueProjection.lean}{119}{tailDifference_not_integral_of_separation}, \lproof{ErdosProblems/Erdos249/PaperCompleteR21/PhaseEnergyAndForeignResidueProjection.lean}{103}{projectedForeignDefect_paper}, \lproof{ErdosProblems/Erdos249/PaperCompleteR21/PhaseEnergyAndForeignResidueProjection.lean}{110}{foreignComplementBound_paper}.}
\begin{proof}
If $R_{2H}-R_H=z\in\Z$, the two hypotheses give opposite inequalities
for $|A_H+P_{H,D}-z|$.
\end{proof}

\paragraph{Why the error hypothesis holds for $D\ge2H$.}
Every divisor of $H$ then lies below the cutoff. For $d>D$, both
$H$ and $2H$ are smaller than $d$, so the omitted increment simplifies to
\[
 \kappa_d(2H)-\kappa_d(H)
 =\mu(d)2^H(2^H-1)
       \left(\frac1{2^d-1}+\frac1{(2^d-1)^2}\right).
\]
Since $|\mu(d)|\le1$ and $2^d-1\ge2^{d-1}$,
\[
\begin{aligned}
 \bigl|(R_{2H}-R_H)-(A_H+P_{H,D})\bigr|
 &\le 2^H(2^H-1)\sum_{d>D}\left(\frac2{2^d}+\frac4{4^d}\right)\\
 &=\varepsilon_{H,D}.
\end{aligned}
\]
Thus for $H>0$ and $D\ge2H$ the error bound needs no additional
convergence hypothesis. The same estimate is proved in the pinned source
(\lean{Erdos249257.TotientShiftedMobiusForeignBridge.controlledForeignProjection_of_twice_le}{TotientShiftedMobiusForeignBridge.lean:465});
together with the displayed separation inequality, it gives the cited
nonintegrality consequence
(\lean{Erdos249257.TotientShiftedMobiusForeignBridge.scaleFullTarget_miss_of_projected_separation_of_twice_le_of_forall_int}{TotientShiftedMobiusForeignBridge.lean:483}).

For example, at $H=1$ and $D=5$ the rational centre and radius are
\[
 A_1+P_{1,5}=\frac{304282}{423801},\qquad
 \varepsilon_{1,5}=\frac{49}{384},\qquad
 \frac12<A_1+P_{1,5}-\varepsilon_{1,5}
          <A_1+P_{1,5}+\varepsilon_{1,5}<1.
\]
Thus this finite calculation excludes an integral value of $R_2-R_1$.
The cutoff restriction alone does not ensure success: at $H=1,D=4$,
the centre $346/441$ is within the radius $25/96$ of $1$.
For fixed $H$ the radius tends to zero as $D\to\infty$, but an
irrationality argument still needs the required family of separated
centres. Neither a single example nor convergence of the error supplies it.

When $D\ge2H$, the first-order part of the omitted sum can also be
included exactly:
$\sum_{d>D}\mu(d)/(2^d-1)=1/2-\sum_{d\le D}\mu(d)/(2^d-1)$.
Adding $2^H(2^H-1)$ times this rational correction to $A_H+P_{H,D}$
gives the squared-Mersenne approximation above
(\lean{Erdos249257.TotientShiftedMobiusForeignBridge.renormalizedResidueAgreement_of_twice_le}{TotientShiftedMobiusForeignBridge.lean:338}). This explains the
relation between the two approximations and the improved squared-denominator
remainder; it does not provide their integer separation.

\subsubsection{Coprime-pair sums}

Counting coprime pairs gives another expression for the totient
series. The adjacent Lambert identity is classical and uses different
weights. Its sum is an elementary rational function of the parameter;
it does not decide irrationality of $S$.


\begin{thm}[Coprime-pair counting and the totient]
For every $n\in\mathbb N$,
$\#\{(a,b)\in\mathbb N^2:a+b=n,\ a>0,\ \gcd(a,b)=1\}=\varphi(n)$.
Here $b=0$ is allowed: the boundary pair $(1,0)$ accounts for the value
$\varphi(1)=1$.
\lean{GeometricCoprimality.card_antidiagonal_filter_pos_coprime}{GeometricCoprimality.lean:56}
\quad \ev{Lean} \quad \scale{uniform} \quad \coord{other:coprime-lattice}.
\end{thm}
\leannote{Lean: \lproof{ErdosProblems/Erdos249/PaperCompleteR21/PhaseEnergyAndForeignResidueProjection.lean}{147}{card_coprime_antidiagonal}, \lproof{ErdosProblems/Erdos249/PaperCompleteR21/PhaseEnergyAndForeignResidueProjection.lean}{154}{boundary_pair_at_one}.}

\begin{prop}[Two lattice sums]
For $0\le r<1$, the two choices of boundary give
\[
\begin{aligned}
 \sum_{\substack{a\ge1,\ b\ge0\\\gcd(a,b)=1}}r^{a+b}
   &=\sum_{n\ge1}\varphi(n)r^n,\\
 \sum_{\substack{a,b\ge1\\\gcd(a,b)=1}}r^{a+b}
   &=\sum_{n\ge1}\varphi(n)r^n-r.
\end{aligned}
\]
The removed pair is $(1,0)$. Partitioning all strictly positive pairs
by their greatest common divisor gives
\[
 \sum_{g\ge1}\ \sum_{\substack{a,b\ge1\\\gcd(a,b)=1}}
       r^{g(a+b)}=\left(\frac r{1-r}\right)^2.
\]
This total is $1$ exactly when $r=1/2$.
\lean{GeometricCoprimality.tsum_coprime_pair_pow_eq_tsum_totient_mul_pow}{GeometricCoprimality.lean:120}
\quad \lean{GeometricCoprimality.tsum_pos_coprime_pair_pow}{GeometricCoprimality.lean:182}
\quad \lean{GeometricCoprimality.tsum_gcd_layer_pos_coprime_pow}{GeometricCoprimality.lean:328}
\quad \ev{Lean} \quad \scale{uniform} \quad \coord{other:coprime-lattice}.
\end{prop}
\leannote{Lean: \lproof{ErdosProblems/Erdos249/PaperCompleteR21/CoprimeLatticeSumsAndLambert.lean}{27}{tsum_totient_pow_shift}, \lproof{ErdosProblems/Erdos249/PaperCompleteR21/CoprimeLatticeSumsAndLambert.lean}{38}{coprimeLattice_halfOpen_sum}, \lproof{ErdosProblems/Erdos249/PaperCompleteR21/CoprimeLatticeSumsAndLambert.lean}{46}{coprimeLattice_positive_sum}, \lproof{ErdosProblems/Erdos249/PaperCompleteR21/CoprimeLatticeSumsAndLambert.lean}{53}{coprimeLattice_removed_pair}, and 2 further declarations in the \hyperref[sec:coverage]{coverage section}.}

\begin{thm}[The classical coprime-pair Lambert identity]
For every $0\le r<1$, $\sum_{(a,b)\ \mathrm{coprime},\, a,b\ge 1}
\dfrac{r^{a+b}}{1-r^{a+b}} = \Bigl(\dfrac{r}{1-r}\Bigr)^2$, an elementary
rational function of $r$, hence rational at every rational $r$ including
$r=1/2$. This is the classical visible-point identity, and the Lean
declaration is a formalisation of it rather than a new result: writing each
pair $(A,B)$ of positive integers uniquely as $g\cdot(a,b)$ with
$\gcd(a,b)=1$ converts the quadrant sum $\sum_{A,B\ge1}r^{A+B}=(r/(1-r))^2$
into the displayed sum over visible points. With the plain weight
$r^{a+b}$, the same strictly positive index set instead sums to
$\sum_{n\ge1}\varphi(n)r^n-r$. Thus at $r=1/2$ the Lambert-weighted
sum is $1$, whereas the plain-weight sum is $S-1/2$, not $S$.
Adding the boundary pair $(1,0)$ recovers $S$ in the plain-weight sum.
\lean{GeometricCoprimality.tsum_pos_coprime_lambert_eq_sq}{GeometricCoprimality.lean:407}
\quad \ev{Lean} \quad \scale{uniform} \quad \coord{other:coprime-lattice}.
\end{thm}
\leannote{Lean: \lproof{ErdosProblems/Erdos249/PaperCompleteR21/CoprimeLatticeSumsAndLambert.lean}{106}{coprimeLattice_lambert_identity}, \lproof{ErdosProblems/Erdos249/PaperCompleteR21/CoprimeLatticeSumsAndLambert.lean}{114}{coprimeLattice_lambert_rational}, \lproof{ErdosProblems/Erdos249/PaperCompleteR21/CoprimeLatticeSumsAndLambert.lean}{127}{coprimeLattice_lambert_half_eq_one}, \lproof{ErdosProblems/Erdos249/PaperCompleteR21/CoprimeLatticeSumsAndLambert.lean}{46}{coprimeLattice_positive_sum}, and 2 further declarations in the \hyperref[sec:coverage]{coverage section}.}

\begin{obs}
The Lambert-weighted sum is rational at every rational $r\in[0,1)$.
Consequently, coprimality of the indices, the decomposition by greatest
common divisor and geometric decay alone do not imply its irrationality.
This comparison does not decide the arithmetic nature of the differently
weighted sum $S-1/2$, or of $S$. A proof using the plain-weight sum must
exploit information not contained in those three properties.
\end{obs}

\subsubsection{The single-parameter diagonal condition}

The period can be chosen to be a least common multiple. This reduces
the required certificate condition to a single scale $t$. The reduction
does not remove the demand for arbitrarily large scales, which is the
remaining arithmetic assertion.


\begin{obs}[The diagonal implication]
The irrationality of $S$ follows if, for every integer $t_0\ge0$, there
are $t\ge t_0$ and $L\ge0$ such that $\mathcal C(H_t,H_t,L)$, where
$H_t=\operatorname{lcm}(1,\ldots,t)$ and $H_0=1$. The declaration below
gives this implication. Theorem~\ref{catalogue:cert:b3} also records the
converse. The quantified condition is therefore equivalent to
irrationality, not a weaker hypothesis.
\lean{Erdos249257.TotientTailPeriodKiller.irrational_totient_series_of_lcm_diagonal_certificate_supply}{LcmDiagonalReduction.lean:112}
\quad \ev{Lean} (implication proved; antecedent \ev{Open})
\quad \scale{cofinal} \quad \coord{other:lcm-diagonal}.
\end{obs}

\begin{thm}[Nondivisors in a short LCM window]
Let $t\ge1$. If $1\le j<2t$ and $j\nmid H_t$, then
$j=p^a>t$ for a prime $p$ and an integer $a\ge1$. Every integer
$1\le j\le t$ divides $H_t$.
To see the first claim, some prime-power divisor $p^a$ of $j$ exceeds $t$;
otherwise every prime-power divisor would divide $H_t$. Since
$j<2t<2p^a$, its remaining cofactor is $1$.

If $j\mid H_t$ and every prime divisor of $j$ also divides $H_t/j$, then,
for every integer $q\ge0$,
\[
 \varphi(qH_t+j)=\varphi(j)\varphi\bigl(q(H_t/j)+1\bigr).
\]
Indeed, the second factor's argument is coprime to $j$, so totient
multiplicativity applies. The hypothesis is essential to this
factorisation: at $t=2$, $j=2$, $q=1$ the left side is $\varphi(4)=2$
but the displayed product would be $1$.
\lean{Erdos249257.TotientTailPeriodKiller.eq_prime_pow_of_not_dvd_periodLcm}{LcmDiagonalReduction.lean:137}
\quad \lean{Erdos249257.TotientTailPeriodKiller.totient_periodLcm_ray_split}{LcmDiagonalReduction.lean:208}
\quad \ev{Lean} \quad \scale{uniform} \quad \coord{other:lcm-diagonal}.
\end{thm}
\leannote{Lean: \lproof{ErdosProblems/Erdos249/PaperCompleteR20/LcmGridCorrespondence.lean}{57}{short_lcm_window_nondivisor}, \lproof{ErdosProblems/Erdos249/PaperCompleteR20/LcmGridCorrespondence.lean}{67}{clean_lcm_ray_factorisation}, \lproof{ErdosProblems/Erdos249/PaperCompleteR20/LcmGridCorrespondence.lean}{76}{unclean_lcm_ray_counterexample}, \lproof{Erdos249257/CarrySurvivorExtinction.lean}{525}{dvd_periodLcm}.}

\begin{obs}[Finite diagonal examples]
The cited file proves $\mathcal C(H_t,H_t,L)$ at the following depths:
\[
 \begin{array}{c|rrrrrrrr}
 t&1&2&3&4&5&6&7&8\\
 H_t&1&2&6&12&60&60&420&840\\
 L&6&5&7&7&9&9&14&15
 \end{array}
\]
The largest totient argument is $2H_6+9=129$ for $t\le6$, and
$2H_8+15=1695$ for the full table.
\lean{Erdos249257.TotientTailPeriodKiller.certifiedKill_periodLcm_diagonal_upto_six}{LcmDiagonalReduction.lean:257}
\quad \lean{Erdos249257.TotientTailPeriodKiller.certifiedKill_periodLcm_diagonal_at_eight}{LcmDiagonalReduction.lean:266}
\quad \ev{Cert} \quad \scale{fixed} \quad \coord{other:lcm-diagonal}; a
finite verification, not a proof that such certificates occur at
arbitrarily large scales.
\end{obs}

\subsubsection{Prime divisors of Mersenne factors and finite exclusions}

For the polynomial $X-2$, elementary order arguments give unbounded
prime support in the Mersenne factors. This settles the particular abstract
prime-support hypothesis used below. It does not give the tail-residue
inequality needed for irrationality. The separate finite exclusions retain
their stated periods and basepoints.


\begin{thm}[Unbounded prime support in Mersenne factors]
Every prime divisor $p$ of $2^q-1$, for prime $q$, satisfies $q\mid p-1$
(the order of $2$ mod $p$ is exactly $q$, by Fermat/Lagrange in
$(\mathbb{Z}/p)^\times$); consequently the prime divisors appearing in the
layers $\{2^n-1\}$ are unbounded, unconditionally, with no cyclotomic
resultant hypothesis left open.
\lean{ErdosProblems.Erdos249.CyclotomicAnchoredKill.prime_index_dvd_pred}{ErdosProblems/Erdos249/CyclotomicAnchoredKill.lean:33}
\quad \lean{ErdosProblems.Erdos249.CyclotomicAnchoredKill.mersenneLayer_unboundedPrimeDivisorSupply}{ErdosProblems/Erdos249/CyclotomicAnchoredKill.lean:105}
\quad \ev{Lean} \quad \scale{uniform} \quad \coord{other:cyclotomic-anchor}.
\end{thm}
\leannote{Lean: \lproof{ErdosProblems/Erdos249/PaperCompleteR21/MersennePrimeSupportAnchors.lean}{33}{mersenneLayer_prime_divisor_order}, \lproof{ErdosProblems/Erdos249/PaperCompleteR21/MersennePrimeSupportAnchors.lean}{63}{mersenneLayer_unbounded_prime_support}.}

\begin{thm}[A prime satisfying the stated cyclotomic conditions]
For every period $h>0$ and threshold $N_0$, there exist a prime $q$ and a
prime factor $p$ of $|\Phi_{hq}(2)|$ (the binary cyclotomic layer) with $p$
coprime to $hq$, $hq\mid p-1$, and $p-1\ge N_0$. The characteristic-prime
exceptional case in the cyclotomic order decomposition is eliminated
directly, by choosing $q>2^h$ (rules out $p=q$) and $q>h$ (rules out
$p\mid h$).
\lean{ErdosProblems.Erdos249.CyclotomicAnchoredKill.exists_clean_binaryCyclotomicAnchor}{ErdosProblems/Erdos249/CyclotomicAnchoredKill.lean:1607}
\quad \lean{ErdosProblems.Erdos249.CyclotomicAnchoredKill.binaryCyclotomicLayer_prime_order_decomposition}{ErdosProblems/Erdos249/CyclotomicAnchoredKill.lean:1431}
\quad \ev{Lean} \quad \scale{uniform} \quad \coord{other:cyclotomic-anchor}.
\end{thm}
\leannote{Lean: \lproof{ErdosProblems/Erdos249/PaperCompleteR21/MersennePrimeSupportAnchors.lean}{73}{exists_clean_cyclotomic_anchor_paper}, \lproof{ErdosProblems/Erdos249/PaperCompleteR21/MersennePrimeSupportAnchors.lean}{84}{cyclotomic_layer_prime_order_decomposition_paper}.}

\begin{obs}
The first theorem uses the multiplicative order of $2$ modulo a prime.
For the second, choose a prime $q>\max(h,N_0,2^h)$ and then a prime
divisor of $|\Phi_{hq}(2)|>1$. The cyclotomic order formula, together
with the stated inequalities on $q$, excludes prime divisors of $hq$
and gives $\operatorname{ord}_p(2)=hq$. Existence of arbitrarily large
primes suffices here; no prime-in-progression theorem is needed for this
construction.

These results supply prime divisors with the required orders. They do
not supply the tail-residue inequality at the resulting periods. The
all-prime and eventual forms of the order hypothesis must also be
distinguished, as explained in \S\ref{prime-ray-curvature}.
\end{obs}

\begin{obs}[Finite exclusions at period 30]
The kernel decides two independent certificates at the composite cyclotomic
anchor period $h=30$ (natural basepoint $N=300$, and the prime-anchored
basepoint $N=330$). The certificate at $N=300$ gives the denominator exclusion
$S\ne r$ for every $r\in\mathbb{Q}$ with $r.\mathrm{den}\mid 2^{300}(2^{30}-1)$.
\lean{ErdosProblems.Erdos249.CyclotomicAnchoredKill.certifiedKill_cyclotomic_30_331_natural}{ErdosProblems/Erdos249/CyclotomicAnchoredKill.lean:3356}
\quad \lean{ErdosProblems.Erdos249.CyclotomicAnchoredKill.totient_series_ne_rat_of_den_dvd_30_300}{ErdosProblems/Erdos249/CyclotomicAnchoredKill.lean:3377}
\quad \ev{Cert} \quad \scale{fixed} \quad \coord{other:cyclotomic-anchor}; a
finite exclusion, not a proof of the corresponding condition at
arbitrarily large scales.
\end{obs}

\subsubsection{A sufficient hypothesis for unbounded prime support}
\label{prime-ray-curvature}

The residue group matters in this criterion. The bounded exponent $k$
in $p^k\equiv1\pmod{mq}$ is a finite-field extension degree, or
equivalently a bound for the order of $p$ modulo $mq$. It is not a bound
for the order of $2$ modulo $p$ in the binary examples. The displayed
condition forces a prime divisor of $C(mq)$ to grow with $q$. A separate
assumption that these values exceed $1$ ensures that prime divisors exist.

\begin{thm}[A sufficient order hypothesis for unbounded prime support]
Let $C:\N\to\N$, and fix integers $m\ge1$ and $d\ge0$. Suppose that
for every pair of primes $q,p$ with $p\mid C(mq)$ there is an integer
$k$ such that
\[
 1\le k\le d,\qquad mq\mid p^k-1.
\]
The divisibility forces $\gcd(p,mq)=1$. For such $p$, existence of
$k$ in the displayed range is equivalent to
$\operatorname{ord}_{mq}(p)\le d$. It also gives $mq<p^d$.
Hence every fixed finite set of primes is disjoint from
the prime divisors of $C(mq)$ for all sufficiently large prime $q$.
If, in addition, for all sufficiently large prime $q$ one has
$C(mq)>1$ and $\gcd(C(mq),mq)=1$, then for every $B,N_0$ there are
primes $q\ge N_0$ and $p>B$ with $p\mid C(mq)$.

Indeed, the divisibility gives $mq\le p^k-1<p^k\le p^d$.
For a finite set of primes, take $q$ larger than all their $d$-th powers;
for the last assertion, take a prime divisor of the nontrivial value
$C(mq)$ after this threshold. The stated coprimality is included in the
formal source's layer hypothesis, although this last extraction argument
uses only $C(mq)>1$.
\lean{ErdosProblems.Erdos249.PrimeRayCyclotomicCurvature.finitePrimeSupportEscape_of_orderConsumer}{ErdosProblems/Erdos249/PrimeRayCyclotomicCurvature.lean:165}
\quad \lean{ErdosProblems.Erdos249.PrimeRayCyclotomicCurvature.unboundedPrimeDivisorSupply_of_orderConsumer}{ErdosProblems/Erdos249/PrimeRayCyclotomicCurvature.lean:242}
\quad \ev{Lean} \quad \scale{uniform} \quad \coord{other:cyclotomic-anchor}.
For $C(n)=2^n-1$, the first theorem above supplies the hypothesis with
$m=d=1$: here $\operatorname{ord}_q(p)=1$, whereas
$\operatorname{ord}_p(2)=q$. For $C(n)=|\Phi_n(2)|$ and a general fixed $m>0$, the source
uses the \emph{eventual} version: the divisibility is required only for
prime $q$ above a fixed threshold. The same proof then applies after that
threshold. It is not valid to replace this by an all-prime assertion:
$\Phi_6(2)=3$, but $6\nmid3-1$. For the binary cyclotomic family, choosing
$q>\max(m,2^m)$ removes these exceptional indices. Other families require
their own proof of the divisibility and nontriviality assumptions.
\end{thm}
\leannote{Lean: \lproof{ErdosProblems/Erdos249/PaperCompleteR21/MersennePrimeSupportAnchors.lean}{94}{boundedDegreeOrderConsumer_unfolded}, \lproof{ErdosProblems/Erdos249/PaperCompleteR21/MersennePrimeSupportAnchors.lean}{101}{coprime_of_dvd_pow_sub_one}, \lproof{ErdosProblems/Erdos249/PaperCompleteR21/MersennePrimeSupportAnchors.lean}{113}{dvd_pow_sub_one_iff_orderOf_dvd}, \lproof{ErdosProblems/Erdos249/PaperCompleteR21/MersennePrimeSupportAnchors.lean}{136}{boundedOrder_witness_iff_orderOf_le}, and 8 further declarations in the \hyperref[sec:coverage]{coverage section}.}

\subsubsection{The lower bound for rank-one quotients}
\label{rank-one-subrank}

The following finite-sum quotients do not approach $\Theta_2=S-1/2$.
Their lower bound is uniform in both the exponent and the truncation
length. This matters for rational approximation: varying either parameter
cannot make the associated linear form tend to zero.

\begin{thm}[A uniform positive gap for rank-one quotients]
For integers $r\ge2$ and $Y\ge1$, write
\[
 \Theta_r=\sum_{d\ge1}\frac{\mu(d)}{(2^d-1)^r},
 \qquad t(Y,r)=\sum_{d=1}^{Y}\frac{\mu(d)}{(2^d-1)^r}.
\]
For every $e\ge1$ and $Y\ge4$,
\[
 \frac{t(Y,e+2)^2}{t(Y,2e+2)}-\Theta_2>\frac1{480}.
\]
Here $\Theta_2=S-1/2$. For $r\ge3$, the proof uses
$1429/1512\le\Theta_r<1$ and
$|t(Y,r)-\Theta_r|\le1/3584$ when $Y\ge4$.
In particular, these inequalities make the denominator positive.
\lean{ErdosProblems.Erdos249.RankOneSubrankObstruction.rankOneSubrankQuotient_sub_theta_two_gt}{ErdosProblems/Erdos249/RankOneSubrankObstruction.lean:238}
\quad \lean{ErdosProblems.Erdos249.RankOneSubrankObstruction.mobiusMersenneTheta_ge_alpha}{ErdosProblems/Erdos249/RankOneSubrankObstruction.lean:163}
\quad \ev{Lean} \quad \scale{uniform} \quad \coord{other:rank-one-subrank}.
\end{thm}
\leannote{Lean: \lproof{ErdosProblems/Erdos249/RankOneSubrankObstruction.lean}{238}{rankOneSubrankQuotient_sub_theta_two_gt}, \lproof{ErdosProblems/Erdos249/RankOneSubrankObstruction.lean}{163}{mobiusMersenneTheta_ge_alpha}, \lproof{ErdosProblems/Erdos249/RankOneSubrankObstruction.lean}{196}{mobiusMersenneTheta_lt_one}, \lproof{ErdosProblems/Erdos249/RankOneSubrankObstruction.lean}{73}{abs_mobiusMersenneTheta_sub_prefix_le}, and 2 further declarations in the \hyperref[sec:coverage]{coverage section}.}

\begin{prop}[The bound is preserved by positive normalised averaging]
For a nonempty finite family of admissible quotients with positive
weights summing to $1$, the weighted average still exceeds $\Theta_2$
by more than $1/480$. If an admissible quotient is $p/q$ in lowest terms,
with $q>0$, then $|q\Theta_2-p|>q/480$.
\lean{ErdosProblems.Erdos249.RankOneSubrankObstruction.positive_direct_sum_sub_theta_two_gt}{ErdosProblems/Erdos249/RankOneSubrankObstruction.lean:300}
\quad \lean{ErdosProblems.Erdos249.RankOneSubrankObstruction.primitive_form_abs_gt}{ErdosProblems/Erdos249/RankOneSubrankObstruction.lean:341}
\quad \ev{Lean} \quad \scale{uniform} \quad \coord{other:rank-one-subrank}.
\end{prop}

\begin{obs}
Writing an admissible quotient in lowest terms as $p/q$, with $q>0$,
gives $|q\Theta_2-p|>q/480\ge1/480$. Such forms cannot tend to zero.
The conclusion concerns exactly the displayed quotients and their positive
normalised averages. It makes no assertion about other rational
approximants or about signed averages.
\end{obs}

\subsubsection{Multiples of a period}

Allowing positive multiples of a prescribed period gives another
criterion equivalent to irrationality of $S$. Sufficiency follows by
telescoping the tail-period identity, and necessity follows from
completeness of the residue test. The equivalence is proved, but its
quantified arithmetic condition is not established. The finite examples
below give denominator exclusions whose overlap with the diagonal
certificates is described explicitly.

\begin{thm}[Concatenation and specified period multiples]
For nonnegative integers $a,b,N$, define the integer block sum
\[
 Q_{a,N}=\sum_{j=1}^{a}\varphi(N+j)2^{a-j}.
\]
Splitting a block after its first $a$ terms gives
\[
 Q_{a+b,N}=2^bQ_{a,N}+Q_{b,N+a}.
\]
In particular, $Q_{2h,N}=2^hQ_{h,N}+Q_{h,N+h}$.
The corresponding cyclotomic identity is
$\Phi_4(2^h)=2^{2h}+1=\Phi_2(2^{2h})$: the order-four factor at height
$h$ is the order-two factor at height $2h$.

The order-three factor is different:
$\Phi_3(2^h)=2^{2h}+2^h+1$ comes from tripling the period.
It is not generally part of the doubling sequence. For example,
$\Phi_3(2)=7$ divides none of $2^{2^j}-1$, because the order of $2$
modulo $7$ is $3$, which does not divide $2^j$. Thus the block identities
relate specified period multiples; they do not put all the order-two,
order-three and order-four factors into one doubling chain.
\lean{ErdosProblems.Erdos249.PeriodMultipleEscape.totientBlock_add}{ErdosProblems/Erdos249/PeriodMultipleEscape.lean:77}
\quad \lean{ErdosProblems.Erdos249.PeriodMultipleEscape.totientBlock_two_mul}{ErdosProblems/Erdos249/PeriodMultipleEscape.lean:373}
\quad \ev{Lean} \quad \scale{uniform} \quad \coord{other:lcm-period-multiple}.
\end{thm}
\leannote{Lean: \lproof{ErdosProblems/Erdos249/PaperCompleteR21/TotientBlockConcatenation.lean}{23}{totientBlock_eq_paper_indexed_sum}, \lproof{ErdosProblems/Erdos249/PaperCompleteR21/TotientBlockConcatenation.lean}{41}{totientBlock_concatenation}, \lproof{ErdosProblems/Erdos249/PaperCompleteR21/TotientBlockConcatenation.lean}{46}{totientBlock_doubling}, \lproof{ErdosProblems/Erdos249/PaperCompleteR21/TotientBlockConcatenation.lean}{53}{cyclotomic_four_eval}, and 5 further declarations in the \hyperref[sec:coverage]{coverage section}.}

\begin{obs}[An equivalent quantified condition]
\lean{ErdosProblems.Erdos249.PeriodMultipleEscape.periodMultipleKillSupply_iff_irrational}{ErdosProblems/Erdos249/PeriodMultipleEscape.lean:432}
\quad \ev{Lean} (both implications proved; the quantified condition itself
is not established)
\quad \scale{cofinal} \quad \coord{other:lcm-period-multiple}. The paper's
current aggregate diagonal table gives certificates at every $H_t$ for $t\le82$;
this module's equivalence contributes no new information about $t=83$ or any
condition at arbitrarily large scales.
\lean{ErdosProblems.Skip.LadderT67.exists_diagonalKill_le_82}{ErdosProblems/Skip/LadderT67.lean:71264}
\end{obs}

\begin{obs}[Eight finite period certificates]
The cited declarations give certificates at basepoint $N=300$ for
\[
 h\in\{67,81,97,101,121,125,127,128\}.
\]
Each excludes every reduced fraction $a/b$, with $b>0$ and
$b\mid2^{300}(2^h-1)$. These exclusions overlap with those already
obtained from the diagonal. In particular, $67$ and $81$ divide $H_{82}$,
and $H_{82}>300$, so
\[
 2^{300}(2^h-1)\mid2^{H_{82}}(2^{H_{82}}-1)
 \qquad(h=67,81).
\]
Their entire denominator classes are already covered by
Proposition~\ref{prop:B11-inv}. The eight cases therefore cannot all
be described as new denominator classes relative to that proposition.
The exact comparison of the odd factors is
\[
 \gcd(2^h-1,2^{H_{82}}-1)=2^{\gcd(h,H_{82})}-1,
\]
by the Euclidean algorithm on the exponents. None of the remaining
six periods divides $H_{82}$, so its full odd factor is not covered
by that diagonal denominator. This compares with the $t=82$ certificate,
not with every other exclusion in the record; the case $h=127$ includes
the Mersenne prime $2^{127}-1$.

The case $h=67$ also succeeds at depth $L=67$. It illustrates the
restriction $L=h$ in the full-depth condition, while retaining a small
basepoint. That universally quantified condition is equivalent to
irrationality; one finite instance is not.
\lean{ErdosProblems.Erdos249.PeriodMultipleEscape.certifiedKill_67_300}{ErdosProblems/Erdos249/PeriodMultipleEscape.lean:471}
\quad \lean{ErdosProblems.Erdos249.PeriodMultipleEscape.certifiedKill_127_300}{ErdosProblems/Erdos249/PeriodMultipleEscape.lean:490}
\quad \lean{ErdosProblems.Erdos249.PeriodMultipleEscape.certifiedKill_67_300_fullDepth}{ErdosProblems/Erdos249/PeriodMultipleEscape.lean:498}
\quad \lean{ErdosProblems.Erdos249.PeriodMultipleEscape.totient_series_ne_rat_of_den_dvd_127_300}{ErdosProblems/Erdos249/PeriodMultipleEscape.lean:534}
\quad \ev{Cert} \quad \scale{fixed} \quad \coord{other:lcm-period-multiple}; eight finite exclusions, not a proof of the condition at arbitrarily large scales.
\end{obs}

\subsubsection{The strict prime-position criterion}

A sufficient prime-index condition is that, for every integer $h\ge1$
and every $N_0\ge0$, there is a prime
$p\ge\max(N_0+h+1,h+5)$ such that
\[
 \operatorname{Re}e\bigl(2^{p-h-1}(2^h-1)S\bigr)<\frac9{10}.
\]
Here $e(x)=\exp(2\pi i x)$. The positive difference between $9/10$ and
this real part may depend on the chosen prime; the truncation depth is
then chosen so that its error fits within that difference. This is a
weaker phase requirement than the earlier bound $4/5$, not a larger
saving from $1$. The implication to irrationality is proved, but the
displayed prime-index condition remains unproved.
\lean{ErdosProblems.Erdos249.irrational_totient_series_of_naturalPrimeTailOrbitStrictGap}{ErdosProblems/Erdos249/TotientStrictPrimeEscape.lean:524}
\quad \ev{Lean} (implication proved; antecedent
\ev{Open})
\quad \scale{cofinal} \quad \coord{other:first-harmonic}.

\subsubsection{Finite Euler-product identities}

These are identities for the Dirichlet-convolution square $\mu*\mu$,
not the pointwise square $\mu(n)^2$. For each prime $p$, its coefficients
at $1,p,p^2$ are $1,-2,1$, and the higher prime-power coefficients vanish.
The local factors at exponents $1$ and $2$ are
\[
 1-\frac2p+\frac1{p^2}=\left(1-\frac1p\right)^2,
 \qquad
 1-\frac2{p^2}+\frac1{p^4}=\left(1-\frac1{p^2}\right)^2.
\]
For a prime $p$ and $e\ge0$, the corresponding convolution with the
divisor-sum function $\sigma$ gives
\[
 \sigma(p^{e+2})-2\sigma(p^{e+1})+\sigma(p^e)
       =p^{e+1}(p-1)=\varphi(p^{e+2}).
\]
This follows by subtracting the finite geometric sums. These finite
identities do not assert a transcendence result.
\lean{ErdosProblems.Erdos249.FiniteEulerSieve.muSqEulerFactor_one}{ErdosProblems/Erdos249/FiniteEulerSieve.lean:28}
\quad \lean{ErdosProblems.Erdos249.FiniteEulerSieve.sigmaPrimePow_secondDiff}{ErdosProblems/Erdos249/FiniteEulerSieve.lean:51}
\quad \ev{Lean} \quad \scale{uniform} \quad \coord{other:euler-sieve}.
% ; - part a249_p2 ; -
\section{Parameter ranges and quantifiers}\label{sec:249-scale-ladder}

The table records parameter ranges and quantifiers. ``Fixed'' refers
to listed examples, ``bounded'' to a finite range, and ``cofinal'' to a
statement of the form $\forall N_0\,\exists N\ge N_0$. A bound of the form
$\forall a\ge8$ is not finite-range evidence: it gives a conclusion at all
sufficiently large parameters. ``Uniform'' describes the quantification of
a statement, not whether all hypotheses needed to apply it are proved.
A bound on a secondary variable, such as $j<2\cdot2^a$, does not
restrict the scale $a$ to a finite range. Such parameter-uniform results
are grouped with the uniform statements, not with finite calculations.


Some identities have no distinguished scale parameter. They are listed
with the general identities. An equivalence can hold for all parameters
without establishing either of the equivalent arithmetic assertions.


The LCM formulation of Problem~249 requires arbitrarily large scales:
\lean{irrational\_totientSeries\_iff\_actualLcmOrbitNonintegralitySupply}{TotientActualLcmOrbitNonintegrality.lean:37}
\[
  \mathrm{Irrational}\Bigl(\sum_{n\ge 1}\varphi(n)/2^n\Bigr) \iff \forall a_0\ \exists a\ge a_0,\ \Omega_a \notin \mathbb Z,
\]
a cofinal statement in the exponent $a$ indexing $H=H(2^a)$.
A finite calculation alone does not establish it; a further theorem
could propagate finite information to arbitrarily large scales.

The table groups statements by their parameter ranges, not by proof
status. In particular, some assertions with cofinal quantifiers are
already proved, such as the prime constructions and rational comparison
sequences. Others are implications whose quantified hypotheses remain
open. The horizontal rule separates the cofinal group for ease of
reference; it does not mark a boundary between proved and unproved
mathematics. The descriptions identify this distinction in that group.
The prefixes summarise the quantifiers; complete hypotheses are given in
the cited statements. A definition records an object or condition; it does not by itself prove
the associated soundness or existence assertion.

\begingroup
In the periodic-weight rows, write
$w:\mathbb N\to\mathbb Z$ and
$c_w(n)=\sum_{d\mid n}w(d)$ for $n\ge1$. Positivity and nonvanishing
refer to these divisor coefficients, not to the weight $w$ itself.
Here ``nonzero arbitrarily far out'' means that for every $N$ there is
$n>N$ with $c_w(n)\ne0$; it does not mean nonvanishing on every
arithmetic progression. The nonprincipal character $\chi_4$ modulo four
is $0$ on even integers and $(-1)^{(n-1)/2}$ on odd integers.

\begin{landscape}
\fontsize{8}{9.8}\selectfont
\setlength{\tabcolsep}{3pt}
\begin{longtable}{@{}>{\raggedright\arraybackslash}p{0.22\linewidth}>{\raggedright\arraybackslash}p{0.18\linewidth}>{\raggedright\arraybackslash}p{0.08\linewidth}>{\raggedright\arraybackslash}p{0.34\linewidth}>{\raggedright\arraybackslash}p{0.12\linewidth}@{}}
\caption{Results and unproved conditions, grouped by their quantifiers.
The dividing rule separates the cofinal group, not proof status.
Complete hypotheses are given in the cited statements.}\\
\toprule
\papertablehead
\textbf{Result} & \textbf{Source} & \textbf{Scale} & \textbf{Quantifier prefix} & \textbf{Representation}\\
\midrule
\endfirsthead
\toprule
\papertablehead
\textbf{Result} & \textbf{Source} & \textbf{Scale} & \textbf{Quantifier prefix} & \textbf{Representation}\\
\midrule
\endhead
\midrule \multicolumn{5}{r}{\emph{continued on next page}} \\
\endfoot
\bottomrule
\endlastfoot

\papertablehead\multicolumn{5}{l}{\emph{Uniform statements, including unconditional identities and implications}}\\
scaled totient tail (definition) & \lean{totientTail}{TotientTailPeriodKiller.lean:57} & uniform & $\forall N$ & binary-weighted series\\
integer prefix and real tail & \lean{two\_pow\_mul\_totient\_series\_eq}{TotientTailPeriodKiller.lean:150} & uniform & $\forall N$ & binary-weighted series\\
$D(h,N,L)$ (definition) & \lean{windowDiscrepancy}{TotientTailPeriodKiller.lean:66} & uniform & $\forall h\,\forall N\,\forall L$ & finite binary sums\\
$\mathcal C(h,N,L)$ (definition) & \lean{certifiedKill}{TotientTailPeriodKiller.lean:72} & uniform & $\forall h\,\forall N\,\forall L$ & finite binary sums\\
necessary depth inequality & \lean{certifiedKill\_depth\_floor}{TotientTailPeriodKiller.lean:79} & uniform & $\forall h\,\forall N\,\forall L$, given \ensuremath{\mathcal{C}}$\,h\,N\,L$ & finite binary sums\\
residue test soundness & \lean{tail\_diff\_notMem\_int\_of\_certifiedKill}{TotientTailPeriodKiller.lean:262} & uniform & $\forall h\,\forall N\,\forall L$, given \ensuremath{\mathcal{C}} & finite binary sums\\
completeness of the residue test & \lean{exists\_certifiedKill\_iff\_tail\_diff\_notMem\_int}{LcmConeFlatness.lean:316} & uniform & $\forall h\,\forall N$ & finite binary sums\\
integrality from denominator divisibility & \lean{tail\_diff\_int\_of\_den\_dvd}{TotientTailPeriodKiller.lean:327} & uniform & $\forall N\,\forall h\,\forall r{:}\mathbb{Q}$, given $S=r$, $r.\mathrm{den}\mid 2^N(2^h{-}1)$ & binary-weighted series\\
tail-period law & \lean{eventual\_period\_of\_not\_irrational}{TotientTailPeriodKiller.lean:358} & uniform & $\lnot\mathrm{Irrational}\,S\to\exists h{>}0\,\exists N_0\,\forall N{\ge}N_0$ & binary-weighted series\\
least common multiple $H(t)$ & \lean{periodLcm}{CarrySurvivorExtinction.lean:515} & uniform & $\forall t$ & LCM periods\\
a nondivisor below $2t$ is a prime power & \lean{eq\_prime\_pow\_of\_not\_dvd\_periodLcm}{LcmDiagonalReduction.lean:137} & uniform & $\forall t\,\forall j\,(0{<}j{<}2t)$, given $j\nmid{H}\,t$ & LCM offsets\\
totient factorisation at divisor offsets & \lean{totient\_periodLcm\_ray\_split}{LcmDiagonalReduction.lean:208} & uniform & $j\mid H(t)$, with the prime-divisor hypotheses of Proposition~\ref{catalogue:cert:b2} & totient factorisation\\
integrality on the LCM grid under rationality & \lean{rational\_totient\_series\_forces\_lcm\_cone\_flatness}{CertificateKernel.lean:19002} & uniform & $\lnot\mathrm{Irrational}\,S\to\exists t_1\,\forall t{\ge}t_1\,\forall q{>}0\,\forall m$ & LCM grids\\
second-difference residue test & \lean{second\_diff\_notMem\_int\_of\_certifiedRank2Kill}{LcmConeFlatness.lean:493} & uniform & $\forall h\,\forall N\,\forall L$, given \ensuremath{\text{second-difference residue condition}} & finite binary sums\\
a nonintegral pair on a finite grid & \lean{exists\_nonintegral\_pair\_of\_coneNonflatCert}{LcmConeNonflat.lean:126} & uniform & finite $Q\ne\emptyset$; all size and residue hypotheses of Theorem~\ref{catalogue:cert:b10a} & finite LCM grids\\
finite exclusion of integer carries & \lean{tail\_diff\_notMem\_int\_of\_survivorKill}{CarrySurvivorExtinction.lean:428} & uniform & $\forall h\,\forall N\,\forall K$, given \ensuremath{\text{finite carry exclusion}} & integer recurrences\\
Farey neighbour denominator law & \lean{farey\_gap}{GapFareyBound.lean:51} & uniform & $a,b,c,d,r,s\in\mathbb Z$; $b,d>0$, $bc-ad=1$, $as<br$, $dr<cs$ & rational intervals\\
rational-approximation criterion & \lean{irrational\_of\_den\_mul\_abs\_sub\_tendsto\_zero}{CertificateKernel.lean:5371} & uniform & $u_j\in\mathbb Q$, $x\in\mathbb R$; eventually $u_j\ne x$, and $\operatorname{den}(u_j)|u_j-x|\to0$ & general criterion\\
near-integer irrationality criterion & \lean{irrational\_of\_int\_mul\_near\_int}{CertificateKernel.lean:6120} & uniform & $\forall \xi{:}\mathbb{R}$, $\forall q{>}0\,\exists m\,z\,0{<}|m\xi{-}z|{<}1/q$ & general criterion\\
independence from a nonsingular evaluation matrix & \lean{linearIndependent\_of\_separatedMinorCertificate}{TotientMahlerDefect.lean:91} & uniform & a finite family with a square evaluation matrix of nonzero determinant & general criterion\\
dyadic totient-kernel rank $=2^e{+}1$ & \lean{linearIndependent\_canonicalTotientKernelFamily}{TotientMahlerDefect.lean:935} & uniform & $\forall e\ge1$ & dyadic sections\\
rationality forces unbounded carry-kernel rank & \lean{not\_irrational\_totientSeries\_implies\_unbounded\_carryRank\_unconditional}{TotientCarryKernelRigidity.lean:300} & uniform & $\lnot\mathrm{Irrational}\,S\to\exists v{>}0\,\exists u\,\forall e$ & carry sections\\
no fixed $D$ clears all $A(n)/n$ for $A=\varphi*\mu$ & \lean{no\_fixed\_integral\_primitive\_euler\_index}{PrimitiveDeterminantLift.lean:169} & uniform & $\forall D{>}0$ & M\"obius sums\\
Lambert-series value identities & \lean{tsum\_totient\_div\_pow\_two\_eq\_pnat\_half\_pow}{CertificateKernel.lean:18415} & uniform & value identities, $L(\mu){=}\tfrac12$, $L(\varphi){=}2$, $L(1){=}$Erd\H{o}s--Borwein & M\"obius sums\\
prime-power unit-gap criterion & \lean{reducedDenominatorUnitGapCert\_prime\_pow\_iff}{PrimitiveWeightCertificate.lean:54} & uniform & $\forall p\ \text{prime},\ e>0,\ N,K\ge0$ & rational intervals\\
positive rational-difference lower bound & \lean{positive\_rational\_difference\_lower\_bound}{PrimitiveRationalGapSupply.lean:31} & uniform & $\forall\ \text{whole}\,\text{pfx}{:}\mathbb{Q},\ \text{pfx}{<}\text{whole}$ & general criterion\\
$L_2(\mu)=S-\tfrac12$ & \lean{tsum\_totient\_half\_pow\_eq\_half\_add\_moebius\_sq}{MersenneLambertLadder.lean:665} & uniform & unconditional & M\"obius sums\\
weighted squared-denominator identity & \lean{tsum\_lambert\_linear\_weight\_sq\_pure}{GcdMomentCalculus.lean:105} & uniform & $\forall w{:}\mathbb{N}\to\mathbb{R}\,(|w(d)|{\le}d)\,\forall r\in[0,1)$ & M\"obius sums\\
$L_2(\mu)=S-\tfrac12$ & Proposition~\ref{prop:mobsq} & uniform & unconditional & M\"obius sums\\
$L_2(1)=\zeta_q(2)-\zeta_q(1)$ identity & \lean{tsum\_one\_div\_mersenne\_sq\_eq\_sigma\_sub\_tau\_series}{GcdMomentCalculus.lean:216} & uniform & unconditional & M\"obius sums\\
$L_2(\varphi)$ = gcd moment & \lean{tsum\_totient\_div\_mersenne\_sq\_eq\_gcd\_moment\_series}{GcdMomentCalculus.lean:235} & uniform & unconditional & gcd moments\\
comparison of the weights & Propositions~\ref{prop:mobsq}--\ref{prop:pillai} & uniform & unconditional & M\"obius sums\\
gcd-divisibility factorizes & \lean{tsum\_pos\_pair\_both\_dvd\_half\_eq\_inv\_mersenne\_sq}{GcdMomentCalculus.lean:266} & uniform & $\forall d{:}\mathbb{N},\,d{>}0$ & geometric weights\\
sum over positive coprime pairs equals $1$ & \lean{tsum\_pos\_coprime\_inv\_mersenne\_eq\_one}{GcdMomentCalculus.lean:349} & uniform & unconditional (sum over all coprime pairs) & geometric weights\\
Stern--Brocot recursion with stopping & \lean{cylinderMass\_split}{GcdMomentCalculus.lean:474} & uniform & positive integers $a,b$ & Stern--Brocot tree\\
positive numerator coefficients & \lean{mobiusNumeratorPolynomial\_eq\_gcdWord}{RepunitMobiusNumerator.lean:205} & uniform & $\forall r\,(\mathrm{Squarefree}\,r)\,\forall k{<}r$ & divisor and cyclotomic sums\\
evaluation at $2$ gives the numerator & \lean{mobiusNumeratorPolynomial\_eval\_two}{RepunitMobiusNumerator.lean:445} & uniform & $\forall r\,(\mathrm{Squarefree}\,r)$ & divisor and cyclotomic sums\\
reduction to the radical of $H$ & \lean{scaledMobiusShadow\_eq\_radicalBase}{RadicalMobiusShadow.lean:125} & uniform & $\forall H{>}0\,\forall r{>}0$ & M\"obius sums\\
numerator congruence modulo $\Phi_m$ & \lean{cyclotomic\_dvd\_mobiusNumeratorPolynomial\_sub}{CyclotomicProjectionOfShadow.lean:299} & uniform & $\forall r\,(\mathrm{Squarefree}\,r)\,\forall m\,(m{\mid}r)$ & divisor and cyclotomic sums\\
gcd of the numerator and $\Phi_r(2)$ & \lean{mobiusNumerator\_gcd\_cyclotomicValue}{CyclotomicProjectionOfShadow.lean:350} & uniform & $\forall r\,(\mathrm{Squarefree}\,r)$ & divisor and cyclotomic sums\\
Mersenne factors retained in the denominator & \lean{upperHalfChannel\_product\_dvd\_den\_of\_scale\_primeFactors\_le}{MersenneShadowCyclotomicNoncollapse.lean:809} & uniform & $t\ge5$, squarefree $r$; finite $P$ of primes dividing $r$ in $(t/2,t]$; prime factors of $r$ and the multiplier are at most $t$ & divisor and cyclotomic sums\\
denominator lower bound / exact value & \lean{lcmHeight\_scaledMobiusShadow\_den\_exact}{MersenneShadowDenominatorGrowth.lean:147} & uniform & $\forall t$ (exact-value); $t{\ge}5$ (lower bound) & M\"obius sums\\
numerator congruence after adjoining a prime & \lean{cyclotomic\_dvd\_primeJump\_new\_fibre}{PrimePowerJumpDynamics.lean:281} & uniform & $r\ge1$, $p$ prime, $p\nmid r$, $m\mid r$ & divisor and cyclotomic sums\\
squared-Mersenne series: partial sum plus tail & \lean{mobius\_square\_series\_eq\_partial\_add\_tail}{SquaredMersenneDiagonalEnclosure.lean:94} & uniform & $\forall D{:}\mathbb{N}$ & M\"obius sums\\
cancellation of the affine divisor terms & \lean{joint35\_oldChannel\_zero}{JointExponentTransport.lean:71} & uniform & $\forall d{>}0\,(d{\mid}H)$; general: finite affine annihilator & LCM diagonal\\
new-divisor counting identity & \lean{supportCoeff\_mul\_eq\_add\_defect}{CompositeDilationDefect.lean:30} & uniform & $\forall a{\in}A\,(a{>}0)\,\forall x{>}0$ & divisor counts\\
bound on zero blocks in divisor counts & \lean{supportCoeffZeroWindow\_length\_le\_eps\_logb\_add}{SublogDivisorCoverage.lean:392} & uniform & $\forall \varepsilon{>}0\,\exists B$, given rationality of the base-2 support series & divisor counts\\
equivalent shift condition for the target $1/2$ (Problem~257) & \lean{shiftWindowZero\_iff\_half\_mem\_mersenneAchievementSet}{CampbellShiftSynchronization.lean:488} & uniform & iff, no extra parameter & Mersenne subsums\\
totient multiplication with common prime factors & \lean{totient\_mul\_eq\_overlapFactor\_mul}{TotientActualLcmOrbitArithmetic.lean:43} & uniform & $\forall x{>}0$ ($j$ fixed) & totient identities\\
totient relative-Euler-product split & \lean{totient\_mul\_eq\_totient\_mul\_relativeEulerProduct}{TotientActualLcmOrbitArithmetic.lean:100} & uniform & $\forall j\,x{>}0$ & totient identities\\
formula for $\varphi(2H(t)+j)-\varphi(H(t)+j)$ & \lean{lcmRayArithmeticLetter\_eq\_deltaTotient}{TotientActualLcmOrbitArithmetic.lean:298} & uniform & $\forall t\,\forall j$ & M\"obius sums\\
the finite diagonal sum equals $D(H(t),H(t),L)$ & \lean{lcmDiagonalArithmeticWord\_eq\_windowDiscrepancy}{TotientActualLcmOrbitArithmetic.lean:1998} & uniform & $\forall t\,\forall L$ & M\"obius sums\\
\#249 $\iff$ cofinal actual-orbit nonintegrality & \lean{irrational\_totientSeries\_iff\_actualLcmOrbitNonintegralitySupply}{TotientActualLcmOrbitNonintegrality.lean:37} & uniform & iff & M\"obius sums\\
the LCM tail difference is an integer translate of a multiple of $S$ & \lean{actualLcmTailOrbit\_eq\_scaled\_totientSeries\_sub\_prefix}{TotientActualLcmOrbitSeparation.lean:68} & uniform & $\forall a$ & M\"obius sums\\
explicit finite-block approximation & \lean{abs\_actualLcmTailOrbit\_sub\_rawApprox\_lt}{TotientActualLcmOrbitSeparation.lean:141} & uniform & $\forall a\,\forall q$ & M\"obius sums\\
an equivalent residue-interval condition & \lean{exists\_certifiedKill\_iff\_guardCylinderWitness}{TotientActualLcmTopEdgeStaircase.lean:844} & uniform & $\forall h\,N,\ \text{iff}$ & integer recurrences\\
sign of the three-term totient difference under the stated middle-value bound & \lean{fixedRankSecondDifference}{TotientFixedRankLcmAsymptotic.lean:256} & uniform & $\forall H\,j$, given MiddleRankTotientExtremal & totient differences\\
exact value on a specified dyadic family & \lean{fixedRankSecondDifference\_dyadic\_fixture\_exact}{TotientFixedRankLcmAsymptotic.lean:462} & uniform & $\forall n$ (closed-form) & totient differences\\
the rational parity comparison (definition) & \lean{parityCoboundaryWeight}{TotientParityCoboundaryCountermodel.lean:59} & uniform & $\forall n$ (definition) & binary-weighted series\\
$R_N$ as a sum over positive divisor indices & \lean{totientTail\_eq\_tsum\_positiveForeignResidueKernel}{TotientShiftedMobiusForeignBridge.lean:61} & uniform & $\forall N$ & divisor and cyclotomic sums\\
$R_N$ as a shifted M\"obius sum & \lean{totientTail\_eq\_tsum\_shiftedMobiusPulse}{TotientShiftedMobiusPulse.lean:547} & uniform & $\forall N$ & divisor and cyclotomic sums\\
absolute convergence and regrouping of Lambert sums & \lean{tsum\_lambert\_pair\_regroup\_if}{TotientShiftedMobiusPulse.lean:124} & uniform & $\forall w\,v{:}\mathbb{N}\to\mathbb{R}\,(|w(d)|{\le}d,\,|v(m)|{\le}m)$ & divisor and cyclotomic sums\\
carry divisibility and integral tail differences & \lean{carryShift\_dvd\_iff\_tailDiff\_mem\_int}{TotientTailCarryPeriod.lean:77} & uniform & $\forall N\,k$, given tempered orbit & integer recurrences\\
rationality $\Rightarrow$ periodicity \& unbounded rank & \lean{not\_irrational\_totientSeries\_implies\_mod\_period\_and\_unbounded\_rank}{TotientTailCarryPeriod.lean:224} & uniform & $\lnot\mathrm{Irrational}\,S\to\exists v\,u\,\forall e$ & integer recurrences\\
the stated reconstruction cost is at least $3$ & \lean{three\_le\_totientAdjugateTailCost}{TotientTailCarryPeriod.lean:266} & uniform & $\forall\ \text{finite}\ w{:}\iota\to\mathbb{Q}\,x{:}\iota\to\mathbb{N}$, given $\sum w_i\varphi(x_i){=}1$ & weighted tail bounds\\
directed residue interval (definition) & \lean{directedCertifiedKill}{TotientTailCarryPeriod.lean:766} & uniform & $h,N,L\ge0$; soundness and the existential-depth equivalence are proved in Proposition~\ref{prop:CP-06} & integer recurrences\\
rationality and an integer carry with vanishing boundary term & \lean{binaryCoeffSeries\_rational\_iff\_exists\_temperedBinaryOrbit}{GenericTailOrbitRigidity.lean:426} & uniform & $\forall c{:}\mathbb{N}\to\mathbb{N}\,(c(n){\le}n)$, iff & binary-weighted series\\
doubling-orbit rigidity & \lean{doublingOrbit\_eq\_zero\_of\_tempered}{GenericTailOrbitRigidity.lean:315} & uniform & $\forall d{:}\mathbb{N}\to\mathbb{R}$, doubling $+$ tempered $\to d{\equiv}0$ & integer recurrences\\
number of labels needed for exact tail recovery & \lean{balancedPulse\_no\_autonomous\_decoder}{GenericTailOrbitRigidity.lean:247} & uniform & $\forall m\,\forall\mathrm{State}$, collapsing-state $\to$ no decoder & binary-weighted series\\
endpoint residues do not determine the initial carry & \lean{affineBinaryOrbit\_mod\_twoPow\_eq}{GenericTailOrbitRigidity.lean:307} & uniform & $\forall a{:}\mathbb{N}\to\mathbb{Z}\,\forall u_0\,v_0\,\forall L$ & binary-weighted series\\
integer-coefficient block-certificate dichotomy & \lean{int\_coeff\_series\_irrational\_or\_bpow\_mul\_eq\_intCast\_of\_full\_block\_certificates}{CertificateKernel.lean:13934} & uniform & $\forall b{\ge}2\,\forall c\,\forall q$, given full-block certificate at $q$ & binary-weighted series\\
divisibility under iterated forward differences & \lean{dvd\_iteratedForwardDifference}{LcmFactorIdealPulseObstruction.lean:211} & uniform & $d\mid f(n)$ for every $n$; every finite list of shifts & integer sequences\\
totient series as a coprime-pair sum & \lean{tsum\_coprime\_pair\_pow\_eq\_tsum\_totient\_mul\_pow}{GeometricCoprimality.lean:120} & uniform & $\forall n$; $\forall r\in[0,1)$ & geometric weights\\
factorisation through a rational-valued linear functional & \lean{linearDescender\_eq\_smul\_eval}{AdelicHeightObstruction.lean:120} & uniform & rational vector spaces $V,W$; surjective linear $\mathrm{ev}:V\to\Q$, linear $L:V\to W$, $\ker(\mathrm{ev})\subseteq\ker L$ & divisibility and linear algebra\\
divisibility of a reduced denominator & \lean{divisor\_dvd\_divInt\_den}{RationalDenominatorSurvival.lean:17} & uniform & $\forall D{>}0\,\forall m{\mid}D\,\forall a$ & divisibility and linear algebra\\
the complementary denominator divides the integer multiplier & \lean{scalarLocalization\_complement\_dvd}{AdelicHeightObstruction.lean:23} & uniform & $\forall x{:}\mathbb{Q}\,\forall H\,(H{\mid}x.\mathrm{den})\,\forall c\,((cx).\mathrm{den}{\mid}H)$ & divisibility and linear algebra\\
largest-exponent nonvanishing criterion & \lean{scaled\_dyadic\_sum\_ne\_zero}{SignedQMomentObstruction.lean:96} & uniform & $\exists m{\in}s$ (strictly maximal, odd coeff), $\forall i{\ne}m\,e(i){<}e(m)$ & divisibility and linear algebra\\
column rescaling and determinant nonvanishing & \lean{det\_residualMonomialMatrix}{ResidualGaugeObstruction.lean:47} & uniform & $\forall d\,\forall e{:}\mathrm{Fin}\,d{\to}\mathbb{N}\,\forall z{:}\mathrm{Fin}\,d{\to}\mathbb{C}$ ($z$ nonzero) & complex matrices\\
determinant $1$ for the divisor-incidence matrix & \lean{incidenceMobius\_det\_eq\_one}{IncidenceQuotientHermitePade.lean:56} & uniform & $\forall N{:}\mathbb{N}$ & M\"obius sums\\

error bound from convergence & \lean{controlledForeignProjection\_of\_tail\_limit}{LambertDiagonalEnclosure.lean:31} & uniform & $2H{\le}D$; convergence is supplied for $H>0$ in Section~\ref{sec:nondivisor-enclosure} & M\"obius sums\\
carries realising fixed-precision unit residues & \lean{fixedPrecisionTropicalNoGo}{TropicalCurvatureCarry.lean:137} & uniform & $\forall u{>}0$ (fixed precision), $\forall$ finite carry word & divisibility and linear algebra\\
vanishing of iterated prime dilation differences & \lean{iteratedDilationDifference\_eq\_zero\_of\_cardFactors\_le}{SupportDilationDifferences.lean:250} & uniform & distinct primes $p_1,\ldots,p_r$, $r>K$, $\gcd(n,\prod p_i)=1$; $\Omega(a)\le K$ on $A$ & divisor counts\\
$\varphi(n)>\tfrac34 n$ for the stated rough integers & \lean{three\_quarters\_mul\_lt\_totient\_of\_rough\_lt\_two\_pow}{TotientActualLcmOrbitArithmetic.lean:768} & uniform & $\forall a{\ge}8\,\forall n{>}0$ ($n$ is $t$-rough, $t{=}2^a$, $n{<}2^{2\cdot2^a}$) & M\"obius sums\\
positivity of short-window totient differences & \lean{lcmRayArithmeticLetter\_pos\_of\_lt\_two\_mul}{TotientActualLcmOrbitArithmetic.lean:1703} & uniform & $\forall a{\ge}8\,\forall j\,(0{<}j{<}2\cdot2^a)$ & M\"obius sums\\
the negative representative of the true endpoint & \lean{actualLcm\_trueEndpointSurvivor\_neg}{TotientActualLcmOrbitSign.lean:172} & uniform & $\forall a{\ge}8\,\forall J\,K$ (room bound) & M\"obius sums\\
integrality forces upper-endpoint residue & \lean{actualLcm\_integral\_forces\_topEdgeResidue}{TotientActualLcmOrbitSign.lean:211} & uniform & $\forall a{\ge}8\,\forall J\,K$ (room, $2H{+}J{+}K{+}2{<}2^K$) & M\"obius sums\\
terminal dyadic staircase impossible with room & \lean{not\_actualLcmTerminalDyadicStaircase\_of\_room}{TotientActualLcmTopEdgeStaircase.lean:251} & uniform & $\forall a{\ge}8\,\forall J\,K\,m$ (room, modulus-wide) & M\"obius sums\\
penultimate difference in a partial divisibility pattern & \lean{puncturedDyadicStaircase\_penultimate\_eq\_half}{TotientActualLcmTopEdgeStaircase.lean:1187} & uniform & $\forall a{\ge}8\,\forall J\,K\,m$ (hypotheses of Proposition~\ref{prop:TE-02-inv}) & M\"obius sums\\
upper-endpoint residue-gap $\Rightarrow$ nonintegral & \lean{actualLcmTailDiff\_notMem\_int\_of\_topEdgeResidueGap}{TotientActualLcmTopEdgeStaircase.lean:1260} & uniform & $\forall a{\ge}8\,\forall J\,K\,m$ (room, one-sided) & M\"obius sums\\
relation between the centred residue and the actual carry & \lean{two\_mul\_actualOddHalfCenteredLift\_eq\_terminal\_sub\_trueCarry}{TotientActualLcmTopEdgeStaircase.lean:1678} & uniform & $\forall a{\ge}8\,\forall q$ (room, fit bound $2(H{+}q{+}2){\le}4^q$) & M\"obius sums\\
divisibility of the three-term totient difference & \lean{two\_mul\_totient\_dvd\_fixedRankSecondDifference}{TotientFixedRankLcmAsymptotic.lean:519} & uniform & $\forall a{\ge}4\,\forall j{>}0\,(j^2{\le}2^a)$ & totient differences\\
equality of two finite rational approximations & \lean{renormalizedResidueAgreement\_of\_twice\_le}{TotientShiftedMobiusForeignBridge.lean:338} & uniform & $H>0$, $D\ge2H$; no separation hypothesis for this equality & M\"obius sums\\
a Mersenne divisor need not divide a tail difference & \lean{idCoeff\_mersenneModulus\_does\_not\_annihilate\_tailShift}{TotientTailCarryPeriod.lean:342} & uniform & $\forall K,q\,(K{>}0,K{<}q,q{\mid}2^K{-}1)$ & M\"obius sums\\
finite numerator separation implies nonintegrality & \lean{scaleFullTarget\_miss\_of\_lambert\_projected\_num\_gap}{LambertProjectedNumeratorGap.lean:23} & uniform & $\forall H,D$ fixed, given uniform numerator-gap bound & M\"obius sums\\
periodic-weight Lambert dichotomy & \lean{irrational\_or\_bpow\_mul\_eq\_intCast\_intWeightedErdosSeries\_periodic}{CertificateKernel.lean:14175} & uniform & $\forall b{\ge}2\,\forall m{>}0\,\forall w$ ($w(n{+}m){=}w(n)$) & M\"obius sums\\
periodic weight with nonnegative divisor coefficients & \lean{irrational\_intWeightedErdosSeries\_periodic\_of\_coeff\_nonneg\_of\_frequently\_ne\_zero}{CertificateKernel.lean:14583} & uniform & $b\ge2$, $m>0$, $w(n+m)=w(n)$; $c_w(n)\ge0$ and $c_w(n)\ne0$ arbitrarily far out & M\"obius sums\\

zeros on one progression for the character modulo four & \lean{intWeightedCoeff\_periodFourSignWeight\_eq\_zero\_of\_mod\_four\_eq\_three}{CertificateKernel.lean:14226} & uniform & $c_{\chi_4}(n)=0$ for every $n\equiv3\pmod4$; the weight is fixed & M\"obius sums\\

positivity of the translated LCM tail difference & \lean{actualLcmTailDiff\_shift\_pos}{TotientActualLcmOrbitSign.lean:39} & uniform$^\dagger$ & $\forall a{\ge}8\,\forall J\,(J{+}(a{+}6){<}2\cdot2^a)$, no rationality hypothesis & M\"obius sums\\

\papertablehead\multicolumn{5}{l}{\emph{Finite ranges and parameter-bounded checks}}\\
Farey window $K{=}240$ denominator bound & \lean{gap\_check\_window\_1\_240\_le\_79639646646701375323355774875831053}{GapFareyBound.lean:176} & bounded &
$\forall q{:}\mathbb{N},\,0{<}q{\le}$\newline
$79{,}639{,}646{,}646{,}701{,}$\newline
$375{,}323{,}355{,}774{,}875{,}831{,}053$ & rational intervals\\
denominator lower bound $S{\ne}p/q$, $q{\le}Q_0$ & \lean{tsum\_totient\_div\_pow\_two\_ne\_ratCast\_of\_den\_le\_79639646646701375323355774875831053}{CertificateKernel.lean:18384} & bounded & $\forall p{:}\mathbb{Q},\,p.\mathrm{den}{\le}Q_0$ & rational intervals\\
finite LCM certificates covering thresholds $a_0\le6$ & \lean{powerTwoActualLcmShortArithmeticKillSupply\_through\_six}{TotientActualLcmShortKill.lean:54} & bounded & $\forall a_0{\le}6\,\exists a\,L$ & M\"obius sums\\

contiguous diagonal certificate band through $t{\le}82$ & \lean{ErdosProblems.Skip.LadderT67.exists_diagonalKill_le_82}{ErdosProblems/Skip/LadderT67.lean:71264} & bounded & $\forall t{\le}82\,\exists L$; no $t{=}83$ certificate claimed & LCM diagonal\\

\papertablehead\multicolumn{5}{l}{\emph{Listed examples}}\\
residue certificates for $1\le h\le8$, $N=12$, $L=16$ & \lean{certifiedKill\_all\_small}{TotientTailPeriodKiller.lean:404} & fixed & $\forall h\in[1,8]$, \ensuremath{\mathcal{C}}\,$h$\,12\,16 (by \texttt{decide}) & finite binary sums\\
denominator exclusions for $1\le h\le16$, binary exponent $14$ & \lean{totient\_series\_ne\_rat\_of\_den\_dvd\_pow\_two\_mul\_mersenne\_upto\_sixteen}{CertificateKernel.lean:18890} & fixed & $\forall h\in[1,16]$, denominator exponent 14 (by \texttt{decide}) & finite binary sums\\
the second-difference test uses more depth at $(h,N)=(1,8)$ & \lean{totient\_tail\_rank\_two\_kill\_sound\_but\_not\_shallower\_cell}{CertificateKernel.lean:19090} & fixed & first-difference certificate at $L=8$; no second-difference certificate for $1\le L\le8$, and one at $L=9$ & finite binary sums\\
28 listed diagonal certificates with $t\le64$ & \lean{certifiedKill\_diagonal\_all\_imported}{DiagonalPincerCertificates.lean:2870} & fixed & $t\in\{1,2,3,4,5,7,8,9,11,13,16,17,\dots\}$ (28 explicit values) & LCM diagonal\\
LCM certificates at $a=4$ and $a=6$ & \lean{lcmDiagonalArithmeticKill\_two\_pow\_four}{TotientActualLcmShortKill.lean:18} & fixed & $a{=}4$, $a{=}6$ (concrete, formally checked) & M\"obius sums\\
square-CRT coprime-factor examples: zero / nonzero differences & \lean{squareCRT\_clean\_block\_can\_vanish}{SquareCRTCube.lean:459} & fixed & witness $n{=}52$ (vanishes); witness $n{=}27$ (nonzero) & totient factorisation\\

\midrule
\papertablehead\multicolumn{5}{c}{\emph{Statements with cofinal quantifiers: proved constructions and conditional implications}}\\
\midrule

\papertablehead\multicolumn{5}{l}{\emph{Cofinal statements; status specified in the description}}\\
certificate criterion (proved implication; unproved premise) & \lean{irrational\_totient\_series\_of\_certificate\_supply}{TotientTailPeriodKiller.lean:394} & \textbf{cofinal} & $\forall h{\ge}1\,\forall N_0{\ge}0\,\exists N{\ge}N_0\,\exists L,\ \ensuremath{\mathcal{C}}\,h\,N\,L$ & finite binary sums\\
period-multiple criterion (equivalent; unproved condition) & \lean{irrational\_totient\_series\_of\_multiple\_certificate\_supply}{CarrySurvivorExtinction.lean:502} & cofinal & $\forall h_0{>}0\,\forall N_0\,\exists m{>}0\,\exists N{\ge}N_0\,\exists L$ & multiples of periods\\
LCM-diagonal criterion (unproved condition) & \lean{irrational\_totient\_series\_of\_lcm\_diagonal\_certificate\_supply}{LcmDiagonalReduction.lean:112} & cofinal & $\forall t_0\,\exists t{\ge}t_0\,\exists L$ & LCM diagonal\\
LCM-grid certificate criterion (unproved premise) & \lean{irrational\_totient\_series\_of\_lcm\_cone\_window\_kill\_supply}{CertificateKernel.lean:19014} & cofinal & $\forall t_0\,\exists t{\ge}t_0\,\exists q{>}0\,\exists m\,\exists L$ & LCM grids\\
LCM nonintegrality criterion (unproved condition) & \lean{irrational\_totient\_series\_of\_lcm\_diagonal\_nonintegrality\_supply}{CertificateKernel.lean:19034} & cofinal & $\forall t_0\,\exists t{\ge}t_0$, $\ensuremath{{R}}(2H_t){-}\ensuremath{{R}}(H_t)\notin\mathbb{Z}$ & real tail differences\\
finite-grid residue criterion (unproved premise) & \lean{irrational\_totient\_series\_of\_lcm\_cone\_nonflat\_supply}{CertificateKernel.lean:19140} & cofinal & for every $t_0$, data $t\ge t_0,L,Q$ satisfying Theorem~\ref{catalogue:cert:b10b} & finite LCM grids\\
first-harmonic norm criterion (unproved premise) & \lean{irrational\_totient\_series\_of\_first\_harmonic\_norm\_gap}{FirstHarmonicPivot.lean:83} & cofinal & $\forall h{>}0\,\forall X_0\,\exists X{\ge}\max(X_0,1)\,\exists L$, room $16(2X{+}h{+}L{+}2){\le}2^L$ & binary-weighted series\\
four-tail residue criterion (unproved premise) & \lean{irrational\_totient\_series\_of\_primeJumpSharpKill\_supply}{PrimeJumpWindow.lean:193} & cofinal & $\forall t_0\,\exists t{\ge}t_0\,\exists p,L{>}0$ & binary-weighted series\\
short LCM-window criterion (unproved premise) & \lean{PowerTwoActualLcmShortArithmeticKillSupply}{TotientActualLcmOrbitArithmetic.lean:2107} & cofinal & $\forall a_0\,\exists a\,L,\ a_0{\le}a\land L{<}2\cdot2^a$ & M\"obius sums\\
LCM-tail separation criterion (unproved premise) & \lean{irrational\_totientSeries\_of\_actualLcmOrbitSeparationSupply}{TotientActualLcmOrbitSeparation.lean:305} & cofinal & cofinal $1/32{+}$error separation at guarded odd ranks & M\"obius sums\\
five endpoint criteria (unproved premises) & \lean{irrational\_totientSeries\_of\_flexibleActualTerminalDominanceSupply}{TotientActualLcmTopEdgeStaircase.lean:2221} & cofinal & $\forall a_0\,\exists a{\ge}a_0\,\exists\langle\text{supply predicate}\rangle$ & M\"obius sums\\
proved rational counterexample to bound, parity and aperiodicity & \lean{exists\_totientParity\_arbitrarilyManySeparatedCarry\_rational\_countermodel}{TotientParityCoboundaryCountermodel.lean:637} & cofinal & $\exists c{:}\mathbb{N}\to\mathbb{N}$, $c{\le}6$, $c{\le}n$, $c{\equiv}\varphi\!\!\pmod2$, arbitrarily-separated blocks & binary-weighted series\\
proved mod-4 prime construction (CRT and Dirichlet) & \lean{exists\_prime\_deltaTotient\_four\_mul\_mod\_four\_two}{TotientTailCarryPeriod.lean:384} & cofinal & $\forall h{>}0\,\forall B\,\exists\text{prime}\,p{>}B$ & M\"obius sums\\
sufficient mod-4 condition (unproved premise) & \lean{irrational\_totientSeries\_of\_cofinal\_modFourPulseSurvivorKill}{TotientTailCarryPeriod.lean:1066} & cofinal & the displayed mod-four residue condition at arbitrarily large scales & M\"obius sums\\
proved two-adic prime construction & \lean{exists\_prime\_totient\_twoAdic\_pulse\_divisors}{TotientTwoAdicPulseBlock.lean:54} & cofinal & $\forall K{\ge}2\,\forall H{>}K\,\forall B\,\exists\text{cofinally many primes}\,p{>}B$ & divisibility and linear algebra\\
proved implication from eventual integrality & \lean{eventual\_integral\_tailDiff\_has\_cofinal\_twoAdic\_half\_pulse}{TotientTwoAdicPulseBlock.lean:327} & cofinal & $\forall K{\ge}2$, given the stated two-adic prime construction and eventual integrality & divisibility and linear algebra\\
\end{longtable}
\end{landscape}
\endgroup

\noindent $^\dagger$ In the positivity statement, only the window parameter
$J$ is bounded relative to $2\cdot2^a$. The exponent $a$ ranges over all
$a\ge8$, and the constants $4$, $8$ and $32$ are independent of $a$.
The bound on $J$ is a hypothesis of a uniform theorem; it is not a
finite verification of the exponent range.

\subsection{What extending a finite calculation would require}

A finite computation can be extended by further computation or by a
uniform argument. The following examples distinguish these possibilities
from parameter-uniform theorems that are already proved. They do not
assert that a computational bound is intrinsic to a method.

\begin{enumerate}[leftmargin=*]
\item \textbf{Finite certificate tables.} Results such as
\lean{certifiedKill\_diagonal\_all\_imported}{DiagonalPincerCertificates.lean:2870} (28 explicit
values of $t$ through $t{=}64$, now a strict subset of the aggregate
$\forall t\le82$ band),
\lean{certifiedKill\_all\_small}{TotientTailPeriodKiller.lean:404} ($h\in[1,8]$ at fixed depth
16), and
\lean{totient\_series\_ne\_rat\_of\_den\_dvd\_pow\_two\_mul\_mersenne\_upto\_sixteen}{CertificateKernel.lean:18890}
($h\in[1,16]$ at exponent 14) are each a finite \emph{list} of independently verified rows, not a
single argument instantiated at a free parameter. The existing lists do not
establish further rows; an extension needs either additional computations
or an argument for a family of parameters.
\item \textbf{A bounded window inside a uniform theorem.} The results
\lean{three\_quarters\_mul\_lt\_totient\_of\_rough\_lt\_two\_pow}{TotientActualLcmOrbitArithmetic.lean:768}
and
\lean{actualLcmTailDiff\_shift\_pos}{TotientActualLcmOrbitSign.lean:39}
hold for every $a\ge8$, subject to their stated bounds on the other
parameters. In the sign theorem, for example, $J+a+6<2\cdot2^a$.
Their proofs are therefore not merely enumerations of a fixed set of
exponents. The missing conclusion is a suitable residue inequality at
arbitrarily large scales, not extension of the range of $a$ in the
positivity theorem. No complexity lower bound follows from this audit.
\item \textbf{Parameter dependence in an approximation estimate.}
For reduced approximants $p_K/q_K$ to $S$, the conditions
$0<|S-p_K/q_K|\le\varepsilon_K$ and
$q_K\varepsilon_K\to0$ imply irrationality. A growing denominator
alone does not control this product. In particular, the exact
reduced-denominator formula
\lean{lcmHeight\_scaledMobiusShadow\_den\_exact}{MersenneShadowDenominatorGrowth.lean:147}
concerns a specified finite M\"obius sum; it does not estimate the
complementary tail. Likewise, a bound proved with a constant depending
on a fixed auxiliary parameter is not automatically uniform when that
parameter grows. Both the dependence of the constant and the error
estimate must be retained in any passage to an unbounded family.
\end{enumerate}

The missing arithmetic inputs include the stated constant-saving
exponential-sum estimate and an upper-endpoint residue inequality.
For the positivity theorem, the unresolved step is exclusion of that
residue interval, not positivity at more exponents.

\section{Relations between the representations}

The formulas below express the same series through binary tails,
M\"obius sums, rational approximations and exponential sums. Each
identity permits a change of variables or representation. A limitation
of one proposed argument does not automatically apply after that change:
its hypotheses must be checked for the new objects. The last subsection
illustrates this distinction with finite-state decoding and polynomial
denominator identities.

\begin{description}[leftmargin=*,style=nextline]

\item[Binary-weighted series and tails.]
The coefficients in $S=\sum_{n\ge1}\varphi(n)2^{-n}$ are unbounded;
this expression is not the base-two digit expansion of $S$.
The tail $R_N$, finite difference numerator $D(h,N,L)$ and test
$\mathcal C(h,N,L)$ retain this coefficient information. The identity
\[
 S=\frac12+\sum_{d\ge1}\frac{\mu(d)}{(2^d-1)^2}
\]
\lean{tsum\_totient\_half\_pow\_eq\_half\_add\_moebius\_sq}{MersenneLambertLadder.lean:665}
and its partial-sum version
\lean{mobius\_square\_series\_eq\_partial\_add\_tail}{SquaredMersenneDiagonalEnclosure.lean:94}
express the same value by M\"obius sums. An equality of values does not
by itself transfer a claim about the original coefficient subsequences.

\item[Lambert sums and finite numerator polynomials.]
The Lambert identities for $\mu$, $1$ and $\varphi$ give the series
relations in Section~\ref{sec:series}. For a finite divisor sum, the
polynomial evaluation
\lean{mobiusNumeratorPolynomial\_eval\_two}{RepunitMobiusNumerator.lean:445}
and the definition
\lean{mobiusNumerator}{RadicalMobiusShadow.lean:101}
put its numerator into integer polynomial arithmetic. This evaluates a
finite contribution, not the complementary infinite tail.

\item[Coprimality probability.]
Let $X,Y$ be independent with $\Pr(X=n)=\Pr(Y=n)=2^{-n}$ for $n\ge1$.
Then $S-\tfrac12=\Pr(\gcd(X,Y)=1)$. The divisibility probabilities are
$\Pr(d\mid X,\ d\mid Y)=(2^d-1)^{-2}$
\lean{tsum\_pos\_pair\_both\_dvd\_half\_eq\_inv\_mersenne\_sq}{GcdMomentCalculus.lean:266}.
M\"obius inversion converts these to the coprimality probability.
The decomposition by the greatest common divisor also gives the weighted
coprime-pair identity
\lean{tsum\_pos\_coprime\_inv\_mersenne\_eq\_one}{GcdMomentCalculus.lean:349}.
Counting coprime pairs by their sum gives the totient expression directly
\lean{tsum\_coprime\_pair\_pow\_eq\_tsum\_totient\_mul\_pow}{GeometricCoprimality.lean:120}.

\item[Cyclotomic factors of finite denominators.]
The coefficient formula for $P_r$
\lean{mobiusNumeratorPolynomial\_coeff}{RepunitMobiusNumerator.lean:217}
and the cyclotomic divisibility statement
\lean{cyclotomicValue\_dvd\_baseMobiusShadow\_den}{CyclotomicProjectionOfShadow.lean:389}
identify factors in the reduced denominators of the specified finite sums.
The product-divisibility theorem
\lean{upperHalfChannel\_product\_dvd\_den\_of\_scale\_primeFactors\_le}{MersenneShadowCyclotomicNoncollapse.lean:809}
is used in the exact denominator formula
\lean{lcmHeight\_scaledMobiusShadow\_den\_exact}{MersenneShadowDenominatorGrowth.lean:147}.
Growing finite denominators alone do not exclude cancellation with the
remaining tail.

\item[Numerator divisibility and denominator bounds.]
Write $x=a/b$ in lowest terms, with $b>0$. For $c\in\Z$ and a
positive divisor $H$ of $b$, if the reduced denominator of $cx$ divides
$H$, then $b/H\mid c$
\lean{scalarLocalization\_complement\_dvd}{AdelicHeightObstruction.lean:23}.
In particular, $|c|\ge b/H$ when $c\ne0$. A related bound is
\[
 0<\frac ab<\frac2{2^n-1},\quad 2^r\mid a,\quad n\ge1
 \quad\Longrightarrow\quad 2^r(2^n-1)<2b
\]
for integers $r\ge0$
\lean{positiveRat\_mersenne\_height}{AdelicHeightObstruction.lean:103}.
Here positivity gives $a\ge2^r$; the size inequality then gives the
conclusion. These are arithmetic bounds on a reduced fraction, not
irrationality criteria without an additional approximation estimate.

\item[Farey intervals.]
The denominator bound
\lean{farey\_gap}{GapFareyBound.lean:51}
applied to the explicit window $(N,K)=(1,240)$ gives
\lean{gap\_check\_window\_1\_240\_le\_79639646646701375323355774875831053}{GapFareyBound.lean:176}
and hence $S\ne p/q$ for integers $p$ and $1\le q\le Q_0$.
The finite input is a totient numerator modulo $2^{240}$.
Its exact first failing denominator is $Q_0+1$, the denominator of the
relevant Farey mediant
\lean{gap\_check\_window\_1\_240\_first\_failure}{GapFareyBound.lean:225}.
This last assertion concerns this fixed window.

\item[Differences at least-common-multiple shifts.]
For $H_t=\operatorname{lcm}(1,\ldots,t)$
\lean{periodLcm}{CarrySurvivorExtinction.lean:515},
a jump in $H_t$ can occur only at a prime power
\lean{eq\_prime\_pow\_of\_not\_dvd\_periodLcm}{LcmDiagonalReduction.lean:137}.
The identities
\lean{lcmRayArithmeticLetter\_eq\_deltaTotient}{TotientActualLcmOrbitArithmetic.lean:298}
and
\lean{diagonalWindowIncrement\_eq\_lcmRayArithmeticLetter}{TotientActualLcmOrbitArithmetic.lean:310}
identify the coefficient difference as
$\delta_t(j)=\varphi(2H_t+j)-\varphi(H_t+j)$ and give
$D(H_t,H_t,L)=\sum_{j=1}^L\delta_t(j)2^{L-j}$.
Thus these arithmetic formulas concern particular instances of the same
finite residue test; additional window restrictions remain in force.

\item[Exponential sums of the window residues.]
For integers $h,X\ge1$, the relevant block statistic is
\[
 \sum_{N=X}^{2X-1}\cos\!\left(2\pi\frac{D(h,N,L)}{2^L}\right).
\]
A bound by $9X/10$, together with $L\ge h$ and
$16(2X+h+L+2)\le2^L$, yields a certificate in the block
\lean{exists\_certifiedKill\_of\_first\_harmonic\_gap}{FirstHarmonicGap.lean:128}.
The linked irrationality implication uses the stronger complex norm
bound $21X/25$, for every $h$ on arbitrarily late blocks satisfying
the same size inequality
\lean{irrational\_totient\_series\_of\_first\_harmonic\_norm\_gap}{FirstHarmonicPivot.lean:83}.
The finite implication uses $\cos(\pi/8)>9/10$; the required cancellation
bound is still unproved. The norm bound on the complex exponential sum
is sufficient for the displayed real-part bound, but is stronger.

\end{description}

\subsection{Binary tails and M\"obius sums}

Erd\H{o}s's argument for $\sum_n1/(2^n-1)$ chooses congruences so
that divisor counts along a block are divisible by increasing powers of
the base. A tail bound then forces a long run of zero digits
\cite[pp.~63--66]{erdos1948lambert}. It uses no finite automaton.

A different proposal would recover each later block from a fixed finite
summary of the preceding history. The balanced-pulse counterexample
\lean{balancedPulse\_no\_autonomous\_decoder}{GenericTailOrbitRigidity.lean:247}
shows why that proposal fails for the family constructed there. For
$m\ge2$, its members have the same earlier history but different
successors. A state that identifies all those histories cannot recover
every successor. More generally, an exact label-and-decoder pair for
that family needs at least $\lfloor(m+1)/2\rfloor+1$ distinct labels.
This is a lower bound on the number of labels, not a linear lower bound
on the number of bits, and it is not a state-complexity theorem for the
actual totient sequence.
\ev{Lean} \scale{uniform} \coord{binary-digit}.

The cofinal certificate implication
\lean{irrational\_totient\_series\_of\_certificate\_supply}{TotientTailPeriodKiller.lean:394}
is a separate theorem. The counterexample above neither proves its
hypothesis nor makes it the only possible route to irrationality.
Likewise,
\lean{det\_residualMonomialMatrix}{ResidualGaugeObstruction.lean:47}
shows that inverse-phase column weights can preserve a determinant, and
\lean{linearIndependent\_canonicalTotientKernelFamily}{TotientMahlerDefect.lean:935}
gives the exact growing ranks of totient sections. These are different
statements with different hypotheses; none is a general prohibition on
other representations of $S$.

The M\"obius representation instead allows the identity
\lean{tsum\_lambert\_linear\_weight\_sq\_pure}{GcdMomentCalculus.lean:105}
to be applied to $\mu$, since $|\mu(d)|\le d$ for $d\ge1$. The
polynomial identities from
\lean{mobiusNumeratorPolynomial\_eq\_gcdWord}{RepunitMobiusNumerator.lean:205}
to
\lean{cyclotomic\_dvd\_primeJump\_new\_fibre}{PrimePowerJumpDynamics.lean:281}
operate at each finite radical without an autonomous decoder. The exact
reduced denominator and its lower bound are then given by
\lean{lcmHeight\_scaledMobiusShadow\_den\_exact}{MersenneShadowDenominatorGrowth.lean:147}
and
\lean{upperHalfMersenneProduct\_lower\_bound}{MersenneShadowDenominatorGrowth.lean:60}.
Those formulas concern a finite rational contribution. They do not
exclude cancellation by the complementary tail and hence do not settle
irrationality. The representations are equal expressions for $S$;
what differs is the information supplied by the particular arguments.
% ; - part a249_p3 ; -
\section{Index of unproved conditions}

This section collects the unproved hypotheses used by the conditional
results above. For each one, it states the additional arithmetic
assertion and the implication to $S\notin\mathbb Q$. A proof of an
implication is not a proof of its hypothesis. The record does not settle
whether $S=\sum_{n\ge0}\varphi(n)2^{-n}$ is irrational.

\subsection{How to use the index}

Each entry separates the proved implication from the additional
hypothesis needed to apply it. The comparison concerns the displayed
statements, not the prospects of an entire method.

The range tags describe quantification, not a ranking of mathematical
strength. \scale{fixed} records finitely many explicit instances;
\scale{bounded} records the displayed bounded parameter range;
\scale{uniform} ranges over every allowed value of the stated parameters;
and \scale{cofinal} includes a quantifier of the form
$\forall a_0\ \exists a\ge a_0$. All other hypotheses and quantifiers
remain part of the statement. A cofinal assertion about a different
quantity, or with a different order of quantifiers, need not supply the
irrationality condition. \scale{n/a} denotes the absence of a relevant
growing parameter. Proof status is recorded separately by
\ev{Lean}, \ev{Cert}, \ev{Math}, \ev{Cited} and \ev{Open}, as in the
reading guide; a uniform or cofinal statement may itself be proved or
may remain an unproved hypothesis.

The quantified residue condition for \#249, stated in Definition~\ref{defn:sep} and proved sufficient
(\lean{irrational_totient_series_of_certificate_supply}{TotientTailPeriodKiller.lean:394}), is the
following statement:
\[
  \forall h \ge 1,\ \forall N_0,\ \exists N \ge N_0,\ \exists L,\ \mathcal{C}\ h\ N\ L .
\]
The following catalogue compares fourteen proposed approaches to this
condition. It starts with the exponential-sum estimate, then gives thirteen
further approaches with their individual hypotheses and source references.

\subsection{The exponential-sum criterion}
\label{ssec:headline}

The next three definitions recall the finite sum, its residue test and
its phase. The following theorems explain which estimates for that phase
would imply the quantified condition in Definition~\ref{defn:sep}.

\begin{defn}[The truncated difference]
\lean{windowDiscrepancy}{TotientTailPeriodKiller.lean:66}. For $h,N,L\in\mathbb N$,
\[
  D(h,N,L) \;=\; \sum_{j=0}^{L-1} \bigl(\varphi(N+h+1+j) - \varphi(N+1+j)\bigr)\cdot 2^{\,L-1-j}
  \ \in \mathbb Z .
\]
This is the depth-$L$ truncation of $2^L\cdot(R_{N+h}-R_N)$, where
$R_N = \sum_{j\ge 1}\varphi(N+j)/2^j$ is the local totient tail
(\lean{totientTail}{TotientTailPeriodKiller.lean:57}).
\end{defn}

\begin{defn}[The central-residue condition]
\lean{certifiedKill}{TotientTailPeriodKiller.lean:72}.
\[
  \mathcal{C}\ h\ N\ L \;:\Leftrightarrow\;
  (N{+}h{+}L{+}2) < D(h,N,L)\bmod 2^L < 2^L - (N{+}h{+}L{+}2).
\]
That is: the residue of $D(h,N,L)$ modulo $2^L$ avoids the radius-$(N{+}h{+}L{+}2)$ neighbourhood
of $0$. \lean{certifiedKill_depth_floor}{TotientTailPeriodKiller.lean:79} records the necessary
room condition this forces: $2(N{+}h{+}L{+}2) < 2^L$.
\end{defn}

\begin{defn}[The complex phase and its real part]
\lean{windowFirstCos}{FirstHarmonicGap.lean:27} and
\lean{windowFirstExp}{FirstHarmonicPivot.lean:38}:
\[
  \operatorname{Re}E(h,N,L) = \cos\!\Bigl(2\pi\,\tfrac{D(h,N,L)\bmod 2^L}{2^L}\Bigr),\qquad
  E(h,N,L) = \exp\!\Bigl(i\,2\pi\,\tfrac{D(h,N,L)\bmod 2^L}{2^L}\Bigr).
\]
\lean{windowFirstExp_re}{FirstHarmonicPivot.lean:42} records $\mathrm{Re}({E})
= \operatorname{Re}E$ and \lean{norm_windowFirstExp}{FirstHarmonicPivot.lean:48} that
$|E(h,N,L)| = 1$: this is literally the first additive character of the
endpoint discrepancy modulo $2^L$, a standard exponential-sum object.
\end{defn}

\begin{thm}[A real-part bound gives a certificate]
\label{thm:hgap-real}
\lean{exists_certifiedKill_of_first_harmonic_gap}{FirstHarmonicGap.lean:128}. \ev{Lean}.
\scale{uniform}. \coord{first-harmonic}.
For all $h,X,L$ with $0<X$ and the room condition $16(2X{+}h{+}L{+}2)\le 2^L$, if
\[
  \sum_{N=X}^{2X-1} \operatorname{Re}E(h,N,L) \;\le\; \tfrac{9}{10}\,X ,
\]
then $\exists N\in[X,2X)$ with $\mathcal{C}\ h\ N\ L$.
\end{thm}
\leannote{Lean: \lproof{ErdosProblems/Erdos249/PaperCompleteR21/FirstHarmonicBlockCriteria.lean}{23}{windowFirstCos_unfolded}, \lproof{ErdosProblems/Erdos249/PaperCompleteR21/FirstHarmonicBlockCriteria.lean}{34}{exists_certifiedKill_of_block_real_part_bound}.}

\begin{thm}[The same implication for a nonempty subset]
\label{thm:hgap-subset}
\lean{exists_certifiedKill_of_first_harmonic_gap_subset}{FirstHarmonicGap.lean:104}.
\ev{Lean}. \scale{uniform}. \coord{first-harmonic}.
For any nonempty finite $T\subseteq\mathbb N$ with $T\subset[0,2X)$ and the same room condition, if
\[
  \sum_{N\in T} \operatorname{Re}E(h,N,L) \;\le\; \tfrac{9}{10}\,|T| ,
\]
then $\exists N\in T$ with $\mathcal{C}\ h\ N\ L$. This generalises
Theorem~\ref{thm:hgap-real}: $T$ can be \emph{any} explicitly chosen nonempty finite subset of the
dyadic block, not the whole block, and the proof (an averaging pigeonhole, \S below) never uses
that $T$ has positive density or comes from a partition.
\end{thm}
\leannote{Lean: \lproof{ErdosProblems/Erdos249/PaperCompleteR21/FirstHarmonicBlockCriteria.lean}{50}{exists_certifiedKill_of_subset_real_part_bound}, \lproof{ErdosProblems/Erdos249/PaperCompleteR21/FirstHarmonicBlockCriteria.lean}{65}{block_real_part_bound_of_subset_form}.}

\begin{thm}[A norm bound gives the real-part criterion]
\label{thm:hgap-norm}
\lean{exists_certifiedKill_of_first_harmonic_norm_gap}{FirstHarmonicPivot.lean:53} and
\lean{irrational_totient_series_of_first_harmonic_norm_gap}{FirstHarmonicPivot.lean:83}.
\ev{Lean} for both implications; the hypothesis \m{block norm condition} is \ev{Open}.
The complex norm bound implies the real-part bound
(\m{|z|\ge\mathrm{Re}(z)}, and $21/25 < 9/10$ absorbs the slack), so it composes through
Theorem~\ref{thm:hgap-real} to the same certificate. Define
\[
  \text{block norm condition} :\Leftrightarrow\;
  \forall h{>}0\ \forall X_0\ \exists X, L,\ \max(X_0,1)\le X\ \wedge\ 16(2X{+}h{+}L{+}2)\le 2^L\ \wedge
\]
\[
  \Bigl\|\ \sum_{N=X}^{2X-1}E(h,N,L)\ \Bigr\| \;\le\; \tfrac{21}{25}\,X .
\]
\coord{first-harmonic}. \scale{cofinal} (this is the open target itself). Then
\[
  \text{block norm condition} \;\Longrightarrow\; \mathrm{Irrational}\Bigl(\sum_{n\ge 0}
  \tfrac{\varphi(n)}{2^n}\Bigr),
\]
proved in full, with no further gap, by chaining
Theorem~\ref{thm:hgap-norm}$\to$Theorem~\ref{thm:hgap-real}$\to$
\lean{irrational_totient_series_of_certificate_supply}{TotientTailPeriodKiller.lean:394}.
\end{thm}
\leannote{Lean: \lproof{ErdosProblems/Erdos249/PaperCompleteR21/FirstHarmonicBlockCriteria.lean}{86}{blockNormCondition_unfolded}, \lproof{ErdosProblems/Erdos249/PaperCompleteR21/FirstHarmonicBlockCriteria.lean}{97}{real_part_bound_of_norm_bound}, \lproof{ErdosProblems/Erdos249/PaperCompleteR21/FirstHarmonicBlockCriteria.lean}{110}{exists_certifiedKill_of_block_norm_bound}, \lproof{ErdosProblems/Erdos249/PaperCompleteR21/FirstHarmonicBlockCriteria.lean}{118}{irrational_of_blockNormCondition}.}

\begin{obs}[What remains unproved]
\label{obs:no-instance}
The implications above do not establish the cofinal block estimate.
The following callers retain a separation or gap hypothesis: \lean{AdjacentPhaseSeparation.lean}{AdjacentPhaseSeparation.lean:267}
uses the two-element set $\{N,M\}$, while
\lean{PivotAntiReconstruction.lean}{PivotAntiReconstruction.lean:882}
and
\lean{PivotAntiReconstruction.lean}{PivotAntiReconstruction.lean:1572}
use, respectively, a supplied finite set $T$ and the singleton $\{N\}$.
They prove what follows from the relevant estimate; they do not provide
that estimate at arbitrarily large scales. The finite numerical
observations in Table~\ref{tab:blocks} are a different kind of evidence
and do not discharge this quantified hypothesis. \ev{Open}
\end{obs}

\begin{rem}[What the reduction does and does not prove]
The required savings are explicit: a real-part bound of $9X/10$, or
the stronger norm bound of $21X/25$, on blocks at arbitrarily large
scales for every positive shift $h$. An unspecified nonzero saving is
not the stated hypothesis. A bound $o(X)$ would suffice, but it is more
than the implication needs.

The inequality $16(2X+h+L+2)\le2^L$ can be met by increasing $L$ for
fixed $X,h$. Doing so changes the exponential sum, and hence does not
preserve an estimate proved at another depth. The arithmetic estimate
must hold at a depth satisfying this inequality.

The subset theorem permits any \emph{nonempty} finite sample inside
$[0,2X)$; no density assumption is required. A prime-defined sample is
one possible choice. The decomposition in
\lean{FirstHarmonicPivot.lean}{FirstHarmonicPivot.lean} isolates one prime
factor and the remaining totient-dependent phase, but does not prove
the required cancellation between them.
\end{rem}

\subsection{Proved results and the additional assumptions they need}

Section~\ref{ssec:headline} gives the exponential-sum condition, listed
as \textbf{e1-companion} in the source catalogue. The remaining entries
compare tail signs, endpoint residues, finite approximations, short windows,
Farey bounds, divisibility constructions, denominator growth and carry
rank. Each states the extra assumption needed for its proposed application.
The entries are not a classification of all possible approaches to \#249.

\subsubsection{Positivity and the remaining residue inequality}

\begin{defn}
The \emph{actual LCM tail orbit} at exponent $a$ is
$\Omega_a = R_{2H}-R_H
=R_{2H}-R_{H}$, where $H = {H}(2^a)$
(\lean{actualLcmTailOrbit_pos}{TotientActualLcmOrbitSign.lean:144}).
\end{defn}

\begin{thm}
\lean{actualLcmTailDiff_shift_pos}{TotientActualLcmOrbitSign.lean:39} and
\lean{actualLcmTailOrbit_pos}{TotientActualLcmOrbitSign.lean:144}. \ev{Lean}. \scale{uniform}.
For every $a\ge 8$ and every $J$ with $J+(a{+}6) < 2\cdot 2^a$,
\[
  0 \;<\; {R}(2H{+}J) - {R}(H{+}J), \qquad H={H}(2^a).
\]
Hence $0 < \Omega_a$ for every $a\ge 8$. The proof is genuinely uniform in
$a$: it rests on two facts proved for all $a\ge 8$ by a two-case structural split (divisor letter
vs.\ foreign prime power), with absolute constants $4,8,32$ and no lookup table; \lean{lcmRayArithmeticLetter_pos_of_lt_two_mul}{TotientActualLcmOrbitArithmetic.lean:1703} and
\lean{periodLcm_lt_eight_mul_t_mul_lcmRayArithmeticLetter}{TotientActualLcmOrbitArithmetic.lean:1906}.
\end{thm}
\leannote{Lean: \lproof{ErdosProblems/Erdos249/PaperCompleteR21/TopEdgeCorridorAndSeparation.lean}{39}{actualLcm_corridor_pos}, \lproof{ErdosProblems/Erdos249/PaperCompleteR21/TopEdgeCorridorAndSeparation.lean}{49}{corridor_letter_pos}, \lproof{ErdosProblems/Erdos249/PaperCompleteR21/TopEdgeCorridorAndSeparation.lean}{56}{corridor_height_lt_letter}.}

\begin{rem}[Which endpoint interval must be excluded]
Under the hypotheses of Proposition~\ref{prop:SGN-03}, integrality would
force the residue to be $2^K-e$, where
$e=R_{2H+J+K}-R_{H+J+K}$ is a positive integer smaller than
$2H+J+K+2$. Thus the residue lies near the upper endpoint, not near zero
in its least nonnegative representation. The representative $-e$ is
negative; the carry $e$ is positive.

The one-sided test below excludes precisely that upper interval. This
uses the bound $J+K+(a+6)<2\cdot2^a$, as well as the modulus bound.
It does not follow from positivity that the residue already satisfies
half of a nonintegrality certificate without those assumptions.
\par\noindent\ev{Lean} \scale{uniform}
\lean{actualLcm_integral_forces_topEdgeResidue}{TotientActualLcmOrbitSign.lean:211}
\lean{carryOrbit}{CarrySurvivorExtinction.lean:387}
\end{rem}

\subsubsection{An exact endpoint identity}

The sign estimate permits a one-sided residue test under the bounds
below. This is a sufficient criterion, with additional arithmetic
hypotheses, rather than a claimed strict weakening of irrationality.

\begin{thm}[The proved implications]
\lean{actualLcmTailDiff_notMem_int_of_topEdgeResidueGap}{TotientActualLcmTopEdgeStaircase.lean:1260},
\lean{actualLcmTailOrbit_notMem_int_of_topEdgeResidueGap}{TotientActualLcmTopEdgeStaircase.lean:1314},
\lean{irrational_totientSeries_of_topEdgeResidueGapSupply}{TotientActualLcmTopEdgeStaircase.lean:2184}.
\ev{Lean}. \scale{uniform}. For every $a\ge 8$ and every $J,K,m$ inside the sign corridor
($J{+}K{+}(a{+}6) < 2\cdot 2^a$), a one-sided residue gap at precision $m\le K$; room $2H{+}J{+}K{+}2 < 2^m$ and $D(H,H+J,K)\bmod 2^m \le 2^m - (2H{+}J{+}K{+}2)$; already
forces ${R}(2H{+}J) - {R}(H{+}J)\notin\mathbb Z$. No lower margin
at all is demanded. The proof chain is complete: the theorem holds for every $a\ge 8$, and
\emph{\m{\text{cofinal upper-endpoint condition}} $\Rightarrow$ Irrational $S$} is proved.
\end{thm}
\leannote{Lean: \lproof{ErdosProblems/Erdos249/PaperCompleteR21/TopEdgeCorridorAndSeparation.lean}{64}{topEdgeResidueGap_unfolded}, \lproof{ErdosProblems/Erdos249/PaperCompleteR21/TopEdgeCorridorAndSeparation.lean}{75}{topEdgeResidueGap_forces_nonintegral}, \lproof{ErdosProblems/Erdos249/PaperCompleteR21/TopEdgeCorridorAndSeparation.lean}{83}{topEdgeResidueGap_orbit_nonintegral}, \lproof{ErdosProblems/Erdos249/PaperCompleteR21/TopEdgeCorridorAndSeparation.lean}{92}{irrational_of_topEdgeResidueGapSupply}.}

\begin{defn}[The unproved quantified condition]
\lean{PowerTwoActualLcmTopEdgeResidueGapSupply}{TotientActualLcmTopEdgeStaircase.lean:1325}.
\ev{Open}. \scale{cofinal}.
\[
  \forall a_0\ \exists a,K,m,\ a_0\le a\ \wedge\ 8\le a\ \wedge\ K{+}(a{+}6)<2\cdot 2^a\ \wedge\
  \mathcal G_a(0,K,m) .
\]
No unconditional statement anywhere in the record locates the residue of the diagonal word
$D(H,H,K)$ (where $H={H}(2^a)$) below the top strip at even one large $a$.
\lean{actualLcmTopEdgeResidueGap_iff_terminal}{TotientActualLcmTopEdgeStaircase.lean:909} further
reduces the condition to the last $m$ arithmetic letters of the word alone.
\end{defn}

\begin{prop}[Five sufficient conditions]
\label{prop:te-chain}
Each of the five conditions below suffices for irrationality. The first
four imply the upper-endpoint condition
\lean{PowerTwoActualLcmTopEdgeResidueGapSupply}{TotientActualLcmTopEdgeStaircase.lean:1325};
the fifth gives nonintegrality directly by the endpoint identity. These
are not asserted to form a linear hierarchy. Here $H=H(2^a)$, and every
condition quantifies over arbitrarily large exponents $a$.
\begin{enumerate}[leftmargin=*]
  \item \lean{ErdosProblems.Erdos249.PaperCompleteR21.PaperAdjacentSuffixMidbandSupply}{lean/ErdosProblems/Erdos249/PaperCompleteR21/TopEdgeChainPaperBand.lean:46}
    (\ev{Open}, \scale{cofinal}):
    \[
      \forall a_0\ \exists a,m,\ a_0\le a \wedge 8\le a \wedge m{+}1{+}(a{+}6)<2\cdot 2^a \wedge
      2H{+}m{+}3 < 2^m \wedge
    \]
    \[
      2H{+}m{+}2 \le \text{suffix residue}(2^a)\,0\,m \le 2^m - (2H{+}m{+}2)
    \]
    (a two-sided band on the adjacent-suffix residue directly, one candidate depth $m$).
  \item \lean{PowerTwoOddGuardTopEdgeHalfWordBandSupply}{TotientActualLcmTopEdgeStaircase.lean:1349}
    (\ev{Open}, \scale{cofinal}): at the odd guarded depth $2q{+}1$, a two-sided band of half-width
    $H{+}q{+}2$ on the half-word residue modulo $4^q$. The depth is fixed
    by the stated guard; it cannot be chosen independently of $a$.
  \item \lean{PowerTwoActualFinalTopEdgeMagnitudeSupply}{TotientActualLcmTopEdgeStaircase.lean:1471}
    (\ev{Open}, \scale{cofinal}), \emph{proved equivalent} to the previous one via
    \lean{topEdgeHalfWordBandSupply_iff_actualFinalCenteredMagnitudeSupply}{TotientActualLcmTopEdgeStaircase.lean:1479}:
    \[
      \forall a_0\ \exists a,q,\ \max(14,a_0)\le a \wedge q=q_a
      \wedge H{+}q{+}2 \le |u_{a,q}| .
    \]
  \item \lean{PowerTwoFlexibleActualTopEdgeMagnitudeSupply}{TotientActualLcmTopEdgeStaircase.lean:1927}
    (\ev{Open}, \scale{cofinal}): the same magnitude bound with the depth restriction relaxed from
    ``canonical guarded'' to any odd $2q{+}1$ satisfying the half-cell fit
    $2(H{+}q{+}2)\le 4^q$ and the sign-corridor room $2q{+}2{+}(a{+}6)<2\cdot 2^a$.
  \item \lean{PowerTwoFlexibleActualTerminalDominanceSupply}{TotientActualLcmTopEdgeStaircase.lean:1908}: under the same bounds as item 4,
    $\delta_{2^a}(2q+2)\le2u_{a,q}$. This is a different one-sided
    comparison, treated immediately below.
\end{enumerate}
For the first four conditions, the relevant implications into
\lean{PowerTwoActualLcmTopEdgeResidueGapSupply}{TotientActualLcmTopEdgeStaircase.lean:1325}
are as follows:
midband $\to$ residue-gap at \lean{}{TotientActualLcmTopEdgeStaircase.lean:2109}; half-word-band
$\to$ midband at \lean{}{TotientActualLcmTopEdgeStaircase.lean:2058}; final-magnitude
$\Leftrightarrow$ half-word-band at \lean{}{TotientActualLcmTopEdgeStaircase.lean:1479};
flexible-magnitude $\to$ midband at \lean{}{TotientActualLcmTopEdgeStaircase.lean:2008};
final-magnitude $\to$ flexible-magnitude at \lean{}{TotientActualLcmTopEdgeStaircase.lean:1993}.
\end{prop}
\leannote{Lean: \lproof{ErdosProblems/Erdos249/PaperCompleteR21/TopEdgeChainPaperBand.lean}{56}{paperTeChain_item_one_unfolded}, \lproof{ErdosProblems/Erdos249/PaperCompleteR21/TopEdgeCorridorAndSeparation.lean}{113}{teChain_item_two_unfolded}, \lproof{ErdosProblems/Erdos249/PaperCompleteR21/TopEdgeCorridorAndSeparation.lean}{124}{teChain_item_three_unfolded}, \lproof{ErdosProblems/Erdos249/PaperCompleteR21/TopEdgeCorridorAndSeparation.lean}{133}{teChain_item_four_unfolded}, and 11 further declarations in the \hyperref[sec:coverage]{coverage section}.}

\subsubsection{A bound on the last totient difference}

The next condition compares one totient difference with a centred
residue. Its usefulness would depend on proving that comparison for
arbitrarily large $a$ within the required range of $q$.

\begin{defn}
\lean{PowerTwoFlexibleActualTerminalDominanceSupply}{TotientActualLcmTopEdgeStaircase.lean:1908}.
\ev{Open}. \scale{cofinal}. Writing $H={H}(2^a)$,
\[
  \forall a_0\ \exists a,q,\ a_0\le a \wedge 8\le a \wedge 2q{+}2{+}(a{+}6)<2\cdot 2^a \wedge
  2(H{+}q{+}2)\le 4^q \wedge
\]
\[
  \delta_{2^a}(2q{+}2) \;\le\; 2\cdot u_{a,q} .
\]
\end{defn}

\begin{thm}[The proved implication]
\lean{actualLcmTailOrbit_notMem_int_of_actualTerminalDominance}{TotientActualLcmTopEdgeStaircase.lean:1880}
and
\lean{irrational_totientSeries_of_flexibleActualTerminalDominanceSupply}{TotientActualLcmTopEdgeStaircase.lean:2221}.
\ev{Lean}. \scale{uniform}. The dominance hypothesis at a single odd rank already excludes
integrality of \m{\Omega_a}, and the supply predicate composes to
\m{S\notin\mathbb Q}.
\end{thm}
\leannote{Lean: \lproof{ErdosProblems/Erdos249/PaperCompleteR21/TopEdgeCorridorAndSeparation.lean}{217}{terminalDominance_orbit_nonintegral}, \lproof{ErdosProblems/Erdos249/PaperCompleteR21/TopEdgeCorridorAndSeparation.lean}{228}{irrational_of_terminalDominanceSupply}.}

\begin{obs}[How the sign and size bounds enter the contradiction]\label{obs:staircase-tension}
Use the hypotheses and notation of Proposition~\ref{prop:TE-06}. Under
the assumption $z=\Omega_a\in\mathbb Z$, the exact identity is
\[
 2u_{a,q}=\delta_{2^a}(2q+2)-c_{H,H,z}(2q+1),
 \qquad 0<c_{H,H,z}(2q+1)<2H+2q+3.
\]
Terminal dominance $\delta_{2^a}(2q+2)\le2u_{a,q}$ would contradict the
positive lower bound on the carry. The other sufficient inequality,
$2u_{a,q}\le\delta_{2^a}(2q+2)-(2H+2q+3)$, would contradict its upper
bound. Both are valid ways to contradict integrality.

These bounds do not eliminate either approach: they hold under the
integrality assumption that the approach is intended to refute. They
neither prove nor disprove an unconditional estimate for $u_{a,q}$.
Finite examples lying between the two boundaries do not establish an
asymptotic exclusion of either inequality.
\par\noindent\ev{Lean} (identities and conditional bounds)
\lean{two_mul_actualOddHalfCenteredLift_eq_terminal_sub_trueCarry}{TotientActualLcmTopEdgeStaircase.lean:1678}
\lean{actualLcm_trueEndpointSurvivor_neg}{TotientActualLcmOrbitSign.lean:172}
\lean{actualLcmTailOrbit_notMem_int_of_actualTerminalCarryCorridorEscape}{TotientActualLcmTopEdgeStaircase.lean:1824}
\end{obs}

\begin{prop}[A sufficient extension]
Either (i) the lower-escape branch cofinally; as displayed just above; or (ii) the two-sided
magnitude form \lean{PowerTwoFlexibleActualTopEdgeMagnitudeSupply}{TotientActualLcmTopEdgeStaircase.lean:1927}
(item 4 of Proposition~\ref{prop:te-chain}), which asks only $H{+}q{+}2 \le
|u_{a,q}|$ and which the file proves
(\lean{flexibleActualTerminalCarryCorridorEscapeSupply_of_magnitude}{TotientActualLcmTopEdgeStaircase.lean:1938})
implies corridor escape via a clean sign split (positive branch escapes above the terminal letter,
negative branch escapes below the directed bound), without prescribing the sign of the centred representative. Neither
branch is supplied by the conditional sign identity itself.
\end{prop}
\leannote{Lean: \lproof{ErdosProblems/Erdos249/PaperCompleteR21/TopEdgeCorridorAndSeparation.lean}{236}{irrational_of_lower_escape_supply}, \lproof{ErdosProblems/Erdos249/PaperCompleteR21/TopEdgeCorridorAndSeparation.lean}{252}{corridor_escape_and_irrational_of_magnitude}, \lproof{ErdosProblems/Erdos249/PaperCompleteR21/TopEdgeCorridorAndSeparation.lean}{133}{teChain_item_four_unfolded}.}

\subsubsection{Separation of a rational approximation}

\begin{thm}
\lean{abs_actualLcmTailOrbit_sub_rawApprox_lt}{TotientActualLcmOrbitSeparation.lean:141} and
\lean{irrational_totientSeries_of_actualLcmOrbitSeparationSupply}{TotientActualLcmOrbitSeparation.lean:305}.
\ev{Lean}. \scale{uniform}. Unconditionally, for every $a$ and $q$,
\[
  \bigl|\, \Omega_a - \rho_{a,q} \,\bigr|
  \;<\; \frac{4H + 2(2q{+}1) + 4}{2^{2q+2}} ,
\]
where $\rho_{a,q}$ is an explicit finite computable rational block. For fixed $a$, the error tends to zero as $q$ grows. In the quantified
condition below, however, the depth is prescribed, so increasing $q$
is not a free way to satisfy the separation hypothesis.
\end{thm}
\leannote{Lean: \lproof{ErdosProblems/Erdos249/PaperCompleteR21/TopEdgeCorridorAndSeparation.lean}{276}{abs_actualLcmTailOrbit_sub_rawApprox_lt_paper_form}, \lproof{ErdosProblems/Erdos249/PaperCompleteR21/TopEdgeCorridorAndSeparation.lean}{266}{two_pow_odd_eq}, \lproof{ErdosProblems/Erdos249/PaperCompleteR21/TopEdgeCorridorAndSeparation.lean}{312}{actualLcmRawApprox_isRat}, \lproof{ErdosProblems/Erdos249/PaperCompleteR21/TopEdgeCorridorAndSeparation.lean}{294}{actualLcmRawErrorRadius_tendsto_zero}, and 1 further declaration in the \hyperref[sec:coverage]{coverage section}.}

\begin{rem}[Scope of the implication]
A cofinal $1/32$-separation of the raw approximants at the prescribed
indices $q=q_a$ suffices for \#249. Allowing arbitrary depths would be a
different hypothesis. The cited short-window theorem gives
nonintegrality of $\Omega_a$ at $a=4,6$
(\lean{actualLcmTailOrbit_four_and_six_notMem_int}{TotientActualLcmShortKill.lean:35}, \ev{Cert},
\scale{fixed}, via \lean{certifiedKill_diagonal_t16}{DiagonalPincerCertificates.lean:2842} and
\lean{certifiedKill_diagonal_t64}{DiagonalPincerCertificateT64Endpoint.lean:1928}).
It neither asserts the $1/32$ margin at those exponents nor says that
these are the only successful exponents. Indeed, the diagonal table also
gives short-window certificates $(a,L)=(2,7)$ and $(3,15)$, since
$H_4=12$, $H_8=840$, and $7<8$, $15<16$. All these are finite
nonintegrality examples, not the cofinal approximation-separation condition.
\end{rem}

\begin{defn}[The additional hypothesis]
$\text{cofinal approximation-separation condition}$: \ev{Open}, \scale{cofinal}.
\[
 \begin{gathered}
  \forall a_0\ \exists a\ge\max(2,a_0)\ \exists q,\quad q=q_a\ \wedge\\
  \forall z\in\mathbb Z,\quad \tfrac1{32}+\varepsilon_{a,q}
       \le |\Omega_a-z| .
 \end{gathered}
\]
This hypothesis is about the real quantity $\Omega_a$, not just a finite
rational calculation. It implies
$|\rho_{a,q_a}-z|>1/32$ for every integer $z$, by the displayed error
bound and the triangle inequality. A cofinal $1/32$-separation condition
on the finite approximants is another sufficient input, through the
residue-band implication. These sufficient conditions are not asserted
to be equivalent.
\end{defn}

\subsubsection{The short-window examples through exponent 6}

\begin{thm}
\lean{powerTwoActualLcmShortArithmeticKillSupply_through_six}{TotientActualLcmShortKill.lean:54}
and
\lean{irrational_totientSeries_of_shortArithmeticKillSupply}{TotientActualLcmShortKill.lean:63}.
\ev{Cert} for the witness, \ev{Lean} for the implication. \scale{bounded}.
\[
  \forall a_0\le 6,\ \exists a,L,\ a_0\le a \wedge L < 2\cdot 2^a \wedge
  \mathcal C(H_{2^a},H_{2^a},L)
\]
The single witness $(a,L)=(6,93)$ works for every threshold $a_0\le6$.
Its certificate is verified by exact integer arithmetic on the finite
totient windows in
\lean{certifiedKill_diagonal_t64}{DiagonalPincerCertificateT64Endpoint.lean:1928}.
This proves the displayed bounded statement. It does not construct
witnesses for unbounded thresholds; further finite checks would extend
the verified range, not establish the universally quantified condition below.
\end{thm}
\leannote{Lean: \lproof{ErdosProblems/Erdos249/PaperCompleteR21/ShortWindowSupplyAndSixteenShifts.lean}{29}{shortWindowSupply_through_six_paper}, \lproof{ErdosProblems/Erdos249/PaperCompleteR21/ShortWindowSupplyAndSixteenShifts.lean}{38}{shortWindowSupply_single_witness_six_ninetyThree}, \lproof{ErdosProblems/Erdos249/PaperCompleteR21/ShortWindowSupplyAndSixteenShifts.lean}{47}{shortWindowSupply_witness_eq_t64_certificate}.}

\begin{defn}[The additional hypothesis]
$\text{cofinal short-window condition}$: \ev{Open}, \scale{cofinal}.
\[
  \forall a_0\ \exists a,L,\ a_0\le a \wedge L < 2\cdot 2^a \wedge
  \mathcal C(H_{2^a},H_{2^a},L) .
\]
The depth restriction gives a useful description of the offsets:
if $1\le j<2\cdot2^a$ and $j\nmid H_{2^a}$, then $j$ is a prime
power greater than $2^a$. Indeed, some prime-power divisor $p^r$ of
$j$ exceeds $2^a$; if $j/p^r\ge2$, then $j>2\cdot2^a$, a contradiction.
This describes the possible offsets, not the size of their weighted sum.

The condition restricts the depth in the ordinary diagonal certificate
criterion. It is not established by pointwise completeness, which permits
arbitrarily large depths. Nor does a successful short-window certificate
by itself yield the fixed $1/32$ separation required by the preceding
approximation condition. These are distinct sufficient hypotheses;
no implication from this condition to that fixed-margin condition is
proved here.
\end{defn}

\subsubsection{The diagonal certificate table}

\begin{thm}
\lean{certifiedKill_diagonal_all_imported}{DiagonalPincerCertificates.lean:2870}, with
\lean{diagonalPincerCertificateScales}{DiagonalPincerCertificates.lean:2852} and
\lean{diagonalPincerKillDepth}{DiagonalPincerCertificates.lean:2854}. \ev{Cert}. \scale{fixed}.
\[
  \forall t\in\{1,2,3,4,5,7,8,9,11,13,16,17\},\ \exists L,\
  \mathcal{C}\ ({H}\ t)\ ({H}\ t)\ L
\]
at depths $\{6,5,7,7,9,14,15,14,21,22,23,26\}$ respectively, extended by the separate T19\ldots T64
modules to 28 historical values through $t=64$ (the $t=64$ endpoint is
\lean{certifiedKill_diagonal_t64}{DiagonalPincerCertificateT64Endpoint.lean:1928}).
The later aggregate theorem closes every scale $t\le82$ with no holes
(\lean{ErdosProblems.Skip.LadderT67.exists_diagonalKill_le_82}{ErdosProblems/Skip/LadderT67.lean:71264}).
Each witness is checked by exact integer arithmetic on two finite totient
windows. For example, $H_{17}=12252240$ and $L=26$ use
$\varphi(H_{17}+i)$ and $\varphi(2H_{17}+i)$ for $1\le i\le26$,
not every totient between the endpoints $12252241$ and $24504506$.
These computations verify the listed cases; they do not give a formula
producing a certificate for every $t$.
\end{thm}
\leannote{Lean: \lproof{ErdosProblems/Erdos249/PaperCompleteR21/DiagonalCertificateTableScales.lean}{22}{diagonalPincerCertificateScales_list}, \lproof{ErdosProblems/Erdos249/PaperCompleteR21/DiagonalCertificateTableScales.lean}{26}{diagonalPincerKillDepth_list}, \lproof{ErdosProblems/Erdos249/PaperCompleteR21/DiagonalCertificateTableScales.lean}{33}{certifiedKill_diagonal_table}, \lproof{ErdosProblems/Erdos249/PaperCompleteR21/DiagonalCertificateTableScales.lean}{40}{exists_diagonalKill_on_table}, and 3 further declarations in the \hyperref[sec:coverage]{coverage section}.}

\begin{rem}[The observed depths are finite data]
\lean{certifiedKill_depth_floor}{TotientTailPeriodKiller.lean:79} forces $2(2H_t{+}L{+}2)<2^L$,
i.e.\ $L \gtrsim \log_2(4\cdot{H}\ t)$ at any certified depth. Comparing that floor
against $\mathtt{diagonalPincerKillDepth}$ shows every proved certificate fires within a small
additive constant of the theoretical minimum depth, at every tested scale; a numerically
supported (\ev{Cert}, not \ev{Math}) separation observation, not a theorem.
\end{rem}

\begin{prop}[A sufficient extension]
\[
  \exists C\ \forall t_0\ \exists t\ge t_0\ \exists L\le \log_2(4\cdot{H}\ t)+C,\quad
  \mathcal{C}\ ({H}\ t)\ ({H}\ t)\ L .
\]
This asks for a certificate within a fixed additive constant of the
necessary logarithmic depth bound. The uniform constant is an additional
requirement, not a consequence of the finite data.
\end{prop}
\leannote{Lean: \lproof{ErdosProblems/Erdos249/PaperCompleteR21/ShortWindowSupplyAndSixteenShifts.lean}{58}{irrational_of_logarithmicDepth_diagonal_supply}, \lproof{ErdosProblems/Erdos249/PaperCompleteR21/ShortWindowSupplyAndSixteenShifts.lean}{72}{irrational_of_restrictedDepth_diagonal_supply}.}

\subsubsection{One common certificate for sixteen shifts}

\begin{thm}
\lean{certifiedKill_all_upto_sixteen}{CarrySurvivorExtinction.lean:574} and
\lean{totient_series_ne_rat_of_den_dvd_pow_two_mul_mersenne_upto_sixteen}{CertificateKernel.lean:18890}
(the earlier eight-shift computation is \lean{certifiedKill_all_small}{TotientTailPeriodKiller.lean:404}). \ev{Cert}.
\scale{fixed}.
\[
  \mathcal C(h,14,9)\qquad(h\in\{1,\ldots,16\}).
\]
The basepoint $14$ and depth $9$ are common to all sixteen shifts.
Consequently, $S\ne r$ whenever $r\in\mathbb Q$ and
$\operatorname{den}(r)\mid2^{14}(2^h-1)$ for some integer $h$ with
$1\le h\le16$. The exact computation uses two nine-term windows for
each shift and reduces their discrepancy modulo $2^9$. The earlier
eight-shift computation uses basepoint $12$ and depth $16$.
\end{thm}
\leannote{Lean: \lproof{ErdosProblems/Erdos249/PaperCompleteR21/ShortWindowSupplyAndSixteenShifts.lean}{85}{commonCertificate_sixteen_shifts_basepoint_fourteen}, \lproof{ErdosProblems/Erdos249/PaperCompleteR21/ShortWindowSupplyAndSixteenShifts.lean}{91}{totientSeries_ne_rat_of_den_dvd_two_pow_fourteen_mul_mersenne}, \lproof{ErdosProblems/Erdos249/PaperCompleteR21/ShortWindowSupplyAndSixteenShifts.lean}{98}{commonCertificate_eight_shifts_basepoint_twelve}.}

\begin{rem}[Finite simultaneity and the unbounded condition]
The displayed computation concerns only $N=14$, $L=9$ and
$1\le h\le16$. It does not provide certificates beyond every basepoint
threshold. Simultaneous certification of finitely many shifts at a
sufficiently large, unrestricted depth is not an additional obstacle:
under irrationality it follows from the same truncation estimate, as
shown next. The argument below does not prescribe the common depth in terms of
the basepoint and number of shifts.
\end{rem}

\begin{prop}[Simultaneous certificates with unrestricted depth]
The existence of a function $f:\N\to\N$ with $f(N)\to\infty$ such that
\[
 \forall N_0\ \exists N\ge N_0\ \exists L\ 
 \forall h\in\{1,\ldots,f(N)\},\quad \mathcal C(h,N,L)
\]
is equivalent to irrationality of $S$. Under irrationality, one can
take $f(N)=N+1$, and every $N$ has a suitable depth. \ev{Math}
\end{prop}
\leannote{Lean: \lproof{ErdosProblems/Erdos249/PaperCompleteR21/SimultaneousShiftCertificateDepth.lean}{131}{exists_growingShift_simultaneous_certificate_iff_irrational}, \lproof{ErdosProblems/Erdos249/PaperCompleteR21/SimultaneousShiftCertificateDepth.lean}{122}{exists_simultaneous_depth_succ_of_irrational}, \lproof{ErdosProblems/Erdos249/PaperCompleteR21/SimultaneousShiftCertificateDepth.lean}{89}{exists_simultaneous_depth_of_irrational}.}
\begin{proof}
The displayed condition supplies Proposition~\ref{prop:A10}'s
quantified certificate condition: for fixed $h\ge1$, choose the
basepoint threshold so large that $f(N)\ge h$ beyond it. That
proposition gives irrationality.

Conversely, suppose $S$ is irrational. For every $h\ge1$ and $N\ge0$,
\[
 R_{N+h}-R_N
 =2^N(2^h-1)S-(\Phi_{N+h}-\Phi_N)
\]
is nonintegral. Fix $N$ and a positive integer $M$. Choose $L$ so large
that
\[
 \frac{2(N+M+L+2)}{2^L}
 <\min_{1\le h\le M}
   \operatorname{dist}(R_{N+h}-R_N,\Z).
\]
The finite minimum is positive, whereas the left side tends to zero.
For each $h\le M$, the truncation error between $D(h,N,L)/2^L$ and
$R_{N+h}-R_N$ is at most $(N+h+L+2)/2^L$. The triangle inequality
therefore puts $D(h,N,L)/2^L$ at distance greater than
$(N+h+L+2)/2^L$ from every integer, which is precisely
$\mathcal C(h,N,L)$. Taking $M=N+1$ proves the converse.
\end{proof}
This proof gives no quantitative lower bound for the finite minimum
as the basepoint or number of shifts grows. It selects a common
depth after those parameters are fixed; it does not impose depth equal
to the period.

\subsubsection{The fixed Farey bound}

\begin{thm}
\lean{tsum_totient_div_pow_two_ne_ratCast_of_den_le_79639646646701375323355774875831053}{CertificateKernel.lean:18384},
built from the classical mediant lemma \lean{farey_gap}{GapFareyBound.lean:51} and the window
sharpness certificate \lean{gap_check_window_1_240_le_79639646646701375323355774875831053}{GapFareyBound.lean:176}
/ \lean{gap_check_window_1_240_first_failure}{GapFareyBound.lean:225}. \ev{Lean} for the general
mediant lemma; \ev{Cert} for the $K=240$ window instantiation. \scale{fixed}. If $S$ is rational,
its reduced denominator exceeds $7.9639646646701375323355774875831053\times 10^{34}$ (the exact
numeral in the declaration name). This is the finite denominator exclusion in the
corpus, logically independent of the certificate-supply reduction.
\end{thm}
\leannote{Lean: \lproof{ErdosProblems/Erdos249/PaperCompleteR21/FareyDenominatorFloorExtension.lean}{23}{totientSeries_rational_den_gt_fareyBound}, \lproof{ErdosProblems/Erdos249/PaperCompleteR21/FareyDenominatorFloorExtension.lean}{31}{farey_gap_paper}, \lproof{ErdosProblems/Erdos249/PaperCompleteR21/FareyDenominatorFloorExtension.lean}{39}{gapCheck_window_1_240_paper}, \lproof{ErdosProblems/Erdos249/PaperCompleteR21/FareyDenominatorFloorExtension.lean}{48}{gapCheck_window_1_240_first_failure_paper}.}

\begin{rem}[What a larger Farey bound requires]
Unbounded valid denominator exclusions would imply $S\notin\Q$.
For the single window $K=240$, however,
\lean{gap_check_window_1_240_first_failure}{GapFareyBound.lean:225}
proves that the residue inequality fails at the mediant denominator
$b+d=Q_0+1$. This exact first failure prevents a larger cutoff from
that same inequality at that same window. It is not an obstruction to
other windows or to other tests using the finite data.

The existing proofs for individual windows use explicit totient residues
and rational endpoint calculations. Extending that implementation calls
for new finite data; a uniform proof of a growing bound need not use
hard-coded convergents. If $S$ had reduced denominator $q$, every valid
exclusion bound would be smaller than $q$. Thus unbounded exclusions suffice for irrationality; rationality
would bound them, but would not force eventual constancy.
No converse or comparison with the
continued-fraction denominators of $S$ is proved here.
\end{rem}

\begin{obs}[The limited gain from a unit-gap refinement]
The finite condition
\lean{ReducedDenominatorUnitGapCert}{PrimitiveWeightCertificate.lean:30}
and its quantified version
\lean{PrimitiveReducedDenominatorUnitGapSupply}{PrimitiveWeightCertificate.lean:47}
are definitions; the latter is not asserted to hold. The separate results
\lean{reducedDenominatorUnitGapCert_prime_pow_iff}{PrimitiveWeightCertificate.lean:54} and
\lean{reducedDenominatorCandidateCount_prime_pow_le_one}{PrimitiveWeightCertificate.lean:111}
(\ev{Lean}, \scale{uniform}) show that, for $p$ prime and $e\ge1$,
a window at denominator $p^e$ that \emph{passes the refined test}
contains at most one candidate integer. A window with two or more
candidates fails the test. This is a condition for successful
certification, not a bound on all windows or on the improvement
in the denominator cutoff. It does not apply to all odd denominators:
Observation~\ref{prop:D8cert-kill} gives an actual totient window with
two admissible nonunit candidates modulo $15$.
\end{obs}


\begin{prop}[A sufficient extension]
Suppose $g(K)\to\infty$ and, for every $K$, the $(N=1,K)$ gap check
excludes every rational of reduced denominator at most $g(K)$.
Then $S$ is irrational: any rational value of $S$ would have a fixed
finite denominator, contradicted at a sufficiently large $K$.
Only unbounded exclusion bounds along a sequence of windows are
needed for this argument. No equivalence with a bound for the
continued-fraction denominators of $S$ is asserted; such a comparison
would require its own proof connecting those convergents to the
certified intervals.
\end{prop}
\leannote{Lean: \lproof{ErdosProblems/Erdos249/PaperCompleteR21/FareyDenominatorFloorExtension.lean}{68}{irrational_of_unbounded_window_one_gapCheck}, \lproof{ErdosProblems/Erdos249/PaperCompleteR21/FareyDenominatorFloorExtension.lean}{58}{gapCheck_window_one_excludes}.}

\subsubsection{A two-adic congruence that does not give a certificate}

\begin{thm}
\lean{exists_prime_totient_twoAdic_pulse_divisors}{TotientTwoAdicPulseBlock.lean:54},
\lean{exists_prime_deltaTotient_twoAdic_pulseBlock}{TotientTwoAdicPulseBlock.lean:211},
\lean{windowDiscrepancy_modEq_half_of_twoAdic_pulse}{TotientTwoAdicPulseBlock.lean:262},
\lean{eventual_integral_tailDiff_has_cofinal_twoAdic_half_pulse}{TotientTwoAdicPulseBlock.lean:327}.
\ev{Lean}. \scale{uniform} (the theorem is unconditional and holds for \emph{every} $K$, $H$, $B$,
not merely cofinally many). For every $K\ge 2$, every $H>K$, and every bound $B$, there are primes
$p>B$ with a length-$(K{-}1)$ zero prefix and a terminal half-modulus, giving
\[
  D(H,p-K,K) \equiv 2^{K-1} \pmod{2^K}.
\]
Under eventual integrality this transfers to an integer $z$ with $(z:\mathbb R) =
{R}(p{+}H) - R_{p}$ and $z\equiv 2^{K-1}\pmod{2^K}$.
\end{thm}
\leannote{Lean: \lproof{ErdosProblems/Erdos249/PaperCompleteR21/TwoAdicHalfPulseAndAccumulatedResidue.lean}{26}{exists_prime_twoAdic_half_pulse_window}, \lproof{ErdosProblems/Erdos249/PaperCompleteR21/TwoAdicHalfPulseAndAccumulatedResidue.lean}{41}{eventual_integral_tailDiff_twoAdic_half_pulse}.}

\begin{rem}[Why the central residue still fails the test]
This lands the residue at the exact centre of the residue-certificate arc: $2^{K-1}$ is maximally far
from $0$ modulo $2^K$, precisely the coordinate the obligation wants. Taking $N{=}p{-}K$,
$h{=}H$, $L{=}K$ gives radius $N{+}h{+}L{+}2 = p{+}H{+}2$, so $\mathcal{C}$ requires
$2^{K-1} > p{+}H{+}2$. But $z\equiv 2^{K-1}\pmod{2^K}$ forces $|z|\ge 2^{K-1}$, while the directed
tail bound forces $|z| < p{+}H{+}2$; and the construction's own congruence
$p\equiv 1{+}2^{K-1}\pmod{2^K}$ (equivalently $v_2(p{-}1)=K{-}1$) forces $p\ge 1{+}2^{K-1}$. So
$p{+}H{+}2 > 2^{K-1}$ always, and no choice of $p$ rescues it. The quantifier order is inverted:
the theorem gives, for fixed $K$, cofinally many \emph{large} $p$; the certificate needs a $p$
\emph{small} relative to $2^{K-1}$, impossible for a single letter since $|\varphi(N{+}H)-\varphi(N)|
< N{+}H$. Depth-promotion is not the issue: it has already been performed, generalizing an
exponent-2 construction to every $K$ uniformly by CRT-gluing over a
$\mathrm{Unit}\oplus\mathrm{Fin}(K{-}1)\oplus\mathrm{Fin}(K{-}1)$ family.
\end{rem}

\begin{prop}[A sufficient accumulated-residue condition]
Suppose that for every integer $h\ge1$ and every $N_0\in\mathbb N$
there are $N\ge N_0$ and $L\ge1$ such that
\[
 D(h,N,L)\equiv2^{L-1}\pmod{2^L},\qquad
 N+h+L+2<2^{L-1}.
\]
Then $S\notin\mathbb Q$: these two inequalities imply
$\mathcal C(h,N,L)$, and the full quantified certificate condition
applies. The half-modulus condition here is imposed on the accumulated
weighted sum $D(h,N,L)$, not on a single totient difference.
An unbounded set of triples with no control of the shift does not
supply the displayed quantifiers. The required family remains unproved.
\end{prop}
\leannote{Lean: \lproof{ErdosProblems/Erdos249/PaperCompleteR21/TwoAdicHalfPulseAndAccumulatedResidue.lean}{77}{irrational_of_accumulated_halfModulus_supply}, \lproof{ErdosProblems/Erdos249/PaperCompleteR21/TwoAdicHalfPulseAndAccumulatedResidue.lean}{53}{certifiedKill_of_halfModulus_residue}.}

\subsubsection{The residue margin at an LCM jump}

\begin{thm}[Arbitrarily large prime-power LCM jumps]
\lean{exists_periodLcm_strict_jump_ge}{DiagonalFreshLossBridge.lean:2635}. \ev{Lean}.
\scale{cofinal} (proved): for every $t_0\in\mathbb N$ there is
$t\ge t_0$ with $H(t)<H(t+1)$. One may take $t=p-1$ for any
prime $p>t_0$. The positions $t=2^a-1$ are also strict LCM jumps for $a\ge1$:
$2^a$ is the next required power of $2$. The restriction excludes
$a=0$, since $H(0)=H(1)=1$.
The condition
\lean{CanonicalAdjacentSuffixPowerTwoPostJumpSlackSupply}{DiagonalFreshLossBridge.lean:2693}
asks for the additional residue margin along this explicit sequence.
Knowing these jump positions does not establish that margin.
\end{thm}
\leannote{Lean: \lproof{ErdosProblems/Erdos249/PaperCompleteR21/LcmJumpPositionsAndCentralSlack.lean}{26}{exists_periodLcm_strict_jump_ge_paper}, \lproof{ErdosProblems/Erdos249/PaperCompleteR21/LcmJumpPositionsAndCentralSlack.lean}{31}{periodLcm_strict_jump_at_prime_pred}, \lproof{ErdosProblems/Erdos249/PaperCompleteR21/LcmJumpPositionsAndCentralSlack.lean}{46}{periodLcm_strict_jump_at_powerTwo_pred}, \lproof{ErdosProblems/Erdos249/PaperCompleteR21/LcmJumpPositionsAndCentralSlack.lean}{59}{periodLcm_zero_and_one_eq_one}.}

\begin{defn}[The remaining residue margin]
For $t\ge3$, put $H=H_t$ and define
\[
 \begin{aligned}
 m&=\lfloor\log_2 H\rfloor+10,\\
 d&=\bigl(D(H,H+1,m)-D(H,H,m)\bigr)\bmod2^m,\\
 \sigma_t&=\min\{d-2^{m-5},\,2^m-2^{m-5}-d\}.
 \end{aligned}
\]
Thus $d$ is the difference of two adjacent window numerators, reduced
modulo $2^m$, not the residue of $D(H,H,m)$ alone. The integer
$\sigma_t$ is nonnegative exactly when
$2^{m-5}\le d\le2^m-2^{m-5}$. The choice of $m$ makes the error
bound $3H+2m+5$ smaller than $2^{m-5}$; nonnegativity is the
additional arithmetic requirement, not a consequence of that choice.
\lean{canonicalAdjacentSuffixCentralSlack}{DiagonalFreshLossBridge.lean:2651}
\lean{canonicalAdjacentSuffixDepth}{DiagonalFreshLossBridge.lean:1870}
\lean{diagonalAdjacentSuffixResidue}{DiagonalFreshLossBridge.lean:525}

The first unproved condition asks for every $t_0$ for an index
$t\ge\max(3,t_0)$ with $H_t<H_{t+1}$ and $\sigma_{t+1}\ge0$.
The second restricts these endpoints to powers of two: for every $a_0$,
there must be $a\ge\max(2,a_0)$ with $\sigma_{2^a}\ge0$.
These are respectively
\lean{CanonicalAdjacentSuffixPostJumpSlackSupply}{DiagonalFreshLossBridge.lean:2686}
and
\lean{CanonicalAdjacentSuffixPowerTwoPostJumpSlackSupply}{DiagonalFreshLossBridge.lean:2693}.
Each condition implies $S\notin\mathbb Q$, by
\lean{irrational_totientSeries_of_canonicalAdjacentSuffixPostJumpSlackSupply}{DiagonalFreshLossBridge.lean:2941}
and
\lean{irrational_totientSeries_of_canonicalAdjacentSuffixPowerTwoPostJumpSlackSupply}{DiagonalFreshLossBridge.lean:2948}.
\ev{Open} \scale{cofinal} (hypotheses; the implications are proved).
\end{defn}

\begin{rem}[Finite numerical evidence]
\ev{Cert}, \scale{fixed}. All $40$ strict jumps up to endpoint $113$ pass, closest margin
$\approx 0.000221$ of the modulus at $t=100$; the power-of-two endpoints $4,8,16,32$ pass with
slack fractions $0.41,0.036,0.40,0.19$. This is a finite computation, not a cofinal theorem. The identity
$\mathtt{periodLcm\_succ\_eq\_prime\_mul\_of\_strict\_jump}$
(\lean{periodLcm_succ_eq_prime_mul_of_strict_jump}{DiagonalFreshLossBridge.lean:2332}) says a
strict jump has $t+1=p^k$ for a prime $p$ and $k\ge1$, with
$H_{t+1}=pH_t$. Thus $\sigma_{t+1}$ uses the new height $pH_t$,
not the pre-jump height $H_t$. Estimating its change would require
information about the adjacent-window residue, not just the LCM ratio.
\lean{canonicalAdjacentSuffixCentralSlack_eq_of_periodLcm_eq}{DiagonalFreshLossBridge.lean:2658}
proves $\sigma_t=\sigma_u$ whenever $H_t=H_u$. Thus the value changes
only when the LCM height changes. This reduces the number of distinct
values to study; it does not determine their signs.
\end{rem}

\begin{prop}[A sufficient inequality]
If, for every $a_0\in\mathbb N$, there is $a\ge\max(2,a_0)$ with
$\sigma_{2^a}\ge0$, then $S\notin\mathbb Q$.
This is the sufficient condition defined above, not a claim that the
inequality holds at arbitrarily large indices. Equality at the edge
is allowed: the condition is nonnegativity, not strict positivity.
\end{prop}
\leannote{Lean: \lproof{ErdosProblems/Erdos249/PaperCompleteR21/LcmJumpPositionsAndCentralSlack.lean}{90}{irrational_of_powerTwo_postJump_slack_supply}, \lproof{ErdosProblems/Erdos249/PaperCompleteR21/LcmJumpPositionsAndCentralSlack.lean}{68}{canonicalAdjacentSuffixCentralSlack_paper_formula}.}

\subsubsection{A certificate for a four-tail combination}

\begin{thm}
\lean{primeJumpSharpRadius}{PrimeJumpWindow.lean:37} and
\lean{primeJumpTailCommutator_notMem_int_of_sharpKill}{PrimeJumpWindow.lean:178}. \ev{Lean}.
\scale{uniform}. Use $J(H,p)$, $W(H,p,L)$ and $B(H,p,L)$ from
Theorem~\ref{catalogue:mob:e2}. For all $H,p,L\in\mathbb N$,
\[
 B(H,p,L)<W(H,p,L)\bmod2^L<2^L-B(H,p,L)
 \quad\Longrightarrow\quad J(H,p)\notin\mathbb Z.
\]
The bound $B(H,p,L)=3pH+(p+1)(L+2)$ follows by grouping the
nonnegative remainders as in that theorem. Compared with the
$4pH+(p+1)(L+2)$ bound obtained by treating the two diagonal
differences separately, it saves $pH$. This is a strict improvement
when $pH>0$; optimality for totient tails is not asserted.
\end{thm}
\leannote{Lean: \lproof{ErdosProblems/Erdos249/PaperCompleteR21/PrimeJumpWitnessAndMersenneChannels.lean}{36}{primeJumpTailCommutator_notMem_int_of_central_window}, \lproof{ErdosProblems/Erdos249/PaperCompleteR21/PrimeJumpWitnessAndMersenneChannels.lean}{31}{primeJumpSharpRadius_formula}, \lproof{ErdosProblems/Erdos249/PaperCompleteR21/PrimeJumpWitnessAndMersenneChannels.lean}{45}{primeJumpSharpRadius_saves_pH}, \lproof{ErdosProblems/Erdos249/PaperCompleteR21/PrimeJumpWitnessAndMersenneChannels.lean}{52}{primeJumpSharpRadius_lt_twoCellRadius}.}

\begin{thm}[One explicit witness]
\lean{primeJumpSharpKill_twelve_five}{PrimeJumpWindow.lean:186}. \ev{Cert}. \scale{fixed}.
For $H=H(4)=12$, $p=5$ and $L=15$, direct integer evaluation gives
\[
 W(12,5,15)=149906,\qquad
 W(12,5,15)\bmod32768=18834,\qquad B(12,5,15)=282.
\]
Thus $282<18834<32486$, and $J(12,5)\notin\mathbb Z$.
This is the one instance singled out by the cited declaration, not a
claim of uniqueness or of the smallest possible height. The implication
from a quantified family is
\lean{irrational_totient_series_of_primeJumpSharpKill_supply}{PrimeJumpWindow.lean:193}.
\end{thm}
\leannote{Lean: \lproof{ErdosProblems/Erdos249/PaperCompleteR21/PrimeJumpWitnessAndMersenneChannels.lean}{62}{periodLcm_four_eq_twelve}, \lproof{ErdosProblems/Erdos249/PaperCompleteR21/PrimeJumpWitnessAndMersenneChannels.lean}{68}{primeJump_witness_twelve_five_values}, \lproof{ErdosProblems/Erdos249/PaperCompleteR21/PrimeJumpWitnessAndMersenneChannels.lean}{76}{primeJumpTailCommutator_twelve_five_notMem_int}, \lproof{ErdosProblems/Erdos249/PaperCompleteR21/PrimeJumpWitnessAndMersenneChannels.lean}{81}{irrational_of_primeJumpSharp_supply}.}

\begin{rem}[Scope of the witness]
The quantified hypothesis asks, for every $t_0\in\mathbb N$, for
$t\ge t_0$, $p\ge1$ and $L\ge0$ satisfying the displayed residue
condition with $H=H(t)$. The multiplier $p$ need not be prime, and need
not be tied to the next LCM jump. No comparison of the difficulty of
this condition with the prescribed-depth separation condition is proved.

At a strict LCM jump, $H(t+1)=pH(t)$ for a prime $p$. This gives one
natural choice of multiplier, but the prime need not be a new divisor:
$H(4)=2H(3)$ and $2$ already divides $H(3)=6$.
The identity in
\lean{oldChannel_affine_moment_annihilation}{JointExponentTransport.lean:36}
cancels contributions affine in the multiplier. It does not bound the
remaining terms of the four-tail combination. The finite example above
therefore does not establish the quantified condition.
\end{rem}

\subsubsection{Denominators of M\"obius sums}

\begin{thm}
\lean{upperHalfMersenneProduct_lower_bound}{MersenneShadowDenominatorGrowth.lean:60},
\lean{lcmHeight_scaledMobiusShadow_den_lower_bound}{MersenneShadowDenominatorGrowth.lean:85},
\lean{exists_upperHalf_mersenne_channel_dvd_den_with_bounds}{MersenneShadowDenominatorGrowth.lean:99}.
\ev{Lean}. \scale{uniform}. For every integer $t\ge5$, with
$\mathcal P_t=\{p\text{ prime}:t/2<p\le t\}$ as above,
\[
 2^{\lfloor t/2\rfloor}
 \le\prod_{p\in\mathcal P_t}(2^p-1)
 \le\operatorname{den}(H_t\beta_{H_t}).
\]
Moreover, some $p\in\mathcal P_t$ satisfies
\[
 2^{\lfloor t/2\rfloor}\le2^p-1<2^t,
 \qquad 2^p-1\mid\operatorname{den}(H_t\beta_{H_t}).
\]
Indeed, Bertrand's postulate makes $\mathcal P_t$ nonempty. Its members
satisfy $p-1\ge\lfloor t/2\rfloor$ and $p\le t$, giving the size
bounds. The divisibility follows from
Proposition~\ref{catalogue:mob:b6}; divisibility of a positive reduced
denominator also gives the product inequality.
\end{thm}
\leannote{Lean: \lproof{ErdosProblems/Erdos249/PaperCompleteR21/PrimeJumpWitnessAndMersenneChannels.lean}{110}{upperHalfMersenneProduct_between_bounds}, \lproof{ErdosProblems/Erdos249/PaperCompleteR21/PrimeJumpWitnessAndMersenneChannels.lean}{123}{exists_upperHalf_channel_paper}, \lproof{ErdosProblems/Erdos249/PaperCompleteR21/PrimeJumpWitnessAndMersenneChannels.lean}{98}{upperHalfPrimes_nonempty_paper}, \lproof{ErdosProblems/Erdos249/PaperCompleteR21/PrimeJumpWitnessAndMersenneChannels.lean}{103}{upperHalfPrimes_member_bounds}, and 2 further declarations in the \hyperref[sec:coverage]{coverage section}.}

\begin{rem}[Growth of these denominators does not give approximation of $S$]
The number $H_t\beta_{H_t}$ is the finite rational contribution to a
totient-tail decomposition. Its complementary term can still cancel
its fractional part. The stated denominator bound supplies neither an
error estimate for approximation of $S$ nor nonintegrality of
$R_{2H_t}-R_{H_t}$.

To apply the diagonal criterion in Theorem~\ref{catalogue:cert:b3},
one must establish this tail nonintegrality at arbitrarily large LCM
scales. Denominator growth alone does not supply that condition.
\end{rem}

\begin{prop}[The additional approximation hypothesis]
A sequence of rationals $u_t$ proves irrationality of $S$ if
\[
 u_t\ne S\quad\text{for all sufficiently large }t,
 \qquad \operatorname{den}(u_t)|S-u_t|\longrightarrow0.
\]
This is the rational-separation criterion in
\lean{irrational_of_den_mul_abs_sub_tendsto_zero}{CertificateKernel.lean:5371}:
if $S=a/b$ were reduced, every unequal $u_t$ would instead satisfy
$\operatorname{den}(u_t)|S-u_t|\ge1/b$. An upper bound tending to zero
for this product is needed. A lower bound for the denominator cannot
establish it; improving a lower bound does not change the actual
approximation error.

A different argument could deduce a residue certificate from the
surviving divisor and further information about the complementary
term. The denominator theorem alone contains no such separation
statement. Neither sufficient argument is established by its lower
bound, and these are not claimed to exhaust possible approaches.
\end{prop}
\leannote{Lean: \lproof{ErdosProblems/Erdos249/PaperCompleteR21/PrimeJumpWitnessAndMersenneChannels.lean}{137}{irrational_totientSeries_of_rational_separation}, \lproof{ErdosProblems/Erdos249/PaperCompleteR21/PrimeJumpWitnessAndMersenneChannels.lean}{149}{den_mul_abs_sub_ge_one_div_den}.}

\subsubsection{Rationality and carry rank}

The next two theorems concern the totient sections and, under rationality,
the sections of an integral carry. Both give lower bounds for dimension.
A rank argument for irrationality would also need an incompatible upper
bound. Rationality alone does not provide one.

\begin{thm}
\lean{linearIndependent_canonicalTotientKernelFamily}{TotientMahlerDefect.lean:935}. \ev{Lean}.
\scale{uniform}. For every $e\ge1$, the retained dyadic totient sections form a basis for
all sections through level $e$, of dimension $2^e+1$. The proof uses the
Chinese remainder theorem and Dirichlet's theorem. The complete level-zero
truncation, consisting only of $\varphi$, has dimension one.
\end{thm}
\leannote{Lean: \lproof{ErdosProblems/Erdos249/PaperCompleteR21/DyadicSectionBasisAndRationalCarry.lean}{79}{retainedSections_basis_and_rank}, \lproof{ErdosProblems/Erdos249/PaperCompleteR21/DyadicSectionBasisAndRationalCarry.lean}{93}{completeLevelZeroTruncation_is_totient_and_rank_one}, \lproof{ErdosProblems/Erdos249/PaperCompleteR21/DyadicSectionBasisAndRationalCarry.lean}{33}{canonicalTotientKernelFamily_entries}.}

\begin{thm}
\lean{not_irrational_totientSeries_implies_unbounded_carryRank_unconditional}{TotientCarryKernelRigidity.lean:300}.
\ev{Lean}. \scale{uniform}. If $S\in\Q$, choose an integer $v>0$
with $vS\in\Z$ and put $u_N=vR_N$. Then $u_N\in\Z$,
$u_{N+1}=2u_N-v\varphi(N+1)$, and $0\le u_N\le v(N+2)$. For every $e\ge0$,
\[
 \dim_{\Q}\operatorname{span}
 \{n\mapsto u_{2^j n+r}:1\le j\le e,\ 0\le r<2^j\}\ge2^e-1.
\]
The proof is given in Section~\ref{sec:mahler-defect}.
\end{thm}
\leannote{Lean: \lproof{ErdosProblems/Erdos249/PaperCompleteR21/DyadicSectionBasisAndRationalCarry.lean}{129}{rationalValue_integral_carry_and_rank_floor}.}

\begin{rem}[The generic upper-bound inference is false]
The bound $2^e-1$ holds at every level. An upper bound contradicting it
would need additional information about the totient, since the rational
$5/4$ comparison has the same lower bound. One proposed system of integer
equalities and strict inequalities is inconsistent:
\lean{false_of_compressedAdjointCertificate}{TotientMahlerDefect.lean:1183} (\ev{Lean},
\scale{n/a}) shows no
\lean{CompressedAdjointCertificate}{TotientMahlerDefect.lean:1171} with $Q,v>0$, $A\ne0$, $QvA=\mathrm{boundary}$ and
$|\mathrm{boundary}|<Qv$ can exist: a nonzero integer multiple of $Qv$
has absolute value at least $Qv$. This excludes those simultaneous
requirements, not every argument using a matrix or its adjugate.
Separately, infinite-dimensionality of the totient kernel span is the
conclusion of the source theorem immediately preceding it.
\end{rem}

\begin{prop}[An additional hypothesis for a totient-specific rank argument]
A rank bound that contradicts the lower bound $2^e-1$ would have to use
additional arithmetic of the actual totient coefficients. The generic
proposal that every rational coefficient series has bounded tempered-carry
rank is ruled out by the rational control recorded in the short paper:
its carry rank is at least $2^e-1$ at every level. It is therefore not a
remaining general lemma from which totient irrationality follows.
The conditional theorem in the totient tail carry period source, line~224,
records eventual periodicity modulo $v$ together with unbounded rational
rank. Periodicity of the residue sequences concerns values in a finite
quotient; it does not bound the rational span of the integer-valued
sections. The $5/4$ comparison shows that this distinction persists for
a rational coefficient series.
\end{prop}
\leannote{Lean: \lproof{ErdosProblems/Erdos249/PaperCompleteR21/DyadicSectionBasisAndRationalCarry.lean}{238}{rationalControl_periodic_with_unbounded_carry_rank}, \lproof{ErdosProblems/Erdos249/PaperCompleteR21/DyadicSectionBasisAndRationalCarry.lean}{256}{rationality_gives_mod_period_and_unbounded_rank}.}

\subsection{What the preceding comparisons show}

The comparison separates three issues. A block estimate with a fixed
saving in the real part or norm of an exponential sum would provide
certificates. The upper-endpoint conditions instead ask for a particular
residue or size estimate along an LCM sequence. Their sufficient forms are
related by Proposition~\ref{prop:te-chain}, but those implications do not
prove any of the arithmetic hypotheses.

The generic bounded-rank argument fails for a separate reason: the rational
comparison sequence has unbounded section rank. A totient-specific upper
bound would require hypotheses not shared by that sequence.

These are statements about the displayed conditions and comparisons, not
an exhaustive classification of possible proofs. Some conditions are
equivalent to irrationality; others are only sufficient. No relative
difficulty, truth probability, or independent formal verification follows
from their position in the table.

\section{Counterexamples to proposed deductions}

The following examples disprove particular deductions proposed for
$S=\sum_{n\ge1}\varphi(n)2^{-n}$. Each example specifies the assumptions
that it satisfies and the conclusion that fails. Some concern other
coefficient sequences; others concern a particular way of estimating a
totient sum. Neither kind decides the irrationality of $S$, and neither
excludes arguments using additional properties of $\varphi$.


\subsection{Changing the coefficients after a finite prefix}

\textbf{(a) Proposed argument.} The quantified residue condition requires
certificates beyond every starting index. A finite computation can check
many instances, but one fixed inspection bound cannot establish these
quantifiers. The following comparison explains this limitation without
making any claim about all possible irrationality arguments.
\lean{Erdos249257.irrational\_totient\_series\_iff\_certificate\_supply}{LcmConeFlatness.lean:412}


\textbf{(b) Counterexample.} Theorem~\ref{thm:gamma} gives an explicit splice
construction; its series value, exact reduced denominator and failure of the
separation condition are checked in Lean (\ev{Lean}). Choose an integer \(P>B\), set
\(\gamma(n)=\varphi(n)\) for \(n\le B\), and \(\gamma(n)=n-\mathbf 1_{P\mid n}\)
for \(n>B\). Then \(0\le\gamma(n)\le n\), and the binary series
\(\sum\gamma(n)/2^n\) is \emph{rational}: it is a finite rational prefix, plus
the tail of \(\sum_n n2^{-n}\), minus a geometric sum over the multiples of
\(P\) beyond \(B\). Because \(\gamma\) agrees with \(\varphi\) on every index
\(\le B\), \(\gamma\) passes every finite certificate that only inspects
indices \(\le B\); in particular it passes or fails every instance of
\(\mathcal{C}(h,N,L)\) with \(N+h+L\le B\) exactly as \(\varphi\) does. Yet
\(\sum\gamma(n)/2^n\in\mathbb Q\) by construction.

\textbf{(c) Scope.} For each fixed bound $B$, a different rational
comparison sequence can agree with $\varphi$ on all inspected coefficients.
This rules out a deduction from that finite prefix alone. It does not rule
out a proof that establishes the required instances for arbitrarily large
$B$ with the correct quantifiers. Such a proof would no longer be a single
finite inspection. The comparison also does not preserve multiplicativity
or every other arithmetic property of $\varphi$.


\textbf{(d) Consequence.} The finite-prefix example and the parity example
below are distinct constructions. The former preserves a prescribed finite
prefix; the latter preserves parity and has separated nonperiodic blocks.
Each shows why its own list of assumptions is insufficient. The recorded
finite diagonal examples remain useful individual certificates, but their
number alone gives no assertion at arbitrarily large scales.
\lean{Erdos249257.certifiedKill\_diagonal\_all\_imported\_through\_t64}{DiagonalPincerCertificatesT64.lean:1967}


\subsection{A rational series with matching totient parity}
\label{sec:parity-countermodel}

\textbf{(a) Proposed argument.} Consider the claim that a uniformly
bounded sequence $c:\N\to\N$, with $c(n)\le n$, exact totient parity
and no eventual period, must have an irrational binary sum. The next
example refutes this claim. It also has arbitrarily many paired changes
beyond every threshold, with an arbitrarily prescribed minimum separation.

\textbf{(b) Construction.} Start with the sequence
$0,1,1,2,4,4,4,\ldots$. At each position $m=2^{k+3}$, $k\ge0$, replace
the pair $(4,4)$ at $m,m+1$ by $(6,0)$. Equivalently,
\[
 c(n)=
 \begin{cases}
  0,&n=0,\\
  1,&n=1,2,\\
  2,&n=3,\\
  6,&n=2^{k+3}\text{ for some }k\ge0,\\
  0,&n=2^{k+3}+1\text{ for some }k\ge0,\\
  4,&\text{otherwise}.
 \end{cases}
\]
Each replacement leaves the binary sum unchanged because
$2\cdot2^{-m}-4\cdot2^{-(m+1)}=0$. The increasingly large gaps between
replacements prevent eventual periodicity. The strongest recorded result
also quantifies their separation:
\begin{prop}[A rational series preserving totient parity and the stated separation properties]
There exists \(c:\mathbb N\to\mathbb N\) such that: \(c(n)\le6\) for all
\(n\); \(c(n)\le n\); \(c(n)\equiv\varphi(n)\pmod2\) for every \(n\); for
every \(N,G,K\) there is a block of \(K\) explicit \((6,0)\) carry-pulse
pairs beyond \(N\), each pair separated by more than \(G\); and
\(\sum_n c(n)/2^n = 3/2\).
\end{prop}
\leannote{Lean: \lproof{Erdos249257/TotientParityCoboundaryCountermodel.lean}{227}{parityCoboundaryWeight_le_six}, \lproof{Erdos249257/TotientParityCoboundaryCountermodel.lean}{155}{parityCoboundaryWeight_le_self}, \lproof{Erdos249257/TotientParityCoboundaryCountermodel.lean}{194}{parityCoboundaryWeight_mod_two_eq_totient}, \lproof{Erdos249257/TotientParityCoboundaryCountermodel.lean}{457}{exists_later_arbitrarily_many_separated_parityCoboundaryWeight_carry_pairs}, and 1 further declaration in the \hyperref[sec:coverage]{coverage section}.}
\lean{Erdos257PeriodNoncollapse.exists\_totientParity\_arbitrarilyManySeparatedCarry\_rational\_countermodel}{TotientParityCoboundaryCountermodel.lean:637}
(sum computation \lean{tsum\_parityCoboundaryWeight\_eq\_three\_halves}{TotientParityCoboundaryCountermodel.lean:359}; non-periodicity
\lean{parityCoboundaryWeight\_not\_eventually\_periodic}{TotientParityCoboundaryCountermodel.lean:487}; exact parity match
\lean{parityCoboundaryWeight\_mod\_two\_eq\_totient}{TotientParityCoboundaryCountermodel.lean:194}). \ev{Lean} \scale{cofinal}
\coord{binary-digit}.

\textbf{(c) Scope.} The displayed bounds, parity agreement and
aperiodicity, including the stated number and separation of the edits,
are jointly compatible with a rational binary sum. They therefore do not
imply irrationality. The sequence $c$ differs from $\varphi$: the example
preserves its parity, not all its arithmetic properties. It excludes an
argument using only the listed assumptions, not every strengthening of
an aperiodicity condition or every use of totient coefficients.

\textbf{(d) Consequence.} The paired edits have zero total contribution,
since $2\cdot2^{-m}-4\cdot2^{-(m+1)}=0$. Applied at the specified sparse
positions of this particular base sequence, they preserve the rational
sum and the required bounds while preventing eventual periodicity. This
does not construct a rational series for every prescribed parity pattern,
and it is not the finite-prefix construction of Theorem~\ref{thm:gamma}.
A proposed criterion using only the assumptions of this example must fail;
a criterion using further arithmetic information is not addressed.


\subsection{What the square-CRT construction proves}

For a prime $p$ and $m\ge1$,
\[
 \varphi(pm)=(p-1)\varphi(m)+\mathbf1_{p\mid m}\varphi(m).
\]
A congruence modulo $p^2$ can prescribe the cofactor modulo $p$
and hence remove the second term over a chosen finite range. This
makes each product formula simpler; it does not imply that differences
of the resulting values are nonzero.

The examples in Observation~\ref{prop:B8-kill} show both possibilities.
At $n=52$, the offset pairs $(13,13)$ and $(26,18)$ give respectively
$\varphi(65)-\varphi(65)=0$ and
$\varphi(78)-\varphi(70)=0$
(\lean{squareCRT\_clean\_block\_can\_vanish}{SquareCRTCube.lean:459}).
All four prime--cofactor products are coprime
(\lean{squareCRT\_vanishing\_countermodel\_is\_clean}{SquareCRTCube.lean:468}).
At $n=27$, the offset pairs $(3,3)$ and $(6,8)$ instead give
$\varphi(30)-\varphi(30)=0$ and
$\varphi(33)-\varphi(35)=-4$
(\lean{squareCRT\_clean\_block\_can\_be\_nonzero}{SquareCRTCube.lean:477}),
again with coprime factors
(\lean{squareCRT\_nonzero\_countermodel\_is\_clean}{SquareCRTCube.lean:485}).
These finite equalities are checked by exact arithmetic. No minimality
claim about either witness is needed. The general existence theorem
\lean{exists\_squareCRT\_clean\_horizon}{SquareCRTCube.lean:327}
supplies a common base for the stated finite congruence requirements,
not a nonvanishing assertion. \ev{Cert} \scale{fixed}
\coord{other:square-crt-cocycle}.

The vanishing example disproves the implication from correction-free
factorisation to nonvanishing; the other example shows that this
factorisation does not force vanishing either. This does not establish
formal independence of two propositions, or exclude an argument using
additional hypotheses. The prime-dilation identity
\lean{totient\_mul\_prime\_cocycle}{SquareCRTCube.lean:38}
and the CRT constructions
\lean{exists\_squareCRT\_base}{SquareCRTCube.lean:230},
\lean{exists\_large\_squareCRT\_base}{SquareCRTCube.lean:263}
remain available with their own prime and residue hypotheses.

A separate elementary identity in the same file concerns alternating
sums. Let $I$ be finite, let $i\notin I$, and suppose an
integer-valued function $f$ on finite subsets satisfies
$f(T\cup\{i\})=f(T)$ for every $T\subseteq I$. Pairing $T$ with
$T\cup\{i\}$ gives
\[
 \sum_{T\subseteq I\cup\{i\}}(-1)^{|T|}f(T)=0.
\]
This is the content of
\lean{alternating\_powerset\_sum\_eq\_zero\_of\_insert\_invariant}{SquareCRTCube.lean:363}.
It has no totient hypothesis, but an application must still verify
invariance under the indicated insertion. Neither it nor the congruence
construction supplies the missing nonzero totient difference at a
prescribed LCM height.

\subsection{The cost of clearing a rational denominator}
\label{sec:adelic-height}

For example, multiplying $5/12$ by $4$ gives $5/3$: the factor $4$
removed from the denominator appears in the multiplier. The following
lemma gives the exact divisibility behind this calculation. It concerns
ordinary reduced fractions, without assumptions on totients.

\begin{lem}[Complement divisibility after multiplication]
Write $x=a/b$ in lowest terms, where $a\in\Z$ and $b\ge1$. Let
$c\in\Z$ and let $H$ be a positive divisor of $b$. If the reduced
denominator of $cx$ divides $H$, then
\[
 b/H\mid c.
\]
\end{lem}
\leannote{Lean: \lproof{ErdosProblems/Erdos249/PaperCompleteR21/ScalarLocalisationAndInversePhaseGauge.lean}{29}{complementDenominator_dvd_scalar}.}
\lean{Erdos249257.AdelicHeightObstruction.scalarLocalization\_complement\_dvd}{AdelicHeightObstruction.lean:23}.
Indeed, the hypothesis makes $Hcx=Hca/b$ integral. Since
$\gcd(a,b)=1$, it follows that $b\mid Hc$, hence $b/H\mid c$.
Writing $c=(b/H)t$ also gives $Hcx=ta$; this equality is recorded in
\lean{Erdos249257.AdelicHeightObstruction.scalarLocalization\_integer\_eq\_mul\_num}{AdelicHeightObstruction.lean:56}.
If $c\ne0$, the divisibility implies $|c|\ge b/H$. The nonzero
hypothesis is essential for this size bound: $c=0$ has reduced
denominator one and satisfies the divisibility, but has no positive
size cost. The lemma is about multiplication of one rational value,
not arbitrary cancellations among several values.

A related elementary inequality compares a positive numerator with its
denominator. If $a>0$, $r\ge0$, $n\ge1$, $2^r\mid a$, and
$a/b<2/(2^n-1)$, then
\[
 2^r(2^n-1)\le a(2^n-1)<2b.
\]
The source gives this Mersenne instance
\lean{Erdos249257.AdelicHeightObstruction.positiveRat\_mersenne\_height}{AdelicHeightObstruction.lean:103}
and the version with general positive denominator parameter
\lean{positiveRat\_numDivisor\_mul\_lt\_two\_mul\_den}{AdelicHeightObstruction.lean:77}.
Positivity of $a$ is what turns numerator divisibility into a lower
bound; it cannot be omitted.

The accompanying linear-algebra statement has a one-dimensional
quotient. Let $V,W$ be vector spaces over $\Q$, let
$\mathrm{ev}:V\to\Q$ and $\Lambda:V\to W$ be linear, and suppose
$\ker\mathrm{ev}\subseteq\ker\Lambda$. Choose $e\in V$ with
$\mathrm{ev}(e)=1$. Then, for every $v\in V$,
\[
 v-\mathrm{ev}(v)e\in\ker\mathrm{ev},\qquad
 \Lambda(v)=\mathrm{ev}(v)\Lambda(e).
\]
This is factorisation through the scalar quotient $V/\ker\mathrm{ev}$
(\lean{Erdos249257.AdelicHeightObstruction.linearDescender\_eq\_smul\_eval}{AdelicHeightObstruction.lean:120}).
A surjection onto a higher-dimensional space would imply only a linear
factorisation through that space, not multiplication by one fixed vector.

The denominator calculations use
\lean{Erdos249257.RationalDenominatorSurvival.divisor\_dvd\_divInt\_den}{RationalDenominatorSurvival.lean:17}
and
\lean{survivingDivisor\_dvd\_scaled\_divInt\_den}{RationalDenominatorSurvival.lean:38}.
For integers $a$ and $D\ge1$, the reduced denominator of $a/D$ is
$D/\gcd(|a|,D)$. Hence a divisor of $D$ coprime to $a$ divides the
reduced denominator. Coprimality is sufficient, not necessary: $2/8=1/4$
retains the divisor $2$. More generally, if $C$ is a positive divisor of $D$, $\gcd(C,a)=1$,
and $h\ge0$ is an integer, then $C/\gcd(C,h)$ divides the reduced
denominator of $ha/D$. At $h=0$ this assertion retains only the divisor
one. \ev{Lean} \scale{uniform} \coord{other:rational-height}.

For a nonzero multiplier, a lower bound for $b/H$ therefore gives a
lower bound for its absolute value. For an integer $t\ge5$, the product estimate
\(2^{\lfloor t/2\rfloor}\le\prod_{p\in\mathcal P_t}(2^p-1)\)
(\lean{upperHalfMersenneProduct\_lower\_bound}{MersenneShadowDenominatorGrowth.lean:60})
can be used when that product is shown to divide $b/H$.
Neither the product estimate nor the scalar lemma proves a quadratic
growth law for arbitrary approximation schemes.

The finite Farey exclusion for $S$ is a separate use of reduced
denominators. Its certified interval contains no rational whose reduced
denominator is at most
\[
 79{,}639{,}646{,}646{,}701{,}375{,}323{,}355{,}774{,}875{,}831{,}053
\]
(\lean{tsum\_totient\_div\_pow\_two\_ne\_ratCast\_of\_den\_le\_79639646646701375323355774875831053}{CertificateKernel.lean:18384}).
That numerical exclusion requires the interval and Farey calculation;
it is not a consequence of the scalar lemma alone.

\subsection{Signed moments and parity}

A determinant construction needs a nonzero numerator as well as a
small approximation error. The following parity test can establish the
first requirement for a specified finite sum.

\begin{lem}[Nonvanishing from a unique largest denominator exponent]
Let $I$ be a finite set, let $u_i\in\Z$ and $e_i\in\N$ for $i\in I$,
and suppose that $m\in I$ satisfies $e_i<e_m$ for every $i\ne m$.
If $u_m$ is odd, then
\[
 2^{e_m}\sum_{i\in I}\frac{u_i}{2^{e_i}}
   =\sum_{i\in I}u_i2^{e_m-e_i}\equiv1\pmod2.
\]
In particular, the sum is nonzero.
\end{lem}
\leannote{Lean: \lproof{ErdosProblems/Erdos249/PaperCompleteR20/SignedDyadicClearing.lean}{8}{signed_dyadic_clearing}, \lproof{ErdosProblems/Erdos249/PaperCompleteR20/SignedDyadicClearing.lean}{28}{signed_dyadic_sum_ne_zero}, \lproof{Erdos249257/SignedQMomentObstruction.lean}{78}{scaled_dyadic_sum_odd}.}
Every term except $u_m$ is even, which proves the assertion. For
example, $1/2-1/8=3/8$ has odd cleared numerator. Uniqueness is
necessary for this argument: two odd contributions with the same
largest exponent can cancel. The formal statement is
\lean{scaled\_dyadic\_sum\_ne\_zero}{SignedQMomentObstruction.lean:96}.

The same file gives a separate determinant expansion from the Leibniz
formula
\lean{det\_mul\_rectangular}{SignedQMomentObstruction.lean:29},
together with the truncated signed moments
$\sum_{N=1}^Y a_N2^{-N(k+1)}$ and their Hankel determinants
(\lean{truncatedMoment}{SignedQMomentObstruction.lean:108} and
\lean{hankelDet}{SignedQMomentObstruction.lean:117}). The parity lemma
is not derived from that determinant identity. A determinant application
must first identify a sum to which the parity hypotheses apply.
\ev{Lean} \scale{uniform} \coord{p-adic}.

For the totient series, one would still need suitable configurations at
arbitrarily large scales and a denominator-weighted error tending to
zero. Nonvanishing alone provides neither. No such family is supplied
by these two algebraic statements; compare the separate rank discussion
in \S\ref{sec:mahler-defect}. In another base, the analogous argument
would require divisibility by a prime factor of that base rather than
an unchanged appeal to parity.

\subsection{Column rescaling does not certify cancellation}

A nonzero determinant need not prevent one row from being constant.
For example,
\[
 \begin{pmatrix}1&1\\1&-1\end{pmatrix}
 \begin{pmatrix}1&0\\0&-1\end{pmatrix}
 =\begin{pmatrix}1&-1\\1&1\end{pmatrix}.
\]
The first and last matrices have determinants $-2$ and $2$, respectively;
the second row of the resulting matrix is constant. This simple column rescaling
matters when a determinant is proposed as evidence for cancellation in
the first-harmonic sums
(\lean{FirstHarmonicPivot.lean}{FirstHarmonicPivot.lean}).

\begin{prop}[A nonzero minor survives inverse-phase column weights]
Let $d\ge1$ be an integer, let $e_0,\ldots,e_{d-1}$ be nonnegative integers, and let
$z_0,\ldots,z_{d-1}$ be nonzero complex numbers. Write
$P_{ij}=z_j^{e_i}$, and suppose $e_{i_0}=1$ for some $i_0$.
Multiplying column $j$ by $W_j=z_j^{-1}$ gives
\[
 (P\operatorname{diag}(W))_{i_0j}=1,\qquad
 \det(P\operatorname{diag}(W))=\det(P)\prod_{j=0}^{d-1}z_j^{-1}.
\]
Thus a nonzero determinant remains nonzero while row $i_0$ becomes
constant. If all $|z_j|=1$, its absolute value is also unchanged.
\end{prop}
\leannote{Lean: \lproof{ErdosProblems/Erdos249/PaperCompleteR21/ScalarLocalisationAndInversePhaseGauge.lean}{51}{inversePhaseGauge_locks_row_and_preserves_minor}.}
The entry formula and the determinant of a product prove the statement
(\lean{Erdos249257.locked\_reconstruction\_preserves\_nonzero\_minor}{ResidualGaugeObstruction.lean:94},
\lean{det\_residualMonomialMatrix}{ResidualGaugeObstruction.lean:47}).
If row-dependent weights are allowed, the choice
$W_{ij}=z_j^i/z_j^{e_i}$ gives $W_{ij}P_{ij}=z_j^i$ and makes row zero
constant. It produces the ordinary Vandermonde matrix, which is
nonsingular when the $z_j$ are pairwise distinct
(\lean{rowResidualMonomialMatrix\_reconstructs\_phasePowerMatrix}{ResidualGaugeObstruction.lean:124},
\lean{rowReconstruction\_zero\_row\_one}{ResidualGaugeObstruction.lean:145}).
\ev{Lean} \scale{uniform} \coord{other:first-harmonic-phase}.

These are statements about complex matrices, not arithmetic properties
of the totient sequence. In particular, determinant nonvanishing, even
with unit-modulus weights, does not by itself exclude complete
alignment in the indicated row. An application of the first-harmonic
criterion
(\lean{irrational\_totient\_series\_of\_first\_harmonic\_norm\_gap}{FirstHarmonicPivot.lean:83})
needs additional information that excludes such weights or controls the
resulting weighted sum. The example does not require that information
to take the form of an identity coupling different rows, and it does
not rule out determinant arguments that use appropriate additional
hypotheses.

\subsection{Tail integrality on an LCM grid}

\textbf{(a) The comparison.} The diagonal condition
(\lean{Erdos249257.irrational\_totient\_series\_iff\_lcm\_diagonal\_certificate\_supply}{LcmConeFlatness.lean:426})
uses the pair $(H_t,2H_t)$. Under rationality, the same fractional part
occurs at every positive multiple of $H_t$ once $t$ is sufficiently
large. A finite set of these multiples therefore gives another possible
place to test for a nonintegral difference. The criterion below does not
assert a smaller successful depth than the diagonal test.

\textbf{(b) Eventual equality of fractional parts.}
\begin{prop}[Tail integrality on an LCM grid]
\(S\in\mathbb Q \implies \exists t_1,\ \forall t\ge t_1,\
\forall q,m:\mathbb N,\ 0<q \implies {R}(q\cdot
{H}(t)+m\cdot{H}(t)) -
{R}(q\cdot{H}(t)) \in
\mathbb Z\).
\end{prop}
\leannote{Lean: \lproof{ErdosProblems/Erdos249/PaperCompleteR20/LcmGridCorrespondence.lean}{8}{lcm_grid_flatness}.}
\lean{rational\_totient\_series\_forces\_lcm\_cone\_flatness}{CertificateKernel.lean:19002} (body
lcm cone flatness source). \ev{Lean} \scale{uniform} (for all
sufficiently large \(t\), \emph{given} rationality) \coord{other:lcm-ray-totient}.
Rationality makes the fractional parts of $R_{kH_t}$ equal for all
$k\ge1$ at each scale $t\ge t_1$. The corresponding implication
\lean{irrational\_totient\_series\_of\_lcm\_cone\_window\_kill\_supply}{CertificateKernel.lean:19014} shows that a certificate $\mathcal C(mH_t,qH_t,L)$ with $q\ge1$
at arbitrarily large $t$ implies $S\notin\mathbb Q$.


\textbf{(b\('\)) A finite test on the grid.}
\begin{prop}[A finite-grid certificate gives a nonintegral pair]
Let $Q\subseteq\mathbb N_{>0}$ be finite and nonempty. For each $q\in Q$,
put $A_q=\sum_{j=1}^{L}\varphi(qH+j)2^{L-j}$ and $B_q=qH+L+2$.
Suppose $B_q<2^L$ for every $q\in Q$ and
\[
 \forall q_i\in Q\ \exists q_j\in Q,\qquad
 B_{q_j}<(A_{q_i}-A_{q_j})\bmod2^L.
\]
Then $R_{q_jH}-R_{q_iH}\notin\mathbb Z$ for some $q_i,q_j\in Q$.
The proof and a four-point example are given in
Theorem~\ref{catalogue:cert:b10a}.
\end{prop}
\leannote{Lean: \lproof{ErdosProblems/Erdos249/PaperCompleteR20/FiniteGridCorrespondence.lean}{12}{paperGridNumerator_eq}, \lproof{ErdosProblems/Erdos249/PaperCompleteR20/FiniteGridCorrespondence.lean}{36}{finite_grid_nonintegral_pair}.}
\lean{exists\_nonintegral\_pair\_of\_coneNonflatCert}{LcmConeNonflat.lean:126} (quantified assertion
\lean{irrational\_totient\_series\_of\_lcm\_cone\_nonflat\_supply}{CertificateKernel.lean:19140}). \ev{Lean}
\scale{bounded} (proved for each finite set satisfying the hypotheses;
the required family at arbitrarily large $t$ is not established)
\coord{other:lcm-ray-totient}.

\textbf{(c) Scope.} Eventual integrality on the LCM grid is a necessary
consequence of rationality, not a claim about the actual series made
without that hypothesis. The finite-grid test uses the individual
one-sided bounds $B_q<2^L$, rather than charging the larger radius to
both ends of a pairwise interval. This compares modulus bounds, not
binary depths. A finite nonempty set and every stated size inequality
are needed; three or more points and pairwise compatibility are not.

Neither statement supplies the unbounded family required to decide
\#249. Exhibiting a nonintegral tail difference on the LCM grid at
arbitrarily large $t$ would refute rationality. The finite-grid test is a
way to certify such a difference, not a proof that these tests succeed
at arbitrarily large scales.

The implication from rationality to grid flatness
(\lean{rational\_totient\_series\_forces\_lcm\_cone\_flatness}{LcmConeFlatness.lean:90})
should not be confused with an assertion of flatness for the actual series.
The exact diagonal integrality reformulation and its quantified converse
(\lean{diagonal\_int\_iff\_foreignDiagonalDefect\_hits\_fullTarget}{DiagonalPincerDecomposition.lean:215},
\lean{irrational\_totientSeries\_of\_full\_target\_avoidance\_supply}{DiagonalPincerDecomposition.lean:290})
give other equivalent forms of the problem. They are not the only such
forms; Proposition~\ref{prop:iffs} lists further equivalent certificate
conditions.

\textbf{(d) Consequence.} Choosing the least deep tail is the step that
turns a common-fractional-part assumption into an intersection of arcs.
A related finite comparison concerns first and second differences: at
$(h,N)=(1,8)$ the first-difference test holds at depth $8$, while no
second-difference certificate exists at depth at most $8$
(\lean{totient\_tail\_rank\_two\_kill\_sound\_but\_not\_shallower\_cell}{CertificateKernel.lean:19090}).
This disproves a universal claim of smaller depth for second differences;
it does not show that they can never improve a certificate.

\subsection{Exact dyadic rank and the limits of a rank argument}
\label{sec:mahler-defect}

\textbf{(a) The question.} Which linear relations hold among the
sections $n\mapsto\ph(2^j n+r)$, and how does their span grow as the
level increases? Coons had already proved nonregularity in every base
$k\ge2$~\cite[Theorem~3.2]{coons} \ev{Cited}. The result below identifies
a basis and gives the exact rank at each dyadic level. Martin's
Theorem~1 implies the independence of the positive-residue affine
forms~\cite[Theorem~1]{martin-phi-inequalities}; the two sections with
residue zero need a separate argument. The short note gives the
all-base proof and the integral relations.

\textbf{(b) Construction.}
\begin{thm}[Exact dyadic rank and infinite-dimensionality]
For every $e\ge1$, the family
\[
 \{n\mapsto\ph(n),\ n\mapsto\ph(2n)\}
 \ \cup\ \{n\mapsto\ph(2^j n+r):1\le j\le e,\ 0<r<2^j,\ r\text{ odd}\}
\]
is linearly independent over $\Q$. It contains $2^e+1$ sequences
(\lean{card\_totientCanonicalIndex}{TotientMahlerDefect.lean:61}) and spans
the sections through level $e$ by the zero- and even-residue reductions.
\end{thm}
\leannote{Lean: \lproof{ErdosProblems/Erdos249/PaperCompleteR21/DyadicSectionBasisAndRationalCarry.lean}{33}{canonicalTotientKernelFamily_entries}, \lproof{ErdosProblems/Erdos249/PaperCompleteR21/DyadicSectionBasisAndRationalCarry.lean}{53}{canonicalTotientKernelFamily_independent_card_and_span}.}
The forms $n$ and $2n$ are proportional, so their coefficients need a
separate argument. Restrict a proposed relation to $n=2m+1$, where $\ph(2n)=\ph(n)$.
The odd-residue sections and $\ph(n)$ then correspond to pairwise
nonproportional forms. CRT and Dirichlet's theorem separate them by
an evaluation matrix: make the primitive part of one form prime and
force suitable prime divisors into the others. This kills the
positive-residue coefficients and gives $A+B=0$ for the two remaining
coefficients. Evaluating their relation at $n=2$ gives $A+2B=0$.
Hence both vanish. This gives the dyadic specialization of the proof
in the short note.
The formal independence statement is
\lean{linearIndependent\_canonicalTotientKernelFamily}{TotientMahlerDefect.lean:935}; consequently the full
infinite family spans an infinite-dimensional \(\mathbb Q\)-space,
\lean{not\_finiteDimensional\_span\_fullTotientKernel}{TotientMahlerDefect.lean:1145}. A separate elementary observation concerns a proposed integer identity.
There are no $Q,v\in\N_{>0}$ and $A,b\in\Z$ with
$QvA=b$, $A\ne0$, and $|b|<Qv$: the identity forces
$|b|=Qv|A|\ge Qv$.
The source packages these incompatible requirements as one structure,
\lean{false\_of\_compressedAdjointCertificate}{TotientMahlerDefect.lean:1183}. \ev{Lean}
\scale{uniform} (holds for every \(e\ge1\); existence side is unconditional, via
Mathlib's CRT + primes-in-AP machinery) \coord{binary-digit} (dyadic
totient-kernel; \emph{not} the Möbius coordinate).

The Lean independence theorem is also stated for \(e=0\), but its auxiliary
canonical index then still contains the two zero-residue sections
\(\varphi(n)\) and \(\varphi(2n)\).  It is therefore not the actual truncation
through level zero, which consists only of \(\varphi(n)\) and has dimension one.

\paragraph{Integral relations.}
The independence statement also determines the arithmetic of the section
lattice.  Fix $e\ge1$, write $F_{j,r}(n)=\varphi(2^j n+r)$, and let
$I_e=\{(j,r):0\le j\le e,\ 0\le r<2^j\}$.  Retain $F_{0,0}$,
$F_{1,0}$, and the odd-residue sections at levels $1,\ldots,e$; call
this retained index set $J_e$.  The totient identities reduce every omitted
section to $F_i=a_iF_{j(i)}$ for a retained section $j(i)$ and an integer
$a_i\ge0$.  Thus the retained family generates the lattice
$L_e=\sum_{i\in I_e}\mathbb Z F_i$.  Its rational independence gives
integer independence as well: every element of $L_e$ has unique integer
coordinates in these $2^e+1$ retained sections.

To describe \emph{all} integral relations, let $E_i$ denote the formal
coordinate vector in $\mathbb Z^{I_e}$ and use the evaluation map
\[\adjustbox{max width=\linewidth}{$\displaystyle
 \operatorname{ev}_e:\mathbb Z^{I_e}\longrightarrow L_e,
 \qquad (c_i)\longmapsto\sum_{i\in I_e}c_iF_i.
$}\]
For each omitted index put $R_i=E_i-a_iE_{j(i)}$.  Then
\[\adjustbox{max width=\linewidth}{$\displaystyle
 \ker(\operatorname{ev}_e)
   =\bigoplus_{i\in I_e\mathbin{\backslash} J_e}\mathbb Z R_i,
 \qquad
 \operatorname{rank}_{\mathbb Z}\ker(\operatorname{ev}_e)=2^e-2.
$}\]
Indeed, subtracting $\sum_{i\notin J_e}c_iR_i$ from a relation removes
all omitted coordinates.  The remainder is a relation among retained
sections, so independence makes it zero.  Conversely, the coefficient of
$E_i$ in $\sum_{h\notin J_e}b_hR_h$ is exactly $b_i$, proving uniqueness.
The decisive feature is the unit pivot $1$ in each omitted coordinate
(or $-1$ if a relation is oriented oppositely).  Elimination uses no division;
it therefore identifies the integral relation lattice, not merely its
rational span.  The depth-two example in the short note illustrates these
two-term rows.  The distinction between the retained basis and the redundant
spanning family of all sections already matters for $k$-regular sequences
that are not eventually zero: their growth exponent is the base-$k$
logarithm of the joint spectral radius of their section matrices in a
basis, whereas a redundant spanning family gives only an upper bound
(Coons~\cite[Theorem~1, Proposition~4 and Corollary~7]{coons-jsr}) \ev{Cited}.  The totient is not regular, so that theorem does not
apply; the short note records how the section maps act on the retained bases
from one level to the next.

The \href{https://github.com/wcook04/plectis-erdos/blob/25ef6245d15a47548c6926369ae8f1a0f0a14a80/ErdosProblems/Erdos249/PaperCompleteR8/KernelRelationBasis.lean#L398}{integral normal form} and the \href{https://github.com/wcook04/plectis-erdos/blob/25ef6245d15a47548c6926369ae8f1a0f0a14a80/ErdosProblems/Erdos249/PaperCompleteR8/FullKernelAssemblies.lean#L156}{full dyadic rational basis} are formally checked in Lean.
The formal statement gives the same construction for every integer base
$k\ge2$, with relation rank $\sum_{j=1}^{e-1}k^j$.  The full dyadic
assembly separately gives a basis for the infinite rational span and for
its finitely supported rational relations.  These finite and infinite
statements should not be confused with a claim about infinite sums of
relations.

\textbf{(c) What the rank theorem does not prove.}
Infinite section rank does not imply irrationality of $S$. The result
quantifies the nonregularity already cited above: it gives the rank,
the retained sections and their integral relations. Martin's simultaneous
inequalities supply another proof for pairwise nonproportional affine
forms with positive slopes~\cite[Theorem~1]{martin-phi-inequalities}
\ev{Cited}. His theorem does not separate $n$ from $2n$; the two
evaluations above do that. The linked Lean proof is independent of
Martin's theorem.

The same independence argument also bounds the rank of an integral carry.
If $S\in\Q$, choose an integer $v>0$ with $vS\in\Z$ and put $u_N=vR_N$,
where $R_N=\sum_{j\ge1}\ph(N+j)2^{-j}$. Then $u_N\in\Z$,
\[
 u_{N+1}=2u_N-v\ph(N+1),\qquad
 0\le u_N\le v(N+2),\qquad u_N/2^N\longrightarrow0.
\]
For every $e\ge0$,
\[
 \dim_{\Q}\operatorname{span}
 \{n\mapsto u_{2^j n+r}:1\le j\le e,\ 0\le r<2^j\}
 \ \ge\ 2^e-1.
\]
This is the finite family in the linked statement
(\lean{not\_irrational\_totientSeries\_implies\_unbounded\_carryRank\_unconditional}{TotientCarryKernelRigidity.lean:300});
it contains all carry residues at levels $1,\ldots,e$, not just odd
residues or an unspecified retained subfamily. Indeed, for each odd
$r$ with $0<r<2^j$, the recurrence gives
\[
 v\ph(2^j n+r)=2u_{2^j n+r-1}-u_{2^j n+r}.
\]
These $2^e-1$ totient sections are independent and belong to the span
of the displayed carry sections after multiplication by $v\ne0$.
The rank bound follows. The boundary condition $u_N/2^N\to0$ identifies
the canonical carry, but is not used in this rank-transfer argument.

To contradict rationality by rank, one would also need an incompatible
upper bound. The comparison series with value $5/4$ in
Section~\ref{sub:generic-carry-rank} shows why
rationality, a binary carry recurrence and linear coefficient growth
alone do not supply such a bound. Its coefficients agree with $\ph$
at odd indices but differ at even indices; they are not the totient
sequence. Thus this comparison does not exclude an argument using
additional totient identities.

A separate finite-dimensional calculation concerns the matrix
$U_N(i,j)=\mu((i+1)/(j+1))$ when $(j+1)\mid(i+1)$ and zero otherwise,
for $0\le i,j<N$. It is lower triangular with diagonal entries one.
Hence $\det U_N=1$ and the map $x\mapsto U_Nx$ is injective
(\lean{incidenceMobius\_det\_eq\_one}{IncidenceQuotientHermitePade.lean:56}).
This proves injectivity of that particular map, not the impossibility
of every finite-dimensional argument for irrationality.

\textbf{(d) The evaluation-matrix argument.}
For functions $f_1,\ldots,f_s$ with values in a field, points
$n_1,\ldots,n_s$ for which $\det(f_j(n_i))\ne0$ prove linear
independence: evaluating a relation gives an invertible linear system
for its coefficients. This elementary criterion is the content of
\lean{SeparatedMinorCertificate}{TotientMahlerDefect.lean:83} and
\lean{linearIndependent\_of\_separatedMinorCertificate}{TotientMahlerDefect.lean:91}.
It does not, by itself, construct the points. Here the totient formula,
CRT and Dirichlet's theorem provide them for nonproportional affine
forms, followed by the separate zero-residue argument. Another
arithmetic function would require its own argument for nonvanishing
of the evaluation determinant.

\subsection{Further counterexamples and sharp bounds}

The following table records further counterexamples and bounds. Each row
states the argument, identifies its source and specifies what remains
possible beyond the assumptions it tests.

\begin{landscape}
\small
\begin{longtable}{@{}>{\raggedright\arraybackslash}p{0.20\linewidth}>{\raggedright\arraybackslash}p{0.38\linewidth}>{\raggedright\arraybackslash}p{0.36\linewidth}@{}}
\toprule
\papertablehead
\textbf{Approach} & \textbf{Argument and source} & \textbf{Scope and remaining question} \\
\midrule
\endfirsthead
\toprule
\papertablehead
\textbf{Approach} & \textbf{Argument and source} & \textbf{Scope and remaining question} \\
\midrule
\endhead
\bottomrule
\endlastfoot
Complete terminal divisibility pattern &
Under $a\ge8$, $m\le K$, $J+K+a+6<2\cdot2^a$ and
$2H+J+K+2<2^m$, the full set of divisibilities in
Observation~\ref{prop:TE-01-killer} is impossible: the last difference
would be a positive multiple of $2^m$ smaller than $2^m$.
\lean{not\_actualLcmTerminalDyadicStaircase\_of\_room}{TotientActualLcmTopEdgeStaircase.lean:251}.
\ev{Lean} \scale{bounded} \coord{binary-digit} &
The obstruction is specific to that pattern and those bounds.
If the last divisibility is omitted, the additional terminal bound of
Proposition~\ref{prop:TE-02-inv} forces the penultimate difference to be
$2^{m-1}$. It does not prove that such partial patterns exist.
\lean{puncturedDyadicStaircase\_penultimate\_eq\_half}{TotientActualLcmTopEdgeStaircase.lean:1187}. \\
Finite linear combinations of shifted carry recurrences &
The explicit sequence in Observation~\ref{prop:B4b-kill} has
$a_{(q-1)H-1}=\varphi(H)$ for $2\le q<t$ and a bounded state satisfying
$a_i=2c_i-c_{i+1}$. Every finite integer combination of its shifts
retains a recurrence of this form, including products of differences
along multiples of $H$. Its bounds acquire the factor $\sum_j|\lambda_j|$;
they are not uniform in the coefficients of the linear combination.
The original module introduction is at
\lean{LcmFactorIdealPulseObstruction.lean:1--24}{LcmFactorIdealPulseObstruction.lean:1-24};
the explicit theorem and ordinary proof are given in
Observation~\ref{prop:B4b-kill}. \ev{Math; formal source cited} \scale{uniform}
\coord{other:lcm-ray-totient} &
The constructed coefficients are not asserted to be actual totient
differences. Preserving the displayed recurrence under shifts therefore
does not distinguish them from the coefficients required in the proposed
argument. The full totient divisor identity
\lean{totient\_eq\_sum\_mobiusTotientChannel}{LcmFactorIdealPulseObstruction.lean:328}
retains information absent from the recurrence alone. The counterexample
does not require every future proof to use that particular identity. \\
Adding three multiplicative scales to an integrality test &
For $H\ge0$ and positive integers $p,q$, requiring
$R_{2kH}-R_{kH}\in\Z$ for every $k\in\{1,p,q,pq\}$ is equivalent to
requiring only $R_{2H}-R_H\in\Z$. Neither multiplier need be prime.
The affine transport formula along $H\mapsto kH$ carries integrality at
the first scale to each of the other three.
\lean{Erdos249257.fullTarget\_primeAdjunction\_diamond\_iff\_root}{FullTargetPrimeAdjunctionNoGo.lean:196}.
\ev{Lean} \scale{uniform} \coord{other:mobius-mersenne} &
Adding these three integrality assertions does not strengthen the first
one. A different test involving part of the tail would need an estimate
for the part it discards; this equivalence provides no such estimate.
The result concerns the displayed four conditions, not every test using
several scales. \\
Prescribing valuations and finite unit residues &
For the valuations $\nu_j$ and odd residues $a_j$ of
Proposition~\ref{prop:b5}(iii), the free integers $z_j$ can always be
chosen so that $e_{j+1}=2e_j+2^{\nu_j}(a_j+2^u z_j)$ lies in the
stated centred interval at every step.
\lean{Erdos249257.TotientTailPeriodKiller.fixedPrecisionTropicalNoGo}{TropicalCurvatureCarry.lean:137}.
\ev{Lean} \scale{bounded} \coord{p-adic} &
The specified residues do not exclude all finite carry sequences while
the higher quotients $z_j$ remain unrestricted. An argument would need
further constraints on those quotients or on the actual coefficients.
The lemma does not assert that its freely chosen coefficients are totient
differences. \\
The information retained by an endpoint residue &
For an integer sequence $a$ and integer initial values $u_0,v_0$,
let $u_{n+1}=2u_n-a(n+1)$ and $v_{n+1}=2v_n-a(n+1)$.
Induction gives $u_L-v_L=2^L(u_0-v_0)$ for every $L\ge0$.
\lean{Erdos249257.affineBinaryOrbit\_mod\_twoPow\_eq}{GenericTailOrbitRigidity.lean:307}. \ev{Lean}
\scale{uniform} \coord{binary-digit} &
The endpoints agree modulo $2^L$, but they are distinct integers
whenever $u_0\ne v_0$. Indeed, if the forcing terms are known, then
\[
 u_0=\frac{u_L+\sum_{j=1}^{L}2^{L-j}a(j)}{2^L}.
\]
Only reduction modulo $2^L$ discards the initial value. The identity
does not exclude an argument that retains the full endpoint or other
information about the sequence. \\
Bounded finite-state / autonomous successor decoder for the binary carry &
For the family in Proposition~\ref{prop:b5}(i), with $m\ge2$, all
histories before $m$ agree but the next tail takes
$\lfloor(m+1)/2\rfloor+1$ distinct values. A common state cannot determine
that tail; a set of labels permitting exact recovery must have at least
that many elements.
\lean{Erdos249257.balancedPulse\_no\_autonomous\_decoder}{GenericTailOrbitRigidity.lean:247}. \ev{Lean}
\scale{uniform} \coord{binary-digit} &
This counterexample concerns states determined only by the common
history of the displayed family. It does not exclude a state using
additional information specific to the totient sequence. The bound counts
labels, not bits; it is not a complexity theorem about the binary digits
of $S$. \\
Zeros on a progression versus nonzero coefficients arbitrarily far out &
Let $\chi_4(n)$ be $0$ for even $n$, $1$ for $n\equiv1\pmod4$,
and $-1$ for $n\equiv3\pmod4$, and put
$c(n)=\sum_{d\mid n}\chi_4(d)$ for $n\ge1$.
Divisor pairing gives $c(n)=0$ when $n\equiv3\pmod4$.
\lean{Erdos249257.intWeightedCoeff\_periodFourSignWeight\_eq\_zero\_of\_mod\_four\_eq\_three}{CertificateKernel.lean:14226}.
Nevertheless $c(2^k)=1$ for every $k\ge0$, since $1$ is the only
odd divisor of $2^k$.
\ev{Lean (progression identity)} \scale{uniform} \coord{mobius-mersenne} &
This rules out nonvanishing on an arbitrarily prescribed progression,
not nonzero coefficients arbitrarily far out. In fact $c$ is
multiplicative. Its prime-power values are $1$ at powers of $2$,
$a+1$ at $p^a$ with $p\equiv1\pmod4$, and $1$ or $0$ at $p^a$
with $p\equiv3\pmod4$ according as $a$ is even or odd.
Thus $c(n)\ge0$ for every $n\ge1$. These elementary calculations
verify both coefficient hypotheses of
\lean{irrational\_intWeightedErdosSeries\_periodic\_of\_coeff\_nonneg\_of\_frequently\_ne\_zero}{CertificateKernel.lean:14583}:
$\sum_{a\ge1}\chi_4(a)/(b^a-1)$ is irrational for every integer
$b\ge2$. This application concerns the displayed weighted Lambert
series, not the totient series $S$. \\
Using a finite exclusion test to prove membership &
\(\text{finite Mersenne-weight exclusion}(x,\mathrm{level},\mathrm{lookahead})\)
is a decidable, finite-depth certificate proving non-membership only
(\lean{greedyMersenneDeath\_not\_mem}{GreedyAchievementSet.lean:1747}; witness
\lean{three\_fourths\_certifiedGreedyMersenneDeath}{GreedyAchievementSet.lean:1778}, \(3/4\) excluded at
level 1). \ev{Cert} \scale{bounded} \coord{greedy-orbit} &
The cited test has a sound exclusion implication. The analogous
soundness statements in
\lean{TotientTailPeriodKiller.lean}{TotientTailPeriodKiller.lean} and
\lean{LcmConeFlatness.lean}{LcmConeFlatness.lean} also do not make a
finite unsuccessful search a proof of integrality. This is a limit of
what follows from those tests, not an impossibility theorem for
rationality or membership proofs using additional identities. \\
Reconstructing one totient value from a two-tail absolute-value linear
combination &
For any finite \(w:\iota\to\mathbb Q\), \(x:\iota\to\mathbb N\) with
\(\sum w_i\varphi(x_i)=1\) (exact isolation), the crude two-tail cost
\(\sum|w_i|(2(x_i+1)+(x_i+2))\ge3\), using only \(\varphi(x)\le x\).
\lean{Erdos257PeriodNoncollapse.three\_le\_totientAdjugateTailCost}{TotientTailCarryPeriod.lean:266}
(closure \lean{not\_totientAdjugateTailCost\_lt\_one}{TotientTailCarryPeriod.lean:307}). \ev{Lean}
\scale{uniform} \coord{other:totient-adjugate-linear-algebra} &
The displayed termwise absolute-value estimate cannot be less than $1$.
The same inequality holds for any coefficients $0\le c(n)\le n$ with
$\sum_iw_ic(x_i)=1$. It does not rule out a smaller error bound that uses
cancellation between the tails, or a different approximation identity. \\
Homogeneous Mersenne-primitive prime factor alone forces a contradiction &
For the completely multiplicative sequence $c(n)=n$, if
$0<K<q$ and $q\mid2^K-1$, then the scaled-tail difference is the
integer $K$ at every starting index, but $q\nmid K$. \lean{idCoeff\_mersenneModulus\_does\_not\_annihilate\_tailShift}{TotientTailCarryPeriod.lean:342}.
\ev{Lean} \scale{n/a} \coord{other:mersenne-modulus} &
A factor of $2^K-1$ does not, from the recurrence alone, have to
divide the inhomogeneous tail difference. The example $K=3$, $q=7$
already has a primitive prime factor. An application to another series
would need a correspondence between its coefficients and this recurrence;
Mersenne denominators alone do not provide one. \\
The unit-gap test at a prime-power denominator &
For $p$ prime and $e\ge1$, if every candidate integer in the interval
at denominator $p^e$ is a nonunit modulo $p^e$, then there is at most
one candidate. Equivalently, the notation of
Observation~\ref{prop:D8cert-kill} satisfies $L=0$, or $L=1$ with
$p\mid J+1$. \lean{reducedDenominatorCandidateCount\_prime\_pow\_le\_one}{PrimitiveWeightCertificate.lean:111}
(iff-characterisation \lean{reducedDenominatorUnitGapCert\_prime\_pow\_iff}{PrimitiveWeightCertificate.lean:54}). \ev{Lean}
\scale{uniform} \coord{other:farey-gap} &
This candidate-count bound does not apply to every odd denominator;
the example modulo $15$ in Observation~\ref{prop:D8cert-kill} has two
candidates. It also does not bound the denominator cutoff obtainable
by changing the window or approximation.
It supplies no proof that the fixed exclusion in
Section~\ref{sec:adelic-height} extends to arbitrarily large denominators. \\
Extra divisibility of a three-term totient difference &
At $H=2^{n+2}$, $j=3\cdot2^{n+1}$,
$\varphi(3H+j)-2\varphi(2H+j)+\varphi(H+j)=-2^{n+1}$ for every $n\ge0$.
\lean{fixedRankSecondDifference\_dyadic\_fixture\_exact}{TotientFixedRankLcmAsymptotic.lean:462}. \ev{Lean}
\scale{uniform} \coord{binary-digit} &
The exact valuation concerns this scaled family, not the LCM family
with $j^2\le2^a$; see Observation~\ref{prop:FR-03-kill}.
The independent coefficient calculation says that integer solutions of
$c_1+c_2+c_3=c_1+2c_2+3c_3=0$ are exactly the multiples of $(1,-2,1)$
(\lean{affine\_rank\_three\_kernel\_classification}{TotientFixedRankLcmAsymptotic.lean:338}).
It classifies these three coefficients; it does not determine the
valuation of their totient combination on every restricted family. \\
\end{longtable}
\end{landscape}

\section{Why the stated hypotheses cannot be omitted}
\label{sec:promotion-audit}
% The replaced, uncaptioned longtable advanced this counter. Preserve the
% established number of the later labelled numerical table.
\stepcounter{table}

A conditional theorem can hold at every scale while its arithmetic
hypothesis is known at only a few scales, or at none. The distinctions
below concern the specific implications already stated. They are not
claims that every argument using the same objects must fail.

\subsection*{Implications concerning the totient series}

\paragraph{Cancellation versus a certificate.}
The first-harmonic implication converts the stated cosine bound on one
block into a finite residue certificate
(\lean{exists\_certifiedKill\_of\_first\_harmonic\_gap}{FirstHarmonicGap.lean:128}). To prove irrationality, the bound must hold with the
quantifiers in Section~\ref{sec:survivors}: for every shift and beyond
every block threshold. A finite identity for the sum, or a computation
at one block, does not establish this assertion.

\paragraph{Sign versus distance from an integer.}
The positivity theorem for the LCM tail difference
(\lean{actualLcmTailDiff\_shift\_pos}{TotientActualLcmOrbitSign.lean:39}) identifies the relevant endpoint under the stated
room and modulus bounds. Nonintegrality follows only after excluding
the corresponding residue interval
(\lean{actualLcmTailDiff\_notMem\_int\_of\_topEdgeResidueGap}{TotientActualLcmTopEdgeStaircase.lean:1260}).
Likewise, the terminal-dominance condition is a sufficient hypothesis,
not a construction of the required parameters
(\lean{PowerTwoFlexibleActualTerminalDominanceSupply}{TotientActualLcmTopEdgeStaircase.lean:1908}).
Under integrality its inequality contradicts positivity of the terminal
carry; the contradiction establishes sufficiency, not the inequality itself.

\paragraph{An error bound versus a separated approximation.}
The finite-tail approximation has a uniform error radius
(\lean{abs\_actualLcmTailOrbit\_sub\_rawApprox\_lt}{TotientActualLcmOrbitSeparation.lean:141}). The cited exact check supplies nonintegrality at the two
exponents $a=4,6$
(\lean{actualLcmTailOrbit\_four\_and\_six\_notMem\_int}{TotientActualLcmShortKill.lean:35}); it is not an assertion that no other exponent works.
To obtain irrationality by this criterion, separation must occur at
arbitrarily large LCM scales. In the nondivisor approximation of
Section~\ref{sec:nondivisor-enclosure}, the error bound is also proved
when $D\ge2H$. Its remaining requirement is finite rational separation,
not an additional convergence hypothesis.

\paragraph{Denominator growth versus approximation quality.}
The lower bound for the product of the upper-half Mersenne factors
(\lean{upperHalfMersenneProduct\_lower\_bound}{MersenneShadowDenominatorGrowth.lean:60}) concerns a denominator expression. An irrationality argument
needs a comparison with the error of the corresponding approximation.
Growth of that denominator expression alone does not give either a
successful residue test or a nonintegral totient tail.

\paragraph{A rank lower bound versus a contradiction.}
The coefficient independence and hypothetical-carry rank bounds hold
uniformly in the level
(\lean{linearIndependent\_canonicalTotientKernelFamily}{TotientMahlerDefect.lean:935}, \lean{not\_irrational\_totientSeries\_implies\_unbounded\_carryRank\_unconditional}{TotientCarryKernelRigidity.lean:300}).
A contradiction would require an upper bound for the actual carry that
is strictly smaller than the lower bound at some level. The rational
$5/4$ comparison disproves the proposed generic upper bound. The separate integer boundary conditions in the cited
adjugate calculation are inconsistent
(\lean{false\_of\_compressedAdjointCertificate}{TotientMahlerDefect.lean:1183}), but that calculation does not exclude matrix arguments
with different boundary conditions.

\subsection*{Related source comparisons for Problem~\#257}

The following related statements concern Problem~\#257. They are not
inputs to any proof about $S$; the comparisons below do not extend the
cited statements beyond their hypotheses.

The comparison of the middle-carry estimate
(\lean{one\_third\_lt\_scaledMersenneTail\_floorWall}{HalfCylinderMiddleCarryLowerBound.lean:2308}) with the remainder estimate
(\lean{SeamGreedyRemainderGapBound}{HalfCylinderSeamLimit.lean:192}) distinguishes the direction of an inequality from control
of the required deviation. The associated sign claim in the earlier
working notes had no unconditional proof supplied. A sign alone would
not give a lower bound for the magnitude in any event.

The row-capacity and later-skip implications retain their respective
local deficit and skip hypotheses
(\lean{exists\_exactRowStrictUpperFill\_of\_skippedCoreSharpCapacity}{BooleanMobiusSkipRow.lean:238}, \lean{halfGreedy\_precriticalSuffix\_lt\_of\_future\_skip\_after\_takenBlock}{BooleanMobiusCriticalCapacityCofinal.lean:603}).
The finite-shadow comparison distinguished an existential horizon,
obtained from convergence, from the particular horizon $J=k-2$
needed in that application
(\lean{exists\_greedyHalfFrozenMargin\_nonneg\_of\_excess\_neg}{HalfCylinderFiniteShadow.lean:1198}).
The terminal-strip formulation also imposes its own support and carry
restrictions
(\lean{HalfCarryCofinalTerminalOnlyStrip}{TerminalOnlyCofinal.lean:34}); a bound on carry size does not automatically replace a
restriction on the positions of the nonzero coefficients.

For finite sums of $1/(2^n-1)$, the common denominator is odd.
Their reduced denominators are therefore odd, so they cannot equal a
rational number with even reduced denominator. This is the elementary
reason behind the finite-support exclusion and its denominator source
(\lean{finite\_boolSupport\_ne\_half}{HalfCarryReachability.lean:589}, \lean{finiteErdosSum\_den\_odd}{DyadicPrefixCompression.lean:1513}).
The cited infinite-skipping theorem is stated for the target $1/2$
(\lean{infinite\_greedyMersenneSkippedSupport\_of\_half\_mem}{GreedyAchievementSet.lean:2547}). A proposed extension to other targets requires checking
that theorem's hypotheses and proof; the finite-denominator observation
alone does not prove such an extension or the needed infinite-orbit estimate.

Finally, the support-fraction statement and the more general collision
bound have distinct hypotheses
(\lean{dyadic\_support\_fraction\_reciprocalMass\_diverges\_or\_gt\_one}{RationalSupportCarrySkeleton.lean:2210}, \lean{shiftedBooleanCollision\_wrap\_bound\_of\_common\_multiple}{RationalSupportCarrySkeleton.lean:1968}).
The earlier review proposed a broader use but did not supply its additional
mean-residue estimate. It also left the endpoint proof bodies in the
following dichotomy unverified for a changed target
(\lean{half\_mem\_mersenneAchievementSet\_or\_exists\_fatal\_gap}{HalfCutLocator.lean:623}). No such extension is used here.

\paragraph{Quantifiers must be checked independently.}
For example, a statement of the form
$\exists N\,\exists L\,\forall h$ does not provide the required
$\forall h\,\forall N_0\,\exists N\ge N_0\,\exists L$.
The first assertion controls many shifts at one location; the second
requires arbitrarily late locations for each fixed shift. Relabelling a
source statement does not bridge this difference.

% ; - part a249_invent ; -
\section{Additional mathematics required by the sufficient conditions}

The preceding sections contain identities, conditional implications,
comparison sequences and finite audits. This section develops the
arithmetic estimates needed by the sufficient conditions in
Section~\ref{sec:survivors}. It separates proved reductions from proposed
applications of analytic methods; a limitation of one proposed application
is not a theorem excluding the method.

The section contains a doubling-orbit formulation with explicit thresholds,
a lacunary comparison sequence for the strength of a block estimate,
and sufficient estimates for the prime-factor decomposition. Each result
retains its stated evidence label. The discussion of a rank upper bound
must also be read with the rational $5/4$ comparison sequence: rationality
alone does not give bounded rank. These are the scope distinctions needed
to interpret the individual statements, not priority claims.


Where a claim is proved it is marked \ev{Math}; where it is checked by exact
or floating-point computation, \ev{Cert}; where it is quoted, \ev{Cited};
where it is a proposal, \ev{Open}. The proved reductions identify explicit sufficient hypotheses. Those
hypotheses remain to be established; the reductions do not prove them.

\begin{rem}
Declarations in the shared tree are cited as
\texttt{Erdos249257.name} with file and line. Declarations in the newer
per-problem tree are cited with their path from the repository root,
erdos problems/erdos249/file source:line; the hyperlink target for
those is the per-problem directory, not \texttt{Erdos249257/}.
\end{rem}

\subsection{The doubling identity}

Write $S=\sum_{n\ge 1}\varphi(n)/2^n$, $R_N=\sum_{m\ge 1}\varphi(N+m)/2^m$
for the local tail
(\lean{Erdos249257.\allowbreak TotientTailPeriodKiller.\allowbreak totientTail}{TotientTailPeriodKiller.lean:57}),
and $\Phi_N=\sum_{n\le N}\varphi(n)2^{N-n}\in\Z$ for the integer prefix, so
that $2^N S=\Phi_N+R_N$. For $h\ge 1$ set
\[
  \alpha_h \;:=\; (2^h-1)\,S .
\]

\begin{lem}[The doubling identity]\label{lem:orbit}
For all $N\ge 0$ and $h\ge 1$,
\[
  R_{N+1}=2R_N-\varphi(N+1),
  \qquad
  R_{N+h}-R_N \;=\; 2^N\alpha_h-\bigl(\Phi_{N+h}-\Phi_N\bigr),
\]
so $R_{N+h}-R_N \equiv 2^N\alpha_h \pmod 1$, and consequently the exact
first character of the tail difference is the $\times 2$ orbit of one real
number:
\[
  e(R_{N+h}-R_N)
  \;=\; e\bigl(2^N\alpha_h\bigr),
  \qquad e(x):=\exp(2\pi i x).
\]
Hence, for fixed $h$, the phases $\{\,R_{N+h}-R_N \bmod 1\,\}_{N\ge 0}$ are
the forward orbit of $\alpha_h \bmod 1$ under $x\mapsto 2x$.
\end{lem}
\leannote{Lean: \lproof{ErdosProblems/Erdos249/PaperCompleteR21/DoublingOrbitTransferAndFullDepthPhase.lean}{26}{orbit_tail_recurrence}, \lproof{ErdosProblems/Erdos249/PaperCompleteR21/DoublingOrbitTransferAndFullDepthPhase.lean}{32}{orbit_tail_diff_eq}, \lproof{ErdosProblems/Erdos249/PaperCompleteR21/DoublingOrbitTransferAndFullDepthPhase.lean}{41}{orbit_tail_diff_sub_scaled_is_int}, \lproof{ErdosProblems/Erdos249/PaperCompleteR21/DoublingOrbitTransferAndFullDepthPhase.lean}{53}{orbit_tail_diff_firstChar_eq}, and 2 further declarations in the \hyperref[sec:coverage]{coverage section}.}

\noindent
The second identity and the character form are proved:
\lean{Erdos249257.\allowbreak TotientTailPeriodKiller.\allowbreak tail_diff_eq_scaled_\allowbreak totient_series_sub_prefix}{PivotAntiReconstruction.lean:54}
and
\lean{Erdos249257.\allowbreak TotientTailPeriodKiller.\allowbreak tailOrbitFirstExp_eq_scaledTotientSeriesFirstExp}{PivotAntiReconstruction.lean:66}
\quad \ev{Lean} \quad \scale{uniform} \quad \coord{other:binary-window}.
The recurrence $R_{N+1}=2R_N-\varphi(N+1)$ is one line from the definition
and is used below only for intuition \ev{Math}.

These identities express the phases as a lacunary exponential sum in
one real variable. The next proposition transfers an estimate for that
orbit to the truncated phases, accounting for the truncation error.
The identities themselves are already in the cited sources.


\begin{prop}[Transfer to the doubling orbit]\label{prop:transfer}
Suppose that
\[
  \forall h\ge 1\ \forall X_0\ \exists X\ge \max(X_0,1):\quad
  \sum_{N=X}^{2X-1}\cos\bigl(2\pi\,2^{N}\alpha_h\bigr)\;\le\;\tfrac{89}{100}\,X .
\]
This hypothesis implies that $S$ is irrational.
\end{prop}
\leannote{Lean: \lproof{ErdosProblems/Erdos249/PaperCompleteR21/DoublingOrbitTransferAndFullDepthPhase.lean}{92}{irrational_totientSeries_of_block_cosine_gap}.}

\begin{proof}
Fix $h\ge 1$ and $X_0$, and take $X$ as supplied. Choose $L\ge h$ with
$2^{L}\ge 1024\,(2X+h+L+2)$; such $L$ exists because $2^L/L\to\infty$, and it
satisfies the room inequality $16(2X+h+L+2)\le 2^L$ a fortiori. For every
$N<2X$ the proved truncation bound
\lean{Erdos249257.\allowbreak TotientTailPeriodKiller.\allowbreak norm_windowFirstExp_\allowbreak sub_tailOrbitFirstExp_lt}{PivotAntiReconstruction.lean:202}
gives
$\bigl\|\,{E}(h,N,L)-e(2^N\alpha_h)\,\bigr\|
 < 2\pi(N+L+h+2)/2^{L} \le 2\pi/1024 < 1/100$,
so by
\lean{Erdos249257.\allowbreak TotientTailPeriodKiller.\allowbreak windowFirstCos_le_\allowbreak add_of_tailOrbitGap}{PivotAntiReconstruction.lean:220},
$\operatorname{Re}E(h,N,L)\le \cos(2\pi 2^N\alpha_h)+1/100$. Summing over
the $X$ values $N\in[X,2X)$ gives
$\sum_N \operatorname{Re}E(h,N,L)\le \tfrac{89}{100}X+\tfrac{1}{100}X
=\tfrac{9}{10}X$, so
\lean{Erdos249257.\allowbreak TotientTailPeriodKiller.\allowbreak exists_certifiedKill_of_\allowbreak first_harmonic_gap}{FirstHarmonicGap.lean:128}
produces $N\in[X,2X)$ with $\mathcal{C}(h,N,L)$, and $N\ge X\ge X_0$.
This is exactly the quantified certificate condition consumed by
\lean{Erdos249257.\allowbreak TotientTailPeriodKiller.\allowbreak irrational_totient_series_\allowbreak of_certificate_supply}{TotientTailPeriodKiller.lean:394}.
\end{proof}

\noindent
\ev{Math} \quad \scale{cofinal} \quad \coord{other:binary-window}. The proof
above makes the truncation margin explicit. The full implication is also
recorded in the formal source
\lean{ErdosProblems.Erdos249.irrational_totient_series_of_tailOrbitBlockGap}{ErdosProblems/Erdos249/TotientStrictPrimeEscape.lean:96}.
It uses a real-part estimate for the exact phases and spends $1/100$
on truncation. Compare the sufficient norm condition
\lean{Erdos249257.\allowbreak irrational_totient_series_of_\allowbreak first_harmonic_norm_gap}{FirstHarmonicPivot.lean:83}.
This is a separately sufficient condition. It is not asserted
to be equivalent to the $21/25$ norm bound or to the four separate
estimates stated above. The constants account for different
quantities and different error margins.

\begin{cor}[A digit version of the analytic condition]\label{cor:digitform}
Let $\rho_h(X)$ denote the proportion of $N\in[X,2X)$ with
$\|2^N\alpha_h\|_{\R/\Z}\ge 1/4$. If for every $h\ge 1$ there are cofinally
many $X$ with $\rho_h(X)\ge 11/100$, then $S$ is irrational.
For nondyadic $\alpha_h$, the condition counted by $\rho_h(X)$ is
exactly a change between binary digits $N+1$ and $N+2$. Thus the same
sufficient hypothesis asks for at least $11X/100$ such changes, counted
with $X\le N<2X$, on arbitrarily large blocks for every $h$.
\end{cor}
\leannote{Lean: \lproof{ErdosProblems/Erdos249/PaperCompleteR21/BinaryDigitChangeDensity.lean}{121}{irrational_totientSeries_of_quarterFarPhase_proportion}, \lproof{ErdosProblems/Erdos249/PaperCompleteR21/BinaryDigitChangeDensity.lean}{146}{quarterFarFromInt_iff_binaryDigitAt_change}, \lproof{ErdosProblems/Erdos249/PaperCompleteR21/BinaryDigitChangeDensity.lean}{233}{irrational_totientSeries_of_digitChange_count}, \lproof{ErdosProblems/Erdos249/PaperCompleteR21/BinaryDigitChangeDensity.lean}{95}{irrational_totientSeries_of_quarterFarPhase_count}, and 3 further declarations in the \hyperref[sec:coverage]{coverage section}.}

\begin{proof}
The inequality $\|x\|_{\R/\Z}\ge1/4$ gives $\cos(2\pi x)\le0$.
Every remaining term is at most $1$, so
\[
 \sum_{N=X}^{2X-1}\cos(2\pi 2^N\alpha_h)
          \le(1-\rho_h(X))X\le\frac{89}{100}X.
\]
Proposition~\ref{prop:transfer} now applies. For the digit
reading, for a nondyadic $\alpha_h$ the fractional part of $2^N\alpha_h$
is in $[1/4,3/4]$ exactly when those two binary digits differ.
The nondyadic qualification excludes the endpoint ambiguity at $3/4$.
\end{proof}

\noindent
\ev{Math}. The exact hypothesis is the digit-change inequality in the
corollary. It is not a claim about an average run length: the two boundary
runs of a finite block need not be complete, and their lengths are not
specified by the number of interior digit changes.
The finite observations below satisfy the displayed bound at their sampled
parameters. They do not establish it on arbitrarily large blocks.


\begin{table}[h]\centering
\begin{tabular}{r r r r r}
\hline
$h$ & $X$ & $\rho_h(X)$ & $X^{-1}\sum_{N\in[X,2X)}\cos(2\pi 2^N\alpha_h)$ & required \\
\hline
$1$ & $256$  & $0.4766$ & $\phantom{-}0.0037$ & $\le 0.89$ \\
$1$ & $2048$ & $0.5073$ & $-0.0043$ & $\le 0.89$ \\
$3$ & $256$  & $0.4648$ & $\phantom{-}0.0296$ & $\le 0.89$ \\
$3$ & $2048$ & $0.5020$ & $\phantom{-}0.0027$ & $\le 0.89$ \\
$8$ & $2048$ & $0.5112$ & $-0.0088$ & $\le 0.89$ \\
$13$& $2048$ & $0.5049$ & $-0.0195$ & $\le 0.89$ \\
\hline
\end{tabular}
\caption{Finite block statistics for $\alpha_h=(2^h-1)S$.
The exact digit-change counts in the six rows are $122$, $1039$, $119$,
$1028$, $1047$ and $1034$; the proportions and cosine averages are
rounded to four decimal places. The interval calculation described
below certifies these counts and displayed roundings. \ev{Cert}}
\label{tab:blocks}
\end{table}

For the numerical enclosure, compute
$P=\sum_{n=1}^{6080}\varphi(n)2^{6080-n}$. The elementary tail bound gives
\[
 P\le2^{6080}S\le P+6082.
\]
The endpoints have the same quotient on division by $2^{80}$, which
certifies $\lfloor2^{6000}S\rfloor$; an error bound smaller than one
would not alone certify a floor near an integer boundary.
Multiplying this interval by $(2^h-1)2^{N-6080}$ encloses each phase.
None of the resulting intervals crosses a digit-test boundary, and
interval evaluation of the cosine certifies the rounded averages.
Among the first $5800$ fractional digits, the numbers of runs for
$h=1,3,8$ are $2864$, $2964$, $2927$, and the longest runs have lengths
$13$, $12$, $14$. These count the two boundary runs as observed, whether
or not they continue outside the window. No randomness assumption or
asymptotic conclusion is part of these finite calculations.

\begin{obs}\label{obs:allroutes}
Lemma~\ref{lem:orbit} expresses each tail phase in terms of the doubling
orbit of $\alpha_h=(2^h-1)S$. The block conditions average those phases,
the prime condition samples prescribed prime-indexed positions, and the
full-depth condition requires a finite residue certificate for each
$(d,N)$ at some multiple of $d$. Corollary~\ref{cor:digitform} gives one
sufficient digit-change test for the real-part estimate; it is not an
equivalence with the norm bound or with the four separate budgets.
Likewise, a bound on a binary run length is not the exact prime-gap
condition. Section~\ref{sub:shape} compares the original inequalities.
\end{obs}

\subsection{The full-block exponential-sum estimate}

\paragraph{The required inequality.}
\[
  \forall h\ge 1\ \forall X_0\ \exists X,L:\quad
  \max(X_0,1)\le X,\quad 16(2X+h+L+2)\le 2^{L},\quad
  \Bigl\|\sum_{N=X}^{2X-1} e\bigl(D(h,N,L)/2^{L}\bigr)\Bigr\|\le \tfrac{21}{25}X,
\]
where $D(h,N,L)=\sum_{j<L}\bigl(\varphi(N+h+1+j)-\varphi(N+1+j)\bigr)2^{L-1-j}$.
\lean{Erdos249257.\allowbreak TotientTailPeriodKiller.\allowbreak DTWFirstHarmonicNormGap}{FirstHarmonicPivot.lean:75},
implication
\lean{Erdos249257.\allowbreak irrational_totient_series_of_\allowbreak first_harmonic_norm_gap}{FirstHarmonicPivot.lean:83}
\quad \ev{Open} \quad \scale{cofinal}. By Proposition~\ref{prop:transfer} the
real-part bound with constant $89/100$ for the \emph{exact doubling
phases} suffices; its transfer to the truncated phases uses the separate
error estimate in that proposition.

\paragraph{The discrepancy in terms of the doubling map.}
For $L\ge h$, dividing by $2^L$ and reindexing the two ranges gives
\begin{equation}\label{eq:window}
  \frac{D(h,N,L)}{2^{L}}
  =\underbrace{-\sum_{t=1}^{h}\frac{\varphi(N+t)}{2^{t}}}_{\text{head}}
  \;+\;\underbrace{(2^{h}-1)\sum_{t=h+1}^{L}\frac{\varphi(N+t)}{2^{t}}}_{\text{body}}
  \;+\;\underbrace{\sum_{t=L+1}^{L+h}\frac{\varphi(N+t)}{2^{t-h}}}_{\text{tail}} .
\end{equation}

The phase is a geometrically weighted sum of totients at $L+h$
consecutive arguments. In the first two ranges, $1\le t\le L$, the
coefficient has denominator $2^t$ and odd numerator. Its contribution is
therefore integral exactly when $v_2(\varphi(N+t))\ge t$. In the last range,
$L<t\le L+h$, the denominator is instead $2^{t-h}$, so the corresponding
threshold is $t-h$. Here $v_2(a)$ denotes the exponent of $2$ in the
positive integer $a$. The selected offset in the four-term decomposition
lies in the middle range, where the odd multiplier is $2^h-1$.

\paragraph{Relevant comparison estimates.}
The following comparisons identify the estimates still needed; they do
not rule out these methods.

\emph{Weyl differencing and van der Corput.} The van der Corput inequality
applies to arbitrary bounded complex sequences~\cite[Theorem~2.1]{edeko-vdc},
but its use here requires estimates for shifted correlations. For the exact
phases $x_N=2^N\alpha_h$ of Lemma~\ref{lem:orbit},
\[
 x_{N+q}-x_N=(2^q-1)x_N.
\]
Thus differencing gives another doubling orbit, not a polynomial phase of
lower degree. The finite phase $D(h,N,L)/2^L$ approximates this orbit modulo
$1$; it is not the exact orbit itself. No bound for the resulting
correlations is established here.

\emph{The large sieve.} A proposed application would need a suitable
family of separated frequencies and a way to recover a bound for each
fixed $h$. Averaging only over $h$ would not suffice: the condition asks
for a block estimate for every $h$, not for most of them.

\emph{Bilinear sums over prime factors.} Factoring $N+t=mp$ and using
$\varphi(mp)=\varphi(m)(p-1)$ makes the selected totient contribution linear
in the prime $p$. Section~\ref{sub:pivot} gives this decomposition and
identifies the remaining weight. A bound for that weight is still needed
before a Vaughan- or Vinogradov-type argument can be applied.

\emph{Erd\H{o}s's 1948 digit method.} The comparison series is
$E=\sum_{n\ge1}1/(2^n-1)=\sum_{N\ge1}d(N)/2^N$, whose irrationality
Erd\H{o}s proved~\cite{erdos1948lambert}. For the totient series, the
corresponding M\"obius expansion is exact: $\varphi=\mu * \mathrm{id}$ gives
\coord{mobius-mersenne}
\begin{equation}\label{eq:mobmers}
  S=\sum_{d\ge 1}\mu(d)\,\frac{2^{d}}{(2^{d}-1)^{2}}
\end{equation}
The identity is labelled \ev{Math}; the recorded numerical check gives
$S=1.3676308019850223\ldots$ to sixteen digits \ev{Cert}.

Erd\H{o}s combines congruences forcing increasing powers of $2$ to divide
a block of divisor counts with an averaged tail bound, producing long
zero blocks in the binary expansion~\cite[pp.~63--66]{erdos1948lambert}.
To adapt that argument here, one would need congruences for the actual
totient coefficients together with a compatible tail estimate. The
identities alone provide neither combination. In \eqref{eq:mobmers},
the coefficients $\mu(d)$ can have either sign, so positivity of the
original totient series does not give positivity term by term in this
expansion. This is a limitation of the proposed direct transfer, not a
proof that a digit argument or a different use of the M\"obius expansion
cannot work.

\paragraph{The unresolved estimate.}
For general coefficients $0\le c(n)\le n$, irrationality alone does not
imply the block norm condition. The following example has very sparse
nonzero coefficients. It is irrational, but almost all phases in each
sufficiently large block lie close to $1$.

\begin{thm}[The block norm condition is stronger in the class $0\le c(n)\le n$]\label{thm:lacunary}
Let $c(n)=1$ if $n=k!$ for some $k\ge 1$ and $c(n)=0$ otherwise, so
$0\le c(n)\le n$ for all $n\ge 1$, and let $\beta=\sum_{n\ge1}c(n)/2^{n}
=\sum_{k\ge 1}2^{-k!}$. Then $\beta$ is irrational, and for every $h\ge 1$
and every $X\ge 81(h+5)$,
\[
  \sum_{N=X}^{2X-1}\cos\bigl(2\pi\,2^{N}(2^{h}-1)\beta\bigr) \;>\; \tfrac{9}{10}X .
\]
Consequently the block norm condition fails for $\beta$ at every
sufficiently large scale, for every admissible truncation depth, even though
$\beta$ is irrational. Thus irrationality, nonnegative integer coefficients
and the bound $c(n)\le n$ do not imply the block gap. A proof for $S$
needs an additional property not shared by this example; the comparison
does not specify which additional property will suffice.
\end{thm}
\leannote{Lean: \lproof{ErdosProblems/Erdos249/PaperCompleteR21/LacunaryFactorialBlockNorm.lean}{68}{lacCoef_bounds}, \lproof{ErdosProblems/Erdos249/PaperCompleteR21/LacunaryFactorialBlockNorm.lean}{88}{lacBeta_eq_factorial_series}, \lproof{ErdosProblems/Erdos249/PaperCompleteR21/LacunaryFactorialBlockNorm.lean}{121}{irrational_lacBeta}, \lproof{ErdosProblems/Erdos249/PaperCompleteR21/LacunaryFactorialBlockNorm.lean}{333}{lacunary_block_cos_gap}, and 2 further declarations in the \hyperref[sec:coverage]{coverage section}.}

\begin{proof}
The binary expansion of $\beta$ has infinitely many ones separated by
unbounded strings of zeros. It is not eventually periodic, so $\beta$ is
irrational. Write $\gamma=(2^h-1)\beta$, and let $k!$ be the least
factorial greater than $N$. The earlier terms become integers after
multiplication by $2^N$, so
\[
 2^N\gamma\equiv(2^h-1)\sum_{j\ge k}2^{N-j!}\pmod{\mathbb Z}.
\]
If $N\le k!-h-5$, the expression on the right lies between $0$ and
$2\cdot2^{N+h-k!}\le1/16$. Thus $\|2^N\gamma\|_{\mathbb R/\mathbb Z}\le1/16$ and
$\cos(2\pi 2^{N}\gamma)\ge\cos(\pi/8)>9238/10000$. The remaining $N$
lie in $(k!-h-5,\,k!]$, at most $h+5$ values per factorial. The relevant
factorials have ratio at least $3$, and $h+5<X$, so $[X,2X)$ meets at
most one such interval. Hence the sum is greater than
\[
 \frac{9238}{10000}(X-h-5)-(h+5)\ge\frac9{10}X+\frac{h+5}{250}
\]
when $X\ge81(h+5)$.

To pass from these exact phases to the finite residue test, note that
$0\le c(n)\le1$ puts each scaled tail in $[0,1]$. The error in the
truncated tail difference is therefore at most $2^{-L}$, so each complex
phase changes by at most $2\pi2^{-L}$. At an admissible depth,
$2^L\ge16(2X+h+L+2)>32X$, and the total change in the block sum is less
than $\pi/16$. Its real part is consequently greater than
$9X/10-\pi/16>21X/25$ for $X\ge81(h+5)$. Its norm cannot satisfy the
$21X/25$ bound either.
\end{proof}

\noindent
\ev{Math} \quad \scale{cofinal} \quad \coord{other:binary-window}.
This example rules out the implication from irrationality and the stated
coefficient bound to the block norm condition. It does not say that the
totient sum is Liouville-like, or that every argument based on its digits
fails. A proof of the block estimate must use information that distinguishes
the actual totient phases from this comparison sequence.

\paragraph{The size of the required saving.}
Under Lebesgue measure on the phase circle, the binary digit maps are
independent fair bits. The digit-change proportion therefore tends to
$1/2$ for almost every phase. For almost every phase, the real-part block
sum is also $O(\sqrt{X\log\log X})$, by the law of the iterated logarithm of
Erd\H{o}s and G\'al for lacunary series~\cite[p.~65]{erdos-gal}.
The block bound follows by subtracting two initial partial sums.
An almost-everywhere theorem supplies no conclusion at the particular
phase $\alpha_h=(2^h-1)S$ \ev{Cited}. The requirement is $\rho_h\ge 0.11$ and a saving of a constant
factor. The finite values in Table~\ref{tab:blocks} have $\rho_h$ near $0.50$
and small block averages. They are compatible with square-root
cancellation, but do not establish its asymptotic order. The sufficient
condition asks only for the stated fixed saving on arbitrarily large
blocks. No estimate in this record supplies it for the totient phases.

\subsection{The four-term decomposition}\label{sub:pivot}

\paragraph{The simultaneous inequalities.}
Use the four finite sums defined in Section~\ref{sub:four-sums}; in
particular, each group mean averages the selected prime phase, not the
residual weight. For every $h\ge 1$ there must exist $s\ge 1$ and $\eta\in(0,1)$ such that for
every $X_0$ there are $X\ge\max(X_0,1)$ and $L$ with $h\le L-s$,
$16(2X+h+L+2)\le 2^{L}$, and all four of
\[
\begin{aligned}
 \operatorname{Re}\sum_{N\in\mathcal G}w_N(z_N-\bar z_{m_N})
   &\le\tfrac{14}{25}X,\\
 \left\lVert\sum_{N\in\mathcal G}w_N\bar z_{m_N}\right\rVert
   &\le\tfrac{1}{100}X,\\
 \left\lVert\sum_{N\in\mathcal A\smallsetminus\mathcal G}E(h,N,L)\right\rVert
   &\le\tfrac{1}{100}X,\\
 \left\lVert\sum_{\substack{X\le N<2X\\N\notin\mathcal A}}E(h,N,L)\right\rVert
   &\le\tfrac{8}{25}X.
\end{aligned}
\]
\lean{Erdos249257.\allowbreak TotientTailPeriodKiller.\allowbreak DTWPivotResidualDecorrelation}{FirstHarmonicPivot.lean:569},
implication
\lean{Erdos249257.\allowbreak irrational_totient_\allowbreak series_of_pivotResidualDecorrelation}{FirstHarmonicPivot.lean:577}
\quad \ev{Open} \quad \scale{cofinal}. The decomposition is an exact identity,
\lean{Erdos249257.\allowbreak TotientTailPeriodKiller.\allowbreak windowFirstExp_sum_\allowbreak eq_pivot_decomposition}{FirstHarmonicPivot.lean:514}
\quad \ev{Lean}. For $X\ge4$ and $s=L-h$, the single cofactor group
$m_N=1$, not the full assigned set, consists of the indices for which
$N+h+1$ is prime,
\lean{Erdos249257.\allowbreak TotientTailPeriodKiller.\allowbreak mem_pivotFiber_\allowbreak one_overlap_iff}{FirstHarmonicPivot.lean:423}
\quad \ev{Lean}.

\paragraph{An upper bound for the unassigned terms.}
An unassigned index has a smooth shifted argument, but the converse is
not asserted. The factorisation gives the containment needed for an upper
bound on the number of unassigned indices.

\begin{prop}[A one-sided bound for the unassigned terms]\label{prop:dickman}
Fix $h,s$ and choose the admissible depth $L$ minimally for each large $X$.
Put $t=L-s+1=O_{h,s}(\log X)$ and $y_X=4\sqrt X+2t/\sqrt X$.
If $n=N+t$ is unassigned, then its largest prime factor satisfies
$P(n)\le y_X$. Consequently
\[
\begin{aligned}
 \#\{N\in[X,2X):N\notin\mathcal A\}
 &\le \Psi(2X+t-1,y_X)-\Psi(X+t-1,y_X)\\
 &=\bigl(1-\log2+o(1)\bigr)X<\tfrac8{25}X
\end{aligned}
\]
for all sufficiently large $X$. Here $\Psi(x,y)$ counts the positive
integers at most $x$ whose prime factors are all at most $y$.
This is an upper bound for the unassigned count, not an asymptotic equality
for that count.
\end{prop}
\leannote{Not formalised: the bound on the largest prime factor at an unassigned index is checked, and the passage from it to the count of smooth integers in the two intervals, with the asymptotic evaluating that count, is not. See the \hyperref[sec:coverage]{coverage section}.}

\begin{proof}
Let $p=P(n)$ and $m=n/p$. Since $n>1$, failure of the assignment condition
means either $p\le2\lfloor\sqrt X\rfloor$ or
$m>\lfloor\sqrt X/2\rfloor$. In the second case $m>\sqrt X/2$, so
\[
 p=\frac{n}{m}<\frac{2(2X+t)}{\sqrt X}=y_X.
\]
The first case also gives $p\le y_X$. Counting the resulting smooth
integers in $[X+t,2X+t)$ proves the first inequality.

For completeness, the classical fixed-parameter asymptotic
$\Psi(x,x^{1/u})\sim x\rho(u)$ suffices; no additional short-interval
estimate is needed. Indeed, for either $x=X+t-1$ or $x=2X+t-1$,
$\log y_X/\log x\to1/2$. For every fixed $0<\epsilon<1$, eventually
\[
 x^{1/(2+\epsilon)}\le y_X\le x^{1/(2-\epsilon)}.
\]
Monotonicity in $y$ bounds $\Psi(x,y_X)/x$ between quantities tending to
$\rho(2+\epsilon)$ and $\rho(2-\epsilon)$. Letting $\epsilon\to0$ and
using continuity gives the limit $\rho(2)$. Subtracting the two endpoint
asymptotics yields $(\rho(2)+o(1))X$. Since $\rho(2)=1-\log2<8/25$,
the claimed bound follows.
\end{proof}

\noindent
\ev{Lean} for the whole proposition,
\lean{ErdosProblems.Erdos249.PaperCompleteR21.prop\_dickman}{lean/ErdosProblems/Erdos249/PaperCompleteR21/UnassignedSmoothCount.lean:556}.
The Lean proof uses no smooth-number asymptotic. For large $X$,
$y_X^2>2X+t-1$, so an integer in $(X+t-1,2X+t-1]$ that is not
$y_X$-smooth has exactly one prime factor above $y_X$, and
\[
 \Psi(2X+t-1,y_X)-\Psi(X+t-1,y_X)
 =X-\sum_{y_X<p\le2X+t-1}
 \left(\left\lfloor\frac{2X+t-1}{p}\right\rfloor
      -\left\lfloor\frac{X+t-1}{p}\right\rfloor\right).
\]
Mertens' theorems, proved in the development with explicit constants,
then give $(1-\log2)X+O(X/\log X)$. The written proof above
cites instead the fixed-$u$
asymptotic~\cite[(1.1), p.~268]{granville-smooth} and the value
$\rho(u)=1-\log u$ on $[1,2]$~\cite[(1.3), p.~268]{granville-smooth}.
No effective crossover scale or exact density for the unassigned indices
is asserted. The estimate uses the minimal admissible depth; it does not
supply either of the two remaining correlation bounds at that depth.

\begin{prop}[The excluded-cofactor estimate]\label{prop:badcof}
Fix $h,s$ and use the minimal admissible depth $L$, as in
Proposition~\ref{prop:dickman}; thus $t=L-s+1=O_{h,s}(\log X)$.
For $\eta\in(0,1)$ let $B(\eta)=\{m\ge1:\varphi(m)<\eta m\}$, with natural density
$D(\eta)$; by Schoenberg's theorem $D$ exists, is continuous, and
$D(0+)=0$~\cite[\S17, p.~193]{schoenberg1928}, in the framework
of~\cite[Theorem~1, pp.~318--319, and \S8, p.~323]{schoenberg1936} \ev{Cited}. Then
\[
  \#\{N\in\mathcal A:m_N\in B(\eta)\}
  \;\le\;\bigl(D(\eta)+o(1)\bigr)X ,
\]
so a single choice of $\eta$ with $D(\eta)<1/200$ meets the $\tfrac{1}{100}X$
budget for all large $X$. This choice fixes $\eta$ before $X_0$, as
required. It supplies only the excluded-cofactor bound: the mean and
mean-subtracted estimates must still hold for this same $\eta$, and do
not follow from making $\eta$ smaller.
\end{prop}
\leannote{Not formalised: the counting function of indices with excluded cofactor is carried by the development, and neither Schoenberg's density theorem nor the partial summation over the cofactors is formalised. See the \hyperref[sec:coverage]{coverage section}.}

\begin{proof}
For a fixed cofactor $m\le\sqrt X/2$, each assigned index gives a prime
$p$ with $X+t\le mp<2X+t$. Dropping the other assignment restrictions
bounds the number of indices by the number of primes in this interval.
The prime number theorem, uniformly for $X/m\ge2\sqrt X$, and
$t=O_{h,s}(\log X)$ give the upper bound
\[
 (2+o(1))\frac{X}{m\log X}.
\]
The possible endpoint contribution is $O(1)$ for each $m$, hence $o(X)$
after summing over $m\le\sqrt X/2$. Natural density and partial summation
give
\[
 \sum_{\substack{m\le\sqrt X/2\\m\in B(\eta)}}\frac1m
       =\left(\frac{D(\eta)}2+o(1)\right)\log X.
\]
Multiplying the two estimates proves the stated bound. If
$D(\eta)<1/200$, its right side is less than $X/100$ for all sufficiently
large $X$.
\end{proof}

\noindent
The existence of a usable cutoff can also be made explicit by an
elementary estimate. For $x\ge1$,
\[
\begin{aligned}
 \frac1x\sum_{n\le x}\frac n{\varphi(n)}
 &\le\sum_{d\ge1}\frac{\mu(d)^2}{d\varphi(d)}
   =\prod_p\left(1+\frac1{p(p-1)}\right)\\
 &\le\exp\left(\sum_{n\ge2}\frac1{n(n-1)}\right)
   =\exp(1)<3.
\end{aligned}
\]
The first step uses
$n/\varphi(n)=\sum_{d\mid n}\mu(d)^2/\varphi(d)$ and
$\lfloor x/d\rfloor\le x/d$; the Euler product has nonnegative terms
and converges by the displayed majorant. Markov's inequality gives
$\#(B(\eta)\cap[1,x])<3\eta x$, hence $D(\eta)\le3\eta$.
In particular, $\eta=1/1000$ gives
$D(\eta)\le3/1000<1/200$, an explicit sufficient choice for the
excluded-cofactor budget. This does not establish either weighted phase
estimate for that choice, and it is not claimed to be an optimal cutoff.
\ev{Math}

The first element of $B(1/5)$ is
$30030=2\cdot3\cdot5\cdot7\cdot11\cdot13$, and
$\varphi(30030)/30030=5760/30030\approx0.191808$. Enumeration through
$3\cdot10^5$ gives $38$ elements of $B(1/5)$ and
\[
 \frac{1}{\log(3\cdot10^5)}
 \sum_{\substack{m\le3\cdot10^5\\m\in B(1/5)}}\frac1m
 \approx2.82\cdot10^{-5}.
\]
This finite weighted proportion is not the limiting density $D(1/5)$
\ev{Cert}. Using $\eta=1/5$ in the asymptotic bound would require a
rigorous bound for $D(1/5)$, including the integers beyond the enumeration;
the finite count does not supply one.

\paragraph{The growing dyadic modulus.}
The selected offset is $t=L-s+1$ in \eqref{eq:window}. The four-term
condition fixes $s$ after $h$ but before the threshold $X_0$, whereas
\[
 2^L\ge16(2X+h+L+2)>32X,\qquad 2^t>2^{6-s}X.
\]
On a group with cofactor $m$, the coefficient of $p-1$ in the phase
has reduced denominator $2^{\max(t-v_2(\varphi(m)),0)}$, because $2^h-1$
is odd. Assignment gives
$\varphi(m)\le m\le\sqrt X/2$. Thus, for fixed $s$ and all sufficiently
large $X$, the reduced denominator is greater than
$2^{7-s}\sqrt X$. In particular, the assignment conditions do not reduce
these moduli to a fixed power of $\log X$. This is the modulus for
variation in $p$, not a lower bound for the denominator of the phase
value at each individual prime.

Choosing $t=O(\log\log X)$ would require $s$ to grow with $X$, contrary
to the stated quantifiers. Likewise, the formal observation that the phase
is $1$ when $v_2(\varphi(m))\ge t$ cannot provide such groups at large $X$
under these size bounds. The proposed small-modulus prime-distribution
argument therefore does not establish the mean estimate. This does not
disprove the four-term condition or exclude other estimates for the same
prime sums. \ev{Math} \scale{cofinal}.

\paragraph{The remaining prime-weighted correlation.}
Propositions~\ref{prop:dickman} and \ref{prop:badcof} give the unassigned
and excluded-cofactor bounds at minimal admissible depth. They must be
combined with mean and correlation estimates at the same $X,L,s,\eta$;
bounds obtained at unrelated depths cannot be substituted.

With the notation of Section~\ref{sub:four-sums}, the correlation estimate
still required is
\[
 \operatorname{Re}\sum_{N\in\mathcal G}w_N
 \left(e\!\left(\frac{(2^h-1)\varphi(m_N)(p_N-1)}{2^t}\right)
       -\bar z_{m_N}\right)\le\frac{14}{25}X.
\]
The sum is over the assigned indices with good cofactors, not all primes
and cofactors in the interval. By \eqref{eq:window}, $w_N$ is the phase
with the selected offset $t$ removed.

Regrouping by $m_N=m$ gives weighted sums over the primes $p_N=p$.
A bilinear argument would need to separate the $m$- and $p$-dependence or
estimate the resulting two-variable weight directly. For $N=mp-t$,
that weight depends on $\varphi(mp+j-t)$ at the other offsets
$1\le j\le L+h$, $j\ne t$. The identity
$\varphi(mp)=\varphi(m)(p-1)$ controls the selected value, not these
neighbouring totients. No bound for the resulting correlation is
established here; this is the missing step, not a limitation of every
bilinear method.

A fixed-offset example illustrates the correlation at issue. At $h=1$
and offset $t=2$, a prime $p\ge5$ satisfies $\varphi(p-1)\equiv0\pmod2$
and $\varphi(p)=p-1$. The phase of $R_{p-2}$ is therefore
$\chi_{-4}(p)$ times the phase of the following weighted totients, since
$e((p-1)/4)=i^{p-1}=\chi_{-4}(p)$. One possible fixed-saving estimate,
with a fixed $\delta>0$, would be
\[
  \operatorname{Re}\sum_{X<p\le 2X}\chi_{-4}(p)\,
    e\Bigl(\frac{\varphi(p+1)}{8}+\frac{\varphi(p+2)}{16}+\cdots\Bigr)
  \;\le\;\Bigl(\frac{9}{10}-\delta\Bigr)\,\bigl(\pi(2X)-\pi(X)\bigr).
\]

This fixed-offset example is not Definition~\ref{defn:pivot}: there $s$
is fixed while $t=L-s+1$ grows. The example isolates a correlation between
a quadratic character at $p$ and the binary phases of totients at
$p+1,p+2,\ldots$. It is analogous in this respect to shifted correlations
studied in the Chowla and Elliott problems, but the analogy does not
supply an applicable theorem. The cited theorems have different hypotheses and types of average: Matom\"aki and Radziwi\l{}\l{} compare short and long averages of
real multiplicative functions bounded by $1$, outside an exceptional set of
short intervals~\cite[Theorem~1]{matomaki-radziwill}; Tao proves logarithmically averaged two-point Chowla and a corrected
Elliott statement, the latter under an explicit non-pretentiousness
hypothesis on one of the bounded multiplicative functions
~\cite[Theorems~1.2--1.3]{tao-elliott}; and Tao and Ter\"av\"ainen
prove the odd-order logarithmically averaged Chowla
conjecture~\cite[Theorem~1.1]{tao-teravainen} \ev{Cited}. The phase here
depends on $\varphi$ at a growing window of shifts, so none of these applies
directly.

\paragraph{Averaging over the block parameter.}
For each fixed $h$, an average of the nonnegative, normalised block norms
can establish the full-block condition. It is enough to prove
\[
 \sum_{k=1}^{K}2^{-k}
 \left|\sum_{N=2^k}^{2^{k+1}-1}
       e\bigl(D(h,N,L_k)/2^{L_k}\bigr)\right|=o(K),
\]
where every $L_k$ satisfies the room inequality for $X=2^k$.
Indeed, the number of $k\le K$ for which the normalised norm exceeds
$21/25$ is at most $25/21$ times the displayed sum, hence is $o(K)$.
There are therefore arbitrarily large admissible blocks satisfying
Definition~\ref{defn:fh}. This argument does not establish the four-term
condition: that condition requires all four separate budgets with common
$X,L,s,\eta$.

The displayed average is a Ces\`aro average of norms on dyadic blocks.
It is not the signed logarithmically weighted correlation appearing in
the cited Chowla--Elliott theorems. The position of the absolute value
matters: cancellation between different blocks would not bound this
sum of norms. Invoking entropy decrement would therefore require an
additional argument yielding this particular estimate.

The finite phase depends on the $L+h$ totient values at offsets
$1,\ldots,L+h$, with the coefficients in \eqref{eq:window}.
Although its absolute value is $1$, it is not a multiplicative function
of $N$. Both the number of shifts and their dyadic denominators grow with
$L$. The cited fixed-shift correlation theorems do not supply the uniform
bounds for this growing family. No such extension is proved here. \ev{Open}

\subsection{Separation along prime positions}

\paragraph{The quantified prime condition.}
\[
  \forall h\ge 1\ \forall N_0\ \exists p\ \text{prime}:\quad
  \max(N_0+h+1,\,h+5)\le p
  \ \wedge\
  \operatorname{Re}\bigl(e(R_{p-1}-R_{p-h-1})\bigr)<\tfrac{9}{10}.
\]
\lean{ErdosProblems.\allowbreak Erdos249.\allowbreak DTWNaturalPrimeTailOrbitStrictGap}{ErdosProblems/Erdos249/TotientStrictPrimeEscape.lean:157},
implication
\lean{ErdosProblems.\allowbreak Erdos249.\allowbreak irrational_totient_\allowbreak series_of_naturalPrimeTailOrbitStrictGap}{ErdosProblems/Erdos249/TotientStrictPrimeEscape.lean:524}
\quad \ev{Open} \quad \scale{cofinal}.

\paragraph{Its meaning for the doubling map.}
$\operatorname{Re}e(2^{p-h-1}\alpha_h)<9/10$ says
$\|2^{p-h-1}\alpha_h\|_{\R/\Z}>\arccos(9/10)/2\pi=0.0717831\ldots$, so the
requirement is exactly: \emph{for each $h$, the indicated prime-indexed
$\times2$ orbit point stays more than $0.0717831\ldots$ from the nearest
integer, for cofinally many primes $p$}.  A short-run condition on the binary
digits can be a sufficient proxy for this circle-distance inequality, but it
is not equivalent to it; a run-length description alone does not determine
the residual position inside the dyadic cylinder.

\paragraph{What the prime condition adds.}
The preceding condition asks for a separated point at a prime-indexed
position. The full-block average does not by itself locate such a point
on a prescribed sparse set. Conversely, the prime condition does not
assert cancellation over the whole interval. They are different
sufficient hypotheses; no general comparison of their difficulty is
proved here.

\paragraph{Finite comparison.}
Among the $501$ primes in $[1000,5000)$, the strict inequality holds
at $432$, $426$ and $430$ primes for $h=1,2,5$, respectively. The
corresponding proportions are $0.8623$, $0.8503$ and $0.8583$ to four
decimal places, and each sample has a cosine value below $-0.9999$.
These counts were recomputed using the same interval for $S$ as in
Table~\ref{tab:blocks}, with no interval meeting the threshold $9/10$.
They do not establish separation at arbitrarily large primes.

\paragraph{An average over a prescribed prime sample.}
The subset implication
\lean{Erdos249257.\allowbreak TotientTailPeriodKiller.\allowbreak exists_certifiedKill_of_\allowbreak first_harmonic_gap_subset}{FirstHarmonicGap.lean:104}
allows a nonempty finite set $T$. If $T$ is chosen after a separated
point is known, the subset formulation may add no arithmetic content:
the equivalence proved using selected pairs illustrates this issue
\lean{Erdos249257.\allowbreak dtwWindowSeparatedPairs_iff_\allowbreak irrational_totient_series}{PivotAntiReconstruction.lean:1765}.
To avoid that particular choice, fix the sample in advance by
\[
 T_{h,X}=\{N\in[X,2X):N+h+1\text{ is prime}\}.
\]
For $X\ge4$ and $L\ge h$, this is precisely the prime-factor sample
$\mathrm{pivotFiber}(X,L,L-h,1)$ in the supplied source
\lean{Erdos249257.\allowbreak TotientTailPeriodKiller.\allowbreak mem_pivotFiber_\allowbreak one_overlap_iff}{FirstHarmonicPivot.lean:423}.

A sufficient, unproved condition is
\begin{equation}\label{eq:primefibre}
\begin{aligned}
 &\forall h\ge1\ \forall X_0\ \exists X\ge\max(X_0,4)\ \exists L\ge h:\\
 &\qquad T_{h,X}\ne\varnothing,\qquad
          16(2X+h+L+2)\le2^L,\\
 &\qquad \sum_{N\in T_{h,X}}\operatorname{Re}E(h,N,L)
       \le\frac9{10}|T_{h,X}|.
\end{aligned}
\end{equation}
The subset theorem then gives a certificate at some $N\ge X\ge X_0$.
Nonemptiness matters: for the empty set the displayed sum inequality
would hold automatically and would produce no certificate. The bounds
$X\ge4$ and $L\ge h$ are also the hypotheses of the identification with
the source's prime-factor sample. \ev{Open}

On this sample the selected cofactor is $1$, so its coefficient is the
odd integer $2^h-1$ and its modulus is $2^{h+1}$, fixed once $h$ is fixed.
The unweighted prime average has zero Ramanujan main term, since
$c_{2^{h+1}}(2^h-1)=\mu(2^{h+1})=0$. That observation does not establish
\eqref{eq:primefibre}: the prime phase is multiplied by a residual phase
built from the other totient values. Their correlation is the missing
estimate. Restricting the sample removes the complement and unwanted
cofactor terms, but not this residual correlation.

\begin{rem}
Fixing the sample rules out the particular argument that selects a pair
only after locating separation. It does not prove that
\eqref{eq:primefibre} is inequivalent to irrationality, or that it is a
strictly stronger condition. No such comparison is established here.
\end{rem}

\subsection{Depth equal to a period multiple}

\paragraph{The quantified residue condition.}
\[
  \text{full-depth condition} :\equiv
  \forall d\ge 1\ \forall N\ \exists t\ge 1:\ \mathcal{C}(td,\,N,\,td),
\]
unpacked, $(N+2td+2:\Z)<D(td,N,td)\bmod 2^{td}<2^{td}-(N+2td+2)$.
\lean{ErdosProblems.\allowbreak Erdos249.\allowbreak PeriodMultipleEscape.\allowbreak ApFullDepthEscape}{ErdosProblems/Erdos249/PeriodMultipleEscape.lean:441},
implication
\lean{ErdosProblems.\allowbreak Erdos249.\allowbreak PeriodMultipleEscape.\allowbreak irrational_totient_\allowbreak series_of_apFullDepthEscape}{ErdosProblems/Erdos249/PeriodMultipleEscape.lean:450},
ambient equivalence
\lean{ErdosProblems.\allowbreak Erdos249.\allowbreak PeriodMultipleEscape.\allowbreak periodMultipleKillSupply_iff_irrational}{ErdosProblems/Erdos249/PeriodMultipleEscape.lean:432},
equivalence
\lean{ErdosProblems.\allowbreak Erdos249.\allowbreak FullDepthRayAmplifier.\allowbreak apFullDepthEscape_iff_irrational}{ErdosProblems/Erdos249/FullDepthRayAmplifier.lean:406}
\quad \ev{Open} \quad \scale{cofinal}. The depth-locked condition is
equivalent to irrationality; the remaining gap is a nonintegral seed for every
ray and basepoint, not an unknown converse.

\paragraph{The condition in terms of binary digits.}
Writing $h=td$ and using
$D(h,N,L)/2^{L}=(R_{N+h}-R_N)-(R_{N+L+h}-R_{N+L})/2^{L}$ at $L=h$, together
with the proved strip $|R_{M+h}-R_M|<M+h+2$
(\lean{Erdos249257.\allowbreak TotientTailPeriodKiller.\allowbreak abs_tail_\allowbreak diff_lt}{CarrySurvivorExtinction.lean:79},
\ev{Lean}), applied at $M=N+h$, one gets:

\begin{prop}\label{prop:route4}
$\mathcal{C}(h,N,h)$ holds whenever
$\bigl\|2^{N+h}S-2^{N}S\bigr\|_{\R/\Z}>2(N+2h+2)/2^{h}$.
\end{prop}
\leannote{Lean: \lproof{ErdosProblems/Erdos249/PaperCompleteR21/DoublingOrbitTransferAndFullDepthPhase.lean}{130}{certifiedKill_of_fullDepth_phase_separation}, \lproof{ErdosProblems/Erdos249/PaperCompleteR21/DoublingOrbitTransferAndFullDepthPhase.lean}{119}{bracket_of_two_sided_separation}.}

\begin{proof}
$\|D(h,N,h)/2^{h}\|\ge\|R_{N+h}-R_N\|-|R_{N+2h}-R_{N+h}|/2^{h}
 >2(N+2h+2)/2^{h}-(N+2h+2)/2^{h}=(N+2h+2)/2^{h}$, and
$R_{N+h}-R_N\equiv 2^{N}(2^{h}-1)S=2^{N+h}S-2^{N}S\pmod 1$ by
Lemma~\ref{lem:orbit}. A residue at distance more than $(N+2h+2)/2^{h}$ from
$\Z$ is exactly a residue certificate at depth $h$.
\end{proof}

\noindent
\ev{Math}. Proposition~\ref{prop:route4} gives a sufficient analytic
criterion for Route~4: for each $d\ge1$ and $N\ge0$, it is enough to find
$h=td$, $t\ge1$, for which the displayed distance exceeds the stated
threshold. The criterion is not asserted to be equivalent to each individual
finite certificate. It involves the phase of $S$ itself, and the threshold
tends to zero exponentially along a fixed ray. The earlier full-depth
amplifier explains why a nonintegral seed produces late certificates; it
does not supply seeds for all $d,N$.

\paragraph{Comparison with digit complexity.}
The phase inequality above is a quantitative statement at specified
positions and period multiples. A general block-complexity assertion is a
different quantified statement, of the kind treated by the
Adamczewski--Bugeaud method for algebraic numbers
\cite[Theorem~1, p.~549]{adamczewski-bugeaud}. No implication from such an
assertion to the required inequality is established here. In particular, this record
does not justify describing one existing method as the only possible route,
or algebraicity as the sole missing input. The arithmetic nonintegrality
obligation must be stated and proved in its own right.

\paragraph{Finite examples.}
At $N=300$, the least $t$ with $\mathcal{C}(td,300,td)$ exists
and is small for every ray $d\le 24$: $t=10$ for $d=1$, $t=5$ for $d=2$,
$t=4$ for $d=3$, $t=3$ for $d=4$, $t=2$ for $5\le d\le 9$, and $t=1$ for
$10\le d\le 24$; in every case at or near the first depth admitted by the
room floor. The residues are not marginal: $r/2^{h}=0.3594$ at $d=1$,
$0.5977$ at $d=3$, $0.6853$ at $d=7$, $0.2620$ at $d=17$, $0.5887$ at
$d=23$, against edge radii $3.1\cdot10^{-1}$, $8.0\cdot10^{-2}$,
$2.0\cdot10^{-2}$, $2.6\cdot10^{-3}$, $4.1\cdot10^{-5}$ \ev{Cert}. This is a
finite verification at the displayed basepoint and shifts. It neither
bounds the required multiplier for arbitrary $d,N$ nor establishes the
universally quantified condition.

\paragraph{A possible heuristic, not an estimate.}
If a phase $U$ were uniform on $\R/\Z$, then
$\Pr(\|U\|\le\epsilon)=\min(1,2\epsilon)$ for $\epsilon\ge0$.
At $\epsilon=2(N+2h+2)2^{-h}$ this gives the marginal model
$\min(1,4(N+2h+2)2^{-h})$. The phases
$2^N(2^{td}-1)S$ on a fixed ray are arithmetically dependent; no independence
model or uniform distribution theorem for this orbit is proved here.
Consequently the marginal estimate cannot be multiplied across depths, and
no summability conclusion over all basepoints follows from it. The finite
experiments remain finite evidence, not a proof of the all-ray supply.
\ev{Open} (heuristic only).

\subsection{Why rationality does not bound carry rank}
\label{sub:generic-carry-rank}

For an integer sequence $u$, write $\operatorname{rk}_e(u)$ for the
rational dimension of the span of all sequences
$n\mapsto u_{2^j n+r}$ with $0\le j\le e$ and $0\le r<2^j$.
Rationality does not force these ranks to be bounded, even for one fixed
coefficient sequence. More precisely, the assumptions
\[
 0\le c(n)\le n,\qquad \sum_{n\ge1}c(n)2^{-n}\in\Q,
 \qquad v\in\N_{>0},\qquad u_n\in\Z,
\]
\[
 u_{n+1}=2u_n-vc(n+1),\qquad u_n/2^n\longrightarrow0
\]
do not imply $\sup_e\operatorname{rk}_e(u)<\infty$. Nor do they imply
an upper bound that is $o(2^e)$. Such a bound for the actual totient
carry would contradict the lower bound
\lean{Erdos249257.\allowbreak not_irrational_totientSeries_\allowbreak implies_unbounded_carryRank_unconditional}{TotientCarryKernelRigidity.lean:300}
\quad \ev{Lean}, but the following rational comparison already has
that lower bound.

\paragraph{The comparison with value $5/4$.}
Use a centred base-four expansion. The elementary bounds $9/8\le S\le3/2$ follow respectively from the first four terms
and from $\ph(n)\le n-1$ for $n\ge2$. Put $x_0=5/4-S$ and, for $m\ge1$, define
\[
 d_m=\left\lfloor4x_{m-1}+\frac23\right\rfloor,
 \qquad x_m=4x_{m-1}-d_m.
\]
The interval $[-2/3,1/3)$ contains $x_0$ and is preserved by this
recurrence, so $d_m\in\{-2,-1,0,1\}$. Iteration gives
\[
 x_0=\sum_{m=1}^{M}d_m4^{-m}+4^{-M}x_M,
 \qquad |x_M|\le\frac23.
\]
Consequently $x_0=\sum_{m\ge1}d_m4^{-m}$. Define
\[
 c(0)=0,\qquad c(2m-1)=\ph(2m-1),\qquad
 c(2m)=\ph(2m)+d_m\quad(m\ge1).
\]
These are nonnegative integers: for $m\ge2$, $\ph(2m)\ge2$;
for $m=1$, the sharper bound $x_0\ge-1/4$ gives $d_1\ge-1$.
The inequalities $d_m\le1$ and $\ph(2m)\le2m-1$ also give $c(n)\le n$.
Thus $|c(n)-\ph(n)|\le2$ and
\[
 \sum_{n\ge1}c(n)2^{-n}=S+\sum_{m\ge1}d_m4^{-m}=\frac54.
\]
This is the centred base-four construction in the supplied
\href{https://github.com/wcook04/plectis-erdos-lean/blob/52f29ad173b04e3bac941b3663f2b9aebe5de0bb/ErdosProblems/Erdos249/ParityPerturbedRationalControl.lean#L379}{comparison-series theorem}.

Take $v=4$ and $u_N=4\sum_{j\ge1}c(N+j)2^{-j}$. The rational sum makes
$u_N$ integral; the coefficient bound gives $0\le u_N\le4(N+2)$.
Since $c$ agrees with $\ph$ at every odd index, for odd $r$,
\[
 4\ph(2^j n+r)=2u_{2^j n+r-1}-u_{2^j n+r}.
\]
The same $2^e-1$ independent totient sections therefore lie in the
span of the carry sections at levels $1,\ldots,e$, proving
$\operatorname{rk}_e(u)\ge2^e-1$ for every $e\ge0$.
The corresponding formal statement is the
\href{https://github.com/wcook04/plectis-erdos-lean/blob/52f29ad173b04e3bac941b3663f2b9aebe5de0bb/ErdosProblems/Erdos249/ParityPerturbedRationalControl.lean#L478}{unbounded-rank theorem for this comparison}.
This example is distinct from the bounded $3/2$ parity comparison in
Proposition~\ref{prop:parity}. It disproves an inference from the listed
assumptions alone, not one that also uses identities of the actual
totient coefficients. \ev{Math}

\paragraph{The rank quantities.}
Finite-dimensional span of all dyadic sections is the definition of a
$2$-regular sequence in the sense of Allouche and
Shallit~\cite[Definition~2.1]{allouche-shallit} \ev{Cited}.  For
integer-valued sequences this is equivalent to finite generation of the
$\mathbb Z$-module generated by the sections: if their span has dimension $d$,
some $d$ evaluations embed it in $\mathbb Q^d$ and that module in
$\mathbb Z^d$ \ev{Math}.
Coons~\cite[Theorem~3.2]{coons} proved that
$\varphi$ is not $k$-regular for any $k\ge2$ \ev{Cited}.  The
unconditional theorem
\lean{Erdos249257.\allowbreak not_finiteDimensional_\allowbreak span_fullTotientKernel}{TotientMahlerDefect.lean:1145}
proves the case $k=2$ independently.  For integer coefficient and carry sequences, the recurrence transfers
non-$2$-regularity of the coefficients to non-$2$-regularity of the carry.
A proof must include the zero-residue sections, not just the positive
residues. The following argument gives an explicit dimension bound.

For a rational-valued sequence $f$, let $V_e(f)$ be the span of all its sections
through levels $0,\ldots,e$. Define the backward shift by
$(Bf)(0)=0$ and $(Bf)(n)=f(n-1)$ for $n\ge1$, and let $\delta_0$ be $1$
at $0$ and $0$ elsewhere. A level-$j$ section of $Bf$ is a section of
$f$ if its residue is positive. At residue zero it is the backward shift
of $n\mapsto f(2^jn+2^j-1)$. Hence
\[
 V_e(Bf)\subseteq V_e(f)+B(V_e(f)).
\]
For integer sequences $c,u$ and a positive integer $v$, the recurrence
$v c(n+1)=2u(n)-u(n+1)$ gives
\[
 c=\frac{2Bu-u}{v}+\left(c(0)+\frac{u(0)}v\right)\delta_0.
\]
Since every section of $\delta_0$ is $0$ or $\delta_0$, it follows that
\[
 V_e(c)\subseteq V_e(u)+B(V_e(u))+\Q\delta_0,
 \qquad \dim_{\Q} V_e(c)\le2\dim_{\Q} V_e(u)+1.
\]
In particular, bounded carry ranks would force bounded coefficient ranks.
Non-$2$-regularity of $c$ therefore forces unbounded carry ranks.
This is an ordinary proof, not a claim of a new Lean declaration.
It also gives the generic bound
$\dim_{\Q} V_e(u)\ge2^{e-1}$ for $c=\varphi$ and $e\ge1$, using the
coefficient rank $2^e+1$. The factor $2$ is not an additive $O(1)$ loss,
and this bound concerns all sections, including level zero. \ev{Math}

The source theorem for $\varphi$ is sharper: the carry sections at all
residues of levels $1,\ldots,e$ have rank at least $2^e-1$. It uses the independence of the odd-residue
totient sections. Neither lower bound is an irrationality proof, since
rational coefficient series can have nonregular carries. The separate
$5/4$ comparison verifies this directly. Any proposed rank upper bound must
use additional totient-specific hypotheses not satisfied by the $5/4$
comparison sequence.

\subsection{What the available results suggest}\label{sub:shape}

This subsection compares the stated conditions and the recorded finite
observations. It does not infer an asymptotic digit distribution or rank
the prospects of different proof methods.

\paragraph{Comparison of conditions.}
The following table records the conditions themselves, rather than
sufficient digit proxies. Lemma~\ref{lem:orbit} explains how their phases
are related.

\begin{center}
\setlength{\tabcolsep}{5pt}
\begin{tabular}{>{\raggedright\arraybackslash}p{0.16\linewidth} >{\raggedright\arraybackslash}p{0.44\linewidth} >{\raggedright\arraybackslash}p{0.32\linewidth}}
\hline
Condition & Required assertion & Indices and quantifiers \\
\hline
Block norm & Norm at most $21X/25$ for the truncated phase sum &
All $N\in[X,2X)$ at one depth; arbitrarily large $X$ for each $h$ \\
Four terms & The four separate budgets of Definition~\ref{defn:pivot} &
One common $X,L,s,\eta$; $s,\eta$ fixed before $X_0$ \\
Prime gap & Distance from $\Z$ greater than $\arccos(9/10)/(2\pi)$ &
The phase $2^{p-h-1}\alpha_h$ at arbitrarily large primes for each $h$ \\
Full depth & $\mathcal C(td,N,td)$ for some $t\ge1$ &
Every $d\ge1$ and every $N\ge0$ \\
Generic rank bound (false) & A uniform or subexponential upper bound for carries of rational coefficient series &
All finite dyadic levels; refuted by the $5/4$ comparison \\
\hline
\end{tabular}
\end{center}

\noindent
The first four conditions concern the same constant but impose different
averaging, sampling and separation requirements. The $11/100$ digit-change
test in Corollary~\ref{cor:digitform} is sufficient for a real-part bound;
it is not a replacement for the norm or four-term rows. Only the full-depth
row is proved here to be equivalent to irrationality. The last row records
a false generic upper bound, not an additional open hypothesis.

\paragraph{Quantitative estimates.}
The counterexamples identify missing assumptions in the specific
arguments tested here. They do not imply that the binary digits of $S$
are random, normal, or unusually structured. A constructed counterexample
can disprove an implication without resembling $S$ in other respects.
Likewise, failure of one congruence construction does not exclude every
use of the Chinese remainder theorem.

\paragraph{Additional arithmetic information.}
Table~\ref{tab:blocks} records finite observations: the sampled digit-change
proportions are near $0.50$, above the sufficient threshold $0.11$, and
the displayed normalized cosine sums are below $0.89$. These computations
are compatible with the proposed conditions. They do not establish those
conditions for every shift and arbitrarily large blocks, or prove a
fair-coin model for the digits. The relevant quantified estimates remain
unproved here.

The recorded small margins near a forbidden endpoint also require careful
interpretation. One small margin disproves a proposed lower bound larger
than that margin at that index. It does not rule out every positive margin
on a suitable subsequence, or every estimate with scale-dependent bounds.
No numerical census in this record determines which of these possibilities
holds asymptotically.

\paragraph{Scope of the comparison.}
Corollary~\ref{cor:digitform} specifies one sufficient digit statistic:
for every fixed $h$, the digit-change proportion must reach at least
$11/100$ on arbitrarily large blocks. A statement for one shift, a smaller
positive proportion, or finitely many blocks is not this hypothesis. An
upper bound $o(N)$ on the longest run alone would not imply irrationality;
a periodic alternating expansion has bounded runs. The prime-index
condition in \eqref{eq:primefibre} is another proposed sufficient estimate.
Its exact algebraic decomposition does not establish the required
correlation bound.
\section{Remaining arithmetic questions}

Erdős \#249 asks whether $S = \sum_{n\ge 0} \varphi(n)/2^n$ is irrational. It is
\ev{Open}: nothing in this corpus decides it.

This section collects quantified conditions on the scaled tail
$R_N=\sum_{j\ge0}\varphi(N+1+j)2^{-(j+1)}$ that suffice for irrationality.
The implication theorems and the arithmetic assumptions have different
status: the former are proved in the indicated sources, while the latter
are not supplied for the actual totient sequence.
\lean{Erdos249257.TotientTailPeriodKiller.totientTail}{TotientTailPeriodKiller.lean:57}


Each condition is stated with its quantifiers and linked to its formal
source. The comparison distinguishes proved equivalences from one-way
sufficient conditions. It is not an exhaustive list of possible methods.
A finite list of certificates does not stand in for a condition beyond
every threshold, and an implication does not establish its antecedent.


Recall that $H_t=H(t)=\operatorname{lcm}(1,\dots,t)$
(\lean{Erdos249257.CarrySurvivorExtinction.periodLcm}{CarrySurvivorExtinction.lean:515}).
For nonnegative integers $h,N,L$, put
\[
D(h,N,L)=\sum_{j=0}^{L-1}
 \bigl(\varphi(N+h+1+j)-\varphi(N+1+j)\bigr)2^{L-1-j}.
\]
\lean{Erdos249257.TotientTailPeriodKiller.windowDiscrepancy}{TotientTailPeriodKiller.lean:66}
The certificate $\mathcal C(h,N,L)$ is the pair of strict inequalities
\[
 N+h+L+2<D(h,N,L)\bmod2^L<2^L-(N+h+L+2).
\]
\lean{Erdos249257.TotientTailPeriodKiller.certifiedKill}{TotientTailPeriodKiller.lean:72}
\quad \ev{Lean} (the definition; decidable, instance at line 84)
\quad \scale{n/a} \quad \coord{other:binary-window}.
Thus the least nonnegative residue must lie more than $N+h+L+2$
from either endpoint. For fixed $h,N$, this excluded radius grows
linearly with $L$, while its proportion of the modulus tends to zero.

\subsection{Certificate completeness}

For each fixed $h,N\ge0$, a certificate at some \emph{unrestricted}
depth exists exactly when the tail difference is nonintegral:
\[
 \bigl(\exists L\ge0,\ \mathcal C(h,N,L)\bigr)
 \quad\Longleftrightarrow\quad R_{N+h}-R_N\notin\mathbb Z.
\]
\lean{Erdos249257.exists_certifiedKill_iff_tail_diff_notMem_int}{LcmConeFlatness.lean:316}
(The same equivalence is also stated at
\lean{Erdos249257.totient\_tail\_window\_kill\_exists\_iff\_tail\_diff\_not\_int}{CertificateKernel.lean:19025}.)
\quad \ev{Lean} \quad \scale{uniform} (holds for all $h,N$)
\quad \coord{other:binary-window}.

This equivalence removes the existential depth from the unrestricted
certificate conditions below, including the diagonal and LCM-grid forms.
It does not by itself supply a depth bound, a prescribed separation
margin, a common depth for a family, or an exponential-sum estimate.
Such requirements must be retained unless a separate argument removes
them. In particular, completeness does not establish any of the
quantified nonintegrality assertions needed for irrationality.

\subsection{The original quantified residue condition}

\begin{quote}

\textbf{The quantified residue condition.}
\[
  S\notin\mathbb Q
  \quad\Longleftrightarrow\quad
  \forall h\ge1\ \forall N_0\ \exists N\ge N_0\ \exists L,
  \quad\mathcal C(h,N,L).
\]
\end{quote}
Here $\mathcal C(h,N,L)$ is the finite residue test defined earlier;
the universal quantifiers over $h$ and $N_0$ belong to the displayed
assertion, not to that finite test.
\lean{Erdos249257.irrational_totient_series_iff_certificate_supply}{LcmConeFlatness.lean:412}
\quad \ev{Lean} (the equivalence is proved; the quantified existence assertion is not proved)
\quad \scale{cofinal} \quad \coord{other:binary-window}.

The reverse direction is by contradiction against the tail-period law: if $S$ were
rational, the eventual-period theorem
(\lean{Erdos249257.eventual_period_of_not_irrational}{TotientTailPeriodKiller.lean:358},
\ev{Lean}, unconditional) supplies a period $h>0$ and pre-period $N_0$ with
$R_{N+h}-R_N\in\mathbb{Z}$ for all $N\ge N_0$; the certificate hypothesis
instantiated at that $(h,N_0)$ produces an $N\ge N_0$ where
$R_{N+h}-R_N\notin\mathbb{Z}$, a contradiction via nonintegrality from a certificate
(\lean{Erdos249257.tail_diff_notMem_int_of_certifiedKill}{TotientTailPeriodKiller.lean:262}).

For the forward direction, irrationality gives pointwise non-integrality of
every positive tail shift and certificate completeness supplies a witness
already at $N=N_0$.


For fixed $h$, the required certificate may choose both $N\ge N_0$ and
$L$. It is not required at every $N$ or every depth. Its error radius
is $N+h+L+2$, so for fixed $h,N$ it grows only linearly in $L$, whereas
the modulus is $2^L$. The missing assertion is that an admissible
residue lies outside both endpoint intervals beyond every threshold,
for every positive $h$. The reformulations below specify other ways
to choose such witnesses; none establishes this assertion for the
actual totient sequence.

\subsection{Multiples, diagonals, and LCM grids}


The next forms change which tail differences may provide a witness.
Each follows from irrationality and each is sufficient for
irrationality, by the displayed implications and the preceding
equivalence. Their different quantifiers may make different arithmetic
estimates applicable, but logical equivalence alone does not provide
any of those estimates.

\paragraph{Multiples of a period.}
\[
  \forall h_0>0,\ \forall N_0,\ \exists m>0,\ \exists N\ge N_0,\ \exists L,\
    \mathcal{C}(m\cdot h_0, N, L)
  \ \Longrightarrow\ S\notin\mathbb Q.
\]
\lean{Erdos249257.irrational_totient_series_of_multiple_certificate_supply}{CarrySurvivorExtinction.lean:502}
\quad \ev{Lean} \quad \scale{cofinal} \quad \coord{other:lcm-period-multiple}.

The preceding residue condition implies this one by taking $m=1$.
Conversely, a hypothetical period $h_0$ also makes every positive
multiple $mh_0$ a period of the tail fractional parts. A certificate
for one such multiple contradicts that consequence. The freedom to
choose $m$ changes the allowed witnesses, while the quantified
assertion remains equivalent to irrationality.

\paragraph{The diagonal.}
\[
  S\notin\mathbb Q
  \ \Longleftrightarrow\
  \forall t_0,\ \exists t\ge t_0,\ \exists L,\
    \mathcal{C}(H_t, H_t, L).
\]
\lean{Erdos249257.irrational_totient_series_iff_lcm_diagonal_certificate_supply}{LcmConeFlatness.lean:426}
\quad \ev{Lean} \quad \scale{cofinal} \quad \coord{other:lcm-diagonal}.

The two parameters in a hypothetical rational tail period are absorbed
by one LCM index. For $t\ge\max(h_0,N_0)$, both $h_0\mid H_t$ and
$H_t\ge t\ge N_0$ hold. A diagonal certificate then contradicts
periodicity through the multiple-period argument; the precise
intermediate implication is
\lean{Erdos249257.irrational_totient_series_of_lcm_certificate_supply}{LcmDiagonalReduction.lean:92}.
Conversely, irrationality and pointwise completeness give a diagonal
certificate already at $t=t_0$.

The finite assertion $P(t):\equiv\exists L\,\mathcal C(H_t,H_t,L)$
is verified for every $t\le82$
(\lean{ErdosProblems.Skip.LadderT67.exists_diagonalKill_le_82}{ErdosProblems/Skip/LadderT67.lean:71264}).
The historical bank contained 28 explicit values through $t=64$
(\lean{Erdos249257.certifiedKill_diagonal_all_imported}{DiagonalPincerCertificates.lean:2870},
depths $[6,5,7,7,9,14,15,14,21,22,23,26,\dots]$ for
$t\in\{1,2,3,4,5,7,8,9,11,13,16,17,\dots\}$; endpoint
\lean{Erdos249257.certifiedKill_diagonal_t64}{DiagonalPincerCertificateT64Endpoint.lean:1928})
\ev{Cert} \scale{fixed} \coord{other:lcm-diagonal}).  The current contiguous

range is still finite and does not establish witnesses at arbitrarily
large indices; no certificate at $t=83$ is claimed here.

\paragraph{The LCM grid.}
\[
  \forall t_0,\ \exists t\ge t_0,\ \exists q,m,L,\ 0<q\ \wedge\
    \mathcal{C}(m\cdot H_t,\, q\cdot H_t,\, L)
  \ \Longrightarrow\ S\notin\mathbb Q.
\]
\lean{Erdos249257.irrational_totient_series_of_lcm_cone_window_kill_supply}{CertificateKernel.lean:19014}
\quad \ev{Lean} \quad \scale{cofinal} \quad \coord{other:lcm-cone}. The

diagonal choice is $q=m=1$. Allowing other positive grid points
therefore enlarges the choice of finite witnesses. To prove
sufficiency, suppose rationality supplies a period $h$ and pre-period
$N_0$. For $t\ge\max(h,N_0)$, all positive multiples of $H_t$ are
past the pre-period and their differences are multiples of $h$.
Consequently their tails have the same fractional part:
$R_{qH_t+mH_t}-R_{qH_t}\in\mathbb Z$ for every $q>0$ and $m\ge0$
(\lean{Erdos249257.rational_totient_series_forces_lcm_cone_flatness}{CertificateKernel.lean:19002},
\ev{Lean}). A certificate for any such difference is a contradiction.

The finite-grid condition of Theorem~\ref{catalogue:cert:b10a} offers
a second test. At $H=H_t$, choose a finite nonempty
$Q\subseteq\mathbb N_{>0}$ with $B_q=qH_t+L+2<2^L$ for every $q$.
It requires
\[
 \forall q_i\in Q\ \exists q_j\in Q,\qquad
 B_{q_j}<(A_{q_i}-A_{q_j})\bmod2^L,
 \quad A_q=\sum_{j=1}^{L}\varphi(qH_t+j)2^{L-j}.
\]
For every threshold $t_0$, such data at some $t\ge t_0$ would imply
irrationality. This quantified condition, including the nonemptiness
and size assumptions, remains \ev{Open}, \scale{cofinal};
\lean{Erdos249257.irrational_totient_series_of_lcm_cone_nonflat_supply}{CertificateKernel.lean:19140}. The modulus bound is one-sided at each point; it does not imply
a factor-of-two reduction in the number of binary digits.

\paragraph{Equivalent nonintegrality statements}
Via the completeness iff above, the diagonal and cone forms restate,
\emph{exactly}, with no reference to $\mathcal{C}$ at all:
\[
  \forall t_0,\ \exists t\ge t_0,\ R_{2H_t}-R_{H_t}\notin\mathbb{Z}
  \ \Longleftrightarrow\ \text{diagonal certificate condition}
  \ \Longrightarrow\ S\notin\mathbb Q,
\]
\lean{Erdos249257.irrational_totient_series_of_lcm_diagonal_nonintegrality_supply}{CertificateKernel.lean:19034}
(cone analogue at
\lean{Erdos249257.irrational\_totient\_series\_of\_lcm\_cone\_nonintegrality\_supply}{CertificateKernel.lean:19046})
\quad \ev{Lean} \quad \scale{cofinal} \quad \coord{other:real-analytic-nonintegrality}.

Thus one can ask directly whether $R_{2H_t}-R_{H_t}\notin\mathbb Z$
for arbitrarily large $t$. Finite instances are known in this record;
the universally quantified assertion is not proved here.

\subsection{The block exponential-sum condition}


A bound for an exponential sum is sufficient for the quantified
residue condition. With $e(x)=\exp(2\pi i x)$, the assertion is
\[
  \text{block norm condition} :\equiv\
  \forall h>0,\ \forall X_0,\ \exists X,L,\quad
  \max(X_0,1)\le X\ \wedge\
  16\,(2X+h+L+2)\le 2^L\ \wedge\
  \Bigl\|\sum_{N\in[X,2X)} e\bigl(D(h,N,L)/2^L\bigr)\Bigr\| \le \tfrac{21}{25}X.
\]
\lean{Erdos249257.DTWFirstHarmonicNormGap}{FirstHarmonicPivot.lean:75}

\quad \ev{Open} (the universally quantified assertion is not proved)
\quad \scale{cofinal} \quad \coord{other:first-harmonic}.

\[
  \text{block norm condition} \ \Longrightarrow\ S\notin\mathbb Q.
\]
\lean{Erdos249257.irrational_totient_series_of_first_harmonic_norm_gap}{FirstHarmonicPivot.lean:83}
\quad \ev{Lean} \quad \scale{cofinal} \quad \coord{other:first-harmonic}. The

threshold is $X_0$, while $X$ and $L$ are chosen witnesses and may
depend on both $h$ and $X_0$. Taking $X_0=N_0$ gives a certificate
at some $N\in[X,2X)$, hence at $N\ge N_0$. The elementary reason is
that, if every residue failed the test, each would be within
$(2X+h+L+2)/2^L\le1/16$ of an integer. Every summand would then have
real part at least $\cos(\pi/8)>21/25$, contradicting the norm bound.
The implication is proved; the bound on suitable blocks beyond every
threshold remains to be established.

The underlying finite lemma is more general: it applies to any nonempty
finite $T\subseteq[0,2X)$ with the stated size bound and
$\|\sum_{N\in T}e(D(h,N,L)/2^L)\|\le(21/25)|T|$.
Taking $T=[X,2X)$ gives the full-block implication. Allowing a chosen
subset weakens the hypothesis; it does not strengthen it. The source
\lean{Erdos249257.exists_certifiedKill_of_first_harmonic_norm_gap}{FirstHarmonicPivot.lean:53}
uses the elementary block lemma
\lean{Erdos249257.exists_certifiedKill_of_first_harmonic_gap}{FirstHarmonicGap.lean:128}
(\ev{Lean}, unconditional as an implication). The real-part version uses
$\operatorname{Re}\sum_{N\in T}e(D(h,N,L)/2^L)\le(9/10)|T|$.
Its proof uses $\cos(\pi/8)>9/10$ and averaging. An unspecified saving
below $|T|$ is not the stated hypothesis.

\emph{What supplying this form would take:} genuine
cancellation for the additive character $e(D(h,N,L)/2^L)$ at the
growing dyadic moduli permitted by the size bound. The quantifiers require
suitable blocks beyond every threshold for each fixed $h>0$; they do not
require one threshold that works simultaneously for all $h$. \lean{Erdos249257.ResidualGaugeObstruction.det_residualMonomialMatrix}{ResidualGaugeObstruction.lean:47}
gives a comparison in which every weighted phase equals $1$ while the
specified determinant remains nonzero. Thus nonvanishing of that determinant
alone cannot prove cancellation. A successful use of a determinant would
need further information connecting its entries to the phase sum; the
comparison does not exclude every such argument.

\subsection{The actual tail difference on the LCM diagonal}

The same irrationality question can be expressed on the sparse sequence
$H=H_{2^a}$. The equivalence below is useful because its right-hand side
contains only the actual tail differences $\Omega_a$, not an auxiliary
approximation.

\begin{quote}
\textbf{The exact equivalence.}
\[
  S\notin\mathbb Q
  \iff
  \forall a_0,\ \exists a\ge a_0,\ \Omega_a\notin\mathbb{Z}.
\]
\end{quote}
\lean{Erdos257PeriodNoncollapse.DiagonalFreshLossBridge.PowerTwoOddWindowAffine.irrational_totientSeries_iff_actualLcmOrbitNonintegralitySupply}{TotientActualLcmOrbitNonintegrality.lean:37}
\quad \ev{Lean} \quad \scale{n/a} (an unconditional equivalence, not a
supply) \quad \coord{mobius-mersenne}, where
$\Omega_a := 2^H(2^H-1)S - (\Phi_{2H}-\Phi_H)$,
$H=H_{2^a}$
(\lean{Erdos257PeriodNoncollapse.DiagonalFreshLossBridge.PowerTwoOddWindowAffine.actualLcmTailOrbit_eq_scaled_totientSeries_sub_prefix}{TotientActualLcmOrbitSeparation.lean:68}).
The right-hand side is an equivalent reformulation of irrationality,
not a weaker hypothesis already established by the finite examples.
The following sufficient conditions are attempts to prove it.


\paragraph{An impossible complete divisibility pattern}
Let $a\ge8$, $H=H_{2^a}$ and $J,K,m\in\mathbb N$, with
\[
 m\le K,\qquad J+K+a+6<2\cdot2^a,\qquad
 B:=2H+J+K+2<2^m.
\]
It is impossible that
\[
 2^{r+1}\mid\delta_{2^a}(J+K-m+r+1)
              \quad\text{for every }0\le r<m.
\]
\lean{Erdos257PeriodNoncollapse.DiagonalFreshLossBridge.PowerTwoOddWindowAffine.not_actualLcmTerminalDyadicStaircase_of_room}{TotientActualLcmTopEdgeStaircase.lean:251}
\quad \ev{Lean} \quad \scale{bounded} \quad \coord{mobius-mersenne}.
Indeed, the last requirement would make $\delta_{2^a}(J+K)$ a
multiple of $2^m$, whereas the short-window bounds give
$0<\delta_{2^a}(J+K)<B<2^m$. The positivity input is
\[
 \delta_{2^a}(j)>0\qquad(a\ge8,\ 0<j<2\cdot2^a),
\]
\lean{Erdos257PeriodNoncollapse.DiagonalFreshLossBridge.PowerTwoOddWindowAffine.lcmRayArithmeticLetter_pos_of_lt_two_mul}{TotientActualLcmOrbitArithmetic.lean:1703}, without a rationality hypothesis.

Omitting only the last divisibility requirement leaves the partial
pattern in Proposition~\ref{prop:TE-02-inv}. It additionally requires
$\delta_{2^a}(J+K)\le2^m-B$. The resulting necessary equality is
\emph{penultimate}:
\[
 \delta_{2^a}(J+K-1)=2^{m-1},\qquad B<2^m<2B.
\]
\lean{Erdos257PeriodNoncollapse.DiagonalFreshLossBridge.PowerTwoOddWindowAffine.puncturedDyadicStaircase_penultimate_eq_half}{TotientActualLcmTopEdgeStaircase.lean:1187}
\quad \ev{Lean} \quad \scale{bounded}.
The terminal difference is not pinned to $2^{m-1}$. Neither this
necessary restriction nor the full-pattern impossibility supplies an
unbounded family of successful partial patterns.

\paragraph{The remaining upper-endpoint residue}
\[
  \mathcal G_a(J,K,m) :\equiv
  m\le K\ \wedge\
  2H+J+K+2 < 2^m\ \wedge\
  D(H,\,H+J,\,K) \bmod 2^m \le 2^m - (2H+J+K+2).
\]
\lean{Erdos257PeriodNoncollapse.DiagonalFreshLossBridge.PowerTwoOddWindowAffine.ActualLcmTopEdgeResidueGap}{TotientActualLcmTopEdgeStaircase.lean:900}
\quad (definition; an unbounded family satisfying it is \ev{Open})
\quad \scale{bounded} \quad \coord{mobius-mersenne}.

The positivity theorem gives
\[
 R_{2H+J}-R_{H+J}>0
 \quad\text{if }a\ge8\text{ and }J+a+6<2\cdot2^a.
\]
\lean{Erdos257PeriodNoncollapse.DiagonalFreshLossBridge.PowerTwoOddWindowAffine.actualLcmTailDiff_shift_pos}{TotientActualLcmOrbitSign.lean:39}
For $a\ge8$, the residue test at depth $K$ uses the stronger bound
$J+K+a+6<2\cdot2^a$. Under these hypotheses, an integral initial
tail difference would give an integer terminal carry
$e=R_{2H+J+K}-R_{H+J+K}$ with $0<e<2H+J+K+2$. The tail identity is
\[
 D(H,H+J,K)=2^K(R_{2H+J}-R_{H+J})-e.
\]
Hence, when $2H+J+K+2<2^K$, it is the residue of $D(H,H+J,K)$,
not of $e$, that equals $2^K-e$:
\lean{Erdos257PeriodNoncollapse.DiagonalFreshLossBridge.PowerTwoOddWindowAffine.actualLcm_integral_forces_topEdgeResidue}{TotientActualLcmOrbitSign.lean:211}.
More generally, the bounds $m\le K$ and $2H+J+K+2<2^m$ give
\[
 D(H,H+J,K)\bmod2^m=2^m-e>2^m-(2H+J+K+2),
\]
because $2^m$ divides $2^K$ and $0<e<2^m$. This contradicts
$\mathcal G_a(J,K,m)$. Consequently, for $a\ge8$ and
$J+K+a+6<2\cdot2^a$,
\[
 \mathcal G_a(J,K,m)\quad\Longrightarrow\quad
       R_{2H+J}-R_{H+J}\notin\mathbb Z.
\]
\lean{Erdos257PeriodNoncollapse.DiagonalFreshLossBridge.PowerTwoOddWindowAffine.actualLcmTailDiff_notMem_int_of_topEdgeResidueGap}{TotientActualLcmTopEdgeStaircase.lean:1260}
The one-sided inequality excludes the upper-endpoint interval forced
by integrality. Positivity alone does not establish that inequality.

The
cofinal target built from it:
\[
  \text{cofinal upper-endpoint condition} :\equiv\
  \forall a_0,\ \exists a,K,m,\quad
  a_0\le a\ \wedge\ 8\le a\ \wedge\ K+(a{+}6)<2\cdot 2^a\ \wedge\
  \mathcal G_a(0,K,m).
\]
\lean{Erdos257PeriodNoncollapse.DiagonalFreshLossBridge.PowerTwoOddWindowAffine.PowerTwoActualLcmTopEdgeResidueGapSupply}{TotientActualLcmTopEdgeStaircase.lean:1325}
\quad \ev{Open} \quad \scale{cofinal} \quad \coord{mobius-mersenne}.
\[
  \text{cofinal upper-endpoint condition} \Longrightarrow S\notin\mathbb Q.
\]
\lean{Erdos257PeriodNoncollapse.DiagonalFreshLossBridge.PowerTwoOddWindowAffine.irrational_totientSeries_of_topEdgeResidueGapSupply}{TotientActualLcmTopEdgeStaircase.lean:2184}
\quad \ev{Lean} \quad \scale{cofinal} \quad \coord{mobius-mersenne}.

\paragraph{Five sufficient conditions for that residue inequality}
The next five quantified conditions each imply irrationality. Conditions
(i)--(iv) reach the upper-endpoint test; (v) gives nonintegrality directly.
Their proven relations are summarised in Proposition~\ref{prop:TE-05}.

\begin{enumerate}
\item[(i)] For every $a_0$, there are $a\ge\max(a_0,8)$ and $m\ge0$
with $H=H_{2^a}$ such that
\[
\begin{gathered}
 m+1+a+6<2\cdot2^a,\qquad 2H+m+3<2^m,\\
 2H+m+2\le
 \bigl(D(H,H,m)+\delta_{2^a}(m+1)\bigr)\bmod2^m
 \le2^m-(2H+m+3),
\end{gathered}
\]
where $\delta_{2^a}(j)=\varphi(2H+j)-\varphi(H+j)$.
The size condition includes the larger depth $m+1$, so the conclusion
remains valid whichever of the two adjacent depths gives the certificate. \lean{Erdos257PeriodNoncollapse.DiagonalFreshLossBridge.PowerTwoOddWindowAffine.PowerTwoAdjacentSuffixMidbandSupply}{TotientActualLcmTopEdgeStaircase.lean:1334}
\ev{Open} \scale{cofinal} \coord{mobius-mersenne}. Sufficiency:
\lean{Erdos257PeriodNoncollapse.DiagonalFreshLossBridge.PowerTwoOddWindowAffine.powerTwoActualLcmTopEdgeResidueGapSupply_of_adjacentSuffixMidband}{TotientActualLcmTopEdgeStaircase.lean:2108}, direct endpoint
\lean{Erdos257PeriodNoncollapse.DiagonalFreshLossBridge.PowerTwoOddWindowAffine.irrational_totientSeries_of_adjacentSuffixMidbandSupply}{TotientActualLcmTopEdgeStaircase.lean:2191}.

\item[(ii)] For every $a_0$, there is $a\ge\max(a_0,14)$ with the
following property. Put $H=H_{2^a}$ and choose $q$ so that $2q+1$ is the
least odd integer not smaller than $\lfloor\log_2H\rfloor+10$.
Define integers $W_{a,j}$ recursively by
\[
\begin{aligned}
 W_{a,0}&=0,\\
 W_{a,j+1}&=4W_{a,j}
  +\frac{\delta_{2^a}(2j+3)-2\delta_{2^a}(2j+2)
                  +\delta_{2^a}(2j+4)}{2}.
\end{aligned}
\]
The increments $\delta_{2^a}(s)$ are even, so the recurrence is integral.
The required inequality is
\[
 H+q+2\le W_{a,q}\bmod4^q\le4^q-(H+q+2).
\]
At depth $m=2q+1$, the residue in (i) equals
$2(W_{a,q}\bmod4^q)$. Dividing its two endpoint bounds by $2$ and using
integrality gives exactly the threshold $H+q+2$.
\lean{Erdos257PeriodNoncollapse.DiagonalFreshLossBridge.PowerTwoOddWindowAffine.PowerTwoOddGuardTopEdgeHalfWordBandSupply}{TotientActualLcmTopEdgeStaircase.lean:1349}
\ev{Open} \scale{cofinal} \coord{mobius-mersenne}. The width is $H+q+2$ and the depth is the prescribed guarded depth;
both must be retained when comparing this with a fixed-fraction band. Sufficiency:
\lean{Erdos257PeriodNoncollapse.DiagonalFreshLossBridge.PowerTwoOddWindowAffine.powerTwoAdjacentSuffixMidbandSupply_of_oddGuardTopEdgeHalfWordBand}{TotientActualLcmTopEdgeStaircase.lean:2057}.


\item[(iii)] Use exactly the quantifiers and prescribed depth of (ii).
Let $r=W_{a,q}\bmod4^q$ and define its centred representative by
\[
 u_{a,q}=\begin{cases}
 r,&2r\le4^q,\\
 r-4^q,&2r>4^q.
 \end{cases}
\]
Thus $-4^q/2<u_{a,q}\le4^q/2$; when a midpoint occurs, the positive
representative is used. The equivalent inequality is
\[
 H+q+2\le |u_{a,q}|.
\]
\lean{Erdos257PeriodNoncollapse.DiagonalFreshLossBridge.PowerTwoOddWindowAffine.PowerTwoActualFinalTopEdgeMagnitudeSupply}{TotientActualLcmTopEdgeStaircase.lean:1471}
\ev{Open} \scale{cofinal} \coord{mobius-mersenne}.
At the prescribed depth the size bound $2(H+q+2)\le4^q$ holds.
For a residue in $[0,4^q)$, the band in (ii) is therefore exactly
this absolute-value condition. The equivalence, including its
quantifiers, is proved:
\lean{Erdos257PeriodNoncollapse.DiagonalFreshLossBridge.PowerTwoOddWindowAffine.topEdgeHalfWordBandSupply_iff_actualFinalCenteredMagnitudeSupply}{TotientActualLcmTopEdgeStaircase.lean:1479} \quad \ev{Lean}.

\item[(iv)] Allow the odd depth to vary, but retain both size bounds.
For every $a_0$, require $a\ge\max(a_0,8)$ and $q\in\mathbb N$ with
\[
\begin{gathered}
 2q+2+a+6<2\cdot2^a,\qquad 2(H+q+2)\le4^q,\\
 H+q+2\le |u_{a,q}|,\qquad H=H_{2^a}.
\end{gathered}
\]
\lean{Erdos257PeriodNoncollapse.DiagonalFreshLossBridge.PowerTwoOddWindowAffine.PowerTwoFlexibleActualTopEdgeMagnitudeSupply}{TotientActualLcmTopEdgeStaircase.lean:1927}
\ev{Open} \scale{cofinal} \coord{mobius-mersenne}.
The first bound keeps the adjacent depth inside the short window;
the second makes the required residue band nonempty. One implication
passes through (i), and another gives the two-sided carry exclusion
below: \lean{Erdos257PeriodNoncollapse.DiagonalFreshLossBridge.PowerTwoOddWindowAffine.powerTwoAdjacentSuffixMidbandSupply_of_flexibleActualTopEdgeMagnitude}{TotientActualLcmTopEdgeStaircase.lean:2007}, direct endpoint \lean{Erdos257PeriodNoncollapse.DiagonalFreshLossBridge.PowerTwoOddWindowAffine.irrational_totientSeries_of_flexibleActualTopEdgeMagnitudeSupply}{TotientActualLcmTopEdgeStaircase.lean:2213}.

\item[(v)] Retain the full quantifiers and the two size bounds in
(iv), but replace its absolute-value inequality by
\[
 \delta_{2^a}(2q+2)\le2u_{a,q}.
\]
This compares the signed centred residue with the terminal totient
difference. It is not an unrestricted assertion at an arbitrary $q$.
\lean{Erdos257PeriodNoncollapse.DiagonalFreshLossBridge.PowerTwoOddWindowAffine.PowerTwoFlexibleActualTerminalDominanceSupply}{TotientActualLcmTopEdgeStaircase.lean:1908}
\ev{Open} \scale{cofinal} \coord{mobius-mersenne}. Direct endpoint
\lean{Erdos257PeriodNoncollapse.DiagonalFreshLossBridge.PowerTwoOddWindowAffine.irrational_totientSeries_of_flexibleActualTerminalDominanceSupply}{TotientActualLcmTopEdgeStaircase.lean:2221}.
\end{enumerate}

\paragraph{A further sufficient terminal-carry condition}
With $H=H_{2^a}$, the unproved quantified condition is
\[
\begin{aligned}
 &\forall a_0\ \exists a\ge\max(a_0,8)\ \exists q:\\
 &\quad 2q+2+(a+6)<2\cdot2^a,\qquad 2(H+q+2)\le4^q,\\
 &\quad 2u_{a,q}\le\delta_{2^a}(2q+2)-(2H+2q+3)
       \quad\hbox{or}\quad \delta_{2^a}(2q+2)\le2u_{a,q}.
\end{aligned}
\]
\lean{Erdos257PeriodNoncollapse.DiagonalFreshLossBridge.PowerTwoOddWindowAffine.PowerTwoFlexibleActualTerminalCarryCorridorEscapeSupply}{TotientActualLcmTopEdgeStaircase.lean:1895}
Either alternative gives nonintegrality at that scale, so the
quantified condition implies irrationality
\lean{Erdos257PeriodNoncollapse.DiagonalFreshLossBridge.PowerTwoOddWindowAffine.irrational_totientSeries_of_flexibleActualTerminalCarryCorridorEscapeSupply}{TotientActualLcmTopEdgeStaircase.lean:2230}.

The proof assumes an integer $z=\Omega_a$ and uses
\[
 2u_{a,q}=\delta_{2^a}(2q+2)-c_{H,H,z}(2q+1).
\]
\lean{Erdos257PeriodNoncollapse.DiagonalFreshLossBridge.PowerTwoOddWindowAffine.two_mul_actualOddHalfCenteredLift_eq_terminal_sub_trueCarry}{TotientActualLcmTopEdgeStaircase.lean:1678}
The carry is then a positive tail difference, strictly smaller than
$2H+2q+3$. The source name referring to a negative endpoint concerns the
representative \emph{minus} this carry
\lean{Erdos257PeriodNoncollapse.DiagonalFreshLossBridge.PowerTwoOddWindowAffine.actualLcm_trueEndpointSurvivor_neg}{TotientActualLcmOrbitSign.lean:172}.
The first alternative contradicts the upper bound; the second
contradicts positivity, which is the terminal-dominance argument
\lean{Erdos257PeriodNoncollapse.DiagonalFreshLossBridge.PowerTwoOddWindowAffine.actualLcmTailOrbit_notMem_int_of_actualTerminalDominance}{TotientActualLcmTopEdgeStaircase.lean:1880}.

Both arguments are conditional contradictions, not evidence against
either unproved estimate. No converse from nonintegrality to either
finite inequality is asserted. In particular, the exact identity must
not be described as making this sufficient condition necessary.


\paragraph{Certificates in short windows}
For every $t,j\in\mathbb N$, the arithmetic difference on the LCM
progression is exactly
\[
 \delta_t(j)=\varphi(2H_t+j)-\varphi(H_t+j).
\]
No coprimality restriction on $j$ is needed
(\lean{Erdos257PeriodNoncollapse.DiagonalFreshLossBridge.PowerTwoOddWindowAffine.lcmRayArithmeticLetter_eq_deltaTotient}{TotientActualLcmOrbitArithmetic.lean:298}, \ev{Lean}). Consequently the corresponding window
numerator is
\[
 D(H_t,H_t,L)=\sum_{j=1}^{L}\delta_t(j)2^{L-j},
\]
and applying the same residue and error bounds gives precisely the
test $\mathcal C(H_t,H_t,L)$ (\lean{Erdos257PeriodNoncollapse.DiagonalFreshLossBridge.PowerTwoOddWindowAffine.lcmDiagonalArithmeticKill_iff_certifiedKill}{TotientActualLcmOrbitArithmetic.lean:2045}, \ev{Lean}).
Thus the quantified assertion
\[
 \forall a_0\ \exists a\ge a_0\ \exists L<2\cdot2^a,
 \qquad\mathcal C(H_{2^a},H_{2^a},L)
\]
\lean{Erdos257PeriodNoncollapse.DiagonalFreshLossBridge.PowerTwoOddWindowAffine.PowerTwoActualLcmShortArithmeticKillSupply}{TotientActualLcmOrbitArithmetic.lean:2107}, \ev{Open}, \scale{cofinal}, \coord{mobius-mersenne},
implies irrationality through \lean{Erdos257PeriodNoncollapse.DiagonalFreshLossBridge.PowerTwoOddWindowAffine.actualLcmOrbitNonintegralitySupply_of_shortArithmeticKill}{TotientActualLcmOrbitArithmetic.lean:2152}.
It restricts the diagonal indices to powers of two and imposes the
additional bound $L<2\cdot2^a$. Pointwise completeness allows an
unspecified finite depth; it does not imply a witness within this
short window.

The displayed certified examples have $a=4$, $L=23$ and $a=6$, $L=93$:
\lean{Erdos257PeriodNoncollapse.DiagonalFreshLossBridge.PowerTwoOddWindowAffine.actualLcmTailOrbit_four_and_six_notMem_int}{TotientActualLcmShortKill.lean:35} \ev{Cert} \scale{fixed}.
The bounded statement \lean{Erdos257PeriodNoncollapse.DiagonalFreshLossBridge.PowerTwoOddWindowAffine.powerTwoActualLcmShortArithmeticKillSupply_through_six}{TotientActualLcmShortKill.lean:54}
(\ev{Lean}, \scale{bounded}) uses the single witness $(a,L)=(6,93)$
to cover all thresholds $a_0\le6$. It makes no assertion for larger
thresholds, or that these are the only successful exponents.


\paragraph{A sufficient approximation bound}
Put $H=H_{2^a}$ and define the finite rational and its error bound by
\[
 \rho_{a,q}=
 \frac{D(H,H,2q+1)+\delta_{2^a}(2q+2)}{2^{2q+1}},\qquad
 \varepsilon_{a,q}=\frac{2H+2q+3}{2^{2q+1}}.
\]
For every $a,q\in\mathbb N$,
\[
 |\Omega_a-\rho_{a,q}|<\varepsilon_{a,q}.
\]
\lean{Erdos257PeriodNoncollapse.DiagonalFreshLossBridge.PowerTwoOddWindowAffine.abs_actualLcmTailOrbit_sub_rawApprox_lt}{TotientActualLcmOrbitSeparation.lean:141}
\quad \ev{Lean} \quad \scale{uniform} \quad \coord{mobius-mersenne}.
Let $2q_a+1$ be the least odd integer not smaller than
$\lfloor\log_2H\rfloor+10$, as in condition (ii). The sufficient
quantified separation condition is
\[
 \forall a_0\ \exists a\ge\max(2,a_0)\ \forall z\in\mathbb Z,
 \qquad \tfrac1{32}+\varepsilon_{a,q_a}\le|\Omega_a-z|.
\]
\ev{Open} \scale{cofinal} \coord{mobius-mersenne}; sufficiency is
proved in \lean{Erdos257PeriodNoncollapse.DiagonalFreshLossBridge.PowerTwoOddWindowAffine.irrational_totientSeries_of_actualLcmOrbitSeparationSupply}{TotientActualLcmOrbitSeparation.lean:305} (\ev{Lean}). The triangle inequality then gives
$|\rho_{a,q_a}-z|>1/32$ for every integer $z$.
The hypothesis concerns the real tail $\Omega_a$, not merely its
computable approximant. The truncation estimate transfers the stated
separation, but does not itself establish it. Also, the depth is
prescribed by $a$; it is not a free parameter in this hypothesis.

\subsection{An equivalent test using two leading residue bits}

Proposition~\ref{prop:TE-03-inv} gives an equivalent certificate
condition: there are $s,b\in\mathbb N$ such that either
$\mathcal C(h,N+s,b+1)$ holds, or both
\[
 N+s+h+b+4<2^b,\qquad
 2^b\le D(h,N+s,b+2)\bmod2^{b+2}<3\cdot2^b.
\]
\lean{Erdos257PeriodNoncollapse.DiagonalFreshLossBridge.PowerTwoOddWindowAffine.exists_certifiedKill_iff_guardCylinderWitness}{TotientActualLcmTopEdgeStaircase.lean:844}
\quad \ev{Lean} \quad \scale{uniform} \quad \coord{seam-integer}.
The forward proof selects $b$ using the depth $L$ of the first
certificate, so this is an encoding of a witness rather than a bound
on how far one must search. Existence for the actual totient sequence
remains the missing assertion. A use for another coefficient sequence
would require its own recurrence and tail bounds.

\subsection{Farey growth and the failed generic rank argument}

The Farey growth law below is a separate unproved assertion.  The
rank statement is retained only to record a counterfactual sufficient input
and the reason that the generic compression route to it is closed; it is not a
second open problem supported by the present argument.

\paragraph{The hypothetical rank bound.}
Unconditionally, for every level $e\ge1$, the dyadic totient basis of
$2^e+1$ sequences is linearly independent over $\mathbb{Q}$ (via CRT and
Dirichlet's theorem on primes in arithmetic progression):
\lean{Erdos249257.linearIndependent_canonicalTotientKernelFamily}{TotientMahlerDefect.lean:935}
\quad \ev{Lean} \quad \scale{uniform} \quad \coord{other:carry-kernel-rank}.
Consequently, if $S\in\Q$, there is an integer $v>0$ with $vS\in\Z$.
Its integral carry $u_N=vR_N$ satisfies, for every $e\ge0$,
\[
 \dim_{\Q}\operatorname{span}
 \{n\mapsto u_{2^j n+r}:1\le j\le e,\ 0\le r<2^j\}\ge2^e-1.
\]
It obeys $u_{N+1}=2u_N-v\ph(N+1)$ and $0\le u_N\le v(N+2)$.
\lean{Erdos249257.not_irrational_totientSeries_implies_unbounded_carryRank_unconditional}{TotientCarryKernelRigidity.lean:300}
\quad \ev{Lean} \quad \scale{uniform} \quad \coord{other:carry-kernel-rank}.

A bound independent of $e$ for these same carry-section ranks would
contradict the displayed lower bound; so would an upper bound that is
$o(2^e)$. Rationality and linear coefficient growth alone do not give
either conclusion: the fixed sequence constructed in
Section~\ref{sub:generic-carry-rank} has sum $5/4$ and the same rank
lower bound. An incompatible upper bound for the actual totient carry
would have to use information not shared by that comparison sequence.

The most direct generic construction is also impossible: no
integer certificate satisfying the following conditions ($Q\cdot v\cdot A =
\mathrm{boundary}$ with $|\mathrm{boundary}|<Q\cdot v$ and $A\ne0$) can
exist; \lean{Erdos249257.false_of_compressedAdjointCertificate}{TotientMahlerDefect.lean:1183}
\quad \ev{Lean}; and the full family is proved infinite-dimensional in
span, the case $k=2$ of Coons's non-regularity theorem
(\S\ref{sec:mahler-defect}):
\lean{Erdos249257.not_finiteDimensional_span_fullTotientKernel}{TotientMahlerDefect.lean:1145}
\quad \ev{Lean}.

The conditional source statement gives a common eventual period
modulo $v$ for the sections of this same $u$, together with their
unbounded rational ranks:
\lean{Erdos249257.not_irrational_totientSeries_implies_mod_period_and_unbounded_rank}{TotientTailCarryPeriod.lean:224}
\quad \ev{Lean}. These conclusions concern different properties of the same carry; they supply no rank
upper bound.

\paragraph{Growth of the Farey denominator bound.}
The finite denominator exclusion for $S$ comes from
a Farey-gap denominator exclusion at a single fixed window $K=240$:
\[
 S\ne a/b\quad\text{for all }a\in\Z,\quad
 1\le b\le79639646646701375323355774875831053,
 \quad\gcd(a,b)=1.
\]
\lean{Erdos249257.tsum_totient_div_pow_two_ne_ratCast_of_den_le_79639646646701375323355774875831053}{CertificateKernel.lean:18384}
\quad \ev{Lean} \quad \scale{bounded} \quad \coord{farey}. This bound is
\emph{sharp} at $K=240$; the mediant denominator
$q=79639646646701375323355774875831054$ is the exact first failing
denominator
(\lean{Erdos249257.gap_check_window_1_240_first_failure}{GapFareyBound.lean:225},
\ev{Lean}); so the same interval cannot exclude every rational at the next
denominator. A different window requires a new certificate or a
uniform argument; the fixed-window calculation supplies neither.


A sufficient extension is a function $g(K)\to\infty$ for which every
$(N=1,K)$ gap check excludes all rationals of reduced denominator
at most $g(K)$. More generally, unbounded valid exclusion bounds
along a sequence of windows suffice. This is a statement about the
certified intervals. The record does not prove it equivalent to a
lower bound on the continued-fraction denominators of $S$.
\ev{Open} \scale{cofinal} \coord{farey}.

For $p$ prime and $e\ge1$, a successful unit-gap certificate at
denominator $p^e$ has zero candidates or a single candidate divisible
by $p$; \lean{Erdos249257.reducedDenominatorCandidateCount_prime_pow_le_one}{PrimitiveWeightCertificate.lean:111}
\quad \ev{Lean}. The success hypothesis is essential: the bound is
not a statement about every window at that denominator. Nor does it
bound the eventual denominator cutoff. The prime-power restriction
cannot be omitted, as the modulo-$15$ example in
Observation~\ref{prop:D8cert-kill} shows. No obstruction to a uniform
estimate as $K$ varies is proved here.

\subsection{Summary of implications}

The second column gives the relation to the quantified residue condition
or to the explicitly named intermediate condition.
\begin{center}
\setlength{\tabcolsep}{5pt}
\begin{longtable}{>{\raggedright\arraybackslash}p{0.29\linewidth} >{\raggedright\arraybackslash}p{0.35\linewidth} >{\footnotesize\raggedright\arraybackslash}p{0.28\linewidth}}
\textbf{Condition} & \textbf{Proved relation} & \textbf{Source} \\
\hline
\endhead
Quantified $\mathcal C(h,N,L)$ condition (the base condition) & \textbf{equivalent} to $S\notin\mathbb Q$ & \nolinkurl{LcmConeFlatness.lean:412} \\
Certificates for multiples of a period & equivalent via the base condition: take multiplier one in one direction and use the cited implication in the other & \nolinkurl{CarrySurvivorExtinction.lean:502} \\
Diagonal certificates beyond every threshold & \textbf{equivalent} to $S\notin\mathbb Q$; one universally quantified threshold & \nolinkurl{LcmConeFlatness.lean:426} \\
Certificates on LCM grids beyond every threshold & equivalent via the base condition; diagonal certificates have $q=m=1$ & \nolinkurl{CertificateKernel.lean:19014} \\
Quantified finite-grid condition & sufficient; requires each one-sided size bound & \nolinkurl{CertificateKernel.lean:19140} \\
Nonintegral diagonal or LCM-grid tail differences & \textbf{equivalent} to the corresponding unrestricted-depth certificate condition & \nolinkurl{CertificateKernel.lean:19034}, \nolinkurl{CertificateKernel.lean:19046} \\
Block norm condition & sufficient for the quantified residue condition; no converse proved & \nolinkurl{FirstHarmonicPivot.lean:83} \\
$\Omega_a\notin\mathbb Z$ beyond every threshold & \textbf{equivalent} to $S\notin\mathbb Q$ itself & \nolinkurl{TotientActualLcmOrbitNonintegrality.lean:37} \\
Short-window certificate condition & sufficient; imposes an additional short-depth restriction at powers of two & \nolinkurl{TotientActualLcmOrbitArithmetic.lean:2107} \\
Upper-endpoint residue condition beyond every threshold & sufficient for $\Omega_a\notin\mathbb Z$ beyond every threshold & \nolinkurl{TotientActualLcmTopEdgeStaircase.lean:1325} \\
Adjacent-depth residue band, condition (i) & sufficient for the upper-endpoint residue gap & \nolinkurl{TotientActualLcmTopEdgeStaircase.lean:1334} \\
Prescribed-depth residue band, condition (ii) & \textbf{equivalent} to condition (iii) & \nolinkurl{TotientActualLcmTopEdgeStaircase.lean:1349} \\
Prescribed-depth centred magnitude, condition (iii) & \textbf{equivalent} to condition (ii) & \nolinkurl{TotientActualLcmTopEdgeStaircase.lean:1471} \\
Variable-depth centred magnitude, condition (iv) & sufficient for both the midband and the carry exclusion & \nolinkurl{TotientActualLcmTopEdgeStaircase.lean:1927} \\
Terminal comparison, condition (v) & sufficient for carry exclusion; one-sided & \nolinkurl{TotientActualLcmTopEdgeStaircase.lean:1908} \\
Either terminal-carry exclusion inequality & sufficient; no converse asserted & \nolinkurl{TotientActualLcmTopEdgeStaircase.lean:1895} \\
Real-tail separation with approximation margin & sufficient; fixes the depth and approximation margin & \nolinkurl{TotientActualLcmOrbitSeparation.lean:305} \\
Carry-rank upper bound & the rational $5/4$ comparison refutes a bound from rationality and the stated coefficient hypotheses alone & \nolinkurl{ParityPerturbedRationalControl.lean:478} \\
Unbounded certified denominator exclusions & sufficient; the cited result is one finite exclusion, not an unbounded family & \nolinkurl{GapFareyBound.lean}, \nolinkurl{CertificateKernel.lean:18384} \\
Shifted short certificate or two-bit residue test & equivalent existence test with unrestricted depth; no a priori search-depth bound follows & \nolinkurl{TotientActualLcmTopEdgeStaircase.lean:844} \\
\end{longtable}
\end{center}

\emph{Reading the table.} ``Sufficient'' records only the indicated
implication. It does not assert a converse or a strict comparison of
strength. ``Equivalent'' records two proved directions, either directly
or through the base equivalence with $S\notin\mathbb Q$. Logical
equivalence does not by itself provide a direct construction from the
finite witnesses in one formulation to those in another.

\subsection{What remains to be proved}

The implications above leave explicit arithmetic assertions to prove
for $\varphi$. Two types of estimate can be stated directly, without
assuming either that they follow from known methods or that they exhaust
the possible approaches.



\paragraph{Exponential sums.} For each fixed $h>0$ and every
threshold $X_0$, the block norm condition asks for some $X\ge\max(X_0,1)$
and $L$ satisfying both the modulus bound and
\[
 \left|\sum_{N\in[X,2X)}e\bigl(D(h,N,L)/2^L\bigr)\right|
       \le\frac{21}{25}X.
\]
The saving $21/25$ is fixed, but the chosen block and depth may depend
on $h$ and $X_0$; uniformity of these choices in $h$ is not required.
The preceding four-term decomposition and the proved counting bounds
identify contributions to this sum. The two remaining weighted
estimates must still be obtained with the same parameters. The record
therefore gives an exact reduction and some of its estimates, not a
proof of the displayed cancellation bound.

\paragraph{The last totient difference.} The endpoint criteria compare
\[
 \delta_{2^a}(2q+2)
   =\varphi(2H+2q+2)-\varphi(H+2q+2),\qquad H=H_{2^a},
\]
with twice the centred residue $u_{a,q}$. An appropriate inequality,
proved for arbitrarily large $a$ with the stated bounds on $q$, would
contradict the carry interval required by integrality. The exact identity
explains why such an inequality would suffice; it does not supply the
inequality. The finite examples through $t=82$, and the short-window
examples at $a=4,6$, do not establish an asymptotic pattern.

\clearpage
\section*{Declarations}

\paragraph{Further background.}
Erick Wong posted a proof that
$\sum_{n\ge1}(\varphi(n)\bmod k)k^{-n}$ is irrational for every integer $k>2$
on Mathematics Stack Exchange on 29 March 2015~\cite{wong2015}; there the base
equals the modulus.  The M\"obius-square form of $S$ and its coprimality
interpretation were posted publicly before this record, as credited in
\S\ref{ssec:coprime}.  Kova\v{c} and Tao show that suitable subsums of Lambert
series in distinct integer bases can be rational~\cite[Theorem~2.3]{kovac-tao};
the freedom to choose those subsums is absent for $S$.  The Lambert-series
factorisations of Merca~\cite[Theorem~1.2]{merca2017} and of Merca and
Schmidt~\cite{merca-schmidt} are transform background for the
M\"obius--Mersenne coordinate.  Balasubramanian, Giri and Srivastav prove
asymptotics for finite shifted convolutions of divisor-convolution
functions, including $\varphi(n)/n$, uniformly for $|h|\le x/2$ at
cutoff $x$~\cite[Theorem~2.2]{bgs2016}. Their partial-summation extension
uses differentiable weights with bounded derivative on the summation
interval~\cite[Rem.~2.9]{bgs2016}.
Applying that observation here would require control of the weighted
error for the present phases; no such estimate is established in this
record. The cited results are not themselves the dyadic residue
separation required by $\mathrm{Sep}$.  The remaining sources are cited where their hypotheses are
used.

At the generating-function level, the published height criterion of
Adamczewski, Bell and Smertnig gives $k$-regularity for a $k$-Mahler
series with integer coefficients of polynomial growth
~\cite[Thm.~1.2(b), p.~2528]{adamczewski-bell-smertnig2023}.
For the totient coefficients the logarithmic heights are
$\log\max(1,\ph(n))\le\log n$ for $n\ge1$, so Coons's nonregularity
already rules out $k$-Mahler equations in every base. The later
multiplicative classification of Bell and Smertnig provides another
route~\cite[Thm.~1.3 and p.~3]{bell-smertnig2026}.
The algebraic and $D$-finite classifications of Bell, Bruin and Coons
belong to the same function-level context~\cite[Thms.~1.5--1.6]{bell-bruin-coons}.
They do not by themselves decide the value at $z=1/2$.

Automatic prime-sampling theorems address a different issue.
Adamczewski, Drmota and M\"ullner identify a computable local coprimality
condition governing the relevant logarithmic densities, with a conditional
ordinary-density transfer~\cite[Thm.~1.4 and Rem.~1.5]{adamczewski-drmota-mullner}.
Their theorem does not assert that every output state has positive density
on a prescribed prime progression, nor does it prove independence with
automatic coefficient weights. Any use for that purpose needs its own
compatibility argument. \ev{Cited}

\paragraph{Sources and data.} The linked Lean declarations,
fixed toolchain, and library manifest are published in the companion
repository. The ordinary arguments in these papers explain the
mathematics; the pinned declarations record the corresponding formal
statements. Their hypotheses must be checked separately when an
implication is applied.

\subsection*{Formalisation coverage and remaining dependencies}
\label{sec:coverage}
\papersectiontarget{coverage}

Every theorem, lemma, proposition and corollary of this record other than the
five listed below has a Lean statement of the same assertion, with the same
hypotheses, checked by the Lean kernel using only the axioms
\texttt{propext}, \texttt{Classical.choice} and \texttt{Quot.sound}, in the
development at revision \texttt{181078b6b009}.

Proposition~\ref{prop:dickman}, one of the five, was proved later.  It has a
Lean statement of the same assertion, checked by the Lean kernel using only
those three axioms, at revision \texttt{e4ba6841dfdb}; the paragraph after
its proof describes the Lean argument.

Propositions~\ref{prop:zetaq} and~\ref{catalogue:mob:a4} share one input. The
displayed three-term identity is unconditional and checked in both. The
development verifies the reduction of the irrationality clause to the theorem
of Postelmans and Van Assche that $1$, $\zeta_q(1)$ and $\zeta_q(2)$ are
linearly independent over $\mathbb{Q}$ for $q=1/p$ with an integer
$p\ge2$~\cite[Theorem~1.3, p.~3]{postelmans-vanassche}. A formal proof of that
input is not included in the verified development, so those two endpoints
remain conditional on it. Proposition~\ref{prop:D7-inv} is the same situation
with a different input. Its four Lambert identities and the identification of
$L(\mathrm{Id})$ with $\sum_{m\ge1}\sigma(m)/2^m$ are checked, and the
development verifies the reduction of the transcendence of that value to
Nesterenko's theorem on the algebraic independence of the three Eisenstein
values at an algebraic point of the punctured unit
disc~\cite[Cor.~2, p.~1320]{nesterenko1996}.

Proposition~\ref{prop:badcof} is incomplete in a different way: part of its
proof is formalised and a named later step is not. The development carries
the counting function of indices with excluded cofactor, and neither
Schoenberg's theorem nor the partial summation over $m\le\sqrt X/2$ is
formalised.

A theorem whose own statement is conditional is formalised exactly as stated.
Propositions~\ref{prop:zetaq}, \ref{catalogue:mob:a4}, \ref{prop:D7-inv}
and~\ref{prop:badcof} are asserted unconditionally in this record; the first
three have Lean counterparts carrying a named input, and for the last only
part of the argument is formalised.

% BEGIN GENERATED CONCORDANCE (claude_lane/build_concordance.py; do not edit by hand)
\paragraph{Concordance of statements and Lean declarations.}
Each result of this record that has a kernel-checked Lean statement of the same assertion is listed
below with the declarations that jointly state it.  Each name links to its declaration at revision
\texttt{181078b6b009}.  Where the Lean statement is stronger than the printed one and implies it by
an immediate specialisation, the entry says so.  An entry marked \emph{compared} was also checked
independently: a restatement of the same declarations against Mathlib alone, together with its
proof, was verified by Comparator (\texttt{leanprover/comparator}) in a clean continuous
integration environment, in the run named by its number.  Comparator trusts the restated
statement, so the correspondence between the printed result and that statement is the one this
concordance records.

\begingroup\small\raggedright
\par\noindent\hangindent=1.5em Theorem~\ref{thm:denom}: \href{https://github.com/wcook04/plectis-erdos/blob/181078b6b009d905809cf2e007ad309386a3d1e0/lean/Erdos249257/CertificateKernel.lean\#L18384}{\nolinkurl{tsum_totient_div_pow_two_ne_ratCast_of_den_le_79639646646701375323355774875831053}}, \href{https://github.com/wcook04/plectis-erdos/blob/181078b6b009d905809cf2e007ad309386a3d1e0/lean/Erdos249257/GapFareyBound.lean\#L225}{\nolinkurl{gap_check_window_1_240_first_failure}}.  \emph{Compared}, runs \href{https://github.com/wcook04/plectis-erdos-lean/actions/runs/35544127144}{\texttt{35544127144}}, \href{https://github.com/wcook04/plectis-erdos-lean/actions/runs/35715779405}{\texttt{35715779405}}.
\par\noindent\hangindent=1.5em Proposition~\ref{prop:deposits}: \href{https://github.com/wcook04/plectis-erdos/blob/181078b6b009d905809cf2e007ad309386a3d1e0/lean/Erdos249257/DiagonalPincerCertificatesT64.lean\#L1967}{\nolinkurl{certifiedKill_diagonal_all_imported_through_t64}}, \href{https://github.com/wcook04/plectis-erdos/blob/181078b6b009d905809cf2e007ad309386a3d1e0/lean/Erdos249257/DiagonalPincerCertificatesT64.lean\#L1928}{\nolinkurl{certifiedKill_diagonal_t64}}, \href{https://github.com/wcook04/plectis-erdos/blob/181078b6b009d905809cf2e007ad309386a3d1e0/lean/Erdos249257/CarrySurvivorExtinction.lean\#L574}{\nolinkurl{certifiedKill_all_upto_sixteen}}, \href{https://github.com/wcook04/plectis-erdos/blob/181078b6b009d905809cf2e007ad309386a3d1e0/lean/ErdosProblems/Erdos249/PeriodMultipleEscape.lean\#L471}{\nolinkurl{certifiedKill_67_300}}, \href{https://github.com/wcook04/plectis-erdos/blob/181078b6b009d905809cf2e007ad309386a3d1e0/lean/ErdosProblems/Erdos249/PeriodMultipleEscape.lean\#L474}{\nolinkurl{certifiedKill_81_300}}, \href{https://github.com/wcook04/plectis-erdos/blob/181078b6b009d905809cf2e007ad309386a3d1e0/lean/ErdosProblems/Erdos249/PeriodMultipleEscape.lean\#L477}{\nolinkurl{certifiedKill_97_300}}, \href{https://github.com/wcook04/plectis-erdos/blob/181078b6b009d905809cf2e007ad309386a3d1e0/lean/ErdosProblems/Erdos249/PeriodMultipleEscape.lean\#L480}{\nolinkurl{certifiedKill_101_300}}, \href{https://github.com/wcook04/plectis-erdos/blob/181078b6b009d905809cf2e007ad309386a3d1e0/lean/ErdosProblems/Erdos249/PeriodMultipleEscape.lean\#L483}{\nolinkurl{certifiedKill_121_300}}, \href{https://github.com/wcook04/plectis-erdos/blob/181078b6b009d905809cf2e007ad309386a3d1e0/lean/ErdosProblems/Erdos249/PeriodMultipleEscape.lean\#L486}{\nolinkurl{certifiedKill_125_300}}, \href{https://github.com/wcook04/plectis-erdos/blob/181078b6b009d905809cf2e007ad309386a3d1e0/lean/ErdosProblems/Erdos249/PeriodMultipleEscape.lean\#L490}{\nolinkurl{certifiedKill_127_300}}, \href{https://github.com/wcook04/plectis-erdos/blob/181078b6b009d905809cf2e007ad309386a3d1e0/lean/ErdosProblems/Erdos249/PeriodMultipleEscape.lean\#L493}{\nolinkurl{certifiedKill_128_300}}, \href{https://github.com/wcook04/plectis-erdos/blob/181078b6b009d905809cf2e007ad309386a3d1e0/lean/ErdosProblems/Skip/LadderT67.lean\#L71264}{\nolinkurl{exists_diagonalKill_le_82}}.  \emph{Compared}, runs \href{https://github.com/wcook04/plectis-erdos-lean/actions/runs/35643815458}{\texttt{35643815458}}, \href{https://github.com/wcook04/plectis-erdos-lean/actions/runs/35674034595}{\texttt{35674034595}}, \href{https://github.com/wcook04/plectis-erdos-lean/actions/runs/35682858662}{\texttt{35682858662}}.
\par\noindent\hangindent=1.5em Proposition~\ref{prop:sign}: \href{https://github.com/wcook04/plectis-erdos/blob/181078b6b009d905809cf2e007ad309386a3d1e0/lean/Erdos249257/TotientActualLcmOrbitSign.lean\#L39}{\nolinkurl{actualLcmTailDiff_shift_pos}}, \href{https://github.com/wcook04/plectis-erdos/blob/181078b6b009d905809cf2e007ad309386a3d1e0/lean/Erdos249257/TotientActualLcmOrbitSign.lean\#L211}{\nolinkurl{actualLcm_integral_forces_topEdgeResidue}}, \href{https://github.com/wcook04/plectis-erdos/blob/181078b6b009d905809cf2e007ad309386a3d1e0/lean/Erdos249257/CarrySurvivorExtinction.lean\#L393}{\nolinkurl{carryOrbit_eq_tail_diff}}.  \emph{Compared}, run \href{https://github.com/wcook04/plectis-erdos-lean/actions/runs/35674034595}{\texttt{35674034595}}.
\par\noindent\hangindent=1.5em Proposition~\ref{prop:rank}: \href{https://github.com/wcook04/plectis-erdos/blob/181078b6b009d905809cf2e007ad309386a3d1e0/lean/Erdos249257/TotientCarryKernelRigidity.lean\#L300}{\nolinkurl{not_irrational_totientSeries_implies_unbounded_carryRank_unconditional}}.  \emph{Compared}, run \href{https://github.com/wcook04/plectis-erdos-lean/actions/runs/35544127144}{\texttt{35544127144}}.
\par\noindent\hangindent=1.5em Proposition~\ref{prop:period-not-rank}: \href{https://github.com/wcook04/plectis-erdos/blob/181078b6b009d905809cf2e007ad309386a3d1e0/lean/Erdos249257/TotientTailCarryPeriod.lean\#L224}{\nolinkurl{not_irrational_totientSeries_implies_mod_period_and_unbounded_rank}}.  \emph{Compared}, run \href{https://github.com/wcook04/plectis-erdos-lean/actions/runs/35544127144}{\texttt{35544127144}}.
\par\noindent\hangindent=1.5em Proposition~\ref{prop:iffs}: \href{https://github.com/wcook04/plectis-erdos/blob/181078b6b009d905809cf2e007ad309386a3d1e0/lean/ErdosProblems/Erdos249/PeriodMultipleEscape.lean\#L432}{\nolinkurl{periodMultipleKillSupply_iff_irrational}}, \href{https://github.com/wcook04/plectis-erdos/blob/181078b6b009d905809cf2e007ad309386a3d1e0/lean/Erdos249257/LcmConeFlatness.lean\#L412}{\nolinkurl{irrational_totient_series_iff_certificate_supply}}, \href{https://github.com/wcook04/plectis-erdos/blob/181078b6b009d905809cf2e007ad309386a3d1e0/lean/Erdos249257/LcmConeFlatness.lean\#L426}{\nolinkurl{irrational_totient_series_iff_lcm_diagonal_certificate_supply}}, \href{https://github.com/wcook04/plectis-erdos/blob/181078b6b009d905809cf2e007ad309386a3d1e0/lean/Erdos249257/TotientTailCarryPeriod.lean\#L875}{\nolinkurl{irrational_totientSeries_iff_cofinalDirectedLcmCertificateSupply}}, \href{https://github.com/wcook04/plectis-erdos/blob/181078b6b009d905809cf2e007ad309386a3d1e0/lean/Erdos249257/PivotAntiReconstruction.lean\#L1765}{\nolinkurl{dtwWindowSeparatedPairs_iff_irrational_totient_series}}.  \emph{Compared}, runs \href{https://github.com/wcook04/plectis-erdos-lean/actions/runs/35643815458}{\texttt{35643815458}}, \href{https://github.com/wcook04/plectis-erdos-lean/actions/runs/35674034595}{\texttt{35674034595}}.
\par\noindent\hangindent=1.5em Proposition~\ref{prop:parity}: \href{https://github.com/wcook04/plectis-erdos/blob/181078b6b009d905809cf2e007ad309386a3d1e0/lean/Erdos249257/TotientParityCoboundaryCountermodel.lean\#L637}{\nolinkurl{exists_totientParity_arbitrarilyManySeparatedCarry_rational_countermodel}}, \href{https://github.com/wcook04/plectis-erdos/blob/181078b6b009d905809cf2e007ad309386a3d1e0/lean/Erdos249257/TotientParityCoboundaryCountermodel.lean\#L359}{\nolinkurl{tsum_parityCoboundaryWeight_eq_three_halves}}.  \emph{Compared}, run \href{https://github.com/wcook04/plectis-erdos-lean/actions/runs/35624228171}{\texttt{35624228171}}.
\par\noindent\hangindent=1.5em Lemma~\ref{lem:gsound}: \href{https://github.com/wcook04/plectis-erdos/blob/181078b6b009d905809cf2e007ad309386a3d1e0/lean/ErdosProblems/Erdos249/PaperCompleteR20/GenericTailCertificates.lean\#L66}{\nolinkurl{certificate_sound}}, \href{https://github.com/wcook04/plectis-erdos/blob/181078b6b009d905809cf2e007ad309386a3d1e0/lean/ErdosProblems/Erdos249/PaperCompleteR20/GenericTailCertificates.lean\#L48}{\nolinkurl{scaled_difference}}, \href{https://github.com/wcook04/plectis-erdos/blob/181078b6b009d905809cf2e007ad309386a3d1e0/lean/ErdosProblems/Erdos249/PaperCompleteR20/GenericTailCertificates.lean\#L55}{\nolinkurl{truncation_error_bound}}.  \emph{Compared}, run \href{https://github.com/wcook04/plectis-erdos-lean/actions/runs/35643815458}{\texttt{35643815458}}.
\par\noindent\hangindent=1.5em Lemma~\ref{lem:gperiod}: \href{https://github.com/wcook04/plectis-erdos/blob/181078b6b009d905809cf2e007ad309386a3d1e0/lean/ErdosProblems/Erdos249/PaperCompleteR20/GenericTailCertificates.lean\#L98}{\nolinkurl{generic_tail_period}}, \href{https://github.com/wcook04/plectis-erdos/blob/181078b6b009d905809cf2e007ad309386a3d1e0/lean/ErdosProblems/Erdos249/PaperCompleteR20/GenericTailCertificates.lean\#L33}{\nolinkurl{scaled_tail_split}}.  \emph{Compared}, run \href{https://github.com/wcook04/plectis-erdos-lean/actions/runs/35643815458}{\texttt{35643815458}}.
\par\noindent\hangindent=1.5em Theorem~\ref{thm:gamma}: \href{https://github.com/wcook04/plectis-erdos/blob/181078b6b009d905809cf2e007ad309386a3d1e0/lean/ErdosProblems/Erdos249/PaperCompleteR20/FinitePrefixCountermodel.lean\#L15}{\nolinkurl{gamma_le}}, \href{https://github.com/wcook04/plectis-erdos/blob/181078b6b009d905809cf2e007ad309386a3d1e0/lean/ErdosProblems/Erdos249/PaperCompleteR20/FinitePrefixCountermodel.lean\#L21}{\nolinkurl{gamma_prefix}}, \href{https://github.com/wcook04/plectis-erdos/blob/181078b6b009d905809cf2e007ad309386a3d1e0/lean/ErdosProblems/Erdos249/PaperCompleteR20/FinitePrefixCountermodel.lean\#L24}{\nolinkurl{discrepancy_prefix}}, \href{https://github.com/wcook04/plectis-erdos/blob/181078b6b009d905809cf2e007ad309386a3d1e0/lean/ErdosProblems/Erdos249/PaperCompleteR20/FinitePrefixCountermodel.lean\#L108}{\nolinkurl{exact_series}}, \href{https://github.com/wcook04/plectis-erdos/blob/181078b6b009d905809cf2e007ad309386a3d1e0/lean/ErdosProblems/Erdos249/PaperCompleteR20/FinitePrefixCountermodelEndpoint.lean\#L31}{\nolinkurl{value_cast}}, \href{https://github.com/wcook04/plectis-erdos/blob/181078b6b009d905809cf2e007ad309386a3d1e0/lean/ErdosProblems/Erdos249/PaperCompleteR20/FinitePrefixCountermodelEndpoint.lean\#L39}{\nolinkurl{exact_denominator}}, \href{https://github.com/wcook04/plectis-erdos/blob/181078b6b009d905809cf2e007ad309386a3d1e0/lean/ErdosProblems/Erdos249/PaperCompleteR20/FinitePrefixCountermodelEndpoint.lean\#L51}{\nolinkurl{no_certificate_after_prefix}}, \href{https://github.com/wcook04/plectis-erdos/blob/181078b6b009d905809cf2e007ad309386a3d1e0/lean/ErdosProblems/Erdos249/PaperCompleteR20/FinitePrefixCountermodelEndpoint.lean\#L69}{\nolinkurl{gamma_not_separation}}.  \emph{Compared}, run \href{https://github.com/wcook04/plectis-erdos-lean/actions/runs/35643815458}{\texttt{35643815458}}.
\par\noindent\hangindent=1.5em Corollary~\ref{cor:b1}: \href{https://github.com/wcook04/plectis-erdos/blob/181078b6b009d905809cf2e007ad309386a3d1e0/lean/ErdosProblems/Erdos249/PaperCompleteR20/FinitePrefixCountermodelEndpoint.lean\#L74}{\nolinkurl{no_uniform_prefix_rule}}, \href{https://github.com/wcook04/plectis-erdos/blob/181078b6b009d905809cf2e007ad309386a3d1e0/lean/ErdosProblems/Erdos249/PaperCompleteR20/FinitePrefixCountermodelEndpoint.lean\#L69}{\nolinkurl{gamma_not_separation}}.  \emph{Compared}, run \href{https://github.com/wcook04/plectis-erdos-lean/actions/runs/35643815458}{\texttt{35643815458}}.
\par\noindent\hangindent=1.5em Proposition~\ref{prop:b2}: \href{https://github.com/wcook04/plectis-erdos/blob/181078b6b009d905809cf2e007ad309386a3d1e0/lean/ErdosProblems/Erdos249/PaperCompleteR21/ThreeParticularEquivalences.lean\#L33}{\nolinkurl{three_particular_equivalences}}, \href{https://github.com/wcook04/plectis-erdos/blob/181078b6b009d905809cf2e007ad309386a3d1e0/lean/ErdosProblems/Erdos249/PaperCompleteR21/ThreeParticularEquivalences.lean\#L55}{\nolinkurl{two_point_sample_numerical_requirement}}.  \emph{Compared}, runs \href{https://github.com/wcook04/plectis-erdos-lean/actions/runs/35643815458}{\texttt{35643815458}}, \href{https://github.com/wcook04/plectis-erdos-lean/actions/runs/35674034595}{\texttt{35674034595}}.
\par\noindent\hangindent=1.5em Proposition~\ref{prop:b4}: \href{https://github.com/wcook04/plectis-erdos/blob/181078b6b009d905809cf2e007ad309386a3d1e0/lean/ErdosProblems/Erdos249/PaperCompleteR21/TwoAdicPulseCertificateFailure.lean\#L29}{\nolinkurl{twoAdic_pulse_construction_never_certifies}}, \href{https://github.com/wcook04/plectis-erdos/blob/181078b6b009d905809cf2e007ad309386a3d1e0/lean/ErdosProblems/Erdos249/PaperCompleteR21/TwoAdicPulseCertificateFailure.lean\#L60}{\nolinkurl{twoAdic_pulse_defining_congruence}}, \href{https://github.com/wcook04/plectis-erdos/blob/181078b6b009d905809cf2e007ad309386a3d1e0/lean/ErdosProblems/Erdos249/PaperCompleteR21/TwoAdicPulseCertificateFailure.lean\#L85}{\nolinkurl{twoAdic_pulse_error_bound}}.  \emph{Compared}, runs \href{https://github.com/wcook04/plectis-erdos-lean/actions/runs/35643815458}{\texttt{35643815458}}, \href{https://github.com/wcook04/plectis-erdos-lean/actions/runs/35674034595}{\texttt{35674034595}}.
\par\noindent\hangindent=1.5em Proposition~\ref{prop:b5}: \href{https://github.com/wcook04/plectis-erdos/blob/181078b6b009d905809cf2e007ad309386a3d1e0/lean/ErdosProblems/Erdos249/PaperCompleteR21/CarryDescriptionInformationLoss.lean\#L93}{\nolinkurl{balancedPulse_common_history}}, \href{https://github.com/wcook04/plectis-erdos/blob/181078b6b009d905809cf2e007ad309386a3d1e0/lean/ErdosProblems/Erdos249/PaperCompleteR21/CarryDescriptionInformationLoss.lean\#L30}{\nolinkurl{balancedPulse_tail_at}}, \href{https://github.com/wcook04/plectis-erdos/blob/181078b6b009d905809cf2e007ad309386a3d1e0/lean/ErdosProblems/Erdos249/PaperCompleteR21/CarryDescriptionInformationLoss.lean\#L76}{\nolinkurl{balancedPulse_series}}, \href{https://github.com/wcook04/plectis-erdos/blob/181078b6b009d905809cf2e007ad309386a3d1e0/lean/ErdosProblems/Erdos249/PaperCompleteR21/CarryDescriptionInformationLoss.lean\#L114}{\nolinkurl{balancedPulse_label_lower_bound}}, \href{https://github.com/wcook04/plectis-erdos/blob/181078b6b009d905809cf2e007ad309386a3d1e0/lean/ErdosProblems/Erdos249/PaperCompleteR21/CarryDescriptionInformationLoss.lean\#L122}{\nolinkurl{balancedPulse_no_decoder_from_common_state}}, \href{https://github.com/wcook04/plectis-erdos/blob/181078b6b009d905809cf2e007ad309386a3d1e0/lean/ErdosProblems/Erdos249/PaperCompleteR21/CarryDescriptionInformationLoss.lean\#L134}{\nolinkurl{affineBinaryOrbit_difference_and_reset}}, \href{https://github.com/wcook04/plectis-erdos/blob/181078b6b009d905809cf2e007ad309386a3d1e0/lean/ErdosProblems/Erdos249/PaperCompleteR21/CarryDescriptionInformationLoss.lean\#L146}{\nolinkurl{fixed_precision_carry_completion}}.
\par\noindent\hangindent=1.5em Proposition~\ref{prop:b6} (the Lean statement is stronger): \href{https://github.com/wcook04/plectis-erdos/blob/181078b6b009d905809cf2e007ad309386a3d1e0/lean/ErdosProblems/Erdos249/PaperCompleteR21/FourLinearConstructionLimits.lean\#L72}{\nolinkurl{b6_retained_dyadic_sections_independent}}, \href{https://github.com/wcook04/plectis-erdos/blob/181078b6b009d905809cf2e007ad309386a3d1e0/lean/ErdosProblems/Erdos249/PaperCompleteR21/FourLinearConstructionLimits.lean\#L80}{\nolinkurl{b6_compressed_adjoint_identity_impossible}}, \href{https://github.com/wcook04/plectis-erdos/blob/181078b6b009d905809cf2e007ad309386a3d1e0/lean/ErdosProblems/Erdos249/PaperCompleteR21/FourLinearConstructionLimits.lean\#L101}{\nolinkurl{b6_mobius_incidence_unimodular_and_injective}}, \href{https://github.com/wcook04/plectis-erdos/blob/181078b6b009d905809cf2e007ad309386a3d1e0/lean/ErdosProblems/Erdos249/PaperCompleteR21/FourLinearConstructionLimits.lean\#L129}{\nolinkurl{b6_adjugate_tail_cost_floor}}, \href{https://github.com/wcook04/plectis-erdos/blob/181078b6b009d905809cf2e007ad309386a3d1e0/lean/ErdosProblems/Erdos249/PaperCompleteR21/FourLinearConstructionLimits.lean\#L165}{\nolinkurl{b6_synthetic_sequence_prescribed_differences}}, \href{https://github.com/wcook04/plectis-erdos/blob/181078b6b009d905809cf2e007ad309386a3d1e0/lean/ErdosProblems/Erdos249/PaperCompleteR21/FourLinearConstructionLimits.lean\#L199}{\nolinkurl{b6_synthetic_shift_combinations_same_form}}, \href{https://github.com/wcook04/plectis-erdos/blob/181078b6b009d905809cf2e007ad309386a3d1e0/lean/ErdosProblems/Erdos249/PaperCompleteR21/FourLinearConstructionLimits.lean\#L236}{\nolinkurl{b6_mobiusMersennePrefix_eq_icc_sum}}, \href{https://github.com/wcook04/plectis-erdos/blob/181078b6b009d905809cf2e007ad309386a3d1e0/lean/ErdosProblems/Erdos249/PaperCompleteR21/FourLinearConstructionLimits.lean\#L253}{\nolinkurl{b6_mobiusMersenne_rung_estimates}}, \href{https://github.com/wcook04/plectis-erdos/blob/181078b6b009d905809cf2e007ad309386a3d1e0/lean/ErdosProblems/Erdos249/PaperCompleteR21/FourLinearConstructionLimits.lean\#L262}{\nolinkurl{b6_mobiusMersenneTheta_two_eq_totientSeries_sub_half}}, \href{https://github.com/wcook04/plectis-erdos/blob/181078b6b009d905809cf2e007ad309386a3d1e0/lean/ErdosProblems/Erdos249/PaperCompleteR21/FourLinearConstructionLimits.lean\#L270}{\nolinkurl{b6_rankOneSubrankQuotient_sub_totientSeries_offset_gt}}.  \emph{Compared}, run \href{https://github.com/wcook04/plectis-erdos-lean/actions/runs/35674034595}{\texttt{35674034595}}.
\par\noindent\hangindent=1.5em Proposition~\ref{prop:b7}: \href{https://github.com/wcook04/plectis-erdos/blob/181078b6b009d905809cf2e007ad309386a3d1e0/lean/ErdosProblems/Erdos249/PaperCompleteR21/RationalParityCountermodelProperties.lean\#L28}{\nolinkurl{exists_rational_parityComparison}}, \href{https://github.com/wcook04/plectis-erdos/blob/181078b6b009d905809cf2e007ad309386a3d1e0/lean/ErdosProblems/Erdos249/PaperCompleteR21/RationalParityCountermodelProperties.lean\#L37}{\nolinkurl{parityComparisonProperties_do_not_imply_irrational}}.  \emph{Compared}, run \href{https://github.com/wcook04/plectis-erdos-lean/actions/runs/35643815458}{\texttt{35643815458}}.
\par\noindent\hangindent=1.5em Proposition~\ref{prop:shift}: \href{https://github.com/wcook04/plectis-erdos/blob/181078b6b009d905809cf2e007ad309386a3d1e0/lean/Erdos249257/TotientTailPeriodKiller.lean\#L150}{\nolinkurl{two_pow_mul_totient_series_eq}}.  \emph{Compared}, run \href{https://github.com/wcook04/plectis-erdos-lean/actions/runs/35643815458}{\texttt{35643815458}}.
\par\noindent\hangindent=1.5em Lemma~\ref{lem:farey}: \href{https://github.com/wcook04/plectis-erdos/blob/181078b6b009d905809cf2e007ad309386a3d1e0/lean/Erdos249257/GapFareyBound.lean\#L51}{\nolinkurl{farey_gap}}.  \emph{Compared}, run \href{https://github.com/wcook04/plectis-erdos-lean/actions/runs/35715779405}{\texttt{35715779405}}.
\par\noindent\hangindent=1.5em Proposition~\ref{prop:gapwindow}: \href{https://github.com/wcook04/plectis-erdos/blob/181078b6b009d905809cf2e007ad309386a3d1e0/lean/Erdos249257/CertificateKernel.lean\#L18310}{\nolinkurl{totient_carry_residue_window_1_240_eq}}, \href{https://github.com/wcook04/plectis-erdos/blob/181078b6b009d905809cf2e007ad309386a3d1e0/lean/Erdos249257/GapFareyBound.lean\#L176}{\nolinkurl{gap_check_window_1_240_le_79639646646701375323355774875831053}}, \href{https://github.com/wcook04/plectis-erdos/blob/181078b6b009d905809cf2e007ad309386a3d1e0/lean/Erdos249257/GapFareyBound.lean\#L225}{\nolinkurl{gap_check_window_1_240_first_failure}}.  \emph{Compared}, runs \href{https://github.com/wcook04/plectis-erdos-lean/actions/runs/35624228171}{\texttt{35624228171}}, \href{https://github.com/wcook04/plectis-erdos-lean/actions/runs/35715779405}{\texttt{35715779405}}.
\par\noindent\hangindent=1.5em Theorem~\ref{thm:denom-record}: \href{https://github.com/wcook04/plectis-erdos/blob/181078b6b009d905809cf2e007ad309386a3d1e0/lean/Erdos249257/CertificateKernel.lean\#L18384}{\nolinkurl{tsum_totient_div_pow_two_ne_ratCast_of_den_le_79639646646701375323355774875831053}}.  \emph{Compared}, run \href{https://github.com/wcook04/plectis-erdos-lean/actions/runs/35544127144}{\texttt{35544127144}}.
\par\noindent\hangindent=1.5em Proposition~\ref{prop:coprime}: \href{https://github.com/wcook04/plectis-erdos/blob/181078b6b009d905809cf2e007ad309386a3d1e0/lean/Erdos249257/CertificateKernel.lean\#L18544}{\nolinkurl{tsum_visible_coprime_pairs_eq_totient_series}}, \href{https://github.com/wcook04/plectis-erdos/blob/181078b6b009d905809cf2e007ad309386a3d1e0/lean/Erdos249257/CertificateKernel.lean\#L18557}{\nolinkurl{totient_series_eq_half_add_visible_coprime_pairs}}.  \emph{Compared}, run \href{https://github.com/wcook04/plectis-erdos-lean/actions/runs/35624228171}{\texttt{35624228171}}.
\par\noindent\hangindent=1.5em Proposition~\ref{prop:gcdlayer}: \href{https://github.com/wcook04/plectis-erdos/blob/181078b6b009d905809cf2e007ad309386a3d1e0/lean/Erdos249257/GeometricCoprimality.lean\#L480}{\nolinkurl{tsum_gcd_layer_pos_coprime_half_eq_one}}, \href{https://github.com/wcook04/plectis-erdos/blob/181078b6b009d905809cf2e007ad309386a3d1e0/lean/Erdos249257/GcdMomentCalculus.lean\#L266}{\nolinkurl{tsum_pos_pair_both_dvd_half_eq_inv_mersenne_sq}}.  \emph{Compared}, run \href{https://github.com/wcook04/plectis-erdos-lean/actions/runs/35715779405}{\texttt{35715779405}}.
\par\noindent\hangindent=1.5em Theorem~\ref{thm:denomcoprime}: \href{https://github.com/wcook04/plectis-erdos/blob/181078b6b009d905809cf2e007ad309386a3d1e0/lean/Erdos249257/CertificateKernel.lean\#L18572}{\nolinkurl{tsum_visible_coprime_pairs_ne_int_div_of_den_le_39819823323350687661677887437915526}}.  \emph{Compared}, run \href{https://github.com/wcook04/plectis-erdos-lean/actions/runs/35624228171}{\texttt{35624228171}}.
\par\noindent\hangindent=1.5em Theorem~\ref{thm:denommobsq}: \href{https://github.com/wcook04/plectis-erdos/blob/181078b6b009d905809cf2e007ad309386a3d1e0/lean/Erdos249257/CertificateKernel.lean\#L18487}{\nolinkurl{tsum_moebius_div_two_pow_sub_one_sq_ne_int_div_of_den_le_39819823323350687661677887437915526}}.  \emph{Compared}, run \href{https://github.com/wcook04/plectis-erdos-lean/actions/runs/35624228171}{\texttt{35624228171}}.
\par\noindent\hangindent=1.5em Proposition~\ref{prop:mobsq}: \href{https://github.com/wcook04/plectis-erdos/blob/181078b6b009d905809cf2e007ad309386a3d1e0/lean/ErdosProblems/Erdos249/PaperCompleteR20/MobiusSquareReduction.lean\#L10}{\nolinkurl{mobius_square_reduction}}, \href{https://github.com/wcook04/plectis-erdos/blob/181078b6b009d905809cf2e007ad309386a3d1e0/lean/ErdosProblems/Erdos249/PaperCompleteR20/MobiusSquareReduction.lean\#L17}{\nolinkurl{irrational_totient_iff_mobius_square}}, \href{https://github.com/wcook04/plectis-erdos/blob/181078b6b009d905809cf2e007ad309386a3d1e0/lean/ErdosProblems/Erdos249/PaperCompleteR20/MobiusSquareReduction.lean\#L25}{\nolinkurl{moebius_three_values}}.  \emph{Compared}, runs \href{https://github.com/wcook04/plectis-erdos-lean/actions/runs/35624228171}{\texttt{35624228171}}, \href{https://github.com/wcook04/plectis-erdos-lean/actions/runs/35643815458}{\texttt{35643815458}}.
\par\noindent\hangindent=1.5em Proposition~\ref{prop:lambertengine}: \href{https://github.com/wcook04/plectis-erdos/blob/181078b6b009d905809cf2e007ad309386a3d1e0/lean/Erdos249257/GcdMomentCalculus.lean\#L105}{\nolinkurl{tsum_lambert_linear_weight_sq_pure}}.  \emph{Compared}, run \href{https://github.com/wcook04/plectis-erdos-lean/actions/runs/35715779405}{\texttt{35715779405}}.
\par\noindent\hangindent=1.5em Proposition~\ref{prop:pillai}: \href{https://github.com/wcook04/plectis-erdos/blob/181078b6b009d905809cf2e007ad309386a3d1e0/lean/ErdosProblems/Erdos249/PaperCompleteR21/PillaiGcdExpectation.lean\#L334}{\nolinkurl{gcd_moment_identity_three_members}}, \href{https://github.com/wcook04/plectis-erdos/blob/181078b6b009d905809cf2e007ad309386a3d1e0/lean/ErdosProblems/Erdos249/PaperCompleteR21/PillaiGcdExpectation.lean\#L307}{\nolinkurl{tsum_pos_pair_gcd_half_eq_totient_div_mersenne_sq}}, \href{https://github.com/wcook04/plectis-erdos/blob/181078b6b009d905809cf2e007ad309386a3d1e0/lean/ErdosProblems/Erdos249/PaperCompleteR21/PillaiGcdExpectation.lean\#L77}{\nolinkurl{sum_gcd_Icc_eq_pillaiP}}, \href{https://github.com/wcook04/plectis-erdos/blob/181078b6b009d905809cf2e007ad309386a3d1e0/lean/ErdosProblems/Erdos249/PaperCompleteR21/PillaiGcdExpectation.lean\#L67}{\nolinkurl{pillaiP_eq_totient_mul_id}}, \href{https://github.com/wcook04/plectis-erdos/blob/181078b6b009d905809cf2e007ad309386a3d1e0/lean/Erdos249257/GcdMomentCalculus.lean\#L235}{\nolinkurl{tsum_totient_div_mersenne_sq_eq_gcd_moment_series}}.  \emph{Compared}, runs \href{https://github.com/wcook04/plectis-erdos-lean/actions/runs/35643815458}{\texttt{35643815458}}, \href{https://github.com/wcook04/plectis-erdos-lean/actions/runs/35715779405}{\texttt{35715779405}}.
\par\noindent\hangindent=1.5em Theorem~\ref{catalogue:cert:a9}: \href{https://github.com/wcook04/plectis-erdos/blob/181078b6b009d905809cf2e007ad309386a3d1e0/lean/ErdosProblems/Erdos249/PaperCompleteR20/SpecifiedTailPeriod.lean\#L44}{\nolinkurl{specified_euler_tail_period}}.  \emph{Compared}, run \href{https://github.com/wcook04/plectis-erdos-lean/actions/runs/35674034595}{\texttt{35674034595}}.
\par\noindent\hangindent=1.5em Theorem~\ref{catalogue:cert:d5}: \href{https://github.com/wcook04/plectis-erdos/blob/181078b6b009d905809cf2e007ad309386a3d1e0/lean/Erdos249257/TotientCarryKernelRigidity.lean\#L300}{\nolinkurl{not_irrational_totientSeries_implies_unbounded_carryRank_unconditional}}.  \emph{Compared}, run \href{https://github.com/wcook04/plectis-erdos-lean/actions/runs/35544127144}{\texttt{35544127144}}.
\par\noindent\hangindent=1.5em Theorem~\ref{catalogue:cert:d4}: \href{https://github.com/wcook04/plectis-erdos/blob/181078b6b009d905809cf2e007ad309386a3d1e0/lean/ErdosProblems/Erdos249/PaperCompleteR8/FullKernelAssemblies.lean\#L177}{\nolinkurl{displayed_canonical_and_full_dyadic}}.  \emph{Compared}, run \href{https://github.com/wcook04/plectis-erdos-lean/actions/runs/35605146048}{\texttt{35605146048}}.
\par\noindent\hangindent=1.5em Proposition~\ref{catalogue:cert:a8}: \href{https://github.com/wcook04/plectis-erdos/blob/181078b6b009d905809cf2e007ad309386a3d1e0/lean/Erdos249257/TotientTailPeriodKiller.lean\#L327}{\nolinkurl{tail_diff_int_of_den_dvd}}.  \emph{Compared}, run \href{https://github.com/wcook04/plectis-erdos-lean/actions/runs/35643815458}{\texttt{35643815458}}.
\par\noindent\hangindent=1.5em Proposition~\ref{catalogue:mob:b5}: \href{https://github.com/wcook04/plectis-erdos/blob/181078b6b009d905809cf2e007ad309386a3d1e0/lean/Erdos249257/RepunitMobiusNumerator.lean\#L445}{\nolinkurl{mobiusNumeratorPolynomial_eval_two}}, \href{https://github.com/wcook04/plectis-erdos/blob/181078b6b009d905809cf2e007ad309386a3d1e0/lean/Erdos249257/CyclotomicProjectionOfShadow.lean\#L350}{\nolinkurl{mobiusNumerator_gcd_cyclotomicValue}}, \href{https://github.com/wcook04/plectis-erdos/blob/181078b6b009d905809cf2e007ad309386a3d1e0/lean/Erdos249257/CyclotomicProjectionOfShadow.lean\#L389}{\nolinkurl{cyclotomicValue_dvd_baseMobiusShadow_den}}.  \emph{Compared}, run \href{https://github.com/wcook04/plectis-erdos-lean/actions/runs/35624228171}{\texttt{35624228171}}.
\par\noindent\hangindent=1.5em Proposition~\ref{catalogue:mob:b6}: \href{https://github.com/wcook04/plectis-erdos/blob/181078b6b009d905809cf2e007ad309386a3d1e0/lean/Erdos249257/MersenneShadowCyclotomicNoncollapse.lean\#L983}{\nolinkurl{lcmHeight_upperHalf_product_dvd_den}}, \href{https://github.com/wcook04/plectis-erdos/blob/181078b6b009d905809cf2e007ad309386a3d1e0/lean/Erdos249257/MersenneShadowCyclotomicNoncollapse.lean\#L795}{\nolinkurl{upperHalfChannel_product_dvd_den_of_coprime_scale}}, \href{https://github.com/wcook04/plectis-erdos/blob/181078b6b009d905809cf2e007ad309386a3d1e0/lean/Erdos249257/MersenneShadowCyclotomicNoncollapse.lean\#L809}{\nolinkurl{upperHalfChannel_product_dvd_den_of_scale_primeFactors_le}}.  \emph{Compared}, run \href{https://github.com/wcook04/plectis-erdos-lean/actions/runs/35643815458}{\texttt{35643815458}}.
\par\noindent\hangindent=1.5em Proposition~\ref{catalogue:mob:b7a}: \href{https://github.com/wcook04/plectis-erdos/blob/181078b6b009d905809cf2e007ad309386a3d1e0/lean/ErdosProblems/Erdos249/PaperCompleteR20/DenominatorBounds.lean\#L10}{\nolinkurl{upper_half_product_denominator_bounds}}.  \emph{Compared}, run \href{https://github.com/wcook04/plectis-erdos-lean/actions/runs/35643815458}{\texttt{35643815458}}.
\par\noindent\hangindent=1.5em Proposition~\ref{catalogue:mob:b7b}: \href{https://github.com/wcook04/plectis-erdos/blob/181078b6b009d905809cf2e007ad309386a3d1e0/lean/Erdos249257/MersenneShadowDenominatorGrowth.lean\#L147}{\nolinkurl{lcmHeight_scaledMobiusShadow_den_exact}}, \href{https://github.com/wcook04/plectis-erdos/blob/181078b6b009d905809cf2e007ad309386a3d1e0/lean/Erdos249257/MersenneShadowDenominatorGrowth.lean\#L187}{\nolinkurl{lcmHeight_five_scaledMobiusShadow_den_exact}}.  \emph{Compared}, run \href{https://github.com/wcook04/plectis-erdos-lean/actions/runs/35643815458}{\texttt{35643815458}}.
\par\noindent\hangindent=1.5em Proposition~\ref{catalogue:mob:d3}: \href{https://github.com/wcook04/plectis-erdos/blob/181078b6b009d905809cf2e007ad309386a3d1e0/lean/ErdosProblems/Erdos249/PaperCompleteR20/SignedDyadicClearing.lean\#L8}{\nolinkurl{signed_dyadic_clearing}}, \href{https://github.com/wcook04/plectis-erdos/blob/181078b6b009d905809cf2e007ad309386a3d1e0/lean/ErdosProblems/Erdos249/PaperCompleteR20/SignedDyadicClearing.lean\#L28}{\nolinkurl{signed_dyadic_sum_ne_zero}}, \href{https://github.com/wcook04/plectis-erdos/blob/181078b6b009d905809cf2e007ad309386a3d1e0/lean/Erdos249257/SignedQMomentObstruction.lean\#L78}{\nolinkurl{scaled_dyadic_sum_odd}}.  \emph{Compared}, runs \href{https://github.com/wcook04/plectis-erdos-lean/actions/runs/35624228171}{\texttt{35624228171}}, \href{https://github.com/wcook04/plectis-erdos-lean/actions/runs/35643815458}{\texttt{35643815458}}.
\par\noindent\hangindent=1.5em Proposition~\ref{catalogue:cert:a11}: \href{https://github.com/wcook04/plectis-erdos/blob/181078b6b009d905809cf2e007ad309386a3d1e0/lean/ErdosProblems/Erdos249/PaperCompleteR20/FiniteCertificateBatch.lean\#L33}{\nolinkurl{small_certificate_windows}}, \href{https://github.com/wcook04/plectis-erdos/blob/181078b6b009d905809cf2e007ad309386a3d1e0/lean/ErdosProblems/Erdos249/PaperCompleteR20/FiniteCertificateBatch.lean\#L45}{\nolinkurl{small_certificates_and_exclusions}}.  \emph{Compared}, run \href{https://github.com/wcook04/plectis-erdos-lean/actions/runs/35674034595}{\texttt{35674034595}}.
\par\noindent\hangindent=1.5em Proposition~\ref{catalogue:cert:a12}: \href{https://github.com/wcook04/plectis-erdos/blob/181078b6b009d905809cf2e007ad309386a3d1e0/lean/ErdosProblems/Erdos249/PaperCompleteR20/FiniteCertificateBatch.lean\#L39}{\nolinkurl{sixteen_certificate_windows}}, \href{https://github.com/wcook04/plectis-erdos/blob/181078b6b009d905809cf2e007ad309386a3d1e0/lean/ErdosProblems/Erdos249/PaperCompleteR20/FiniteCertificateBatch.lean\#L51}{\nolinkurl{sixteen_certificates_and_exclusions}}.  \emph{Compared}, run \href{https://github.com/wcook04/plectis-erdos-lean/actions/runs/35674034595}{\texttt{35674034595}}.
\par\noindent\hangindent=1.5em Proposition~\ref{catalogue:cert:b11}: \href{https://github.com/wcook04/plectis-erdos/blob/181078b6b009d905809cf2e007ad309386a3d1e0/lean/ErdosProblems/Erdos249/PaperCompleteR20/FiniteCertificateBatch.lean\#L57}{\nolinkurl{historical_table_size_and_initial_depths}}, \href{https://github.com/wcook04/plectis-erdos/blob/181078b6b009d905809cf2e007ad309386a3d1e0/lean/ErdosProblems/Erdos249/PaperCompleteR20/FiniteCertificateBatch.lean\#L64}{\nolinkurl{historical_table_and_complete_band}}.  \emph{Compared}, run \href{https://github.com/wcook04/plectis-erdos-lean/actions/runs/35674034595}{\texttt{35674034595}}.
\par\noindent\hangindent=1.5em Proposition~\ref{catalogue:cert:a2}: \href{https://github.com/wcook04/plectis-erdos/blob/181078b6b009d905809cf2e007ad309386a3d1e0/lean/ErdosProblems/Erdos249/PaperCompleteR20/TailDepthCorrespondence.lean\#L8}{\nolinkurl{prefix_fractional_part}}, \href{https://github.com/wcook04/plectis-erdos/blob/181078b6b009d905809cf2e007ad309386a3d1e0/lean/Erdos249257/TotientTailPeriodKiller.lean\#L150}{\nolinkurl{two_pow_mul_totient_series_eq}}, \href{https://github.com/wcook04/plectis-erdos/blob/181078b6b009d905809cf2e007ad309386a3d1e0/lean/Erdos249257/LcmConeFlatness.lean\#L327}{\nolinkurl{tail_diff_mem_int_iff_scaled_series_mem_int}}.  \emph{Compared}, runs \href{https://github.com/wcook04/plectis-erdos-lean/actions/runs/35643815458}{\texttt{35643815458}}, \href{https://github.com/wcook04/plectis-erdos-lean/actions/runs/35674034595}{\texttt{35674034595}}.
\par\noindent\hangindent=1.5em Lemma~\ref{catalogue:cert:a5}: \href{https://github.com/wcook04/plectis-erdos/blob/181078b6b009d905809cf2e007ad309386a3d1e0/lean/ErdosProblems/Erdos249/PaperCompleteR20/TailDepthCorrespondence.lean\#L20}{\nolinkurl{certificate_logarithmic_depth}}, \href{https://github.com/wcook04/plectis-erdos/blob/181078b6b009d905809cf2e007ad309386a3d1e0/lean/ErdosProblems/Erdos249/PaperCompleteR20/TailDepthCorrespondence.lean\#L30}{\nolinkurl{fixed_depth_bounds_indices}}, \href{https://github.com/wcook04/plectis-erdos/blob/181078b6b009d905809cf2e007ad309386a3d1e0/lean/Erdos249257/TotientTailPeriodKiller.lean\#L79}{\nolinkurl{certifiedKill_depth_floor}}.  \emph{Compared}, runs \href{https://github.com/wcook04/plectis-erdos-lean/actions/runs/35643815458}{\texttt{35643815458}}, \href{https://github.com/wcook04/plectis-erdos-lean/actions/runs/35674034595}{\texttt{35674034595}}.
\par\noindent\hangindent=1.5em Proposition~\ref{catalogue:cert:a6}: \href{https://github.com/wcook04/plectis-erdos/blob/181078b6b009d905809cf2e007ad309386a3d1e0/lean/ErdosProblems/Erdos249/PaperCompleteR20/TailDepthCorrespondence.lean\#L13}{\nolinkurl{totient_scaled_truncation_error}}, \href{https://github.com/wcook04/plectis-erdos/blob/181078b6b009d905809cf2e007ad309386a3d1e0/lean/Erdos249257/TotientTailPeriodKiller.lean\#L262}{\nolinkurl{tail_diff_notMem_int_of_certifiedKill}}.  \emph{Compared}, runs \href{https://github.com/wcook04/plectis-erdos-lean/actions/runs/35643815458}{\texttt{35643815458}}, \href{https://github.com/wcook04/plectis-erdos-lean/actions/runs/35674034595}{\texttt{35674034595}}.
\par\noindent\hangindent=1.5em Theorem~\ref{catalogue:cert:a7}: \href{https://github.com/wcook04/plectis-erdos/blob/181078b6b009d905809cf2e007ad309386a3d1e0/lean/Erdos249257/LcmConeFlatness.lean\#L316}{\nolinkurl{exists_certifiedKill_iff_tail_diff_notMem_int}}.  \emph{Compared}, run \href{https://github.com/wcook04/plectis-erdos-lean/actions/runs/35643815458}{\texttt{35643815458}}.
\par\noindent\hangindent=1.5em Proposition~\ref{catalogue:cert:c1}: \href{https://github.com/wcook04/plectis-erdos/blob/181078b6b009d905809cf2e007ad309386a3d1e0/lean/Erdos249257/GapFareyBound.lean\#L51}{\nolinkurl{farey_gap}}.  \emph{Compared}, run \href{https://github.com/wcook04/plectis-erdos-lean/actions/runs/35715779405}{\texttt{35715779405}}.
\par\noindent\hangindent=1.5em Proposition~\ref{catalogue:cert:d1}: \href{https://github.com/wcook04/plectis-erdos/blob/181078b6b009d905809cf2e007ad309386a3d1e0/lean/Erdos249257/CertificateKernel.lean\#L5371}{\nolinkurl{irrational_of_den_mul_abs_sub_tendsto_zero}}.  \emph{Compared}, run \href{https://github.com/wcook04/plectis-erdos-lean/actions/runs/35624228171}{\texttt{35624228171}}.
\par\noindent\hangindent=1.5em Proposition~\ref{catalogue:cert:d2} (the Lean statement is stronger): \href{https://github.com/wcook04/plectis-erdos/blob/181078b6b009d905809cf2e007ad309386a3d1e0/lean/Erdos249257/CertificateKernel.lean\#L6120}{\nolinkurl{irrational_of_int_mul_near_int}}, \href{https://github.com/wcook04/plectis-erdos/blob/181078b6b009d905809cf2e007ad309386a3d1e0/lean/Erdos249257/CertificateKernel.lean\#L6149}{\nolinkurl{irrational_of_pow_mul_near_int}}, \href{https://github.com/wcook04/plectis-erdos/blob/181078b6b009d905809cf2e007ad309386a3d1e0/lean/ErdosProblems/Erdos249/PaperCompleteR21/NearIntegerIrrationalityCriterion.lean\#L16}{\nolinkurl{irrational_of_near_integer_multiples}}, \href{https://github.com/wcook04/plectis-erdos/blob/181078b6b009d905809cf2e007ad309386a3d1e0/lean/ErdosProblems/Erdos249/PaperCompleteR21/NearIntegerIrrationalityCriterion.lean\#L25}{\nolinkurl{irrational_of_near_integer_base_powers}}, \href{https://github.com/wcook04/plectis-erdos/blob/181078b6b009d905809cf2e007ad309386a3d1e0/lean/ErdosProblems/Erdos249/PaperCompleteR21/GeneralIrrationalityCriteriaAndGapBounds.lean\#L33}{\nolinkurl{one_div_den_le_abs_int_combination}}, \href{https://github.com/wcook04/plectis-erdos/blob/181078b6b009d905809cf2e007ad309386a3d1e0/lean/ErdosProblems/Erdos249/PaperCompleteR21/SquareBlockBinaryDilationCountermodel.lean\#L393}{\nolinkurl{digit_block}}, \href{https://github.com/wcook04/plectis-erdos/blob/181078b6b009d905809cf2e007ad309386a3d1e0/lean/ErdosProblems/Erdos249/PaperCompleteR21/SquareBlockBinaryDilationCountermodel.lean\#L424}{\nolinkurl{not_eventually_periodic}}, \href{https://github.com/wcook04/plectis-erdos/blob/181078b6b009d905809cf2e007ad309386a3d1e0/lean/ErdosProblems/Erdos249/PaperCompleteR21/SquareBlockBinaryDilationCountermodel.lean\#L476}{\nolinkurl{irrational_xi}}, \href{https://github.com/wcook04/plectis-erdos/blob/181078b6b009d905809cf2e007ad309386a3d1e0/lean/ErdosProblems/Erdos249/PaperCompleteR21/SquareBlockBinaryDilationCountermodel.lean\#L408}{\nolinkurl{no_three_consecutive_equal}}, \href{https://github.com/wcook04/plectis-erdos/blob/181078b6b009d905809cf2e007ad309386a3d1e0/lean/ErdosProblems/Erdos249/PaperCompleteR21/SquareBlockBinaryDilationCountermodel.lean\#L480}{\nolinkurl{fract_two_pow_mul_xi_mem_Icc}}, \href{https://github.com/wcook04/plectis-erdos/blob/181078b6b009d905809cf2e007ad309386a3d1e0/lean/ErdosProblems/Erdos249/PaperCompleteR21/SquareBlockBinaryDilationCountermodel.lean\#L235}{\nolinkurl{fract_mem_Icc_of_no_three_equal}}, \href{https://github.com/wcook04/plectis-erdos/blob/181078b6b009d905809cf2e007ad309386a3d1e0/lean/ErdosProblems/Erdos249/PaperCompleteR21/SquareBlockBinaryDilationCountermodel.lean\#L486}{\nolinkurl{one_div_eight_le_dist_xi}}, \href{https://github.com/wcook04/plectis-erdos/blob/181078b6b009d905809cf2e007ad309386a3d1e0/lean/ErdosProblems/Erdos249/PaperCompleteR21/SquareBlockBinaryDilationCountermodel.lean\#L492}{\nolinkurl{not_approaches_int}}, \href{https://github.com/wcook04/plectis-erdos/blob/181078b6b009d905809cf2e007ad309386a3d1e0/lean/ErdosProblems/Erdos249/PaperCompleteR21/SquareBlockBinaryDilationCountermodel.lean\#L502}{\nolinkurl{not_near_integer_along_powers_of_two}}, \href{https://github.com/wcook04/plectis-erdos/blob/181078b6b009d905809cf2e007ad309386a3d1e0/lean/ErdosProblems/Erdos249/PaperCompleteR21/SquareBlockBinaryDilationCountermodel.lean\#L515}{\nolinkurl{exists_irrational_basePower_bounded_away}}, \href{https://github.com/wcook04/plectis-erdos/blob/181078b6b009d905809cf2e007ad309386a3d1e0/lean/ErdosProblems/Erdos249/PaperCompleteR21/SquareBlockBinaryDilationCountermodel.lean\#L550}{\nolinkurl{strict_lower_bound_needed}}, \href{https://github.com/wcook04/plectis-erdos/blob/181078b6b009d905809cf2e007ad309386a3d1e0/lean/ErdosProblems/Erdos249/PaperCompleteR21/SquareBlockBinaryDilationCountermodel.lean\#L566}{\nolinkurl{upper_bound_needed_for_every_q}}, \href{https://github.com/wcook04/plectis-erdos/blob/181078b6b009d905809cf2e007ad309386a3d1e0/lean/ErdosProblems/Erdos249/PaperCompleteR21/SquareBlockBinaryDilationCountermodel.lean\#L123}{\nolinkurl{tail_mem_Icc}}, \href{https://github.com/wcook04/plectis-erdos/blob/181078b6b009d905809cf2e007ad309386a3d1e0/lean/ErdosProblems/Erdos249/PaperCompleteR21/SquareBlockBinaryDilationCountermodel.lean\#L269}{\nolinkurl{irrational_tail_zero}}, \href{https://github.com/wcook04/plectis-erdos/blob/181078b6b009d905809cf2e007ad309386a3d1e0/lean/ErdosProblems/Erdos249/PaperCompleteR21/SquareBlockBinaryDilationCountermodel.lean\#L243}{\nolinkurl{one_div_eight_le_abs_sub_int}}.  \emph{Compared}, runs \href{https://github.com/wcook04/plectis-erdos-lean/actions/runs/35624228171}{\texttt{35624228171}}, \href{https://github.com/wcook04/plectis-erdos-lean/actions/runs/35643815458}{\texttt{35643815458}}.
\par\noindent\hangindent=1.5em Proposition~\ref{catalogue:cert:d3}: \href{https://github.com/wcook04/plectis-erdos/blob/181078b6b009d905809cf2e007ad309386a3d1e0/lean/Erdos249257/TotientMahlerDefect.lean\#L91}{\nolinkurl{linearIndependent_of_separatedMinorCertificate}}.  \emph{Compared}, run \href{https://github.com/wcook04/plectis-erdos-lean/actions/runs/35682858662}{\texttt{35682858662}}.
\par\noindent\hangindent=1.5em Proposition~\ref{catalogue:cert:d9}: \href{https://github.com/wcook04/plectis-erdos/blob/181078b6b009d905809cf2e007ad309386a3d1e0/lean/ErdosProblems/Erdos249/PaperCompleteR20/RationalSpacingCorrespondence.lean\#L6}{\nolinkurl{rational_difference_exact}}, \href{https://github.com/wcook04/plectis-erdos/blob/181078b6b009d905809cf2e007ad309386a3d1e0/lean/ErdosProblems/Erdos249/PaperCompleteR20/RationalSpacingCorrespondence.lean\#L17}{\nolinkurl{rational_cross_numerator_positive}}, \href{https://github.com/wcook04/plectis-erdos/blob/181078b6b009d905809cf2e007ad309386a3d1e0/lean/ErdosProblems/Erdos249/PaperCompleteR20/RationalSpacingCorrespondence.lean\#L27}{\nolinkurl{rational_error_denominator_bound}}, \href{https://github.com/wcook04/plectis-erdos/blob/181078b6b009d905809cf2e007ad309386a3d1e0/lean/Erdos249257/PrimitiveRationalGapSupply.lean\#L31}{\nolinkurl{positive_rational_difference_lower_bound}}.  \emph{Compared}, run \href{https://github.com/wcook04/plectis-erdos-lean/actions/runs/35643815458}{\texttt{35643815458}}.
\par\noindent\hangindent=1.5em Proposition~\ref{catalogue:mob:a1a}: \href{https://github.com/wcook04/plectis-erdos/blob/181078b6b009d905809cf2e007ad309386a3d1e0/lean/Erdos249257/MersenneLambertLadder.lean\#L491}{\nolinkurl{tsum_moebius_lambert_sq}}, \href{https://github.com/wcook04/plectis-erdos/blob/181078b6b009d905809cf2e007ad309386a3d1e0/lean/Erdos249257/CertificateKernel.lean\#L18454}{\nolinkurl{totient_series_eq_half_add_moebius_mersenne_square}}.  \emph{Compared}, runs \href{https://github.com/wcook04/plectis-erdos-lean/actions/runs/35624228171}{\texttt{35624228171}}, \href{https://github.com/wcook04/plectis-erdos-lean/actions/runs/35715779405}{\texttt{35715779405}}.
\par\noindent\hangindent=1.5em Corollary~\ref{catalogue:mob:a1b}: \href{https://github.com/wcook04/plectis-erdos/blob/181078b6b009d905809cf2e007ad309386a3d1e0/lean/ErdosProblems/Erdos249/PaperCompleteR7/ArithmeticAssemblies.lean\#L40}{\nolinkurl{irrational_totient_iff_moebius_square}}.  \emph{Compared}, run \href{https://github.com/wcook04/plectis-erdos-lean/actions/runs/35682858662}{\texttt{35682858662}}.
\par\noindent\hangindent=1.5em Proposition~\ref{catalogue:mob:a2}: \href{https://github.com/wcook04/plectis-erdos/blob/181078b6b009d905809cf2e007ad309386a3d1e0/lean/Erdos249257/GcdMomentCalculus.lean\#L105}{\nolinkurl{tsum_lambert_linear_weight_sq_pure}}.  \emph{Compared}, run \href{https://github.com/wcook04/plectis-erdos-lean/actions/runs/35715779405}{\texttt{35715779405}}.
\par\noindent\hangindent=1.5em Proposition~\ref{catalogue:mob:a5}: \href{https://github.com/wcook04/plectis-erdos/blob/181078b6b009d905809cf2e007ad309386a3d1e0/lean/ErdosProblems/Erdos249/PaperCompleteR21/PillaiGcdExpectation.lean\#L334}{\nolinkurl{gcd_moment_identity_three_members}}, \href{https://github.com/wcook04/plectis-erdos/blob/181078b6b009d905809cf2e007ad309386a3d1e0/lean/ErdosProblems/Erdos249/PaperCompleteR21/PillaiGcdExpectation.lean\#L307}{\nolinkurl{tsum_pos_pair_gcd_half_eq_totient_div_mersenne_sq}}, \href{https://github.com/wcook04/plectis-erdos/blob/181078b6b009d905809cf2e007ad309386a3d1e0/lean/ErdosProblems/Erdos249/PaperCompleteR21/PillaiGcdExpectation.lean\#L67}{\nolinkurl{pillaiP_eq_totient_mul_id}}, \href{https://github.com/wcook04/plectis-erdos/blob/181078b6b009d905809cf2e007ad309386a3d1e0/lean/ErdosProblems/Erdos249/PaperCompleteR21/PillaiGcdExpectation.lean\#L77}{\nolinkurl{sum_gcd_Icc_eq_pillaiP}}, \href{https://github.com/wcook04/plectis-erdos/blob/181078b6b009d905809cf2e007ad309386a3d1e0/lean/Erdos249257/GcdMomentCalculus.lean\#L235}{\nolinkurl{tsum_totient_div_mersenne_sq_eq_gcd_moment_series}}.  \emph{Compared}, runs \href{https://github.com/wcook04/plectis-erdos-lean/actions/runs/35643815458}{\texttt{35643815458}}, \href{https://github.com/wcook04/plectis-erdos-lean/actions/runs/35715779405}{\texttt{35715779405}}.
\par\noindent\hangindent=1.5em Proposition~\ref{catalogue:mob:a7}: \href{https://github.com/wcook04/plectis-erdos/blob/181078b6b009d905809cf2e007ad309386a3d1e0/lean/ErdosProblems/Erdos249/PaperCompleteR21/SquaredMersenneDivisorIdentities.lean\#L62}{\nolinkurl{pair_divisibility_mass}}, \href{https://github.com/wcook04/plectis-erdos/blob/181078b6b009d905809cf2e007ad309386a3d1e0/lean/ErdosProblems/Erdos249/PaperCompleteR21/SquaredMersenneDivisorIdentities.lean\#L79}{\nolinkurl{tsum_geometric_multiples}}, \href{https://github.com/wcook04/plectis-erdos/blob/181078b6b009d905809cf2e007ad309386a3d1e0/lean/Erdos249257/GcdMomentCalculus.lean\#L266}{\nolinkurl{tsum_pos_pair_both_dvd_half_eq_inv_mersenne_sq}}.  \emph{Compared}, runs \href{https://github.com/wcook04/plectis-erdos-lean/actions/runs/35643815458}{\texttt{35643815458}}, \href{https://github.com/wcook04/plectis-erdos-lean/actions/runs/35715779405}{\texttt{35715779405}}.
\par\noindent\hangindent=1.5em Proposition~\ref{catalogue:mob:a8}: \href{https://github.com/wcook04/plectis-erdos/blob/181078b6b009d905809cf2e007ad309386a3d1e0/lean/Erdos249257/GcdMomentCalculus.lean\#L349}{\nolinkurl{tsum_pos_coprime_inv_mersenne_eq_one}}.  \emph{Compared}, run \href{https://github.com/wcook04/plectis-erdos-lean/actions/runs/35715779405}{\texttt{35715779405}}.
\par\noindent\hangindent=1.5em Proposition~\ref{catalogue:mob:a9a}: \href{https://github.com/wcook04/plectis-erdos/blob/181078b6b009d905809cf2e007ad309386a3d1e0/lean/ErdosProblems/Erdos249/PaperCompleteR21/SternBrocotStoppingRecursion.lean\#L42}{\nolinkurl{divisibility_mass}}, \href{https://github.com/wcook04/plectis-erdos/blob/181078b6b009d905809cf2e007ad309386a3d1e0/lean/ErdosProblems/Erdos249/PaperCompleteR21/SternBrocotStoppingRecursion.lean\#L86}{\nolinkurl{cylinderMass_eq_divisibility_mass_mul}}, \href{https://github.com/wcook04/plectis-erdos/blob/181078b6b009d905809cf2e007ad309386a3d1e0/lean/ErdosProblems/Erdos249/PaperCompleteR21/SternBrocotStoppingRecursion.lean\#L95}{\nolinkurl{cylinder_mediant_split}}, \href{https://github.com/wcook04/plectis-erdos/blob/181078b6b009d905809cf2e007ad309386a3d1e0/lean/ErdosProblems/Erdos249/PaperCompleteR21/SternBrocotStoppingRecursion.lean\#L103}{\nolinkurl{cylinder_root_values}}, \href{https://github.com/wcook04/plectis-erdos/blob/181078b6b009d905809cf2e007ad309386a3d1e0/lean/ErdosProblems/Erdos249/PaperCompleteR21/SternBrocotStoppingRecursion.lean\#L117}{\nolinkurl{normalised_split_probabilities}}, \href{https://github.com/wcook04/plectis-erdos/blob/181078b6b009d905809cf2e007ad309386a3d1e0/lean/ErdosProblems/Erdos249/PaperCompleteR21/SternBrocotStoppingRecursion.lean\#L147}{\nolinkurl{stopping_transition_probabilities_sum_one}}, \href{https://github.com/wcook04/plectis-erdos/blob/181078b6b009d905809cf2e007ad309386a3d1e0/lean/ErdosProblems/Erdos249/PaperCompleteR21/SternBrocotStoppingRecursion.lean\#L162}{\nolinkurl{stopping_probability_ge_third}}.  \emph{Compared}, runs \href{https://github.com/wcook04/plectis-erdos-lean/actions/runs/35624228171}{\texttt{35624228171}}, \href{https://github.com/wcook04/plectis-erdos-lean/actions/runs/35643815458}{\texttt{35643815458}}.
\par\noindent\hangindent=1.5em Proposition~\ref{catalogue:mob:a9b}: \href{https://github.com/wcook04/plectis-erdos/blob/181078b6b009d905809cf2e007ad309386a3d1e0/lean/Erdos249257/GcdMomentCalculus.lean\#L514}{\nolinkurl{cylinderMass_children_le}}, \href{https://github.com/wcook04/plectis-erdos/blob/181078b6b009d905809cf2e007ad309386a3d1e0/lean/Erdos249257/GcdMomentCalculus.lean\#L525}{\nolinkurl{sternBrocotDepthMass_error}}, \href{https://github.com/wcook04/plectis-erdos/blob/181078b6b009d905809cf2e007ad309386a3d1e0/lean/Erdos249257/GcdMomentCalculus.lean\#L560}{\nolinkurl{tendsto_sternBrocotDepthMass}}.  \emph{Compared}, run \href{https://github.com/wcook04/plectis-erdos-lean/actions/runs/35715779405}{\texttt{35715779405}}.
\par\noindent\hangindent=1.5em Proposition~\ref{catalogue:mob:b1}: \href{https://github.com/wcook04/plectis-erdos/blob/181078b6b009d905809cf2e007ad309386a3d1e0/lean/Erdos249257/RepunitMobiusNumerator.lean\#L217}{\nolinkurl{mobiusNumeratorPolynomial_coeff}}, \href{https://github.com/wcook04/plectis-erdos/blob/181078b6b009d905809cf2e007ad309386a3d1e0/lean/Erdos249257/RepunitMobiusNumerator.lean\#L234}{\nolinkurl{mobiusNumeratorPolynomial_coeff_pos}}.  \emph{Compared}, run \href{https://github.com/wcook04/plectis-erdos-lean/actions/runs/35624228171}{\texttt{35624228171}}.
\par\noindent\hangindent=1.5em Proposition~\ref{catalogue:mob:b2}: \href{https://github.com/wcook04/plectis-erdos/blob/181078b6b009d905809cf2e007ad309386a3d1e0/lean/Erdos249257/RepunitMobiusNumerator.lean\#L445}{\nolinkurl{mobiusNumeratorPolynomial_eval_two}}, \href{https://github.com/wcook04/plectis-erdos/blob/181078b6b009d905809cf2e007ad309386a3d1e0/lean/ErdosProblems/Erdos249/PaperCompleteR20/NumeratorEvaluation.lean\#L32}{\nolinkurl{numerator_eval_two_divisors}}.  \emph{Compared}, run \href{https://github.com/wcook04/plectis-erdos-lean/actions/runs/35624228171}{\texttt{35624228171}}.
\par\noindent\hangindent=1.5em Proposition~\ref{catalogue:mob:b3}: \href{https://github.com/wcook04/plectis-erdos/blob/181078b6b009d905809cf2e007ad309386a3d1e0/lean/ErdosProblems/Erdos249/PaperCompleteR20/RadicalDecomposition.lean\#L50}{\nolinkurl{radical_decomposition}}.  \emph{Compared}, run \href{https://github.com/wcook04/plectis-erdos-lean/actions/runs/35624228171}{\texttt{35624228171}}.
\par\noindent\hangindent=1.5em Proposition~\ref{catalogue:mob:b4}: \href{https://github.com/wcook04/plectis-erdos/blob/181078b6b009d905809cf2e007ad309386a3d1e0/lean/Erdos249257/CyclotomicProjectionOfShadow.lean\#L299}{\nolinkurl{cyclotomic_dvd_mobiusNumeratorPolynomial_sub}}, \href{https://github.com/wcook04/plectis-erdos/blob/181078b6b009d905809cf2e007ad309386a3d1e0/lean/Erdos249257/CyclotomicProjectionOfShadow.lean\#L310}{\nolinkurl{mobiusNumerator_mod_cyclotomicEval}}, \href{https://github.com/wcook04/plectis-erdos/blob/181078b6b009d905809cf2e007ad309386a3d1e0/lean/Erdos249257/CyclotomicProjectionOfShadow.lean\#L67}{\nolinkurl{jordanTotientTwo_eq_prod_primeFactors}}.  \emph{Compared}, run \href{https://github.com/wcook04/plectis-erdos-lean/actions/runs/35624228171}{\texttt{35624228171}}.
\par\noindent\hangindent=1.5em Proposition~\ref{catalogue:mob:b8a}: \href{https://github.com/wcook04/plectis-erdos/blob/181078b6b009d905809cf2e007ad309386a3d1e0/lean/Erdos249257/PrimePowerJumpDynamics.lean\#L281}{\nolinkurl{cyclotomic_dvd_primeJump_new_fibre}}.  \emph{Compared}, run \href{https://github.com/wcook04/plectis-erdos-lean/actions/runs/35624228171}{\texttt{35624228171}}.
\par\noindent\hangindent=1.5em Proposition~\ref{catalogue:mob:b8b}: \href{https://github.com/wcook04/plectis-erdos/blob/181078b6b009d905809cf2e007ad309386a3d1e0/lean/Erdos249257/PrimePowerJumpDynamics.lean\#L310}{\nolinkurl{cyclotomic_dvd_primeJump_old_fibre}}.  \emph{Compared}, run \href{https://github.com/wcook04/plectis-erdos-lean/actions/runs/35624228171}{\texttt{35624228171}}.
\par\noindent\hangindent=1.5em Proposition~\ref{catalogue:cert:d7}: \href{https://github.com/wcook04/plectis-erdos/blob/181078b6b009d905809cf2e007ad309386a3d1e0/lean/ErdosProblems/Erdos249/PaperCompleteR21/LambertDivisorTransform.lean\#L42}{\nolinkurl{lambertValue_eq_divisor_sum_series}}, \href{https://github.com/wcook04/plectis-erdos/blob/181078b6b009d905809cf2e007ad309386a3d1e0/lean/ErdosProblems/Erdos249/PaperCompleteR21/LambertDivisorTransform.lean\#L101}{\nolinkurl{alpha_divisor_sum_eq_totient}}, \href{https://github.com/wcook04/plectis-erdos/blob/181078b6b009d905809cf2e007ad309386a3d1e0/lean/ErdosProblems/Erdos249/PaperCompleteR21/LambertDivisorTransform.lean\#L109}{\nolinkurl{lambertValue_alpha_eq_totientSeries}}.  \emph{Compared}, runs \href{https://github.com/wcook04/plectis-erdos-lean/actions/runs/35624228171}{\texttt{35624228171}}, \href{https://github.com/wcook04/plectis-erdos-lean/actions/runs/35643815458}{\texttt{35643815458}}.
\par\noindent\hangindent=1.5em Proposition~\ref{catalogue:mob:e3}: \href{https://github.com/wcook04/plectis-erdos/blob/181078b6b009d905809cf2e007ad309386a3d1e0/lean/ErdosProblems/Erdos249/PaperCompleteR21/AffineDivisorAnnihilation.lean\#L28}{\nolinkurl{transportResidueKernel_eq_mobiusTermKernel}}, \href{https://github.com/wcook04/plectis-erdos/blob/181078b6b009d905809cf2e007ad309386a3d1e0/lean/ErdosProblems/Erdos249/PaperCompleteR21/AffineDivisorAnnihilation.lean\#L47}{\nolinkurl{mobiusTermKernel_affine_in_multiplier}}, \href{https://github.com/wcook04/plectis-erdos/blob/181078b6b009d905809cf2e007ad309386a3d1e0/lean/ErdosProblems/Erdos249/PaperCompleteR21/AffineDivisorAnnihilation.lean\#L99}{\nolinkurl{joint35ConeWindow_eq}}, \href{https://github.com/wcook04/plectis-erdos/blob/181078b6b009d905809cf2e007ad309386a3d1e0/lean/ErdosProblems/Erdos249/PaperCompleteR21/AffineDivisorAnnihilation.lean\#L119}{\nolinkurl{totientTail_enclosure}}, \href{https://github.com/wcook04/plectis-erdos/blob/181078b6b009d905809cf2e007ad309386a3d1e0/lean/ErdosProblems/Erdos249/PaperCompleteR21/AffineDivisorAnnihilation.lean\#L58}{\nolinkurl{mobiusTermKernel_moment_annihilation}}, \href{https://github.com/wcook04/plectis-erdos/blob/181078b6b009d905809cf2e007ad309386a3d1e0/lean/ErdosProblems/Erdos249/PaperCompleteR21/AffineDivisorAnnihilation.lean\#L76}{\nolinkurl{joint35_mobiusTermKernel_zero}}, \href{https://github.com/wcook04/plectis-erdos/blob/181078b6b009d905809cf2e007ad309386a3d1e0/lean/ErdosProblems/Erdos249/PaperCompleteR21/AffineDivisorAnnihilation.lean\#L87}{\nolinkurl{joint35_coefficient_moments}}, \href{https://github.com/wcook04/plectis-erdos/blob/181078b6b009d905809cf2e007ad309386a3d1e0/lean/ErdosProblems/Erdos249/PaperCompleteR21/AffineDivisorAnnihilation.lean\#L132}{\nolinkurl{joint35_truncation_error}}, \href{https://github.com/wcook04/plectis-erdos/blob/181078b6b009d905809cf2e007ad309386a3d1e0/lean/ErdosProblems/Erdos249/PaperCompleteR21/AffineDivisorAnnihilation.lean\#L183}{\nolinkurl{joint35_nonintegral_of_separated_window}}.  \emph{Compared}, runs \href{https://github.com/wcook04/plectis-erdos-lean/actions/runs/35624228171}{\texttt{35624228171}}, \href{https://github.com/wcook04/plectis-erdos-lean/actions/runs/35643815458}{\texttt{35643815458}}, \href{https://github.com/wcook04/plectis-erdos-lean/actions/runs/35674034595}{\texttt{35674034595}}.
\par\noindent\hangindent=1.5em Proposition~\ref{catalogue:mob:f1}: \href{https://github.com/wcook04/plectis-erdos/blob/181078b6b009d905809cf2e007ad309386a3d1e0/lean/ErdosProblems/Erdos249/PaperCompleteR21/CompositeDilationIdentity.lean\#L27}{\nolinkurl{composite_dilation_divisor_count}}, \href{https://github.com/wcook04/plectis-erdos/blob/181078b6b009d905809cf2e007ad309386a3d1e0/lean/ErdosProblems/Erdos249/PaperCompleteR21/CompositeDilationIdentity.lean\#L36}{\nolinkurl{composite_dilation_defect_eq_zero_of_prime_support}}, \href{https://github.com/wcook04/plectis-erdos/blob/181078b6b009d905809cf2e007ad309386a3d1e0/lean/ErdosProblems/Erdos249/PaperCompleteR21/CompositeDilationIdentity.lean\#L42}{\nolinkurl{composite_dilation_divisor_count_prime_support}}, \href{https://github.com/wcook04/plectis-erdos/blob/181078b6b009d905809cf2e007ad309386a3d1e0/lean/ErdosProblems/Erdos249/PaperCompleteR21/CompositeDilationIdentity.lean\#L49}{\nolinkurl{composite_dilation_defect_univ_six_one}}, \href{https://github.com/wcook04/plectis-erdos/blob/181078b6b009d905809cf2e007ad309386a3d1e0/lean/ErdosProblems/Erdos249/PaperCompleteR21/CompositeDilationIdentity.lean\#L72}{\nolinkurl{composite_dilation_defect_univ_six_one_card}}, \href{https://github.com/wcook04/plectis-erdos/blob/181078b6b009d905809cf2e007ad309386a3d1e0/lean/ErdosProblems/Erdos249/PaperCompleteR21/CompositeDilationIdentity.lean\#L83}{\nolinkurl{lambert_support_series}}, \href{https://github.com/wcook04/plectis-erdos/blob/181078b6b009d905809cf2e007ad309386a3d1e0/lean/ErdosProblems/Erdos249/PaperCompleteR21/CompositeDilationIdentity.lean\#L92}{\nolinkurl{lambert_support_series_restricted}}, \href{https://github.com/wcook04/plectis-erdos/blob/181078b6b009d905809cf2e007ad309386a3d1e0/lean/ErdosProblems/Erdos249/PaperCompleteR21/CompositeDilationIdentity.lean\#L108}{\nolinkurl{totient_convolution_weight_mul_zeta}}, \href{https://github.com/wcook04/plectis-erdos/blob/181078b6b009d905809cf2e007ad309386a3d1e0/lean/ErdosProblems/Erdos249/PaperCompleteR21/CompositeDilationIdentity.lean\#L114}{\nolinkurl{sum_divisors_totient_ne_totient}}, \href{https://github.com/wcook04/plectis-erdos/blob/181078b6b009d905809cf2e007ad309386a3d1e0/lean/ErdosProblems/Erdos249/PaperCompleteR21/CompositeDilationIdentity.lean\#L120}{\nolinkurl{totient_convolution_weight_prime}}, \href{https://github.com/wcook04/plectis-erdos/blob/181078b6b009d905809cf2e007ad309386a3d1e0/lean/ErdosProblems/Erdos249/PaperCompleteR21/CompositeDilationIdentity.lean\#L125}{\nolinkurl{totient_convolution_weight_unbounded}}, \href{https://github.com/wcook04/plectis-erdos/blob/181078b6b009d905809cf2e007ad309386a3d1e0/lean/ErdosProblems/Erdos249/PaperCompleteR21/CompositeDilationIdentity.lean\#L131}{\nolinkurl{totient_convolution_weight_not_periodic}}.  \emph{Compared}, runs \href{https://github.com/wcook04/plectis-erdos-lean/actions/runs/35624228171}{\texttt{35624228171}}, \href{https://github.com/wcook04/plectis-erdos-lean/actions/runs/35643815458}{\texttt{35643815458}}.
\par\noindent\hangindent=1.5em Theorem~\ref{catalogue:cert:a10}: \href{https://github.com/wcook04/plectis-erdos/blob/181078b6b009d905809cf2e007ad309386a3d1e0/lean/Erdos249257/LcmConeFlatness.lean\#L412}{\nolinkurl{irrational_totient_series_iff_certificate_supply}}.  \emph{Compared}, run \href{https://github.com/wcook04/plectis-erdos-lean/actions/runs/35643815458}{\texttt{35643815458}}.
\par\noindent\hangindent=1.5em Theorem~\ref{catalogue:cert:b2}: \href{https://github.com/wcook04/plectis-erdos/blob/181078b6b009d905809cf2e007ad309386a3d1e0/lean/ErdosProblems/Erdos249/PaperCompleteR21/PeriodMultipleAndSecondDifference.lean\#L25}{\nolinkurl{irrational_of_period_multiple_certificate_supply}}, \href{https://github.com/wcook04/plectis-erdos/blob/181078b6b009d905809cf2e007ad309386a3d1e0/lean/ErdosProblems/Erdos249/PaperCompleteR21/PeriodMultipleAndSecondDifference.lean\#L33}{\nolinkurl{period_multiple_certificate_at_one}}, \href{https://github.com/wcook04/plectis-erdos/blob/181078b6b009d905809cf2e007ad309386a3d1e0/lean/ErdosProblems/Erdos249/PaperCompleteR21/PeriodMultipleAndSecondDifference.lean\#L55}{\nolinkurl{rational_forces_period_multiple_integrality}}, \href{https://github.com/wcook04/plectis-erdos/blob/181078b6b009d905809cf2e007ad309386a3d1e0/lean/ErdosProblems/Erdos249/PaperCompleteR21/PeriodMultipleAndSecondDifference.lean\#L41}{\nolinkurl{period_multiple_certificate_supply_iff}}.  \emph{Compared}, run \href{https://github.com/wcook04/plectis-erdos-lean/actions/runs/35674034595}{\texttt{35674034595}}.
\par\noindent\hangindent=1.5em Theorem~\ref{catalogue:cert:b3}: \href{https://github.com/wcook04/plectis-erdos/blob/181078b6b009d905809cf2e007ad309386a3d1e0/lean/Erdos249257/LcmConeFlatness.lean\#L426}{\nolinkurl{irrational_totient_series_iff_lcm_diagonal_certificate_supply}}.  \emph{Compared}, run \href{https://github.com/wcook04/plectis-erdos-lean/actions/runs/35674034595}{\texttt{35674034595}}.
\par\noindent\hangindent=1.5em Lemma~\ref{catalogue:cert:b4}: \href{https://github.com/wcook04/plectis-erdos/blob/181078b6b009d905809cf2e007ad309386a3d1e0/lean/Erdos249257/LcmDiagonalReduction.lean\#L137}{\nolinkurl{eq_prime_pow_of_not_dvd_periodLcm}}.  \emph{Compared}, run \href{https://github.com/wcook04/plectis-erdos-lean/actions/runs/35674034595}{\texttt{35674034595}}.
\par\noindent\hangindent=1.5em Proposition~\ref{catalogue:cert:b5}: \href{https://github.com/wcook04/plectis-erdos/blob/181078b6b009d905809cf2e007ad309386a3d1e0/lean/ErdosProblems/Erdos249/PaperCompleteR20/LcmGridCorrespondence.lean\#L67}{\nolinkurl{clean_lcm_ray_factorisation}}, \href{https://github.com/wcook04/plectis-erdos/blob/181078b6b009d905809cf2e007ad309386a3d1e0/lean/ErdosProblems/Erdos249/PaperCompleteR20/LcmGridCorrespondence.lean\#L76}{\nolinkurl{unclean_lcm_ray_counterexample}}.  \emph{Compared}, run \href{https://github.com/wcook04/plectis-erdos-lean/actions/runs/35674034595}{\texttt{35674034595}}.
\par\noindent\hangindent=1.5em Theorem~\ref{catalogue:cert:b6}: \href{https://github.com/wcook04/plectis-erdos/blob/181078b6b009d905809cf2e007ad309386a3d1e0/lean/ErdosProblems/Erdos249/PaperCompleteR20/LcmGridCorrespondence.lean\#L8}{\nolinkurl{lcm_grid_flatness}}, \href{https://github.com/wcook04/plectis-erdos/blob/181078b6b009d905809cf2e007ad309386a3d1e0/lean/ErdosProblems/Erdos249/PaperCompleteR20/LcmGridCorrespondence.lean\#L15}{\nolinkurl{lcm_grid_fractional_parts}}, \href{https://github.com/wcook04/plectis-erdos/blob/181078b6b009d905809cf2e007ad309386a3d1e0/lean/ErdosProblems/Erdos249/PaperCompleteR20/LcmGridCorrespondence.lean\#L27}{\nolinkurl{certificate_denominator_exclusion}}, \href{https://github.com/wcook04/plectis-erdos/blob/181078b6b009d905809cf2e007ad309386a3d1e0/lean/ErdosProblems/Erdos249/PaperCompleteR20/LcmGridCorrespondence.lean\#L40}{\nolinkurl{lcm_grid_supply_iff}}.  \emph{Compared}, run \href{https://github.com/wcook04/plectis-erdos-lean/actions/runs/35674034595}{\texttt{35674034595}}.
\par\noindent\hangindent=1.5em Theorem~\ref{catalogue:cert:b7}: \href{https://github.com/wcook04/plectis-erdos/blob/181078b6b009d905809cf2e007ad309386a3d1e0/lean/ErdosProblems/Erdos249/PaperCompleteR20/LcmGridCorrespondence.lean\#L40}{\nolinkurl{lcm_grid_supply_iff}}, \href{https://github.com/wcook04/plectis-erdos/blob/181078b6b009d905809cf2e007ad309386a3d1e0/lean/ErdosProblems/Erdos249/PaperCompleteR20/LcmGridCorrespondence.lean\#L50}{\nolinkurl{lcm_grid_multiplier_positive}}.  \emph{Compared}, run \href{https://github.com/wcook04/plectis-erdos-lean/actions/runs/35674034595}{\texttt{35674034595}}.
\par\noindent\hangindent=1.5em Corollary~\ref{catalogue:cert:b8}: \href{https://github.com/wcook04/plectis-erdos/blob/181078b6b009d905809cf2e007ad309386a3d1e0/lean/Erdos249257/LcmConeFlatness.lean\#L386}{\nolinkurl{irrational_totient_series_iff_all_tail_diffs_nonintegral}}, \href{https://github.com/wcook04/plectis-erdos/blob/181078b6b009d905809cf2e007ad309386a3d1e0/lean/Erdos249257/LcmConeFlatness.lean\#L439}{\nolinkurl{periodLcm_diagonal_kill_iff_tail_diff_notMem_int}}, \href{https://github.com/wcook04/plectis-erdos/blob/181078b6b009d905809cf2e007ad309386a3d1e0/lean/Erdos249257/LcmConeFlatness.lean\#L451}{\nolinkurl{irrational_totient_series_of_lcm_diagonal_nonintegrality_supply}}, \href{https://github.com/wcook04/plectis-erdos/blob/181078b6b009d905809cf2e007ad309386a3d1e0/lean/Erdos249257/LcmConeFlatness.lean\#L462}{\nolinkurl{irrational_totient_series_of_lcm_cone_nonintegrality_supply}}.  \emph{Compared}, runs \href{https://github.com/wcook04/plectis-erdos-lean/actions/runs/35643815458}{\texttt{35643815458}}, \href{https://github.com/wcook04/plectis-erdos-lean/actions/runs/35674034595}{\texttt{35674034595}}.
\par\noindent\hangindent=1.5em Proposition~\ref{catalogue:cert:b9a}: \href{https://github.com/wcook04/plectis-erdos/blob/181078b6b009d905809cf2e007ad309386a3d1e0/lean/ErdosProblems/Erdos249/PaperCompleteR21/PeriodMultipleAndSecondDifference.lean\#L65}{\nolinkurl{abs_tail_diff_scaled_sub_window_le}}, \href{https://github.com/wcook04/plectis-erdos/blob/181078b6b009d905809cf2e007ad309386a3d1e0/lean/ErdosProblems/Erdos249/PaperCompleteR21/PeriodMultipleAndSecondDifference.lean\#L104}{\nolinkurl{second_difference_error_bound}}, \href{https://github.com/wcook04/plectis-erdos/blob/181078b6b009d905809cf2e007ad309386a3d1e0/lean/ErdosProblems/Erdos249/PaperCompleteR21/PeriodMultipleAndSecondDifference.lean\#L121}{\nolinkurl{second_difference_certificate_sound}}, \href{https://github.com/wcook04/plectis-erdos/blob/181078b6b009d905809cf2e007ad309386a3d1e0/lean/ErdosProblems/Erdos249/PaperCompleteR21/PeriodMultipleAndSecondDifference.lean\#L136}{\nolinkurl{second_difference_cell_one_eight}}.  \emph{Compared}, run \href{https://github.com/wcook04/plectis-erdos-lean/actions/runs/35674034595}{\texttt{35674034595}}.
\par\noindent\hangindent=1.5em Theorem~\ref{catalogue:cert:b10a}: \href{https://github.com/wcook04/plectis-erdos/blob/181078b6b009d905809cf2e007ad309386a3d1e0/lean/ErdosProblems/Erdos249/PaperCompleteR20/FiniteGridCorrespondence.lean\#L12}{\nolinkurl{paperGridNumerator_eq}}, \href{https://github.com/wcook04/plectis-erdos/blob/181078b6b009d905809cf2e007ad309386a3d1e0/lean/ErdosProblems/Erdos249/PaperCompleteR20/FiniteGridCorrespondence.lean\#L36}{\nolinkurl{finite_grid_nonintegral_pair}}.  \emph{Compared}, run \href{https://github.com/wcook04/plectis-erdos-lean/actions/runs/35674034595}{\texttt{35674034595}}.
\par\noindent\hangindent=1.5em Theorem~\ref{catalogue:cert:b10b}: \href{https://github.com/wcook04/plectis-erdos/blob/181078b6b009d905809cf2e007ad309386a3d1e0/lean/ErdosProblems/Erdos249/PaperCompleteR20/FiniteGridCorrespondence.lean\#L12}{\nolinkurl{paperGridNumerator_eq}}, \href{https://github.com/wcook04/plectis-erdos/blob/181078b6b009d905809cf2e007ad309386a3d1e0/lean/ErdosProblems/Erdos249/PaperCompleteR20/FiniteGridCorrespondence.lean\#L49}{\nolinkurl{finite_grid_supply_irrational}}.  \emph{Compared}, run \href{https://github.com/wcook04/plectis-erdos-lean/actions/runs/35674034595}{\texttt{35674034595}}.
\par\noindent\hangindent=1.5em Proposition~\ref{catalogue:cert:b12}: \href{https://github.com/wcook04/plectis-erdos/blob/181078b6b009d905809cf2e007ad309386a3d1e0/lean/ErdosProblems/Erdos249/PaperCompleteR20/FiniteCarryCorrespondence.lean\#L8}{\nolinkurl{finite_carry_test_sound}}, \href{https://github.com/wcook04/plectis-erdos/blob/181078b6b009d905809cf2e007ad309386a3d1e0/lean/ErdosProblems/Erdos249/PaperCompleteR20/FiniteCarryCorrespondence.lean\#L22}{\nolinkurl{finite_carry_candidate_count}}, \href{https://github.com/wcook04/plectis-erdos/blob/181078b6b009d905809cf2e007ad309386a3d1e0/lean/ErdosProblems/Erdos249/PaperCompleteR20/FiniteCarryCorrespondence.lean\#L27}{\nolinkurl{finite_carry_true_orbit}}.  \emph{Compared}, runs \href{https://github.com/wcook04/plectis-erdos-lean/actions/runs/35643815458}{\texttt{35643815458}}, \href{https://github.com/wcook04/plectis-erdos-lean/actions/runs/35674034595}{\texttt{35674034595}}.
\par\noindent\hangindent=1.5em Theorem~\ref{catalogue:mob:e1}: \href{https://github.com/wcook04/plectis-erdos/blob/181078b6b009d905809cf2e007ad309386a3d1e0/lean/ErdosProblems/Erdos249/PaperCompleteR21/HarmonicGapAndFourTail.lean\#L71}{\nolinkurl{irrational_of_first_harmonic_norm_gap}}, \href{https://github.com/wcook04/plectis-erdos/blob/181078b6b009d905809cf2e007ad309386a3d1e0/lean/ErdosProblems/Erdos249/PaperCompleteR21/HarmonicGapAndFourTail.lean\#L40}{\nolinkurl{first_harmonic_re_bound_of_norm_bound}}, \href{https://github.com/wcook04/plectis-erdos/blob/181078b6b009d905809cf2e007ad309386a3d1e0/lean/ErdosProblems/Erdos249/PaperCompleteR21/HarmonicGapAndFourTail.lean\#L54}{\nolinkurl{exists_certificate_of_first_harmonic_real_bound}}, \href{https://github.com/wcook04/plectis-erdos/blob/181078b6b009d905809cf2e007ad309386a3d1e0/lean/ErdosProblems/Erdos249/PaperCompleteR21/HarmonicGapAndFourTail.lean\#L61}{\nolinkurl{exists_certificate_of_first_harmonic_norm_bound}}, \href{https://github.com/wcook04/plectis-erdos/blob/181078b6b009d905809cf2e007ad309386a3d1e0/lean/ErdosProblems/Erdos249/PaperCompleteR21/HarmonicGapAndFourTail.lean\#L28}{\nolinkurl{nine_tenths_lt_cos_pi_div_eight}}.  \emph{Compared}, runs \href{https://github.com/wcook04/plectis-erdos-lean/actions/runs/35643815458}{\texttt{35643815458}}, \href{https://github.com/wcook04/plectis-erdos-lean/actions/runs/35674034595}{\texttt{35674034595}}.
\par\noindent\hangindent=1.5em Theorem~\ref{catalogue:mob:e2}: \href{https://github.com/wcook04/plectis-erdos/blob/181078b6b009d905809cf2e007ad309386a3d1e0/lean/ErdosProblems/Erdos249/PaperCompleteR21/HarmonicGapAndFourTail.lean\#L103}{\nolinkurl{four_tail_combination_eq}}, \href{https://github.com/wcook04/plectis-erdos/blob/181078b6b009d905809cf2e007ad309386a3d1e0/lean/ErdosProblems/Erdos249/PaperCompleteR21/HarmonicGapAndFourTail.lean\#L93}{\nolinkurl{four_tail_window_eq}}, \href{https://github.com/wcook04/plectis-erdos/blob/181078b6b009d905809cf2e007ad309386a3d1e0/lean/ErdosProblems/Erdos249/PaperCompleteR21/HarmonicGapAndFourTail.lean\#L81}{\nolinkurl{windowDiscrepancy_diagonal_eq}}, \href{https://github.com/wcook04/plectis-erdos/blob/181078b6b009d905809cf2e007ad309386a3d1e0/lean/ErdosProblems/Erdos249/PaperCompleteR21/HarmonicGapAndFourTail.lean\#L113}{\nolinkurl{four_tail_error_bound}}, \href{https://github.com/wcook04/plectis-erdos/blob/181078b6b009d905809cf2e007ad309386a3d1e0/lean/ErdosProblems/Erdos249/PaperCompleteR21/HarmonicGapAndFourTail.lean\#L129}{\nolinkurl{totientTail_bounds}}, \href{https://github.com/wcook04/plectis-erdos/blob/181078b6b009d905809cf2e007ad309386a3d1e0/lean/ErdosProblems/Erdos249/PaperCompleteR21/HarmonicGapAndFourTail.lean\#L152}{\nolinkurl{four_tail_criterion_sound}}, \href{https://github.com/wcook04/plectis-erdos/blob/181078b6b009d905809cf2e007ad309386a3d1e0/lean/ErdosProblems/Erdos249/PaperCompleteR21/HarmonicGapAndFourTail.lean\#L169}{\nolinkurl{irrational_of_four_tail_supply}}, \href{https://github.com/wcook04/plectis-erdos/blob/181078b6b009d905809cf2e007ad309386a3d1e0/lean/ErdosProblems/Erdos249/PaperCompleteR21/HarmonicGapAndFourTail.lean\#L135}{\nolinkurl{rational_forces_four_tail_diagonals_integral}}, \href{https://github.com/wcook04/plectis-erdos/blob/181078b6b009d905809cf2e007ad309386a3d1e0/lean/ErdosProblems/Erdos249/PaperCompleteR21/HarmonicGapAndFourTail.lean\#L187}{\nolinkurl{four_tail_checked_instance}}.  \emph{Compared}, run \href{https://github.com/wcook04/plectis-erdos-lean/actions/runs/35674034595}{\texttt{35674034595}}.
\par\noindent\hangindent=1.5em Proposition~\ref{prop:NI-01}: \href{https://github.com/wcook04/plectis-erdos/blob/181078b6b009d905809cf2e007ad309386a3d1e0/lean/ErdosProblems/Erdos249/PaperCompleteR21/ActualLcmDiagonalConditions.lean\#L29}{\nolinkurl{irrational_iff_diagonal_orbit_nonintegrality}}.  \emph{Compared}, run \href{https://github.com/wcook04/plectis-erdos-lean/actions/runs/35674034595}{\texttt{35674034595}}.
\par\noindent\hangindent=1.5em Proposition~\ref{prop:AR-07}: \href{https://github.com/wcook04/plectis-erdos/blob/181078b6b009d905809cf2e007ad309386a3d1e0/lean/ErdosProblems/Erdos249/PaperCompleteR21/ActualLcmDiagonalConditions.lean\#L56}{\nolinkurl{irrational_of_short_window_diagonal_supply}}, \href{https://github.com/wcook04/plectis-erdos/blob/181078b6b009d905809cf2e007ad309386a3d1e0/lean/ErdosProblems/Erdos249/PaperCompleteR21/ActualLcmDiagonalConditions.lean\#L40}{\nolinkurl{diagonal_certificate_unfolded}}, \href{https://github.com/wcook04/plectis-erdos/blob/181078b6b009d905809cf2e007ad309386a3d1e0/lean/ErdosProblems/Erdos249/PaperCompleteR21/ActualLcmDiagonalConditions.lean\#L68}{\nolinkurl{pointwise_completeness_supplies_some_depth}}.  \emph{Compared}, run \href{https://github.com/wcook04/plectis-erdos-lean/actions/runs/35674034595}{\texttt{35674034595}}.
\par\noindent\hangindent=1.5em Proposition~\ref{prop:SEP-03}: \href{https://github.com/wcook04/plectis-erdos/blob/181078b6b009d905809cf2e007ad309386a3d1e0/lean/ErdosProblems/Erdos249/PaperCompleteR21/ActualLcmDiagonalConditions.lean\#L145}{\nolinkurl{irrational_of_diagonal_orbit_separation_supply}}, \href{https://github.com/wcook04/plectis-erdos/blob/181078b6b009d905809cf2e007ad309386a3d1e0/lean/ErdosProblems/Erdos249/PaperCompleteR21/ActualLcmDiagonalConditions.lean\#L77}{\nolinkurl{prescribedOddIndex}}, \href{https://github.com/wcook04/plectis-erdos/blob/181078b6b009d905809cf2e007ad309386a3d1e0/lean/ErdosProblems/Erdos249/PaperCompleteR21/ActualLcmDiagonalConditions.lean\#L80}{\nolinkurl{oddGuarded_depth_eq_prescribed}}, \href{https://github.com/wcook04/plectis-erdos/blob/181078b6b009d905809cf2e007ad309386a3d1e0/lean/ErdosProblems/Erdos249/PaperCompleteR21/ActualLcmDiagonalConditions.lean\#L98}{\nolinkurl{abs_orbit_sub_rawApprox_lt}}, \href{https://github.com/wcook04/plectis-erdos/blob/181078b6b009d905809cf2e007ad309386a3d1e0/lean/ErdosProblems/Erdos249/PaperCompleteR21/ActualLcmDiagonalConditions.lean\#L112}{\nolinkurl{rawApprox_separation_of_orbit_separation}}.
\par\noindent\hangindent=1.5em Proposition~\ref{prop:TE-04}: \href{https://github.com/wcook04/plectis-erdos/blob/181078b6b009d905809cf2e007ad309386a3d1e0/lean/ErdosProblems/Erdos249/PaperCompleteR21/TopEdgeStaircaseConditions.lean\#L29}{\nolinkurl{upper_endpoint_condition_iff}}, \href{https://github.com/wcook04/plectis-erdos/blob/181078b6b009d905809cf2e007ad309386a3d1e0/lean/ErdosProblems/Erdos249/PaperCompleteR21/TopEdgeStaircaseConditions.lean\#L38}{\nolinkurl{upper_endpoint_gap_nonintegral}}, \href{https://github.com/wcook04/plectis-erdos/blob/181078b6b009d905809cf2e007ad309386a3d1e0/lean/ErdosProblems/Erdos249/PaperCompleteR21/TopEdgeStaircaseConditions.lean\#L50}{\nolinkurl{integral_tail_forces_upper_endpoint_residue}}, \href{https://github.com/wcook04/plectis-erdos/blob/181078b6b009d905809cf2e007ad309386a3d1e0/lean/ErdosProblems/Erdos249/PaperCompleteR21/TopEdgeStaircaseConditions.lean\#L66}{\nolinkurl{irrational_of_upper_endpoint_gap_supply}}.  \emph{Compared}, run \href{https://github.com/wcook04/plectis-erdos-lean/actions/runs/35674034595}{\texttt{35674034595}}.
\par\noindent\hangindent=1.5em Proposition~\ref{prop:TE-05}: \href{https://github.com/wcook04/plectis-erdos/blob/181078b6b009d905809cf2e007ad309386a3d1e0/lean/ErdosProblems/Erdos249/PaperCompleteR21/TopEdgeStaircaseConditions.lean\#L80}{\nolinkurl{te_chain_relations}}.  \emph{Compared}, run \href{https://github.com/wcook04/plectis-erdos-lean/actions/runs/35674034595}{\texttt{35674034595}}.
\par\noindent\hangindent=1.5em Proposition~\ref{prop:TE-06}: \href{https://github.com/wcook04/plectis-erdos/blob/181078b6b009d905809cf2e007ad309386a3d1e0/lean/ErdosProblems/Erdos249/PaperCompleteR21/TopEdgeStaircaseConditions.lean\#L142}{\nolinkurl{endpoint_identity}}, \href{https://github.com/wcook04/plectis-erdos/blob/181078b6b009d905809cf2e007ad309386a3d1e0/lean/ErdosProblems/Erdos249/PaperCompleteR21/TopEdgeStaircaseConditions.lean\#L158}{\nolinkurl{integral_carry_strictly_between}}, \href{https://github.com/wcook04/plectis-erdos/blob/181078b6b009d905809cf2e007ad309386a3d1e0/lean/ErdosProblems/Erdos249/PaperCompleteR21/TopEdgeStaircaseConditions.lean\#L187}{\nolinkurl{endpoint_criterion_nonintegral}}, \href{https://github.com/wcook04/plectis-erdos/blob/181078b6b009d905809cf2e007ad309386a3d1e0/lean/ErdosProblems/Erdos249/PaperCompleteR21/TopEdgeStaircaseConditions.lean\#L116}{\nolinkurl{oddHalfCenteredLift_spec}}, \href{https://github.com/wcook04/plectis-erdos/blob/181078b6b009d905809cf2e007ad309386a3d1e0/lean/ErdosProblems/Erdos249/PaperCompleteR21/TopEdgeStaircaseConditions.lean\#L106}{\nolinkurl{centeredLift_range}}.  \emph{Compared}, runs \href{https://github.com/wcook04/plectis-erdos-lean/actions/runs/35624228171}{\texttt{35624228171}}, \href{https://github.com/wcook04/plectis-erdos-lean/actions/runs/35674034595}{\texttt{35674034595}}.
\par\noindent\hangindent=1.5em Proposition~\ref{prop:SK-02}: \href{https://github.com/wcook04/plectis-erdos/blob/181078b6b009d905809cf2e007ad309386a3d1e0/lean/ErdosProblems/Erdos249/PaperCompleteR21/ActualLcmDiagonalConditions.lean\#L161}{\nolinkurl{short_window_diagonal_through_six}}, \href{https://github.com/wcook04/plectis-erdos/blob/181078b6b009d905809cf2e007ad309386a3d1e0/lean/ErdosProblems/Erdos249/PaperCompleteR21/ActualLcmDiagonalConditions.lean\#L169}{\nolinkurl{short_window_diagonal_witnesses}}.  \emph{Compared}, run \href{https://github.com/wcook04/plectis-erdos-lean/actions/runs/35674034595}{\texttt{35674034595}}.
\par\noindent\hangindent=1.5em Proposition~\ref{prop:FR-01}: \href{https://github.com/wcook04/plectis-erdos/blob/181078b6b009d905809cf2e007ad309386a3d1e0/lean/ErdosProblems/Erdos249/PaperCompleteR21/ExtremalOrderDirectedAndPulse.lean\#L27}{\nolinkurl{extremal_order_curvature_pos}}, \href{https://github.com/wcook04/plectis-erdos/blob/181078b6b009d905809cf2e007ad309386a3d1e0/lean/ErdosProblems/Erdos249/PaperCompleteR21/ExtremalOrderDirectedAndPulse.lean\#L36}{\nolinkurl{extremal_order_curvature_neg}}, \href{https://github.com/wcook04/plectis-erdos/blob/181078b6b009d905809cf2e007ad309386a3d1e0/lean/ErdosProblems/Erdos249/PaperCompleteR21/ExtremalOrderDirectedAndPulse.lean\#L45}{\nolinkurl{extremal_order_curvature_ne_zero}}.  \emph{Compared}, run \href{https://github.com/wcook04/plectis-erdos-lean/actions/runs/35643815458}{\texttt{35643815458}}.
\par\noindent\hangindent=1.5em Proposition~\ref{prop:CP-06}: \href{https://github.com/wcook04/plectis-erdos/blob/181078b6b009d905809cf2e007ad309386a3d1e0/lean/ErdosProblems/Erdos249/PaperCompleteR21/ExtremalOrderDirectedAndPulse.lean\#L59}{\nolinkurl{directed_certificate_iff}}, \href{https://github.com/wcook04/plectis-erdos/blob/181078b6b009d905809cf2e007ad309386a3d1e0/lean/ErdosProblems/Erdos249/PaperCompleteR21/ExtremalOrderDirectedAndPulse.lean\#L71}{\nolinkurl{directed_certificate_example}}.  \emph{Compared}, run \href{https://github.com/wcook04/plectis-erdos-lean/actions/runs/35674034595}{\texttt{35674034595}}.
\par\noindent\hangindent=1.5em Proposition~\ref{prop:CP-07}: \href{https://github.com/wcook04/plectis-erdos/blob/181078b6b009d905809cf2e007ad309386a3d1e0/lean/ErdosProblems/Erdos249/PaperCompleteR21/ExtremalOrderDirectedAndPulse.lean\#L105}{\nolinkurl{irrational_of_modFour_pulse_supply}}, \href{https://github.com/wcook04/plectis-erdos/blob/181078b6b009d905809cf2e007ad309386a3d1e0/lean/ErdosProblems/Erdos249/PaperCompleteR21/ExtremalOrderDirectedAndPulse.lean\#L88}{\nolinkurl{rational_forces_pulse_class_integrality}}.  \emph{Compared}, run \href{https://github.com/wcook04/plectis-erdos-lean/actions/runs/35674034595}{\texttt{35674034595}}.
\par\noindent\hangindent=1.5em Proposition~\ref{prop:A10}: \href{https://github.com/wcook04/plectis-erdos/blob/181078b6b009d905809cf2e007ad309386a3d1e0/lean/Erdos249257/LcmConeFlatness.lean\#L412}{\nolinkurl{irrational_totient_series_iff_certificate_supply}}.  \emph{Compared}, run \href{https://github.com/wcook04/plectis-erdos-lean/actions/runs/35643815458}{\texttt{35643815458}}.
\par\noindent\hangindent=1.5em Proposition~\ref{prop:B3}: \href{https://github.com/wcook04/plectis-erdos/blob/181078b6b009d905809cf2e007ad309386a3d1e0/lean/Erdos249257/LcmConeFlatness.lean\#L426}{\nolinkurl{irrational_totient_series_iff_lcm_diagonal_certificate_supply}}.  \emph{Compared}, run \href{https://github.com/wcook04/plectis-erdos-lean/actions/runs/35674034595}{\texttt{35674034595}}.
\par\noindent\hangindent=1.5em Proposition~\ref{prop:B7}: \href{https://github.com/wcook04/plectis-erdos/blob/181078b6b009d905809cf2e007ad309386a3d1e0/lean/Erdos249257/CertificateKernel.lean\#L19014}{\nolinkurl{irrational_totient_series_of_lcm_cone_window_kill_supply}}.  \emph{Compared}, run \href{https://github.com/wcook04/plectis-erdos-lean/actions/runs/35674034595}{\texttt{35674034595}}.
\par\noindent\hangindent=1.5em Proposition~\ref{prop:B10}: \href{https://github.com/wcook04/plectis-erdos/blob/181078b6b009d905809cf2e007ad309386a3d1e0/lean/ErdosProblems/Erdos249/PaperCompleteR20/FiniteGridCorrespondence.lean\#L12}{\nolinkurl{paperGridNumerator_eq}}, \href{https://github.com/wcook04/plectis-erdos/blob/181078b6b009d905809cf2e007ad309386a3d1e0/lean/ErdosProblems/Erdos249/PaperCompleteR20/FiniteGridCorrespondence.lean\#L36}{\nolinkurl{finite_grid_nonintegral_pair}}, \href{https://github.com/wcook04/plectis-erdos/blob/181078b6b009d905809cf2e007ad309386a3d1e0/lean/ErdosProblems/Erdos249/PaperCompleteR20/FiniteGridCorrespondence.lean\#L49}{\nolinkurl{finite_grid_supply_irrational}}.  \emph{Compared}, run \href{https://github.com/wcook04/plectis-erdos-lean/actions/runs/35674034595}{\texttt{35674034595}}.
\par\noindent\hangindent=1.5em Proposition~\ref{prop:C2sup}: \href{https://github.com/wcook04/plectis-erdos/blob/181078b6b009d905809cf2e007ad309386a3d1e0/lean/ErdosProblems/Erdos249/PaperCompleteR21/FareyGapDenominatorExclusion.lean\#L41}{\nolinkurl{gapCertificate_window_1_240}}, \href{https://github.com/wcook04/plectis-erdos/blob/181078b6b009d905809cf2e007ad309386a3d1e0/lean/ErdosProblems/Erdos249/PaperCompleteR21/FareyGapDenominatorExclusion.lean\#L29}{\nolinkurl{irrational_totientSeries_of_fareyGapExclusionUnbounded}}, \href{https://github.com/wcook04/plectis-erdos/blob/181078b6b009d905809cf2e007ad309386a3d1e0/lean/ErdosProblems/Erdos249/PaperCompleteR21/FareyGapDenominatorExclusion.lean\#L51}{\nolinkurl{gapFareyBound_window_1_240}}.  \emph{Compared}, run \href{https://github.com/wcook04/plectis-erdos-lean/actions/runs/35643815458}{\texttt{35643815458}}.
\par\noindent\hangindent=1.5em Proposition~\ref{prop:D5cons}: \href{https://github.com/wcook04/plectis-erdos/blob/181078b6b009d905809cf2e007ad309386a3d1e0/lean/ErdosProblems/Erdos249/PaperCompleteR21/GenericCarryRankCeilingCounterexample.lean\#L143}{\nolinkurl{rank_floor_and_false_proposed_carryRank_ceiling}}, \href{https://github.com/wcook04/plectis-erdos/blob/181078b6b009d905809cf2e007ad309386a3d1e0/lean/ErdosProblems/Erdos249/PaperCompleteR21/GenericCarryRankCeilingCounterexample.lean\#L64}{\nolinkurl{fiveQuarter_comparison_rational_with_carryRank_floor}}, \href{https://github.com/wcook04/plectis-erdos/blob/181078b6b009d905809cf2e007ad309386a3d1e0/lean/ErdosProblems/Erdos249/PaperCompleteR21/GenericCarryRankCeilingCounterexample.lean\#L100}{\nolinkurl{no_generic_rationality_carryRank_ceiling}}.  \emph{Compared}, runs \href{https://github.com/wcook04/plectis-erdos-lean/actions/runs/35624228171}{\texttt{35624228171}}, \href{https://github.com/wcook04/plectis-erdos-lean/actions/runs/35674034595}{\texttt{35674034595}}.
\par\noindent\hangindent=1.5em Proposition~\ref{prop:B12cons}: \href{https://github.com/wcook04/plectis-erdos/blob/181078b6b009d905809cf2e007ad309386a3d1e0/lean/ErdosProblems/Erdos249/PaperCompleteR20/FiniteCarryCorrespondence.lean\#L8}{\nolinkurl{finite_carry_test_sound}}, \href{https://github.com/wcook04/plectis-erdos/blob/181078b6b009d905809cf2e007ad309386a3d1e0/lean/ErdosProblems/Erdos249/PaperCompleteR20/FiniteCarryCorrespondence.lean\#L22}{\nolinkurl{finite_carry_candidate_count}}, \href{https://github.com/wcook04/plectis-erdos/blob/181078b6b009d905809cf2e007ad309386a3d1e0/lean/ErdosProblems/Erdos249/PaperCompleteR20/FiniteCarryCorrespondence.lean\#L27}{\nolinkurl{finite_carry_true_orbit}}.  \emph{Compared}, runs \href{https://github.com/wcook04/plectis-erdos-lean/actions/runs/35643815458}{\texttt{35643815458}}, \href{https://github.com/wcook04/plectis-erdos-lean/actions/runs/35674034595}{\texttt{35674034595}}.
\par\noindent\hangindent=1.5em Proposition~\ref{prop:AR-04-inv}: \href{https://github.com/wcook04/plectis-erdos/blob/181078b6b009d905809cf2e007ad309386a3d1e0/lean/ErdosProblems/Erdos249/PaperCompleteR21/ActualLcmShortWindowArithmetic.lean\#L56}{\nolinkurl{lcmRayArithmeticLetter_eq_totient_difference}}, \href{https://github.com/wcook04/plectis-erdos/blob/181078b6b009d905809cf2e007ad309386a3d1e0/lean/ErdosProblems/Erdos249/PaperCompleteR21/ActualLcmShortWindowArithmetic.lean\#L70}{\nolinkurl{totient_mul_eq_totient_mul_gcd_div_totient_gcd}}, \href{https://github.com/wcook04/plectis-erdos/blob/181078b6b009d905809cf2e007ad309386a3d1e0/lean/ErdosProblems/Erdos249/PaperCompleteR21/ActualLcmShortWindowArithmetic.lean\#L86}{\nolinkurl{lcmRay_divisor_denominators_pos}}, \href{https://github.com/wcook04/plectis-erdos/blob/181078b6b009d905809cf2e007ad309386a3d1e0/lean/ErdosProblems/Erdos249/PaperCompleteR21/ActualLcmShortWindowArithmetic.lean\#L104}{\nolinkurl{lcmRay_divisor_product_formula}}, \href{https://github.com/wcook04/plectis-erdos/blob/181078b6b009d905809cf2e007ad309386a3d1e0/lean/ErdosProblems/Erdos249/PaperCompleteR21/ActualLcmShortWindowArithmetic.lean\#L156}{\nolinkurl{lcmRay_divisor_clean_formula}}, \href{https://github.com/wcook04/plectis-erdos/blob/181078b6b009d905809cf2e007ad309386a3d1e0/lean/ErdosProblems/Erdos249/PaperCompleteR21/ActualLcmShortWindowArithmetic.lean\#L144}{\nolinkurl{lcmRay_nondivisor_literal}}, \href{https://github.com/wcook04/plectis-erdos/blob/181078b6b009d905809cf2e007ad309386a3d1e0/lean/ErdosProblems/Erdos249/PaperCompleteR21/ActualLcmShortWindowArithmetic.lean\#L180}{\nolinkurl{lcmRay_divisor_gcd_example}}.  \emph{Compared}, runs \href{https://github.com/wcook04/plectis-erdos-lean/actions/runs/35643815458}{\texttt{35643815458}}, \href{https://github.com/wcook04/plectis-erdos-lean/actions/runs/35674034595}{\texttt{35674034595}}.
\par\noindent\hangindent=1.5em Proposition~\ref{prop:AR-03-inv}: \href{https://github.com/wcook04/plectis-erdos/blob/181078b6b009d905809cf2e007ad309386a3d1e0/lean/ErdosProblems/Erdos249/PaperCompleteR21/ActualLcmShortWindowArithmetic.lean\#L215}{\nolinkurl{rough_integer_prime_count_and_totient_bound}}.  \emph{Compared}, run \href{https://github.com/wcook04/plectis-erdos-lean/actions/runs/35643815458}{\texttt{35643815458}}.
\par\noindent\hangindent=1.5em Proposition~\ref{prop:AR-05-inv}: \href{https://github.com/wcook04/plectis-erdos/blob/181078b6b009d905809cf2e007ad309386a3d1e0/lean/ErdosProblems/Erdos249/PaperCompleteR21/ActualLcmShortWindowArithmetic.lean\#L229}{\nolinkurl{shortWindow_totient_difference_pos}}.  \emph{Compared}, run \href{https://github.com/wcook04/plectis-erdos-lean/actions/runs/35674034595}{\texttt{35674034595}}.
\par\noindent\hangindent=1.5em Proposition~\ref{prop:AR-06-inv}: \href{https://github.com/wcook04/plectis-erdos/blob/181078b6b009d905809cf2e007ad309386a3d1e0/lean/ErdosProblems/Erdos249/PaperCompleteR21/ActualLcmShortWindowArithmetic.lean\#L241}{\nolinkurl{weighted_shortWindow_sum_eq_windowDiscrepancy}}, \href{https://github.com/wcook04/plectis-erdos/blob/181078b6b009d905809cf2e007ad309386a3d1e0/lean/ErdosProblems/Erdos249/PaperCompleteR21/ActualLcmShortWindowArithmetic.lean\#L262}{\nolinkurl{weighted_shortWindow_band_iff_certifiedKill}}.  \emph{Compared}, run \href{https://github.com/wcook04/plectis-erdos-lean/actions/runs/35674034595}{\texttt{35674034595}}.
\par\noindent\hangindent=1.5em Proposition~\ref{prop:SEP-01-inv}: \href{https://github.com/wcook04/plectis-erdos/blob/181078b6b009d905809cf2e007ad309386a3d1e0/lean/ErdosProblems/Erdos249/PaperCompleteR21/ActualLcmSeparationAndSign.lean\#L24}{\nolinkurl{actualLcmTailOrbit_global_to_local}}.  \emph{Compared}, run \href{https://github.com/wcook04/plectis-erdos-lean/actions/runs/35674034595}{\texttt{35674034595}}.
\par\noindent\hangindent=1.5em Proposition~\ref{prop:SEP-02-inv}: \href{https://github.com/wcook04/plectis-erdos/blob/181078b6b009d905809cf2e007ad309386a3d1e0/lean/ErdosProblems/Erdos249/PaperCompleteR21/ActualLcmSeparationAndSign.lean\#L39}{\nolinkurl{abs_actualLcmTailOrbit_sub_rawApprox_lt_explicit}}, \href{https://github.com/wcook04/plectis-erdos/blob/181078b6b009d905809cf2e007ad309386a3d1e0/lean/ErdosProblems/Erdos249/PaperCompleteR21/ActualLcmSeparationAndSign.lean\#L49}{\nolinkurl{tendsto_actualLcmRawErrorRadius_atTop_nhds_zero}}.  \emph{Compared}, run \href{https://github.com/wcook04/plectis-erdos-lean/actions/runs/35674034595}{\texttt{35674034595}}.
\par\noindent\hangindent=1.5em Proposition~\ref{prop:SGN-01}: \href{https://github.com/wcook04/plectis-erdos/blob/181078b6b009d905809cf2e007ad309386a3d1e0/lean/ErdosProblems/Erdos249/PaperCompleteR21/ActualLcmSeparationAndSign.lean\#L82}{\nolinkurl{actualLcm_tailDiff_shift_pos_paper}}, \href{https://github.com/wcook04/plectis-erdos/blob/181078b6b009d905809cf2e007ad309386a3d1e0/lean/ErdosProblems/Erdos249/PaperCompleteR21/ActualLcmSeparationAndSign.lean\#L89}{\nolinkurl{actualLcmTailOrbit_pos}}.  \emph{Compared}, run \href{https://github.com/wcook04/plectis-erdos-lean/actions/runs/35674034595}{\texttt{35674034595}}.
\par\noindent\hangindent=1.5em Proposition~\ref{prop:SGN-03}: \href{https://github.com/wcook04/plectis-erdos/blob/181078b6b009d905809cf2e007ad309386a3d1e0/lean/ErdosProblems/Erdos249/PaperCompleteR21/ActualLcmSeparationAndSign.lean\#L105}{\nolinkurl{actualLcm_integral_forces_topEdgeResidue_paper}}.  \emph{Compared}, run \href{https://github.com/wcook04/plectis-erdos-lean/actions/runs/35674034595}{\texttt{35674034595}}.
\par\noindent\hangindent=1.5em Proposition~\ref{prop:TE-02-inv}: \href{https://github.com/wcook04/plectis-erdos/blob/181078b6b009d905809cf2e007ad309386a3d1e0/lean/ErdosProblems/Erdos249/PaperCompleteR21/PenultimateStaircaseAndRankCurvature.lean\#L37}{\nolinkurl{penultimate_shortWindow_difference_eq_half}}, \href{https://github.com/wcook04/plectis-erdos/blob/181078b6b009d905809cf2e007ad309386a3d1e0/lean/ErdosProblems/Erdos249/PaperCompleteR21/PenultimateStaircaseAndRankCurvature.lean\#L80}{\nolinkurl{dyadicScale_unique_in_open_interval}}.  \emph{Compared}, runs \href{https://github.com/wcook04/plectis-erdos-lean/actions/runs/35643815458}{\texttt{35643815458}}, \href{https://github.com/wcook04/plectis-erdos-lean/actions/runs/35674034595}{\texttt{35674034595}}.
\par\noindent\hangindent=1.5em Proposition~\ref{prop:TE-03-inv}: \href{https://github.com/wcook04/plectis-erdos/blob/181078b6b009d905809cf2e007ad309386a3d1e0/lean/ErdosProblems/Erdos249/PaperCompleteR21/PenultimateStaircaseAndRankCurvature.lean\#L129}{\nolinkurl{exists_certifiedKill_iff_twoBitResidueTest}}, \href{https://github.com/wcook04/plectis-erdos/blob/181078b6b009d905809cf2e007ad309386a3d1e0/lean/ErdosProblems/Erdos249/PaperCompleteR21/PenultimateStaircaseAndRankCurvature.lean\#L104}{\nolinkurl{dyadicMixedGuard_iff_twoBitBand}}.  \emph{Compared}, runs \href{https://github.com/wcook04/plectis-erdos-lean/actions/runs/35643815458}{\texttt{35643815458}}, \href{https://github.com/wcook04/plectis-erdos-lean/actions/runs/35674034595}{\texttt{35674034595}}.
\par\noindent\hangindent=1.5em Proposition~\ref{prop:FR-02-inv}: \href{https://github.com/wcook04/plectis-erdos/blob/181078b6b009d905809cf2e007ad309386a3d1e0/lean/ErdosProblems/Erdos249/PaperCompleteR21/PenultimateStaircaseAndRankCurvature.lean\#L219}{\nolinkurl{two_mul_totient_dvd_totient_second_difference}}, \href{https://github.com/wcook04/plectis-erdos/blob/181078b6b009d905809cf2e007ad309386a3d1e0/lean/ErdosProblems/Erdos249/PaperCompleteR21/PenultimateStaircaseAndRankCurvature.lean\#L169}{\nolinkurl{fixedRank_cleanWindow_structure}}.  \emph{Compared}, run \href{https://github.com/wcook04/plectis-erdos-lean/actions/runs/35674034595}{\texttt{35674034595}}.
\par\noindent\hangindent=1.5em Proposition~\ref{prop:CP-01-inv}: \href{https://github.com/wcook04/plectis-erdos/blob/181078b6b009d905809cf2e007ad309386a3d1e0/lean/ErdosProblems/Erdos249/PaperCompleteR21/TailCarryPeriodAndRankFloor.lean\#L31}{\nolinkurl{carryShift_dvd_iff_tailDiff_integral}}, \href{https://github.com/wcook04/plectis-erdos/blob/181078b6b009d905809cf2e007ad309386a3d1e0/lean/ErdosProblems/Erdos249/PaperCompleteR21/TailCarryPeriodAndRankFloor.lean\#L20}{\nolinkurl{temperedCarry_eq_scaledTail_and_shift}}.  \emph{Compared}, run \href{https://github.com/wcook04/plectis-erdos-lean/actions/runs/35624228171}{\texttt{35624228171}}.
\par\noindent\hangindent=1.5em Proposition~\ref{prop:CP-02}: \href{https://github.com/wcook04/plectis-erdos/blob/181078b6b009d905809cf2e007ad309386a3d1e0/lean/ErdosProblems/Erdos249/PaperCompleteR21/TailCarryPeriodAndRankFloor.lean\#L46}{\nolinkurl{rationality_forces_mod_period_and_unbounded_rank}}.  \emph{Compared}, run \href{https://github.com/wcook04/plectis-erdos-lean/actions/runs/35624228171}{\texttt{35624228171}}.
\par\noindent\hangindent=1.5em Proposition~\ref{prop:CP-05-inv}: \href{https://github.com/wcook04/plectis-erdos/blob/181078b6b009d905809cf2e007ad309386a3d1e0/lean/ErdosProblems/Erdos249/PaperCompleteR21/TwoAdicPulseBlockAndMobiusInversion.lean\#L28}{\nolinkurl{exists_prime_totient_shift_four_mul_congr_two_mod_four}}.  \emph{Compared}, run \href{https://github.com/wcook04/plectis-erdos-lean/actions/runs/35643815458}{\texttt{35643815458}}.
\par\noindent\hangindent=1.5em Proposition~\ref{prop:TA-inv}: \href{https://github.com/wcook04/plectis-erdos/blob/181078b6b009d905809cf2e007ad309386a3d1e0/lean/ErdosProblems/Erdos249/PaperCompleteR21/TwoAdicPulseBlockAndMobiusInversion.lean\#L81}{\nolinkurl{exists_prime_twoAdic_pulse_block}}, \href{https://github.com/wcook04/plectis-erdos/blob/181078b6b009d905809cf2e007ad309386a3d1e0/lean/ErdosProblems/Erdos249/PaperCompleteR21/TwoAdicPulseBlockAndMobiusInversion.lean\#L41}{\nolinkurl{pulse_delta_of_divisor_data}}, \href{https://github.com/wcook04/plectis-erdos/blob/181078b6b009d905809cf2e007ad309386a3d1e0/lean/ErdosProblems/Erdos249/PaperCompleteR21/TwoAdicPulseBlockAndMobiusInversion.lean\#L100}{\nolinkurl{exists_prime_integral_tailDiff_half_pulse}}.  \emph{Compared}, run \href{https://github.com/wcook04/plectis-erdos-lean/actions/runs/35674034595}{\texttt{35674034595}}.
\par\noindent\hangindent=1.5em Proposition~\ref{prop:MP-01-inv}: \href{https://github.com/wcook04/plectis-erdos/blob/181078b6b009d905809cf2e007ad309386a3d1e0/lean/ErdosProblems/Erdos249/PaperCompleteR21/TwoAdicPulseBlockAndMobiusInversion.lean\#L114}{\nolinkurl{forwardMultiple_spec}}, \href{https://github.com/wcook04/plectis-erdos/blob/181078b6b009d905809cf2e007ad309386a3d1e0/lean/ErdosProblems/Erdos249/PaperCompleteR21/TwoAdicPulseBlockAndMobiusInversion.lean\#L125}{\nolinkurl{forwardMultipleShift_least}}, \href{https://github.com/wcook04/plectis-erdos/blob/181078b6b009d905809cf2e007ad309386a3d1e0/lean/ErdosProblems/Erdos249/PaperCompleteR21/TwoAdicPulseBlockAndMobiusInversion.lean\#L139}{\nolinkurl{forwardMultiple_enumeration}}, \href{https://github.com/wcook04/plectis-erdos/blob/181078b6b009d905809cf2e007ad309386a3d1e0/lean/ErdosProblems/Erdos249/PaperCompleteR21/TwoAdicPulseBlockAndMobiusInversion.lean\#L151}{\nolinkurl{totientTail_eq_tsum_mobius_inversion}}.  \emph{Compared}, run \href{https://github.com/wcook04/plectis-erdos-lean/actions/runs/35624228171}{\texttt{35624228171}}.
\par\noindent\hangindent=1.5em Proposition~\ref{prop:SK-01-inv}: \href{https://github.com/wcook04/plectis-erdos/blob/181078b6b009d905809cf2e007ad309386a3d1e0/lean/ErdosProblems/Erdos249/PaperCompleteR21/RationalTailPeriodWitnesses.lean\#L36}{\nolinkurl{shortWindow_certificates_kill_omega_four_and_six}}, \href{https://github.com/wcook04/plectis-erdos/blob/181078b6b009d905809cf2e007ad309386a3d1e0/lean/ErdosProblems/Erdos249/PaperCompleteR21/RationalTailPeriodWitnesses.lean\#L49}{\nolinkurl{shortWindow_inequalities}}, \href{https://github.com/wcook04/plectis-erdos/blob/181078b6b009d905809cf2e007ad309386a3d1e0/lean/ErdosProblems/Erdos249/PaperCompleteR21/RationalTailPeriodWitnesses.lean\#L27}{\nolinkurl{actualLcmTailOrbit_eq_tail_difference}}.  \emph{Compared}, runs \href{https://github.com/wcook04/plectis-erdos-lean/actions/runs/35643815458}{\texttt{35643815458}}, \href{https://github.com/wcook04/plectis-erdos-lean/actions/runs/35674034595}{\texttt{35674034595}}.
\par\noindent\hangindent=1.5em Proposition~\ref{prop:A5-inv}: \href{https://github.com/wcook04/plectis-erdos/blob/181078b6b009d905809cf2e007ad309386a3d1e0/lean/ErdosProblems/Erdos249/PaperCompleteR20/TailDepthCorrespondence.lean\#L20}{\nolinkurl{certificate_logarithmic_depth}}, \href{https://github.com/wcook04/plectis-erdos/blob/181078b6b009d905809cf2e007ad309386a3d1e0/lean/ErdosProblems/Erdos249/PaperCompleteR20/TailDepthCorrespondence.lean\#L30}{\nolinkurl{fixed_depth_bounds_indices}}, \href{https://github.com/wcook04/plectis-erdos/blob/181078b6b009d905809cf2e007ad309386a3d1e0/lean/Erdos249257/TotientTailPeriodKiller.lean\#L79}{\nolinkurl{certifiedKill_depth_floor}}.  \emph{Compared}, runs \href{https://github.com/wcook04/plectis-erdos-lean/actions/runs/35643815458}{\texttt{35643815458}}, \href{https://github.com/wcook04/plectis-erdos-lean/actions/runs/35674034595}{\texttt{35674034595}}.
\par\noindent\hangindent=1.5em Proposition~\ref{prop:A9-inv}: \href{https://github.com/wcook04/plectis-erdos/blob/181078b6b009d905809cf2e007ad309386a3d1e0/lean/ErdosProblems/Erdos249/PaperCompleteR21/RationalTailPeriodWitnesses.lean\#L65}{\nolinkurl{rational_tail_period_explicit_witnesses}}, \href{https://github.com/wcook04/plectis-erdos/blob/181078b6b009d905809cf2e007ad309386a3d1e0/lean/ErdosProblems/Erdos249/PaperCompleteR21/RationalTailPeriodWitnesses.lean\#L56}{\nolinkurl{totient_one_eq_one}}, \href{https://github.com/wcook04/plectis-erdos/blob/181078b6b009d905809cf2e007ad309386a3d1e0/lean/ErdosProblems/Erdos249/PaperCompleteR21/RationalTailPeriodWitnesses.lean\#L111}{\nolinkurl{eventual_tail_period_of_not_irrational}}, \href{https://github.com/wcook04/plectis-erdos/blob/181078b6b009d905809cf2e007ad309386a3d1e0/lean/ErdosProblems/Erdos249/PaperCompleteR21/RationalTailPeriodWitnesses.lean\#L119}{\nolinkurl{irrational_of_certificate_supply}}, \href{https://github.com/wcook04/plectis-erdos/blob/181078b6b009d905809cf2e007ad309386a3d1e0/lean/ErdosProblems/Erdos249/PaperCompleteR21/RationalTailPeriodWitnesses.lean\#L126}{\nolinkurl{tail_diff_notMem_int_of_irrational}}, \href{https://github.com/wcook04/plectis-erdos/blob/181078b6b009d905809cf2e007ad309386a3d1e0/lean/ErdosProblems/Erdos249/PaperCompleteR21/RationalTailPeriodWitnesses.lean\#L163}{\nolinkurl{exists_certifiedKill_iff_tail_diff_nonintegral}}.  \emph{Compared}, runs \href{https://github.com/wcook04/plectis-erdos-lean/actions/runs/35643815458}{\texttt{35643815458}}, \href{https://github.com/wcook04/plectis-erdos-lean/actions/runs/35674034595}{\texttt{35674034595}}.
\par\noindent\hangindent=1.5em Proposition~\ref{prop:B6-inv}: \href{https://github.com/wcook04/plectis-erdos/blob/181078b6b009d905809cf2e007ad309386a3d1e0/lean/ErdosProblems/Erdos249/PaperCompleteR20/LcmGridCorrespondence.lean\#L8}{\nolinkurl{lcm_grid_flatness}}, \href{https://github.com/wcook04/plectis-erdos/blob/181078b6b009d905809cf2e007ad309386a3d1e0/lean/ErdosProblems/Erdos249/PaperCompleteR20/LcmGridCorrespondence.lean\#L27}{\nolinkurl{certificate_denominator_exclusion}}, \href{https://github.com/wcook04/plectis-erdos/blob/181078b6b009d905809cf2e007ad309386a3d1e0/lean/ErdosProblems/Erdos249/PaperCompleteR20/LcmGridCorrespondence.lean\#L40}{\nolinkurl{lcm_grid_supply_iff}}.  \emph{Compared}, run \href{https://github.com/wcook04/plectis-erdos-lean/actions/runs/35674034595}{\texttt{35674034595}}.
\par\noindent\hangindent=1.5em Proposition~\ref{prop:D4-inv}: \href{https://github.com/wcook04/plectis-erdos/blob/181078b6b009d905809cf2e007ad309386a3d1e0/lean/ErdosProblems/Erdos249/PaperCompleteR21/CanonicalDyadicSectionRank.lean\#L21}{\nolinkurl{card_canonical_dyadic_index}}, \href{https://github.com/wcook04/plectis-erdos/blob/181078b6b009d905809cf2e007ad309386a3d1e0/lean/ErdosProblems/Erdos249/PaperCompleteR21/CanonicalDyadicSectionRank.lean\#L26}{\nolinkurl{linearIndependent_canonical_dyadic_family}}, \href{https://github.com/wcook04/plectis-erdos/blob/181078b6b009d905809cf2e007ad309386a3d1e0/lean/ErdosProblems/Erdos249/PaperCompleteR21/CanonicalDyadicSectionRank.lean\#L33}{\nolinkurl{canonical_family_basis_through_level}}, \href{https://github.com/wcook04/plectis-erdos/blob/181078b6b009d905809cf2e007ad309386a3d1e0/lean/ErdosProblems/Erdos249/PaperCompleteR21/CanonicalDyadicSectionRank.lean\#L46}{\nolinkurl{canonical_level_zero_channels}}, \href{https://github.com/wcook04/plectis-erdos/blob/181078b6b009d905809cf2e007ad309386a3d1e0/lean/ErdosProblems/Erdos249/PaperCompleteR21/CanonicalDyadicSectionRank.lean\#L59}{\nolinkurl{finrank_throughLevel_zero_eq_one}}.  \emph{Compared}, run \href{https://github.com/wcook04/plectis-erdos-lean/actions/runs/35624228171}{\texttt{35624228171}}.
\par\noindent\hangindent=1.5em Proposition~\ref{prop:C2-inv}: \href{https://github.com/wcook04/plectis-erdos/blob/181078b6b009d905809cf2e007ad309386a3d1e0/lean/ErdosProblems/Erdos249/PaperCompleteR21/FareyExactRangeAndDenominatorExclusion.lean\#L26}{\nolinkurl{farey_window_1_240_exact_range}}, \href{https://github.com/wcook04/plectis-erdos/blob/181078b6b009d905809cf2e007ad309386a3d1e0/lean/ErdosProblems/Erdos249/PaperCompleteR21/FareyExactRangeAndDenominatorExclusion.lean\#L39}{\nolinkurl{farey_window_1_240_range_magnitude}}.  \emph{Compared}, run \href{https://github.com/wcook04/plectis-erdos-lean/actions/runs/35643815458}{\texttt{35643815458}}.
\par\noindent\hangindent=1.5em Proposition~\ref{prop:C3-inv}: \href{https://github.com/wcook04/plectis-erdos/blob/181078b6b009d905809cf2e007ad309386a3d1e0/lean/ErdosProblems/Erdos249/PaperCompleteR21/FareyExactRangeAndDenominatorExclusion.lean\#L49}{\nolinkurl{totient_series_ne_reduced_fraction_of_small_denominator}}, \href{https://github.com/wcook04/plectis-erdos/blob/181078b6b009d905809cf2e007ad309386a3d1e0/lean/ErdosProblems/Erdos249/PaperCompleteR21/FareyExactRangeAndDenominatorExclusion.lean\#L57}{\nolinkurl{totient_series_ne_int_div_of_small_denominator}}.  \emph{Compared}, run \href{https://github.com/wcook04/plectis-erdos-lean/actions/runs/35643815458}{\texttt{35643815458}}.
\par\noindent\hangindent=1.5em Proposition~\ref{prop:D1D2-inv} (the Lean statement is stronger): \href{https://github.com/wcook04/plectis-erdos/blob/181078b6b009d905809cf2e007ad309386a3d1e0/lean/Erdos249257/CertificateKernel.lean\#L5371}{\nolinkurl{irrational_of_den_mul_abs_sub_tendsto_zero}}, \href{https://github.com/wcook04/plectis-erdos/blob/181078b6b009d905809cf2e007ad309386a3d1e0/lean/Erdos249257/CertificateKernel.lean\#L5342}{\nolinkurl{one_div_den_mul_den_le_abs_sub}}, \href{https://github.com/wcook04/plectis-erdos/blob/181078b6b009d905809cf2e007ad309386a3d1e0/lean/Erdos249257/CertificateKernel.lean\#L6120}{\nolinkurl{irrational_of_int_mul_near_int}}, \href{https://github.com/wcook04/plectis-erdos/blob/181078b6b009d905809cf2e007ad309386a3d1e0/lean/Erdos249257/CertificateKernel.lean\#L6149}{\nolinkurl{irrational_of_pow_mul_near_int}}, \href{https://github.com/wcook04/plectis-erdos/blob/181078b6b009d905809cf2e007ad309386a3d1e0/lean/ErdosProblems/Erdos249/PaperCompleteR21/GeneralIrrationalityCriteriaAndGapBounds.lean\#L66}{\nolinkurl{irrational_of_den_mul_error_tendsto_zero}}, \href{https://github.com/wcook04/plectis-erdos/blob/181078b6b009d905809cf2e007ad309386a3d1e0/lean/ErdosProblems/Erdos249/PaperCompleteR21/GeneralIrrationalityCriteriaAndGapBounds.lean\#L75}{\nolinkurl{irrational_of_dirichlet_gap}}, \href{https://github.com/wcook04/plectis-erdos/blob/181078b6b009d905809cf2e007ad309386a3d1e0/lean/ErdosProblems/Erdos249/PaperCompleteR21/GeneralIrrationalityCriteriaAndGapBounds.lean\#L27}{\nolinkurl{one_div_den_mul_den_le_abs_diff}}, \href{https://github.com/wcook04/plectis-erdos/blob/181078b6b009d905809cf2e007ad309386a3d1e0/lean/ErdosProblems/Erdos249/PaperCompleteR21/GeneralIrrationalityCriteriaAndGapBounds.lean\#L33}{\nolinkurl{one_div_den_le_abs_int_combination}}, \href{https://github.com/wcook04/plectis-erdos/blob/181078b6b009d905809cf2e007ad309386a3d1e0/lean/ErdosProblems/Erdos249/PaperCompleteR21/GeneralIrrationalityCriteriaAndGapBounds.lean\#L92}{\nolinkurl{irrational_of_basePower_dilation}}, \href{https://github.com/wcook04/plectis-erdos/blob/181078b6b009d905809cf2e007ad309386a3d1e0/lean/ErdosProblems/Erdos249/PaperCompleteR21/SquareBlockBinaryDilationCountermodel.lean\#L529}{\nolinkurl{basePower_dilation_not_universal}}, \href{https://github.com/wcook04/plectis-erdos/blob/181078b6b009d905809cf2e007ad309386a3d1e0/lean/ErdosProblems/Erdos249/PaperCompleteR21/SquareBlockBinaryDilationCountermodel.lean\#L476}{\nolinkurl{irrational_xi}}, \href{https://github.com/wcook04/plectis-erdos/blob/181078b6b009d905809cf2e007ad309386a3d1e0/lean/ErdosProblems/Erdos249/PaperCompleteR21/SquareBlockBinaryDilationCountermodel.lean\#L486}{\nolinkurl{one_div_eight_le_dist_xi}}.  \emph{Compared}, runs \href{https://github.com/wcook04/plectis-erdos-lean/actions/runs/35624228171}{\texttt{35624228171}}, \href{https://github.com/wcook04/plectis-erdos-lean/actions/runs/35643815458}{\texttt{35643815458}}.
\par\noindent\hangindent=1.5em Proposition~\ref{prop:D9-inv}: \href{https://github.com/wcook04/plectis-erdos/blob/181078b6b009d905809cf2e007ad309386a3d1e0/lean/ErdosProblems/Erdos249/PaperCompleteR21/GeneralIrrationalityCriteriaAndGapBounds.lean\#L138}{\nolinkurl{rational_gap_lower_bound}}, \href{https://github.com/wcook04/plectis-erdos/blob/181078b6b009d905809cf2e007ad309386a3d1e0/lean/ErdosProblems/Erdos249/PaperCompleteR21/GeneralIrrationalityCriteriaAndGapBounds.lean\#L146}{\nolinkurl{den_lower_bound_of_positive_error}}, \href{https://github.com/wcook04/plectis-erdos/blob/181078b6b009d905809cf2e007ad309386a3d1e0/lean/ErdosProblems/Erdos249/PaperCompleteR21/GeneralIrrationalityCriteriaAndGapBounds.lean\#L164}{\nolinkurl{denominator_bound_tendsto_atTop_iff}}, \href{https://github.com/wcook04/plectis-erdos/blob/181078b6b009d905809cf2e007ad309386a3d1e0/lean/ErdosProblems/Erdos249/PaperCompleteR21/GeneralIrrationalityCriteriaAndGapBounds.lean\#L189}{\nolinkurl{irrational_of_den_mul_error_product_tendsto_zero}}, \href{https://github.com/wcook04/plectis-erdos/blob/181078b6b009d905809cf2e007ad309386a3d1e0/lean/ErdosProblems/Erdos249/PaperCompleteR21/GeneralIrrationalityCriteriaAndGapBounds.lean\#L201}{\nolinkurl{dyadic_prefix_denominator_bound_vacuous}}, \href{https://github.com/wcook04/plectis-erdos/blob/181078b6b009d905809cf2e007ad309386a3d1e0/lean/ErdosProblems/Erdos249/PaperCompleteR21/DyadicPrefixTailBound.lean\#L14}{\nolinkurl{dyadic_prefix_den_dvd}}, \href{https://github.com/wcook04/plectis-erdos/blob/181078b6b009d905809cf2e007ad309386a3d1e0/lean/ErdosProblems/Erdos249/PaperCompleteR21/DyadicPrefixTailBound.lean\#L22}{\nolinkurl{totientTail_le_add_two}}, \href{https://github.com/wcook04/plectis-erdos/blob/181078b6b009d905809cf2e007ad309386a3d1e0/lean/ErdosProblems/Erdos249/PaperCompleteR21/DyadicPrefixTailBound.lean\#L28}{\nolinkurl{dyadic_prefix_tail_le}}.  \emph{Compared}, runs \href{https://github.com/wcook04/plectis-erdos-lean/actions/runs/35643815458}{\texttt{35643815458}}, \href{https://github.com/wcook04/plectis-erdos-lean/actions/runs/35674034595}{\texttt{35674034595}}.
\par\noindent\hangindent=1.5em Proposition~\ref{prop:B11-inv}: \href{https://github.com/wcook04/plectis-erdos/blob/181078b6b009d905809cf2e007ad309386a3d1e0/lean/ErdosProblems/Erdos249/PaperCompleteR20/FiniteCertificateBatch.lean\#L57}{\nolinkurl{historical_table_size_and_initial_depths}}, \href{https://github.com/wcook04/plectis-erdos/blob/181078b6b009d905809cf2e007ad309386a3d1e0/lean/ErdosProblems/Erdos249/PaperCompleteR20/FiniteCertificateBatch.lean\#L64}{\nolinkurl{historical_table_and_complete_band}}.  \emph{Compared}, run \href{https://github.com/wcook04/plectis-erdos-lean/actions/runs/35674034595}{\texttt{35674034595}}.
\par\noindent\hangindent=1.5em Proposition~\ref{prop:C1-inv}: \href{https://github.com/wcook04/plectis-erdos/blob/181078b6b009d905809cf2e007ad309386a3d1e0/lean/ErdosProblems/Erdos249/PaperCompleteR21/GeneralIrrationalityCriteriaAndGapBounds.lean\#L218}{\nolinkurl{card_visible_antidiagonal}}, \href{https://github.com/wcook04/plectis-erdos/blob/181078b6b009d905809cf2e007ad309386a3d1e0/lean/ErdosProblems/Erdos249/PaperCompleteR21/GeneralIrrationalityCriteriaAndGapBounds.lean\#L225}{\nolinkurl{totient_zero_eq_zero}}, \href{https://github.com/wcook04/plectis-erdos/blob/181078b6b009d905809cf2e007ad309386a3d1e0/lean/ErdosProblems/Erdos249/PaperCompleteR21/GeneralIrrationalityCriteriaAndGapBounds.lean\#L228}{\nolinkurl{visible_antidiagonal_one}}, \href{https://github.com/wcook04/plectis-erdos/blob/181078b6b009d905809cf2e007ad309386a3d1e0/lean/ErdosProblems/Erdos249/PaperCompleteR21/GeneralIrrationalityCriteriaAndGapBounds.lean\#L234}{\nolinkurl{positive_antidiagonal_one}}, \href{https://github.com/wcook04/plectis-erdos/blob/181078b6b009d905809cf2e007ad309386a3d1e0/lean/ErdosProblems/Erdos249/PaperCompleteR21/GeneralIrrationalityCriteriaAndGapBounds.lean\#L241}{\nolinkurl{tsum_pos_coprime_pairs_eq_series_sub_half}}, \href{https://github.com/wcook04/plectis-erdos/blob/181078b6b009d905809cf2e007ad309386a3d1e0/lean/ErdosProblems/Erdos249/PaperCompleteR21/GeneralIrrationalityCriteriaAndGapBounds.lean\#L266}{\nolinkurl{tsum_pos_coprime_pairs_product_form}}.  \emph{Compared}, run \href{https://github.com/wcook04/plectis-erdos-lean/actions/runs/35643815458}{\texttt{35643815458}}.
\par\noindent\hangindent=1.5em Theorem (The values at exponents one and two): \href{https://github.com/wcook04/plectis-erdos/blob/181078b6b009d905809cf2e007ad309386a3d1e0/lean/ErdosProblems/Erdos249/PaperCompleteR21/MobiusMersenneLadderLogConcavity.lean\#L23}{\nolinkurl{mobiusMersenneTheta_one_and_two}}.  \emph{Compared}, run \href{https://github.com/wcook04/plectis-erdos-lean/actions/runs/35643815458}{\texttt{35643815458}}.
\par\noindent\hangindent=1.5em Theorem (The first two summands give a positive Hankel gap): \href{https://github.com/wcook04/plectis-erdos/blob/181078b6b009d905809cf2e007ad309386a3d1e0/lean/ErdosProblems/Erdos249/PaperCompleteR21/MobiusMersenneLadderLogConcavity.lean\#L35}{\nolinkurl{twoAtom_hankel_gap}}, \href{https://github.com/wcook04/plectis-erdos/blob/181078b6b009d905809cf2e007ad309386a3d1e0/lean/ErdosProblems/Erdos249/PaperCompleteR21/MobiusMersenneLadderLogConcavity.lean\#L43}{\nolinkurl{twoAtom_strict_logConcave}}.  \emph{Compared}, run \href{https://github.com/wcook04/plectis-erdos-lean/actions/runs/35643815458}{\texttt{35643815458}}.
\par\noindent\hangindent=1.5em Theorem (Strict log-concavity for all integer $r\ge1$): \href{https://github.com/wcook04/plectis-erdos/blob/181078b6b009d905809cf2e007ad309386a3d1e0/lean/ErdosProblems/Erdos249/PaperCompleteR21/MobiusMersenneLadderLogConcavity.lean\#L52}{\nolinkurl{theta_strict_logConcave}}, \href{https://github.com/wcook04/plectis-erdos/blob/181078b6b009d905809cf2e007ad309386a3d1e0/lean/ErdosProblems/Erdos249/PaperCompleteR21/MobiusMersenneLadderLogConcavity.lean\#L58}{\nolinkurl{theta_hankel_det_neg}}, \href{https://github.com/wcook04/plectis-erdos/blob/181078b6b009d905809cf2e007ad309386a3d1e0/lean/ErdosProblems/Erdos249/PaperCompleteR21/MobiusMersenneLadderLogConcavity.lean\#L66}{\nolinkurl{theta_hankel_two_neg}}.  \emph{Compared}, run \href{https://github.com/wcook04/plectis-erdos-lean/actions/runs/35643815458}{\texttt{35643815458}}.
\par\noindent\hangindent=1.5em Theorem~\ref{thm:binary-carry-criterion}: \href{https://github.com/wcook04/plectis-erdos/blob/181078b6b009d905809cf2e007ad309386a3d1e0/lean/Erdos249257/GenericTailOrbitRigidity.lean\#L426}{\nolinkurl{binaryCoeffSeries_rational_iff_exists_temperedBinaryOrbit}}, \href{https://github.com/wcook04/plectis-erdos/blob/181078b6b009d905809cf2e007ad309386a3d1e0/lean/Erdos249257/GenericTailOrbitRigidity.lean\#L339}{\nolinkurl{temperedBinaryOrbit_eq_scaledTail}}.  \emph{Compared}, run \href{https://github.com/wcook04/plectis-erdos-lean/actions/runs/35643815458}{\texttt{35643815458}}.
\par\noindent\hangindent=1.5em Proposition (Uniqueness under the growth condition): \href{https://github.com/wcook04/plectis-erdos/blob/181078b6b009d905809cf2e007ad309386a3d1e0/lean/ErdosProblems/Erdos249/PaperCompleteR21/TemperedOrbitAndSquaredMersenneTail.lean\#L28}{\nolinkurl{doubling_tempered_sequence_eq_zero}}.  \emph{Compared}, run \href{https://github.com/wcook04/plectis-erdos-lean/actions/runs/35643815458}{\texttt{35643815458}}.
\par\noindent\hangindent=1.5em Proposition (An exact rational approximation formula): \href{https://github.com/wcook04/plectis-erdos/blob/181078b6b009d905809cf2e007ad309386a3d1e0/lean/ErdosProblems/Erdos249/PaperCompleteR21/TemperedOrbitAndSquaredMersenneTail.lean\#L118}{\nolinkurl{tailDifference_sub_rationalApproximation}}, \href{https://github.com/wcook04/plectis-erdos/blob/181078b6b009d905809cf2e007ad309386a3d1e0/lean/ErdosProblems/Erdos249/PaperCompleteR21/TemperedOrbitAndSquaredMersenneTail.lean\#L88}{\nolinkurl{tailDifference_eq_coefficient_mul_series}}, \href{https://github.com/wcook04/plectis-erdos/blob/181078b6b009d905809cf2e007ad309386a3d1e0/lean/ErdosProblems/Erdos249/PaperCompleteR21/TemperedOrbitAndSquaredMersenneTail.lean\#L99}{\nolinkurl{totientSeries_pnat_form}}, \href{https://github.com/wcook04/plectis-erdos/blob/181078b6b009d905809cf2e007ad309386a3d1e0/lean/ErdosProblems/Erdos249/PaperCompleteR21/TemperedOrbitAndSquaredMersenneTail.lean\#L106}{\nolinkurl{totientSeries_eq_half_add_moebius_sq}}.  \emph{Compared}, runs \href{https://github.com/wcook04/plectis-erdos-lean/actions/runs/35643815458}{\texttt{35643815458}}, \href{https://github.com/wcook04/plectis-erdos-lean/actions/runs/35674034595}{\texttt{35674034595}}.
\par\noindent\hangindent=1.5em Proposition (A geometric tail bound): \href{https://github.com/wcook04/plectis-erdos/blob/181078b6b009d905809cf2e007ad309386a3d1e0/lean/ErdosProblems/Erdos249/PaperCompleteR21/TemperedOrbitAndSquaredMersenneTail.lean\#L155}{\nolinkurl{abs_mobiusSquareTail_le_paper}}, \href{https://github.com/wcook04/plectis-erdos/blob/181078b6b009d905809cf2e007ad309386a3d1e0/lean/ErdosProblems/Erdos249/PaperCompleteR21/TemperedOrbitAndSquaredMersenneTail.lean\#L138}{\nolinkurl{abs_moebius_cast_le_one}}, \href{https://github.com/wcook04/plectis-erdos/blob/181078b6b009d905809cf2e007ad309386a3d1e0/lean/ErdosProblems/Erdos249/PaperCompleteR21/TemperedOrbitAndSquaredMersenneTail.lean\#L144}{\nolinkurl{mersenne_geometric_shift}}, \href{https://github.com/wcook04/plectis-erdos/blob/181078b6b009d905809cf2e007ad309386a3d1e0/lean/ErdosProblems/Erdos249/PaperCompleteR21/TemperedOrbitAndSquaredMersenneTail.lean\#L149}{\nolinkurl{tsum_quarter_geometric}}.  \emph{Compared}, run \href{https://github.com/wcook04/plectis-erdos-lean/actions/runs/35643815458}{\texttt{35643815458}}.
\par\noindent\hangindent=1.5em Proposition (The sum over divisor indices): \href{https://github.com/wcook04/plectis-erdos/blob/181078b6b009d905809cf2e007ad309386a3d1e0/lean/ErdosProblems/Erdos249/PaperCompleteR21/DivisorChannelSplitAndSeamDoubling.lean\#L68}{\nolinkurl{sum_divisorIndices_mobius}}, \href{https://github.com/wcook04/plectis-erdos/blob/181078b6b009d905809cf2e007ad309386a3d1e0/lean/ErdosProblems/Erdos249/PaperCompleteR21/DivisorChannelSplitAndSeamDoubling.lean\#L98}{\nolinkurl{sum_divisors_moebius_div_eq_totient_div}}, \href{https://github.com/wcook04/plectis-erdos/blob/181078b6b009d905809cf2e007ad309386a3d1e0/lean/ErdosProblems/Erdos249/PaperCompleteR21/DivisorChannelSplitAndSeamDoubling.lean\#L121}{\nolinkurl{sum_divisorIndices_radical_form}}, \href{https://github.com/wcook04/plectis-erdos/blob/181078b6b009d905809cf2e007ad309386a3d1e0/lean/ErdosProblems/Erdos249/PaperCompleteR21/DivisorChannelSplitAndSeamDoubling.lean\#L117}{\nolinkurl{squarefreeKernel_eq_prod_primeFactors}}.  \emph{Compared}, runs \href{https://github.com/wcook04/plectis-erdos-lean/actions/runs/35624228171}{\texttt{35624228171}}, \href{https://github.com/wcook04/plectis-erdos-lean/actions/runs/35643815458}{\texttt{35643815458}}.
\par\noindent\hangindent=1.5em Proposition (The divisor sum and its complement): \href{https://github.com/wcook04/plectis-erdos/blob/181078b6b009d905809cf2e007ad309386a3d1e0/lean/ErdosProblems/Erdos249/PaperCompleteR21/DivisorChannelSplitAndSeamDoubling.lean\#L217}{\nolinkurl{totientDifference_eq_divisorPart_add_complement}}, \href{https://github.com/wcook04/plectis-erdos/blob/181078b6b009d905809cf2e007ad309386a3d1e0/lean/ErdosProblems/Erdos249/PaperCompleteR21/DivisorChannelSplitAndSeamDoubling.lean\#L134}{\nolinkurl{totient_eq_mobius_divisor_sum}}, \href{https://github.com/wcook04/plectis-erdos/blob/181078b6b009d905809cf2e007ad309386a3d1e0/lean/ErdosProblems/Erdos249/PaperCompleteR21/DivisorChannelSplitAndSeamDoubling.lean\#L141}{\nolinkurl{divisorIndex_endpoint_behaviour}}, \href{https://github.com/wcook04/plectis-erdos/blob/181078b6b009d905809cf2e007ad309386a3d1e0/lean/ErdosProblems/Erdos249/PaperCompleteR21/DivisorChannelSplitAndSeamDoubling.lean\#L159}{\nolinkurl{complementSummand_eq_phaseTerm}}.  \emph{Compared}, run \href{https://github.com/wcook04/plectis-erdos-lean/actions/runs/35643815458}{\texttt{35643815458}}.
\par\noindent\hangindent=1.5em Theorem (Doubling the full totient difference): \href{https://github.com/wcook04/plectis-erdos/blob/181078b6b009d905809cf2e007ad309386a3d1e0/lean/ErdosProblems/Erdos249/PaperCompleteR21/DivisorChannelSplitAndSeamDoubling.lean\#L262}{\nolinkurl{totientDifference_doubling_seam}}, \href{https://github.com/wcook04/plectis-erdos/blob/181078b6b009d905809cf2e007ad309386a3d1e0/lean/ErdosProblems/Erdos249/PaperCompleteR21/DivisorChannelSplitAndSeamDoubling.lean\#L250}{\nolinkurl{totient_two_mul_even}}, \href{https://github.com/wcook04/plectis-erdos/blob/181078b6b009d905809cf2e007ad309386a3d1e0/lean/ErdosProblems/Erdos249/PaperCompleteR21/DivisorChannelSplitAndSeamDoubling.lean\#L255}{\nolinkurl{totient_two_mul_odd}}.  \emph{Compared}, run \href{https://github.com/wcook04/plectis-erdos-lean/actions/runs/35643815458}{\texttt{35643815458}}.
\par\noindent\hangindent=1.5em Proposition (A doubling identity for two portions of the sum): \href{https://github.com/wcook04/plectis-erdos/blob/181078b6b009d905809cf2e007ad309386a3d1e0/lean/ErdosProblems/Erdos249/PaperCompleteR21/DivisorChannelSplitAndSeamDoubling.lean\#L289}{\nolinkurl{complementSummand_low_double_echo}}.  \emph{Compared}, run \href{https://github.com/wcook04/plectis-erdos-lean/actions/runs/35643815458}{\texttt{35643815458}}.
\par\noindent\hangindent=1.5em Theorem (Positive coefficients of the numerator polynomial): \href{https://github.com/wcook04/plectis-erdos/blob/181078b6b009d905809cf2e007ad309386a3d1e0/lean/ErdosProblems/Erdos249/PaperCompleteR21/NumeratorPolynomialAndMersenneRemainder.lean\#L67}{\nolinkurl{paperNumerator_eq_gcdWordForm}}, \href{https://github.com/wcook04/plectis-erdos/blob/181078b6b009d905809cf2e007ad309386a3d1e0/lean/ErdosProblems/Erdos249/PaperCompleteR21/NumeratorPolynomialAndMersenneRemainder.lean\#L82}{\nolinkurl{paperNumerator_coeff_pos}}, \href{https://github.com/wcook04/plectis-erdos/blob/181078b6b009d905809cf2e007ad309386a3d1e0/lean/ErdosProblems/Erdos249/PaperCompleteR21/NumeratorPolynomialAndMersenneRemainder.lean\#L88}{\nolinkurl{paperNumerator_coeff_eq_zero}}, \href{https://github.com/wcook04/plectis-erdos/blob/181078b6b009d905809cf2e007ad309386a3d1e0/lean/ErdosProblems/Erdos249/PaperCompleteR21/NumeratorPolynomialAndMersenneRemainder.lean\#L94}{\nolinkurl{paperNumerator_coeff_eq_gcd_divisor_sum}}, \href{https://github.com/wcook04/plectis-erdos/blob/181078b6b009d905809cf2e007ad309386a3d1e0/lean/ErdosProblems/Erdos249/PaperCompleteR21/NumeratorPolynomialAndMersenneRemainder.lean\#L57}{\nolinkurl{paperNumeratorPolynomial_eq}}.  \emph{Compared}, run \href{https://github.com/wcook04/plectis-erdos-lean/actions/runs/35624228171}{\texttt{35624228171}}.
\par\noindent\hangindent=1.5em Proposition (The geometric remainder bound): \href{https://github.com/wcook04/plectis-erdos/blob/181078b6b009d905809cf2e007ad309386a3d1e0/lean/ErdosProblems/Erdos249/PaperCompleteR21/NumeratorPolynomialAndMersenneRemainder.lean\#L166}{\nolinkurl{mersenneRemainder_identity_and_bound}}, \href{https://github.com/wcook04/plectis-erdos/blob/181078b6b009d905809cf2e007ad309386a3d1e0/lean/ErdosProblems/Erdos249/PaperCompleteR21/NumeratorPolynomialAndMersenneRemainder.lean\#L156}{\nolinkurl{one_sub_half_pow_ge}}, \href{https://github.com/wcook04/plectis-erdos/blob/181078b6b009d905809cf2e007ad309386a3d1e0/lean/ErdosProblems/Erdos249/PaperCompleteR21/NumeratorPolynomialAndMersenneRemainder.lean\#L208}{\nolinkurl{mersenneRemainderTail_le}}, \href{https://github.com/wcook04/plectis-erdos/blob/181078b6b009d905809cf2e007ad309386a3d1e0/lean/ErdosProblems/Erdos249/PaperCompleteR21/NumeratorPolynomialAndMersenneRemainder.lean\#L182}{\nolinkurl{tsum_eighth_pow_tail}}, \href{https://github.com/wcook04/plectis-erdos/blob/181078b6b009d905809cf2e007ad309386a3d1e0/lean/ErdosProblems/Erdos249/PaperCompleteR21/NumeratorPolynomialAndMersenneRemainder.lean\#L200}{\nolinkurl{four_thirds_tsum_eighth_pow_tail}}.  \emph{Compared}, run \href{https://github.com/wcook04/plectis-erdos-lean/actions/runs/35624228171}{\texttt{35624228171}}.
\par\noindent\hangindent=1.5em Proposition (Squared distance from the phase one): \href{https://github.com/wcook04/plectis-erdos/blob/181078b6b009d905809cf2e007ad309386a3d1e0/lean/ErdosProblems/Erdos249/PaperCompleteR21/PhaseEnergyAndForeignResidueProjection.lean\#L31}{\nolinkurl{sum_sq_dist_from_phase_one}}.  \emph{Compared}, run \href{https://github.com/wcook04/plectis-erdos-lean/actions/runs/35674034595}{\texttt{35674034595}}.
\par\noindent\hangindent=1.5em Lemma (Squared-distance bound for separated pairs): \href{https://github.com/wcook04/plectis-erdos/blob/181078b6b009d905809cf2e007ad309386a3d1e0/lean/ErdosProblems/Erdos249/PaperCompleteR21/PhaseEnergyAndForeignResidueProjection.lean\#L40}{\nolinkurl{card_mul_sq_le_pairwise_energy}}.  \emph{Compared}, run \href{https://github.com/wcook04/plectis-erdos-lean/actions/runs/35643815458}{\texttt{35643815458}}.
\par\noindent\hangindent=1.5em Proposition (The finite divisor sum): \href{https://github.com/wcook04/plectis-erdos/blob/181078b6b009d905809cf2e007ad309386a3d1e0/lean/ErdosProblems/Erdos249/PaperCompleteR21/PhaseEnergyAndForeignResidueProjection.lean\#L69}{\nolinkurl{divisorChannels_sum_eq}}, \href{https://github.com/wcook04/plectis-erdos/blob/181078b6b009d905809cf2e007ad309386a3d1e0/lean/ErdosProblems/Erdos249/PaperCompleteR21/PhaseEnergyAndForeignResidueProjection.lean\#L50}{\nolinkurl{residueOffset_of_dvd}}, \href{https://github.com/wcook04/plectis-erdos/blob/181078b6b009d905809cf2e007ad309386a3d1e0/lean/ErdosProblems/Erdos249/PaperCompleteR21/PhaseEnergyAndForeignResidueProjection.lean\#L58}{\nolinkurl{residueKernel_increment_of_dvd}}, \href{https://github.com/wcook04/plectis-erdos/blob/181078b6b009d905809cf2e007ad309386a3d1e0/lean/ErdosProblems/Erdos249/PaperCompleteR21/PhaseEnergyAndForeignResidueProjection.lean\#L84}{\nolinkurl{scaleExplicitShadow_eq_divisorChannels}}.  \emph{Compared}, runs \href{https://github.com/wcook04/plectis-erdos-lean/actions/runs/35624228171}{\texttt{35624228171}}, \href{https://github.com/wcook04/plectis-erdos-lean/actions/runs/35643815458}{\texttt{35643815458}}.
\par\noindent\hangindent=1.5em Proposition (Separation larger than the error implies exclusion): \href{https://github.com/wcook04/plectis-erdos/blob/181078b6b009d905809cf2e007ad309386a3d1e0/lean/ErdosProblems/Erdos249/PaperCompleteR21/PhaseEnergyAndForeignResidueProjection.lean\#L119}{\nolinkurl{tailDifference_not_integral_of_separation}}, \href{https://github.com/wcook04/plectis-erdos/blob/181078b6b009d905809cf2e007ad309386a3d1e0/lean/ErdosProblems/Erdos249/PaperCompleteR21/PhaseEnergyAndForeignResidueProjection.lean\#L103}{\nolinkurl{projectedForeignDefect_paper}}, \href{https://github.com/wcook04/plectis-erdos/blob/181078b6b009d905809cf2e007ad309386a3d1e0/lean/ErdosProblems/Erdos249/PaperCompleteR21/PhaseEnergyAndForeignResidueProjection.lean\#L110}{\nolinkurl{foreignComplementBound_paper}}.  \emph{Compared}, runs \href{https://github.com/wcook04/plectis-erdos-lean/actions/runs/35624228171}{\texttt{35624228171}}, \href{https://github.com/wcook04/plectis-erdos-lean/actions/runs/35643815458}{\texttt{35643815458}}.
\par\noindent\hangindent=1.5em Theorem (Coprime-pair counting and the totient): \href{https://github.com/wcook04/plectis-erdos/blob/181078b6b009d905809cf2e007ad309386a3d1e0/lean/ErdosProblems/Erdos249/PaperCompleteR21/PhaseEnergyAndForeignResidueProjection.lean\#L147}{\nolinkurl{card_coprime_antidiagonal}}, \href{https://github.com/wcook04/plectis-erdos/blob/181078b6b009d905809cf2e007ad309386a3d1e0/lean/ErdosProblems/Erdos249/PaperCompleteR21/PhaseEnergyAndForeignResidueProjection.lean\#L154}{\nolinkurl{boundary_pair_at_one}}.  \emph{Compared}, run \href{https://github.com/wcook04/plectis-erdos-lean/actions/runs/35643815458}{\texttt{35643815458}}.
\par\noindent\hangindent=1.5em Proposition (Two lattice sums): \href{https://github.com/wcook04/plectis-erdos/blob/181078b6b009d905809cf2e007ad309386a3d1e0/lean/ErdosProblems/Erdos249/PaperCompleteR21/CoprimeLatticeSumsAndLambert.lean\#L27}{\nolinkurl{tsum_totient_pow_shift}}, \href{https://github.com/wcook04/plectis-erdos/blob/181078b6b009d905809cf2e007ad309386a3d1e0/lean/ErdosProblems/Erdos249/PaperCompleteR21/CoprimeLatticeSumsAndLambert.lean\#L38}{\nolinkurl{coprimeLattice_halfOpen_sum}}, \href{https://github.com/wcook04/plectis-erdos/blob/181078b6b009d905809cf2e007ad309386a3d1e0/lean/ErdosProblems/Erdos249/PaperCompleteR21/CoprimeLatticeSumsAndLambert.lean\#L46}{\nolinkurl{coprimeLattice_positive_sum}}, \href{https://github.com/wcook04/plectis-erdos/blob/181078b6b009d905809cf2e007ad309386a3d1e0/lean/ErdosProblems/Erdos249/PaperCompleteR21/CoprimeLatticeSumsAndLambert.lean\#L53}{\nolinkurl{coprimeLattice_removed_pair}}, \href{https://github.com/wcook04/plectis-erdos/blob/181078b6b009d905809cf2e007ad309386a3d1e0/lean/ErdosProblems/Erdos249/PaperCompleteR21/CoprimeLatticeSumsAndLambert.lean\#L77}{\nolinkurl{coprimeLattice_gcd_layer_total}}, \href{https://github.com/wcook04/plectis-erdos/blob/181078b6b009d905809cf2e007ad309386a3d1e0/lean/ErdosProblems/Erdos249/PaperCompleteR21/CoprimeLatticeSumsAndLambert.lean\#L84}{\nolinkurl{coprimeLattice_gcd_layer_total_eq_one_iff}}.  \emph{Compared}, run \href{https://github.com/wcook04/plectis-erdos-lean/actions/runs/35643815458}{\texttt{35643815458}}.
\par\noindent\hangindent=1.5em Theorem (The classical coprime-pair Lambert identity): \href{https://github.com/wcook04/plectis-erdos/blob/181078b6b009d905809cf2e007ad309386a3d1e0/lean/ErdosProblems/Erdos249/PaperCompleteR21/CoprimeLatticeSumsAndLambert.lean\#L106}{\nolinkurl{coprimeLattice_lambert_identity}}, \href{https://github.com/wcook04/plectis-erdos/blob/181078b6b009d905809cf2e007ad309386a3d1e0/lean/ErdosProblems/Erdos249/PaperCompleteR21/CoprimeLatticeSumsAndLambert.lean\#L114}{\nolinkurl{coprimeLattice_lambert_rational}}, \href{https://github.com/wcook04/plectis-erdos/blob/181078b6b009d905809cf2e007ad309386a3d1e0/lean/ErdosProblems/Erdos249/PaperCompleteR21/CoprimeLatticeSumsAndLambert.lean\#L127}{\nolinkurl{coprimeLattice_lambert_half_eq_one}}, \href{https://github.com/wcook04/plectis-erdos/blob/181078b6b009d905809cf2e007ad309386a3d1e0/lean/ErdosProblems/Erdos249/PaperCompleteR21/CoprimeLatticeSumsAndLambert.lean\#L46}{\nolinkurl{coprimeLattice_positive_sum}}, \href{https://github.com/wcook04/plectis-erdos/blob/181078b6b009d905809cf2e007ad309386a3d1e0/lean/ErdosProblems/Erdos249/PaperCompleteR21/CoprimeLatticeSumsAndLambert.lean\#L136}{\nolinkurl{coprimeLattice_plain_half_eq_series_sub_half}}, \href{https://github.com/wcook04/plectis-erdos/blob/181078b6b009d905809cf2e007ad309386a3d1e0/lean/ErdosProblems/Erdos249/PaperCompleteR21/CoprimeLatticeSumsAndLambert.lean\#L155}{\nolinkurl{coprimeLattice_halfOpen_half_eq_series}}.  \emph{Compared}, run \href{https://github.com/wcook04/plectis-erdos-lean/actions/runs/35643815458}{\texttt{35643815458}}.
\par\noindent\hangindent=1.5em Theorem (Nondivisors in a short LCM window): \href{https://github.com/wcook04/plectis-erdos/blob/181078b6b009d905809cf2e007ad309386a3d1e0/lean/ErdosProblems/Erdos249/PaperCompleteR20/LcmGridCorrespondence.lean\#L57}{\nolinkurl{short_lcm_window_nondivisor}}, \href{https://github.com/wcook04/plectis-erdos/blob/181078b6b009d905809cf2e007ad309386a3d1e0/lean/ErdosProblems/Erdos249/PaperCompleteR20/LcmGridCorrespondence.lean\#L67}{\nolinkurl{clean_lcm_ray_factorisation}}, \href{https://github.com/wcook04/plectis-erdos/blob/181078b6b009d905809cf2e007ad309386a3d1e0/lean/ErdosProblems/Erdos249/PaperCompleteR20/LcmGridCorrespondence.lean\#L76}{\nolinkurl{unclean_lcm_ray_counterexample}}, \href{https://github.com/wcook04/plectis-erdos/blob/181078b6b009d905809cf2e007ad309386a3d1e0/lean/Erdos249257/CarrySurvivorExtinction.lean\#L525}{\nolinkurl{dvd_periodLcm}}.  \emph{Compared}, run \href{https://github.com/wcook04/plectis-erdos-lean/actions/runs/35674034595}{\texttt{35674034595}}.
\par\noindent\hangindent=1.5em Theorem (Unbounded prime support in Mersenne factors): \href{https://github.com/wcook04/plectis-erdos/blob/181078b6b009d905809cf2e007ad309386a3d1e0/lean/ErdosProblems/Erdos249/PaperCompleteR21/MersennePrimeSupportAnchors.lean\#L33}{\nolinkurl{mersenneLayer_prime_divisor_order}}, \href{https://github.com/wcook04/plectis-erdos/blob/181078b6b009d905809cf2e007ad309386a3d1e0/lean/ErdosProblems/Erdos249/PaperCompleteR21/MersennePrimeSupportAnchors.lean\#L63}{\nolinkurl{mersenneLayer_unbounded_prime_support}}.  \emph{Compared}, run \href{https://github.com/wcook04/plectis-erdos-lean/actions/runs/35643815458}{\texttt{35643815458}}.
\par\noindent\hangindent=1.5em Theorem (A prime satisfying the stated cyclotomic conditions): \href{https://github.com/wcook04/plectis-erdos/blob/181078b6b009d905809cf2e007ad309386a3d1e0/lean/ErdosProblems/Erdos249/PaperCompleteR21/MersennePrimeSupportAnchors.lean\#L73}{\nolinkurl{exists_clean_cyclotomic_anchor_paper}}, \href{https://github.com/wcook04/plectis-erdos/blob/181078b6b009d905809cf2e007ad309386a3d1e0/lean/ErdosProblems/Erdos249/PaperCompleteR21/MersennePrimeSupportAnchors.lean\#L84}{\nolinkurl{cyclotomic_layer_prime_order_decomposition_paper}}.  \emph{Compared}, run \href{https://github.com/wcook04/plectis-erdos-lean/actions/runs/35643815458}{\texttt{35643815458}}.
\par\noindent\hangindent=1.5em Theorem (A sufficient order hypothesis for unbounded prime support): \href{https://github.com/wcook04/plectis-erdos/blob/181078b6b009d905809cf2e007ad309386a3d1e0/lean/ErdosProblems/Erdos249/PaperCompleteR21/MersennePrimeSupportAnchors.lean\#L94}{\nolinkurl{boundedDegreeOrderConsumer_unfolded}}, \href{https://github.com/wcook04/plectis-erdos/blob/181078b6b009d905809cf2e007ad309386a3d1e0/lean/ErdosProblems/Erdos249/PaperCompleteR21/MersennePrimeSupportAnchors.lean\#L101}{\nolinkurl{coprime_of_dvd_pow_sub_one}}, \href{https://github.com/wcook04/plectis-erdos/blob/181078b6b009d905809cf2e007ad309386a3d1e0/lean/ErdosProblems/Erdos249/PaperCompleteR21/MersennePrimeSupportAnchors.lean\#L113}{\nolinkurl{dvd_pow_sub_one_iff_orderOf_dvd}}, \href{https://github.com/wcook04/plectis-erdos/blob/181078b6b009d905809cf2e007ad309386a3d1e0/lean/ErdosProblems/Erdos249/PaperCompleteR21/MersennePrimeSupportAnchors.lean\#L136}{\nolinkurl{boundedOrder_witness_iff_orderOf_le}}, \href{https://github.com/wcook04/plectis-erdos/blob/181078b6b009d905809cf2e007ad309386a3d1e0/lean/ErdosProblems/Erdos249/PaperCompleteR21/MersennePrimeSupportAnchors.lean\#L160}{\nolinkurl{orderConsumer_index_lt_pow}}, \href{https://github.com/wcook04/plectis-erdos/blob/181078b6b009d905809cf2e007ad309386a3d1e0/lean/ErdosProblems/Erdos249/PaperCompleteR21/MersennePrimeSupportAnchors.lean\#L168}{\nolinkurl{orderConsumer_finite_prime_escape}}, \href{https://github.com/wcook04/plectis-erdos/blob/181078b6b009d905809cf2e007ad309386a3d1e0/lean/ErdosProblems/Erdos249/PaperCompleteR21/MersennePrimeSupportAnchors.lean\#L178}{\nolinkurl{orderConsumer_unbounded_prime_divisors}}, \href{https://github.com/wcook04/plectis-erdos/blob/181078b6b009d905809cf2e007ad309386a3d1e0/lean/ErdosProblems/Erdos249/PaperCompleteR21/MersennePrimeSupportAnchors.lean\#L190}{\nolinkurl{unbounded_prime_divisors_of_escape_of_nontrivial}}, \href{https://github.com/wcook04/plectis-erdos/blob/181078b6b009d905809cf2e007ad309386a3d1e0/lean/ErdosProblems/Erdos249/PaperCompleteR21/MersennePrimeSupportAnchors.lean\#L212}{\nolinkurl{mersenneLayer_orderConsumer_instance}}, \href{https://github.com/wcook04/plectis-erdos/blob/181078b6b009d905809cf2e007ad309386a3d1e0/lean/ErdosProblems/Erdos249/PaperCompleteR21/MersennePrimeSupportAnchors.lean\#L232}{\nolinkurl{eventual_orderConsumer_conclusions}}, \href{https://github.com/wcook04/plectis-erdos/blob/181078b6b009d905809cf2e007ad309386a3d1e0/lean/ErdosProblems/Erdos249/PaperCompleteR21/MersennePrimeSupportAnchors.lean\#L241}{\nolinkurl{binaryCyclotomicLayer_eventual_instance}}, \href{https://github.com/wcook04/plectis-erdos/blob/181078b6b009d905809cf2e007ad309386a3d1e0/lean/ErdosProblems/Erdos249/PaperCompleteR21/MersennePrimeSupportAnchors.lean\#L247}{\nolinkurl{binaryCyclotomic_allPrime_form_fails}}.  \emph{Compared}, run \href{https://github.com/wcook04/plectis-erdos-lean/actions/runs/35643815458}{\texttt{35643815458}}.
\par\noindent\hangindent=1.5em Theorem (A uniform positive gap for rank-one quotients): \href{https://github.com/wcook04/plectis-erdos/blob/181078b6b009d905809cf2e007ad309386a3d1e0/lean/ErdosProblems/Erdos249/RankOneSubrankObstruction.lean\#L238}{\nolinkurl{rankOneSubrankQuotient_sub_theta_two_gt}}, \href{https://github.com/wcook04/plectis-erdos/blob/181078b6b009d905809cf2e007ad309386a3d1e0/lean/ErdosProblems/Erdos249/RankOneSubrankObstruction.lean\#L163}{\nolinkurl{mobiusMersenneTheta_ge_alpha}}, \href{https://github.com/wcook04/plectis-erdos/blob/181078b6b009d905809cf2e007ad309386a3d1e0/lean/ErdosProblems/Erdos249/RankOneSubrankObstruction.lean\#L196}{\nolinkurl{mobiusMersenneTheta_lt_one}}, \href{https://github.com/wcook04/plectis-erdos/blob/181078b6b009d905809cf2e007ad309386a3d1e0/lean/ErdosProblems/Erdos249/RankOneSubrankObstruction.lean\#L73}{\nolinkurl{abs_mobiusMersenneTheta_sub_prefix_le}}, \href{https://github.com/wcook04/plectis-erdos/blob/181078b6b009d905809cf2e007ad309386a3d1e0/lean/Erdos249257/SignedQMomentObstruction.lean\#L214}{\nolinkurl{mobiusMersenneTheta_two_eq_totient_offset}}, \href{https://github.com/wcook04/plectis-erdos/blob/181078b6b009d905809cf2e007ad309386a3d1e0/lean/ErdosProblems/Erdos249/PaperCompleteR7/ShortNoteAssemblies.lean\#L58}{\nolinkurl{rankOne_denominator_pos}}.  \emph{Compared}, runs \href{https://github.com/wcook04/plectis-erdos-lean/actions/runs/35624228171}{\texttt{35624228171}}, \href{https://github.com/wcook04/plectis-erdos-lean/actions/runs/35643815458}{\texttt{35643815458}}, \href{https://github.com/wcook04/plectis-erdos-lean/actions/runs/35674034595}{\texttt{35674034595}}.
\par\noindent\hangindent=1.5em Proposition (The bound is preserved by positive normalised averaging) (the Lean statement is stronger): \href{https://github.com/wcook04/plectis-erdos/blob/181078b6b009d905809cf2e007ad309386a3d1e0/lean/ErdosProblems/Erdos249/RankOneSubrankObstruction.lean\#L300}{\nolinkurl{positive_direct_sum_sub_theta_two_gt}}, \href{https://github.com/wcook04/plectis-erdos/blob/181078b6b009d905809cf2e007ad309386a3d1e0/lean/ErdosProblems/Erdos249/RankOneSubrankObstruction.lean\#L341}{\nolinkurl{primitive_form_abs_gt}}.  \emph{Compared}, run \href{https://github.com/wcook04/plectis-erdos-lean/actions/runs/35643815458}{\texttt{35643815458}}.
\par\noindent\hangindent=1.5em Theorem (Concatenation and specified period multiples): \href{https://github.com/wcook04/plectis-erdos/blob/181078b6b009d905809cf2e007ad309386a3d1e0/lean/ErdosProblems/Erdos249/PaperCompleteR21/TotientBlockConcatenation.lean\#L23}{\nolinkurl{totientBlock_eq_paper_indexed_sum}}, \href{https://github.com/wcook04/plectis-erdos/blob/181078b6b009d905809cf2e007ad309386a3d1e0/lean/ErdosProblems/Erdos249/PaperCompleteR21/TotientBlockConcatenation.lean\#L41}{\nolinkurl{totientBlock_concatenation}}, \href{https://github.com/wcook04/plectis-erdos/blob/181078b6b009d905809cf2e007ad309386a3d1e0/lean/ErdosProblems/Erdos249/PaperCompleteR21/TotientBlockConcatenation.lean\#L46}{\nolinkurl{totientBlock_doubling}}, \href{https://github.com/wcook04/plectis-erdos/blob/181078b6b009d905809cf2e007ad309386a3d1e0/lean/ErdosProblems/Erdos249/PaperCompleteR21/TotientBlockConcatenation.lean\#L53}{\nolinkurl{cyclotomic_four_eval}}, \href{https://github.com/wcook04/plectis-erdos/blob/181078b6b009d905809cf2e007ad309386a3d1e0/lean/ErdosProblems/Erdos249/PaperCompleteR21/TotientBlockConcatenation.lean\#L64}{\nolinkurl{cyclotomic_four_two_pow_eq_cyclotomic_two}}, \href{https://github.com/wcook04/plectis-erdos/blob/181078b6b009d905809cf2e007ad309386a3d1e0/lean/ErdosProblems/Erdos249/PaperCompleteR21/TotientBlockConcatenation.lean\#L73}{\nolinkurl{cyclotomic_three_two_pow}}, \href{https://github.com/wcook04/plectis-erdos/blob/181078b6b009d905809cf2e007ad309386a3d1e0/lean/ErdosProblems/Erdos249/PaperCompleteR21/TotientBlockConcatenation.lean\#L83}{\nolinkurl{orderOf_two_mod_seven}}, \href{https://github.com/wcook04/plectis-erdos/blob/181078b6b009d905809cf2e007ad309386a3d1e0/lean/ErdosProblems/Erdos249/PaperCompleteR21/TotientBlockConcatenation.lean\#L92}{\nolinkurl{three_not_dvd_two_pow}}, \href{https://github.com/wcook04/plectis-erdos/blob/181078b6b009d905809cf2e007ad309386a3d1e0/lean/ErdosProblems/Erdos249/PaperCompleteR21/TotientBlockConcatenation.lean\#L100}{\nolinkurl{cyclotomic_three_eval_two_not_dvd_doubling_chain}}.  \emph{Compared}, runs \href{https://github.com/wcook04/plectis-erdos-lean/actions/runs/35643815458}{\texttt{35643815458}}, \href{https://github.com/wcook04/plectis-erdos-lean/actions/runs/35674034595}{\texttt{35674034595}}.
\par\noindent\hangindent=1.5em Theorem~\ref{thm:hgap-real}: \href{https://github.com/wcook04/plectis-erdos/blob/181078b6b009d905809cf2e007ad309386a3d1e0/lean/ErdosProblems/Erdos249/PaperCompleteR21/FirstHarmonicBlockCriteria.lean\#L23}{\nolinkurl{windowFirstCos_unfolded}}, \href{https://github.com/wcook04/plectis-erdos/blob/181078b6b009d905809cf2e007ad309386a3d1e0/lean/ErdosProblems/Erdos249/PaperCompleteR21/FirstHarmonicBlockCriteria.lean\#L34}{\nolinkurl{exists_certifiedKill_of_block_real_part_bound}}.  \emph{Compared}, run \href{https://github.com/wcook04/plectis-erdos-lean/actions/runs/35674034595}{\texttt{35674034595}}.
\par\noindent\hangindent=1.5em Theorem~\ref{thm:hgap-subset}: \href{https://github.com/wcook04/plectis-erdos/blob/181078b6b009d905809cf2e007ad309386a3d1e0/lean/ErdosProblems/Erdos249/PaperCompleteR21/FirstHarmonicBlockCriteria.lean\#L50}{\nolinkurl{exists_certifiedKill_of_subset_real_part_bound}}, \href{https://github.com/wcook04/plectis-erdos/blob/181078b6b009d905809cf2e007ad309386a3d1e0/lean/ErdosProblems/Erdos249/PaperCompleteR21/FirstHarmonicBlockCriteria.lean\#L65}{\nolinkurl{block_real_part_bound_of_subset_form}}.  \emph{Compared}, run \href{https://github.com/wcook04/plectis-erdos-lean/actions/runs/35674034595}{\texttt{35674034595}}.
\par\noindent\hangindent=1.5em Theorem~\ref{thm:hgap-norm}: \href{https://github.com/wcook04/plectis-erdos/blob/181078b6b009d905809cf2e007ad309386a3d1e0/lean/ErdosProblems/Erdos249/PaperCompleteR21/FirstHarmonicBlockCriteria.lean\#L86}{\nolinkurl{blockNormCondition_unfolded}}, \href{https://github.com/wcook04/plectis-erdos/blob/181078b6b009d905809cf2e007ad309386a3d1e0/lean/ErdosProblems/Erdos249/PaperCompleteR21/FirstHarmonicBlockCriteria.lean\#L97}{\nolinkurl{real_part_bound_of_norm_bound}}, \href{https://github.com/wcook04/plectis-erdos/blob/181078b6b009d905809cf2e007ad309386a3d1e0/lean/ErdosProblems/Erdos249/PaperCompleteR21/FirstHarmonicBlockCriteria.lean\#L110}{\nolinkurl{exists_certifiedKill_of_block_norm_bound}}, \href{https://github.com/wcook04/plectis-erdos/blob/181078b6b009d905809cf2e007ad309386a3d1e0/lean/ErdosProblems/Erdos249/PaperCompleteR21/FirstHarmonicBlockCriteria.lean\#L118}{\nolinkurl{irrational_of_blockNormCondition}}.  \emph{Compared}, run \href{https://github.com/wcook04/plectis-erdos-lean/actions/runs/35674034595}{\texttt{35674034595}}.
\par\noindent\hangindent=1.5em Theorem (The proved implications): \href{https://github.com/wcook04/plectis-erdos/blob/181078b6b009d905809cf2e007ad309386a3d1e0/lean/ErdosProblems/Erdos249/PaperCompleteR21/TopEdgeCorridorAndSeparation.lean\#L64}{\nolinkurl{topEdgeResidueGap_unfolded}}, \href{https://github.com/wcook04/plectis-erdos/blob/181078b6b009d905809cf2e007ad309386a3d1e0/lean/ErdosProblems/Erdos249/PaperCompleteR21/TopEdgeCorridorAndSeparation.lean\#L75}{\nolinkurl{topEdgeResidueGap_forces_nonintegral}}, \href{https://github.com/wcook04/plectis-erdos/blob/181078b6b009d905809cf2e007ad309386a3d1e0/lean/ErdosProblems/Erdos249/PaperCompleteR21/TopEdgeCorridorAndSeparation.lean\#L83}{\nolinkurl{topEdgeResidueGap_orbit_nonintegral}}, \href{https://github.com/wcook04/plectis-erdos/blob/181078b6b009d905809cf2e007ad309386a3d1e0/lean/ErdosProblems/Erdos249/PaperCompleteR21/TopEdgeCorridorAndSeparation.lean\#L92}{\nolinkurl{irrational_of_topEdgeResidueGapSupply}}.  \emph{Compared}, run \href{https://github.com/wcook04/plectis-erdos-lean/actions/runs/35674034595}{\texttt{35674034595}}.
\par\noindent\hangindent=1.5em Proposition~\ref{prop:te-chain}: \href{https://github.com/wcook04/plectis-erdos/blob/181078b6b009d905809cf2e007ad309386a3d1e0/lean/ErdosProblems/Erdos249/PaperCompleteR21/TopEdgeChainPaperBand.lean\#L56}{\nolinkurl{paperTeChain_item_one_unfolded}}, \href{https://github.com/wcook04/plectis-erdos/blob/181078b6b009d905809cf2e007ad309386a3d1e0/lean/ErdosProblems/Erdos249/PaperCompleteR21/TopEdgeCorridorAndSeparation.lean\#L113}{\nolinkurl{teChain_item_two_unfolded}}, \href{https://github.com/wcook04/plectis-erdos/blob/181078b6b009d905809cf2e007ad309386a3d1e0/lean/ErdosProblems/Erdos249/PaperCompleteR21/TopEdgeCorridorAndSeparation.lean\#L124}{\nolinkurl{teChain_item_three_unfolded}}, \href{https://github.com/wcook04/plectis-erdos/blob/181078b6b009d905809cf2e007ad309386a3d1e0/lean/ErdosProblems/Erdos249/PaperCompleteR21/TopEdgeCorridorAndSeparation.lean\#L133}{\nolinkurl{teChain_item_four_unfolded}}, \href{https://github.com/wcook04/plectis-erdos/blob/181078b6b009d905809cf2e007ad309386a3d1e0/lean/ErdosProblems/Erdos249/PaperCompleteR21/TopEdgeCorridorAndSeparation.lean\#L142}{\nolinkurl{teChain_item_five_unfolded}}, \href{https://github.com/wcook04/plectis-erdos/blob/181078b6b009d905809cf2e007ad309386a3d1e0/lean/ErdosProblems/Erdos249/PaperCompleteR21/TopEdgeChainPaperBand.lean\#L214}{\nolinkurl{paperTeChain_five_sufficient_for_irrationality}}, \href{https://github.com/wcook04/plectis-erdos/blob/181078b6b009d905809cf2e007ad309386a3d1e0/lean/ErdosProblems/Erdos249/PaperCompleteR21/TopEdgeChainPaperBand.lean\#L233}{\nolinkurl{paperTeChain_first_four_imply_topEdgeSupply}}, \href{https://github.com/wcook04/plectis-erdos/blob/181078b6b009d905809cf2e007ad309386a3d1e0/lean/ErdosProblems/Erdos249/PaperCompleteR21/TopEdgeChainPaperBand.lean\#L271}{\nolinkurl{paperTeChain_fifth_gives_nonintegrality}}, \href{https://github.com/wcook04/plectis-erdos/blob/181078b6b009d905809cf2e007ad309386a3d1e0/lean/ErdosProblems/Erdos249/PaperCompleteR21/TopEdgeChainPaperBand.lean\#L254}{\nolinkurl{paperTeChain_relations}}, \href{https://github.com/wcook04/plectis-erdos/blob/181078b6b009d905809cf2e007ad309386a3d1e0/lean/ErdosProblems/Erdos249/PaperCompleteR21/TopEdgeChainPaperBand.lean\#L87}{\nolinkurl{topEdgeResidueGap_or_of_paperAdjacentSuffixMidband}}, \href{https://github.com/wcook04/plectis-erdos/blob/181078b6b009d905809cf2e007ad309386a3d1e0/lean/ErdosProblems/Erdos249/PaperCompleteR21/TopEdgeChainPaperBand.lean\#L179}{\nolinkurl{powerTwoActualLcmTopEdgeResidueGapSupply_of_paperAdjacentSuffixMidband}}, \href{https://github.com/wcook04/plectis-erdos/blob/181078b6b009d905809cf2e007ad309386a3d1e0/lean/ErdosProblems/Erdos249/PaperCompleteR21/TopEdgeChainPaperBand.lean\#L190}{\nolinkurl{irrational_of_paperAdjacentSuffixMidbandSupply}}, \href{https://github.com/wcook04/plectis-erdos/blob/181078b6b009d905809cf2e007ad309386a3d1e0/lean/ErdosProblems/Erdos249/PaperCompleteR21/TopEdgeChainPaperBand.lean\#L197}{\nolinkurl{paperAdjacentSuffixMidbandSupply_of_oddGuardTopEdgeHalfWordBand}}, \href{https://github.com/wcook04/plectis-erdos/blob/181078b6b009d905809cf2e007ad309386a3d1e0/lean/ErdosProblems/Erdos249/PaperCompleteR21/TopEdgeChainPaperBand.lean\#L204}{\nolinkurl{paperAdjacentSuffixMidbandSupply_of_flexibleActualTopEdgeMagnitude}}, \href{https://github.com/wcook04/plectis-erdos/blob/181078b6b009d905809cf2e007ad309386a3d1e0/lean/ErdosProblems/Erdos249/PaperCompleteR21/TopEdgeChainPaperBand.lean\#L69}{\nolinkurl{paperAdjacentSuffixMidbandSupply_of_adjacentSuffixMidband}}.  \emph{Compared}, run \href{https://github.com/wcook04/plectis-erdos-lean/actions/runs/35674034595}{\texttt{35674034595}}.
\par\noindent\hangindent=1.5em Theorem (The proved implication): \href{https://github.com/wcook04/plectis-erdos/blob/181078b6b009d905809cf2e007ad309386a3d1e0/lean/ErdosProblems/Erdos249/PaperCompleteR21/TopEdgeCorridorAndSeparation.lean\#L217}{\nolinkurl{terminalDominance_orbit_nonintegral}}, \href{https://github.com/wcook04/plectis-erdos/blob/181078b6b009d905809cf2e007ad309386a3d1e0/lean/ErdosProblems/Erdos249/PaperCompleteR21/TopEdgeCorridorAndSeparation.lean\#L228}{\nolinkurl{irrational_of_terminalDominanceSupply}}.  \emph{Compared}, run \href{https://github.com/wcook04/plectis-erdos-lean/actions/runs/35674034595}{\texttt{35674034595}}.
\par\noindent\hangindent=1.5em Proposition (A sufficient extension): \href{https://github.com/wcook04/plectis-erdos/blob/181078b6b009d905809cf2e007ad309386a3d1e0/lean/ErdosProblems/Erdos249/PaperCompleteR21/TopEdgeCorridorAndSeparation.lean\#L236}{\nolinkurl{irrational_of_lower_escape_supply}}, \href{https://github.com/wcook04/plectis-erdos/blob/181078b6b009d905809cf2e007ad309386a3d1e0/lean/ErdosProblems/Erdos249/PaperCompleteR21/TopEdgeCorridorAndSeparation.lean\#L252}{\nolinkurl{corridor_escape_and_irrational_of_magnitude}}, \href{https://github.com/wcook04/plectis-erdos/blob/181078b6b009d905809cf2e007ad309386a3d1e0/lean/ErdosProblems/Erdos249/PaperCompleteR21/TopEdgeCorridorAndSeparation.lean\#L133}{\nolinkurl{teChain_item_four_unfolded}}.  \emph{Compared}, run \href{https://github.com/wcook04/plectis-erdos-lean/actions/runs/35674034595}{\texttt{35674034595}}.
\par\noindent\hangindent=1.5em Proposition (A sufficient extension): \href{https://github.com/wcook04/plectis-erdos/blob/181078b6b009d905809cf2e007ad309386a3d1e0/lean/ErdosProblems/Erdos249/PaperCompleteR21/ShortWindowSupplyAndSixteenShifts.lean\#L58}{\nolinkurl{irrational_of_logarithmicDepth_diagonal_supply}}, \href{https://github.com/wcook04/plectis-erdos/blob/181078b6b009d905809cf2e007ad309386a3d1e0/lean/ErdosProblems/Erdos249/PaperCompleteR21/ShortWindowSupplyAndSixteenShifts.lean\#L72}{\nolinkurl{irrational_of_restrictedDepth_diagonal_supply}}.  \emph{Compared}, run \href{https://github.com/wcook04/plectis-erdos-lean/actions/runs/35674034595}{\texttt{35674034595}}.
\par\noindent\hangindent=1.5em Proposition (Simultaneous certificates with unrestricted depth): \href{https://github.com/wcook04/plectis-erdos/blob/181078b6b009d905809cf2e007ad309386a3d1e0/lean/ErdosProblems/Erdos249/PaperCompleteR21/SimultaneousShiftCertificateDepth.lean\#L131}{\nolinkurl{exists_growingShift_simultaneous_certificate_iff_irrational}}, \href{https://github.com/wcook04/plectis-erdos/blob/181078b6b009d905809cf2e007ad309386a3d1e0/lean/ErdosProblems/Erdos249/PaperCompleteR21/SimultaneousShiftCertificateDepth.lean\#L122}{\nolinkurl{exists_simultaneous_depth_succ_of_irrational}}, \href{https://github.com/wcook04/plectis-erdos/blob/181078b6b009d905809cf2e007ad309386a3d1e0/lean/ErdosProblems/Erdos249/PaperCompleteR21/SimultaneousShiftCertificateDepth.lean\#L89}{\nolinkurl{exists_simultaneous_depth_of_irrational}}.  \emph{Compared}, run \href{https://github.com/wcook04/plectis-erdos-lean/actions/runs/35674034595}{\texttt{35674034595}}.
\par\noindent\hangindent=1.5em Proposition (A sufficient extension): \href{https://github.com/wcook04/plectis-erdos/blob/181078b6b009d905809cf2e007ad309386a3d1e0/lean/ErdosProblems/Erdos249/PaperCompleteR21/FareyDenominatorFloorExtension.lean\#L68}{\nolinkurl{irrational_of_unbounded_window_one_gapCheck}}, \href{https://github.com/wcook04/plectis-erdos/blob/181078b6b009d905809cf2e007ad309386a3d1e0/lean/ErdosProblems/Erdos249/PaperCompleteR21/FareyDenominatorFloorExtension.lean\#L58}{\nolinkurl{gapCheck_window_one_excludes}}.  \emph{Compared}, run \href{https://github.com/wcook04/plectis-erdos-lean/actions/runs/35643815458}{\texttt{35643815458}}.
\par\noindent\hangindent=1.5em Proposition (A sufficient accumulated-residue condition): \href{https://github.com/wcook04/plectis-erdos/blob/181078b6b009d905809cf2e007ad309386a3d1e0/lean/ErdosProblems/Erdos249/PaperCompleteR21/TwoAdicHalfPulseAndAccumulatedResidue.lean\#L77}{\nolinkurl{irrational_of_accumulated_halfModulus_supply}}, \href{https://github.com/wcook04/plectis-erdos/blob/181078b6b009d905809cf2e007ad309386a3d1e0/lean/ErdosProblems/Erdos249/PaperCompleteR21/TwoAdicHalfPulseAndAccumulatedResidue.lean\#L53}{\nolinkurl{certifiedKill_of_halfModulus_residue}}.  \emph{Compared}, run \href{https://github.com/wcook04/plectis-erdos-lean/actions/runs/35674034595}{\texttt{35674034595}}.
\par\noindent\hangindent=1.5em Theorem (Arbitrarily large prime-power LCM jumps): \href{https://github.com/wcook04/plectis-erdos/blob/181078b6b009d905809cf2e007ad309386a3d1e0/lean/ErdosProblems/Erdos249/PaperCompleteR21/LcmJumpPositionsAndCentralSlack.lean\#L26}{\nolinkurl{exists_periodLcm_strict_jump_ge_paper}}, \href{https://github.com/wcook04/plectis-erdos/blob/181078b6b009d905809cf2e007ad309386a3d1e0/lean/ErdosProblems/Erdos249/PaperCompleteR21/LcmJumpPositionsAndCentralSlack.lean\#L31}{\nolinkurl{periodLcm_strict_jump_at_prime_pred}}, \href{https://github.com/wcook04/plectis-erdos/blob/181078b6b009d905809cf2e007ad309386a3d1e0/lean/ErdosProblems/Erdos249/PaperCompleteR21/LcmJumpPositionsAndCentralSlack.lean\#L46}{\nolinkurl{periodLcm_strict_jump_at_powerTwo_pred}}, \href{https://github.com/wcook04/plectis-erdos/blob/181078b6b009d905809cf2e007ad309386a3d1e0/lean/ErdosProblems/Erdos249/PaperCompleteR21/LcmJumpPositionsAndCentralSlack.lean\#L59}{\nolinkurl{periodLcm_zero_and_one_eq_one}}.  \emph{Compared}, run \href{https://github.com/wcook04/plectis-erdos-lean/actions/runs/35674034595}{\texttt{35674034595}}.
\par\noindent\hangindent=1.5em Proposition (A sufficient inequality): \href{https://github.com/wcook04/plectis-erdos/blob/181078b6b009d905809cf2e007ad309386a3d1e0/lean/ErdosProblems/Erdos249/PaperCompleteR21/LcmJumpPositionsAndCentralSlack.lean\#L90}{\nolinkurl{irrational_of_powerTwo_postJump_slack_supply}}, \href{https://github.com/wcook04/plectis-erdos/blob/181078b6b009d905809cf2e007ad309386a3d1e0/lean/ErdosProblems/Erdos249/PaperCompleteR21/LcmJumpPositionsAndCentralSlack.lean\#L68}{\nolinkurl{canonicalAdjacentSuffixCentralSlack_paper_formula}}.  \emph{Compared}, run \href{https://github.com/wcook04/plectis-erdos-lean/actions/runs/35674034595}{\texttt{35674034595}}.
\par\noindent\hangindent=1.5em Theorem (One explicit witness): \href{https://github.com/wcook04/plectis-erdos/blob/181078b6b009d905809cf2e007ad309386a3d1e0/lean/ErdosProblems/Erdos249/PaperCompleteR21/PrimeJumpWitnessAndMersenneChannels.lean\#L62}{\nolinkurl{periodLcm_four_eq_twelve}}, \href{https://github.com/wcook04/plectis-erdos/blob/181078b6b009d905809cf2e007ad309386a3d1e0/lean/ErdosProblems/Erdos249/PaperCompleteR21/PrimeJumpWitnessAndMersenneChannels.lean\#L68}{\nolinkurl{primeJump_witness_twelve_five_values}}, \href{https://github.com/wcook04/plectis-erdos/blob/181078b6b009d905809cf2e007ad309386a3d1e0/lean/ErdosProblems/Erdos249/PaperCompleteR21/PrimeJumpWitnessAndMersenneChannels.lean\#L76}{\nolinkurl{primeJumpTailCommutator_twelve_five_notMem_int}}, \href{https://github.com/wcook04/plectis-erdos/blob/181078b6b009d905809cf2e007ad309386a3d1e0/lean/ErdosProblems/Erdos249/PaperCompleteR21/PrimeJumpWitnessAndMersenneChannels.lean\#L81}{\nolinkurl{irrational_of_primeJumpSharp_supply}}.  \emph{Compared}, run \href{https://github.com/wcook04/plectis-erdos-lean/actions/runs/35674034595}{\texttt{35674034595}}.
\par\noindent\hangindent=1.5em Proposition (The additional approximation hypothesis): \href{https://github.com/wcook04/plectis-erdos/blob/181078b6b009d905809cf2e007ad309386a3d1e0/lean/ErdosProblems/Erdos249/PaperCompleteR21/PrimeJumpWitnessAndMersenneChannels.lean\#L137}{\nolinkurl{irrational_totientSeries_of_rational_separation}}, \href{https://github.com/wcook04/plectis-erdos/blob/181078b6b009d905809cf2e007ad309386a3d1e0/lean/ErdosProblems/Erdos249/PaperCompleteR21/PrimeJumpWitnessAndMersenneChannels.lean\#L149}{\nolinkurl{den_mul_abs_sub_ge_one_div_den}}.  \emph{Compared}, run \href{https://github.com/wcook04/plectis-erdos-lean/actions/runs/35643815458}{\texttt{35643815458}}.
\par\noindent\hangindent=1.5em Proposition (An additional hypothesis for a totient-specific rank argument) (the Lean statement is stronger): \href{https://github.com/wcook04/plectis-erdos/blob/181078b6b009d905809cf2e007ad309386a3d1e0/lean/ErdosProblems/Erdos249/PaperCompleteR21/DyadicSectionBasisAndRationalCarry.lean\#L238}{\nolinkurl{rationalControl_periodic_with_unbounded_carry_rank}}, \href{https://github.com/wcook04/plectis-erdos/blob/181078b6b009d905809cf2e007ad309386a3d1e0/lean/ErdosProblems/Erdos249/PaperCompleteR21/DyadicSectionBasisAndRationalCarry.lean\#L256}{\nolinkurl{rationality_gives_mod_period_and_unbounded_rank}}.  \emph{Compared}, run \href{https://github.com/wcook04/plectis-erdos-lean/actions/runs/35624228171}{\texttt{35624228171}}.
\par\noindent\hangindent=1.5em Proposition (A rational series preserving totient parity and the stated separation properties): \href{https://github.com/wcook04/plectis-erdos/blob/181078b6b009d905809cf2e007ad309386a3d1e0/lean/Erdos249257/TotientParityCoboundaryCountermodel.lean\#L227}{\nolinkurl{parityCoboundaryWeight_le_six}}, \href{https://github.com/wcook04/plectis-erdos/blob/181078b6b009d905809cf2e007ad309386a3d1e0/lean/Erdos249257/TotientParityCoboundaryCountermodel.lean\#L155}{\nolinkurl{parityCoboundaryWeight_le_self}}, \href{https://github.com/wcook04/plectis-erdos/blob/181078b6b009d905809cf2e007ad309386a3d1e0/lean/Erdos249257/TotientParityCoboundaryCountermodel.lean\#L194}{\nolinkurl{parityCoboundaryWeight_mod_two_eq_totient}}, \href{https://github.com/wcook04/plectis-erdos/blob/181078b6b009d905809cf2e007ad309386a3d1e0/lean/Erdos249257/TotientParityCoboundaryCountermodel.lean\#L457}{\nolinkurl{exists_later_arbitrarily_many_separated_parityCoboundaryWeight_carry_pairs}}, \href{https://github.com/wcook04/plectis-erdos/blob/181078b6b009d905809cf2e007ad309386a3d1e0/lean/Erdos249257/TotientParityCoboundaryCountermodel.lean\#L359}{\nolinkurl{tsum_parityCoboundaryWeight_eq_three_halves}}.  \emph{Compared}, run \href{https://github.com/wcook04/plectis-erdos-lean/actions/runs/35624228171}{\texttt{35624228171}}.
\par\noindent\hangindent=1.5em Lemma (Complement divisibility after multiplication): \href{https://github.com/wcook04/plectis-erdos/blob/181078b6b009d905809cf2e007ad309386a3d1e0/lean/ErdosProblems/Erdos249/PaperCompleteR21/ScalarLocalisationAndInversePhaseGauge.lean\#L29}{\nolinkurl{complementDenominator_dvd_scalar}}.  \emph{Compared}, run \href{https://github.com/wcook04/plectis-erdos-lean/actions/runs/35643815458}{\texttt{35643815458}}.
\par\noindent\hangindent=1.5em Lemma (Nonvanishing from a unique largest denominator exponent): \href{https://github.com/wcook04/plectis-erdos/blob/181078b6b009d905809cf2e007ad309386a3d1e0/lean/ErdosProblems/Erdos249/PaperCompleteR20/SignedDyadicClearing.lean\#L8}{\nolinkurl{signed_dyadic_clearing}}, \href{https://github.com/wcook04/plectis-erdos/blob/181078b6b009d905809cf2e007ad309386a3d1e0/lean/ErdosProblems/Erdos249/PaperCompleteR20/SignedDyadicClearing.lean\#L28}{\nolinkurl{signed_dyadic_sum_ne_zero}}, \href{https://github.com/wcook04/plectis-erdos/blob/181078b6b009d905809cf2e007ad309386a3d1e0/lean/Erdos249257/SignedQMomentObstruction.lean\#L78}{\nolinkurl{scaled_dyadic_sum_odd}}.  \emph{Compared}, runs \href{https://github.com/wcook04/plectis-erdos-lean/actions/runs/35624228171}{\texttt{35624228171}}, \href{https://github.com/wcook04/plectis-erdos-lean/actions/runs/35643815458}{\texttt{35643815458}}.
\par\noindent\hangindent=1.5em Proposition (A nonzero minor survives inverse-phase column weights): \href{https://github.com/wcook04/plectis-erdos/blob/181078b6b009d905809cf2e007ad309386a3d1e0/lean/ErdosProblems/Erdos249/PaperCompleteR21/ScalarLocalisationAndInversePhaseGauge.lean\#L51}{\nolinkurl{inversePhaseGauge_locks_row_and_preserves_minor}}.  \emph{Compared}, run \href{https://github.com/wcook04/plectis-erdos-lean/actions/runs/35624228171}{\texttt{35624228171}}.
\par\noindent\hangindent=1.5em Proposition (Tail integrality on an LCM grid): \href{https://github.com/wcook04/plectis-erdos/blob/181078b6b009d905809cf2e007ad309386a3d1e0/lean/ErdosProblems/Erdos249/PaperCompleteR20/LcmGridCorrespondence.lean\#L8}{\nolinkurl{lcm_grid_flatness}}.  \emph{Compared}, run \href{https://github.com/wcook04/plectis-erdos-lean/actions/runs/35674034595}{\texttt{35674034595}}.
\par\noindent\hangindent=1.5em Proposition (A finite-grid certificate gives a nonintegral pair): \href{https://github.com/wcook04/plectis-erdos/blob/181078b6b009d905809cf2e007ad309386a3d1e0/lean/ErdosProblems/Erdos249/PaperCompleteR20/FiniteGridCorrespondence.lean\#L12}{\nolinkurl{paperGridNumerator_eq}}, \href{https://github.com/wcook04/plectis-erdos/blob/181078b6b009d905809cf2e007ad309386a3d1e0/lean/ErdosProblems/Erdos249/PaperCompleteR20/FiniteGridCorrespondence.lean\#L36}{\nolinkurl{finite_grid_nonintegral_pair}}.  \emph{Compared}, run \href{https://github.com/wcook04/plectis-erdos-lean/actions/runs/35674034595}{\texttt{35674034595}}.
\par\noindent\hangindent=1.5em Theorem (Exact dyadic rank and infinite-dimensionality): \href{https://github.com/wcook04/plectis-erdos/blob/181078b6b009d905809cf2e007ad309386a3d1e0/lean/ErdosProblems/Erdos249/PaperCompleteR21/DyadicSectionBasisAndRationalCarry.lean\#L33}{\nolinkurl{canonicalTotientKernelFamily_entries}}, \href{https://github.com/wcook04/plectis-erdos/blob/181078b6b009d905809cf2e007ad309386a3d1e0/lean/ErdosProblems/Erdos249/PaperCompleteR21/DyadicSectionBasisAndRationalCarry.lean\#L53}{\nolinkurl{canonicalTotientKernelFamily_independent_card_and_span}}.  \emph{Compared}, run \href{https://github.com/wcook04/plectis-erdos-lean/actions/runs/35624228171}{\texttt{35624228171}}.
\par\noindent\hangindent=1.5em Lemma~\ref{lem:orbit} (the Lean statement is stronger): \href{https://github.com/wcook04/plectis-erdos/blob/181078b6b009d905809cf2e007ad309386a3d1e0/lean/ErdosProblems/Erdos249/PaperCompleteR21/DoublingOrbitTransferAndFullDepthPhase.lean\#L26}{\nolinkurl{orbit_tail_recurrence}}, \href{https://github.com/wcook04/plectis-erdos/blob/181078b6b009d905809cf2e007ad309386a3d1e0/lean/ErdosProblems/Erdos249/PaperCompleteR21/DoublingOrbitTransferAndFullDepthPhase.lean\#L32}{\nolinkurl{orbit_tail_diff_eq}}, \href{https://github.com/wcook04/plectis-erdos/blob/181078b6b009d905809cf2e007ad309386a3d1e0/lean/ErdosProblems/Erdos249/PaperCompleteR21/DoublingOrbitTransferAndFullDepthPhase.lean\#L41}{\nolinkurl{orbit_tail_diff_sub_scaled_is_int}}, \href{https://github.com/wcook04/plectis-erdos/blob/181078b6b009d905809cf2e007ad309386a3d1e0/lean/ErdosProblems/Erdos249/PaperCompleteR21/DoublingOrbitTransferAndFullDepthPhase.lean\#L53}{\nolinkurl{orbit_tail_diff_firstChar_eq}}, \href{https://github.com/wcook04/plectis-erdos/blob/181078b6b009d905809cf2e007ad309386a3d1e0/lean/ErdosProblems/Erdos249/PaperCompleteR21/DoublingOrbitTransferAndFullDepthPhase.lean\#L75}{\nolinkurl{orbit_tail_diff_fract_eq_doubling_orbit}}, \href{https://github.com/wcook04/plectis-erdos/blob/181078b6b009d905809cf2e007ad309386a3d1e0/lean/ErdosProblems/Erdos249/PaperCompleteR21/DoublingOrbitTransferAndFullDepthPhase.lean\#L64}{\nolinkurl{doublingMap_iterate_apply}}.  \emph{Compared}, runs \href{https://github.com/wcook04/plectis-erdos-lean/actions/runs/35643815458}{\texttt{35643815458}}, \href{https://github.com/wcook04/plectis-erdos-lean/actions/runs/35674034595}{\texttt{35674034595}}.
\par\noindent\hangindent=1.5em Proposition~\ref{prop:transfer}: \href{https://github.com/wcook04/plectis-erdos/blob/181078b6b009d905809cf2e007ad309386a3d1e0/lean/ErdosProblems/Erdos249/PaperCompleteR21/DoublingOrbitTransferAndFullDepthPhase.lean\#L92}{\nolinkurl{irrational_totientSeries_of_block_cosine_gap}}.  \emph{Compared}, run \href{https://github.com/wcook04/plectis-erdos-lean/actions/runs/35643815458}{\texttt{35643815458}}.
\par\noindent\hangindent=1.5em Corollary~\ref{cor:digitform}: \href{https://github.com/wcook04/plectis-erdos/blob/181078b6b009d905809cf2e007ad309386a3d1e0/lean/ErdosProblems/Erdos249/PaperCompleteR21/BinaryDigitChangeDensity.lean\#L121}{\nolinkurl{irrational_totientSeries_of_quarterFarPhase_proportion}}, \href{https://github.com/wcook04/plectis-erdos/blob/181078b6b009d905809cf2e007ad309386a3d1e0/lean/ErdosProblems/Erdos249/PaperCompleteR21/BinaryDigitChangeDensity.lean\#L146}{\nolinkurl{quarterFarFromInt_iff_binaryDigitAt_change}}, \href{https://github.com/wcook04/plectis-erdos/blob/181078b6b009d905809cf2e007ad309386a3d1e0/lean/ErdosProblems/Erdos249/PaperCompleteR21/BinaryDigitChangeDensity.lean\#L233}{\nolinkurl{irrational_totientSeries_of_digitChange_count}}, \href{https://github.com/wcook04/plectis-erdos/blob/181078b6b009d905809cf2e007ad309386a3d1e0/lean/ErdosProblems/Erdos249/PaperCompleteR21/BinaryDigitChangeDensity.lean\#L95}{\nolinkurl{irrational_totientSeries_of_quarterFarPhase_count}}, \href{https://github.com/wcook04/plectis-erdos/blob/181078b6b009d905809cf2e007ad309386a3d1e0/lean/ErdosProblems/Erdos249/PaperCompleteR21/BinaryDigitChangeDensity.lean\#L45}{\nolinkurl{quarterFarFromInt_iff_floor_bounds}}, \href{https://github.com/wcook04/plectis-erdos/blob/181078b6b009d905809cf2e007ad309386a3d1e0/lean/ErdosProblems/Erdos249/PaperCompleteR21/BinaryDigitChangeDensity.lean\#L70}{\nolinkurl{cos_nonpos_of_quarterFarFromInt}}, \href{https://github.com/wcook04/plectis-erdos/blob/181078b6b009d905809cf2e007ad309386a3d1e0/lean/ErdosProblems/Erdos249/PaperCompleteR21/BinaryDigitChangeDensity.lean\#L81}{\nolinkurl{tailOrbitFirstExp_re_eq}}.  \emph{Compared}, run \href{https://github.com/wcook04/plectis-erdos-lean/actions/runs/35624228171}{\texttt{35624228171}}.
\par\noindent\hangindent=1.5em Theorem~\ref{thm:lacunary} (the Lean statement is stronger): \href{https://github.com/wcook04/plectis-erdos/blob/181078b6b009d905809cf2e007ad309386a3d1e0/lean/ErdosProblems/Erdos249/PaperCompleteR21/LacunaryFactorialBlockNorm.lean\#L68}{\nolinkurl{lacCoef_bounds}}, \href{https://github.com/wcook04/plectis-erdos/blob/181078b6b009d905809cf2e007ad309386a3d1e0/lean/ErdosProblems/Erdos249/PaperCompleteR21/LacunaryFactorialBlockNorm.lean\#L88}{\nolinkurl{lacBeta_eq_factorial_series}}, \href{https://github.com/wcook04/plectis-erdos/blob/181078b6b009d905809cf2e007ad309386a3d1e0/lean/ErdosProblems/Erdos249/PaperCompleteR21/LacunaryFactorialBlockNorm.lean\#L121}{\nolinkurl{irrational_lacBeta}}, \href{https://github.com/wcook04/plectis-erdos/blob/181078b6b009d905809cf2e007ad309386a3d1e0/lean/ErdosProblems/Erdos249/PaperCompleteR21/LacunaryFactorialBlockNorm.lean\#L333}{\nolinkurl{lacunary_block_cos_gap}}, \href{https://github.com/wcook04/plectis-erdos/blob/181078b6b009d905809cf2e007ad309386a3d1e0/lean/ErdosProblems/Erdos249/PaperCompleteR21/LacunaryFactorialBlockNorm.lean\#L587}{\nolinkurl{lacunary_block_norm_fails}}, \href{https://github.com/wcook04/plectis-erdos/blob/181078b6b009d905809cf2e007ad309386a3d1e0/lean/ErdosProblems/Erdos249/PaperCompleteR21/LacunaryFactorialBlockNorm.lean\#L317}{\nolinkurl{cos_pi_div_eight_gt}}.  \emph{Compared}, run \href{https://github.com/wcook04/plectis-erdos-lean/actions/runs/35624228171}{\texttt{35624228171}}.
\par\noindent\hangindent=1.5em Proposition~\ref{prop:route4}: \href{https://github.com/wcook04/plectis-erdos/blob/181078b6b009d905809cf2e007ad309386a3d1e0/lean/ErdosProblems/Erdos249/PaperCompleteR21/DoublingOrbitTransferAndFullDepthPhase.lean\#L130}{\nolinkurl{certifiedKill_of_fullDepth_phase_separation}}, \href{https://github.com/wcook04/plectis-erdos/blob/181078b6b009d905809cf2e007ad309386a3d1e0/lean/ErdosProblems/Erdos249/PaperCompleteR21/DoublingOrbitTransferAndFullDepthPhase.lean\#L119}{\nolinkurl{bracket_of_two_sided_separation}}.  \emph{Compared}, runs \href{https://github.com/wcook04/plectis-erdos-lean/actions/runs/35643815458}{\texttt{35643815458}}, \href{https://github.com/wcook04/plectis-erdos-lean/actions/runs/35674034595}{\texttt{35674034595}}.
\par\endgroup
% END GENERATED CONCORDANCE

\paragraph{Verification scope.} Kernel checking validates a formal proof of
its stated proposition relative to the declared axioms. It does not check
that every prose summary has the same hypotheses or verify the literature
comparisons. In particular, a checked implication with an unproved
hypothesis is not a proof that Problem 249 is solved.

\paragraph*{Acknowledgements.} I thank Wouter van Doorn for advice on
writing for a first-time reader and on explaining the strength of a
hypothesis. His advice concerned a different note; it was not a
mathematical review or endorsement of the results presented here.

\paragraph{Funding and competing interests.} This work received no external
funding. The author declares no competing interests.

\clearpage
\begin{thebibliography}{99}
\small
\raggedright
\setlength{\itemsep}{0pt}
\bibitem{allouche-shallit}
J.-P.~Allouche and J.~Shallit, \emph{The ring of $k$-regular sequences},
Theoret.\ Comput.\ Sci.\ \textbf{98} (1992), no.~2, 163--197,
doi:\href{https://doi.org/10.1016/0304-3975(92)90001-V}{10.1016/0304-3975(92)90001-V}.
Definitions and early statements are cited with the author preprint's numbering.
\bibitem{coons}
M.~Coons, \emph{(Non)Automaticity of number theoretic functions},
J.\ Th\'eor.\ Nombres Bordeaux \textbf{22} (2010), no.~2, 339--352,
doi:\href{https://doi.org/10.5802/jtnb.718}{10.5802/jtnb.718}.
Theorem~3.2, p.~348.
\bibitem{coons-jsr}
M.~Coons, \emph{Regular sequences and the joint spectral radius},
Internat.\ J.\ Found.\ Comput.\ Sci.\ \textbf{28} (2017), no.~2, 135--140,
doi:\href{https://doi.org/10.1142/S0129054117500095}{10.1142/S0129054117500095};
arXiv:\href{https://arxiv.org/abs/1511.07535v1}{1511.07535v1}.
Theorem~1, Proposition~4 and Corollary~7 are cited with the arXiv v1 numbering.
\bibitem{martin-phi-inequalities}
G.~Martin, \emph{Simultaneous inequalities among values of the Euler
phi-function}, arXiv:\href{https://arxiv.org/abs/math/0603053v1}{math/0603053v1}, 2006.
Theorem~1, pp.~1--2.
\bibitem{wong2015}
E.~Wong, \href{https://math.stackexchange.com/a/1211557}{answer 1211557} to
\emph{An infinite sum based on the mod-parity of Euler's totient function},
Mathematics Stack Exchange, 29 March 2015, accessed 15 September 2026.
\bibitem{eldar2026oeis}
A.~Eldar, comment and formula added to OEIS A256936 (revisions 28 and 31),
15 March 2026, \url{https://oeis.org/history?seq=A256936},
accessed 16 September 2026.
\bibitem{fan2026totient}
S.~Fan, comment on Erd\H{o}s Problem \#249, 16 May 2026, 19:01,
\href{https://www.erdosproblems.com/forum/thread/249}{Erd\H{o}s Problems
discussion thread}, accessed 16 September 2026.
\bibitem{erdosgraham1980}
P.~Erd\H{o}s and R.~L.~Graham, \emph{Old and New Problems and Results in
Combinatorial Number Theory}, Monographies de L'Enseignement
Math\'ematique \textbf{28}, 1980, p.~61.
\bibitem{erdos1948lambert}
P.~Erd\H{o}s, \emph{On arithmetical properties of Lambert series},
J.\ Indian Math.\ Soc.\ (N.S.) \textbf{12} (1948), 63--66,
\url{https://users.renyi.hu/~p_erdos/1948-04.pdf}.
\bibitem{postelmans-vanassche}
K.~Postelmans and W.~Van Assche, \emph{Irrationality of $\zeta_q(1)$ and
$\zeta_q(2)$}, J.\ Number Theory \textbf{126} (2007), no.~1, 119--154,
doi:\href{https://doi.org/10.1016/j.jnt.2006.11.011}{10.1016/j.jnt.2006.11.011};
arXiv:\href{https://arxiv.org/abs/math/0604312v1}{math/0604312v1}.
Theorem~1.3 is cited with the arXiv v1 pagination, p.~3.
\bibitem{nesterenko1996}
Yu.~V.~Nesterenko, \emph{Modular functions and transcendence questions},
Sb.\ Math.\ \textbf{187} (1996), no.~9, 1319--1348,
doi:\href{https://doi.org/10.1070/SM1996v187n09ABEH000158}{10.1070/SM1996v187n09ABEH000158}.
Corollary~2, p.~1320.
\bibitem{schramm-gcd}
W.~Schramm, \emph{The Fourier transform of functions of the greatest common
divisor}, Integers \textbf{8} (2008), \#A50,
\url{https://math.colgate.edu/~integers/i50/i50.pdf}.
Theorem and equation~(2), p.~2.
\bibitem{toth-gcd}
L.~T\'oth, \emph{A survey of gcd-sum functions}, J.\ Integer Seq.\ \textbf{13}
(2010), Article~10.8.1,
\url{https://cs.uwaterloo.ca/journals/JIS/VOL13/Toth/toth10.pdf}.
\S1, equations (1)--(2).
\bibitem{merca-schmidt}
M.~Merca and M.~D.~Schmidt, \emph{Generating special arithmetic functions by
Lambert series factorizations}, Contrib.\ Discrete Math.\ \textbf{14} (2019),
no.~1, 31--45,
doi:\href{https://doi.org/10.55016/ojs/cdm.v14i1.62425}{10.55016/ojs/cdm.v14i1.62425};
arXiv:\href{https://arxiv.org/abs/1706.00393v2}{1706.00393v2}.
Equation~(1) is cited with the arXiv v2 pagination, p.~2.
\bibitem{merca2017}
M.~Merca, \emph{The Lambert series factorization theorem},
Ramanujan J.\ \textbf{44} (2017), 417--435,
doi:\href{https://doi.org/10.1007/s11139-016-9856-3}{10.1007/s11139-016-9856-3}.
Theorem~1.2, p.~420.
\bibitem{kovac-tao}
V.~Kova\v{c} and T.~Tao, \emph{On several irrationality problems for Ahmes
series}, Acta Math.\ Hungar.\ \textbf{175} (2025), no.~2, 572--608,
doi:\href{https://doi.org/10.1007/s10474-025-01528-0}{10.1007/s10474-025-01528-0};
arXiv:\href{https://arxiv.org/abs/2406.17593v4}{2406.17593v4}.
Theorem~2.3 is cited with the arXiv v4 pagination, p.~5.
\bibitem{bgs2016}
R.~Balasubramanian, S.~Giri and P.~Srivastav, \emph{On correlations of certain
multiplicative functions}, J.\ Number Theory \textbf{174} (2017), 221--238,
doi:\href{https://doi.org/10.1016/j.jnt.2016.10.001}{10.1016/j.jnt.2016.10.001};
arXiv:\href{https://arxiv.org/abs/1511.02221v3}{1511.02221v3}.
Theorem~2.2 is cited with the arXiv v3 pagination, p.~2.
\bibitem{granville-smooth}
A.~Granville, \emph{Smooth numbers: computational number theory and beyond},
in \emph{Algorithmic Number Theory}, MSRI Publ.\ \textbf{44} (2008), 267--323,
\url{https://library.slmath.org/books/Book44/files/09andrew.pdf}.
Equations (1.1) and (1.3), p.~268.
\bibitem{schoenberg1928}
I.~Schoenberg, \emph{\"Uber die asymptotische Verteilung reeller Zahlen
mod~1}, Math.\ Z.\ \textbf{28} (1928), 171--199,
doi:\href{https://doi.org/10.1007/BF01181156}{10.1007/BF01181156}.
\S17, p.~193, for the distribution of $\varphi(n)/n$.
\bibitem{schoenberg1936}
I.~J.~Schoenberg, \emph{On asymptotic distributions of arithmetical
functions}, Trans.\ Amer.\ Math.\ Soc.\ \textbf{39} (1936), no.~2, 315--330.
Theorem~1, pp.~318--319; \S8, p.~323.
\bibitem{edeko-vdc}
N.~Edeko, H.~Kreidler and R.~Nagel, \emph{A dynamical proof of the van der
Corput inequality}, Dyn.\ Syst.\ \textbf{37} (2022), no.~4, 648--665,
doi:\href{https://doi.org/10.1080/14689367.2022.2100244}{10.1080/14689367.2022.2100244};
arXiv:\href{https://arxiv.org/abs/2106.11835v3}{2106.11835v3}.
Theorem~2.1 is cited with the arXiv v3 pagination, p.~4.
\bibitem{matomaki-radziwill}
K.~Matom\"aki and M.~Radziwi\l\l, \emph{Multiplicative functions in short
intervals}, Ann.\ of Math.\ \textbf{183} (2016), no.~3, 1015--1056,
doi:\href{https://doi.org/10.4007/annals.2016.183.3.6}{10.4007/annals.2016.183.3.6};
arXiv:\href{https://arxiv.org/abs/1501.04585v4}{1501.04585v4}.
Theorem~1 is cited with the arXiv v4 pagination, pp.~1--2.
\bibitem{tao-elliott}
T.~Tao, \emph{The logarithmically averaged Chowla and Elliott conjectures for
two-point correlations}, Forum Math.\ Pi \textbf{4} (2016), e8,
doi:\href{https://doi.org/10.1017/fmp.2016.6}{10.1017/fmp.2016.6};
arXiv:\href{https://arxiv.org/abs/1509.05422v4}{1509.05422v4}.
Theorems~1.2--1.3 are cited with the arXiv v4 pagination, pp.~2 and~5.
\bibitem{tao-teravainen}
T.~Tao and J.~Ter\"av\"ainen, \emph{Odd order cases of the logarithmically
averaged Chowla conjecture}, J.\ Th\'eor.\ Nombres Bordeaux \textbf{30}
(2018), no.~3, 997--1015,
doi:\href{https://doi.org/10.5802/jtnb.1062}{10.5802/jtnb.1062};
arXiv:\href{https://arxiv.org/abs/1710.02112v1}{1710.02112v1}.
Theorem~1.1 is cited with the arXiv v1 pagination, p.~2.
\bibitem{adamczewski-bugeaud}
B.~Adamczewski and Y.~Bugeaud, \emph{On the complexity of algebraic numbers
I. Expansions in integer bases}, Ann.\ of Math.\ \textbf{165} (2007), no.~2,
547--565. Theorem~1, p.~549.
\bibitem{everest-periodic}
G.~Everest, A.~J.~van der Poorten, Y.~Puri and T.~Ward, \emph{Integer
sequences and periodic points}, J.\ Integer Seq.\ \textbf{5} (2002),
Article~02.2.3. Lemma~3.1 is cited with the preprint pagination, p.~5.
\bibitem{erdos-gal}
P.~Erd\H{o}s and I.~S.~G\'al, \emph{On the law of the iterated logarithm.~I},
Proc.\ Kon.\ Ned.\ Akad.\ Wetensch.\ Ser.~A \textbf{58} $=$ Indag.\ Math.\
\textbf{17} (1955), 65--76, \url{https://www.renyi.hu/~p_erdos/1955-06.pdf}.
\bibitem{yazdani2001} S.~Yazdani,
\emph{Multiplicative functions and $k$-automatic sequences},
J.\ Th\'eor.\ Nombres Bordeaux \textbf{13} (2001), no.~2, 651--658,
\url{https://www.numdam.org/item/JTNB_2001__13_2_651_0/}.
Theorem~2 and its proof, pp.~652--653; Corollary~4, p.~654.
\bibitem{allouche-shallit-yassawi} J.-P.~Allouche, J.~Shallit and R.~Yassawi,
\emph{How to prove that a sequence is not automatic},
Exposition.\ Math.\ \textbf{40} (2022), no.~1, 1--22,
doi:\href{https://doi.org/10.1016/j.exmath.2021.08.001}{10.1016/j.exmath.2021.08.001};
arXiv:\href{https://arxiv.org/abs/2104.13072v1}{2104.13072v1}.
Theorem and example numbering follows the arXiv v1 (2021).
\bibitem{bell-smertnig2026} J.~Bell and D.~Smertnig,
\emph{Mahler series with multiplicative coefficient sequences},
arXiv:\href{https://arxiv.org/abs/2603.23456v1}{2603.23456v1},
24 March 2026, preprint. Theorem~1.3, p.~2; Lemma~3.3, p.~9.
\bibitem{bell-bruin-coons} J.~P.~Bell, N.~Bruin and M.~Coons,
\emph{Transcendence of generating functions whose coefficients are multiplicative},
Trans.\ Amer.\ Math.\ Soc.\ \textbf{364} (2012), no.~2, 933--959,
doi:\href{https://doi.org/10.1090/S0002-9947-2011-05479-6}{10.1090/S0002-9947-2011-05479-6};
arXiv:\href{https://arxiv.org/abs/1003.2221v2}{1003.2221v2}.
Theorems~1.5--1.6 are cited with the arXiv v2 pagination, p.~2.
\bibitem{adamczewski-drmota-mullner} B.~Adamczewski, M.~Drmota and C.~M\"ullner,
\emph{(Logarithmic) densities for automatic sequences along primes and squares},
Trans.\ Amer.\ Math.\ Soc.\ \textbf{375} (2022), no.~1, 455--499,
doi:\href{https://doi.org/10.1090/tran/8476}{10.1090/tran/8476};
arXiv:\href{https://arxiv.org/abs/2009.14773v2}{2009.14773v2}.
Theorem~1.4 and Remark~1.5, p.~4, and the proof at the end of Section~7,
p.~19, refer to the arXiv v2 (2021).
\bibitem{adamczewski-bell-smertnig2023}
B.~Adamczewski, J.~Bell and D.~Smertnig,
\emph{A height gap theorem for coefficients of Mahler functions},
J.~Eur.~Math.~Soc.\ \textbf{25} (2023), no.~7, 2525--2571,
doi:\href{https://doi.org/10.4171/JEMS/1244}{10.4171/JEMS/1244}.
Theorem~1.2(b), p.~2528; logarithmic height in Section~1.1.1, p.~2527.
\end{thebibliography}

% ---- part a249_p0 ----
% ---- part a249_p1a ----
% ---- part a249_p1b ----
% ---- part a249_newdecls ----
% ---- part a249_p2 ----
% ---- part a249_p3 ----
% ---- part a249_p4 ----
% ---- part a249_invent ----
% ---- part a249_p5 ----
\end{document}
